% Auto-generated from data/segments.json + layout.json (+ transforms) — do not edit by hand.
% volume: 29Abhi01
% mode: printing
% segments: 2408
% transform_rules: 8
% toc_marks: 266

% Volume layout defaults (layout.json ``layout``).
\csromanlayoutapply{1}{1.5}{21.6pt}{8pt}{8pt}{65pt}{2.5em}%

\pitaka{อภิธมฺม\-ปิฏก}
\csromanmatikaanchor{mtk.0}
\csromantocmarknopage{chapter}{ธมฺมสงฺคณีปาฬิ}{mtk.0}
\bookmark[dest={mtk.0},level=0]{ธมฺมสงฺคณีปาฬิ}
\gambhira{ธมฺมสงฺคณี\-ปาฬิ}
\namakkaram{นโม ตสฺส ภควโต อรหโต สมฺมา\-สมฺพุทฺธสฺส.}
\csromanmatikaanchor{mtk.1}
\csromantocmarkref{section}{ติกมาติกา}{mtk.1}
\bookmark[dest={mtk.1},level=1]{ติกมาติกา}
\csromanheader{1}{\textbf{ติกมาติกา}}
\proseitem{1}{\csromanfolio{1}กุสลา ธมฺมา. (363, 985, 1384)}

\prose{อกุสลา ธมฺมา. (365, 427, 986, 1385)}

\prose{อพฺยากตา ธมฺมา. (431, 583, 987, 1386)}

\proseitem{2}{สุขาย เวทนาย สมฺปยุตฺตา ธมฺมา. \mbox{(988, 1387)}}

\prose{ทุกฺขาย เวทนาย สมฺปยุตฺตา ธมฺมา. \mbox{(989, 1388)}}

\prose{อทุกฺขม\-สุขาย เวทนาย สมฺปยุตฺตา ธมฺมา. \mbox{(990, 1389)}}

\proseitem{3}{วิปากา ธมฺมา. \mbox{(991, 1390)}}

\prose{วิปาก\-ธมฺม\-ธมฺมา. \mbox{(992, 1391)}}

\prose{เนว\-วิปาก\-น\-วิปากธมฺม\-ธมฺมา. \mbox{(993, 1392)}}

\proseitem{4}{อุปาทิณฺณุ\-ปาทานิยา\csromansharedfootnote{fe49cad4036be5d3}{อุปาทินฺนุปาทานิยา (สฺยา)} ธมฺมา. \mbox{(994, 1393)}}

\prose{อนุปาทิณฺณุ\-ปาทานิยา ธมฺมา. \mbox{(995, 1394)}}

\prose{อนุปาทิณฺณ\-อนุปาทานิยา\csromansharedfootnote{8e79a70ebae08ad1}{อนุปาทินฺนานุปาทานิยา (สฺยา)} ธมฺมา. \mbox{(996, 1395)}}

\proseitem{5}{สํกิลิฏฺฐ\-สํกิเลสิกา ธมฺมา. \mbox{(997, 1396)}}

\prose{อสํกิลิฏฺฐ\-สํกิเลสิกา ธมฺมา. \mbox{(998, 1397)}}

\prose{อสํกิลิฏฺฐ\-อสํกิเลสิกา\csromansharedfootnote{a5c63ad2b9adbeee}{อสํกิลิฏฺฐาสํกิเลสิกา (สฺยา)} ธมฺมา. \mbox{(999, 1398)}}

\csromanhangingbegin
\hangingheaditem{6}{\csromanfolio{2}สวิตกฺก\-สวิจารา ธมฺมา. \mbox{(1000, 1399)}}
\hangingbody{อวิตกฺกวิจารมตฺตา ธมฺมา. \mbox{(1001, 1400)}}
\hangingbody{อวิตกฺกอวิจารา1 ธมฺมา. \mbox{(1002, 1401)}}
\csromanhangingend

\csromanhangingbegin
\hangingheaditem{7}{ปีติสหคตา ธมฺมา. \mbox{(1003, 1402)}}
\hangingbody{สุขสหคตา ธมฺมา. \mbox{(1004, 1403)}}
\hangingbody{อุเปกฺขาสหคตา ธมฺมา. \mbox{(1005, 1404)}}
\csromanhangingend

\csromanhangingbegin
\hangingheaditem{8}{ทสฺสเนน ปหาตพฺพา ธมฺมา. \mbox{(1006, 1405)}}
\hangingbody{ภาวนาย ปหาตพฺพา ธมฺมา. \mbox{(1011, 1406)}}
\hangingbody{เนว ทสฺสเนน น ภาวนาย ปหาตพฺพา ธมฺมา. \mbox{(1012, 1407)}}
\csromanhangingend

\csromanhangingbegin
\hangingheaditem{9}{ทสฺสเนน ปหาตพฺพ\-เหตุกา ธมฺมา. \mbox{(1013, 1408)}}
\hangingbody{ภาวนาย ปหาตพฺพเหตุกา ธมฺมา. \mbox{(1018, 1409)}}
\hangingbody{เนว ทสฺสเนน น ภาวนาย ปหาตพฺพเหตุกา ธมฺมา.}
\hangingbody{\mbox{(1019, 1410)}}
\csromanhangingend

\csromanhangingbegin
\hangingheaditem{10}{อาจยคามิโน ธมฺมา. \mbox{(1020, 1411)}}
\hangingbody{อปจยคามิโน ธมฺมา. \mbox{(1021, 1412)}}
\hangingbody{เนวาจยคามิ นาปจยคามิโน2 ธมฺมา. \mbox{(1022, 1413)}}
\csromanhangingend

\csromanhangingbegin
\hangingheaditem{11}{เสกฺขา ธมฺมา. \mbox{(1023, 1414)}}
\hangingbody{อเสกฺขา ธมฺมา. \mbox{(1024, 1415)}}
\hangingbody{เนวเสกฺขนาเสกฺขา3 ธมฺมา. \mbox{(1025, 1416)}}
\csromanhangingend

\csromanhangingbegin
\hangingheaditem{12}{ปริตฺตา ธมฺมา. \mbox{(1026, 1417)}}
\hangingbody{มหคฺคตา ธมฺมา. \mbox{(1027, 1418)}}
\hangingbody{อปฺปมาณา ธมฺมา. \mbox{(1028, 1419)}}
\csromanhangingend

\csromanhangingbegin
\hangingheaditem{13}{ปริตฺตารมฺมณา ธมฺมา. \mbox{(1029, 1420)}}
\hangingbody{มหคฺคตารมฺมณา ธมฺมา. \mbox{(1030, 1421)}}
\hangingbody{อปฺปมาณารมฺมณา ธมฺมา. \mbox{(1031, 1422)}}
\csromanhangingend

\csromanhangingbegin
\hangingheaditem{14}{\csromanfolio{3}หีนา ธมฺมา. \mbox{(1032, 1423)}}
\hangingbody{มชฺฌิมา ธมฺมา. \mbox{(1033, 1424)}}
\hangingbody{ปณีตา ธมฺมา. \mbox{(1034, 1425)}}
\csromanhangingend

\csromanhangingbegin
\hangingheaditem{15}{มิจฺฉตฺต\-นิยตา ธมฺมา. \mbox{(1035, 1426)}}
\hangingbody{สมฺมตฺตนิยตา ธมฺมา. \mbox{(1036, 1427)}}
\hangingbody{อนิยตา ธมฺมา. \mbox{(1037, 1428)}}
\csromanhangingend

\csromanhangingbegin
\hangingheaditem{16}{มคฺคารมฺมณา ธมฺมา. \mbox{(1038, 1429)}}
\hangingbody{มคฺคเหตุกา ธมฺมา. \mbox{(1039, 1429)}}
\hangingbody{มคฺคาธิปติโน ธมฺมา. \mbox{(1040, 1429)}}
\csromanhangingend

\csromanhangingbegin
\hangingheaditem{17}{อุปฺปนฺนา ธมฺมา. \mbox{(1041, 1430)}}
\hangingbody{อนุปฺปนฺนา ธมฺมา. \mbox{(1042, 1430)}}
\hangingbody{อุปฺปาทิโน ธมฺมา. \mbox{(1043, 1430)}}
\csromanhangingend

\csromanhangingbegin
\hangingheaditem{18}{อตีตา ธมฺมา. \mbox{(1044, 1431)}}
\hangingbody{อนาคตา ธมฺมา. \mbox{(1045, 1431)}}
\hangingbody{ปจฺจุปฺปนฺนา ธมฺมา. \mbox{(1046, 1431)}}
\csromanhangingend

\csromanhangingbegin
\hangingheaditem{19}{อตีตารมฺมณา ธมฺมา. \mbox{(1047, 1432)}}
\hangingbody{อนาคตารมฺมณา ธมฺมา. \mbox{(1048, 1433)}}
\hangingbody{ปจฺจุปฺปนฺนารมฺมณา ธมฺมา. \mbox{(1049, 1434)}}
\csromanhangingend

\csromanhangingbegin
\hangingheaditem{20}{อชฺฌตฺตา ธมฺมา. \mbox{(1050, 1435)}}
\hangingbody{พหิทฺธา ธมฺมา. \mbox{(1051, 1435)}}
\hangingbody{อชฺฌตฺตพหิทฺธา ธมฺมา. \mbox{(1052, 1435)}}
\csromanhangingend

\csromanhangingbegin
\hangingheaditem{21}{อชฺฌตฺตา\-รมฺมณา ธมฺมา. \mbox{(1053, 1436)}}
\hangingbody{พหิทฺธารมฺมณา ธมฺมา. \mbox{(1054, 1437)}}
\hangingbody{อชฺฌตฺตพหิทฺธารมฺมณา ธมฺมา. \mbox{(1055, 1437)}}
\csromanhangingend

\csromanhangingbegin
\hangingheaditem{22}{สนิทสฺสน\-สปฺปฏิฆา ธมฺมา. \mbox{(1056, 1438)}}
\hangingbody{อนิทสฺสนสปฺปฏิฆา ธมฺมา. \mbox{(1057, 1439)}}
\hangingbody{อนิทสฺสนอปฺปฏิฆา1 ธมฺมา. \mbox{(1058, 1440)}}
\csromanhangingend

\csromanheader{2}{ติกมาติกา.}
\csromanmatikaanchor{mtk.2}
\csromantocmarkref{section}{ทุกมาติกา}{mtk.2}
\bookmark[dest={mtk.2},level=1]{ทุกมาติกา}
\csromanheader{1}{\textbf{ทุกมาติกา}}
\csromanmatikaanchor{mtk.3}
\csromantocmarkref{section}{เหตุโคจฺฉก}{mtk.3}
\bookmark[dest={mtk.3},level=1]{เหตุโคจฺฉก}
\csromanheader{1}{\textbf{เหตุโคจฺฉก}}
\csromanhangingbegin
\hangingheaditem{1}{\csromanfolio{4}เหตู ธมฺมา. (1059, 1077,1441)}
\hangingbody{น เหตู ธมฺมา. \mbox{(1078, 1442)}}
\csromanhangingend

\csromanhangingbegin
\hangingheaditem{2}{สเหตุกา ธมฺมา. \mbox{(1079, 1443)}}
\hangingbody{อเหตุกา ธมฺมา. \mbox{(1080, 1444)}}
\csromanhangingend

\csromanhangingbegin
\hangingheaditem{3}{เหตุสมฺปยุตฺตา ธมฺมา. \mbox{(1081, 1445)}}
\hangingbody{เหตุวิปฺปยุตฺตา ธมฺมา. \mbox{(1082, 1446)}}
\csromanhangingend

\csromanhangingbegin
\hangingheaditem{4}{เหตู เจว ธมฺมา สเหตุกา จ. \mbox{(1083, 1447)}}
\hangingbody{สเหตุกา เจว ธมฺมา น จ เหตู. \mbox{(1084, 1448)}}
\csromanhangingend

\csromanhangingbegin
\hangingheaditem{5}{เหตู เจว ธมฺมา เหตุสมฺปยุตฺตา จ. \mbox{(1085, 1449)}}
\hangingbody{เหตุสมฺปยุตฺตา เจว ธมฺมา น จ เหตู. \mbox{(1086, 1450)}}
\csromanhangingend

\csromanhangingbegin
\hangingheaditem{6}{น เหตู โข ปน ธมฺมา สเหตุกาปิ. \mbox{(1087, 1451)}}
\hangingbody{อเหตุกาปิ. \mbox{(1088, 1452)}}
\csromanhangingend

\vspace{\tipitakaparskip}
\csromancenter{\textbf{เหตุโคจฺฉก}ํ.}
\csromanmatikaanchor{mtk.4}
\csromantocmarkref{section}{จูฬนฺตรทุก}{mtk.4}
\bookmark[dest={mtk.4},level=1]{จูฬนฺตรทุก}
\csromanheader{1}{\textbf{จูฬนฺตรทุก}}
\csromanhangingbegin
\hangingheaditem{7}{สปฺปจฺจยา ธมฺมา. \mbox{(1089, 1453)}}
\hangingbody{อปฺปจฺจยา ธมฺมา. \mbox{(1090, 1454)}}
\csromanhangingend

\csromanhangingbegin
\hangingheaditem{8}{สงฺขตา ธมฺมา. \mbox{(1091, 1455)}}
\hangingbody{อสงฺขตา ธมฺมา. \mbox{(1092, 1456)}}
\csromanhangingend

\csromanhangingbegin
\hangingheaditem{9}{สนิทสฺสนา ธมฺมา. \mbox{(1093, 1457)}}
\hangingbody{อนิทสฺสนา ธมฺมา. \mbox{(1094, 1458)}}
\csromanhangingend

\csromanhangingbegin
\hangingheaditem{10}{สปฺปฏิฆา ธมฺมา. \mbox{(1095, 1459)}}
\hangingbody{อปฺปฏิฆา ธมฺมา. \mbox{(1096, 1460)}}
\csromanhangingend

\csromanhangingbegin
\hangingheaditem{11}{\csromanfolio{5}รูปิโน ธมฺมา. \mbox{(1097, 1461)}}
\hangingbody{อรูปิโน ธมฺมา. \mbox{(1098, 1462)}}
\csromanhangingend

\csromanhangingbegin
\hangingheaditem{12}{โลกิยา ธมฺมา. \mbox{(1099, 1463)}}
\hangingbody{โลกุตฺตรา ธมฺมา. \mbox{(1100, 1464)}}
\csromanhangingend

\proseitem{13}{เกนจิ วิญฺเญยฺยา ธมฺมา. \mbox{(1101, 1464)}}

\setlength{\csromangathapagewidth}{0pt}
\setlength{\csromangathawakwidth}{0pt}
\csromangathameasuregroup{}{เกนจิ น วิญฺเญยฺยา ธมฺมา. \\ (1101, \\ 1464) จูฬนฺตรทุกํ.}
\csromangathameasurewak{เกนจิ น วิญฺเญยฺยา ธมฺมา. \\ (1101, \\ 1464) จูฬนฺตรทุกํ.}
\setlength{\csromangathaleftcol}{0pt}
\csromangathagroup{\csromangathastanza{\csromangathawak{เกนจิ น วิญฺเญยฺยา ธมฺมา. \\ (1101, \\ 1464) จูฬนฺตรทุกํ.}}}

\csromanmatikaanchor{mtk.5}
\csromantocmarkref{section}{อาสวโคจฺฉก}{mtk.5}
\bookmark[dest={mtk.5},level=1]{อาสวโคจฺฉก}
\csromanheader{1}{\textbf{อาสวโคจฺฉก}}
\csromanhangingbegin
\hangingheaditem{14}{อาสวา ธมฺมา. \mbox{(1102, 1465)}}
\hangingbody{โน อาสวา ธมฺมา. \mbox{(1107, 1466)}}
\csromanhangingend

\csromanhangingbegin
\hangingheaditem{15}{สาสวา ธมฺมา. \mbox{(1108, 1467)}}
\hangingbody{อนาสวา ธมฺมา. \mbox{(1109, 1468)}}
\csromanhangingend

\csromanhangingbegin
\hangingheaditem{16}{อาสว\-สมฺปยุตฺตา ธมฺมา. \mbox{(1110, 1469)}}
\hangingbody{อาสววิปฺปยุตฺตา ธมฺมา. \mbox{(1111, 1470)}}
\csromanhangingend

\csromanhangingbegin
\hangingheaditem{17}{อาสวา เจว ธมฺมา สาสวา จ. \mbox{(1112, 1471)}}
\hangingbody{สาสวา เจว ธมฺมา โน จ อาสวา. \mbox{(1113, 1472)}}
\csromanhangingend

\csromanhangingbegin
\hangingheaditem{18}{อาสวา เจว ธมฺมา อาสว\-สมฺปยุตฺตา จ. \mbox{(1114, 1473)}}
\hangingbody{อาสวสมฺปยุตฺตา เจว ธมฺมา โน จ อาสวา. \mbox{(1115, 1474)}}
\csromanhangingend

\proseitem{19}{อาสว\-วิปฺปยุตฺตา โข ปน ธมฺมา สาสวาปิ. \mbox{(1116, 1475)}}

\setlength{\csromangathapagewidth}{0pt}
\setlength{\csromangathawakwidth}{0pt}
\csromangathameasuregroup{}{อนาสวาปิ. \\ (1117, \\ 1476) อาสวโคจฺฉกํ.}
\csromangathameasurewak{อนาสวาปิ. \\ (1117, \\ 1476) อาสวโคจฺฉกํ.}
\setlength{\csromangathaleftcol}{0pt}
\csromangathagroup{\csromangathastanza{\csromangathawak{อนาสวาปิ. \\ (1117, \\ 1476) อาสวโคจฺฉกํ.}}}

\csromanmatikaanchor{mtk.6}
\csromantocmarkref{section}{สญฺโญชนโคจฺฉก}{mtk.6}
\bookmark[dest={mtk.6},level=1]{สญฺโญชนโคจฺฉก}
\csromanheader{1}{สญฺโญชน\-โคจฺฉก}
\csromanhangingbegin
\hangingheaditem{20}{สญฺโญชนา ธมฺมา. \mbox{(1118, 1477)}}
\hangingbody{โน สญฺโญชนา ธมฺมา. \mbox{(1129, 1478)}}
\csromanhangingend

\csromanhangingbegin
\hangingheaditem{21}{\csromanfolio{6}สญฺโญชนิยา ธมฺมา. \mbox{(1130, 1479)}}
\hangingbody{อสญฺโญชนิยา ธมฺมา. \mbox{(1131, 1480)}}
\csromanhangingend

\csromanhangingbegin
\hangingheaditem{22}{สญฺโญชน\-สมฺปยุตฺตา ธมฺมา. \mbox{(1132, 1481)}}
\hangingbody{สญฺโญชนวิปฺปยุตฺตา ธมฺมา. \mbox{(1133, 1482)}}
\csromanhangingend

\csromanhangingbegin
\hangingheaditem{23}{สญฺโญชนา เจว ธมฺมา สญฺโญชนิยา จ. \mbox{(1134, 1483)}}
\hangingbody{สญฺโญชนิยา เจว ธมฺมา โน จ สญฺโญชนา. \mbox{(1135, 1484)}}
\csromanhangingend

\csromanhangingbegin
\hangingheaditem{24}{สญฺโญชนา เจว ธมฺมา สญฺโญชน\-สมฺปยุตฺตา จ. \mbox{(1136, 1485)}}
\hangingbody{สญฺโญชนสมฺปยุตฺตา เจว ธมฺมา โน จ สญฺโญชนา. \mbox{(1137, 1486)}}
\csromanhangingend

\csromanhangingbegin
\hangingheaditem{25}{สญฺโญชน\-วิปฺปยุตฺตา โข ปน ธมฺมา สญฺโญชนิยาปิ. \mbox{(1138, 1487)}}
\hangingbody{อสญฺโญชนิยาปิ. \mbox{(1139, 1488)}}
\csromanhangingend

\vspace{\tipitakaparskip}
\csromancenter{สญฺโญชน\-โคจฺฉกํ.}
\csromanmatikaanchor{mtk.7}
\csromantocmarkref{section}{คนฺถโคจฺฉก}{mtk.7}
\bookmark[dest={mtk.7},level=1]{คนฺถโคจฺฉก}
\csromanheader{1}{\textbf{คนฺถโคจฺฉก}}
\csromanhangingbegin
\hangingheaditem{26}{คนฺถา ธมฺมา. \mbox{(1140, 1489)}}
\hangingbody{โน คนฺถา ธมฺมา. \mbox{(1145, 1490)}}
\csromanhangingend

\csromanhangingbegin
\hangingheaditem{27}{คนฺถนิยา ธมฺมา. \mbox{(1146, 1491)}}
\hangingbody{อคนฺถนิยา ธมฺมา. \mbox{(1147, 1492)}}
\csromanhangingend

\csromanhangingbegin
\hangingheaditem{28}{คนฺถ\-สมฺปยุตฺตา ธมฺมา. \mbox{(1148, 1493)}}
\hangingbody{คนฺถวิปฺปยุตฺตา ธมฺมา. \mbox{(1149, 1494)}}
\csromanhangingend

\csromanhangingbegin
\hangingheaditem{29}{คนฺถา เจว ธมฺมา คนฺถนิยา จ. \mbox{(1150, 1495)}}
\hangingbody{คนฺถนิยา เจว ธมฺมา โน จ คนฺถา. \mbox{(1151, 1496)}}
\csromanhangingend

\csromanhangingbegin
\hangingheaditem{30}{คนฺถา เจว ธมฺมา คนฺถ\-สมฺปยุตฺตา จ. \mbox{(1152, 1497)}}
\hangingbody{คนฺถสมฺปยุตฺตา เจว ธมฺมา โน จ คนฺถา. \mbox{(1153, 1498)}}
\csromanhangingend

\proseitem{31}{\csromanfolio{7}คนฺถ\-วิปฺปยุตฺตา โข ปน ธมฺมา คนฺถนิยาปิ. \mbox{(1154, 1499)}}

\setlength{\csromangathapagewidth}{0pt}
\setlength{\csromangathawakwidth}{0pt}
\csromangathameasuregroup{}{อคนฺถนิยาปิ. \\ (1155, \\ 1500) คนฺถโคจฺฉกํ.}
\csromangathameasurewak{อคนฺถนิยาปิ. \\ (1155, \\ 1500) คนฺถโคจฺฉกํ.}
\setlength{\csromangathaleftcol}{0pt}
\csromangathagroup{\csromangathastanza{\csromangathawak{อคนฺถนิยาปิ. \\ (1155, \\ 1500) คนฺถ\-โคจฺฉกํ.}}}

\csromanmatikaanchor{mtk.8}
\csromantocmarkref{section}{โอฆโคจฺฉก}{mtk.8}
\bookmark[dest={mtk.8},level=1]{โอฆโคจฺฉก}
\csromanheader{1}{\textbf{โอฆโคจฺฉก}}
\csromanhangingbegin
\hangingheaditem{32}{โอฆา ธมฺมา. \mbox{(1156, 1501)}}
\hangingbody{โน โอฆา ธมฺมา.}
\csromanhangingend

\setlength{\csromangathapagewidth}{0pt}
\setlength{\csromangathawakwidth}{0pt}
\csromangathameasuregroup{}{โอฆนิยา ธมฺมา. \\ อโนฆนิยา ธมฺมา. \\ โอฆสมฺปยุตฺตา ธมฺมา. \\ โอฆวิปฺปยุตฺตา ธมฺมา. \\ โอฆา เจว ธมฺมา โอฆนิยา จ. \\ โอฆนิยา เจว ธมฺมา โน จ โอฆา. \\ โอฆา เจว ธมฺมา โอฆสมฺปยุตฺตา จ. \\ โอฆสมฺปยุตฺตา เจว ธมฺมา โน จ โอฆา.}
\csromangathameasurewak{โอฆนิยา ธมฺมา. \\ อโนฆนิยา ธมฺมา. \\ โอฆสมฺปยุตฺตา ธมฺมา. \\ โอฆวิปฺปยุตฺตา ธมฺมา. \\ โอฆา เจว ธมฺมา โอฆนิยา จ. \\ โอฆนิยา เจว ธมฺมา โน จ โอฆา. \\ โอฆา เจว ธมฺมา โอฆสมฺปยุตฺตา จ. \\ โอฆสมฺปยุตฺตา เจว ธมฺมา โน จ โอฆา.}
\setlength{\csromangathaleftcol}{0pt}
\csromangathagroup{\csromangathastanza{\csromangathawak{\csromangathaitemline{33}{โอฆนิยา ธมฺมา.} \\ \csromangathacontline{อโนฆนิยา ธมฺมา.}}}\csromangathastanza{\csromangathawak{\csromangathaitemline{34}{โอฆสมฺปยุตฺตา ธมฺมา.} \\ \csromangathacontline{โอฆวิปฺปยุตฺตา ธมฺมา.}}}\csromangathastanza{\csromangathawak{\csromangathaitemline{35}{โอฆา เจว ธมฺมา โอฆนิยา จ.} \\ \csromangathacontline{โอฆนิยา เจว ธมฺมา โน จ โอฆา.}}}\csromangathastanza{\csromangathawak{\csromangathaitemline{36}{โอฆา เจว ธมฺมา โอฆสมฺปยุตฺตา จ.} \\ \csromangathacontline{โอฆสมฺปยุตฺตา เจว ธมฺมา โน จ โอฆา.}}}}

\proseitem{37}{โอฆวิปฺปยุตฺตา โข ปน ธมฺมา โอฆนิยาปิ.}

\prose{อโนฆนิยาปิ. โอฆโคจฺฉกํ.}

\csromanmatikaanchor{mtk.9}
\csromantocmarkref{section}{โยคโคจฺฉก}{mtk.9}
\bookmark[dest={mtk.9},level=1]{โยคโคจฺฉก}
\csromanheader{1}{\textbf{โยคโคจฺฉก}}
\csromanhangingbegin
\hangingheaditem{38}{โยคา ธมฺมา. \mbox{(1157, 1502)}}
\hangingbody{โน โยคา ธมฺมา.}
\csromanhangingend

\setlength{\csromangathapagewidth}{0pt}
\setlength{\csromangathawakwidth}{0pt}
\csromangathameasuregroup{}{โยคนิยา ธมฺมา. \\ อโยคนิยา ธมฺมา. \\ โยคสมฺปยุตฺตา ธมฺมา. \\ โยควิปฺปยุตฺตา ธมฺมา. \\ โยคา เจว ธมฺมา โยคนิยา จ. \\ โยคนิยา เจว ธมฺมา โน จ โยคา. \\ โยคา เจว ธมฺมา โยคสมฺปยุตฺตา จ. \\ โยคสมฺปยุตฺตา เจว ธมฺมา โน จ โยคา. \\ โยควิปฺปยุตฺตา โข ปน ธมฺมา โยคนิยาปิ. \\ อโยคนิยาปิ.}
\csromangathameasurewak{โยคนิยา ธมฺมา. \\ อโยคนิยา ธมฺมา. \\ โยคสมฺปยุตฺตา ธมฺมา. \\ โยควิปฺปยุตฺตา ธมฺมา. \\ โยคา เจว ธมฺมา โยคนิยา จ. \\ โยคนิยา เจว ธมฺมา โน จ โยคา. \\ โยคา เจว ธมฺมา โยคสมฺปยุตฺตา จ. \\ โยคสมฺปยุตฺตา เจว ธมฺมา โน จ โยคา. \\ โยควิปฺปยุตฺตา โข ปน ธมฺมา โยคนิยาปิ. \\ อโยคนิยาปิ.}
\setlength{\csromangathaleftcol}{0pt}
\csromangathagroup{\csromangathastanza{\csromangathawak{\csromangathaitemline{39}{โยคนิยา ธมฺมา.} \\ \csromangathacontline{อโยคนิยา ธมฺมา.}}}\csromangathastanza{\csromangathawak{\csromangathaitemline{40}{โยคสมฺปยุตฺตา ธมฺมา.} \\ \csromangathacontline{โยควิปฺปยุตฺตา ธมฺมา.}}}\csromangathastanza{\csromangathawak{\csromanfolio{8}\csromangathaitemline{41}{โยคา เจว ธมฺมา โยคนิยา จ.} \\ \csromangathacontline{โยคนิยา เจว ธมฺมา โน จ โยคา.}}}\csromangathastanza{\csromangathawak{\csromangathaitemline{42}{โยคา เจว ธมฺมา โยคสมฺปยุตฺตา จ.} \\ \csromangathacontline{โยคสมฺปยุตฺตา เจว ธมฺมา โน จ โยคา.}}}\csromangathastanza{\csromangathawak{\csromangathaitemline{43}{โยควิปฺปยุตฺตา โข ปน ธมฺมา โยคนิยาปิ.} \\ \csromangathacontline{อโยคนิยาปิ.}}}}

\vspace{\tipitakaparskip}
\csromancenter{โยคโคจฺฉกํ.}
\csromanmatikaanchor{mtk.10}
\csromantocmarkref{section}{นีวรณโคจฺฉก}{mtk.10}
\bookmark[dest={mtk.10},level=1]{นีวรณโคจฺฉก}
\csromanheader{1}{นีวรณ\-โคจฺฉก}
\csromanhangingbegin
\hangingheaditem{44}{นีวรณา ธมฺมา. \mbox{(1158, 1503)}}
\hangingbody{โน นีวรณา ธมฺมา. \mbox{(1169, 1504)}}
\csromanhangingend

\csromanhangingbegin
\hangingheaditem{45}{นีวรณิยา ธมฺมา. \mbox{(1170, 1505)}}
\hangingbody{อนีวรณิยา ธมฺมา. \mbox{(1171, 1506)}}
\csromanhangingend

\csromanhangingbegin
\hangingheaditem{46}{นีวรณ\-สมฺปยุตฺตา ธมฺมา. \mbox{(1172, 1507)}}
\hangingbody{นีวรณวิปฺปยุตฺตา ธมฺมา. \mbox{(1173, 1508)}}
\csromanhangingend

\csromanhangingbegin
\hangingheaditem{47}{นีวรณา เจว ธมฺมา นีวรณิยา จ. \mbox{(1174, 1509)}}
\hangingbody{นีวรณิยา เจว ธมฺมา โน จ นีวรณา. \mbox{(1175, 1510)}}
\csromanhangingend

\csromanhangingbegin
\hangingheaditem{48}{นีวรณา เจว ธมฺมา นีวรณ\-สมฺปยุตฺตา จ. \mbox{(1176, 1511)}}
\hangingbody{นีวรณสมฺปยุตฺตา เจว ธมฺมา โน จ นีวรณา. \mbox{(1177, 1512)}}
\csromanhangingend

\csromanhangingbegin
\hangingheaditem{49}{นีวรณ\-วิปฺปยุตฺตา โข ปน ธมฺมา นีวรณิยาปิ. \mbox{(1178, 1513)}}
\hangingbody{อนีวรณิยาปิ. \mbox{(1179, 1514)}}
\csromanhangingend

\vspace{\tipitakaparskip}
\csromancenter{นีวรณ\-โคจฺฉกํ.}
\csromanmatikaanchor{mtk.11}
\csromantocmarkref{section}{ปรามาสโคจฺฉก}{mtk.11}
\bookmark[dest={mtk.11},level=1]{ปรามาสโคจฺฉก}
\csromanheader{1}{ปรามาส\-โคจฺฉก}
\csromanhangingbegin
\hangingheaditem{50}{\csromanfolio{9}ปรามาสา ธมฺมา. \mbox{(1180, 1515)}}
\hangingbody{โน ปรามาสา ธมฺมา. \mbox{(1182, 1516)}}
\csromanhangingend

\csromanhangingbegin
\hangingheaditem{51}{ปรามฏฺฐา ธมฺมา. \mbox{(1183, 1517)}}
\hangingbody{อปรามฏฺฐา ธมฺมา. \mbox{(1184, 1518)}}
\csromanhangingend

\csromanhangingbegin
\hangingheaditem{52}{ปรามาส\-สมฺปยุตฺตา ธมฺมา. \mbox{(1185, 1519)}}
\hangingbody{ปรามาสวิปฺปยุตฺตา ธมฺมา. \mbox{(1186, 1520)}}
\csromanhangingend

\csromanhangingbegin
\hangingheaditem{53}{ปรามาสา เจว ธมฺมา ปรามฏฺฐา จ. \mbox{(1187, 1521)}}
\hangingbody{ปรามฏฺฐา เจว ธมฺมา โน จ ปรามาสา. \mbox{(1188, 1522)}}
\csromanhangingend

\proseitem{54}{ปรามาส\-วิปฺปยุตฺตา โข ปน ธมฺมา ปรามฏฺฐาปิ. \mbox{(1189, 1523)}}

\setlength{\csromangathapagewidth}{0pt}
\setlength{\csromangathawakwidth}{0pt}
\csromangathameasuregroup{}{อปรามฏฺฐาปิ. \\ (1190, \\ 1524) ปรามาสโคจฺฉกํ.}
\csromangathameasurewak{อปรามฏฺฐาปิ. \\ (1190, \\ 1524) ปรามาสโคจฺฉกํ.}
\setlength{\csromangathaleftcol}{0pt}
\csromangathagroup{\csromangathastanza{\csromangathawak{อปรามฏฺฐาปิ. \\ (1190, \\ 1524) ปรามาส\-โคจฺฉกํ.}}}

\csromanmatikaanchor{mtk.12}
\csromantocmarkref{section}{มหนฺตรทุก}{mtk.12}
\bookmark[dest={mtk.12},level=1]{มหนฺตรทุก}
\csromanheader{1}{\textbf{มหนฺตรทุก}}
\csromanhangingbegin
\hangingheaditem{55}{สารมฺมณา ธมฺมา. \mbox{(1191, 1525)}}
\hangingbody{อนารมฺมณา ธมฺมา. \mbox{(1192, 1526)}}
\csromanhangingend

\csromanhangingbegin
\hangingheaditem{56}{จิตฺตา ธมฺมา. \mbox{(1193, 1527)}}
\hangingbody{โน จิตฺตา ธมฺมา. \mbox{(1194, 1528)}}
\csromanhangingend

\csromanhangingbegin
\hangingheaditem{57}{เจตสิกา ธมฺมา. \mbox{(1195, 1529)}}
\hangingbody{อเจตสิกา ธมฺมา. \mbox{(1196, 1530)}}
\csromanhangingend

\csromanhangingbegin
\hangingheaditem{58}{จิตฺต\-สมฺปยุตฺตา ธมฺมา. \mbox{(1197, 1531)}}
\hangingbody{จิตฺตวิปฺปยุตฺตา ธมฺมา. \mbox{(1198, 1532)}}
\csromanhangingend

\csromanhangingbegin
\hangingheaditem{59}{จิตฺตสํสฏฺฐา ธมฺมา. \mbox{(1199, 1533)}}
\hangingbody{จิตฺตวิสํสฏฺฐา ธมฺมา. \mbox{(1200, 1534)}}
\csromanhangingend

\csromanhangingbegin
\hangingheaditem{60}{จิตฺต\-สมุฏฺฐานา ธมฺมา. \mbox{(1201, 1535)}}
\hangingbody{โน จิตฺตสมุฏฺฐานา ธมฺมา. \mbox{(1202, 1536)}}
\csromanhangingend

\csromanhangingbegin
\hangingheaditem{61}{\csromanfolio{10}จิตฺตสหภุโน ธมฺมา. \mbox{(1203, 1537)}}
\hangingbody{โน จิตฺตสหภุโน ธมฺมา. \mbox{(1204, 1538)}}
\csromanhangingend

\csromanhangingbegin
\hangingheaditem{62}{จิตฺตา\-นุปริวตฺติโน ธมฺมา. \mbox{(1205, 1539)}}
\hangingbody{โน จิตฺตานุปริวตฺติโน ธมฺมา. \mbox{(1206, 1540)}}
\csromanhangingend

\csromanhangingbegin
\hangingheaditem{63}{จิตฺต\-สํสฏฺฐ\-สมุฏฺฐานา ธมฺมา. \mbox{(1207, 1541)}}
\hangingbody{โน จิตฺตสํสฏฺฐสมุฏฺฐานา ธมฺมา. \mbox{(1208, 1542)}}
\csromanhangingend

\csromanhangingbegin
\hangingheaditem{64}{จิตฺต\-สํสฏฺฐ\-สมุฏฺฐาน\-สหภุโน ธมฺมา. \mbox{(1209, 1543)}}
\hangingbody{โน จิตฺตสํสฏฺฐสมุฏฺฐานสหภุโน ธมฺมา. \mbox{(1210, 1544)}}
\csromanhangingend

\csromanhangingbegin
\hangingheaditem{65}{จิตฺต\-สํสฏฺฐ\-สมุฏฺฐานา\-นุปริวตฺติโน ธมฺมา. \mbox{(1211, 1545)}}
\hangingbody{โน จิตฺตสํสฏฺฐสมุฏฺฐานานุปริวตฺติโน ธมฺมา. \mbox{(1212, 1546)}}
\csromanhangingend

\csromanhangingbegin
\hangingheaditem{66}{อชฺฌตฺติกา ธมฺมา. \mbox{(1213, 1547)}}
\hangingbody{พาหิรา ธมฺมา. \mbox{(1214, 1548)}}
\csromanhangingend

\csromanhangingbegin
\hangingheaditem{67}{อุปาทา ธมฺมา. \mbox{(1215, 1549)}}
\hangingbody{โน อุปาทา ธมฺมา. \mbox{(1216, 1550)}}
\csromanhangingend

\proseitem{68}{อุปาทิณฺณา\csromansharedfootnote{4bde756e3115a3e3}{อุปาทินฺนา (สี, สฺยา)} ธมฺมา. \mbox{(1217, 1551)}}

\setlength{\csromangathapagewidth}{0pt}
\setlength{\csromangathawakwidth}{0pt}
\csromangathameasuregroup{}{อนุปาทิณฺณา ธมฺมา. \\ (1218, \\ 1552) มหนฺตรทุกํ.}
\csromangathameasurewak{อนุปาทิณฺณา ธมฺมา. \\ (1218, \\ 1552) มหนฺตรทุกํ.}
\setlength{\csromangathaleftcol}{0pt}
\csromangathagroup{\csromangathastanza{\csromangathawak{อนุปาทิณฺณา ธมฺมา. \\ (1218, \\ 1552) มหนฺตรทุกํ.}}}

\csromanmatikaanchor{mtk.13}
\csromantocmarkref{section}{อุปาทานโคจฺฉก}{mtk.13}
\bookmark[dest={mtk.13},level=1]{อุปาทานโคจฺฉก}
\csromanheader{1}{อุปาทาน\-โคจฺฉก}
\csromanhangingbegin
\hangingheaditem{69}{อุปาทานา ธมฺมา. \mbox{(1219, 1553)}}
\hangingbody{โน อุปาทานา ธมฺมา. \mbox{(1224, 1554)}}
\csromanhangingend

\csromanhangingbegin
\hangingheaditem{70}{อุปาทานิยา ธมฺมา. \mbox{(1225, 1555)}}
\hangingbody{อนุปาทานิยา ธมฺมา. \mbox{(1226, 1556)}}
\csromanhangingend

\csromanhangingbegin
\hangingheaditem{71}{อุปาทาน\-สมฺปยุตฺตา ธมฺมา. \mbox{(1227, 1557)}}
\hangingbody{อุปาทานวิปฺปยุตฺตา ธมฺมา. \mbox{(1228, 1558)}}
\csromanhangingend

\csromanhangingbegin
\hangingheaditem{72}{อุปาทานา เจว ธมฺมา อุปาทานิยา จ. \mbox{(1229, 1559)}}
\hangingbody{อุปาทานิยา เจว ธมฺมา โน จ อุปาทานา. \mbox{(1230, 1560)}}
\csromanhangingend

\csromanhangingbegin
\hangingheaditem{73}{\csromanfolio{11}อุปาทานา เจว ธมฺมา อุปาทาน\-สมฺปยุตฺตา จ. \mbox{(1231, 1561)}}
\hangingbody{อุปาทานสมฺปยุตฺตา เจว ธมฺมา โน จ อุปาทานา. \mbox{(1232, 1562)}}
\csromanhangingend

\csromanhangingbegin
\hangingheaditem{74}{อุปาทาน\-วิปฺปยุตฺตา โข ปน ธมฺมา อุปาทานิยาปิ. \mbox{(1233, 1563)}}
\hangingbody{อนุปาทานิยาปิ. \mbox{(1234, 1564)}}
\csromanhangingend

\vspace{\tipitakaparskip}
\csromancenter{อุปาทาน\-โคจฺฉกํ.}
\csromanmatikaanchor{mtk.14}
\csromantocmarkref{section}{กิเลสโคจฺฉก}{mtk.14}
\bookmark[dest={mtk.14},level=1]{กิเลสโคจฺฉก}
\csromanheader{1}{\textbf{กิเลสโคจฺฉก}}
\csromanhangingbegin
\hangingheaditem{75}{กิเลสา ธมฺมา. \mbox{(1235, 1565)}}
\hangingbody{โน กิเลสา ธมฺมา. \mbox{(1246, 1566)}}
\csromanhangingend

\csromanhangingbegin
\hangingheaditem{76}{สํกิเลสิกา ธมฺมา. \mbox{(1247, 1567)}}
\hangingbody{อสํกิเลสิกา ธมฺมา. \mbox{(1248, 1568)}}
\csromanhangingend

\csromanhangingbegin
\hangingheaditem{77}{สํกิลิฏฺฐา ธมฺมา. \mbox{(1249, 1569)}}
\hangingbody{อสํกิลิฏฺฐา ธมฺมา. \mbox{(1250, 1570)}}
\csromanhangingend

\csromanhangingbegin
\hangingheaditem{78}{กิเลส\-สมฺปยุตฺตา ธมฺมา. \mbox{(1251, 1571)}}
\hangingbody{กิเลสวิปฺปยุตฺตา ธมฺมา. \mbox{(1252, 1572)}}
\csromanhangingend

\csromanhangingbegin
\hangingheaditem{79}{กิเลสา เจว ธมฺมา สํกิเลสิกา จ. \mbox{(1253, 1573)}}
\hangingbody{สํกิเลสิกา เจว ธมฺมา โน จ กิเลสา. \mbox{(1254, 1574)}}
\csromanhangingend

\csromanhangingbegin
\hangingheaditem{80}{กิเลสา เจว ธมฺมา สํกิลิฏฺฐา จ. \mbox{(1255, 1575)}}
\hangingbody{สํกิลิฏฺฐา เจว ธมฺมา โน จ กิเลสา. \mbox{(1256, 1576)}}
\csromanhangingend

\csromanhangingbegin
\hangingheaditem{81}{กิเลสา เจว ธมฺมา กิเลส\-สมฺปยุตฺตา จ. \mbox{(1257, 1577)}}
\hangingbody{กิเลสสมฺปยุตฺตา เจว ธมฺมา โน จ กิเลสา. \mbox{(1258, 1578)}}
\csromanhangingend

\csromanhangingbegin
\hangingheaditem{82}{กิเลส\-วิปฺปยุตฺตา โข ปน ธมฺมา สํกิเลสิกาปิ. \mbox{(1259, 1579)}}
\hangingbody{อสํกิเลสิกาปิ. \mbox{(1260, 1580)}}
\csromanhangingend

\vspace{\tipitakaparskip}
\csromancenter{กิเลส\-โคจฺฉกํ.}
\csromanmatikaanchor{mtk.15}
\csromantocmarkref{section}{ปิฏฺฐิทุก}{mtk.15}
\bookmark[dest={mtk.15},level=1]{ปิฏฺฐิทุก}
\csromanheader{1}{\textbf{ปิฏฺฐิทุก}}
\csromanhangingbegin
\hangingheaditem{83}{\csromanfolio{12}ทสฺสเนน ปหาตพฺพา ธมฺมา. \mbox{(1261, 1581)}}
\hangingbody{น ทสฺสเนน ปหาตพฺพา ธมฺมา. \mbox{(1265, 1582)}}
\csromanhangingend

\csromanhangingbegin
\hangingheaditem{84}{ภาวนาย ปหาตพฺพา ธมฺมา. \mbox{(1266, 1583)}}
\hangingbody{น ภาวนาย ปหาตพฺพา ธมฺมา. \mbox{(1267, 1584)}}
\csromanhangingend

\csromanhangingbegin
\hangingheaditem{85}{ทสฺสเนน ปหาตพฺพ\-เหตุกา ธมฺมา. \mbox{(1268, 1585)}}
\hangingbody{น ทสฺสเนน ปหาตพฺพเหตุกา ธมฺมา. \mbox{(1272, 1586)}}
\csromanhangingend

\csromanhangingbegin
\hangingheaditem{86}{ภาวนาย ปหาตพฺพ\-เหตุกา ธมฺมา. \mbox{(1273, 1587)}}
\hangingbody{น ภาวนาย ปหาตพฺพเหตุกา ธมฺมา. \mbox{(1274, 1588)}}
\csromanhangingend

\csromanhangingbegin
\hangingheaditem{87}{สวิตกฺกา ธมฺมา. \mbox{(1275, 1589)}}
\hangingbody{อวิตกฺกา ธมฺมา. \mbox{(1276, 1590)}}
\csromanhangingend

\csromanhangingbegin
\hangingheaditem{88}{สวิจารา ธมฺมา. \mbox{(1277, 1591)}}
\hangingbody{อวิจารา ธมฺมา. \mbox{(1278, 1592)}}
\csromanhangingend

\csromanhangingbegin
\hangingheaditem{89}{สปฺปีติกา ธมฺมา. \mbox{(1279, 1593)}}
\hangingbody{อปฺปีติกา ธมฺมา. \mbox{(1280, 1594)}}
\csromanhangingend

\csromanhangingbegin
\hangingheaditem{90}{ปีติสหคตา ธมฺมา. \mbox{(1281, 1595)}}
\hangingbody{น ปีติสหคตา ธมฺมา. \mbox{(1282, 1596)}}
\csromanhangingend

\csromanhangingbegin
\hangingheaditem{91}{สุขสหคตา ธมฺมา. \mbox{(1283, 1597)}}
\hangingbody{น สุขสหคตา ธมฺมา. \mbox{(1284, 1598)}}
\csromanhangingend

\csromanhangingbegin
\hangingheaditem{92}{อุเปกฺขา\-สหคตา ธมฺมา. \mbox{(1285, 1599)}}
\hangingbody{น อุเปกฺขาสหคตา ธมฺมา. \mbox{(1286, 1600)}}
\csromanhangingend

\csromanhangingbegin
\hangingheaditem{93}{กามาวจรา ธมฺมา. \mbox{(1287, 1601)}}
\hangingbody{น กามาวจรา ธมฺมา. \mbox{(1288, 1602)}}
\csromanhangingend

\csromanhangingbegin
\hangingheaditem{94}{รูปาวจรา ธมฺมา. \mbox{(1289, 1603)}}
\hangingbody{น รูปาวจรา ธมฺมา. \mbox{(1290, 1604)}}
\csromanhangingend

\csromanhangingbegin
\hangingheaditem{95}{อรูปาวจรา ธมฺมา. \mbox{(1291, 1605)}}
\hangingbody{น อรูปาวจรา ธมฺมา. \mbox{(1292, 1606)}}
\csromanhangingend

\csromanhangingbegin
\hangingheaditem{96}{\csromanfolio{13}ปริยาปนฺนา ธมฺมา. \mbox{(1293, 1607)}}
\hangingbody{อปริยาปนฺนา ธมฺมา. \mbox{(1294, 1608)}}
\csromanhangingend

\csromanhangingbegin
\hangingheaditem{97}{นิยฺยานิกา ธมฺมา. \mbox{(1295, 1609)}}
\hangingbody{อนิยฺยานิกา ธมฺมา. \mbox{(1296, 1610)}}
\csromanhangingend

\csromanhangingbegin
\hangingheaditem{98}{นิยตา ธมฺมา. \mbox{(1297, 1611)}}
\hangingbody{อนิยตา ธมฺมา. \mbox{(1298, 1612)}}
\csromanhangingend

\csromanhangingbegin
\hangingheaditem{99}{สอุตฺตรา ธมฺมา. \mbox{(1299, 1613)}}
\hangingbody{อนุตฺตรา ธมฺมา. \mbox{(1300, 1614)}}
\csromanhangingend

\proseitem{100}{สรณา ธมฺมา. \mbox{(1301, 1615)}}

\csromanheader{2}{อรณา ธมฺมา. \mbox{(1302, 1616)} ปิฏฺฐิทุกํ. อภิธมฺม\-ทุก\-มาติกา.}
\csromanmatikaanchor{mtk.16}
\csromantocmarkref{section}{สุตฺตนฺติกทุกมาติกา}{mtk.16}
\bookmark[dest={mtk.16},level=1]{สุตฺตนฺติกทุกมาติกา}
\csromanheader{1}{สุตฺตนฺติก\-ทุก\-มาติกา}
\prose{101. วิชฺชาภาคิโน ธมฺมา. (1303)}

\prose{อวิชฺชาภาคิโน ธมฺมา. (1304)}

\prose{102. วิชฺชูปมา ธมฺมา. (1305)}

\prose{วชิรูปมา ธมฺมา. (1306)}

\prose{103. พาลา ธมฺมา. (1307)}

\prose{ปณฺฑิตา ธมฺมา. (1308)}

\prose{104. กณฺหา ธมฺมา. (1309)}

\prose{สุกฺกา ธมฺมา. (1310)}

\prose{105. ตปนียา ธมฺมา. (1311)}

\prose{อตปนียา ธมฺมา. (1312)}

\prose{106. อธิวจนา ธมฺมา. (1313)}

\prose{อธิวจน\-ปถา ธมฺมา. (1313)}

\prose{\csromanfolio{14}107. นิรุตฺติ ธมฺมา. (1314)}

\prose{นิรุตฺติปถา ธมฺมา. (1314)}

\prose{108. ปญฺญตฺติ ธมฺมา. (1315)}

\prose{ปญฺญตฺติปถา ธมฺมา. (1315)}

\prose{109. นามญฺจ. (1316)}

\prose{รูปญฺจ. (1317)}

\prose{110. อวิชฺชา จ. (1318)}

\prose{ภวตณฺหา จ. (1319)}

\prose{111. ภวทิฏฺฐิ จ. (1320)}

\prose{วิภวทิฏฺฐิ จ. (1321)}

\prose{112. สสฺสตทิฏฺฐิ จ. (1322)}

\prose{อุจฺเฉททิฏฺฐิ จ. (1323)}

\prose{113. อนฺตวา ทิฏฺฐิ จ. (1324)}

\prose{อนนฺตวา ทิฏฺฐิ จ. (1325)}

\prose{114. ปุพฺพนฺตา\-นุทิฏฺฐิ จ. (1326)}

\prose{อปรนฺตา\-นุทิฏฺฐิ จ. (1327)}

\prose{115. อหิริกญฺจ. (1328)}

\prose{อโนตฺตปฺปญฺจ. (1329)}

\prose{116. หิรี จ. (1330)}

\prose{โอตฺตปฺปญฺจ. (1331)}

\prose{117. โทวจสฺสตา จ. (1332)}

\prose{ปาปมิตฺตตา จ. (1333)}

\prose{118. โสวจสฺสตา จ. (1334)}

\prose{กลฺยาณมิตฺตตา จ. (1335)}

\prose{119. อาปตฺติกุสลตา จ. (1336)}

\prose{อาปตฺติวุฏฺฐาน\-กุสลตา จ. (1337)}

\prose{120. สมาปตฺติ\-กุสลตา จ. (1338)}

\prose{สมาปตฺติ\-วุฏฺฐาน\-กุสลตา จ. (1339)}

\prose{\csromanfolio{15}121. ธาตุกุสลตา จ. (1340)}

\prose{มนสิการ\-กุสลตา จ. (1341)}

\prose{122. อายตน\-กุสลตา จ. (1342)}

\prose{ปฏิจฺจสมุปฺปาท\-กุสลตา จ. (1343)}

\prose{123. ฐานกุสลตา จ. (1344)}

\prose{อฏฺฐาน\-กุสลตา จ. (1345)}

\prose{124. อชฺชโว จ. (1346)}

\prose{มทฺทโว จ. (1347)}

\prose{125. ขนฺติ จ. (1348)}

\prose{โสรจฺจญฺจ. (1349)}

\prose{126. สาขลฺยญฺจ. (1350)}

\prose{ปฏิสนฺถาโร จ\csromansharedfootnote{b318131acb1918a9}{ปฏิสนฺธาโร จ (ก)}. (1351)}

\prose{127. อินฺทฺริเยสุ อคุตฺตทฺวารตา จ. (1352)}

\prose{โภชเน อมตฺตญฺญุตา จ. (1353)}

\prose{128. อินฺทฺริเยสุ คุตฺตทฺวารตา จ. (1354)}

\prose{โภชเน มตฺตญฺญุตา จ. (1355)}

\prose{129. มุฏฺฐสจฺจญฺจ. (1356)}

\prose{อสมฺปชญฺญญฺจ. (1357)}

\prose{130. สติ จ. (1358)}

\prose{สมฺปชญฺญญฺจ. (1359)}

\prose{131. ปฏิสงฺขานพลญฺจ. (1360)}

\prose{ภาวนาพลญฺจ. (1361)}

\prose{132. สมโถ จ. (1362)}

\prose{วิปสฺสนา จ. (1363)}

\prose{133. สมถนิมิตฺตญฺจ. (1364)}

\prose{ปคฺคาห\-นิมิตฺตญฺจ. (1365)}

\prose{\csromanfolio{16}134. ปคฺคาโห จ. (1366)}

\prose{อวิกฺเขโป จ. (1367)}

\prose{135. สีลวิปตฺติ จ. (1368)}

\prose{ทิฏฺฐิวิปตฺติ จ. (1369)}

\prose{136. สีลสมฺปทา จ. (1370)}

\prose{ทิฏฺฐิสมฺปทา จ. (1371)}

\prose{137. สีลวิสุทฺธิ จ. (1372)}

\prose{ทิฏฺฐิวิสุทฺธิ จ. (1373)}

\prose{138. ทิฏฺฐิวิสุทฺธิ โข ปน. (1374)}

\prose{ยถาทิฏฺฐิสฺส จ ปธานํ. (1375)}

\prose{139. สํเวโค จ สํเวชนิเยสุ ฐาเนสุ. (1376)}

\prose{สํวิคฺคสฺส จ โยนิโส ปธานํ. (1377)}

\prose{140. อสนฺตุฏฺฐิตา จ กุสเลสุ ธมฺเมสุ. (1378)}

\prose{อปฺปฏิวานิตา จ ปธานสฺมิํ. (1379)}

\prose{141. วิชฺชา จ. (1380)}

\prose{วิมุตฺติ จ. (1381)}

\prose{142. ขเย ญาณํ. (1382)}

\prose{อนุปฺปาเท ญาณนฺติ. (1383)}

\vspace{\tipitakaparskip}
\csromancenter{สุตฺตนฺติก\-ทุก\-มาติกา\csromansharedfootnote{e18ae947fe43c690}{สุตฺตนฺตมาติกา (สฺยา)}.}
\nitthitam{\textbf{มาติกา นิฏฺฐิตา.}}
\csromanmatikaanchor{mtk.17}
\csromantocmarknopage{chapter}{1. จิตฺตุปฺปาทกณฺฑ}{mtk.17}
\bookmark[dest={mtk.17},level=0]{1. จิตฺตุปฺปาทกณฺฑ}
\chapterheadpageread{1. จิตฺตุปฺปาท\-กณฺฑ}
\csromanmatikaanchor{mtk.18}
\csromantocmarkref{section}{กามาวจรกุสลปทภาชนี}{mtk.18}
\bookmark[dest={mtk.18},level=1]{กามาวจรกุสลปทภาชนี}
\csromanheader{1}{กามาวจร\-กุสล}
\csromanheader{2}{\textbf{ปทภาชนี}}
\proseitem{1}{\csromanfolio{17}กตเม ธมฺมา กุสลา. ยสฺมิํ สมเย กามาวจรํ กุสลํ จิตฺตํ อุปฺปนฺนํ โหติ โสมนสฺส\-สหคตํ ญาณ\-สมฺปยุตฺตํ รูปารมฺมณํ วา สทฺทารมฺมณํ วา คนฺธารมฺมณํ วา รสารมฺมณํ วา โผฏฺฐพฺพา\-รมฺมณํ วา ธมฺมารมฺมณํ วา ยํ ยํ วา ปนารพฺภ. ตสฺมิํ สมเย ผสฺโส โหติ, เวทนา โหติ, สญฺญา โหติ, เจตนา โหติ, จิตฺตํ โหติ, วิตกฺโก โหติ, วิจาโร โหติ, ปีติ โหติ, สุขํ โหติ, จิตฺตสฺเส\-กคฺคตา โหติ, สทฺธินฺทฺริยํ โหติ, วีริยินฺทฺริยํ\csromansharedfootnote{bba6a531d75f224c}{วิริยินฺทฺริยํ (สี, สฺยา)} โหติ, สตินฺทฺริยํ โหติ, สมาธินฺทฺริยํ โหติ, ปญฺญินฺทฺริยํ โหติ, มนินฺทฺริยํ โหติ, โสมนสฺสิ\-นฺทฺริยํ โหติ, ชีวิตินฺทฺริยํ โหติ, สมฺมาทิฏฺฐิ โหติ, สมฺมาสงฺกปฺโป โหติ, สมฺมาวายาโม โหติ, สมฺมาสติ โหติ, สมฺมาสมาธิ โหติ, สทฺธาพลํ โหติ, วีริยพลํ\csromansharedfootnote{12f5d13d0c63c27e}{วิริยพลํ (สี, สฺยา)} โหติ, สติพลํ โหติ, สมาธิพลํ โหติ, ปญฺญาพลํ โหติ, หิริพลํ โหติ, โอตฺตปฺปพลํ โหติ, อโลโภ โหติ, อโทโส โหติ, อโมโห โหติ, อนภิชฺฌา โหติ, อพฺยาปาโท โหติ, สมฺมาทิฏฺฐิ โหติ, หิรี โหติ, โอตฺตปฺปํ โหติ, กายปสฺสทฺธิ\csromansharedfootnote{5c1732c31ef2ba04}{กายปฺปสฺสทฺธิ (สฺยา)} โหติ, จิตฺตปสฺสทฺธิ\csromansharedfootnote{4601fcd7db2d7f99}{จิตฺตปฺปสฺสทฺธิ (สฺยา)} โหติ, กายลหุตา โหติ, จิตฺตลหุตา โหติ, กายมุทุตา โหติ, จิตฺตมุทุตา โหติ, กายกมฺมญฺญตา โหติ, จิตฺต\-กมฺมญฺญตา โหติ, กายปาคุญฺญตา โหติ, จิตฺตปาคุญฺญตา โหติ, กายุชุกตา\csromansharedfootnote{489e6ef1d077ccdb}{กายุชฺชุกตา (สี, ก)} โหติ, จิตฺตุชุกตา\csromansharedfootnote{4e63b79de9e550f6}{จิตฺตุชฺชุกตา (สี, ก)} โหติ, สติ โหติ, สมฺปชญฺญํ โหติ, สมโถ โหติ, วิปสฺสนา โหติ, ปคฺคาโห โหติ, อวิกฺเขโป โหติ. เย วา ปน ตสฺมิํ สมเย อญฺเญปิ อตฺถิ ปฏิจฺจ\-สมุปฺปนฺนา อรูปิโน ธมฺมา. อิเม ธมฺมา กุสลา.}

\proseitem{2}{\csromanfolio{18}กตโม ตสฺมิํ สมเย ผสฺโส โหติ. โย ตสฺมิํ สมเย ผสฺโส ผุสนา สมฺผุสนา สมฺผุสิตตฺตํ. อยํ ตสฺมิํ สมเย ผสฺโส โหติ.}

\proseitem{3}{กตมา ตสฺมิํ สมเย เวทนา โหติ. ยํ ตสฺมิํ สมเย ตชฺชา\-มโนวิญฺญาณธาตุ\-สมฺผสฺสชํ เจตสิกํ สาตํ เจตสิกํ สุขํ เจโต\-สมฺผสฺสชํ สาตํ สุขํ เวทยิตํ เจโต\-สมฺผสฺสชา สาตา สุขา เวทนา. อยํ ตสฺมิํ สมเย เวทนา โหติ.}

\proseitem{4}{กตมา ตสฺมิํ สมเย สญฺญา โหติ. ยา ตสฺมิํ สมเย ตชฺชา\-มโนวิญฺญาณธาตุ\-สมฺผสฺสชา สญฺญา สญฺชานนา สญฺชานิตตฺตํ. อยํ ตสฺมิํ สมเย สญฺญา โหติ.}

\proseitem{5}{กตมา ตสฺมิํ สมเย เจตนา โหติ. ยา ตสฺมิํ สมเย ตชฺชา\-มโนวิญฺญาณธาตุ\-สมฺผสฺสชา เจตนา สญฺเจตนา เจตยิตตฺตํ\csromansharedfootnote{72ae37dc0cd467c0}{สญฺเจตยิตตฺตํ (สฺยา)}. อยํ ตสฺมิํ สมเย เจตนา โหติ.}

\proseitem{6}{กตมํ ตสฺมิํ สมเย จิตฺตํ โหติ. ยํ ตสฺมิํ สมเย จิตฺตํ มโน มานสํ หทยํ ปณฺฑรํ มโน มนายตนํ มนินฺทฺริยํ วิญฺญาณํ วิญฺญาณ\-กฺขนฺโธ ตชฺชา\-มโนวิญฺญาณ\-ธาตุ. อิทํ ตสฺมิํ สมเย จิตฺตํ โหติ.}

\proseitem{7}{กตโม ตสฺมิํ สมเย วิตกฺโก โหติ. โย ตสฺมิํ สมเย ตกฺโก วิตกฺโก สงฺกปฺโป อปฺปนา พฺยปฺปนา เจตโส อภินิโรปนา สมฺมาสงฺกปฺโป. อยํ ตสฺมิํ สมเย วิตกฺโก โหติ.}

\proseitem{8}{กตโม ตสฺมิํ สมเย วิจาโร โหติ. โย ตสฺมิํ สมเย จาโร วิจาโร อนุวิจาโร อุปวิจาโร จิตฺตสฺส อนุสนฺธานตา อนุเปกฺขนตา. อยํ ตสฺมิํ สมเย วิจาโร โหติ.}

\proseitem{9}{กตมา ตสฺมิํ สมเย ปีติ โหติ. ยา ตสฺมิํ สมเย ปีติ ปาโมชฺชํ อาโมทนา ปโมทนา หาโส ปหาโส วิตฺติ โอทคฺยํ อตฺตมนตา จิตฺตสฺส. อยํ ตสฺมิํ สมเย ปีติ โหติ.}

\chapterheadpageread{1. จิตฺตุปฺปาท\-กณฺฑ}
\proseitem{10}{\csromanfolio{19}กตมํ ตสฺมิํ สมเย สุขํ โหติ. ยํ ตสฺมิํ สมเย เจตสิกํ สาตํ เจตสิกํ สุขํ เจโต\-สมฺผสฺสชํ สาตํ สุขํ เวทยิตํ เจโต\-สมฺผสฺสชา สาตา สุขา เวทนา. อิทํ ตสฺมิํ สมเย สุขํ โหติ.}

\proseitem{11}{กตมา ตสฺมิํ สมเย จิตฺตสฺเส\-กคฺคตา โหติ. ยา ตสฺมิํ สมเย จิตฺตสฺส ฐิติ สณฺฐิติ อวฏฺฐิติ อวิสาหาโร อวิกฺเขโป อวิสาหฏ\-มานสตา สมโถ สมาธินฺทฺริยํ สมาธิพลํ สมฺมาสมาธิ. อยํ ตสฺมิํ สมเย จิตฺตสฺเส\-กคฺคตา โหติ.}

\proseitem{12}{กตมํ ตสฺมิํ สมเย สทฺธินฺทฺริยํ โหติ. ยา ตสฺมิํ สมเย สทฺธา สทฺทหนา โอกปฺปนา อภิปฺปสาโท สทฺธา สทฺธินฺทฺริยํ สทฺธาพลํ. อิทํ ตสฺมิํ สมเย สทฺธินฺทฺริยํ โหติ.}

\proseitem{13}{กตมํ ตสฺมิํ สมเย วีริยินฺทฺริยํ โหติ. โย ตสฺมิํ สมเย เจตสิโก วีริยารมฺโภ\csromansharedfootnote{17be1ae89afc2d81}{วิริยารมฺโภ (สี, สฺยา)} นิกฺกโม ปรกฺกโม อุยฺยาโม วายาโม อุสฺสาโห อุสฺโสฬฺหี\csromansharedfootnote{821e414ba93cdbad}{อุสฺโสฬฺหิ (สี)} ถาโม ธิติ อสิถิล\-ปรกฺกมตา อนิกฺขิตฺต\-ฉนฺทตา อนิกฺขิตฺตธุรตา ธุรสมฺปคฺคาโห วีริยํ วีริยินฺทฺริยํ วีริยพลํ สมฺมาวายาโม. อิทํ ตสฺมิํ สมเย วีริยินฺทฺริยํ โหติ.}

\proseitem{14}{กตมํ ตสฺมิํ สมเย สตินฺทฺริยํ โหติ. ยา ตสฺมิํ สมเย สติ อนุสฺสติ ปฏิสฺสติ สติ สรณตา ธารณตา อปิลาปนตา อสมฺมุสฺสนตา สติ สตินฺทฺริยํ สติพลํ สมฺมาสติ. อิทํ ตสฺมิํ สมเย สตินฺทฺริยํ โหติ.}

\proseitem{15}{กตมํ ตสฺมิํ สมเย สมาธินฺทฺริยํ โหติ. ยา ตสฺมิํ สมเย จิตฺตสฺส ฐิติ สณฺฐิติ อวฏฺฐิติ อวิสาหาโร อวิกฺเขโป อวิสาหฏ\-มานสตา สมโถ สมาธินฺทฺริยํ สมาธิพลํ สมฺมาสมาธิ. อิทํ ตสฺมิํ สมเย สมาธินฺทฺริยํ โหติ.}

\proseitem{16}{กตมํ ตสฺมิํ สมเย ปญฺญินฺทฺริยํ โหติ. ยา ตสฺมิํ สมเย ปญฺญา ปชานนา วิจโย ปวิจโย ธมฺมวิจโย สลฺลกฺขณา อุปลกฺขณา ปจฺจุ\-ปลกฺขณา ปณฺฑิจฺจํ โกสลฺลํ เนปุญฺญํ เวภพฺยา จินฺตา อุปปริกฺขา ภูรี \csromanfolio{20}เมธา ปริณายิกา วิปสฺสนา สมฺปชญฺญํ ปโตโท ปญฺญา ปญฺญินฺทฺริยํ ปญฺญาพลํ ปญฺญาสตฺถํ ปญฺญาปาสาโท ปญฺญาอาโลโก ปญฺญาโอภาโส ปญฺญาปชฺโชโต ปญฺญารตนํ อโมโห ธมฺมวิจโย สมฺมาทิฏฺฐิ. อิทํ ตสฺมิํ สมเย ปญฺญินฺทฺริยํ โหติ.}

\proseitem{17}{กตมํ ตสฺมิํ สมเย มนินฺทฺริยํ โหติ. ยํ ตสฺมิํ สมเย จิตฺตํ มโน มานสํ หทยํ ปณฺฑรํ มโน มนายตนํ มนินฺทฺริยํ วิญฺญาณํ วิญฺญาณ\-กฺขนฺโธ ตชฺชา\-มโนวิญฺญาณ\-ธาตุ. อิทํ ตสฺมิํ สมเย มนินฺทฺริยํ โหติ.}

\proseitem{18}{กตมํ ตสฺมิํ สมเย โสมนสฺสิ\-นฺทฺริยํ โหติ. ยํ ตสฺมิํ สมเย เจตสิกํ สาตํ เจตสิกํ สุขํ เจโต\-สมฺผสฺสชํ สาตํ สุขํ เวทยิตํ เจโต\-สมฺผสฺสชา สาตา สุขา เวทนา. อิทํ ตสฺมิํ สมเย โสมนสฺสิ\-นฺทฺริยํ โหติ.}

\proseitem{19}{กตมํ ตสฺมิํ สมเย ชีวิตินฺทฺริยํ โหติ. โย เตสํ อรูปีนํ ธมฺมานํ อายุ ฐิติ ยปนา ยาปนา อิริยนา\csromansharedfootnote{86dbed6c33f17c90}{อิรียนา (สี)} วตฺตนา ปาลนา ชีวิตํ ชีวิตินฺทฺริยํ. อิทํ ตสฺมิํ สมเย ชีวิตินฺทฺริยํ โหติ.}

\proseitem{20}{กตมา ตสฺมิํ สมเย สมฺมาทิฏฺฐิ โหติ. ยา ตสฺมิํ สมเย ปญฺญา ปชานนา วิจโย ปวิจโย ธมฺมวิจโย สลฺลกฺขณา อุปลกฺขณา ปจฺจุ\-ปลกฺขณา ปณฺฑิจฺจํ โกสลฺลํ เนปุญฺญํ เวภพฺยา จินฺตา อุปปริกฺขา ภูรี เมธา ปริณายิกา วิปสฺสนา สมฺปชญฺญํ ปโตโท ปญฺญา ปญฺญินฺทฺริยํ ปญฺญาพลํ ปญฺญาสตฺถํ ปญฺญาปาสาโท ปญฺญาอาโลโก ปญฺญาโอภาโส ปญฺญาปชฺโชโต ปญฺญารตนํ อโมโห ธมฺมวิจโย สมฺมาทิฏฺฐิ. อยํ ตสฺมิํ สมเย สมฺมาทิฏฺฐิ โหติ.}

\proseitem{21}{กตโม ตสฺมิํ สมเย สมฺมาสงฺกปฺโป โหติ. โย ตสฺมิํ สมเย ตกฺโก วิตกฺโก สงฺกปฺโป อปฺปนา พฺยปฺปนา เจตโส อภินิโรปนา สมฺมาสงฺกปฺโป. อยํ ตสฺมิํ สมเย สมฺมาสงฺกปฺโป โหติ.}

\chapterheadpageread{1. จิตฺตุปฺปาท\-กณฺฑ}
\proseitem{22}{\csromanfolio{21}กตโม ตสฺมิํ สมเย สมฺมาวายาโม โหติ. โย ตสฺมิํ สมเย เจตสิโก วีริยารมฺโภ นิกฺกโม ปรกฺกโม อุยฺยาโม วายาโม อุสฺสาโห อุสฺโสฬฺหี ถาโม ธิติ อสิถิล\-ปรกฺกมตา อนิกฺขิตฺต\-ฉนฺทตา อนิกฺขิตฺตธุรตา ธุรสมฺปคฺคาโห วีริยํ วีริยินฺทฺริยํ วีริยพลํ สมฺมาวายาโม. อยํ ตสฺมิํ สมเย สมฺมาวายาโม โหติ.}

\proseitem{23}{กตมา ตสฺมิํ สมเย สมฺมาสติ โหติ. ยา ตสฺมิํ สมเย สติ อนุสฺสติ ปฏิสฺสติ สติ สรณตา ธารณตา อปิลาปนตา อสมฺมุสฺสนตา สติ สตินฺทฺริยํ สติพลํ สมฺมาสติ. อยํ ตสฺมิํ สมเย สมฺมาสติ โหติ.}

\proseitem{24}{กตโม ตสฺมิํ สมเย สมฺมาสมาธิ โหติ. ยา ตสฺมิํ สมเย จิตฺตสฺส ฐิติ สณฺฐิติ อวฏฺฐิติ อวิสาหาโร อวิกฺเขโป อวิสาหฏ\-มานสตา สมโถ สมาธินฺทฺริยํ สมาธิพลํ สมฺมาสมาธิ. อยํ ตสฺมิํ สมเย สมฺมาสมาธิ โหติ.}

\proseitem{25}{กตมํ ตสฺมิํ สมเย สทฺธาพลํ โหติ. ยา ตสฺมิํ สมเย สทฺธา สทฺทหนา โอกปฺปนา อภิปฺปสาโท สทฺธา สทฺธินฺทฺริยํ สทฺธาพลํ. อิทํ ตสฺมิํ สมเย สทฺธาพลํ โหติ.}

\proseitem{26}{กตมํ ตสฺมิํ สมเย วีริยพลํ โหติ. โย ตสฺมิํ สมเย เจตสิโก วีริยารมฺโภ นิกฺกโม ปรกฺกโม อุยฺยาโม วายาโม อุสฺสาโห อุสฺโสฬฺหี ถาโม ธิติ อสิถิล\-ปรกฺกมตา อนิกฺขิตฺตฉนฺทตฺตา อนิกฺขิตฺตธุรตา ธุรสมฺปคฺคาโห วีริยํ วีริยินฺทฺริยํ วีริยพลํ สมฺมาวายาโม. อิทํ ตสฺมิํ สมเย วีริยพลํ โหติ.}

\proseitem{27}{กตมํ ตสฺมิํ สมเย สติพลํ โหติ. ยา ตสฺมิํ สมเย สติ อนุสฺสติ ปฏิสฺสติ สติ สรณตา ธารณตา อปิลาปนตา อสมฺมุสฺสนตา สติ สตินฺทฺริยํ สติพลํ สมฺมาสติ. อิทํ ตสฺมิํ สมเย สติพลํ โหติ.}

\proseitem{28}{กตมํ ตสฺมิํ สมเย สมาธิพลํ โหติ. ยา ตสฺมิํ สมเย จิตฺตสฺส ฐิติ สณฺฐิติ อวฏฺฐิติ อวิสาหาโร อวิกฺเขโป อวิสาหฏ\-มานสตา สมโถ สมาธินฺทฺริยํ สมาธิพลํ สมฺมาสมาธิ. อิทํ ตสฺมิํ สมเย สมาธิพลํ โหติ.}

\proseitem{29}{\csromanfolio{22}กตมํ ตสฺมิํ สมเย ปญฺญาพลํ โหติ. ยา ตสฺมิํ สมเย ปญฺญา ปชานนา วิจโย ปวิจโย ธมฺมวิจโย สลฺลกฺขณา อุปลกฺขณา ปจฺจุ\-ปลกฺขณา ปณฺฑิจฺจํ โกสลฺลํ เนปุญฺญํ เวภพฺยา จินฺตา อุปปริกฺขา ภูรี เมธา ปริณายิกา วิปสฺสนา สมฺปชญฺญํ ปโตโท ปญฺญา ปญฺญินฺทฺริยํ ปญฺญาพลํ ปญฺญาสตฺถํ ปญฺญาปาสาโท ปญฺญาอาโลโก ปญฺญาโอภาโส ปญฺญาปชฺโชโต ปญฺญารตนํ อโมโห ธมฺมวิจโย สมฺมาทิฏฺฐิ. อิทํ ตสฺมิํ สมเย ปญฺญาพลํ โหติ.}

\proseitem{30}{กตมํ ตสฺมิํ สมเย หิริพลํ โหติ. ยํ ตสฺมิํ สมเย หิรียติ หิริยิตพฺเพน หิรียติ ปาปกานํ อกุสลานํ ธมฺมานํ สมาปตฺติยา. อิทํ ตสฺมิํ สมเย หิริพลํ โหติ.}

\proseitem{31}{กตมํ ตสฺมิํ สมเย โอตฺตปฺปพลํ โหติ. ยํ ตสฺมิํ สมเย โอตฺตปฺปติ โอตฺตปฺปิตพฺเพน โอตฺตปฺปติ ปาปกานํ อกุสลานํ ธมฺมานํ สมาปตฺติยา. อิทํ ตสฺมิํ สมเย โอตฺตปฺปพลํ โหติ.}

\proseitem{32}{กตโม ตสฺมิํ สมเย อโลโภ โหติ. โย ตสฺมิํ สมเย อโลโภ อลุพฺภนา อลุพฺภิตตฺตํ อสาราโค อสารชฺชนา อสารชฺชิตตฺตํ อนภิชฺฌา อโลโภ กุสลมูลํ. อยํ ตสฺมิํ สมเย อโลโภ โหติ.}

\proseitem{33}{กตโม ตสฺมิํ สมเย อโทโส โหติ. โย ตสฺมิํ สมเย อโทโส อทุสฺสนา อทุสฺสิตตฺตํ\csromansharedfootnote{db3f66cd7e5443dd}{อทูสนา อทูสิตตฺตํ (สฺยา)} อพฺยาปาโท อพฺยาปชฺโช อโทโส กุสลมูลํ. อยํ ตสฺมิํ สมเย อโทโส โหติ.}

\proseitem{34}{กตโม ตสฺมิํ สมเย อโมโห โหติ. ยา ตสฺมิํ สมเย ปญฺญา ปชานนา วิจโย ปวิจโย ธมฺมวิจโย สลฺลกฺขณา อุปลกฺขณา ปจฺจุ\-ปลกฺขณา ปณฺฑิจฺจํ โกสลฺลํ เนปุญฺญํ เวภพฺยา จินฺตา อุปปริกฺขา ภูรี เมธา ปริณายิกา วิปสฺสนา สมฺปชญฺญํ ปโตโท ปญฺญา ปญฺญินฺทฺริยํ ปญฺญาพลํ ปญฺญาสตฺถํ ปญฺญาปาสาโท ปญฺญาอาโลโก ปญฺญาโอภาโส ปญฺญาปชฺโชโต ปญฺญารตนํ อโมโห ธมฺมวิจโย สมฺมาทิฏฺฐิ อโมโห กุสลมูลํ. อยํ ตสฺมิํ สมเย อโมโห โหติ.}

\chapterheadpageread{1. จิตฺตุปฺปาท\-กณฺฑ}
\proseitem{35}{\csromanfolio{23}กตมา ตสฺมิํ สมเย อนภิชฺฌา โหติ. โย ตสฺมิํ สมเย อโลโภ อลุพฺภนา อลุพฺภิตตฺตํ อสาราโค อสารชฺชนา อสารชฺชิตตฺตํ อนภิชฺฌา อโลโภ กุสลมูลํ. อยํ ตสฺมิํ สมเย อนภิชฺฌา โหติ.}

\proseitem{36}{กตโม ตสฺมิํ สมเย อพฺยาปาโท โหหิ. โย ตสฺมิํ สมเย อโทโส อทุสฺสนา อทุสฺสิตตฺตํ อพฺยาปาโท อพฺยาปชฺโช อโทโส กุสลมูลํ. อยํ ตสฺมิํ สมเย อพฺยาปาโท โหติ.}

\proseitem{37}{กตมา ตสฺมิํ สมเย สมฺมาทิฏฺฐิ โหติ. ยา ตสฺมิํ สมเย ปญฺญา ปชานนา วิจโย ปวิจโย ธมฺมวิจโย สลฺลกฺขณา อุปลกฺขณา ปจฺจุ\-ปลกฺขณา ปณฺฑิจฺจํ โกสลฺลํ เนปุญฺญํ เวภพฺยา จินฺตา อุปปริกฺขา ภูรี เมธา ปริณายิกา วิปสฺสนา สมฺปชญฺญํ ปโตโท ปญฺญา ปญฺญินฺทฺริยํ ปญฺญาพลํ ปญฺญาสตฺถํ ปญฺญาปาสาโท ปญฺญาอาโลโก ปญฺญาโอภาโส ปญฺญาปชฺโชโต ปญฺญารตนํ อโมโห ธมฺมวิจโย สมฺมาทิฏฺฐิ. อยํ ตสฺมิํ สมเย สมฺมาทิฏฺฐิ โหติ.}

\proseitem{38}{กตมา ตสฺมิํ สมเย หิรี โหติ. ยํ ตสฺมิํ สมเย หิรียติ หิริยิตพฺเพน หิรียติ ปาปกานํ อกุสลานํ ธมฺมานํ สมฺมาปตฺติยา. อยํ ตสฺมิํ สมเย หิรี โหติ.}

\proseitem{39}{กตมํ ตสฺมิํ สมเย โอตฺตปฺปํ โหติ. ยํ ตสฺมิํ สมเย โอตฺตปฺปติ โอตฺตปฺปิตพฺเพน โอตฺตปฺปติ ปาปกานํ อกุสลานํ ธมฺมานํ สมาปตฺติยา. อิทํ ตสฺมิํ สมเย โอตฺตปฺปํ โหติ.}

\proseitem{40}{กตมา ตสฺมิํ สมเย กายปสฺสทฺธิ โหติ. ยา ตสฺมิํ สมเย เวทนา\-กฺขนฺธสฺส สญฺญา\-กฺขนฺธสฺส สงฺขาร\-กฺขนฺธสฺส ปสฺสทฺธิ ปฏิปสฺสทฺธิ\csromansharedfootnote{6c673304f4dc7690}{ปฏิปฺปสฺสทฺธิ (สี, สฺยา)} ปสฺสมฺภนา ปฏิปสฺสมฺภนา\csromansharedfootnote{d39ceb36deddc031}{ปฏิปฺปสฺสมฺภนา (สี, สฺยา)} ปฏิปสฺสมฺภิ\-ตตฺตํ\csromansharedfootnote{02d77a2d33f3139a}{ปฏิปฺปสฺสมฺภิตตฺตํ (สี, สฺยา)}. อยํ ตสฺมิํ สมเย กายปสฺสทฺธิ โหติ.}

\proseitem{41}{กตมา ตสฺมิํ สมเย จิตฺตปสฺสทฺธิ โหติ. ยา ตสฺมิํ สมเย วิญฺญาณ\-กฺขนฺธสฺส ปสฺสทฺธิ ปฏิปสฺสทฺธิ ปสฺสมฺภนา ปฏิปสฺสมฺภนา ปฏิปสฺสมฺภิ\-ตตฺตํ. อยํ ตสฺมิํ สมเย จิตฺตปสฺสทฺธิ โหติ.}

\proseitem{42}{\csromanfolio{24}กตมา ตสฺมิํ สมเย กายลหุตา โหติ. ยา ตสฺมิํ สมเย เวทนา\-กฺขนฺธสฺส สญฺญา\-กฺขนฺธสฺส สงฺขาร\-กฺขนฺธสฺส ลหุตา ลหุปริณามตา อทนฺธนตา อวิตฺถนตา. อยํ ตสฺมิํ สมเย กายลหุตา โหติ.}

\proseitem{43}{กตมา ตสฺมิํ สมเย จิตฺตลหุตา โหติ. ยา ตสฺมิํ สมเย วิญฺญาณ\-กฺขนฺธสฺส ลหุตา ลหุปริณามตา อทนฺธนตา อวิตฺถนตา. อยํ ตสฺมิํ สมเย จิตฺตลหุตา โหติ.}

\proseitem{44}{กตมา ตสฺมิํ สมเย กายมุทุตา โหติ. ยา ตสฺมิํ สมเย เวทนา\-กฺขนฺธสฺส สญฺญา\-กฺขนฺธสฺส สงฺขาร\-กฺขนฺธสฺส มุทุตา มทฺทวตา อกกฺขฬตา อกถินตา. อยํ ตสฺมิํ สมเย กายมุทุตา โหติ.}

\proseitem{45}{กตมา ตสฺมิํ สมเย จิตฺตมุทุตา โหติ. ยา ตสฺมิํ สมเย วิญฺญาณ\-กฺขนฺธสฺส มุทุตา มทฺทวตา อกกฺขฬตา อกถินตา. อยํ ตสฺมิํ สมเย จิตฺตมุทุตา โหติ.}

\proseitem{46}{กตมา ตสฺมิํ สมเย กายกมฺมญฺญตา โหติ. ยา ตสฺมิํ สมเย เวทนา\-กฺขนฺธสฺส สญฺญา\-กฺขนฺธสฺส สงฺขาร\-กฺขนฺธสฺส กมฺมญฺญตา กมฺมญฺญตฺตํ กมฺมญฺญภาโว. อยํ ตสฺมิํ สมเย กายกมฺมญฺญตา โหติ.}

\proseitem{47}{กตมา ตสฺมิํ สมเย จิตฺต\-กมฺมญฺญตา โหติ. ยา ตสฺมิํ สมเย วิญฺญาณ\-กฺขนฺธสฺส กมฺมญฺญตา กมฺมญฺญตฺตํ กมฺมญฺญภาโว. อยํ ตสฺมิํ สมเย จิตฺต\-กมฺมญฺญตา โหติ.}

\proseitem{48}{กตมา ตสฺมิํ สมเย กายปาคุญฺญตา โหติ. ยา ตสฺมิํ สมเย เวทนา\-กฺขนฺธสฺส สญฺญา\-กฺขนฺธสฺส สงฺขาร\-กฺขนฺธสฺส ปคุณตา ปคุณตฺตํ ปคุณภาโว. อยํ ตสฺมิํ สมเย กายปาคุญฺญตา โหติ.}

\proseitem{49}{กตมา ตสฺมิํ สมเย จิตฺตปาคุญฺญตา โหติ. ยา ตสฺมิํ สมเย วิญฺญาณ\-กฺขนฺธสฺส ปคุณตา ปคุณตฺตํ ปคุณภาโว. อยํ ตสฺมิํ สมเย จิตฺตปาคุญฺญตา โหติ.}

\proseitem{50}{กตมา ตสฺมิํ สมเย กายุชุกตา โหติ. ยา ตสฺมิํ สมเย เวทนา\-กฺขนฺธสฺส สญฺญา\-กฺขนฺธสฺส สงฺขาร\-กฺขนฺธสฺส อุชุตา อุชุกตา อชิมฺหตา อวงฺกตา อกุฏิลตา. อยํ ตสฺมิํ สมเย กายุชุกตา โหติ.}

\chapterheadpageread{1. จิตฺตุปฺปาท\-กณฺฑ}
\proseitem{51}{\csromanfolio{25}กตมา ตสฺมิํ สมเย จิตฺตุชุกตา โหติ. ยา ตสฺมิํ สมเย วิญฺญาณ\-กฺขนฺธสฺส อุชุตา อุชุกตา อชิมฺหตา อวงฺกตา อกุฏิลตา. อยํ ตสฺมิํ สมเย จิตฺตุชุกตา โหติ.}

\proseitem{52}{กตมา ตสฺมิํ สมเย สติ โหติ. ยา ตสฺมิํ สมเย สติ อนุสฺสติ ปฏิสฺสติ สติ สรณตา ธารณตา อปิลาปนตา อสมฺมุสฺสนตา สติ สตินฺทฺริยํ สติพลํ สมฺมาสติ. อยํ ตสฺมิํ สมเย สติ โหติ.}

\proseitem{53}{กตมํ ตสฺมิํ สมเย สมฺปชญฺญํ โหติ. ยา ตสฺมิํ สมเย ปญฺญา ปชานนา วิจโย ปวิจโย ธมฺมวิจโย สลฺลกฺขณา อุปลกฺขณา ปจฺจุ\-ปลกฺขณา ปณฺฑิจฺจํ โกสลฺลํ เนปุญฺญํ เวภพฺยา จินฺตา อุปปริกฺขา ภูรี เมธา ปริณายิกา วิปสฺสนา สมฺปชญฺญํ ปโตโท ปญฺญา ปญฺญินฺทฺริยํ ปญฺญาพลํ ปญฺญาสตฺถํ ปญฺญาปาสาโท ปญฺญาอาโลโก ปญฺญาโอภาโส ปญฺญาปชฺโชโต ปญฺญารตนํ อโมโห ธมฺมวิจโย สมฺมาทิฏฺฐิ. อิทํ ตสฺมิํ สมเย สมฺปชญฺญํ โหติ.}

\proseitem{54}{กตโม ตสฺมิํ สมเย สมโถ โหติ. ยา ตสฺมิํ สมเย จิตฺตสฺส ฐิติ สณฺฐิติ อวฏฺฐิติ อวิสาหาโร อวิกฺเขโป อวิสาหฏ\-มานสตา สมโถ สมาธินฺทฺริยํ สมาธิพลํ สมฺมาสมาธิ. อยํ ตสฺมิํ สมเย สมโถ โหติ.}

\proseitem{55}{กตมา ตสฺมิํ สมเย วิปสฺสนา โหติ. ยา ตสฺมิํ สมเย ปญฺญา ปชานนา วิจโย ปวิจโย ธมฺมวิจโย สลฺลกฺขณา อุปลกฺขณา ปจฺจุ\-ปลกฺขณา ปณฺฑิจฺจํ โกสลฺลํ เนปุญฺญํ เวภพฺยา จินฺตา อุปปริกฺขา ภูรี เมธา ปริณายิกา วิปสฺสนา สมฺปชญฺญํ ปโตโท ปญฺญา ปญฺญินฺทฺริยํ ปญฺญาพลํ ปญฺญาสตฺถํ ปญฺญาปาสาโท ปญฺญาอาโลโก ปญฺญาโอภาโส ปญฺญาปชฺโชโต ปญฺญารตนํ อโมโห ธมฺมวิจโย สมฺมาทิฏฺฐิ. อยํ ตสฺมิํ สมเย วิปสฺสนา โหติ.}

\proseitem{56}{กตโม ตสฺมิํ สมเย ปคฺคาโห โหติ. โย ตสฺมิํ สมเย เจตสิโก วีริยารมฺโภ นิกฺกโม ปรกฺกโม อุยฺยาโม วายาโม อุสฺสาโห อุสฺโสฬฺหี ถาโม ธิติ อสิถิล\-ปรกฺกมตา \csromanfolio{26}อนิกฺขิตฺต\-ฉนฺทตา อนิกฺขิตฺตธุรตา ธุรสมฺปคฺคาโห วีริยํ วีริยินฺทฺริยํ วีริยพลํ สมฺมาวายาโม. อยํ ตสฺมิํ สมเย ปคฺคาโห โหติ.}

\proseitem{57}{กตโม ตสฺมิํ สมเย อวิกฺเขโป โหติ. ยา ตสฺมิํ สมเย จิตฺตสฺส ฐิติ สณฺฐิติ อวฏฺฐิติ อวิสาหาโร อวิกฺเขโป อวิสาหฏ\-มานสตา สมโถ สมาธินฺทฺริยํ สมาธิพลํ สมฺมาสมาธิ. อยํ ตสฺมิํ สมเย อวิกฺเขโป โหติ.}

\prosewithcloser{เย วา ปน ตสฺมิํ สมเย อญฺเญปิ อตฺถิ ปฏิจฺจ\-สมุปฺปนฺนา อรูปิโน ธมฺมา. อิเม ธมฺมา กุสลา.}{ปทภาชนียํ นิฏฺฐิตํ.}
\csromancenter{ปฐม\-ภาณวาโร.}
\csromanmatikaanchor{mtk.19}
\csromantocmarkref{section}{โกฏฺฐาสวาร}{mtk.19}
\bookmark[dest={mtk.19},level=1]{โกฏฺฐาสวาร}
\csromanheader{1}{\textbf{โกฏฺฐาสวาร}}
\proseitem{58}{ตสฺมิํ โข ปน สมเย จตฺตาโร ขนฺธา โหนฺติ, ทฺวายตนานิ โหนฺติ, เทฺว ธาตุโย โหนฺติ, ตโย อาหารา โหนฺติ, อฏฺฐินฺทฺริยานิ โหนฺติ, ปญฺจงฺคิกํ ฌานํ โหติ, ปญฺจงฺคิโก มคฺโค โหติ, สตฺต พลานิ โหนฺติ, ตโย เหตู โหนฺติ, เอโก ผสฺโส โหติ, เอกา เวทนา โหติ, เอกา สญฺญา โหติ, เอกา เจตนา โหติ, เอกํ จิตฺตํ โหติ, เอโก เวทนากฺขนฺโธ โหติ, เอโก สญฺญากฺขนฺโธ โหติ, เอโก สงฺขาร\-กฺขนฺโธ โหติ, เอโก วิญฺญาณ\-กฺขนฺโธ โหติ, เอกํ มนายตนํ โหติ, เอกํ มนินฺทฺริยํ โหติ, เอกา มโนวิญฺญาณ\-ธาตุ โหติ, เอกํ ธมฺมายตนํ โหติ, เอกา ธมฺมธาตุ โหติ. เย วา ปน ตสฺมิํ สมเย อญฺเญปิ อตฺถิ ปฏิจฺจ\-สมุปฺปนฺนา อรูปิโน ธมฺมา. อิเม ธมฺมา กุสลา.}

\proseitem{59}{กตเม ตสฺมิํ สมเย จตฺตาโร ขนฺธา โหนฺติ. เวทนากฺขนฺโธ สญฺญากฺขนฺโธ สงฺขาร\-กฺขนฺโธ วิญฺญาณ\-กฺขนฺโธ.}

\proseitem{60}{กตโม ตสฺมิํ สมเย เวทนากฺขนฺโธ โหติ. ยํ ตสฺมิํ สมเย เจตสิกํ สาตํ เจตสิกํ สุขํ เจโต\-สมฺผสฺสชํ สาตํ \csromanfolio{27}สุขํ เวทยิตํ เจโต\-สมฺผสฺสชา สาตา สุขา เวทนา. อยํ ตสฺมิํ สมเย เวทนากฺขนฺโธ โหติ.}

\proseitem{61}{กตโม ตสฺมิํ สมเย สญฺญากฺขนฺโธ โหติ. ยา ตสฺมิํ สมเย สญฺญา สญฺชานนา สญฺชานิตตฺตํ. อยํ ตสฺมิํ สมเย สญฺญากฺขนฺโธ โหติ.}

\proseitem{62}{กตโม ตสฺมิํ สมเย สงฺขาร\-กฺขนฺโธ โหติ. ผสฺโส เจตนา วิตกฺโก วิจาโร ปีติ จิตฺตสฺเส\-กคฺคตา สทฺธินฺทฺริยํ วีริยินฺทฺริยํ สตินฺทฺริยํ สมาธินฺทฺริยํ ปญฺญินฺทฺริยํ ชีวิตินฺทฺริยํ สมฺมาทิฏฺฐิ สมฺมาสงฺกปฺโป สมฺมาวายาโม สมฺมาสติ สมฺมาสมาธิ สทฺธาพลํ วีริยพลํ สติพลํ สมาธิพลํ ปญฺญาพลํ หิริพลํ โอตฺตปฺปพลํ อโลโภ อโทโส อโมโห อนภิชฺฌา อพฺยาปาโท สมฺมาทิฏฺฐิ หิรี โอตฺตปฺปํ กายปสฺสทฺธิ จิตฺตปสฺสทฺธิ กายลหุตา จิตฺตลหุตา กายมุทุตา จิตฺตมุทุตา กายกมฺมญฺญตา จิตฺต\-กมฺมญฺญตา กายปาคุญฺญตา จิตฺตปาคุญฺญตา กายุชุกตา จิตฺตุชุกตา สติ สมฺปชญฺญํ สมโถ วิปสฺสนา ปคฺคาโห อวิกฺเขโป. เย วา ปน ตสฺมิํ สมเย อญฺเญปิ อตฺถิ ปฏิจฺจ\-สมุปฺปนฺนา อรูปิโน ธมฺมา ฐเปตฺวา\csromansharedfootnote{81cb95bd9ce77774}{ถเปตฺวา (ก)} เวทนา\-กฺขนฺธํ ฐเปตฺวา สญฺญากฺขนฺธํ ฐเปตฺวา วิญฺญาณ\-กฺขนฺธํ. อยํ ตสฺมิํ สมเย สงฺขาร\-กฺขนฺโธ โหติ.}

\proseitem{63}{กตโม ตสฺมิํ สมเย วิญฺญาณ\-กฺขนฺโธ โหติ. ยํ ตสฺมิํ สมเย จิตฺตํ มโน มานสํ หทยํ ปณฺฑรํ มโน มนายตนํ มนินฺทฺริยํ วิญฺญาณํ วิญฺญาณ\-กฺขนฺโธ ตชฺชา\-มโนวิญฺญาณ\-ธาตุ. อยํ ตสฺมิํ สมเย วิญฺญาณ\-กฺขนฺโธ โหติ. อิเม ตสฺมิํ สมเย จตฺตาโร ขนฺธา โหนฺติ.}

\proseitem{64}{กตมานิ ตสฺมิํ สมเย ทฺวายตนานิ โหนฺติ. มนายตนํ ธมฺมายตนํ.}

\proseitem{65}{กตมํ ตสฺมิํ สมเย มนายตนํ โหติ. ยํ ตสฺมิํ สมเย จิตฺตํ มโน มานสํ หทยํ ปณฺฑรํ มโน มนายตนํ มนินฺทฺริยํ วิญฺญาณํ วิญฺญาณ\-กฺขนฺโธ ตชฺชา\-มโนวิญฺญาณ\-ธาตุ. อิทํ ตสฺมิํ สมเย มนายตนํ โหติ.}

\proseitem{66}{\csromanfolio{28}กตมํ ตสฺมิํ สมเย ธมฺมายตนํ โหติ. เวทนากฺขนฺโธ สญฺญากฺขนฺโธ สงฺขาร\-กฺขนฺโธ. อิทํ ตสฺมิํ สมเย ธมฺมายตนํ โหติ. อิมานิ ตสฺมิํ สมเย ทฺวายตนานิ โหนฺติ.}

\proseitem{67}{กตมา ตสฺมิํ สมเย เทฺว ธาตุโย โหนฺติ. มโนวิญฺญาณ\-ธาตุ ธมฺมธาตุ.}

\proseitem{68}{กตมา ตสฺมิํ สมเย มโนวิญฺญาณ\-ธาตุ โหติ. ยํ ตสฺมิํ สมเย จิตฺตํ มโน มานสํ หทยํ ปณฺฑรํ มโน มนายตนํ มนินฺทฺริยํ วิญฺญาณํ วิญฺญาณ\-กฺขนฺโธ ตชฺชา\-มโนวิญฺญาณ\-ธาตุ. อยํ ตสฺมิํ สมเย มโนวิญฺญาณ\-ธาตุ โหติ.}

\proseitem{69}{กตมา ตสฺมิํ สมเย ธมฺมธาตุ โหติ. เวทนากฺขนฺโธ สญฺญากฺขนฺโธ สงฺขาร\-กฺขนฺโธ. อยํ ตสฺมิํ สมเย ธมฺมธาตุ โหติ. อิมา ตสฺมิํ สมเย เทฺว ธาตุโย โหนฺติ.}

\proseitem{70}{กตเม ตสฺมิํ สมเย ตโย อาหารา โหนฺติ. ผสฺสหาโร มโนสญฺเจตนา\-หาโร วิญฺญาณาหาโร.}

\proseitem{71}{กตโม ตสฺมิํ สมเย ผสฺสาหาโร โหติ. โย ตสฺมิํ สมเย ผสฺโส ผุสนา สมฺผุสนา สมฺผุสิตตฺตํ. อยํ ตสฺมิํ สมเย ผสฺสาหาโร โหติ.}

\proseitem{72}{กตโม ตสฺมิํ สมเย มโนสญฺเจตนา\-หาโร โหติ. ยา ตสฺมิํ สมเย เจตนา สญฺเจตนา เจตยิตตฺตํ. อยํ ตสฺมิํ สมเย มโนสญฺเจตนา\-หาโร โหติ.}

\proseitem{73}{กตโม ตสฺมิํ สมเย วิญฺญาณาหาโร โหติ. ยํ ตสฺมิํ สมเย จิตฺตํ มโน มานสํ หทยํ ปณฺฑรํ มโน มนายตนํ มนินฺทฺริยํ วิญฺญาณํ วิญฺญาณ\-กฺขนฺโธ ตชฺชา\-มโนวิญฺญาณ\-ธาตุ. อยํ ตสฺมิํ สมเย วิญฺญาณาหาโร โหติ. อิเม ตสฺมิํ สมเย ตโย อาหารา โหนฺติ.}

\chapterheadpageread{1. จิตฺตุปฺปาท\-กณฺฑ}
\proseitem{74}{\csromanfolio{29}กตมานิ ตสฺมิํ สมเย อฏฺฐินฺทฺริยานิ โหนฺติ. สทฺธินฺทฺริยํ วีริยินฺทฺริยํ สตินฺทฺริยํ สมาธินฺทฺริยํ ปญฺญินฺทฺริยํ มนินฺทฺริยํ โสมนสฺสิ\-นฺทฺริยํ ชีวิตินฺทฺริยํ.}

\proseitem{75}{กตมํ ตสฺมิํ สมเย สทฺธินฺทฺริยํ โหติ. ยา ตสฺมิํ สมเย สทฺธา สทฺทหนา โอกปฺปนา อภิปฺปสาโท สทฺธา สทฺธินฺทฺริยํ สทฺธาพลํ. อิทํ ตสฺมิํ สมเย สทฺธินฺทฺริยํ โหติ.}

\proseitem{76}{กตมํ ตสฺมิํ สมเย วีริยินฺทฺริยํ โหติ. โย ตสฺมิํ สมเย เจตสิโก วีริยารมฺโภ นิกฺกโม ปรกฺกโม อุยฺยาโม วายาโม อุสฺสาโห อุสฺโสฬฺหี ถาโม ธิติ อสิถิล\-ปรกฺกมตา อนิกฺขิตฺต\-ฉนฺทตา อนิกฺขิตฺตธุรตา ธุรสมฺปคฺคาโห วีริยํ วีริยินฺทฺริยํ วีริยพลํ สมฺมาวายาโม. อิทํ ตสฺมิํ สมเย วีริยินฺทฺริยํ โหติ.}

\proseitem{77}{กตมํ ตสฺมิํ สมเย สตินฺทฺริยํ โหติ. ยา ตสฺมิํ สมเย สติ อนุสฺสติ ปฏิสฺสติ สติ สรณตา ธารณตา อปิลาปนตา อสมฺมุสฺสนตา สติ สตินฺทฺริยํ สติพลํ สมฺมาสติ. อิทํ ตสฺมิํ สมเย สตินฺทฺริยํ โหติ.}

\proseitem{78}{กตมํ ตสฺมิํ สมเย สมาธินฺทฺริยํ โหติ. ยา ตสฺมิํ สมเย จิตฺตสฺส ฐิติ สณฺฐิติ อวฏฺฐิติ อวิสาหาโร อวิกฺเขโป อวิสาหฏ\-มานสตา สมโถ สมาธินฺทฺริยํ สมาธิพลํ สมฺมาสมาธิ. อิทํ ตสฺมิํ สมเย สมาธินฺทฺริยํ โหติ.}

\proseitem{79}{กตมํ ตสฺมิํ สมเย ปญฺญินฺทฺริยํ โหติ. ยา ตสฺมิํ สมเย ปญฺญา ปชานนา วิจโย ปวิจโย ธมฺมวิจโย สลฺลกฺขณา อุปลกฺขณา ปจฺจุ\-ปลกฺขณา ปณฺฑิจฺจํ โกสลฺลํ เนปุญฺญํ เวภพฺยา จินฺตา อุปปริกฺขา ภูรี เมธา ปริณายิกา วิปสฺสนา สมฺปชญฺญํ ปโตโท ปญฺญา ปญฺญินฺทฺริยํ ปญฺญาพลํ ปญฺญาสตฺถํ ปญฺญาปาสาโท ปญฺญาอาโลโก ปญฺญาโอภาโส ปญฺญาปชฺโชโต ปญฺญารตนํ อโมโห ธมฺมวิจโย สมฺมาทิฏฺฐิ. อิทํ ตสฺมิํ สมเย ปญฺญินฺทฺริยํ โหติ.}

\proseitem{80}{กตมํ ตสฺมิํ สมเย มนินฺทฺริยํ โหติ. ยํ ตสฺมิํ สมเย จิตฺตํ มโน มานสํ หทยํ ปณฺฑรํ มโน มนายตนํ มนินฺทฺริยํ วิญฺญาณํ วิญฺญาณ\-กฺขนฺโธ ตชฺชา\-มโนวิญฺญาณ\-ธาตุ. อิทํ ตสฺมิํ สมเย มนินฺทฺริยํ โหติ.}

\proseitem{81}{\csromanfolio{30}กตมํ ตสฺมิํ สมเย โสมนสฺสิ\-นฺทฺริยํ โหติ. ยํ ตสฺมิํ สมเย เจตสิกํ สาตํ เจตสิกํ สุขํ เจโต\-สมฺผสฺสชํ สาตํ สุขํ เวทยิตํ เจโต\-สมฺผสฺสชา สาตา สุขา เวทนา. อิทํ ตสฺมิํ สมเย โสมนสฺสิ\-นฺทฺริยํ โหติ.}

\proseitem{82}{กตมํ ตสฺมิํ สมเย ชีวิตินฺทฺริยํ โหติ. โย เตสํ อรูปีนํ ธมฺมานํ อายุ ฐิติ ยปนา ยาปนา อิริยนา วตฺตนา ปาลนา ชีวิตํ ชีวิตินฺทฺริยํ. อิทํ ตสฺมิํ สมเย ชีวิตินฺทฺริยํ โหติ. อิมานิ ตสฺมิํ สมเย อฏฺฐินฺทฺริยานิ โหนฺติ.}

\proseitem{83}{กตมํ ตสฺมิํ สมเย ปญฺจงฺคิกํ ฌานํ โหติ. วิตกฺโก วิจาโร ปีติ สุขํ จิตฺตสฺเส\-กคฺคตา.}

\proseitem{84}{กตโม ตสฺมิํ สมเย วิตกฺโก โหติ. โย ตสฺมิํ สมเย ตกฺโก วิตกฺโก สงฺกปฺโป อปฺปนา พฺยปฺปนา เจตโส อภินิโรปนา สมาสงฺกปฺโป. อยํ ตสฺมิํ สมเย วิตกฺโก โหติ.}

\proseitem{85}{กตโม ตสฺมิํ สมเย วิจาโร โหติ. โย ตสฺมิํ สมเย จาโร วิจาโร อนุวิจาโร อุปวิจาโร จิตฺตสฺส อนุสนฺธานตา อนุเปกฺขนตา. อยํ ตสฺมิํ สมเย วิจาโร โหติ.}

\proseitem{86}{กตมา ตสฺมิํ สมเย ปีติ โหติ. ยา ตสฺมิํ สมเย ปีติ ปาโมชฺชํ อาโมทนา ปโมทนา หาโส ปหาโส วิตฺติ โอทคฺยํ อตฺตมนตา จิตฺตสฺส. อยํ ตสฺมิํ สมเย ปีติ โหติ.}

\proseitem{87}{กตมํ ตสฺมิํ สมเย สุขํ โหติ. ยํ ตสฺมิํ สมเย เจตสิกํ สาตํ เจตสิกํ สุขํ เจโต\-สมฺผสฺสชํ สาตํ สุขํ เวทยิตํ เจโต\-สมฺผสฺสชา สาตา สุขา เวทนา. อิทํ ตสฺมิํ สมเย สุขํ โหติ.}

\proseitem{88}{กตมา ตสฺมิํ สมเย จิตฺตสฺเส\-กคฺคตา โหติ. ยา ตสฺมิํ สมเย จิตฺตสฺส ฐิติ สณฺฐิติ อวฏฺฐิติ อวิสาหาโร อวิกฺเขโป อวิสาหฏ\-มานสตา สมโถ สมาธินฺทฺริยํ สมาธิพลํ สมฺมาสมาธิ. อยํ ตสฺมิํ สมเย จิตฺตสฺเส\-กคฺคตา โหติ. อิทํ ตสฺมิํ สมเย ปญฺจงฺคิกํ ฌานํ โหติ.}

\chapterheadpageread{1. จิตฺตุปฺปาท\-กณฺฑ}
\proseitem{89}{\csromanfolio{31}กตโม ตสฺมิํ สมเย ปญฺจงฺคิโก มคฺโค โหติ. สมฺมาทิฏฺฐิ สมฺมาสงฺกปฺโป สมฺมาวายาโม สมฺมาสติ สมฺมาสมาธิ.}

\proseitem{90}{กตมา ตสฺมิํ สมเย สมฺมาทิฏฺฐิ โหติ. ยา ตสฺมิํ สมเย ปญฺญา ปชานนา วิจโย ปวิจโย ธมฺมวิจโย สลฺลกฺขณา อุปลกฺขณา ปจฺจุ\-ปลกฺขณา ปณฺฑิจฺจํ โกสลฺลํ เนปุญฺญํ เวภพฺยา จินฺตา อุปปริกฺขา ภูรี เมธา ปริณายิกา วิปสฺสนา สมฺปชญฺญํ ปโตโท ปญฺญา ปญฺญินฺทฺริยํ ปญฺญาพลํ ปญฺญาสตฺถํ ปญฺญาปาสาโท ปญฺญาอาโลโก ปญฺญาโอภาโส ปญฺญาปชฺโชโต ปญฺญารตนํ อโมโห ธมฺมวิจโย สมฺมาทิฏฺฐิ. อยํ ตสฺมิํ สมเย สมฺมาทิฏฺฐิ โหติ.}

\proseitem{91}{กตโม ตสฺมิํ สมเย สมฺมาสงฺกปฺโป โหติ. โย ตสฺมิํ สมเย ตกฺโก วิตกฺโก สงฺกปฺโป อปฺปนา พฺยปฺปนา เจตโส อภินิโรปนา สมฺมาสงฺกปฺโป. อยํ ตสฺมิํ สมเย สมฺมาสงฺกปฺโป โหติ.}

\proseitem{92}{กตโม ตสฺมิํ สมเย สมฺมาวายาโม โหติ. โย ตสฺมิํ สมเย เจตสิโก วีริยารมฺโภ นิกฺกโม ปรกฺกโม อุยฺยาโม วายาโม อุสฺสาโห อุสฺโสฬฺหี ถาโม ธิติ อสิถิล\-ปรกฺกมตา อนิกฺขิตฺต\-ฉนฺทตา อนิกฺขิตฺตธุรตา ธุรสมฺปคฺคาโห วีริยํ วีริยินฺทฺริยํ วีริยพลํ สมฺมาวายาโม. อยํ ตสฺมิํ สมเย สมฺมาวายาโม โหติ.}

\proseitem{93}{กตมา ตสฺมิํ สมเย สมฺมาสติ โหติ. ยา ตสฺมิํ สมเย สติ อนุสฺสติ ปฏิสฺสติ สติ สรณตา ธารณตา อปิลาปนตา อสมฺมุสฺสนตา สติ สตินฺทฺริยํ สติพลํ สมฺมาสติ. อยํ ตสฺมิํ สมเย สมฺมาสติ โหติ.}

\proseitem{94}{กตโม ตสฺมิํ สมเย สมฺมาสมาธิ โหติ. ยา ตสฺมิํ สมเย จิตฺตสฺส ฐิติ สณฺฐิติ อวฏฺฐิติ อวิสาหาโร อวิกฺเขโป อวิสาหฏ\-มานสตา สมโถ สมาธินฺทฺริยํ สมาธิพลํ สมฺมาสมาธิ. อยํ ตสฺมิํ สมเย สมฺมาสมาธิ โหติ. อยํ ตสฺมิํ สมเย ปญฺจงฺคิโก มคฺโค โหติ.}

\proseitem{95}{\csromanfolio{32}กตมานิ ตสฺมิํ สมเย สตฺต พลานิ โหนฺติ. สทฺธาพลํ วีริยพลํ สติพลํ สมาธิพลํ ปญฺญาพลํ หิริพลํ โอตฺตปฺปพลํ.}

\proseitem{96}{กตมํ ตสฺมิํ สมเย สทฺธาพลํ โหติ. ยา ตสฺมิํ สมเย สทฺธา สทฺทหนา โอกปฺปนา อภิปฺปสาโท สทฺธา สทฺธินฺทฺริยํ สทฺธาพลํ. อิทํ ตสฺมิํ สมเย สทฺธาพลํ โหติ.}

\proseitem{97}{กตมํ ตสฺมิํ สมเย วีริยพลํ โหติ. โย ตสฺมิํ สมเย เจตสิโก วีริยารมฺโภ นิกฺกโม ปรกฺกโม อุยฺยาโม วายาโม อุสฺสาโห อุสฺโสฬฺหี ถาโม ธิติ อสิถิล\-ปรกฺกมตา อนิกฺขิตฺต\-ฉนฺทตา อนิกฺขิตฺตธุรตา ธุรสมฺปคฺคาโห วีริยํ วีริยินฺทฺริยํ วีริยพลํ สมฺมาวายาโม. อิทํ ตสฺมิํ สมเย วีริยพลํ โหติ.}

\proseitem{98}{กตมํ ตสฺมิํ สมเย สติพลํ โหติ. ยา ตสฺมิํ สมเย สติ อนุสฺสติ ปฏิสฺสติ สติ สรณตา ธารณตา อปิลาปนตา อสมฺมุสฺสนตา สติ สตินฺทฺริยํ สติพลํ สมฺมาสติ. อิทํ ตสฺมิํ สมเย สติพลํ โหติ.}

\proseitem{99}{กตมํ ตสฺมิํ สมเย สมาธิพลํ โหติ. ยา ตสฺมิํ สมเย จิตฺตสฺส ฐิติ สณฺฐิติ อวฏฺฐิติ อวิสาหาโร อวิกฺเขโป อวิสาหฏ\-มานสตา สมโถ สมาธินฺทฺริยํ สมาธิพลํ สมฺมาสมาธิ. อิทํ ตสฺมิํ สมเย สมาธิพลํ โหติ.}

\proseitem{100}{กตมํ ตสฺมิํ สมเย ปญฺญาพลํ โหติ. ยา ตสฺมิํ สมเย ปญฺญา ปชานนา วิจโย ปวิจโย ธมฺมวิจโย สลฺลกฺขณา อุปลกฺขณา ปจฺจุ\-ปลกฺขณา ปณฺฑิจฺจํ โกสลฺลํ เนปุญฺญํ เวภพฺยา จินฺตา อุปปริกฺขา ภูรี เมธา ปริณายิกา วิปสฺสนา สมฺปชญฺญํ ปโตโท ปญฺญา ปญฺญินฺทฺริยํ ปญฺญาพลํ ปญฺญาสตฺถํ ปญฺญาปาสาโท ปญฺญาอาโลโก ปญฺญาโอภาโส ปญฺญาปชฺโชโต ปญฺญารตนํ อโมโห ธมฺมวิจโย สมฺมาทิฏฺฐิ. อิทํ ตสฺมิํ สมเย ปญฺญาพลํ โหติ.}

\proseitem{101}{กตมํ ตสฺมิํ สมเย หิริพลํ โหติ. ยํ ตสฺมิํ สมเย หิรียติ หิริยิตพฺเพน หิรียติ ปาปกานํ อกุสลานํ ธมฺมานํ สมาปตฺติยา. อิทํ ตสฺมิํ สมเย หิริพลํ โหติ.}

\chapterheadpageread{1. จิตฺตุปฺปาท\-กณฺฑ}
\proseitem{102}{\csromanfolio{33}กตมํ ตสฺมิํ สมเย โอตฺตปฺปพลํ โหติ. ยํ ตสฺมิํ สมเย โอตฺตปฺปติ โอตฺตปฺปิตพฺเพน โอตฺตปฺปติ ปาปกานํ อกุสลานํ ธมฺมานํ สมาปตฺติยา. อิทํ ตสฺมิํ สมเย โอตฺตปฺปพลํ โหติ. อิมานิ ตสฺมิํ สมเย สตฺต พลานิ โหนฺติ.}

\proseitem{103}{กตเม ตสฺมิํ สมเย ตโย เหตู โหนฺติ. อโลโภ อโทโส อโมโห.}

\proseitem{104}{กตโม ตสฺมิํ สมเย อโลโภ โหติ. โย ตสฺมิํ สมเย อโลโภ อลุพฺภนา อลุพฺภิตตฺตํ อสาราโค อสารชฺชนา อสารชฺชิตตฺตํ อนภิชฺฌา อโลโภ กุสลมูลํ. อยํ ตสฺมิํ สมเย อโลโภ โหติ.}

\proseitem{105}{กตโม ตสฺมิํ สมเย อโทโส โหติ. โย ตสฺมิํ สมเย อโทโส อทุสฺสนา อทุสฺสิตตฺตํ อพฺยาปาโท อพฺยาปชฺโช อโทโส กุสลมูลํ. อยํ ตสฺมิํ สมเย อโทโส โหติ.}

\proseitem{106}{กตโม ตสฺมิํ สมเย อโมโห โหติ. ยา ตสฺมิํ สมเย ปญฺญา ปชานนา ฯเปฯ อโมโห ธมฺมวิจโย สมฺมาทิฏฺฐิ. อยํ ตสฺมิํ สมเย อโมโห โหติ. อิมํ ตสฺมิํ สมเย ตโย เหตู โหนฺติ.}

\proseitem{107}{กตโม ตสฺมิํ สมเย เอโก ผสฺโส โหติ. โย ตสฺมิํ สมเย ผสฺโส ผุสนา สมฺผุสนา สมฺผุสิตตฺตํ. อยํ ตสฺมิํ สมเย เอโก ผสฺโส โหติ.}

\proseitem{108}{กตมา ตสฺมิํ สมเย เอกา เวทนา โหติ. ยํ ตสฺมิํ สมเย เจตสิกํ สาตํ เจตสิกํ สุขํ เจโต\-สมฺผสฺสชํ สาตํ สุขํ เวทยิตํ เจโต\-สมฺผสฺสชา สาตา สุขา เวทนา. อยํ ตสฺมิํ สมเย เอกา เวทนา โหติ.}

\proseitem{109}{กตมา ตสฺมิํ สมเย เอกา สญฺญา โหติ. ยา ตสฺมิํ สมเย สญฺญา สญฺชานนา สญฺชานิตตฺตํ. อยํ ตสฺมิํ สมเย เอกา สญฺญา โหติ.}

\proseitem{110}{\csromanfolio{34}กตมา ตสฺมิํ สมเย เอกา เจตนา โหติ. ยา ตสฺมิํ สมเย เจตนา สญฺเจตนา เจตยิตตฺตํ. อยํ ตสฺมิํ สมเย เอกา เจตนา โหติ.}

\proseitem{111}{กตมํ ตสฺมิํ สมเย เอกํ จิตฺตํ โหติ. ยํ ตสฺมิํ สมเย จิตฺตํ มโน มานสํ หทยํ ปณฺฑรํ มโน มนายตนํ มนินฺทฺริยํ วิญฺญาณํ วิญฺญาณ\-กฺขนฺโธ ตชฺชา\-มโนวิญฺญาณ\-ธาตุ. อิทํ ตสฺมิํ สมเย เอกํ จิตฺตํ โหติ.}

\proseitem{112}{กตโม ตสฺมิํ สมเย เอโก เวทนากฺขนฺโธ โหติ. ยํ ตสฺมิํ สมเย เจตสิกํ สาตํ เจตสิกํ สุขํ เจโต\-สมฺผสฺสชํ สาตํ สุขํ เวทยิตํ เจโต\-สมฺผสฺสชา สาตา สุขา เวทนา. อยํ ตสฺมิํ สมเย เอโก เวทนากฺขนฺโธ โหติ.}

\proseitem{113}{กตโม ตสฺมิํ สมเย เอโก สญฺญากฺขนฺโธ โหติ. ยา ตสฺมิํ สมเย สญฺญา สญฺชานนา สญฺชานิตตฺตํ. อยํ ตสฺมิํ สมเย เอโก สญฺญากฺขนฺโธ โหติ.}

\proseitem{114}{กตโม ตสฺมิํ สมเย เอโก สงฺขาร\-กฺขนฺโธ โหติ. ผสฺโส เจตนา วิตกฺโก วิจาโร ปีติ จิตฺตสฺเส\-กคฺคตา สทฺธินฺทฺริยํ วีริยินฺทฺริยํ สตินฺทฺริยํ สมาธินฺทฺริยํ ปญฺญินฺทฺริยํ ชีวิตินฺทฺริยํ สมฺมาทิฏฺฐิ สมฺมาสงฺกปฺโป สมฺมาวายาโม สมฺมาสติ สมฺมาสมาธิ สทฺธาพลํ วีริยพลํ สติพลํ สมาธิพลํ ปญฺญาพลํ หิริพลํ โอตฺตปฺปพลํ อโลโภ อโทโส อโมโห อนภิชฺฌา อพฺยาปาโท สมฺมาทิฏฺฐิ หิรี โอตฺตปฺปํ กายปสฺสทฺธิ จิตฺตปสฺสทฺธิ กายลหุตา จิตฺตลหุตา กายมุทุตา จิตฺตมุทุตา กายกมฺมญฺญตา จิตฺต\-กมฺมญฺญตา กายปาคุญฺญตา จิตฺตปาคุญฺญตา กายุชุกตา จิตฺตุชุกตา สติ สมฺปชญฺญํ สมโถ วิปสฺสนา ปคฺคาโห อวิกฺเขโป. เย วา ปน ตสฺมิํ สมเย อญฺเญปิ อตฺถิ ปฏิจฺจ\-สมุปฺปนฺนา อรูปิโน ธมฺมา ฐเปตฺวา เวทนา\-กฺขนฺธํ ฐเปตฺวา สญฺญากฺขนฺธํ ฐเปตฺวา วิญฺญาณ\-กฺขนฺธํ. อยํ ตสฺมิํ สมเย เอโก สงฺขาร\-กฺขนฺโธ โหติ.}

\proseitem{115}{กตโม ตสฺมิํ สมเย เอโก วิญฺญาณ\-กฺขนฺโธ โหติ. ยํ ตสฺมิํ สมเย จิตฺตํ มโน มานสํ หทยํ ปณฺฑรํ มโน มนายตนํ \csromanfolio{35}มนินฺทฺริยํ วิญฺญาณํ วิญฺญาณ\-กฺขนฺโธ ตชฺชา\-มโนวิญฺญาณ\-ธาตุ. อยํ ตสฺมิํ สมเย เอโก วิญฺญาณ\-กฺขนฺโธ โหติ.}

\proseitem{116}{กตมํ ตสฺมิํ สมเย เอกํ มนายตนํ โหติ. ยํ ตสฺมิํ สมเย จิตฺตํ มโน มานสํ หทยํ ปณฺฑรํ มโน มนายตนํ มนินฺทฺริยํ วิญฺญาณํ วิญฺญาณ\-กฺขนฺโธ ตชฺชา\-มโนวิญฺญาณ\-ธาตุ. อิทํ ตสฺมิํ สมเย เอกํ มนายตนํ โหติ.}

\proseitem{117}{กตมํ ตสฺมิํ สมเย เอกํ มนินฺทฺริยํ โหติ. ยํ ตสฺมิํ สมเย จิตฺตํ มโน มานสํ หทยํ ปณฺฑรํ มโน มนายตนํ มนินฺทฺริยํ วิญฺญาณํ วิญฺญาณ\-กฺขนฺโธ ตชฺชา\-มโนวิญฺญาณ\-ธาตุ. อิทํ ตสฺมิํ สมเย เอกํ มนินฺทฺริยํ โหติ.}

\proseitem{118}{กตมา ตสฺมิํ สมเย เอกา มโนวิญฺญาณ\-ธาตุ โหติ. ยํ ตสฺมิํ สมเย จิตฺตํ มโน มานสํ หทยํ ปณฺฑรํ มโน มนายตนํ มนินฺทฺริยํ วิญฺญาณํ วิญฺญาณ\-กฺขนฺโธ ตชฺชา\-มโนวิญฺญาณ\-ธาตุ. อยํ ตสฺมิํ สมเย เอกา มโนวิญฺญาณ\-ธาตุ โหติ.}

\proseitem{119}{กตมํ ตสฺมิํ สมเย เอกํ ธมฺมายตนํ โหติ. เวทนากฺขนฺโธ สญฺญากฺขนฺโธ สงฺขาร\-กฺขนฺโธ. อิทํ ตสฺมิํ สมเย เอกํ ธมฺมายตนํ โหติ.}

\proseitem{120}{กตมา ตสฺมิํ สมเย เอกา ธมฺมธาตุ โหติ. เวทนากฺขนฺโธ สญฺญากฺขนฺโธ สงฺขาร\-กฺขนฺโธ. อยํ ตสฺมิํ สมเย เอกา ธมฺมธาตุ โหติ.}

\prose{เย วา ปน ตสฺมิํ สมเย อญฺเญปิ อตฺถิ ปฏิจฺจ\-สมุปฺปนฺนา อรูปิโน ธมฺมา. อิเม ธมฺมา กุสลา.}

\vspace{\tipitakaparskip}
\csromancenter{โกฏฺฐาสวาโร.}
\csromanmatikaanchor{mtk.20}
\csromantocmarkref{section}{สุญฺญตวาร}{mtk.20}
\bookmark[dest={mtk.20},level=1]{สุญฺญตวาร}
\csromanheader{1}{\textbf{สุญฺญตวาร}}
\proseitem{121}{ตสฺมิํ โข ปน สมเย ธมฺมา โหนฺติ, ขนฺธา โหนฺติ, อายตนานิ โหนฺติ, ธาตุโย โหนฺติ, อาหารา โหนฺติ, \csromanfolio{36}อินฺทฺริยานิ โหนฺติ, ฌานํ โหติ, มคฺโค โหติ, พลานิ โหนฺติ, เหตู โหนฺติ, ผสฺโส โหติ, เวทนา โหติ, สญฺญา โหติ, เจตนา โหติ, จิตฺตํ โหติ, เวทนากฺขนฺโธ โหติ, สญฺญากฺขนฺโธ โหติ, สงฺขาร\-กฺขนฺโธ โหติ, วิญฺญาณ\-กฺขนฺโธ โหติ, มนายตนํ โหติ, มนินฺทฺริยํ โหติ, มโนวิญฺญาณ\-ธาตุ โหติ, ธมฺมายตนํ โหติ, ธมฺมธาตุ โหติ. เย วา ปน ตสฺมิํ สมเย อญฺเญปิ อตฺถิ ปฏิจฺจ\-สมุปฺปนฺนา อรูปิโน ธมฺมา. อิเม ธมฺมา กุสลา.}

\proseitem{122}{กตเม ตสฺมิํ สมเย ธมฺมา โหนฺติ. เวทนากฺขนฺโธ สญฺญากฺขนฺโธ สงฺขาร\-กฺขนฺโธ วิญฺญาณ\-กฺขนฺโธ. อิเม ตสฺมิํ สมเย ธมฺมา โหนฺติ.}

\proseitem{123}{กตเม ตสฺมิํ สมเย ขนฺธา โหนฺติ. เวทนากฺขนฺโธ สญฺญากฺขนฺโธ สงฺขาร\-กฺขนฺโธ วิญฺญาณ\-กฺขนฺโธ. อิเม ตสฺมิํ สมเย ขนฺธา โหนฺติ.}

\proseitem{124}{กตมานิ ตสฺมิํ สมเย อายตนานิ โหนฺติ. มนายตนํ ธมฺมายตนํ. อิมานิ ตสฺมิํ สมเย อายตนานิ โหนฺติ.}

\proseitem{125}{กตมา ตสฺมิํ สมเย ธาตุโย โหนฺติ. มโนวิญฺญาณ\-ธาตุ ธมฺมธาตุ. อิมา ตสฺมิํ สมเย ธาตุโย โหนฺติ.}

\proseitem{126}{กตเม ตสฺมิํ สมเย อาหารา โหนฺติ. ผสฺสาหาโร มโนสญฺเจตนา\-หาโร วิญฺญาณาหาโร. อิเม ตสฺมิํ สมเย อาหารา โหนฺติ.}

\proseitem{127}{กตมานิ ตสฺมิํ สมเย อินฺทฺริยานิ โหนฺติ. สทฺธินฺทฺริยํ วีริยินฺทฺริยํ สตินฺทฺริยํ สมาธินฺทฺริยํ ปญฺญินฺทฺริยํ มนินฺทฺริยํ โสมนสฺสิ\-นฺทฺริยํ ชีวิตินฺทฺริยํ. อิมานิ ตสฺมิํ สมเย อินฺทฺริยานิ โหนฺติ.}

\proseitem{128}{กตมํ ตสฺมิํ สมเย ฌานํ โหติ. วิตกฺโก วิจาโร ปีติ สุขํ จิตฺตสฺเส\-กคฺคตา. อิทํ ตสฺมิํ สมเย ฌานํ โหติ.}

\proseitem{129}{กตโม ตสฺมิํ สมเย มคฺโค โหติ. สมฺมาทิฏฺฐิ สมฺมาสงฺกปฺโป สมฺมาวายาโม สมฺมาสติ สมฺมาสมาธิ. อยํ ตสฺมิํ สมเย มคฺโค โหติ.}

\chapterheadpageread{1. จิตฺตุปฺปาท\-กณฺฑ}
\proseitem{130}{\csromanfolio{37}กตมานิ ตสฺมิํ สมเย พลานิ โหนฺติ. สทฺธาพลํ วีริยพลํ สติพลํ สมาธิพลํ ปญฺญาพลํ หิริพลํ โอตฺตปฺปพลํ. อิมานิ ตสฺมิํ สมเย พลานิ โหนฺติ.}

\proseitem{131}{กตเม ตสฺมิํ สมเย เหตู โหนฺติ. อโลโภ อโทโส อโมโห. อิเม ตสฺมิํ สมเย เหตู โหนฺติ.}

\proseitem{132}{กตโม ตสฺมิํ สมเย ผสฺโส โหติ ฯเปฯ. อยํ ตสฺมิํ สมเย ผสฺโส โหติ.}

\proseitem{133}{กตมา ตสฺมิํ สมเย เวทนา โหติ ฯเปฯ. อยํ ตสฺมิํ สมเย เวทนา โหติ.}

\proseitem{134}{กตมา ตสฺมิํ สมเย สญฺญา โหติ ฯเปฯ. อยํ ตสฺมิํ สมเย สญฺญา โหติ.}

\proseitem{135}{กตมา ตสฺมิํ สมเย เจตนา โหติ ฯเปฯ. อยํ ตสฺมิํ สมเย เจตนา โหติ.}

\proseitem{136}{กตมํ ตสฺมิํ สมเย จิตฺตํ โหติ ฯเปฯ. อิทํ ตสฺมิํ สมเย จิตฺตํ โหติ.}

\proseitem{137}{กตโม ตสฺมิํ สมเย เวทนากฺขนฺโธ โหติ ฯเปฯ. อยํ ตสฺมิํ สมเย เวทนากฺขนฺโธ โหติ.}

\proseitem{138}{กตโม ตสฺมิํ สมเย สญฺญากฺขนฺโธ โหติ ฯเปฯ. อยํ ตสฺมิํ สมเย สญฺญากฺขนฺโธ โหติ.}

\proseitem{139}{กตโม ตสฺมิํ สมเย สงฺขาร\-กฺขนฺโธ โหติ ฯเปฯ. อยํ ตสฺมิํ สมเย สงฺขาร\-กฺขนฺโธ โหติ.}

\proseitem{140}{กตโม ตสฺมิํ สมเย วิญฺญาณ\-กฺขนฺโธ โหติ ฯเปฯ. อยํ ตสฺมิํ สมเย วิญฺญาณ\-กฺขนฺโธ โหติ.}

\proseitem{141}{กตมํ ตสฺมิํ สมเย มนายตนํ โหติ ฯเปฯ. อิทํ ตสฺมิํ สมเย มนายตนํ โหติ.}

\proseitem{142}{\csromanfolio{38}กตมํ ตสฺมิํ สมเย มนินฺทฺริยํ โหติ ฯเปฯ. อิทํ ตสฺมิํ สมเย มนินฺทฺริยํ โหติ.}

\proseitem{143}{กตมา ตสฺมิํ สมเย มโนวิญฺญาณ\-ธาตุ โหติ ฯเปฯ. อยํ ตสฺมิํ สมเย มโนวิญฺญาณ\-ธาตุ โหติ.}

\proseitem{144}{กตมํ ตสฺมิํ สมเย ธมฺมายตนํ โหติ. เวทนากฺขนฺโธ สญฺญากฺขนฺโธ สงฺขาร\-กฺขนฺโธ. อิทํ ตสฺมิํ สมเย ธมฺมายตนํ โหติ.}

\proseitem{145}{กตมา ตสฺมิํ สมเย ธมฺมธาตุ โหติ. เวทนากฺขนฺโธ สญฺญากฺขนฺโธ สงฺขาร\-กฺขนฺโธ. อยํ ตสฺมิํ สมเย ธมฺมธาตุ โหติ.}

\prose{เย วา ปน ตสฺมิํ สมเย อญฺเญปิ อตฺถิ ปฏิจฺจ\-สมุปฺปนฺนา อรูปิโน ธมฺมา. อิเม ธมฺมา กุสลา.}

\vspace{\tipitakaparskip}
\csromancenter{สุญฺญตวาโร.}
\csromancenter{ปฐมํ จิตฺตํ.}
\proseitem{146}{กตเม ธมฺมา กุสลา. ยสฺมิํ สมเย กามาวจรํ กุสลํ จิตฺตํ อุปฺปนฺนํ โหติ โสมนสฺส\-สหคตํ ญาณ\-สมฺปยุตฺตํ สสงฺขาเรน รูปารมฺมณํ วา ฯเปฯ ธมฺมารมฺมณํ วา ยํ ยํ วา ปนารพฺภ. ตสฺมิํ สมเย ผสฺโส โหติ ฯเปฯ อวิกฺเขโป โหติ ฯเปฯ. อิเม ธมฺมา กุสลา.}

\vspace{\tipitakaparskip}
\csromancenter{ทุติยํ จิตฺตํ.}
\proseitem{147}{กตเม ธมฺมา กุสลา. ยสฺมิํ สมเย กามาวจรํ กุสลํ จิตฺตํ อุปฺปนฺนํ โหติ โสมนสฺส\-สหคตํ ญาณวิปฺปยุตฺตํ รูปารมฺมณํ วา สทฺทารมฺมณํ วา คนฺธารมฺมณํ วา รสารมฺมณํ วา โผฏฺฐพฺพา\-รมฺมณํ วา ธมฺมารมฺมณํ วา ยํ ยํ วา ปนารพฺภ. ตสฺมิํ สมเย ผสฺโส โหติ, เวทนา โหติ, สญฺญา โหติ, เจตนา โหติ, จิตฺตํ โหติ, วิตกฺโก โหติ, วิจาโร โหติ, ปีติ โหติ, สุขํ โหติ, จิตฺตสฺเส\-กคฺคตา โหติ, สทฺธินฺทฺริยํ โหติ, วีริยินฺทฺริยํ โหติ, สตินฺทฺริยํ โหติ, สมาธินฺทฺริยํ \csromanfolio{39}โหติ, มนินฺทฺริยํ โหติ, โสมนสฺสิ\-นฺทฺริยํ โหติ, ชีวิตินฺทฺริยํ โหติ, สมฺมาสงฺกปฺโป โหติ, สมฺมาวายาโม โหติ, สมฺมาสติ โหติ, สมฺมาสมาธิ โหติ, สทฺธาพลํ โหติ, วีริยพลํ โหติ, สติพลํ โหติ, สมาธิพลํ โหติ, หิริพลํ โหติ, โอตฺตปฺปพลํ โหติ, อโลโภ โหติ, อโทโส โหติ, อนภิชฺฌา โหติ, อพฺยาปาโท โหติ, หิรี โหติ, โอตฺตปฺปํ โหติ, กายปสฺสทฺธิ โหติ, จิตฺตปสฺสทฺธิ โหติ, กายลหุตา โหติ, จิตฺตลหุตา โหติ, กายมุทุตา โหติ, จิตฺตมุทุตา โหติ, กายกมฺมญฺญตา โหติ, จิตฺต\-กมฺมญฺญตา โหติ, กายปาคุญฺญตา โหติ, จิตฺตปาคุญฺญตา โหติ, กายุชุกตา โหติ, จิตฺตุชุกตา โหติ, สติ โหติ, สมโถ โหติ, ปคฺคาโห โหติ, อวิกฺเขโป โหติ. เย วา ปน ตสฺมิํ สมเย อญฺเญปิ อตฺถิ ปฏิจฺจ\-สมุปฺปนฺนา อรูปิโน ธมฺมา. อิเม ธมฺมา กุสลา ฯเปฯ.}

\prose{ตสฺมิํ โข ปน สมเย จตฺตาโร ขนฺธา โหนฺติ, ทฺวายตนานิ โหนฺติ, เทฺว ธาตุโย โหนฺติ, ตโย อาหารา โหนฺติ, สตฺตินฺทฺริยานิ โหนฺติ, ปญฺจงฺคิกํ ฌานํ โหติ, จตุรงฺคิโก มคฺโค โหติ, ฉ พลานิ โหนฺติ, เทฺว เหตู โหนฺติ, เอโก ผสฺโส โหติ ฯเปฯ เอกํ ธมฺมายตนํ โหติ, เอกา ธมฺมธาตุ โหติ. เย วา ปน ตสฺมิํ สมเย อญฺเญปิ อตฺถิ ปฏิจฺจ\-สมุปฺปนฺนา อรูปิโน ธมฺมา. อิเม ธมฺมา กุสลา ฯเปฯ.}

\proseitem{148}{กตโม ตสฺมิํ สมเย สงฺขาร\-กฺขนฺโธ โหติ. ผสฺโส เจตนา วิตกฺโก วิจาโร ปีติ จิตฺตสฺเส\-กคฺคตา สทฺธินฺทฺริยํ วีริยินฺทฺริยํ สตินฺทฺริยํ สมาธินฺทฺริยํ ชีวิตินฺทฺริยํ สมฺมาสงฺกปฺโป สมฺมาวายาโม สมฺมาสติ สมฺมาสมาธิ สทฺธาพลํ วีริยพลํ สติพลํ สมาธิพลํ หิริพลํ โอตฺตปฺปพลํ อโลโภ อโทโส อนภิชฺฌา อพฺยาปาโท หิรี โอตฺตปฺปํ กายปสฺสทฺธิ จิตฺตปสฺสทฺธิ กายลหุตา จิตฺตลหุตา กายมุทุตา จิตฺตมุทุตา กายกมฺมญฺญตา จิตฺต\-กมฺมญฺญตา กายปาคุญฺญตา จิตฺตปาคุญฺญตา กายุชุกตา จิตฺตุชุกตา สติ สมโถ ปคฺคาโห อวิกฺเขโป. เย วา ปน ตสฺมิํ สมเย อญฺเญปิ อตฺถิ ปฏิจฺจ\-สมุปฺปนฺนา อรูปิโน ธมฺมา ฐเปตฺวา เวทนา\-กฺขนฺธํ ฐเปตฺวา สญฺญากฺขนฺธํ ฐเปตฺวา วิญฺญาณ\-กฺขนฺธํ. อยํ ตสฺมิํ สมเย สงฺขาร\-กฺขนฺโธ โหติ ฯเปฯ. อิเม ธมฺมา กุสลา.}

\vspace{\tipitakaparskip}
\csromancenter{ตติยํ จิตฺตํ.}
\proseitem{149}{\csromanfolio{40}กตเม ธมฺมา กุสลา. ยสฺมิํ สมเย กามาวจรํ กุสลํ จิตฺตํ อุปฺปนฺนํ โหติ โสมนสฺส\-สหคตํ ญาณวิปฺปยุตฺตํ สสงฺขาเรน รูปารมฺมณํ วา ฯเปฯ ธมฺมารมฺมณํ วา ยํ ยํ วา ปนารพฺภ. ตสฺมิํ สมเย ผสฺโส โหติ ฯเปฯ อวิกฺเขโป โหติ ฯเปฯ. อิเม ธมฺมา กุสลา.}

\vspace{\tipitakaparskip}
\csromancenter{จตุตฺถํ จิตฺตํ.}
\proseitem{150}{กตเม ธมฺมา กุสลา. ยสฺมิํ สมเย กามาวจรํ กุสลํ จิตฺตํ อุปฺปนฺนํ โหติ อุเปกฺขา\-สหคตํ ญาณ\-สมฺปยุตฺตํ รูปารมฺมณํ วา สทฺทารมฺมณํ วา คนฺธารมฺมณํ วา รสารมฺมณํ วา โผฏฺฐพฺพารมฺมาณํ วา ธมฺมารมฺมณํ วา ยํ ยํ วา ปนารพฺภ. ตสฺมิํ สมเย ผสฺโส โหติ, เวทนา โหติ, สญฺญา โหติ, เจตนา โหติ, จิตฺตํ โหติ, วิตกฺโก โหติ, วิจาโร โหติ, อุเปกฺขา โหติ, จิตฺตสฺเส\-กคฺคตา โหติ, สทฺธินฺทฺริยํ โหติ, วีริยินฺทฺริยํ โหติ, สตินฺทฺริยํ โหติ, สมาธินฺทฺริยํ โหติ, ปญฺญินฺทฺริยํ โหติ, มนินฺทฺริยํ โหติ, อุเปกฺขินฺทฺริยํ โหติ, ชีวิตินฺทฺริยํ โหติ, สมฺมาทิฏฺฐิ โหติ, สมฺมาสงฺกปฺโป โหติ, สมฺมาวายาโม โหติ, สมฺมาสติ โหติ, สมฺมาสมาธิ โหติ, สทฺธาพลํ โหติ, วีริยพลํ โหติ, สติพลํ โหติ, สมาธิพลํ โหติ, ปญฺญาพลํ โหติ, หิริพลํ โหติ, โอตฺตปฺปพลํ โหติ, อโลโภ โหติ, อโทโส โหติ, อโมโห โหติ, อนภิชฺฌา โหติ, อพฺยาปาโท โหติ, สมฺมาทิฏฺฐิ โหติ, หิรี โหติ, โอตฺตปฺปํ โหติ, กายปสฺสทฺธิ โหติ, จิตฺตปสฺสทฺธิ โหติ, กายลหุตา โหติ, จิตฺตลหุตา โหติ, กายมุทุตา โหติ, จิตฺตมุทุตา โหติ, กายกมฺมญฺญตา โหติ, จิตฺต\-กมฺมญฺญตา โหติ, กายปาคุญฺญตา โหติ, จิตฺตปาคุญฺญตา โหติ, กายุชุกตา โหติ, จิตฺตุชุกตา โหติ, สติ โหติ, สมฺปชญฺญํ โหติ, สมโถ โหติ, วิปสฺสนา โหติ, ปคฺคาโห โหติ, อวิกฺเขโป โหติ. เย วา ปน ตสฺมิํ สมเย อญฺเญปิ อตฺถิ ปฏิจฺจ\-สมุปฺปนฺนา อรูปิโน ธมฺมา. อิเม ธมฺมา กุสลา.}

\proseitem{151}{กตโม ตสฺมิํ สมเย ผสฺโส โหติ. โย ตสฺมิํ สมเย ผสฺโส ผุสนา สมฺผุสนา สมฺผุสิตตฺตํ. อยํ ตสฺมิํ สมเย ผสฺโส โหติ.}

\chapterheadpageread{1. จิตฺตุปฺปาท\-กณฺฑ}
\proseitem{152}{\csromanfolio{41}กตมา ตสฺมิํ สมเย เวทนา โหติ. ยํ ตสฺมิํ สมเย ตชฺชา\-มโนวิญฺญาณธาตุ\-สมฺผสฺสชํ เจตสิกํ เนว สาตํ นาสาตํ เจโต\-สมฺผสฺสชํ อทุกฺขมสุขํ เวทยิตํ เจโต\-สมฺผสฺสชา อทุกฺขมสุขา เวทนา. อยํ ตสฺมิํ สมเย เวทนา โหติ ฯเปฯ.}

\proseitem{153}{กตมา ตสฺมิํ สมเย อุเปกฺขา โหติ. ยํ ตสฺมิํ สมเย เจตสิกํ เนว สาตํ นาสาตํ เจโต\-สมฺผสฺสชํ อทุกฺขมสุขํ เวทยิตํ เจโต\-สมฺผสฺสชา อทุกฺขมสุขา เวทนา. อยํ ตสฺมิํ สมเย อุเปกฺขา โหติ ฯเปฯ.}

\proseitem{154}{กตมํ ตสฺมิํ สมเย อุเปกฺขินฺทฺริยํ โหติ. ยํ ตสฺมิํ สมเย เจตสิกํ เนว สาตํ นาสาตํ เจโต\-สมฺผสฺสชํ อทุกฺขมสุขํ เวทยิตํ เจโต\-สมฺผสฺสชา อทุกฺขมสุขา เวทนา. อิทํ ตสฺมิํ สมเย อุเปกฺขินฺทฺริยํ โหติ ฯเปฯ. เย วา ปน ตสฺมิํ สมเย อญฺเญปิ อตฺถิ ปฏิจฺจ\-สมุปฺปนฺนา อรูปิโน ธมฺมา. อิเม ธมฺมา กุสลา ฯเปฯ.}

\prose{ตสฺมิํ โข ปน สมเย จตฺตาโร ขนฺธา โหนฺติ, ทฺวายตนานิ โหนฺติ, เทฺว ธาตุโย โหนฺติ, ตโย อาหาโร โหนฺติ, อฏฺฐินฺทฺริยานิ โหนฺติ, จตุรงฺคิกํ ฌานํ โหติ, ปญฺจงฺคิโก มคฺโค โหติ, สตฺต พลานิ โหนฺติ, ตโย เหตู โหนฺติ, เอโก ผสฺโส โหติ ฯเปฯ เอกํ ธมฺมายตนํ โหติ, เอกา ธมฺมธาตุ โหติ. เย วา ปน ตสฺมิํ สมเย อญฺเญปิ อตฺถิ ปฏิจฺจ\-สมุปฺปนฺนา อรูปิโน ธมฺมา. อิเม ธมฺมา กุสลา ฯเปฯ.}

\proseitem{155}{กตโม ตสฺมิํ สมเย สงฺขาร\-กฺขนฺโธ โหติ. ผสฺโส เจตนา วิตกฺโก วิจาโร จิตฺตสฺเส\-กคฺคตา สทฺธินฺทฺริยํ วีริยินฺทฺริยํ สตินฺทฺริยํ สมาธินฺทฺริยํ ปญฺญินฺทฺริยํ ชีวิตินฺทฺริยํ สมฺมาทิฏฺฐิ สมฺมาสงฺกปฺโป สมฺมาวายาโม สมฺมาสติ สมฺมาสมาธิ สทฺธาพลํ วีริยพลํ สติพลํ สมาธิพลํ ปญฺญาพลํ หิริพลํ โอตฺตปฺปพลํ อโลโภ อโทโส อโมโห อนภิชฺฌา อพฺยาปาโท สมฺมาทิฏฺฐิ หิรี โอตฺตปฺปํ กายปสฺสทฺธิ จิตฺตปสฺสทฺธิ กายลหุตา จิตฺตลหุตา กายมุทุตา จิตฺตมุทุตา กายกมฺมญฺญตา จิตฺต\-กมฺมญฺญตา กายปาคุญฺญตา จิตฺตปาคุญฺญตา กายุชุกตา จิตฺตุชุกตา สติ สมฺปชญฺญํ สมโถ วิปสฺสนา ปคฺคาโห อวิกฺเขโป.}

\prose{\csromanfolio{42}เย วา ปน ตสฺมิํ สมเย อญฺเญปิ อตฺถิ ปฏิจฺจ\-สมุปฺปนฺนา อรูปิโน ธมฺมา ฐเปตฺวา เวทนา\-กฺขนฺธํ ฐเปตฺวา สญฺญากฺขนฺธํ ฐเปตฺวา วิญฺญาณ\-กฺขนฺธํ. อยํ ตสฺมิํ สมเย สงฺขาร\-กฺขนฺโธ โหติ ฯเปฯ. อิเม ธมฺมา กุสลา.}

\vspace{\tipitakaparskip}
\csromancenter{ปญฺจมํ จิตฺตํ.}
\proseitem{156}{กตเม ธมฺมา กุสลา. ยสฺมิํ สมเย กามาวจรํ กุสลํ จิตฺตํ อุปฺปนฺนํ โหติ อุเปกฺขา\-สหคตํ ญาณ\-สมฺปยุตฺตํ สสงฺขาเรน รูปารมฺมณํ วา ฯเปฯ ธมฺมารมฺมณํ วา ยํ ยํ วา ปนารพฺภ. ตสฺมิํ สมเย ผสฺโส โหติ ฯเปฯ อวิกฺเขโป โหติ ฯเปฯ. อิเม ธมฺมา กุสลา.}

\vspace{\tipitakaparskip}
\csromancenter{ฉฏฺฐํ จิตฺตํ.}
\proseitem{157}{กตเม ธมฺมา กุสลา. ยสฺมิํ สมเย กามาวจรํ กุสลํ จิตฺตํ อุปฺปนฺนํ โหติ อุเปกฺขา\-สหคตํ ญาณวิปฺปยุตฺตํ รูปารมฺมณํ วา สทฺธารมฺมณํ วา คนฺธารมฺมณํ วา รสารมฺมณํ วา โผฏฺฐพฺพา\-รมฺมณํ วา ธมฺมารมฺมณํ วา ยํ ยํ วา ปนารพฺภ. ตสฺมิํ สมเย ผสฺโส โหติ, เวทนา โหติ, สญฺญา โหติ, เจตนา โหติ, จิตฺตํ โหติ, วิตกฺโก โหติ, วิจาโร โหติ, อุเปกฺขา โหติ, จิตฺตสฺเส\-กคฺคตา โหติ, สทฺธินฺทฺริยํ โหติ, วีริยินฺทฺริยํ โหติ, สตินฺทฺริยํ โหติ, สมาธินฺทฺริยํ โหติ, มนินฺทฺริยํ โหติ, อุเปกฺขินฺทฺริยํ โหติ, ชีวิตินฺทฺริยํ โหติ, สมฺมาสงฺกปฺโป โหติ, สมฺมาวายาโม โหติ, สมฺมาสติ โหติ, สมฺมาสมาธิ โหติ, สทฺธาพลํ โหติ, วีริยพลํ โหติ, สติพลํ โหติ, สมาธิพลํ โหติ, หิริพลํ โหติ, โอตฺตปฺปพลํ โหติ, อโลโภ โหติ, อโทโส โหติ, อนภิชฺฌา โหติ, อพฺยาปาโท โหติ, หิรี โหติ, โอตฺตปฺปํ โหติ, กายปสฺสทฺธิ โหติ, จิตฺตปสฺสทฺธิ โหติ, กายลหุตา โหติ, จิตฺตลหุตา โหติ, กายมุทุตา โหติ, จิตฺตมุทุตา โหติ, กายกมฺมญฺญตา โหติ, จิตฺต\-กมฺมญฺญตา โหติ, กายปาคุญฺญตา โหติ, จิตฺตปาคุญฺญตา โหติ, กายุชุกตา โหติ, จิตฺตุชุกตา โหติ, สติ โหติ, สมโถ โหติ, ปคฺคาโห โหติ, \csromanfolio{43}อวิกฺเขโป โหติ. เย วา ปน ตสฺมิํ สมเย อญฺเญปิ อตฺถิ ปฏิจฺจ\-สมุปฺปนฺนา อรูปิโน ธมฺมา. อิเม ธมฺมา กุสลา ฯเปฯ.}

\prose{ตสฺมิํ โข ปน สมเย จตฺตาโร ขนฺธา โหนฺติ, ทฺวายตนานิ โหนฺติ, เทฺว ธาตุโย โหนฺติ, ตโย อาหารา โหนฺติ, สตฺตินฺทฺริยานิ โหนฺติ, จตุรงฺคิกํ ฌานํ โหติ, จตุรงฺคิโก มคฺโค โหติ, ฉ พลานิ โหนฺติ, เทฺว เหตู โหนฺติ, เอโก ผสฺโส โหติ ฯเปฯ เอกํ ธมฺมายตนํ โหติ, เอกา ธมฺมธาตุ โหติ. เย วา ปน ตสฺมิํ สมเย อญฺเญปิ อตฺถิ ปฏิจฺจ\-สมุปฺปนฺนา อรูปิโน ธมฺมา. อิเม ธมฺมา กุสลา ฯเปฯ.}

\proseitem{158}{กตโม ตสฺมิํ สมเย สงฺขาร\-กฺขนฺโธ โหติ. ผสฺโส เจตนา วิตกฺโก วิจาโร จิตฺตสฺเส\-กคฺคตา สทฺธินฺทฺริยํ วีริยินฺทฺริยํ สตินฺทฺริยํ สมาธินฺทฺริยํ ชีวิตินฺทฺริยํ สมฺมาสงฺกปฺโป สมฺมาวายาโม สมฺมาสติ สมฺมาสมาธิ สทฺธาพลํ วีริยพลํ สติพลํ สมาธิพลํ หิริพลํ โอตฺตปฺปพลํ อโลโภ อโทโส อนภิชฺฌา อพฺยาปาโท หิรี โอตฺตปฺปํ กายปสฺสทฺธิ จิตฺตปสฺสทฺธิ กายลหุตา จิตฺตลหุตา กายมุทุตา จิตฺตมุทุตา กายกมฺมญฺญตา จิตฺต\-กมฺมญฺญตา กายปาคุญฺญตา จิตฺตปาคุญฺญตา กายุชุกตา จิตฺตุชุกตา สติ สมโถ ปคฺคาโห อวิกฺเขโป.}

\prose{เย วา ปน ตสฺมิํ สมเย อญฺเญปิ อตฺถิ ปฏิจฺจ\-สมุปฺปนฺนา อรูปิโน ธมฺมา ฐเปตฺวา เวทนา\-กฺขนฺธํ ฐเปตฺวา สญฺญากฺขนฺธํ ฐเปตฺวา วิญฺญาณ\-กฺขนฺธํ. อยํ ตสฺมิํ สมเย สงฺขาร\-กฺขนฺโธ โหติ ฯเปฯ. อิเม ธมฺมา กุสลา.}

\vspace{\tipitakaparskip}
\csromancenter{สตฺตมํ จิตฺตํ.}
\proseitem{159}{กตเม ธมฺมา กุสลา. ยสฺมิํ สมเย กามาวจรํ กุสลํ จิตฺตํ อุปฺปนฺนํ โหติ อุเปกฺขา\-สหคตํ ญาณวิปฺปยุตฺตํ สสงฺขาเรน รูปารมฺมณํ วา ฯเปฯ ธมฺมารมฺมณํ วา ยํ ยํ วา ปนารพฺภ. ตสฺมิํ สมเย ผสฺโส โหติ ฯเปฯ อวิกฺเขโป โหติ ฯเปฯ. อิเม ธมฺมา กุสลา.}

\vspace{\tipitakaparskip}
\csromancenter{อฏฺฐมํ จิตฺตํ.}
\csromancenter{อฏฺฐ กามาวจร\-มหา\-กุสลจิตฺตานิ.}
\csromancenter{ทุติย\-ภาณวาโร.}
\csromanmatikaanchor{mtk.21}
\csromantocmarkref{section}{รูปาวจรกุสลจตุกฺกนย}{mtk.21}
\bookmark[dest={mtk.21},level=1]{รูปาวจรกุสลจตุกฺกนย}
\csromanheader{1}{รูปาวจร\-กุสล}
\csromanheader{2}{\textbf{จตุกฺกนย}}
\proseitem{160}{\csromanfolio{44}กตเม ธมฺมา กุสลา. ยสฺมิํ สมเย รูปูปปตฺติยา มคฺคํ ภาเวติ วิวิจฺเจว กาเมหิ วิวิจฺจ อกุสเลหิ ธมฺเมหิ สวิตกฺกํ สวิจารํ วิเวกชํ ปีติสุขํ ปฐมํ ฌานํ\csromansharedfootnote{6a703b82ede4a4c7}{ปฐมชฺฌานํ (สี)} อุปสมฺปชฺช วิหรติ ปถวีกสิณํ. ตสฺมิํ สมเย ผสฺโส โหติ ฯเปฯ อวิกฺเขโป โหติ ฯเปฯ. อิเม ธมฺมา กุสลา.}

\proseitem{161}{กตเม ธมฺมา กุสลา. ยสฺมิํ สมเย รูปูปปตฺติยา มคฺคํ ภาเวติ วิตกฺก\-วิจารานํ วูปสมา อชฺฌตฺตํ สมฺปสาทนํ เจตโส เอโกทิภาวํ อวิตกฺกํ อวิจารํ สมาธิชํ ปีติสุขํ ทุติยํ ฌานํ\csromansharedfootnote{9dc8cc7fcd57cc1f}{ทุติยชฺฌานํ (สี)} อุปสมฺปชฺช วิหรติ ปถวีกสิณํ. ตสฺมิํ สมเย ผสฺโส โหติ, เวทนา โหติ, สญฺญา โหติ, เจตนา โหติ, จิตฺตํ โหติ, ปีติ โหติ, สุขํ โหติ, จิตฺตสฺเส\-กคฺคตา โหติ, สทฺธินฺทฺริยํ โหติ, วีริยินฺทฺริยํ โหติ, สตินฺทฺริยํ โหติ, สมาธินฺทฺริยํ โหติ, ปญฺญินฺทฺริยํ โหติ, มนินฺทฺริยํ โหติ, โสมนสฺสิ\-นฺทฺริยํ โหติ, ชีวิตินฺทฺริยํ โหติ, สมฺมาทิฏฺฐิ โหติ, สมฺมาวายาโม โหติ ฯเปฯ ปคฺคาโห โหติ, อวิกฺเขโป โหติ. เย วา ปน ตสฺมิํ สมเย อญฺเญปิ อตฺถิ ปฏิจฺจ\-สมุปฺปนฺนา อรูปิโน ธมฺมา. อิเม ธมฺมา กุสลา ฯเปฯ.}

\prose{ตสฺมิํ โข ปน สมเย จตฺตาโร ขนฺธา โหนฺติ, ทฺวายตนานิ โหนฺติ, เทฺว ธาตุโย โหนฺติ, ตโย อาหารา โหนฺติ, อฏฺฐินฺทฺริยานิ โหนฺติ, ติวงฺคิกํ ฌานํ โหติ, จตุรงฺคิโก มคฺโค โหติ, สตฺต พลานิ โหนฺติ, ตโย เหตู โหนฺติ, เอโก ผสฺโส โหติ ฯเปฯ เอกํ ธมฺมายตนํ โหติ, เอกา ธมฺมธาตุ โหติ. เย วา ปน ตสฺมิํ สมเย อญฺเญปิ อตฺถิ ปฏิจฺจ\-สมุปฺปนฺนา อรูปิโน ธมฺมา. อิเม ธมฺมา กุสลา ฯเปฯ.}

\proseitem{162}{กตโม ตสฺมิํ สมเย สงฺขาร\-กฺขนฺโธ โหติ. ผสฺโส เจตนา ปีติ จิตฺตสฺเส\-กคฺคตา สทฺธินฺทฺริยํ วีริยินฺทฺริยํ สตินฺทฺริยํ สมาธินฺทฺริยํ ปญฺญินฺทฺริยํ ชีวิตินฺทฺริยํ สมฺมาทิฏฺฐิ สมฺมาวายาโม ฯเปฯ ปคฺคาโห อวิกฺเขโป. เย วา ปน ตสฺมิํ สมเย อญฺเญปิ อตฺถิ ปฏิจฺจ\-สมุปฺปนฺนา อรูปิโน ธมฺมา \csromanfolio{45}ฐเปตฺวา เวทนา\-กฺขนฺธํ ฐเปตฺวา สญฺญากฺขนฺธํ ฐเปตฺวา วิญฺญาณ\-กฺขนฺธํ. อยํ ตสฺมิํ สมเย สงฺขาร\-กฺขนฺโธ โหติ ฯเปฯ. อิเม ธมฺมา กุสลา.}

\proseitem{163}{กตเม ธมฺมา กุสลา. ยสฺมิํ สมเย รูปูปปตฺติยา มคฺคํ ภาเวติ ปีติยา จ วิราคา อุเปกฺขโก จ วิหรติ สโต จ สมฺปชาโน สุขญฺจ กาเยน ปฏิสํเวเทติ. ยํ ตํ อริยา อาจิกฺขนฺติ อุเปกฺขโก สติมา สุขวิหารีติ ตติยํ ฌานํ\csromansharedfootnote{1a0d5e790a5f1f73}{ตติยชฺฌานํ (สี)} อุปสมฺปชฺช วิหรติ ปถวีกสิณํ. ตสฺมิํ สมเย ผสฺโส โหติ, เวทนา โหติ, สญฺญา โหติ, เจตนา โหติ, จิตฺตํ โหติ, สุขํ โหติ, จิตฺตสฺเส\-กคฺคตา โหติ, สทฺธินฺทฺริยํ โหติ, วีริยินฺทฺริยํ โหติ, สตินฺทฺริยํ โหติ, สมาธินฺทฺริยํ โหติ, ปญฺญินฺทฺริยํ โหติ, มนินฺทฺริยํ โหติ, โสมนสฺสิ\-นฺทฺริยํ โหติ, ชีวิตินฺทฺริยํ โหติ, สมฺมาทิฏฺฐิ โหติ, สมฺมาวายาโม โหติ ฯเปฯ ปคฺคาโห โหติ, อวิกฺเขโป โหติ. เย วา ปน ตสฺมิํ สมเย อญฺเญปิ อตฺถิ ปฏิจฺจ\-สมุปฺปนฺนา อรูปิโน ธมฺมา. อิเม ธมฺมา กุสลา ฯเปฯ.}

\prose{ตสฺมิํ โข ปน สมเย จตฺตาโร ขนฺธา โหนฺติ, ทฺวายตนานิ โหนฺติ, เทฺว ธาตุโย โหนฺติ, ตโย อาหารา โหนฺติ, อฏฺฐินฺทฺริยานิ โหนฺติ, ทุวงฺคิกํ ฌานํ โหติ, จตุรงฺคิโก มคฺโค โหติ, สตฺต พลานิ โหนฺติ, ตโย เหตู โหนฺติ, เอโก ผสฺโส โหติ ฯเปฯ เอกํ ธมฺมายตนํ โหติ, เอกา ธมฺมธาตุ โหติ. เย วา ปน ตสฺมิํ สมเย อญฺเญปิ อตฺถิ ปฏิจฺจ\-สมุปฺปนฺนา อรูปิโน ธมฺมา. อิเม ธมฺมา กุสลา ฯเปฯ.}

\proseitem{164}{กตโม ตสฺมิํ สมเย สงฺขาร\-กฺขนฺโธ โหติ. ผสฺโส เจตนา จิตฺตสฺเส\-กคฺคตา สทฺธินฺทฺริยํ วีริยินฺทฺริยํ สตินฺทฺริยํ สมาธินฺทฺริยํ ปญฺญินฺทฺริยํ ชีวิตินฺทฺริยํ สมฺมาทิฏฺฐิ สมฺมาวายาโม ฯเปฯ ปคฺคาโห อวิกฺเขโป. เย วา ปน ตสฺมิํ สมเย อญฺเญปิ อตฺถิ ปฏิจฺจ\-สมุปฺปนฺนา อรูปิโน ธมฺมา ฐเปตฺวา เวทนา\-กฺขนฺธํ ฐเปตฺวา สญฺญากฺขนฺธํ ฐเปตฺวา วิญฺญาณ\-กฺขนฺธํ. อยํ ตสฺมิํ สมเย สงฺขาร\-กฺขนฺโธ โหติ ฯเปฯ. อิเม ธมฺมา กุสลา.}

\proseitem{165}{กตเม ธมฺมา กุสลา. ยสฺมิํ สมเย รูปูปปตฺติยา มคฺคํ ภาเวติ สุขสฺส จ ปหานา ทุกฺขสฺส จ ปหานา ปุพฺเพว โสมนสฺส\-โทมนสฺสานํ อตฺถงฺคมา อทุกฺขมสุขํ อุเปกฺขา\-สติ\-ปาริสุทฺธิํ จตุตฺถํ \csromanfolio{46}ฌานํ\csromansharedfootnote{5edf374ad3a4ac94}{จตุตฺถชฺฌานํ (สี)} อุปสมฺปชฺช วิหรติ ปถวีกสิณํ. ตสฺมิํ สมเย ผสฺโส โหติ, เวทนา โหติ, สญฺญา โหติ, เจตนา โหติ, จิตฺตํ โหติ, อุเปกฺขา โหติ, จิตฺตสฺเส\-กคฺคตา โหติ, สทฺธินฺทฺริยํ โหติ, วีริยินฺทฺริยํ โหติ, สตินฺทฺริยํ โหติ, สมาธินฺทฺริยํ โหติ, ปญฺญินฺทฺริยํ โหติ, มนินฺทฺริยํ โหติ, อุเปกฺขินฺทฺริยํ โหติ, ชีวิตินฺทฺริยํ โหติ, สมฺมาทิฏฺฐิ โหติ, สมฺมาวายาโม โหติ ฯเปฯ ปคฺคาโห โหติ, อวิกฺเขโป โหติ. เย วา ปน ตสฺมิํ สมเย อญฺเญปิ อตฺถิ ปฏิจฺจ\-สมุปฺปนฺนา อรูปิโน ธมฺมา. อิเม ธมฺมา กุสลา ฯเปฯ.}

\prose{ตสฺมิํ โข ปน สมเย จตฺตาโร ขนฺธา โหนฺติ, ทฺวายตนานิ โหนฺติ, เทฺว ธาตุโย โหนฺติ, ตโย อาหารา โหนฺติ, อฏฺฐินฺทฺริยานิ โหนฺติ, ทุวงฺคิกํ ฌานํ โหติ, จตุรงฺคิโก มคฺโค โหติ, สตฺต พลานิ โหนฺติ, ตโย เหตู โหนฺติ, เอโก ผสฺโส โหติ ฯเปฯ เอกํ ธมฺมายตนํ โหติ, เอกา ธมฺมธาตุ โหติ. เย วา ปน ตสฺมิํ สมเย อญฺเญปิ อตฺถิ ปฏิจฺจ\-สมุปฺปนฺนา อรูปิโน ธมฺมา. อิเม ธมฺมา กุสลา ฯเปฯ.}

\proseitem{166}{กตโม ตสฺมิํ สมเย สงฺขาร\-กฺขนฺโธ โหติ. ผสฺโส เจตนา จิตฺตสฺเส\-กคฺคตา สทฺธินฺทฺริยํ วีริยินฺทฺริยํ สตินฺทฺริยํ สมาธินฺทฺริยํ ปญฺญินฺทฺริยํ ชีวิตินฺทฺริยํ สมฺมาทิฏฺฐิ สมฺมาวายาโม ฯเปฯ ปคฺคาโห อวิกฺเขโป. เย วา ปน ตสฺมิํ สมเย อญฺเญปิ อตฺถิ ปฏิจฺจ\-สมุปฺปนฺนา อรูปิโน ธมฺมา ฐเปตฺวา เวทนา\-กฺขนฺธํ ฐเปตฺวา สญฺญากฺขนฺธํ ฐเปตฺวา วิญฺญาณ\-กฺขนฺธํ. อยํ ตสฺมิํ สมเย สงฺขาร\-กฺขนฺโธ โหติ ฯเปฯ. อิเม ธมฺมา กุสลา.}

\vspace{\tipitakaparskip}
\csromancenter{จตุกฺกนโย.}
\csromanmatikaanchor{mtk.22}
\csromantocmarkref{section}{ปญฺจกนย}{mtk.22}
\bookmark[dest={mtk.22},level=1]{ปญฺจกนย}
\csromanheader{1}{\textbf{ปญฺจกนย}}
\proseitem{167}{กตเม ธมฺมา กุสลา. ยสฺมิํ สมเย รูปูปปตฺติยา มคฺคํ ภาเวติ วิวิจฺเจว กาเมหิ ฯเปฯ ปฐมํ ฌานํ อุปสมฺปชฺช วิหรติ \csromanfolio{47}ปถวีกสิณํ. ตสฺมิํ สมเย ผสฺโส โหติ ฯเปฯ อวิกฺเขโป โหติ ฯเปฯ. อิเม ธมฺมา กุสลา.}

\proseitem{168}{กตเม ธมฺมา กุสลา. ยสฺมิํ สมเย รูปูปปตฺติยา มคฺคํ ภาเวติ อวิตกฺกํ วิจารมตฺตํ สมาธิชํ ปีติสุขํ ทุติยํ ฌานํ อุปสมฺปชฺช วิหรติ ปถิวีกสิณํ. ตสฺมิํ สมเย ผสฺโส โหติ, เวทนา โหติ, สญฺญา โหติ, เจตนา โหติ, จิตฺตํ โหติ, วิจาโร โหติ, ปีติ โหติ, สุขํ โหติ, จิตฺตสฺเส\-กคฺคตา โหติ, สทฺธินฺทฺริยํ โหติ, วีริยินฺทฺริยํ โหติ, สตินฺทฺริยํ โหติ, สมาธินฺทฺริยํ โหติ, ปญฺญินฺทฺริยํ โหติ, มนินฺทฺริยํ โหติ, โสมนสฺสิ\-นฺทฺริยํ โหติ, ชีวิตินฺทฺริยํ โหติ, สมฺมาทิฏฺฐิ โหติ, สมฺมาวายาโม โหติ ฯเปฯ ปคฺคาโห โหติ, อวิกฺเขโป โหติ. เย วา ปน ตสฺมิํ สมเย อญฺเญปิ อตฺถิ ปฏิจฺจ\-สมุปฺปนฺนา อรูปิโน ธมฺมา. อิเม ธมฺมา กุสลา ฯเปฯ.}

\prose{ตสฺมิํ โข ปน สมเย จตฺตาโร ขนฺธา โหนฺติ, ทฺวายตนานิ โหนฺติ, เทฺว ธาตุโย โหนฺติ, ตโย อาหารา โหนฺติ, อฏฺฐินฺทฺริยานิ โหนฺติ, จตุรงฺคิกํ ฌานํ โหติ, จตุรงฺคิโก มคฺโค โหติ, สตฺต พลานิ โหนฺติ, ตโย เหตู โหนฺติ, เอโก ผสฺโส โหติ ฯเปฯ เอกํ ธมฺมายตนํ โหติ, เอกา ธมฺมธาตุ โหติ. เย วา ปน ตสฺมิํ สมเย อญฺเญปิ อตฺถิ ปฏิจฺจ\-สมุปฺปนฺนา อรูปิโน ธมฺมา. อิเม ธมฺมา กุสลา ฯเปฯ.}

\proseitem{169}{กตโม ตสฺมิํ สมเย สงฺขาร\-กฺขนฺโธ โหติ. ผสฺโส เจตนา วิจาโร ปีติ จิตฺตสฺเส\-กคฺคตา สทฺธินฺทฺริยํ วีริยินฺทฺริยํ สตินฺทฺริยํ สมาธินฺทฺริยํ ปญฺญินฺทฺริยํ ชีวิตินฺทฺริยํ สมฺมาทิฏฺฐิ สมฺมาวายาโม ฯเปฯ ปคฺคาโห อวิกฺเขโป. เย วา ปน ตสฺมิํ สมเย อญฺเญปิ อตฺถิ ปฏิจฺจ\-สมุปฺปนฺนา อรูปิโน ธมฺมา ฐเปตฺวา เวทนา\-กฺขนฺธํ ฐเปตฺวา สญฺญากฺขนฺธํ ฐเปตฺวา วิญฺญาณ\-กฺขนฺธํ. อยํ ตสฺมิํ สมเย สงฺขาร\-กฺขนฺโธ โหติ ฯเปฯ. อิเม ธมฺมา กุสลา.}

\proseitem{170}{กตเม ธมฺมา กุสลา. ยสฺมิํ สมเย รูปูปปตฺติยา มคฺคํ ภาเวติ วิตกฺก\-วิจารานํ วูปสมา ฯเปฯ ตติยํ ฌานํ อุปสมฺปชฺช วิหรติ ปถวีกสิณํ. ตสฺมิํ สมเย ผสฺโส โหติ, เวทนา \csromanfolio{48}โหติ, สญฺญา โหติ, เจตนา โหติ, จิตฺตํ โหติ, ปีติ โหติ, สุขํ โหติ, จิตฺตสฺเส\-กคฺคตา โหติ, สทฺธินฺทฺริยํ โหติ, วีริยินฺทฺริยํ โหติ, สตินฺทฺริยํ โหติ, สมาธินฺทฺริยํ โหติ, ปญฺญินฺทฺริยํ โหติ, มนินฺทฺริยํ โหติ, โสมนสฺสิ\-นฺทฺริยํ โหติ, ชีวิตินฺทฺริยํ โหติ, สมฺมาทิฏฺฐิ โหติ, สมฺมาวายาโม โหติ ฯเปฯ ปคฺคาโห โหติ, อวิกฺเขโป โหติ. เย วา ปน ตสฺมิํ สมเย อญฺเญปิ อตฺถิ ปฏิจฺจ\-สมุปฺปนฺนา อรูปิโน ธมฺมา. อิเม ธมฺมา กุสลา ฯเปฯ.}

\prose{ตสฺมิํ โข ปน สมเย จตฺตาโร ขนฺธา โหนฺติ, ทฺวายตนานิ โหนฺติ, เทฺว ธาตุโย โหนฺติ, ตโย อาหารา โหนฺติ, อฏฺฐินฺทฺริยานิ โหนฺติ, ติวงฺคิกํ ฌานํ โหติ, จตุรงฺคิโก มคฺโค โหติ, สตฺต พลานิ โหนฺติ, ตโย เหตู โหนฺติ, เอโก ผสฺโส โหติ ฯเปฯ เอกํ ธมฺมายตนํ โหติ, เอกา ธมฺมธาตุ โหติ. เย วา ปน ตสฺมิํ สมเย อญฺเญปิ อตฺถิ ปฏิจฺจ\-สมุปฺปนฺนา อรูปิโน ธมฺมา. อิเม ธมฺมา กุสลา ฯเปฯ.}

\proseitem{171}{กตโม ตสฺมิํ สมเย สงฺขาร\-กฺขนฺโธ โหติ. ผสฺโส เจตนา ปีติ จิตฺตสฺเส\-กคฺคตา สทฺธินฺทฺริยํ วีริยินฺทฺริยํ สตินฺทฺริยํ สมาธินฺทฺริยํ ปญฺญินฺทฺริยํ ชีวิตินฺทฺริยํ สมฺมาทิฏฺฐิ สมฺมาวายาโม ฯเปฯ ปคฺคาโห อวิกฺเขโป. เย วา ปน ตสฺมิํ สมเย อญฺเญปิ อตฺถิ ปฏิจฺจ\-สมุปฺปนฺนา อรูปิโน ธมฺมา ฐเปตฺวา เวทนา\-กฺขนฺธํ ฐเปตฺวา สญฺญากฺขนฺธํ ฐเปตฺวา วิญฺญาณ\-กฺขนฺธํ. อยํ ตสฺมิํ สมเย สงฺขาร\-กฺขนฺโธ โหติ ฯเปฯ. อิเม ธมฺมา กุสลา.}

\proseitem{172}{กตเม ธมฺมา กุสลา. ยสฺมิํ สมเย รูปูปปตฺติยา มคฺคํ ภาเวติ ปีติยา จ วิราคา ฯเปฯ จตุตฺถํ ฌานํ อุปสมฺปชฺช วิหรติ ปถวีกสิณํ. ตสฺมิํ สมเย ผสฺโส โหติ, เวทนา โหติ, สญฺญา โหติ, เจตนา โหติ, จิตฺตํ โหติ, สุขํ โหติ, จิตฺตสฺเส\-กคฺคตา โหติ, สทฺธินฺทฺริยํ โหติ, วีริยินฺทฺริยํ โหติ, สตินฺทฺริยํ โหติ, สมาธินฺทฺริยํ โหติ, ปญฺญินฺทฺริยํ โหติ, มนินฺทฺริยํ โหติ, โสมนสฺสิ\-นฺทฺริยํ โหติ, ชีวิตินฺทฺริยํ โหติ, สมฺมาทิฏฺฐิ โหติ, สมฺมาวายาโม โหติ ฯเปฯ ปคฺคาโห โหติ, อวิกฺเขโป โหติ. เย วา ปน ตสฺมิํ สมเย อญฺเญปิ อตฺถิ ปฏิจฺจ\-สมุปฺปนฺนา อรูปิโน ธมฺมา. อิเม ธมฺมา กุสลา ฯเปฯ.}

\chapterheadpageread{1. จิตฺตุปฺปาท\-กณฺฑ}
\prose{\csromanfolio{49}ตสฺมิํ โข ปน สมเย จตฺตาโร ขนฺธา โหนฺติ, ทฺวายตนานิ โหนฺติ, เทฺว ธาตุโย โหนฺติ, ตโย อาหารา โหนฺติ, อฏฺฐินฺทฺริยานิ โหนฺติ, ทุวงฺคิกํ ฌานํ โหติ, จตุรงฺคิโก มคฺโค โหติ, สตฺต พลานิ โหนฺติ, ตโย เหตู โหนฺติ, เอโก ผสฺโส โหติ ฯเปฯ เอกํ ธมฺมายตนํ โหติ, เอกา ธมฺมธาตุ โหติ. เย วา ปน ตสฺมิํ สมเย อญฺเญปิ อตฺถิ ปฏิจฺจ\-สมุปฺปนฺนา อรูปิโน ธมฺมา. อิเม ธมฺมา กุสลา ฯเปฯ.}

\proseitem{173}{กตโม ตสฺมิํ สมเย สงฺขาร\-กฺขนฺโธ โหติ. ผสฺโส เจตนา จิตฺตสฺเส\-กคฺคตา สทฺธินฺทฺริยํ วีริยินฺทฺริยํ สตินฺทฺริยํ สมาธินฺทฺริยํ ปญฺญินฺทฺริยํ ชีวิตินฺทฺริยํ สมฺมาทิฏฺฐิ สมฺมาวายาโม ฯเปฯ ปคฺคาโห อวิกฺเขโป. เย วา ปน ตสฺมิํ สมเย อญฺเญปิ อตฺถิ ปฏิจฺจ\-สมุปฺปนฺนา อรูปิโน ธมฺมา ฐเปตฺวา เวทนา\-กฺขนฺธํ ฐเปตฺวา สญฺญากฺขนฺธํ ฐเปตฺวา วิญฺญาณ\-กฺขนฺธํ. อยํ ตสฺมิํ สมเย สงฺขาร\-กฺขนฺโธ โหติ ฯเปฯ. อิเม ธมฺมา กุสลา.}

\proseitem{174}{กตเม ธมฺมา กุสลา. ยสฺมิํ สมเย รูปูปปตฺติยา มคฺคํ ภาเวติ สุขสฺส จ ปหานา ฯเปฯ ปญฺจมํ ฌานํ\csromansharedfootnote{fba3fbf111c656ba}{ปญฺจมชฺฌานํ (สี)} อุปสมฺปชฺช วิหรติ ปถวีกสิณํ. ตสฺมิํ สมเย ผสฺโส โหติ, เวทนา โหติ, สญฺญา โหติ, เจตนา โหติ, จิตฺตํ โหติ, อุเปกฺขา โหติ, จิตฺตสฺเส\-กคฺคตา โหติ, สทฺธินฺทฺริยํ โหติ, วีริยินฺทฺริยํ โหติ, สตินฺทฺริยํ โหติ, สมาธินฺทฺริยํ โหติ, ปญฺญินฺทฺริยํ โหติ, มนินฺทฺริยํ โหติ, อุเปกฺขินฺทฺริยํ โหติ, ชีวิตินฺทฺริยํ โหติ, สมฺมาทิฏฺฐิ โหติ, สมฺมาวายาโม โหติ ฯเปฯ ปคฺคาโห โหติ, อวิกฺเขโป โหติ. เย วา ปน ตสฺมิํ สมเย อญฺเญปิ อตฺถิ ปฏิจฺจ\-สมุปฺปนฺนา อรูปิโน ธมฺมา. อิเม ธมฺมา กุสลา ฯเปฯ.}

\prose{ตสฺมิํ โข ปน สมเย จตฺตาโร ขนฺธา โหนฺติ, ทฺวายตนานิ โหนฺติ, เทฺว ธาตุโย โหนฺติ, ตโย อาหารา โหนฺติ, อฏฺฐินฺทฺริยานิ โหนฺติ, ทุวงฺคิกํ ฌานํ โหติ, จตุรงฺคิโก มคฺโค โหติ, สตฺต พลานิ โหนฺติ, ตโย เหตู โหนฺติ, เอโก ผสฺโส โหติ ฯเปฯ เอกํ ธมฺมายตนํ โหติ, เอกา ธมฺมธาตุ โหติ. เย \csromanfolio{50}วา ปน ตสฺมิํ สมเย อญฺเญปิ อตฺถิ ปฏิจฺจ\-สมุปฺปนฺนา อรูปิโน ธมฺมา. อิเม ธมฺมา กุสลา ฯเปฯ.}

\proseitem{175}{กตโม ตสฺมิํ สมเย สงฺขาร\-กฺขนฺโธ โหติ. ผสฺโส เจตนา จิตฺตสฺเส\-กคฺคตา สทฺธินฺทฺริยํ วีริยินฺทฺริยํ สตินฺทฺริยํ สมาธินฺทฺริยํ ปญฺญินฺทฺริยํ ชีวิตินฺทฺริยํ สมฺมาทิฏฺฐิ สมฺมาวายาโม ฯเปฯ ปคฺคาโห อวิกฺเขโป. เย วา ปน ตสฺมิํ สมเย อญฺเญปิ อตฺถิ ปฏิจฺจ\-สมุปฺปนฺนา อรูปิโน ธมฺมา ฐเปตฺวา เวทนา\-กฺขนฺธํ ฐเปตฺวา สญฺญากฺขนฺธํ ฐเปตฺวา วิญฺญาณ\-กฺขนฺธํ. อยํ ตสฺมิํ สมเย สงฺขาร\-กฺขนฺโธ โหติ ฯเปฯ. อิเม ธมฺมา กุสลา.}

\vspace{\tipitakaparskip}
\csromancenter{ปญฺจกนโย.}
\csromanmatikaanchor{mtk.23}
\csromantocmarkref{section}{จตสฺโส ปฏิปทา}{mtk.23}
\bookmark[dest={mtk.23},level=1]{จตสฺโส ปฏิปทา}
\csromanheader{1}{\textbf{จตสฺโส ปฏิปทา}}
\proseitem{176}{กตเม ธมฺมา กุสลา. ยสฺมิํ สมเย รูปูปปตฺติยา มคฺคํ ภาเวติ วิวิจฺเจว กาเมหิ ฯเปฯ ปฐมํ ฌานํ อุปสมฺปชฺช วิหรติ ทุกฺข\-ปฏิปทํ\csromansharedfootnote{6c79163cf84c8e14}{ทุกฺขาปฏิปทํ (พหูสุ)} ทนฺธาภิญฺญํ ปถวีกสิณํ. ตสฺมิํ สมเย ผสฺโส โหติ ฯเปฯ อวิกฺเขโป โหติ ฯเปฯ. อิเม ธมฺมา กุสลา.}

\proseitem{177}{กตเม ธมฺมา กุสลา. ยสฺมิํ สมเย รูปูปปตฺติยา มคฺคํ ภาเวติ วิวิจฺเจว กาเมหิ ฯเปฯ ปฐมํ ฌานํ อุปสมฺปชฺช วิหรติ ทุกฺข\-ปฏิปทํ ขิปฺปาภิญฺญํ ปถวีกสิณํ. ตสฺมิํ สมเย ผสฺโส โหติ ฯเปฯ อวิกฺเขโป โหติ ฯเปฯ. อิเม ธมฺมา กุสลา.}

\proseitem{178}{กตเม ธมฺมา กุสลา. ยสฺมิํ สมเย รูปูปปตฺติยา มคฺคํ ภาเวติ วิวิจฺเจว กาเมหิ ฯเปฯ ปฐมํ ฌานํ อุปสมฺปชฺช วิหรติ สุขปฏิปทํ\csromansharedfootnote{351f85dfecf34cfb}{สุขาปฏิปทํ (พหูสุ)} ทนฺธาภิญฺญํ ปถวีกสิณํ. ตสฺมิํ สมเย ผสฺโส โหติ ฯเปฯ อวิกฺเขโป โหติ ฯเปฯ. อิเม ธมฺมา กุสลา.}

\proseitem{179}{กตเม ธมฺมา กุสลา. ยสฺมิํ สมเย รูปูปปตฺติยา มคฺคํ ภาเวติ วิวิจฺเจว กาเมหิ ฯเปฯ. ปฐมํ ฌานํ อุปสมฺปชฺช วิหรติ สุขปฏิปทํ \csromanfolio{51}ขิปฺปาภิญฺญํ ปถวีกสิณํ. ตสฺมิํ สมเย ผสฺโส โหติ ฯเปฯ อวิกฺเขโป โหติ ฯเปฯ. อิเม ธมฺมา กุสลา ฯเปฯ.}

\proseitem{180}{กตเม ธมฺมา กุสลา. ยสฺมิํ สมเย รูปูปปตฺติยา มคฺคํ ภาเวติ วิตกฺก\-วิจารานํ วูปสมา ฯเปฯ ทุติยํ ฌานํ ฯเปฯ ตติยํ ฌานํ ฯเปฯ จตุตฺถํ ฌานํ ฯเปฯ ปฐมํ ฌานํ ฯเปฯ ปญฺจมํ ฌานํ อุปสมฺปชฺช วิหรติ ทุกฺข\-ปฏิปทํ ทนฺธาภิญฺญํ ปถวีกสิณํ ฯเปฯ ทุกฺข\-ปฏิปทํ ขิปฺปาภิญฺญํ ปถวีกสิณํ ฯเปฯ สุขปฏิปทํ ทนฺธาภิญฺญํ ปถวิกสิณํ ฯเปฯ สุขปฏิปทํ ขิปฺปาภิญฺญํ ปถวีกสิณํ. ตสฺมิํ สมเย ผสฺโส โหติ ฯเปฯ อวิกฺเขโป โหติ ฯเปฯ. อิเม ธมฺมา กุสลา.}

\vspace{\tipitakaparskip}
\csromancenter{จตสฺโส ปฏิปทา.}
\csromanmatikaanchor{mtk.24}
\csromantocmarkref{section}{จตฺตาริ อารมฺมณานิ}{mtk.24}
\bookmark[dest={mtk.24},level=1]{จตฺตาริ อารมฺมณานิ}
\csromanheader{1}{\textbf{จตฺตาริ อารมฺมณานิ}}
\proseitem{181}{กตเม ธมฺมา กุสลา. ยสฺมิํ สมเย รูปูปปตฺติยา มคฺคํ ภาเวติ วิวิจฺเจว กาเมหิ ฯเปฯ ปฐมํ ฌานํ อุปสมฺปชฺช วิหรติ ปริตฺตํ ปริตฺตา\-รมฺมณํ ปถวีกสิณํ. ตสฺมิํ สมเย ผสฺโส โหติ ฯเปฯ อวิกฺเขโป โหติ ฯเปฯ. อิเม ธมฺมา กุสลา.}

\proseitem{182}{กตเม ธมฺมา กุสลา. ยสฺมิํ สมเย รูปูปปตฺติยา มคฺคํ ภาเวติ วิวิจฺเจว กาเมหิ ฯเปฯ ปฐมํ ฌานํ อุปสมฺปชฺช วิหรติ ปริตฺตํ อปฺปมาณา\-รมฺมณํ ปถวีกสิณํ. ตสฺมิํ สมเย ผสฺโส โหติ ฯเปฯ อวิกฺเขโป โหติ ฯเปฯ. อิเม ธมฺมา กุสลา.}

\proseitem{183}{กตเม ธมฺมา กุสลา. ยสฺมิํ สมเย รูปูปปตฺติยา มคฺคํ ภาเวติ วิวิจฺเจว กาเมหิ ฯเปฯ ปฐมํ ฌานํ อุปสมฺปชฺช วิหรติ อปฺปมาณํ ปริตฺตา\-รมฺมณํ ปถวีกสิณํ. ตสฺมิํ สมเย ผสฺโส โหติ ฯเปฯ อวิกฺเขโป โหติ ฯเปฯ. อิเม ธมฺมา กุสลา.}

\proseitem{184}{กตเม ธมฺมา กุสลา. ยสฺมิํ สมเย รูปูปปตฺติยา มคฺคํ ภาเวติ วิวิจฺเจว กาเมหิ ฯเปฯ ปฐมํ ฌานํ อุปสมฺปชฺช วิหรติ อปฺปมาณํ อปฺปมาณา\-รมฺมณํ ปถวีกสิณํ. ตสฺมิํ สมเย ผสฺโส โหติ ฯเปฯ อวิกฺเขโป โหติ ฯเปฯ. อิเม ธมฺมา กุสลา.}

\proseitem{185}{\csromanfolio{52}กตเม ธมฺมา กุสลา. ยสฺมิํ สมเย รูปูปปตฺติยา มคฺคํ ภาเวติ วิตกฺก\-วิจารานํ วูปสมา ฯเปฯ ทุติยํ ฌานํ ฯเปฯ ตติยํ ฌานํ ฯเปฯ จตุตฺถํ ฌานํ ฯเปฯ ปฐมํ ฌานํ ฯเปฯ ปญฺจมํ ฌานํ อุปสมฺปชฺช วิหรติ ปริตฺตํ ปริตฺตา\-รมฺมณํ ปถวีกสิณํ ฯเปฯ ปริตฺตํ อปฺปมาณา\-รมฺมณํ ปถวีกสิณํ ฯเปฯ อปฺปมาณํ ปริตฺตา\-รมฺมณํ ปถวีกสิณํ ฯเปฯ อปฺปมาณํ อปฺปมาณา\-รมฺมณํ ปถวีกสิณํ. ตสฺมิํ สมเย ผสฺโส โหติ ฯเปฯ อวิกฺเขโป โหติ ฯเปฯ. อิเม ธมฺมา กุสลา. จตฺตาริ อารมฺมณานิ.}

\csromanmatikaanchor{mtk.25}
\csromantocmarkref{section}{โสฬสกฺขตฺตุกํ}{mtk.25}
\bookmark[dest={mtk.25},level=1]{โสฬสกฺขตฺตุกํ}
\csromanheader{1}{โสฬส\-กฺขตฺตุกํ}
\proseitem{186}{กตเม ธมฺมา กุสลา. ยสฺมิํ สมเย รูปูปปตฺติยา มคฺคํ ภาเวติ วิวิจฺเจว กาเมหิ ฯเปฯ ปฐมํ ฌานํ อุปสมฺปชฺช วิหรติ ทุกฺข\-ปฏิปทํ ทนฺธาภิญฺญํ ปริตฺตํ ปริตฺตา\-รมฺมณํ ปถวีกสิณํ. ตสฺมิํ สมเย ผสฺโส โหติ ฯเปฯ อวิกฺเขโป โหติ ฯเปฯ. อิเม ธมฺมา กุสลา.}

\proseitem{187}{กตเม ธมฺมา กุสลา. ยสฺมิํ สมเย รูปูปปตฺติยา มคฺคํ ภาเวติ วิวิจฺเจว กาเมหิ ฯเปฯ ปฐมํ ฌานํ อุปสมฺปชฺช วิหรติ ทุกฺข\-ปฏิปทํ ทนฺธาภิญฺญํ ปริตฺตํ อปฺปมาณารมฺมาณํ ปถวีกสิณํ. ตสฺมิํ สมเย ผสฺโส โหติ ฯเปฯ อวิกฺเขโป โหติ ฯเปฯ. อิเม ธมฺมา กุสลา.}

\proseitem{188}{กตเม ธมฺมา กุสลา. ยสฺมิํ สมเย รูปูปปตฺติยา มคฺคํ ภาเวติ วิวิจฺเจว กาเมหิ ฯเปฯ ปฐมํ ฌานํ อุปสมฺปชฺช วิหรติ ทุกฺข\-ปฏิปทํ ทนฺธาภิญฺญํ อปฺปมาณํ ปริตฺตา\-รมฺมณํ ปถวีกสิณํ. ตสฺมิํ สมเย ผสฺโส โหติ ฯเปฯ อวิกฺเขโป โหติ ฯเปฯ. อิเม ธมฺมา กุสลา.}

\proseitem{189}{กตเม ธมฺมา กุสลา. ยสฺมิํ สมเย รูปูปปตฺติยา มคฺคํ ภาเวติ วิวิจฺเจว กาเมหิ ฯเปฯ ปฐมํ ฌานํ อุปสมฺปชฺช วิหรติ ทุกฺข\-ปฏิปทํ ทนฺธาภิญฺญํ อปฺปมาณํ อปฺปมาณา\-รมฺมณํ ปถวีกสิณํ. ตสฺมิํ สมเย ผสฺโส โหติ ฯเปฯ อวิกฺเขโป โหติ ฯเปฯ. อิเม ธมฺมา กุสลา.}

\proseitem{190}{กตเม ธมฺมา กุสลา. ยสฺมิํ สมเย รูปูปปตฺติยา มคฺคํ ภาเวติ วิวิจฺเจว กาเมหิ ฯเปฯ ปฐมํ ฌานํ อุปสมฺปชฺช วิหรติ ทุกฺข\-ปฏิปทํ \csromanfolio{53}ขิปฺปาภิญฺญํ ปริตฺตํ ปริตฺตา\-รมฺมณํ ปถวีกสิณํ. ตสฺมิํ สมเย ผสฺโส โหติ ฯเปฯ อวิกฺเขโป โหติ ฯเปฯ. อิเม ธมฺมา กุสลา.}

\proseitem{191}{กตเม ธมฺมา กุสลา. ยสฺมิํ สมเย รูปูปปตฺติยา มคฺคํ ภาเวติ วิวิจฺเจว กาเมหิ ฯเปฯ ปฐมํ ฌานํ อุปสมฺปชฺช วิหรติ ทุกฺข\-ปฏิปทํ ขิปฺปาภิญฺญํ ปริตฺตํ อปฺปมาณา\-รมฺมณํ ปถวีกสิณํ. ตสฺมิํ สมเย ผสฺโส โหติ ฯเปฯ อวิกฺเขโป โหติ ฯเปฯ. อิเม ธมฺมา กุสลา.}

\proseitem{192}{กตเม ธมฺมา กุสลา. ยสฺมิํ สมเย รูปูปปตฺติยา มคฺคํ ภาเวติ วิวิจฺเจว กาเมหิ ฯเปฯ ปฐมํ ฌานํ อุปสมฺปชฺช วิหรติ ทุกฺข\-ปฏิปทํ ขิปฺปาภิญฺญํ อปฺปมาณํ ปริตฺตา\-รมฺมณํ ปถวีกสิณํ. ตสฺมิํ สมเย ผสฺโส โหติ ฯเปฯ อวิกฺเขโป โหติ ฯเปฯ. อิเม ธมฺมา กุสลา.}

\proseitem{193}{กตเม ธมฺมา กุสลา. ยสฺมิํ สมเย รูปูปปตฺติยา มคฺคํ ภาเวติ วิวิจฺเจว กาเมหิ ฯเปฯ ปฐมํ ฌานํ อุปสมฺปชฺช วิหรติ ทุกฺข\-ปฏิปทํ ขิปฺปาภิญฺญํ อปฺปมาณํ อปฺปมาณา\-รมฺมณํ ปถวีกสิณํ. ตสฺมิํ สมเย ผสฺโส โหติ ฯเปฯ อวิกฺเขโป โหติ ฯเปฯ. อิเม ธมฺมา กุสลา.}

\proseitem{194}{กตเม ธมฺมา กุสลา. ยสฺมิํ สมเย รูปูปปตฺติยา มคฺคํ ภาเวติ วิวิจฺเจว กาเมหิ ฯเปฯ ปฐมํ ฌานํ อุปสมฺปชฺช วิหรติ สุขปฏิปทํ ทนฺธาภิญฺญํ ปริตฺตํ ปริตฺตา\-รมฺมณํ ปถวีกสิณํ. ตสฺมิํ สมเย ผสฺโส โหติ ฯเปฯ อวิกฺเขโป โหติ ฯเปฯ. อิเม ธมฺมา กุสลา.}

\proseitem{195}{กตเม ธมฺมา กุสลา. ยสฺมิํ สมเย รูปูปปตฺติยา มคฺคํ ภาเวติ วิวิจฺเจว กาเมหิ ฯเปฯ ปฐมํ ฌานํ อุปสมฺปชฺช วิหรติ สุขปฏิปทํ ทนฺธาภิญฺญํ ปริตฺตํ อปฺปมาณา\-รมฺมณํ ปถวีกสิณํ. ตสฺมิํ สมเย ผสฺโส โหติ ฯเปฯ อวิกฺเขโป โหติ ฯเปฯ. อิเม ธมฺมา กุสลา.}

\proseitem{196}{กตเม ธมฺมา กุสลา. ยสฺมิํ สมเย รูปูปปตฺติยา มคฺคํ ภาเวติ วิวิจฺเจว กาเมหิ ฯเปฯ ปฐมํ ฌานํ อุปสมฺปชฺช วิหรติ สุขปฏิปทํ ทนฺธาภิญฺญํ อปฺปมาณํ ปริตฺตา\-รมฺมณํ ปถวีกสิณํ. ตสฺมิํ สมเย ผสฺโส โหติ ฯเปฯ อวิกฺเขโป โหติ ฯเปฯ. อิเม ธมฺมา กุสลา.}

\proseitem{197}{กตเม ธมฺมา กุสลา. ยสฺมิํ สมเย รูปูปปตฺติยา มคฺคํ ภาเวติ วิวิจฺเจว กาเมหิ ฯเปฯ ปฐมํ ฌานํ อุปสมฺปชฺช วิหรติ สุขปฏิปทํ \csromanfolio{54}ทนฺธาภิญฺญํ อปฺปมาณํ อปฺปมาณา\-รมฺมณํ ปถวีกสิณํ. ตสฺมิํ สมเย ผสฺโส โหติ ฯเปฯ อวิกฺเขโป โหติ ฯเปฯ. อิเม ธมฺมา กุสลา.}

\proseitem{198}{กตเม ธมฺมา กุสลา. ยสฺมิํ สมเย รูปูปปตฺติยา มคฺคํ ภาเวติ วิวิจฺเจว กาเมหิ ฯเปฯ ปฐมํ ฌานํ อุปสมฺปชฺช วิหรติ สุขปฏิปทํ ขิปฺปาภิญฺญํ ปริตฺตํ ปริตฺตา\-รมฺมณํ ปถวีกสิณํ. ตสฺมิํ สมเย ผสฺโส โหติ ฯเปฯ อวิกฺเขโป โหติ ฯเปฯ อิเม ธมฺมา กุสลา.}

\proseitem{199}{กตเม ธมฺมา กุสลา. ยสฺมิํ สมเย รูปูปปตฺติยา มคฺคํ ภาเวติ วิวิจฺเจว กาเมหิ ฯเปฯ ปฐมํ ฌานํ อุปสมฺปชฺช วิหรติ สุขปฏิปทํ ขิปฺปาภิญฺญํ ปริตฺตํ อปฺปมาณา\-รมฺมณํ ปถวีกสิณํ. ตสฺมิํ สมเย ผสฺโส โหติ ฯเปฯ อวิกฺเขโป โหติ ฯเปฯ. อิเม ธมฺมา กุสลา.}

\proseitem{200}{กตเม ธมฺมา กุสลา. ยสฺมิํ สมเย รูปูปปตฺติยา มคฺคํ ภาเวติ วิวิจฺเจว กาเมหิ ฯเปฯ ปฐมํ ฌานํ อุปสมฺปชฺช วิหรติ สุขปฏิปทํ ขิปฺปาภิญฺญํ อปฺปมาณํ ปริตฺตา\-รมฺมณํ ปถวีกสิณํ. ตสฺมิํ สมเย ผสฺโส โหติ ฯเปฯ อวิกฺเขโป โหติ ฯเปฯ. อิเม ธมฺมา กุสลา.}

\proseitem{201}{กตเม ธมฺมา กุสลา. ยสฺมิํ สมเย รูปูปปตฺติยา มคฺคํ ภาเวติ วิวิจฺเจว กาเมหิ ฯเปฯ ปฐมํ ฌานํ อุปสมฺปชฺช วิหรติ สุขปฏิปทํ ขิปฺปาภิญฺญํ อปฺปมาณํ อปฺปมาณา\-รมฺมณํ ปถวีกสิณํ. ตสฺมิํ สมเย ผสฺโส โหติ ฯเปฯ อวิกฺเขโป โหติ ฯเปฯ. อิเม ธมฺมา กุสลา.}

\proseitem{202}{กตเม ธมฺมา กุสลา. ยสฺมิํ สมเย รูปูปปตฺติยา มคฺคํ ภาเวติ วิตกฺก\-วิจารานํ วูปสมา ฯเปฯ ทุติยํ ฌานํ ฯเปฯ ตติยํ ฌานํ ฯเปฯ จตุตฺถํ ฌานํ ฯเปฯ ปฐมํ ฌานํ ฯเปฯ ปญฺจมํ ฌานํ อุปสมฺปชฺช วิหรติ ทุกฺข\-ปฏิปทํ ทนฺธาภิญฺญํ ปริตฺตํ ปริตฺตา\-รมฺมณํ ปถวีกสิณํ ฯเปฯ ทุกฺข\-ปฏิปทํ ทนฺธาภิญฺญํ ปริตฺตํ อปฺปมาณา\-รมฺมณํ ปถวีกสิณํ ฯเปฯ ทุกฺข\-ปฏิปทํ ทนฺธาภิญฺญํ อปฺปมาณํ ปริตฺตา\-รมฺมณํ ปถวีกสิณํ ฯเปฯ ทุกฺข\-ปฏิปทํ ทนฺธาภิญฺญํ อปฺปมาณํ อปฺปมาณา\-รมฺมณํ ปถวีกสิณํ ฯเปฯ ทุกฺข\-ปฏิปทํ ขิปฺปาภิญฺญํ ปริตฺตํ ปริตฺตา\-รมฺมณํ ปถวีกสิณํ ฯเปฯ ทุกฺข\-ปฏิปทํ ขิปฺปาภิญฺญํ ปริตฺตํ อปฺปมาณา\-รมฺมณํ ปถวีกสิณํ ฯเปฯ ทุกฺข\-ปฏิปทํ ขิปฺปาภิญฺญํ อปฺปมาณํ ปริตฺตา\-รมฺมณํ ปถวีกสิณํ ฯเปฯ ทุกฺข\-ปฏิปทํ ขิปฺปาภิญฺญํ อปฺปมาณํ อปฺปมาณา\-รมฺมณํ ปถวีกสิณํ ฯเปฯ สุขปฏิปทํ ทนฺธาภิญฺญํ \csromanfolio{55}ปริตฺตํ ปริตฺตา\-รมฺมณํ ปถวีกสิณํ ฯเปฯ สุขปฏิปทํ ทนฺธาภิญฺญํ ปริตฺตํ อปฺปมาณา\-รมฺมณํ ปถวีกสิณํ ฯเปฯ สุขปฏิปทํ ทนฺธาภิญฺญํ อปฺปมาณํ ปริตฺตา\-รมฺมณํ ปถวีกสิณํ ฯเปฯ สุขปฏิปทํ ทนฺธาภิญฺญํ อปฺปมาณํ อปฺปมาณา\-รมฺมณํ ปถวีกสิณํ ฯเปฯ สุขปฏิปทํ ขิปฺปาภิญฺญํ ปริตฺตํ ปริตฺตา\-รมฺมณํ ปถวีกสิณํ ฯเปฯ สุขปฏิปทํ ขิปฺปาภิญฺญํ ปริตฺตํ อปฺปมาณา\-รมฺมณํ ปถวีกสิณํ ฯเปฯ สุขปฏิปทํ ขิปฺปาภิญฺญํ อปฺปมาณํ ปริตฺตา\-รมฺมณํ ปถวีกสิณํ ฯเปฯ สุขปฏิปทํ ขิปฺปาภิญฺญํ อปฺปมาณํ อปฺปมาณา\-รมฺมณํ ปถวีกสิณํ. ตสฺมิํ สมเย ผสฺโส โหติ ฯเปฯ อวิกฺเขโป โหติ ฯเปฯ. อิเม ธมฺมา กุสลา.}

\vspace{\tipitakaparskip}
\csromancenter{โสฬส\-กฺขตฺตุกํ.}
\csromanmatikaanchor{mtk.26}
\csromantocmarkref{section}{อฏฺฐกสิณํ โสฬสกฺขตฺตุกํ}{mtk.26}
\bookmark[dest={mtk.26},level=1]{อฏฺฐกสิณํ โสฬสกฺขตฺตุกํ}
\csromanheader{1}{อฏฺฐกสิณํ โสฬส\-กฺขตฺตุกํ}
\proseitem{203}{กตเม ธมฺมา กุสลา. ยสฺมิํ สมเย รูปูปปตฺติยา มคฺคํ ภาเวติ วิวิจฺเจว กาเมหิ ฯเปฯ ปฐมํ ฌานํ อุปสมฺปชฺช วิหรติ อาโปกสิณํ ฯเปฯ เตโชกสิณํ ฯเปฯ วาโยกสิณํ ฯเปฯ นีลกสิณํ ฯเปฯ ปีตกสิณํ ฯเปฯ โลหิตกสิณํ ฯเปฯ โอทาตกสิณํ. ตสฺมิํ สมเย ผสฺโส โหติ ฯเปฯ อวิกฺเขโป โหติ ฯเปฯ. อิเม ธมฺมา กุสลา.}

\vspace{\tipitakaparskip}
\csromancenter{อฏฺฐกสิณํ โสฬส\-กฺขตฺตุกํ.}
\csromanmatikaanchor{mtk.27}
\csromantocmarkref{section}{อภิภายตนานิ ปริตฺตานิ}{mtk.27}
\bookmark[dest={mtk.27},level=1]{อภิภายตนานิ ปริตฺตานิ}
\csromanheader{1}{อภิภา\-ยตนานิ ปริตฺตานิ}
\proseitem{204}{กตเม ธมฺมา กุสลา. ยสฺมิํ สมเย รูปูปปตฺติยา มคฺคํ ภาเวติ อชฺฌตฺตํ อรูปสญฺญี พหิทฺธา รูปานิ ปสฺสติ ปริตฺตานิ ตานิ อภิภุยฺย ชานามิ ปสฺสามีติ วิวิจฺเจว กาเมหิ ฯเปฯ ปฐมํ ฌานํ อุปสมฺปชฺช วิหรติ. ตสฺมิํ สมเย ผสฺโส โหติ ฯเปฯ อวิกฺเขโป โหติ ฯเปฯ. อิเม ธมฺมา กุสลา.}

\proseitem{205}{กตเม ธมฺมา กุสลา. ยสฺมิํ สมเย รูปูปปตฺติยา มคฺคํ ภาเวติ อชฺฌตฺตํ อรูปสญฺญี พหิทฺธา รูปานิ ปสฺสติ ปริตฺตานิ ตานิ อภิภุยฺย ชานามิ ปสฺสามีติ วิตกฺก\-วิจารานํ วูปสมา ฯเปฯ ทุติยํ \csromanfolio{56}ฌานํ ฯเปฯ ตติยํ ฌานํ ฯเปฯ จตุตฺถํ ฌานํ ฯเปฯ ปฐมํ ฌานํ ฯเปฯ ปญฺจมํ ฌานํ อุปสมฺปชฺช วิหรติ. ตสฺมิํ สมเย ผสฺโส โหติ ฯเปฯ อวิกฺเขโป โหติ ฯเปฯ. อิเม ธมฺมา กุสลา.}

\csromanmatikaanchor{mtk.28}
\csromantocmarkref{section}{จตสฺโส ปฏิปทา}{mtk.28}
\bookmark[dest={mtk.28},level=1]{จตสฺโส ปฏิปทา}
\csromanheader{1}{\textbf{จตสฺโส ปฏิปทา}}
\proseitem{206}{กตเม ธมฺมา กุสลา. ยสฺมิํ สมเย รูปูปปตฺติยา มคฺคํ ภาเวติ อชฺฌตฺตํ อรูปสญฺญี พหิทฺธา รูปานิ ปสฺสติ ปริตฺตานิ ตานิ อภิภุยฺย ชานามิ ปสฺสามีติ วิวิจฺเจว กาเมหิ ฯเปฯ ปฐมํ ฌานํ อุปสมฺปชฺช วิหรติ ทุกฺข\-ปฏิปทํ ทนฺธาภิญฺญํ. ตสฺมิํ สมเย ผสฺโส โหติ ฯเปฯ อวิกฺเขโป โหติ ฯเปฯ. อิเม ธมฺมา กุสลา.}

\proseitem{207}{กตเม ธมฺมา กุสลา. ยสฺมิํ สมเย รูปูปปตฺติยา มคฺคํ ภาเวติ อชฺฌตฺตํ อรูปสญฺญี พหิทฺธา รูปานิ ปสฺสติ ปริตฺตานิ ตานิ อภิภุยฺย ชานามิ ปสฺสามีติ วิวิจฺเจว กาเมหิ ฯเปฯ ปฐมํ ฌานํ อุปสมฺปชฺช วิหรติ ทุกฺข\-ปฏิปทํ ขิปฺปาภิญฺญํ. ตสฺมิํ สมเย ผสฺโส โหติ ฯเปฯ อวิกฺเขโป โหติ ฯเปฯ. อิเม ธมฺมา กุสลา.}

\proseitem{208}{กตเม ธมฺมา กุสลา. ยสฺมิํ สมเย รูปูปปตฺติยา มคฺคํ ภาเวติ อชฺฌตฺตํ อรูปสญฺญี พหิทฺธา รูปานิ ปสฺสติ ปริตฺตานิ ตานิ อภิภุยฺย ชานามิ ปสฺสามีติ วิวิจฺเจว กาเมหิ ฯเปฯ ปฐมํ ฌานํ อุปสมฺปชฺช วิหรติ สุขปฏิปทํ ทนฺธาภิญฺญํ. ตสฺมิํ สมเย ผสฺโส โหติ ฯเปฯ อวิกฺเขโป โหติ ฯเปฯ. อิเม ธมฺมา กุสลา.}

\proseitem{209}{กตเม ธมฺมา กุสลา. ยสฺมิํ สมเย รูปูปปตฺติยา มคฺคํ ภาเวติ อชฺฌตฺตํ อรูปสญฺญี พหิทฺธา รูปานิ ปสฺสติ ปริตฺตานิ ตานิ อภิภุยฺย ชานามิ ปสฺสามีติ วิวิจฺเจว กาเมหิ ฯเปฯ ปฐมํ ฌานํ อุปสมฺปชฺช วิหรติ สุขปฏิปทํ ขิปฺปาภิญฺญํ. ตสฺมิํ สมเย ผสฺโส โหติ ฯเปฯ อวิกฺเขโป โหติ ฯเปฯ. อิเม ธมฺมา กุสลา.}

\proseitem{210}{กตเม ธมฺมา กุสลา. ยสฺมิํ สมเย รูปูปปตฺติยา มคฺคํ ภาเวติ อชฺฌตฺตํ อรูปสญฺญี พหิทฺธา รูปานิ ปสฺสติ ปริตฺตานิ ตานิ อภิภุยฺย ชานามิ ปสฺสามีติ วิตกฺก\-วิจารานํ วูปสมา ฯเปฯ ทุติยํ ฌานํ ฯเปฯ ตติยํ ฌานํ ฯเปฯ จตุตฺถํ ฌานํ ฯเปฯ ปฐมํ ฌานํ ฯเปฯ ปญฺจมํ ฌานํ อุปสมฺปชฺช \csromanfolio{57}วิหรติ ทุกฺข\-ปฏิปทํ ทนฺธาภิญฺญํ ฯเปฯ ทุกฺข\-ปฏิปทํ ขิปฺปาภิญฺญํ ฯเปฯ สุขปฏิปทํ ทนฺธาภิญฺญํ ฯเปฯ สุขปฏิปทํ ขิปฺปาภิญฺญํ. ตสฺมิํ สมเย ผสฺโส โหติ ฯเปฯ อวิกฺเขโป โหติ ฯเปฯ. อิเม ธมฺมา กุสลา. จตฺตสฺโส ปฏิปทา.}

\csromanmatikaanchor{mtk.29}
\csromantocmarkref{section}{เทฺว อารมฺมณานิ}{mtk.29}
\bookmark[dest={mtk.29},level=1]{เทฺว อารมฺมณานิ}
\csromanheader{1}{\textbf{เทฺว อารมฺมณานิ}}
\proseitem{211}{กตเม ธมฺมา กุสลา. ยสฺมิํ สมเย รูปูปปตฺติยา มคฺคํ ภาเวติ อชฺฌตฺตํ อรูปสญฺญี พหิทฺธา รูปานิ ปสฺสติ ปริตฺตานิ ตานิ อภิภุยฺย ชานามิ ปสฺสามีติ วิวิจฺเจว กาเมหิ ฯเปฯ ปฐมํ ฌานํ อุปสมฺปชฺช วิหรติ ปริตฺตํ ปริตฺตา\-รมฺมณํ. ตสฺมิํ สมเย ผสฺโส โหติ ฯเปฯ อวิกฺเขโป โหติ ฯเปฯ. อิเม ธมฺมา กุสลา.}

\proseitem{212}{กตเม ธมฺมา กุสลา. ยสฺมิํ สมเย รูปูปปตฺติยา มคฺคํ ภาเวติ อชฺฌตฺตํ อรูปสญฺญี พหิทฺธา รูปานิ ปสฺสติ ปริตฺตานิ ตานิ อภิภุยฺย ชานามิ ปสฺสามีติ วิวิจฺเจว กาเมหิ ฯเปฯ ปฐมํ ฌานํ อุปสมฺปชฺช วิหรติ อปฺปมาณํ ปริตฺตา\-รมฺมณํ. ตสฺมิํ สมเย ผสฺโส โหติ ฯเปฯ อวิกฺเขโป โหติ ฯเปฯ. อิเม ธมฺมา กุสลา.}

\proseitem{213}{กตเม ธมฺมา กุสลา. ยสฺมิํ สมเย รูปูปปตฺติยา มคฺคํ ภาเวติ อชฺฌตฺตํ อรูปสญฺญี พหิทฺธา รูปานิ ปสฺสติ ปริตฺตานิ ตานิ อภิภุยฺย ชานามิ ปสฺสามีติ วิตกฺก\-วิจารานํ วูปสมา ฯเปฯ ทุติยํ ฌานํ ฯเปฯ ตติยํ ฌานํ ฯเปฯ จตุตฺถํ ฌานํ ฯเปฯ ปฐมํ ฌานํ ฯเปฯ ปญฺจมํ ฌานํ อุปสมฺปชฺช วิหรติ ปริตฺตํ ปริตฺตา\-รมฺมณํ ฯเปฯ อปฺปมาณํ ปริตฺตา\-รมฺมณํ. ตสฺมิํ สมเย ผสฺโส โหติ ฯเปฯ อวิกฺเขโป โหติ ฯเปฯ. อิเม ธมฺมา กุสลา.}

\vspace{\tipitakaparskip}
\csromancenter{\textbf{เทฺว อารมฺมณานิ}.}
\csromanmatikaanchor{mtk.30}
\csromantocmarkref{section}{อฏฺฐกฺขตฺตุกํ}{mtk.30}
\bookmark[dest={mtk.30},level=1]{อฏฺฐกฺขตฺตุกํ}
\csromanheader{1}{อฏฺฐ\-กฺขตฺตุกํ}
\proseitem{214}{กตเม ธมฺมา กุสลา. ยสฺมิํ สมเย รูปูปปตฺติยา มคฺคํ ภาเวติ อชฺฌตฺตํ อรูปสญฺญี พหิทฺธา รูปานิ ปสฺสติ ปริตฺตานิ ตานิ \csromanfolio{58}อภิภุยฺย ชานามิ ปสฺสามีติ วิวิจฺเจว กาเมหิ ฯเปฯ ปฐมํ ฌานํ อุปสมฺปชฺช วิหรติ ทุกฺข\-ปฏิปทํ ทนฺธาภิญฺญํ ปริตฺตํ ปริตฺตา\-รมฺมณํ. ตสฺมิํ สมเย ผสฺโส โหติ ฯเปฯ อวิกฺเขโป โหติ ฯเปฯ. อิเม ธมฺมา กุสลา.}

\proseitem{215}{กตเม ธมฺมา กุสลา. ยสฺมิํ สมเย รูปูปปตฺติยา มคฺคํ ภาเวติ อชฺฌตฺตํ อรูปสญฺญี พหิทฺธา รูปานิ ปสฺสติ ปริตฺตานิ ตานิ อภิภุยฺย ชานามิ ปสฺสามีติ วิวิจฺเจว กาเมหิ ฯเปฯ ปฐมํ ฌานํ อุปสมฺปชฺช วิหรติ ทุกฺข\-ปฏิปทํ ทนฺธาภิญฺญํ อปฺปมาณํ ปริตฺตา\-รมฺมณํ. ตสฺมิํ สมเย ผสฺโส โหติ ฯเปฯ อวิกฺเขโป โหติ ฯเปฯ. อิเม ธมฺมา กุสลา.}

\proseitem{216}{กตเม ธมฺมา กุสลา. ยสฺมิํ สมเย รูปูปปตฺติยา มคฺคํ ภาเวติ อชฺฌตฺตํ อรูปสญฺญี พหิทฺธา รูปานิ ปสฺสติ ปริตฺตานิ ตานิ อภิภุยฺย ชานามิ ปสฺสามีติ วิวิจฺเจว กาเมหิ ฯเปฯ ปฐมํ ฌานํ อุปสมฺปชฺช วิหรติ ทุกฺข\-ปฏิปทํ ขิปฺปาภิญฺญํ ปริตฺตํ ปริตฺตา\-รมฺมณํ. ตสฺมิํ สมเย ผสฺโส โหติ ฯเปฯ อวิกฺเขโป โหติ ฯเปฯ. อิเม ธมฺมา กุสลา.}

\proseitem{217}{กตเม ธมฺมา กุสลา. ยสฺมิํ สมเย รูปูปปตฺติยา มคฺคํ ภาเวติ อชฺฌตฺตํ อรูปสญฺญี พหิทฺธา รูปานิ ปสฺสติ ปริตฺตานิ ตานิ อภิภุยฺย ชานามิ ปสฺสามีติ วิวิจฺเจว กาเมหิ ฯเปฯ ปฐมํ ฌานํ อุปสมฺปชฺช วิหรติ ทุกฺข\-ปฏิปทํ ขิปฺปาภิญฺญํ อปฺปมาณํ ปริตฺตา\-รมฺมณํ. ตสฺมิํ สมเย ผสฺโส โหติ ฯเปฯ อวิกฺเขโป โหติ ฯเปฯ. อิเม ธมฺมา กุสลา.}

\proseitem{218}{กตเม ธมฺมา กุสลา. ยสฺมิํ สมเย รูปูปปตฺติยา มคฺคํ ภาเวติ อชฺฌตฺตํ อรูปสญฺญี พหิทฺธา รูปานิ ปสฺสติ ปริตฺตานิ ตานิ อภิภุยฺย ชานามิ ปสฺสามีติ วิวิจฺเจว กาเมหิ ฯเปฯ ปฐมํ ฌานํ อุปสมฺปชฺช วิหรติ สุขปฏิปทํ ทนฺธาภิญฺญํ ปริตฺตํ ปริตฺตา\-รมฺมณํ. ตสฺมิํ สมเย ผสฺโส โหติ ฯเปฯ อวิกฺเขโป โหติ ฯเปฯ. อิเม ธมฺมา กุสลา.}

\proseitem{219}{กตเม ธมฺมา กุสลา. ยสฺมิํ สมเย รูปูปปตฺติยา มคฺคํ ภาเวติ อชฺฌตฺตํ อรูปสญฺญี พหิทฺธา รูปานิ ปสฺสติ ปริตฺตานิ ตานิ อภิภุยฺย ชานามิ ปสฺสามีติ วิวิจฺเจว กาเมหิ ฯเปฯ ปฐมํ ฌานํ อุปสมฺปชฺช วิหรติ สุขปฏิปทํ ทนฺธาภิญฺญํ อปฺปมาณํ ปริตฺตา\-รมฺมณํ. ตสฺมิํ สมเย ผสฺโส โหติ ฯเปฯ อวิกฺเขโป โหติ ฯเปฯ. อิเม ธมฺมา กุสลา.}

\chapterheadpageread{1. จิตฺตุปฺปาท\-กณฺฑ}
\proseitem{220}{\csromanfolio{59}กตเม ธมฺมา กุสลา. ยสฺมิํ สมเย รูปูปปตฺติยา มคฺคํ ภาเวติ อชฺฌตฺตํ อรูปสญฺญี พหิทฺธา รูปานิ ปสฺสติ ปริตฺตานิ ตานิ อภิภุยฺย ชานามิ ปสฺสามีติ วิวิจฺเจว กาเมหิ ฯเปฯ ปฐมํ ฌานํ อุปสมฺปชฺช วิหรติ สุขปฏิปทํ ขิปฺปาภิญฺญํ ปริตฺตํ ปริตฺตา\-รมฺมณํ. ตสฺมิํ สมเย ผสฺโส โหติ ฯเปฯ อวิกฺเขโป โหติ ฯเปฯ. อิเม ธมฺมา กุสลา.}

\proseitem{221}{กตเม ธมฺมา กุสลา. ยสฺมิํ สมเย รูปูปปตฺติยา มคฺคํ ภาเวติ อชฺฌตฺตํ อรูปสญฺญี พหิทฺธา รูปานิ ปสฺสติ ปริตฺตานิ ตานิ อภิภุยฺย ชานามิ ปสฺสามีติ วิวิจฺเจว กาเมหิ ฯเปฯ ปฐมํ ฌานํ อุปสมฺปชฺช วิหรติ สุขปฏิปทํ ขิปฺปาภิญฺญํ อปฺปมาณํ ปริตฺตา\-รมฺมณํ. ตสฺมิํ สมเย ผสฺโส โหติ ฯเปฯ อวิกฺเขโป โหติ ฯเปฯ. อิเม ธมฺมา กุสลา.}

\proseitem{222}{กตเม ธมฺมา กุสลา. ยสฺมิํ สมเย รูปูปปตฺติยา มคฺคํ ภาเวติ อชฺฌตฺตํ อรูปสญฺญี พหิทฺธา รูปานิ ปสฺสติ ปริตฺตานิ ตานิ อภิภุยฺย ชานามิ ปสฺสามีติ วิตกฺก\-วิจารานํ วูปสมา ฯเปฯ ทุติยํ ฌานํ ฯเปฯ ตติยํ ฌานํ ฯเปฯ จตุตฺถํ ฌานํ ฯเปฯ ปฐมํ ฌานํ ฯเปฯ ปญฺจมํ ฌานํ อุปสมฺปชฺช วิหรติ ทุกฺข\-ปฏิปทํ ทนฺธาภิญฺญํ ปริตฺตํ ปริตฺตา\-รมฺมณํ ฯเปฯ ทุกฺข\-ปฏิปทํ ทนฺธาภิญฺญํ อปฺปมาณํ ปริตฺตา\-รมฺมณํ ฯเปฯ ทุกฺข\-ปฏิปทํ ขิปฺปาภิญฺญํ ปริตฺตํ ปริตฺตา\-รมฺมณํ ฯเปฯ ทุกฺข\-ปฏิปทํ ขิปฺปาภิญฺญํ อปฺปมาณํ ปริตฺตา\-รมฺมณํ ฯเปฯ สุขปฏิปทํ ทนฺธาภิญฺญํ ปริตฺตํ ปริตฺตา\-รมฺมณํ ฯเปฯ สุขปฏิปทํ ทนฺธาภิญฺญํ อปฺปมาณํ ปริตฺตา\-รมฺมณํ ฯเปฯ สุขปฏิปทํ ขิปฺปาภิญฺญํ ปริตฺตํ ปริตฺตา\-รมฺมณํ ฯเปฯ สุขปฏิปทํ ขิปฺปาภิญฺญํ อปฺปมาณํ ปริตฺตา\-รมฺมณํ. ตสฺมิํ สมเย ผสฺโส โหติ ฯเปฯ อวิกฺเขโป โหติ ฯเปฯ. อิเม ธมฺมา กุสลา. อฏฺฐ\-กฺขตฺตุกํ.}

\csromanmatikaanchor{mtk.31}
\csromantocmarkref{section}{อิทมฺปิ อฏฺฐกฺขตฺตุกํ}{mtk.31}
\bookmark[dest={mtk.31},level=1]{อิทมฺปิ อฏฺฐกฺขตฺตุกํ}
\csromanheader{1}{อิทมฺปิ อฏฺฐ\-กฺขตฺตุกํ}
\proseitem{223}{กตเม ธมฺมา กุสลา. ยสฺมิํ สมเย รูปูปปตฺติยา มคฺคํ ภาเวติ อชฺฌตฺตํ อรูปสญฺญี พหิทฺธา รูปานิ ปสฺสติ ปริตฺตานิ สุวณฺณ\-ทุพฺพณฺณานิ ตานิ อภิภุยฺย ชานามิ ปสฺสามีติ วิวิจฺเจว กาเมหิ ฯเปฯ ปฐมํ ฌานํ อุปสมฺปชฺช วิหรติ. ตสฺมิํ สมเย ผสฺโส โหติ ฯเปฯ อวิกฺเขโป โหติ ฯเปฯ. อิเม ธมฺมา กุสลา.}

\proseitem{224}{\csromanfolio{60}กตเม ธมฺมา กุสลา. ยสฺมิํ สมเย รูปูปปตฺติยา มคฺคํ ภาเวติ อชฺฌตฺตํ อรูปสญฺญี พหิทฺธา รูปานิ ปสฺสติ ปริตฺตานิ สุวณฺณ\-ทุพฺพณฺณานิ ตานิ อภิภุยฺย ชานามิ ปสฺสามีติ วิตกฺก\-วิจารานํ วูปสมา ฯเปฯ ทุติยํ ฌานํ ฯเปฯ ตติยํ ฌานํ ฯเปฯ จตุตฺถํ ฌานํ ฯเปฯ ปฐมํ ฌานํ ฯเปฯ ปญฺจมํ ฌานํ อุปสมฺปชฺช วิหรติ. ตสฺมิํ สมเย ผสฺโส โหติ ฯเปฯ อวิกฺเขโป โหติ ฯเปฯ. อิเม ธมฺมา กุสลา.}

\vspace{\tipitakaparskip}
\csromancenter{อิทมฺปิ อฏฺฐ\-กฺขตฺตุกํ.}
\csromanmatikaanchor{mtk.32}
\csromantocmarkref{section}{อปฺปมาณานิ}{mtk.32}
\bookmark[dest={mtk.32},level=1]{อปฺปมาณานิ}
\csromanheader{1}{\textbf{อปฺปมาณานิ}}
\proseitem{225}{กตเม ธมฺมา กุสลา. ยสฺมิํ สมเย รูปูปปตฺติยา มคฺคํ ภาเวติ อชฺฌตฺตํ อรูปสญฺญี พหิทฺธา รูปานิ ปสฺสติ อปฺปมาณานิ ตานิ อภิภุยฺย ชานามิ ปสฺสามีติ วิวิจฺเจว กาเมหิ ฯเปฯ ปฐมํ ฌานํ อุปสมฺปชฺช วิหรติ. ตสฺมิํ สมเย ผสฺโส โหติ ฯเปฯ อวิกฺเขโป โหติ ฯเปฯ. อิเม ธมฺมา กุสลา.}

\proseitem{226}{กตเม ธมฺมา กุสลา. ยสฺมิํ สมเย รูปูปปตฺติยา มคฺคํ ภาเวติ อชฺฌตฺตํ อรูปสญฺญี พหิทฺธา รูปานิ ปสฺสติ อปฺปมาณานิ ตานิ อภิภุยฺย ชานามิ ปสฺสามีติ วิตกฺก\-วิจารานํ วูปสมา ฯเปฯ ทุติยํ ฌานํ ฯเปฯ ตติยํ ฌานํ ฯเปฯ จตุตฺถํ ฌานํ ฯเปฯ ปฐมํ ฌานํ ฯเปฯ ปญฺจมํ ฌานํ อุปสมฺปชฺช วิหรติ. ตสฺมิํ สมเย ผสฺโส โหติ ฯเปฯ อวิกฺเขโป โหติ ฯเปฯ. อิเม ธมฺมา กุสลา.}

\csromanmatikaanchor{mtk.33}
\csromantocmarkref{section}{จตสฺโส ปฏิปทา}{mtk.33}
\bookmark[dest={mtk.33},level=1]{จตสฺโส ปฏิปทา}
\csromanheader{1}{\textbf{จตสฺโส ปฏิปทา}}
\proseitem{227}{กตเม ธมฺมา กุสลา. ยสฺมิํ สมเย รูปูปปตฺติยา มคฺคํ ภาเวติ อชฺฌตฺตํ อรูปสญฺญี พหิทฺธา รูปานิ ปสฺสติ อปฺปมาณานิ ตานิ อภิภุยฺย ชานามิ ปสฺสามีติ วิวิจฺเจว กิเมหิ ฯเปฯ ปฐมํ ฌานํ อุปสมฺปชฺช วิหรติ ทุกฺข\-ปฏิปทํ ทนฺธาภิญฺญํ. ตสฺมิํ สมเย ผสฺโส โหติ ฯเปฯ อวิกฺเขโป โหติ ฯเปฯ. อิเม ธมฺมา กุสลา.}

\proseitem{228}{กตเม ธมฺมา กุสลา. ยสฺมิํ สมเย รูปูปปตฺติยา มคฺคํ ภาเวติ อชฺฌตฺตํ อรูปสญฺญี พหิทฺธา รูปานิ ปสฺสติ อปฺปมาณานิ ตานิ \csromanfolio{61}อภิภุยฺย ชานามิ ปสฺสามีติ วิวิจฺเจว กาเมหิ ฯเปฯ ปฐมํ ฌานํ อุปสมฺปชฺช วิหรติ ทุกฺข\-ปฏิปทํ ขิปฺปาภิญฺญํ. ตสฺมิํ สมเย ผสฺโส โหติ ฯเปฯ อวิกฺเขโป โหติ ฯเปฯ. อิเม ธมฺมา กุสลา.}

\proseitem{229}{กตเม ธมฺมา กุสลา. ยสฺมิํ สมเย รูปูปปตฺติยา มคฺคํ ภาเวติ อชฺฌตฺตํ อรูปสญฺญี พหิทฺธา รูปานิ ปสฺสติ อปฺปมาณานิ ตานิ อภิภุยฺย ชานามิ ปสฺสามีติ วิวิจฺเจว กาเมหิ ฯเปฯ ปฐมํ ฌานํ อุปสมฺปชฺช วิหรติ สุขปฏิปทํ ทนฺธาภิญฺญํ. ตสฺมิํ สมเย ผสฺโส โหติ ฯเปฯ อวิกฺเขโป โหติ ฯเปฯ. อิเม ธมฺมา กุสลา.}

\proseitem{230}{กตเม ธมฺมา กุสลา. ยสฺมิํ สมเย รูปูปปตฺติยา มคฺคํ ภาเวติ อชฺฌตฺตํ อรูปสญฺญี พหิทฺธา รูปานิ ปสฺสติ อปฺปมาณานิ ตานิ อภิภุยฺย ชานามิ ปสฺสามีติ วิวิจฺเจว กาเมหิ ฯเปฯ ปฐมํ ฌานํ อุปสมฺปชฺช วิหรติ สุขปฏิปทํ ขิปฺปาภิญฺญํ. ตสฺมิํ สมเย ผสฺโส โหติ ฯเปฯ อวิกฺเขโป โหติ ฯเปฯ. อิเม ธมฺมา กุสลา.}

\proseitem{231}{กตเม ธมฺมา กุสลา. ยสฺมิํ สมเย รูปูปปตฺติยา มคฺคํ ภาเวติ อชฺฌตฺตํ อรูปสญฺญี พหิทฺธา รูปานิ ปสฺสติ อปฺปมาณานิ ตานิ อภิภุยฺย ชานามิ ปสฺสามีติ วิตกฺก\-วิจารานํ วูปสมา ฯเปฯ ทุติยํ ฌานํ ฯเปฯ ตติยํ ฌานํ ฯเปฯ จตุตฺถํ ฌานํ ฯเปฯ ปฐมํ ฌานํ ฯเปฯ ปญฺจมํ ฌานํ อุปสมฺปชฺช วิหรติ ทุกฺข\-ปฏิปทํ ทนฺธาภิญฺญํ ฯเปฯ ทุกฺข\-ปฏิปทํ ขิปฺปาภิญฺญํ ฯเปฯ สุขปฏิปทํ ทนฺธาภิญฺญํ ฯเปฯ สุขปฏิปทํ ขิปฺปาภิญฺญํ. ตสฺมิํ สมเย ผสฺโส โหติ ฯเปฯ อวิกฺเขโป โหติ ฯเปฯ. อิเม ธมฺมา กุสลา.}

\vspace{\tipitakaparskip}
\csromancenter{จตสฺโส ปฏิปทา.}
\csromanmatikaanchor{mtk.34}
\csromantocmarkref{section}{เทฺว อารมฺมณานิ}{mtk.34}
\bookmark[dest={mtk.34},level=1]{เทฺว อารมฺมณานิ}
\csromanheader{1}{\textbf{เทฺว อารมฺมณานิ}}
\proseitem{232}{กตเม ธมฺมา กุสลา. ยสฺมิํ สมเย รูปูปปตฺติยา มคฺคํ ภาเวติ อชฺฌตฺตํ อรูปสญฺญี พหิทฺธา รูปานิ ปสฺสติ อปฺปมาณานิ ตานิ อภิภุยฺย ชานามิ ปสฺสามีติ วิวิจฺเจว กาเมหิ ฯเปฯ ปฐมํ ฌานํ อุปสมฺปชฺช วิหรติ ปริตฺตํ อปฺปมาณา\-รมฺมณํ. ตสฺมิํ สมเย ผสฺโส โหติ ฯเปฯ อวิกฺเขโป โหติ ฯเปฯ. อิเม ธมฺมา กุสลา.}

\proseitem{233}{\csromanfolio{62}กตเม ธมฺมา กุสลา. ยสฺมิํ สมเย รูปูปปตฺติยา มคฺคํ ภาเวติ อชฺฌตฺตํ อรูปสญฺญี พหิทฺธา รูปานิ ปสฺสติ อปฺปมาณานิ ตานิ อภิภุยฺย ชานามิ ปสฺสามีติ วิวิจฺเจว กาเมหิ ฯเปฯ ปฐมํ ฌานํ อุปสมฺปชฺช วิหรติ อปฺปมาณํ อปฺปมาณา\-รมฺมณํ. ตสฺมิํ สมเย ผสฺโส โหติ ฯเปฯ อวิกฺเขโป โหติ ฯเปฯ. อิเม ธมฺมา กุสลา.}

\proseitem{234}{กตเม ธมฺมา กุสลา. ยสฺมิํ สมเย รูปูปปตฺติยา มคฺคํ ภาเวติ อชฺฌตฺตํ อรูปสญฺญี พหิทฺธา รูปานิ ปสฺสติ อปฺปมาณานิ ตานิ อภิภุยฺย ชานามิ ปสฺสามีติ วิตกฺก\-วิจารานํ วูปสมา ฯเปฯ ทุติยํ ฌานํ ฯเปฯ ตติยํ ฌานํ ฯเปฯ จตุตฺถํ ฌานํ ฯเปฯ ปฐมํ ฌานํ ฯเปฯ ปญฺจมํ ฌานํ อุปสมฺปชฺช วิหรติ ปริตฺตํ อปฺปมาณา\-รมฺมณํ ฯเปฯ อปฺปมาณํ อปฺปมาณา\-รมฺมณํ. ตสฺมิํ สมเย ผสฺโส โหติ ฯเปฯ อวิกฺเขโป โหติ ฯเปฯ. อิเม ธมฺมา กุสลา.}

\vspace{\tipitakaparskip}
\csromancenter{เทฺว อารมฺมณานิ.}
\csromanmatikaanchor{mtk.35}
\csromantocmarkref{section}{อปรมฺปิ อฏฺฐกฺขตฺตุกํ}{mtk.35}
\bookmark[dest={mtk.35},level=1]{อปรมฺปิ อฏฺฐกฺขตฺตุกํ}
\csromanheader{1}{อปรมฺปิ อฏฺฐ\-กฺขตฺตุกํ}
\proseitem{235}{กตเม ธมฺมา กุสลา. ยสฺมิํ สมเย รูปูปปตฺติยา มคฺคํ ภาเวติ อชฺฌตฺตํ อรูปสญฺญี พหิทฺธา รูปานิ ปสฺสติ อปฺปมาณานิ ตานิ อภิภุยฺย ชานามิ ปสฺสามีติ วิวิจฺเจว กาเมหิ ฯเปฯ ปฐมํ ฌานํ อุปสมฺปชฺช วิหรติ ทุกฺข\-ปฏิปทํ ทนฺธาภิญฺญํ ปริตฺตํ อปฺปมาณา\-รมฺมณํ. ตสฺมิํ สมเย ผสฺโส โหติ ฯเปฯ อวิกฺเขโป โหติ ฯเปฯ. อิเม ธมฺมา กุสลา.}

\proseitem{236}{กตเม ธมฺมา กุสลา. ยสฺมิํ สมเย รูปูปปตฺติยา มคฺคํ ภาเวติ อชฺฌตฺตํ อรูปสญฺญี พหิทฺธา รูปานิ ปสฺสติ อปฺปมาณานิ ตานิ อภิภุยฺย ชานามิ ปสฺสามีติ วิวิจฺเจว กาเมหิ ฯเปฯ ปฐมํ ฌานํ อุปสมฺปชฺช วิหรติ ทุกฺข\-ปฏิปทํ ทนฺธาภิญฺญํ อปฺปมาณํ อปฺปมาณา\-รมฺมณํ. ตสฺมิํ สมเย ผสฺโส โหติ ฯเปฯ อวิกฺเขโป โหติ ฯเปฯ. อิเม ธมฺมา กุสลา.}

\proseitem{237}{กตเม ธมฺมา กุสลา. ยสฺมิํ สมเย รูปูปปตฺติยา มคฺคํ ภาเวติ อชฺฌตฺตํ อรูปสญฺญี พหิทฺธา รูปานิ ปสฺสติ อปฺปมาณานิ ตานิ อภิภุยฺย ชานามิ ปสฺสามีติ วิวิจฺเจว กาเมหิ ฯเปฯ ปฐมํ ฌานํ อุปสมฺปชฺช \csromanfolio{63}วิหรติ ทุกฺข\-ปฏิปทํ ขิปฺปาภิญฺญํ ปริตฺตํ อปฺปมาณา\-รมฺมณํ. ตสฺมิํ สมเย ผสฺโส โหติ ฯเปฯ อวิกฺเขโป โหติ ฯเปฯ. อิเม ธมฺมา กุสลา.}

\proseitem{238}{กตเม ธมฺมา กุสลา. ยสฺมิํ สมเย รูปูปปตฺติยา มคฺคํ ภาเวติ อชฺฌตฺตํ อรูปสญฺญี พหิทฺธา รูปานิ ปสฺสติ อปฺปมาณานิ ตานิ อภิภุยฺย ชานามิ ปสฺสามีติ วิวิจฺเจว กาเมหิ ฯเปฯ ปฐมํ ฌานํ อุปสมฺปชฺช วิหรติ ทุกฺข\-ปฏิปทํ ขิปฺปาภิญฺญํ อปฺปมาณํ อปฺปมาณา\-รมฺมณํ. ตสฺมิํ สมเย ผสฺโส โหติ ฯเปฯ อวิกฺเขโป โหติ ฯเปฯ. อิเม ธมฺมา กุสลา.}

\proseitem{239}{กตเม ธมฺมา กุสลา. ยสฺมิํ สมเย รูปูปปตฺติยา มคฺคํ ภาเวติ อชฺฌตฺตํ อรูปสญฺญี พหิทฺธา รูปานิ ปสฺสติ อปฺปมาณานิ ตานิ อภิภุยฺย ชานามิ ปสฺสามีติ วิวิจฺเจว กาเมหิ ฯเปฯ ปฐมํ ฌานํ อุปสมฺปชฺช วิหรติ สุขปฏิปทํ ทนฺธาภิญฺญํ ปริตฺตํ อปฺปมาณา\-รมฺมณํ. ตสฺมิํ สมเย ผสฺโส โหติ ฯเปฯ อวิกฺเขโป โหติ ฯเปฯ. อิเม ธมฺมา กุสลา.}

\proseitem{240}{กตเม ธมฺมา กุสลา. ยสฺมิํ สมเย รูปูปปตฺติยา มคฺคํ ภาเวติ อชฺฌตฺตํ อรูปสญฺญี พหิทฺธา รูปานิ ปสฺสติ อปฺปมาณานิ ตานิ อภิภุยฺย ชานามิ ปสฺสามีติ วิวิจฺเจว กาเมหิ ฯเปฯ ปฐมํ ฌานํ อุปสมฺปชฺช วิหรติ สุขปฏิปทํ ทนฺธาภิญฺญํ อปฺปมาณํ อปฺปมาณา\-รมฺมณํ. ตสฺมิํ สมเย ผสฺโส โหติ ฯเปฯ อวิกฺเขโป โหติ ฯเปฯ. อิเม ธมฺมา กุสลา.}

\proseitem{241}{กตเม ธมฺมา กุสลา. ยสฺมิํ สมเย รูปูปปตฺติยา มคฺคํ ภาเวติ อชฺฌตฺตํ อรูปสญฺญี พหิทฺธา รูปานิ ปสฺสติ อปฺปมาณานิ ตานิ อภิภุยฺย ชานามิ ปสฺสามีติ วิวิจฺเจว กาเมหิ ฯเปฯ ปฐมํ ฌานํ อุปสมฺปชฺช วิหรติ สุขปฏิปทํ ขิปฺปาภิญฺญํ ปริตฺตํ อปฺปมาณา\-รมฺมณํ. ตสฺมิํ สมเย ผสฺโส โหติ ฯเปฯ อวิกฺเขโป โหติ ฯเปฯ. อิเม ธมฺมา กุสลา.}

\proseitem{242}{กตเม ธมฺมา กุสลา. ยสฺมิํ สมเย รูปูปปตฺติยา มคฺคํ ภาเวติ อชฺฌตฺตํ อรูปสญฺญี พหิทฺธา รูปานิ ปสฺสติ อปฺปมาณานิ ตานิ อภิภุยฺย ชานามิ ปสฺสามีติ วิวิจฺเจว กาเมหิ ฯเปฯ ปฐมํ ฌานํ อุปสมฺปชฺช วิหรติ สุขปฏิปทํ ขิปฺปาภิญฺญํ อปฺปมาณํ อปฺปมาณา\-รมฺมณํ. ตสฺมิํ สมเย ผสฺโส โหติ ฯเปฯ อวิกฺเขโป โหติ ฯเปฯ. อิเม ธมฺมา กุสลา.}

\proseitem{243}{กตเม ธมฺมา กุสลา. ยสฺมิํ สมเย รูปูปปตฺติยา มคฺคํ ภาเวติ อชฺฌตฺตํ อรูปสญฺญี พหิทฺธา รูปานิ ปสฺสติ อปฺปมาณานิ ตานิ \csromanfolio{64}อภิภุยฺย ชานามิ ปสฺสามีติ วิตฺตกฺกวิจารานํ วูปสมา ฯเปฯ ทุติยํ ฌานํ ฯเปฯ ตติยํ ฌานํ ฯเปฯ จตุตฺถํ ฌานํ ฯเปฯ ปฐมํ ฌานํ ฯเปฯ ปญฺจมํ ฌานํ อุปสมฺปชฺช วิหรติ ทุกฺข\-ปฏิปทํ ทนฺธาภิญฺญํ ปริตฺตํ อปฺปมาณา\-รมฺมณํ ฯเปฯ ทุกฺข\-ปฏิปทํ ทนฺธาภิญฺญํ อปฺปมาณํ อปฺปมาณา\-รมฺมณํ ฯเปฯ ทุกฺข\-ปฏิปทํ ขิปฺปาภิญฺญํ ปริตฺตํ อปฺปมาณา\-รมฺมณํ ฯเปฯ ทุกฺข\-ปฏิปทํ ขิปฺปาภิญฺญํ อปฺปมาณํ อปฺปมาณา\-รมฺมณํ ฯเปฯ สุขปฏิปทํ ทนฺธาภิญฺญํ ปริตฺตํ อปฺปมาณา\-รมฺมณํ ฯเปฯ สุขปฏิปทํ ทนฺธาภิญฺญํ อปฺปมาณํ อปฺปมาณา\-รมฺมณํ ฯเปฯ สุขปฏิปทํ ขิปฺปาภิญฺญํ ปริตฺตํ อปฺปมาณา\-รมฺมณํ ฯเปฯ สุขปฏิปทํ ขิปฺปาภิญฺญํ อปฺปมาณํ อปฺปมาณา\-รมฺมณํ. ตสฺมิํ สมเย ผสฺโส โหติ ฯเปฯ อวิกฺเขโป โหติ ฯเปฯ. อิเม ธมฺมา กุสลา.}

\vspace{\tipitakaparskip}
\csromancenter{อปรมฺปิ อฏฺฐ\-กฺขตฺตุกํ.}
\csromanmatikaanchor{mtk.36}
\csromantocmarkref{section}{อิทมฺปิ อฏฺฐกฺขตฺตุกํ}{mtk.36}
\bookmark[dest={mtk.36},level=1]{อิทมฺปิ อฏฺฐกฺขตฺตุกํ}
\csromanheader{1}{อิทมฺปิ อฏฺฐ\-กฺขตฺตุกํ}
\proseitem{244}{กตเม ธมฺมา กุสลา. ยสฺมิํ สมเย รูปูปปตฺติยา มคฺคํ ภาเวติ อชฺฌตฺตํ อรูปสญฺญี พหิทฺธา รูปานิ ปสฺสติ อปฺปมาณานิ สุวณฺณ\-ทุพฺพณฺณานิ ตานิ อภิภุยฺย ชานามิ ปสฺสามีติ วิวิจฺเจว กาเมหิ ฯเปฯ ปฐมํ ฌานํ อุปสมฺปชฺช วิหรติ. ตสฺมิํ สมเย ผสฺโส โหติ ฯเปฯ อวิกฺเขโป โหติ ฯเปฯ. อิเม ธมฺมา กุสลา.}

\proseitem{245}{กตเม ธมฺมา กุสลา. ยสฺมิํ สมเย รูปูปปตฺติยา มคฺคํ ภาเวติ อชฺฌตฺตํ อรูปสญฺญี พหิทฺธา รูปานิ ปสฺสติ อปฺปมาณานิ สุวณฺณ\-ทุพฺพณฺณานิ ตานิ อภิภุยฺย ชานามิ ปสฺสามีติ วิตกฺก\-วิจารานํ วูปสมา ฯเปฯ ทุติยํ ฌานํ ฯเปฯ ตติยํ ฌานํ ฯเปฯ จตุตฺถํ ฌานํ ฯเปฯ ปฐมํ ฌานํ ฯเปฯ ปญฺจมํ ฌานํ อุปสมฺปชฺช วิหรติ. ตสฺมิํ สมเย ผสฺโส โหติ ฯเปฯ อวิกฺเขโป โหติ ฯเปฯ. อิเม ธมฺมา กุสลา.}

\vspace{\tipitakaparskip}
\csromancenter{อิทมฺปิ อฏฺฐ\-กฺขตฺตุกํ.}
\proseitem{246}{กตเม ธมฺมา กุสลา. ยสฺมิํ สมเย รูปูปปตฺติยา มคฺคํ ภาเวติ อชฺฌตฺตํ อรูปสญฺญี พหิทฺธา รูปานิ ปสฺสติ นีลานิ นีลวณฺณานิ นีลนิทสฺสนานิ นีลนิภาสานิ ตานิ อภิภุยฺย ชานามิ ปสฺสามีติ \csromanfolio{65}วิวิจฺเจว กาเมหิ ฯเปฯ ปฐมํ ฌานํ อุปสมฺปชฺช วิหรติ. ตสฺมิํ สมเย ผสฺโส โหติ ฯเปฯ อวิกฺเขโป โหติ ฯเปฯ. อิเม ธมฺมา กุสลา.}

\proseitem{247}{กตเม ธมฺมา กุสลา. ยสฺมิํ สมเย รูปูปปตฺติยา มคฺคํ ภาเวติ อชฺฌตฺตํ อรูปสญฺญี พหิทฺธา รูปานิ ปสฺสติ ปีตานิ ปีตวณฺณานิ ปีตนิทสฺสนานิ ปีตนิภาสานิ ฯเปฯ โลหิตกานิ โลหิตก\-วณฺณานิ โลหิตก\-นิทสฺสนานิ โลหิตก\-นิภาสานิ ฯเปฯ โอทาตานิ โอทาตวณฺณานิ โอทาต\-นิทสฺสนานิ โอทาตนิภาสานิ ตานิ อภิภุยฺย ชานามิ ปสฺสามีติ วิวิจฺเจว กาเมหิ ฯเปฯ ปฐมํ ฌานํ อุปสมฺปชฺช วิหรติ. ตสฺมิํ สมเย ผสฺโส โหติ ฯเปฯ อวิกฺเขโป โหติ ฯเปฯ. อิเม ธมฺมา กุสลา.}

\vspace{\tipitakaparskip}
\csromancenter{อิมานิปิ อภิภา\-ยตนานิ โสฬส\-กฺขตฺตุ\-กานิ.}
\csromanmatikaanchor{mtk.37}
\csromantocmarkref{section}{ตีณิ วิโมกฺขานิ โสฬสกฺขตฺตุกานิ}{mtk.37}
\bookmark[dest={mtk.37},level=1]{ตีณิ วิโมกฺขานิ โสฬสกฺขตฺตุกานิ}
\csromanheader{1}{ตีณิ วิโมกฺขานิ โสฬส\-กฺขตฺตุ\-กานิ}
\proseitem{248}{กตเม ธมฺมา กุสลา. ยสฺมิํ สมเย รูปูปปตฺติยา มคฺคํ ภาเวติ รูปี รูปานิ ปสฺสติ วิวิจฺเจว กาเมหิ ฯเปฯ ปฐมํ ฌานํ อุปสมฺปชฺช วิหรติ. ตสฺมิํ สมเย ผสฺโส โหติ ฯเปฯ อวิกฺเขโป โหติ ฯเปฯ. อิเม ธมฺมา กุสลา.}

\proseitem{249}{กตเม ธมฺมา กุสลา. ยสฺมิํ สมเย รูปูปปตฺติยา มคฺคํ ภาเวติ อชฺฌตฺตํ อรูปสญฺญี พหิทฺธา รูปานิ ปสฺสติ วิวิจฺเจว กาเมหิ ฯเปฯ ปฐมํ ฌานํ อุปสมฺปชฺช วิหรติ. ตสฺมิํ สมเย ผสฺโส โหติ ฯเปฯ อวิกฺเขโป โหติ ฯเปฯ. อิเม ธมฺมา กุสลา.}

\proseitem{250}{กตเม ธมฺมา กุสลา. ยสฺมิํ สมเย รูปูปปตฺติยา มคฺคํ ภาเวติ สุภนฺติ วิวิจฺเจว กาเมหิ ฯเปฯ ปฐมํ ฌานํ อุปสมฺปชฺช วิหรติ. ตสฺมิํ สมเย ผสฺโส โหติ ฯเปฯ อวิกฺเขโป โหติ ฯเปฯ. อิเม ธมฺมา กุสลา.}

\vspace{\tipitakaparskip}
\csromancenter{อิมานิปิ ตีณิ วิโมกฺขานิ โสฬส\-กฺขตฺตุ\-กานิ.}
\csromanmatikaanchor{mtk.38}
\csromantocmarkref{section}{จตฺตาริ พฺรหฺมวิหารฌานานิ โสฬสกฺขตฺตุกานิ}{mtk.38}
\bookmark[dest={mtk.38},level=1]{จตฺตาริ พฺรหฺมวิหารฌานานิ โสฬสกฺขตฺตุกานิ}
\csromanheader{1}{จตฺตาริ พฺรหฺมวิหาร\-ฌานานิ โสฬส\-กฺขตฺตุ\-กานิ}
\proseitem{251}{\csromanfolio{66}กตเม ธมฺมา กุสลา. ยสฺมิํ สมเย รูปูปปตฺติยา มคฺคํ ภาเวติ วิวิจฺเจว กาเมหิ ฯเปฯ ปฐมํ ฌานํ อุปสมฺปชฺช วิหรติ เมตฺตาสหคตํ. ตสฺมิํ สมเย ผสฺโส โหติ ฯเปฯ อวิกฺเขโป โหติ ฯเปฯ. อิเม ธมฺมา กุสลา.}

\proseitem{252}{กตเม ธมฺมา กุสลา. ยสฺมิํ สมเย รูปูปปตฺติยา มคฺคํ ภาเวติ วิตกฺก\-วิจารานํ วูปสมา ฯเปฯ ทุติยํ ฌานํ อุปสมฺปชฺช วิหรติ เมตฺตาสหคตํ. ตสฺมิํ สมเย ผสฺโส โหติ ฯเปฯ อวิกฺเขโป โหติ ฯเปฯ. อิเม ธมฺมา กุสลา.}

\proseitem{253}{กตเม ธมฺมา กุสลา. ยสฺมิํ สมเย รูปูปปตฺติยา มคฺคํ ภาเวติ ปีติยา จ วิราคา ฯเปฯ ตติยํ ฌานํ อุปสมฺปชฺช วิหรติ เมตฺตาสหคตํ. ตสฺมิํ สมเย ผสฺโส โหติ ฯเปฯ อวิกฺเขโป โหติ ฯเปฯ. อิเม ธมฺมา กุสลา.}

\proseitem{254}{กตเม ธมฺมา กุสลา. ยสฺมิํ สมเย รูปูปปตฺติยา มคฺคํ ภาเวติ วิวิจฺเจว กาเมหิ ฯเปฯ ปฐมํ ฌานํ อุปสมฺปชฺช วิหรติ เมตฺตาสหคตํ. ตสฺมิํ สมเย ผสฺโส โหติ ฯเปฯ อวิกฺเขโป โหติ ฯเปฯ. อิเม ธมฺมา กุสลา.}

\proseitem{255}{กตเม ธมฺมา กุสลา. ยสฺมิํ สมเย รูปูปปตฺติยา มคฺคํ ภาเวติ อวิตกฺกํ วิจารมตฺตํ สมาธิชํ ปีติสุขํ ทุติยํ ฌานํ อุปสมฺปชฺช วิหรติ เมตฺตาสหคตํ. ตสฺมิํ สมเย ผสฺโส โหติ ฯเปฯ อวิกฺเขโป โหติ ฯเปฯ. อิเม ธมฺมา กุสลา.}

\proseitem{256}{กตเม ธมฺมา กุสลา. ยสฺมิํ สมเย รูปูปปตฺติยา มคฺคํ ภาเวติ วิตกฺก\-วิจารานํ วูปสมา ฯเปฯ ตติยํ ฌานํ อุปสมฺปชฺช วิหรติ เมตฺตาสหคตํ. ตสฺมิํ สมเย ผสฺโส โหติ ฯเปฯ อวิกฺเขโป โหติ ฯเปฯ. อิเม ธมฺมา กุสลา.}

\proseitem{257}{กตเม ธมฺมา กุสลา. ยสฺมิํ สมเย รูปูปปตฺติยา มคฺคํ ภาเวติ ปีติยา จ วิราคา ฯเปฯ จตุตฺถํ ฌานํ อุปสมฺปชฺช วิหรติ เมตฺตาสหคตํ. ตสฺมิํ สมเย ผสฺโส โหติ ฯเปฯ อวิกฺเขโป โหติ ฯเปฯ. อิเม ธมฺมา กุสลา.}

\chapterheadpageread{1. จิตฺตุปฺปาท\-กณฺฑ}
\proseitem{258}{\csromanfolio{67}กตเม ธมฺมา กุสลา. ยสฺมิํ สมเย รูปูปปตฺติยา มคฺคํ ภาเวติ วิวิจฺเจว กาเมหิ ฯเปฯ ปฐมํ ฌานํ อุปสมฺปชฺช วิหรติ กรุณา\-สหคตํ. ตสฺมิํ สมเย ผสฺโส โหติ ฯเปฯ อวิกฺเขโป โหติ ฯเปฯ. อิเม ธมฺมา กุสลา.}

\proseitem{259}{กตเม ธมฺมา กุสลา. ยสฺมิํ สมเย รูปูปปตฺติยา มคฺคํ ภาเวติ วิตกฺก\-วิจารานํ วูปสมา ฯเปฯ ทุติยํ ฌานํ ฯเปฯ ตติยํ ฌานํ ฯเปฯ ปฐมํ ฌานํ ฯเปฯ จตุตฺถํ ฌานํ อุปสมฺปชฺช วิหรติ กรุณา\-สหคตํ. ตสฺมิํ สมเย ผสฺโส โหติ ฯเปฯ อวิกฺเขโป โหติ ฯเปฯ. อิเม ธมฺมา กุสลา.}

\proseitem{260}{กตเม ธมฺมา กุสลา. ยสฺมิํ สมเย รูปูปปตฺติยา มคฺคํ ภาเวติ วิวิจฺเจว กาเมหิ ฯเปฯ ปฐมํ ฌานํ อุปสมฺปชฺช วิหรติ มุทิตา\-สหคตํ. ตสฺมิํ สมเย ผสฺโส โหติ ฯเปฯ อวิกฺเขโป โหติ ฯเปฯ. อิเม ธมฺมา กุสลา.}

\proseitem{261}{กตเม ธมฺมา กุสลา. ยสฺมิํ สมเย รูปูปปตฺติยา มคฺคํ ภาเวติ วิตกฺก\-วิจารานํ วูปสมา ฯเปฯ ทุติยํ ฌานํ ฯเปฯ ตติยํ ฌานํ ฯเปฯ ปฐมํ ฌานํ ฯเปฯ จตุตฺถํ ฌานํ อุปสมฺปชฺช วิหรติ มุทิตา\-สหคตํ. ตสฺมิํ สมเย ผสฺโส โหติ ฯเปฯ อวิกฺเขโป โหติ ฯเปฯ. อิเม ธมฺมา กุสลา.}

\proseitem{262}{กตเม ธมฺมา กุสลา. ยสฺมิํ สมเย รูปูปปตฺติยา มคฺคํ ภาเวติ สุขสฺส จ ปหานา ฯเปฯ จตุตฺถํ ฌานํ อุปสมฺปชฺช วิหรติ อุเปกฺขา\-สหคตํ. ตสฺมิํ สมเย ผสฺโส โหติ ฯเปฯ อวิกฺเขโป โหติ ฯเปฯ. อิเม ธมฺมา กุสลา. จตฺตาริ พฺรหฺมวิหาร\-ฌานานิ\csromansharedfootnote{d107d19d67d0bacc}{พฺรหฺมวิหารชฺฌานานิ (สี, สฺยา)} โสฬส\-กฺขตฺตุ\-กานิ.}

\csromanmatikaanchor{mtk.39}
\csromantocmarkref{section}{อสุภฌานํ โสฬสกฺขตฺตุกํ}{mtk.39}
\bookmark[dest={mtk.39},level=1]{อสุภฌานํ โสฬสกฺขตฺตุกํ}
\csromanheader{1}{อสุภฌานํ โสฬส\-กฺขตฺตุกํ}
\proseitem{263}{กตเม ธมฺมา กุสลา. ยสฺมิํ สมเย รูปูปปตฺติยา มคฺคํ ภาเวติ วิวิจฺเจว กาเมหิ ฯเปฯ ปฐมํ ฌานํ อุปสมฺปชฺช วิหรติ อุทฺธุมาตก\-สญฺญาสหคตํ. ตสฺมิํ สมเย ผสฺโส โหติ ฯเปฯ อวิกฺเขโป โหติ ฯเปฯ. อิเม ธมฺมา กุสลา.}

\proseitem{264}{\csromanfolio{68}กตเม ธมฺมา กุสลา. ยสฺมิํ สมเย รูปูปปตฺติยา มคฺคํ ภาเวติ วิวิจฺเจว กาเมหิ ฯเปฯ ปฐมํ ฌานํ อุปสมฺปชฺช วิหรติ วินีลกสญฺญา\-สหคตํ ฯเปฯ วิปุพฺพกสญฺญา\-สหคตํ ฯเปฯ วิจฺฉิทฺทก\-สญฺญาสหคตํ ฯเปฯ วิกฺขายิตก\-สญฺญา\-สหคตํ ฯเปฯ วิกฺขิตฺตก\-สญฺญา\-สหคตํ ฯเปฯ หตวิกฺขิตฺตก\-สญฺญา\-สหคตํ ฯเปฯ โลหิตก\-สญฺญา\-สหคตํ ฯเปฯ ปุฬวกสญฺญา\-สหคตํ\csromansharedfootnote{76e0eb3cf9cd8b8c}{ปุฬุวกสญฺญาสหคตํ (ก)} ฯเปฯ อฏฺฐิกสญฺญา\-สหคตํ. ตสฺมิํ สมเย ผสฺโส โหติ ฯเปฯ อวิกฺเขโป โหติ ฯเปฯ. อิเม ธมฺมา กุสลา.}

\vspace{\tipitakaparskip}
\csromancenter{อสุภฌานํ\csromansharedfootnote{2e3d0b6a74c8b7a0}{อสุภชฺฌานํ (สี, สฺยา)} โสฬส\-กฺขตฺตุกํ.}
\csromancenter{รูปาวจร\-กุสลํ.}
\csromanmatikaanchor{mtk.40}
\csromantocmarkref{section}{อรูปาวจรกุสล}{mtk.40}
\bookmark[dest={mtk.40},level=1]{อรูปาวจรกุสล}
\csromanheader{1}{อรูปาวจร\-กุสล}
\csromanheader{2}{จตฺตาริ อรูปฌานานิ โสฬส\-กฺขตฺตุ\-กานิ}
\proseitem{265}{กตเม ธมฺมา กุสลา. ยสฺมิํ สมเย อรูปูปปตฺติยา มคฺคํ ภาเวติ สพฺพโส รูปสญฺญานํ สมติกฺกมา ปฏิฆ\-สญฺญานํ อตฺถงฺคมา นานตฺต\-สญฺญานํ อมนสิการา อากาสา\-นญฺจา\-ยตนสญฺญาสหคตํ สุขสฺส จ ปหานา ฯเปฯ จตุตฺถํ ฌานํ อุปสมฺปชฺช วิหรติ อุเปกฺขา\-สหคตํ. ตสฺมิํ สมเย ผสฺโส โหติ ฯเปฯ อวิกฺเขโป โหติ ฯเปฯ. อิเม ธมฺมา กุสลา.}

\proseitem{266}{กตเม ธมฺมา กุสลา. ยสฺมิํ สมเย อรูปูปปตฺติยา มคฺคํ ภาเวติ สพฺพโส อากาสา\-นญฺจา\-ยตนํ สมติกฺกมฺม วิญฺญาณญฺจายตน\-สญฺญาสหคตํ สุขสฺส จ ปหานา ฯเปฯ จตุตฺถํ ฌานํ อุปสมฺปชฺช วิหรติ อุเปกฺขา\-สหคตํ. ตสฺมิํ สมเย ผสฺโส โหติ ฯเปฯ อวิกฺเขโป โหติ ฯเปฯ. อิเม ธมฺมา กุสลา.}

\proseitem{267}{กตเม ธมฺมา กุสลา. ยสฺมิํ สมเย อรูปูปปตฺติยา มคฺคํ ภาเวติ สพฺพโส วิญฺญาณญฺจา\-ยตนํ สมติกฺกมฺม อากิญฺจญฺญายตน\-สญฺญาสหคตํ สุขสฺส จ ปหานา ฯเปฯ จตุตฺถํ ฌานํ อุปสมฺปชฺช วิหรติ \csromanfolio{69}อุเปกฺขา\-สหคตํ. ตสฺมิํ สมเย ผสฺโส โหติ ฯเปฯ อวิกฺเขโป โหติ ฯเปฯ. อิเม ธมฺมา กุสลา.}

\proseitem{268}{กตเม ธมฺมา กุสลา. ยสฺมิํ สมเย อรูปูปปตฺติยา มคฺคํ ภาเวติ สพฺพโส อากิญฺจญฺญา\-ยตนํ สมติกฺกมฺม เนว\-สญฺญา\-นา\-สญฺญา\-ยตนสญฺญาสหคตํ สุขสฺส จ ปหานา ฯเปฯ จตุตฺถํ ฌานํ อุปสมฺปชฺช วิหรติ อุเปกฺขา\-สหคตํ. ตสฺมิํ สมเย ผสฺโส โหติ ฯเปฯ อวิกฺเขโป โหติ ฯเปฯ. อิเม ธมฺมา กุสลา. จตฺตาริ อรูปฌานานิ\csromansharedfootnote{fd3dea5bb8e28bb2}{อรูปชฺฌานานิ (สี, สฺยา)} โสฬส\-กฺขตฺตุ\-กานิ. อรูปาวจร\-กุสลํ.}

\csromanheader{2}{เตภูมก\-กุสล}
\csromanmatikaanchor{mtk.41}
\csromantocmarkref{section}{เตภูมกกามาวจรกุสล}{mtk.41}
\bookmark[dest={mtk.41},level=1]{เตภูมกกามาวจรกุสล}
\csromanheader{1}{กามาวจร\-กุสล}
\proseitem{269}{กตเม ธมฺมา กุสลา. ยสฺมิํ สมเย กามาวจรํ กุสลํ จิตฺตํ อุปฺปนฺนํ โหติ โสมนสฺส\-สหคตํ ญาณ\-สมฺปยุตฺตํ หีนํ ฯเปฯ มชฺฌิมํ ฯเปฯ ปณีตํ ฯเปฯ ฉนฺทา\-ธิปเตยฺยํ ฯเปฯ วีริยา\-ธิปเตยฺยํ ฯเปฯ จิตฺตา\-ธิปเตยฺยํ ฯเปฯ วีมํสา\-ธิปเตยฺยํ ฯเปฯ ฉนฺทา\-ธิปเตยฺยํ หีนํ ฯเปฯ มชฺฌิมํ ฯเปฯ ปณีตํ ฯเปฯ วีริยา\-ธิปเตยฺยํ หีนํ ฯเปฯ มชฺฌิมํ ฯเปฯ ปณีตํ ฯเปฯ จิตฺตา\-ธิปเตยฺยํ หีนํ ฯเปฯ มชฺฌิมํ ฯเปฯ ปณีตํ ฯเปฯ วีมํสา\-ธิปเตยฺยํ หีนํ ฯเปฯ มชฺฌิมํ ฯเปฯ ปณีตํ. ตสฺมิํ สมเย ผสฺโส โหติ ฯเปฯ อวิกฺเขโป โหติ ฯเปฯ. อิเม ธมฺมา กุสลา.}

\proseitem{270}{กตเม ธมฺมา กุสลา. ยสฺมิํ สมเย กามาวจรํ กุสลํ จิตฺตํ อุปฺปนฺนํ โหติ โสมนสฺส\-สหคตํ ญาณ\-สมฺปยุตฺตํ สสงฺขาเรน ฯเปฯ โสมนสฺส\-สหคตํ ญาณวิปฺปยุตฺตํ ฯเปฯ โสมนสฺส\-สหคตํ ญาณวิปฺปยุตฺตํ สสงฺขาเรน ฯเปฯ อุเปกฺขา\-สหคตํ ญาณ\-สมฺปยุตฺตํ ฯเปฯ อุเปกฺขา\-สหคตํ ญาณ\-สมฺปยุตฺตํ สสงฺขาเรน ฯเปฯ อุเปกฺขา\-สหคตํ ญาณวิปฺปยุตฺตํ ฯเปฯ อุเปกฺขา\-สหคตํ ญาณวิปฺปยุตฺตํ สสงฺขาเรน หีนํ ฯเปฯ มชฺฌิมํ ฯเปฯ ปณีตํ ฯเปฯ ฉนฺทา\-ธิปเตยฺยํ ฯเปฯ วีริยา\-ธิปเตยฺยํ ฯเปฯ จิตฺตา\-ธิปเตยฺยํ ฯเปฯ ฉนฺทา\-ธิปเตยฺยํ หีนํ ฯเปฯ มชฺฌิมํ ฯเปฯ ปณีตํ ฯเปฯ วีริยา\-ธิปเตยฺยํ หีนํ ฯเปฯ มชฺฌิมํ ฯเปฯ \csromanfolio{70}ปณีตํ ฯเปฯ จิตฺตา\-ธิปเตยฺยํ หีนํ ฯเปฯ มชฺฌิมํ ฯเปฯ ปณีตํ. ตสฺมิํ สมเย ผสฺโส โหติ ฯเปฯ อวิกฺเขโป โหติ ฯเปฯ. อิเม ธมฺมา กุสลา.}

\vspace{\tipitakaparskip}
\csromancenter{กามาวจร\-กุสลํ.}
\csromanmatikaanchor{mtk.42}
\csromantocmarkref{section}{รูปาวจรกุสล}{mtk.42}
\bookmark[dest={mtk.42},level=1]{รูปาวจรกุสล}
\csromanheader{1}{รูปาวจร\-กุสล}
\proseitem{271}{กตเม ธมฺมา กุสลา. ยสฺมิํ สมเย รูปูปปตฺติยา มคฺคํ ภาเวติ วิวิจฺเจว กาเมหิ ฯเปฯ ปฐมํ ฌานํ อุปสมฺปชฺช วิหรติ ปถวีกสิณํ หีนํ ฯเปฯ มชฺฌิมํ ฯเปฯ ปณีตํ ฯเปฯ ฉนฺทา\-ธิปเตยฺยํ ฯเปฯ วีริยา\-ธิปเตยฺยํ ฯเปฯ จิตฺตา\-ธิปเตยฺยํ ฯเปฯ วีมํสา\-ธิปเตยฺยํ ฯเปฯ ฉนฺทา\-ธิปเตยฺยํ หีนํ ฯเปฯ มชฺฌิมํ ฯเปฯ ปณีตํ ฯเปฯ วีริยา\-ธิปเตยฺยํ หีนํ ฯเปฯ มชฺฌิมํ ฯเปฯ ปณีตํ ฯเปฯ จิตฺตา\-ธิปเตยฺยํ หีนํ ฯเปฯ มชฺฌิมํ ฯเปฯ ปณีตํ ฯเปฯ วีมํสา\-ธิปเตยฺยํ หีนํ ฯเปฯ มชฺฌิมํ ฯเปฯ ปณีตํ. ตสฺมิํ สมเย ผสฺโส โหติ ฯเปฯ อวิกฺเขโป โหติ ฯเปฯ. อิเม ธมฺมา กุสลา.}

\proseitem{272}{กตเม ธมฺมา กุสลา. ยสฺมิํ สมเย รูปูปปตฺติยา มคฺคํ ภาเวติ วิตกฺก\-วิจารานํ วูปสมา ฯเปฯ ทุติยํ ฌานํ ฯเปฯ ตติยํ ฌานํ ฯเปฯ จตุตฺถํ ฌานํ ฯเปฯ ปฐมํ ฌานํ ฯเปฯ ปญฺจมํ ฌานํ อุปสมฺปชฺช วิหรติ ปถวีกสิณํ หีนํ ฯเปฯ มชฺฌิมํ ฯเปฯ ปณีตํ ฯเปฯ ฉนฺทา\-ธิปเตยฺยํ ฯเปฯ วีริยธิปเตยฺยํ ฯเปฯ จิตฺตา\-ธิปเตยฺยํ ฯเปฯ วีมํสา\-ธิปเตยฺยํ ฯเปฯ ฉนฺทา\-ธิปเตยฺยํ หีนํ ฯเปฯ มชฺฌิมํ ฯเปฯ ปณีตํ ฯเปฯ วีริยา\-ธิปเตยฺยํ หีนํ ฯเปฯ มชฺฌิมํ ฯเปฯ ปณีตํ ฯเปฯ จิตฺตา\-ธิปเตยฺยํ หีนํ ฯเปฯ มชฺฌิมํ ฯเปฯ ปณีตํ ฯเปฯ วีมํสา\-ธิปเตยฺยํ หีนํ ฯเปฯ มชฺฌิมํ ฯเปฯ ปณีตํ. ตสฺมิํ สมเย ผสฺโส โหติ ฯเปฯ อวิกฺเขโป โหติ ฯเปฯ. อิเม ธมฺมา กุสลา. รูปาวจร\-กุสลํ.}

\csromanmatikaanchor{mtk.43}
\csromantocmarkref{section}{อรูปาวจรกุสล}{mtk.43}
\bookmark[dest={mtk.43},level=1]{อรูปาวจรกุสล}
\csromanheader{1}{อรูปาวจร\-กุสล}
\proseitem{273}{กตเม ธมฺมา กุสลา. ยสฺมิํ สมเย อรูปูปปตฺติยา มคฺคํ ภเวติ สพฺพโส รูปสญฺญานํ สมติกฺกมา ปฏิฆ\-สญฺญานํ อตฺถงฺคมา \csromanfolio{71}นานตฺต\-สญฺญานํ อมนสิการา อากาสา\-นญฺจา\-ยตนสญฺญาสหคตํ สุขสฺส จ ปหานา ฯเปฯ จตุตฺถํ ฌานํ อุปสมฺปชฺช วิหรติ หีนํ ฯเปฯ มชฺฌิมํ ฯเปฯ ปณีตํ ฯเปฯ ฉนฺทา\-ธิปเตยฺยํ ฯเปฯ วีริยา\-ธิปเตยฺยํ ฯเปฯ จิตฺตา\-ธิปเตยฺยํ ฯเปฯ วีมํสา\-ธิปเตยฺยํ ฯเปฯ ฉนฺทา\-ธิปเตยฺยํ หีนํ ฯเปฯ มชฺฌิมํ ฯเปฯ ปณีตํ ฯเปฯ วีริยา\-ธิปเตยฺยํ หีนํ ฯเปฯ มชฺฌิมํ ฯเปฯ ปณีตํ ฯเปฯ จิตฺตา\-ธิปเตยฺยํ หีนํ ฯเปฯ มชฺฌิมํ ฯเปฯ ปณีตํ ฯเปฯ วีมํสา\-ธิปเตยฺยํ หีนํ ฯเปฯ มชฺฌิมํ ฯเปฯ ปณีตํ. ตสฺมิํ สมเย ผสฺโส โหติ ฯเปฯ อวิกฺเขโป โหติ ฯเปฯ. อิเม ธมฺมา กุสลา.}

\proseitem{274}{กตเม ธมฺมา กุสลา. ยสฺมิํ สมเย อรูปูปปตฺติยา มคฺคํ ภาเวติ สพฺพโส อากาสา\-นญฺจา\-ยตนํ สมติกฺกมฺม วิญฺญาณญฺจายตน\-สญฺญาสหคตํ สุขสฺส จ ปหานา ฯเปฯ จตุตฺถํ ฌานํ อุปสมฺปชฺช วิหรติ หีนํ ฯเปฯ มชฺฌิมํ ฯเปฯ ปณีตํ ฯเปฯ ฉนฺทา\-ธิปเตยฺยํ ฯเปฯ วีริยา\-ธิปเตยฺยํ ฯเปฯ จิตฺตา\-ธิปเตยฺยํ ฯเปฯ วีมํสา\-ธิปเตยฺยํ ฯเปฯ ฉนฺทา\-ธิปเตยฺยํ หีนํ ฯเปฯ มชฺฌิมํ ฯเปฯ ปณีตํ ฯเปฯ วีริยา\-ธิปเตยฺยํ หีนํ ฯเปฯ มชฺฌิมํ ฯเปฯ ปณีตํ ฯเปฯ จิตฺตา\-ธิปเตยฺยํ หีนํ ฯเปฯ มชฺฌิมํ ฯเปฯ ปณีตํ ฯเปฯ วีมํสา\-ธิปเตยฺยํ หีนํ ฯเปฯ มชฺฌิมํ ฯเปฯ ปณีตํ. ตสฺมิํ สมเย ผสฺโส โหติ ฯเปฯ อวิกฺเขโป โหติ ฯเปฯ. อิเม ธมฺมา กุสลา.}

\proseitem{275}{กตเม ธมฺมา กุสลา. ยสฺมิํ สมเย อรูปูปปตฺติยา มคฺคํ ภาเวติ สพฺพโส วิญฺญาณญฺจา\-ยตนํ สมติกฺกมฺม อากิญฺจญฺญายตน\-สญฺญาสหคตํ สุขสฺส จ ปหานา ฯเปฯ จตุตฺถํ ฌานํ อุปสมฺปชฺช วิหรติ หีนํ ฯเปฯ มชฺฌิมํ ฯเปฯ ปณีตํ ฯเปฯ ฉนฺทธิปเตยฺยํ ฯเปฯ วีริยา\-ธิปเตยฺยํ ฯเปฯ จิตฺตา\-ธิปเตยฺยํ ฯเปฯ วีมํสา\-ธิปเตยฺยํ ฯเปฯ ฉนฺทา\-ธิปเตยฺยํ หีนํ ฯเปฯ มชฺฌิมํ ฯเปฯ ปณีตํ ฯเปฯ วีริยา\-ธิปเตยฺยํ หีนํ ฯเปฯ มชฺฌิมํ ฯเปฯ ปณีตํ ฯเปฯ จิตฺตา\-ธิปเตยฺยํ หีนํ ฯเปฯ มชฺฌิมํ ฯเปฯ ปณีตํ ฯเปฯ วีมํสา\-ธิปเตยฺยํ หีนํ ฯเปฯ มชฺฌิมํ ฯเปฯ ปณีตํ. ตสฺมิํ สมเย ผสฺโส โหติ ฯเปฯ อวิกฺเขโป โหติ ฯเปฯ. อิเม ธมฺมา กุสลา.}

\proseitem{276}{กตเม ธมฺมา กุสลา. ยสฺมิํ สมเย อรูปูปปตฺติยา มคฺคํ ภาเวติ สพฺพโส อากิญฺจญฺญา\-ยตนํ สมติกฺกมฺม เนว\-สญฺญา\-นา\-สญฺญา\-ยตนสญฺญาสหคตํ สุขสฺส จ ปหานา ฯเปฯ จตุตฺถํ ฌานํ อุปสมฺปชฺช วิหรติ หีนํ ฯเปฯ มชฺฌิมํ ฯเปฯ ปณีตํ ฯเปฯ ฉนฺทา\-ธิปเตยฺยํ ฯเปฯ วีริยา\-ธิปเตยฺยํ ฯเปฯ จิตฺตา\-ธิปเตยฺยํ ฯเปฯ วีมํสา\-ธิปเตยฺยํ ฯเปฯ ฉนฺทา\-ธิปเตยฺยํ หีนํ ฯเปฯ มชฺฌิมํ ฯเปฯ \csromanfolio{72}ปณีตํ ฯเปฯ วีริยา\-ธิปเตยฺยํ หีนํ ฯเปฯ มชฺฌิมํ ฯเปฯ ปณีตํ ฯเปฯ จิตฺตา\-ธิปเตยฺยํ หีนํ ฯเปฯ มชฺฌิมํ ฯเปฯ ปณีตํ ฯเปฯ วีมํสา\-ธิปเตยฺยํ หีนํ ฯเปฯ มชฺฌิมํ ฯเปฯ ปณีตํ. ตสฺมิํ สมเย ผสฺโส โหติ ฯเปฯ อวิกฺเขโป โหติ ฯเปฯ. อิเม ธมฺมา กุสลา.}

\vspace{\tipitakaparskip}
\csromancenter{อรูปาวจร\-กุสลํ.}
\csromanmatikaanchor{mtk.44}
\csromantocmarkref{subsection}{โลกุตฺตรกุสลสุทฺธิกปฏิปทา}{mtk.44}
\bookmark[dest={mtk.44},level=2]{โลกุตฺตรกุสลสุทฺธิกปฏิปทา}
\csromanheader{2}{โลกุตฺตร\-กุสล}
\csromanheader{2}{สุทฺธิก\-ปฏิปทา}
\proseitem{277}{กตเม ธมฺมา กุสลา. ยสฺมิํ สมเย โลกุตฺตรํ ฌานํ ภาเวติ นิยฺยานิกํ อปจยคามิํ ทิฏฺฐิคตานํ ปหานาย ปฐมาย ภูมิยา ปตฺติยา วิวิจฺเจว กาเมหิ ฯเปฯ ปฐมํ ฌานํ อุปสมฺปชฺช วิหารติ ทุกฺข\-ปฏิปทํ ทนฺธาภิญฺญํ. ตสฺมิํ สมเย ผสฺโส โหติ, เวทนา โหติ, สญฺญาโหติ, เจตนา โหติ, จิตฺตํ โหติ, วิตกฺโก โหติ, วิจาโร โหติ, ปีติ โหติ, สุขํ โหติ, จิตฺตสฺเส\-กคฺคตา โหติ, สทฺธินฺทฺริยํ โหติ, วีริยินฺทฺริยํ โหติ, สตินฺทฺริยํ โหติ, สมาธินฺทฺริยํ โหติ, ปญฺญินฺทฺริยํ โหติ, มนินฺทฺริยํ โหติ, โสมนสฺสิ\-นฺทฺริยํ โหติ, ชีวิตินฺทฺริยํ โหติ, อนญฺญาต\-ญฺญสฺสามี\-ตินฺทฺริยํ โหติ, สมฺมาทิฏฺฐิ โหติ, สมฺมาสงฺกปฺโป โหติ, สมฺมาวาจา โหติ, สมฺมากมฺมนฺโต โหติ, สมฺมาอาชีโว โหติ, สมฺมาวายาโม โหติ, สมฺมาสติ โหติ, สมฺมาสมาธิ โหติ, สทฺธาพลํ โหติ, วีริยพลํ โหติ, สติพลํ โหติ, สมาธิพลํ โหติ, ปญฺญาพลํ โหติ, หิริพลํ โหติ, โอตฺตปฺปพลํ โหติ, อโลโภ โหติ, อโทโส โหติ, อโมโห โหติ, อนภิชฺฌา โหติ, อพฺยาปาโท โหติ, สมฺมาทิฏฺฐิ โหติ, หิรี โหติ, โอตฺตปฺปํ โหติ, กายปสฺสทฺธิ โหติ, จิตฺตปสฺสทฺธิ โหติ, กายลหุตา โหติ, จิตฺตลหุตา โหติ, กายมุทุตา โหติ, จิตฺตมุทุตา โหติ, กายกมฺมญฺญตา โหติ, จิตฺต\-กมฺมญฺญตา โหติ, กายปาคุญฺญตา โหติ, จิตฺตปาคุญฺญตา โหติ, กายุชุกตา โหติ, จิตฺตุชุกตา โหติ, สติ โหติ, สมฺปชญฺญํ โหติ, สมโถ โหติ, วิปสฺสนา โหติ, \csromanfolio{73}ปคฺคาโห โหติ, อวิกฺเขโป โหติ. เย วา ปน ตสฺมิํ สมเย อญฺเญปิ อตฺถิ ปฏิจฺจ\-สมุปฺปนฺนา อรูปิโน ธมฺมา. อิเม ธมฺมา กุสลา.}

\proseitem{278}{กตโม ตสฺมิํ สมเย ผสฺโส โหติ. โย ตสฺมิํ สมเย ผสฺโส ผุสนา สมฺผุสนา สมฺผุสิตตฺตํ. อยํ ตสฺมิํ สมเย ผสฺโส โหติ.}

\proseitem{279}{กตมา ตสฺมิํ สมเย เวทนา โหติ. ยํ ตสฺมิํ สมเย ตชฺชา\-มโนวิญฺญาณธาตุ\-สมฺผสฺสชํ เจตสิกํ สาตํ เจตสิกํ สุขํ เจโต\-สมฺผสฺสชํ สาตํ สุขํ เวทยิตํ เจโต\-สมฺผสฺสชา สาตา สุขา เวทนา. อยํ ตสฺมิํ สมเย เวทนา โหติ.}

\proseitem{280}{กตมา ตสฺมิํ สมเย สญฺญา โหติ. ยา ตสฺมิํ สมเย ตชฺชา\-มโนวิญฺญาณธาตุ\-สมฺผสฺสชา สญฺญา สญฺชานนา สญฺชานิตตฺตํ. อยํ ตสฺมิํ สมเย สญฺญา โหติ.}

\proseitem{281}{กตมา ตสฺมิํ สมเย เจตนา โหติ. ยา ตสฺมิํ สมเย ตชฺชา\-มโนวิญฺญาณธาตุ\-สมฺผสฺสชา เจตนา สญฺเจตนา เจตยิตตฺตํ. อยํ ตสฺมิํ สมเย เจตนา โหติ.}

\proseitem{282}{กตมํ ตสฺมิํ สมเย จิตฺตํ โหติ. ยํ ตสฺมิํ สมเย จิตฺตํ มโน มานสํ หทยํ ปณฺฑรํ มโน มนายตนํ มนินฺทฺริยํ วิญฺญาณํ วิญฺญาณ\-กฺขนฺโธ ตชฺชา\-มโนวิญฺญาณ\-ธาตุ. อิทํ ตสฺมิํ สมเย จิตฺตํ โหติ.}

\proseitem{283}{กตโม ตสฺมิํ สมเย วิตกฺโก โหติ. โย ตสฺมิํ สมเย ตกฺโก วิตกฺโก สงฺกปฺโป อปฺปนา พฺยปฺปนา เจตโส อภินิโรปนา สมฺมาสงฺกปฺโป มคฺคงฺคํ มคฺค\-ปริยาปนฺนํ. อยํ ตสฺมิํ สมเย วิตกฺโก โหติ.}

\proseitem{284}{กตโม ตสฺมิํ สมเย วิจาโร โหติ. โย ตสฺมิํ สมเย จาโร วิจาโร อนุวิจาโร อุปวิจาโร จิตฺตสฺส อนุสนฺธานตา อนุเปกฺขนตา. อยํ ตสฺมิํ สมเย วิจาโร โหติ.}

\proseitem{285}{\csromanfolio{74}กตมา ตสฺมิํ สมเย ปีติ โหติ. ยา ตสฺมิํ สมเย ปีติ ปาโมชฺชํ อาโมทนา ปโมทนา หาโส ปหาโส วิตฺติ โอทคฺยํ อตฺตมนตา จิตฺตสฺส ปีติสมฺโพชฺฌงฺโค. อยํ ตสฺมิํ สมเย ปีติ โหติ.}

\proseitem{286}{กตมํ ตสฺมิํ สมเย สุขํ โหติ. ยํ ตสฺมิํ สมเย เจตสิกํ สาตํ เจตสิกํ สุขํ เจโต\-สมฺผสฺสชํ สาตํ สุขํ เวทยิตํ เจโต\-สมฺผสฺสชา สาตา สุขา เวทนา. อิทํ ตสฺมิํ สมเย สุขํ โหติ.}

\proseitem{287}{กตมา ตสฺมิํ สมเย จิตฺตสฺเส\-กคฺคตา โหติ. ยา ตสฺมิํ สมเย จิตฺตสฺส ฐิติ สณฺฐิติ อวฏฺฐิติ อวิสาหาโร อวิกฺเขโป อวิสาหฏ\-มานสตา สมโถ สมาธินฺทฺริยํ สมาธิพลํ สมฺมาสมาธิ สมาธิ\-สมฺโพชฺฌงฺโค มคฺคงฺคํ มคฺค\-ปริยาปนฺนํ. อยํ ตสฺมิํ สมเย จิตฺตสฺเส\-กคฺคตา โหติ.}

\proseitem{288}{กตมํ ตสฺมิํ สมเย สทฺธินฺทฺริยํ โหติ. ยา ตสฺมิํ สมเย สทฺธา สทฺทหนา โอกปฺปนา อภิปฺปสาโท สทฺธา สทฺธินฺทฺริยํ สทฺธาพลํ. อิทํ ตสฺมิํ สมเย สทฺธินฺทฺริยํ โหติ.}

\proseitem{289}{กตมํ ตสฺมิํ สมเย วีริยินฺทฺริยํ โหติ. โย ตสฺมิํ สมเย เจตสิโก วีริยารมฺโภ นิกฺกโม ปรกฺกโม อุยฺยาโม วายาโม อุสฺสาโห อุสฺโสฬฺหี ถาโม ธิติ อสิถิล\-ปรกฺกมตา อนิกฺขิตฺต\-ฉนฺทตา อนิกฺขิตฺตธุรตา ธุรสมฺปคฺคาโห วีริยํ วีริยินฺทฺริยํ วีริยพลํ สมฺมาวายาโม วีริย\-สมฺโพชฺฌงฺโค มคฺคงฺคํ มคฺค\-ปริยาปนฺนํ. อิทํ ตสฺมิํ สมเย วีริยินฺทฺริยํ โหติ.}

\proseitem{290}{กตมํ ตสฺมิํ สมเย สตินฺทฺริยํ โหติ. ยา ตสฺมิํ สมเย สติ อนุสฺสติ ปฏิสฺสติ สติ สรณตา ธารณตา อปิลาปนตา อสมฺมุสฺสนตา สติ สตินฺทฺริยํ สติพลํ สมฺมาสติ สติสมฺโพชฺฌงฺโค มคฺคงฺคํ มคฺค\-ปริยาปนฺนํ. อิทํ ตสฺมิํ สมเย สตินฺทฺริยํ โหติ.}

\proseitem{291}{กตมํ ตสฺมิํ สมเย สมาธินฺทฺริยํ โหติ. ยา ตสฺมิํ สมเย จิตฺตสฺส ฐิติ สณฺฐิติ อวฏฺฐิติ อวิสาหาโร อวิกฺเขโป \csromanfolio{75}อวิสาหฏ\-มานสตา สมโถ สมาธินฺทฺริยํ สมาธิพลํ สมฺมาสมาธิ สมาธิ\-สมฺโพชฺฌงฺโค มคฺคงฺคํ มคฺค\-ปริยาปนฺนํ. อิทํ ตสฺมิํ สมาธินฺทฺริยํ โหติ.}

\proseitem{292}{กตมํ ตสฺมิํ สมเย ปญฺญินฺทฺริยํ โหติ. ยา ตสฺมิํ สมเย ปญฺญา ปชานนา วิจโย ปวิจโย ธมฺมวิจโย สลฺลกฺขณา อุปลกฺขณา ปจฺจุ\-ปลกฺขณา ปณฺฑิจฺจํ โกสลฺลํ เนปุญฺญํ เวภพฺยา จินฺตา อุปปริกฺขา ภูรี เมธา ปริณายิกา วิปสฺสนา สมฺปชญฺญํ ปโตโท ปญฺญา ปญฺญินฺทฺริยํ ปญฺญาพลํ ปญฺญาสตฺถํ ปญฺญาปาสาโท ปญฺญาอาโลโก ปญฺญาโอภาโส ปญฺญาปชฺโชโต ปญฺญารตนํ อโมโห ธมฺมวิจโย สมฺมาทิฏฺฐิ ธมฺมวิจย\-สมฺโพชฺฌงฺโค มคฺคงฺคํ มคฺค\-ปริยาปนฺนํ. อิทํ ตสฺมิํ สมเย ปญฺญินฺทฺริยํ โหติ.}

\proseitem{293}{กตมํ ตสฺมิํ สมเย มนินฺทฺริยํ โหติ. ยํ ตสฺมิํ สมเย จิตฺตํ มโน มานสํ หทยํ ปณฺฑรํ มโน มนายตนํ มนินฺทฺริยํ วิญฺญาณํ วิญฺญาณ\-กฺขนฺโธ ตชฺชา\-มโนวิญฺญาณ\-ธาตุ. อิทํ ตสฺมิํ สมเย มนินฺทฺริยํ โหติ.}

\proseitem{294}{กตมํ ตสฺมิํ สมเย โสมนสฺสิ\-นฺทฺริยํ โหติ. ยํ ตสฺมิํ สมเย เจตสิกํ สาตํ เจตสิกํ สุขํ เจโต\-สมฺผสฺสชํ สาตํ สุขํ เวทยิตํ เจโต\-สมฺผสฺสชา สาตา สุขา เวทนา. อิทํ ตสฺมิํ สมเย โสมนสฺสิ\-นฺทฺริยํ โหติ.}

\proseitem{295}{กตมํ ตสฺมิํ สมเย ชีวิตินฺทฺริยํ โหติ. โย เตสํ อรูปีนํ ธมฺมานํ อายุ ฐิติ ยปนา ยาปนา อิริยนา วตฺตนา ปาลนา ชีวิตํ ชีวิตินฺทฺริยํ. อิทํ ตสฺมิํ สมเย ชีวิตินฺทฺริยํ โหติ.}

\proseitem{296}{กตมํ ตสฺมิํ สมเย อนญฺญาต\-ญฺญสฺสามี\-ตินฺทฺริยํ โหติ. ยา เตสํ ธมฺมานํ อนญฺญาตานํ อทิฏฺฐานํ อปฺปตฺตานํ อวิทิตานํ อสจฺฉิกตานํ สจฺฉิกิริยาย ปญฺญา ปชานนา วิจโย ปวิจโย ธมฺมวิจโย สลฺลกฺขณา อุปลกฺขณา ปจฺจุ\-ปลกฺขณา ปณฺฑิจฺจํ โกสลฺลํ เนปุญฺญํ เวภพฺยา จินฺตา อุปปริกฺขา ภูรี เมธา ปริณายิกา วิปสฺสนา สมฺปชญฺญํ ปโตโท \csromanfolio{76}ปญฺญา ปญฺญินฺทฺริยํ ปญฺญาพลํ ปญฺญาสตฺถํ ปญฺญาปาสาโท ปญฺญาอาโลโก ปญฺญาโอภาโส ปญฺญาปชฺโชโต ปญฺญารตนํ อโมโห ธมฺมวิจโย สมฺมาทิฏฺฐิ ธมฺมวิจย\-สมฺโพชฺฌงฺโค มคฺคงฺคํ มคฺค\-ปริยาปนฺนํ. อิทํ ตสฺมิํ สมเย อนญฺญาต\-ญฺญสฺสามี\-ตินฺทฺริยํ โหติ.}

\proseitem{297}{กตมา ตสฺมิํ สมเย สมฺมาทิฏฺฐิ โหติ. ยา ตสฺมิํ สมเย ปญฺญา ปชานนา วิจโย ปวิจโย ธมฺมวิจโย สลฺลกฺขณา อุปลกฺขณา ปจฺจุ\-ปลกฺขณา ปณฺฑิจฺจํ โกสลฺลํ เนปุญฺญํ เวภพฺยา จินฺตา อุปปริกฺขา ภูรี เมธา ปริณายิกา วิปสฺสนา สมฺปชญฺญํ ปโตโท ปญฺญา ปญฺญินฺทฺริยํ ปญฺญาพลํ ปญฺญาสตฺถํ ปญฺญาปาสาโท ปญฺญาอาโลโก ปญฺญาโอภาโส ปญฺญาปชฺโชโต ปญฺญารตนํ อโมโห ธมฺมวิจโย สมฺมาทิฏฺฐิ ธมฺมวิจย\-สมฺโพชฺฌงฺโค มคฺคงฺคํ มคฺค\-ปริยาปนฺนํ. อยํ ตสฺมิํ สมเย สมฺมาทิฏฺฐิ โหติ.}

\proseitem{298}{กตโม ตสฺมิํ สมเย สมฺมาสงฺกปฺโป โหติ. โย ตสฺมิํ สมเย ตกฺโก วิตกฺโก สงฺกปฺโป อปฺปนา พฺยปฺปนา เจตโส อภินิโรปนา สมฺมาสงฺกปฺโป มคฺคงฺคํ มคฺค\-ปริยาปนฺนํ. อยํ ตสฺมิํ สมเย สมฺมาสงฺกปฺโป โหติ.}

\proseitem{299}{กตมา ตสฺมิํ สมเย สมฺมาวาจา โหติ. ยา ตสฺมิํ สมเย จตูหิ วจีทุจฺจริเตหิ อารติ วิรติ ปฏิวิรติ เวรมณี อกิริยา อกรณํ อนชฺฌาปตฺติ เวลาอนติกฺกโม เสตุฆาโต สมฺมาวาจา มคฺคงฺคํ มคฺค\-ปริยาปนฺนํ. อยํ ตสฺมิํ สมเย สมฺมาวาจา โหติ.}

\proseitem{300}{กตโม ตสฺมิํ สมเย สมฺมากมฺมนฺโต โหติ. ยา ตสฺมิํ สมเย ตีหิ กาย\-ทุจฺจริเตหิ อารติ วิรติ ปฏิวิรติ เวรมณี อกิริยา อกรณํ อนชฺฌาปตฺติ เวลาอนติกฺกโม เสตุฆาโต สมฺมากมฺมนฺโต มคฺคงฺคํ มคฺค\-ปริยาปนฺนํ. อยํ ตสฺมิํ สมเย สมฺมากมฺมนฺโต โหติ.}

\proseitem{301}{กตโม ตสฺมิํ สมเย สมฺมาอาชีโว โหติ. ยา ตสฺมิํ สมเย มิจฺฉาอาชีวา อารติ วิรติ ปฏิวิรติ เวรมณี อกิริยา อกรณํ อนชฺฌาปตฺติ เวลาอนติกฺกโม เสตุฆาโต สมฺมาอาชีโว มคฺคงฺคํ มคฺค\-ปริยาปนฺนํ. อยํ ตสฺมิํ สมเย สมฺมาอาชีโว โหติ.}

\chapterheadpageread{1. จิตฺตุปฺปาท\-กณฺฑ}
\proseitem{302}{\csromanfolio{77}กตโม ตสฺมิํ สมเย สมฺมาวายาโม โหติ. โย ตสฺมิํ สมเย เจตสิโก วีริยารมฺโภ นิกฺกโม ปรกฺกโม อุยฺยาโม วายาโม อุสฺสาโห อุสฺโสฬฺหี ถาโม ธิติ อสิถิล\-ปรกฺกมตา อนิกฺขิตฺต\-ฉนฺทตา อนิกฺขิตฺตธุรตา ธุรสมฺปคฺคาโห วีริยํ วีริยินฺทฺริยํ วีริยพลํ สมฺมาวายาโม วีริย\-สมฺโพชฺฌงฺโค มคฺคงฺคํ มคฺค\-ปริยาปนฺนํ. อยํ ตสฺมิํ สมเย สมฺมาวายาโม โหติ.}

\proseitem{303}{กตมา ตสฺมิํ สมเย สมฺมาสติ โหติ. ยา ตสฺมิํ สมเย สติ อนุสฺสติ ปฏิสฺสติ สติ สรณตา ธารณตา อปิลาปนตา อสมฺมุสฺสนตา สติ สตินฺทฺริยํ สติพลํ สมฺมาสติ สติสมฺโพชฺฌงฺโค มคฺคงฺคํ มคฺค\-ปริยาปนฺนํ. อยํ ตสฺมิํ สมเย สมฺมาสติ โหติ.}

\proseitem{304}{กตโม ตสฺมิํ สมเย สมฺมาสมาธิ โหติ. ยา ตสฺมิํ สมเย จิตฺตสฺส ฐิติ สณฺฐิติ อวฏฺฐิติ อวิสาหาโร อวิกฺเขโป อวิสาหฏ\-มานสตา สมโถ สมาธินฺทฺริยํ สมาธิพลํ สมฺมาสมาธิ สมาธิ\-สมฺโพชฺฌงฺโค มคฺคงฺคํ มคฺค\-ปริยาปนฺนํ. อยํ ตสฺมิํ สมเย สมฺมาสมาธิ โหติ.}

\proseitem{305}{กตมํ ตสฺมิํ สมเย สทฺธาพลํ โหติ. ยา ตสฺมิํ สมเย สทฺธา สทฺทหนา โอกปฺปนา อภิปฺปสาโท สทฺธา สทฺธินฺทฺริยํ สทฺธาพลํ. อิทํ ตสฺมิํ สมเย สทฺธาพลํ โหติ.}

\proseitem{306}{กตมํ ตสฺมิํ สมเย วีริยพลํ โหติ. โย ตสฺมิํ สมเย เจตสิโก วีริยารมฺโภ นิกฺกโม ปรกฺกโม อุยฺยาโม วายาโม อุสฺสาโห อุสฺโสฬฺหี ถาโม ธิติ อสิถิล\-ปรกฺกมตา อนิกฺขิตฺต\-ฉนฺทตา อนิกฺขิตฺตธุรตา ธุรสมฺปคฺคาโห วีริยํ วีริยินฺทฺริยํ วีริยพลํ สมฺมาวายาโม วีริย\-สมฺโพชฺฌงฺโค มคฺคงฺคํ มคฺค\-ปริยาปนฺนํ. อิทํ ตสฺมิํ สมเย วีริยพลํ โหติ.}

\proseitem{307}{กตมํ ตสฺมิํ สมเย สติพลํ โหติ. ยา ตสฺมิํ สมเย สติ อนุสฺสติ ปฏิสฺสติ สติ สรณตา ธารณตา อปิลาปนตา อสมฺมุสฺสนตา สติ สตินฺทฺริยํ สติพลํ สมฺมาสติ \csromanfolio{78}สติสมฺโพชฺฌงฺโค มคฺคงฺคํ มคฺค\-ปริยาปนฺนํ. อิทํ ตสฺมิํ สมเย สติพลํ โหติ.}

\proseitem{308}{กตมํ ตสฺมิํ สมเย สมฺมาธิพลํ โหติ. ยา ตสฺมิํ สมเย จิตฺตสฺส ฐิติ สณฺฐิติ อวฏฺฐิติ อวิสาหาโร อวิกฺเขโป อวิสาหฏ\-มานสตา สมโถ สมาธินฺทฺริยํ สมาธิพลํ สมฺมาสมาธิ สมาธิ\-สมฺโพชฺฌงฺโค มคฺคงฺคํ มคฺค\-ปริยาปนฺนํ. อิทํ ตสฺมิํ สมเย สมาธิพลํ โหติ.}

\proseitem{309}{กตมํ ตสฺมิํ สมเย ปญฺญาพลํ โหติ. ยา ตสฺมิํ สมเย ปญฺญา ปชานนา วิจโย ปวิจโย ธมฺมวิจโย สลฺลกฺขณา อุปลกฺขณา ปจฺจุ\-ปลกฺขณา ปณฺฑิจฺจํ โกสลฺลํ เนปุญฺญํ เวภพฺยา จินฺตา อุปปริกฺขา ภูรี เมธา ปริณายิกา วิปสฺสนา สมฺปชญฺญํ ปโตโท ปญฺญา ปญฺญินฺทฺริยํ ปญฺญาพลํ ปญฺญาสตฺถํ ปญฺญาปาสาโท ปญฺญาอาโลโก ปญฺญาโอภาโส ปญฺญาปชฺโชโต ปญฺญารตนํ อโมโห ธมฺมวิจโย สมฺมาทิฏฺฐิ ธมฺมวิจย\-สมฺโพชฺฌงฺโค มคฺคงฺคํ มคฺค\-ปริยาปนฺนํ. อิทํ ตสฺมิํ สมเย ปญฺญาพลํ โหติ.}

\proseitem{310}{กตมํ ตสฺมิํ สมเย หิริพลํ โหติ. ยํ ตสฺมิํ สมเย หีรียติ หิริยิตพฺเพน หิรียติ ปาปกานํ อกุสลานํ ธมฺมานํ สมาปตฺติยา. อิทํ ตสฺมิํ สมเย หิริพลํ โหติ.}

\proseitem{311}{กตมํ ตสฺมิํ สมเย โอตฺตปฺปพลํ โหติ. ยํ ตสฺมิํ สมเย โอตฺตปฺปติ โอตฺตปฺปิตพฺเพน โอตฺตปฺปติ ปาปกานํ อกุสลานํ ธมฺมานํ สมาปตฺติยา. อิทํ ตสฺมิํ สมเย โอตฺตปฺปพลํ โหติ.}

\proseitem{312}{กตโม ตสฺมิํ สมเย อโลโภ โหติ. โย ตสฺมิํ สมเย อโลโภ อลุพฺภนา อลุพฺภิตตฺตํ อสาราโค อสารชฺชนา อสารชฺชิตตฺตํ อนภิชฺฌา อโลโภ กุสลมูลํ. อยํ ตสฺมิํ สมเย อโลโภ โหติ.}

\proseitem{313}{กตโม ตสฺมิํ สมเย อโทโส โหติ. โย ตสฺมิํ สมเย อโทโส อทุสฺสนา อทุสฺสิตตฺตํ อพฺยาปาโท อพฺยาปชฺโช อโทโส กุสลมูลํ. อยํ ตสฺมิํ สมเย อโทโส โหติ.}

\chapterheadpageread{1. จิตฺตุปฺปาท\-กณฺฑ}
\proseitem{314}{\csromanfolio{79}กตโม ตสฺมิํ สมเย อโมโห โหติ. ยา ตสฺมิํ สมเย ปญฺญา ปชานนา วิจโย ปวิจโย ธมฺมวิจโย สลฺลกฺขณา อุปลกฺขณา ปจฺจุ\-ปลกฺขณา ปณฺฑิจฺจํ โกสลฺลํ เนปุญฺญํ เวภพฺยา จินฺตา อุปปริกฺขา ภูรี เมธา ปริณายิกา วิปสฺสนา สมฺปชญฺญํ ปโตโท ปญฺญา ปญฺญินฺทฺริยํ ปญฺญาพลํ ปญฺญาสตฺถํ ปญฺญาปาสาโท ปญฺญาอาโลโก ปญฺญาโอภาโส ปญฺญาปชฺโชโต ปญฺญารตนํ อโมโห ธมฺมวิจโย สมฺมาทิฏฺฐิ ธมฺมวิจย\-สมฺโพชฺฌงฺโค มคฺคงฺคํ มคฺค\-ปริยาปนฺนํ. อยํ ตสฺมิํ สมเย อโมโห โหติ.}

\proseitem{315}{กตมา ตสฺมิํ สมเย อนภิชฺฌา โหติ. โย ตสฺมิํ สมเย อโลโภ อลุพฺภนา อลุพฺภิตตฺตํ อสาราโค อสารชฺชนา อสารชฺชิตตฺตํ อนภิชฺฌา อโลโภ กุสลมูลํ. อยํ ตสฺมิํ สมเย อนภิชฺฌา โหติ.}

\proseitem{316}{กตมา ตสฺมิํ สมเย อพฺยาปาโท โหติ. โย ตสฺมิํ สมเย อโทโส อทุสฺสนา อทุสฺสิตตฺตํ อพฺยาปาโท อพฺยาปชฺโช อโทโส กุสลมูลํ. อยํ ตสฺมิํ สมเย อพฺยาปาโท โหติ.}

\proseitem{317}{กตมา ตสฺมิํ สมเย สมฺมาทิฏฺฐิ โหติ. ยา ตสฺมิํ สมเย ปญฺญา ปชานนา วิจโย ปวิจโย ธมฺมวิจโย สลฺลกฺขณา อุปลกฺขณา ปจฺจุ\-ปลกฺขณา ปณฺฑิจฺจํ โกสลฺลํ เนปุญฺญํ เวภพฺยา จินฺตา อุปปริกฺขา ภูรี เมธา ปริณายิกา วิปสฺสนา สมฺปชญฺญํ ปโตโท ปญฺญา ปญฺญินฺทฺริยํ ปญฺญาพลํ ปญฺญาสตฺถํ ปญฺญาปาสาโท ปญฺญาอาโลโก ปญฺญาโอภาโส ปญฺญาปชฺโชโต ปญฺญารตนํ อโมโห ธมฺมวิจโย สมฺมาทิฏฺฐิ ธมฺมวิจย\-สมฺโพชฺฌงฺโค มคฺคงฺคํ มคฺค\-ปริยาปนฺนํ. อยํ ตสฺมิํ สมเย สมฺมาทิฏฺฐิ โหติ.}

\proseitem{318}{กตมา ตสฺมิํ สมเย หิรี โหติ. ยํ ตสฺมิํ สมเย หิรียติ หิริยิตพฺเพน หิรียติ ปาปกานํ อกุสลานํ ธมฺมานํ สมาปตฺติยา. อยํ ตสฺมิํ สมเย หิรี โหติ.}

\proseitem{319}{กตมํ ตสฺมิํ สมเย โอตฺตปฺปํ โหติ. ยํ ตสฺมิํ สมเย โอตฺตปฺปติ โอตฺตปฺปิตพฺเพน โอตฺตปฺปติ ปาปกานํ อกุสลานํ ธมฺมานํ สมาปตฺติยา. อิทํ ตสฺมิํ สมเย โอตฺตปฺปํ โหติ.}

\proseitem{320}{\csromanfolio{80}กตมา ตสฺมิํ สมเย กายปสฺสทฺธิ โหติ. ยา ตสฺมิํ สมเย เวทนา\-กฺขนฺธสฺส สญฺญา\-กฺขนฺธสฺส สงฺขาร\-กฺขนฺธสฺส ปสฺสทฺธิ ปฏิปสฺสทฺธิ ปสฺสมฺภนา ปฏิปสฺสมฺภนา ปฏิปสฺสมฺภิ\-ตตฺตํ ปสฺสทฺธิ\-สมฺโพชฺฌงฺโค. อยํ ตสฺมิํ สมเย กายปสฺสทฺธิ โหติ.}

\proseitem{321}{กตมา ตสฺมิํ สมเย จิตฺตปสฺสทฺธิ โหติ. ยา ตสฺมิํ สมเย วิญฺญาณ\-กฺขนฺธสฺส ปสฺสทฺธิ ปฏิปสฺสทฺธิ ปสฺสมฺภนา ปฏิปสฺสมฺภนา ปฏิปสฺสมฺภิ\-ตตฺตํ ปสฺสทฺธิ\-สมฺโพชฺฌงฺโค. อยํ ตสฺมิํ สมเย จิตฺตปสฺสทฺธิ โหติ.}

\proseitem{322}{กตมา ตสฺมิํ สมเย กายลหุตา โหติ. ยา ตสฺมิํ สมเย เวทนา\-กฺขนฺธสฺส สญฺญา\-กฺขนฺธสฺส สงฺขาร\-กฺขนฺธสฺส ลหุตา ลหุปริณามตา อทนฺธนตา อวิตฺถนตา. อยํ ตสฺมิํ สมเย กายลหุตา โหติ.}

\proseitem{323}{กตมา ตสฺมิํ สมเย จิตฺตลหุตา โหติ. ยา ตสฺมิํ สมเย วิญฺญาณ\-กฺขนฺธสฺส ลหุตา ลหุปริณามตา อทนฺธนตา อวิตฺถนตา. อยํ ตสฺมิํ สมเย จิตฺตลหุตา โหติ.}

\proseitem{324}{กตมา ตสฺมิํ สมเย กายมุทุตา โหติ. ยา ตสฺมิํ สมเย เวทนา\-กฺขนฺธสฺส สญฺญา\-กฺขนฺธสฺส สงฺขาร\-กฺขนฺธสฺส มุทุตา มทฺทวตา อกกฺขฬตา อกถินตา. อยํ ตสฺมิํ สมเย กายมุทุตา โหติ.}

\proseitem{325}{กตมา ตสฺมิํ สมเย จิตฺตมุทุตา โหติ. ยา ตสฺมิํ สมเย วิญฺญาณ\-กฺขนฺธสฺส มุทุตา มทฺทวตา อกกฺขฬตา อกถินตา อยํ ตสฺมิํ สมเย จิตฺตมุทุตา โหติ.}

\proseitem{326}{กตมา ตสฺมิํ สมเย กายกมฺมญฺญตา โหติ. ยา ตสฺมิํ สมเย เวทนา\-กฺขนฺธสฺส สญฺญา\-กฺขนฺธสฺส สงฺขาร\-กฺขนฺธสฺส กมฺมญฺญตา กมฺมญฺญตฺตํ กมฺมญฺญภาโว. อยํ ตสฺมิํ สมเย กายกมฺมญฺญตา โหติ.}

\proseitem{327}{กตมา ตสฺมิํ สมเย จิตฺต\-กมฺมญฺญตา โหติ. ยา ตสฺมิํ สมเย วิญฺญาณ\-กฺขนฺธสฺส กมฺมญฺญตา กมฺมญฺญตฺตํ กมฺมญฺญภาโว. อยํ ตสฺมิํ สมเย จิตฺต\-กมฺมญฺญตา โหติ.}

\chapterheadpageread{1. จิตฺตุปฺปาท\-กณฺฑ}
\proseitem{328}{\csromanfolio{81}กตมา ตสฺมิํ สมเย กายปาคุญฺญตา โหติ. ยา ตสฺมิํ สมเย เวทนา\-กฺขนฺธสฺส สญฺญา\-กฺขนฺธสฺส สงฺขาร\-กฺขนฺธสฺส ปคุณตา ปคุณตฺตํ ปคุณภาโว. อยํ ตสฺมิํ สมเย กายปาคุญฺญตา โหติ.}

\proseitem{329}{กตมา ตสฺมิํ สมเย จิตฺตปาคุญฺญตา โหติ. ยา ตสฺมิํ สมเย วิญฺญาณ\-กฺขนฺธสฺส ปคุณตา ปคุณตฺตํ ปคุณภาโว. อยํ ตสฺมิํ สมเย จิตฺตปาคุญฺญตา โหติ.}

\proseitem{330}{กตมา ตสฺมิํ สมเย กายุชุกตา โหติ. ยา ตสฺมิํ สมเย เวทนา\-กฺขนฺธสฺส สญฺญา\-กฺขนฺธสฺส สงฺขาร\-กฺขนฺธสฺส อุชุตา อุชุกตา อชิมฺหตา อวงฺกตา อกุฏิลตา. อยํ ตสฺมิํ สมเย กายุชุกตา โหติ.}

\proseitem{331}{กตมา ตสฺมิํ สมเย จิตฺตุชุกตา โหติ. ยา ตสฺมิํ สมเย วิญฺญาณ\-กฺขนฺธสฺส อุชุตา อุชุกตา อชิมฺหตา อวงฺกตา อกุฏิลตา. อยํ ตสฺมิํ สมเย จิตฺตุชุกตา โหติ.}

\proseitem{332}{กตมา ตสฺมิํ สมเย สติ โหติ. ยา ตสฺมิํ สมเย สติ อนุสฺสติ ปฏิสฺสติ สติ สรณตา ธารณตา อปิลาปนตา อสมฺมุสฺสนตา สติ สตินฺทฺริยํ สติพลํ สมฺมาสติ สติสมฺโพชฺฌงฺโค มคฺคงฺคํ มคฺค\-ปริยาปนฺนํ. อยํ ตสฺมิํ สมเย สติ โหติ.}

\proseitem{333}{กตมํ ตสฺมิํ สมเย สมฺปชญฺญํ โหติ. ยา ตสฺมิํ สมเย ปญฺญา ปชานนา วิจโย ปวิจโย ธมฺมวิจโย สลฺลกฺขณา อุปลกฺขณา ปจฺจุ\-ปลกฺขณา ปณฺฑิจฺจํ โกสลฺลํ เนปุญฺญํ เวภพฺยา จินฺตา อุปปริกฺขา ภูรี เมธา ปริณายิกา วิปสฺสนา สมฺปชญฺญํ ปโตโท ปญฺญา ปญฺญินฺทฺริยํ ปญฺญาพลํ ปญฺญาสตฺถํ ปญฺญาปาสาโท ปญฺญาอาโลโก ปญฺญาโอภาโส ปญฺญาปชฺโชโต ปญฺญารตนํ อโมโห ธมฺมวิจโย สมฺมาทิฏฺฐิ ธมฺมวิจย\-สมฺโพชฺฌงฺโค มคฺคงฺคํ มคฺค\-ปริยาปนฺนํ. อิทํ ตสฺมิํ สมเย สมฺปชญฺญํ โหติ.}

\proseitem{334}{กตโม ตสฺมิํ สมเย สมโถ โหติ. ยา ตสฺมิํ สมเย จิตฺตสฺส ฐิติ สณฺฐิติ อวฏฺฐิติ อวิสาหาโร อวิกฺเขโป อวิสาหฏ\-มานสตา สมโถ สมาธินฺทฺริยํ สมาธิพลํ สมฺมาสมาธิ สมาธิ\-สมฺโพชฺฌงฺโค มคฺคงฺคํ มคฺค\-ปริยาปนฺนํ. อยํ ตสฺมิํ สมเย สมโถ โหติ.}

\proseitem{335}{\csromanfolio{82}กตมา ตสฺมิํ สมเย วิปสฺสนา โหติ. ยา ตสฺมิํ สมเย ปญฺญา ปชานนา วิจโย ปวิจโย ธมฺมวิจโย สลฺลกฺขณา อุปลกฺขณา ปจฺจุ\-ปลกฺขณา ปณฺฑิจฺจํ โกสลฺลํ เนปุญฺญํ เวภพฺยา จินฺตา อุปปริกฺขา ภูรี เมธา ปริณายิกา วิปสฺสนา สมฺปชญฺญํ ปโตโท ปญฺญา ปญฺญินฺทฺริยํ ปญฺญาพลํ ปญฺญาสตฺถํ ปญฺญาปาสาโท ปญฺญาอาโลโก ปญฺญาโอภาโส ปญฺญาปชฺโชโต ปญฺญารตนํ อโมโห ธมฺมวิจโย สมฺมาทิฏฺฐิ ธมฺมวิจย\-สมฺโพชฺฌงฺโค มคฺคงฺคํ มคฺค\-ปริยาปนฺนํ. อยํ ตสฺมิํ สมเย วิปสฺสนา โหติ.}

\proseitem{336}{กตโม ตสฺมิํ สมเย ปคฺคาโห โหติ. โย ตสฺมิํ สมเย เจตสิโก วีริยารมฺโภ นิกฺกโม ปรกฺกโม อุยฺยาโม วายาโม อุสฺสาโห อุสฺโสฬฺหี ถาโม ธิติ อสิถิล\-ปรกฺกมตา อนิกฺขิตฺต\-ฉนฺทตา อนิกฺขิตฺตธุรตา ธุรสมฺปคฺคาโห วีริยํ วีริยินฺทฺริยํ วีริยพลํ สมฺมาวายาโม วีริย\-สมฺโพชฺฌงฺโค มคฺคงฺคํ มคฺค\-ปริยาปนฺนํ. อยํ ตสฺมิํ สมเย ปคฺคาโห โหติ.}

\proseitem{337}{กตโม ตสฺมิํ สมเย อวิกฺเขโป โหติ. ยา ตสฺมิํ สมเย จิตฺตสฺส ฐิติ สณฺฐิติ อวฏฺฐิติ อวิสาหาโร อวิกฺเขโป อวิสาหฏ\-มานสตา สมโถ สมาธินฺทฺริยํ สมาธิพลํ สมฺมาสมาธิ สมาธิ\-สมฺโพชฺฌงฺโค มคฺคงฺคํ มคฺค\-ปริยาปนฺนํ. อยํ ตสฺมิํ สมเย อวิกฺเขโป โหติ. เย วา ปน ตสฺมิํ สมเย อญฺเญปิ อตฺถิ ปฏิจฺจ\-สมุปฺปนฺนา อรูปิโน ธมฺมา. อิเม ธมฺมา กุสลา.}

\prose{ตสฺมิํ โข ปน สมเย จตฺตาโร ขนฺธา โหนฺติ, ทฺวายตนานิ โหนฺติ, เทฺว ธาตุโย โหนฺติ, ตโย อาหารา โหนฺติ, นวินฺทฺริยานิ โหนฺติ, ปญฺจงฺคิกํ ฌานํ โหติ, อฏฺฐงฺคิโก มคฺโค โหติ, สตฺต พลานิ โหนฺติ, ตโย เหตู โหนฺติ, เอโก ผสฺโส โหติ, เอกา เวทนา โหติ, เอกา สญฺญา โหติ, เอกา เจตนา โหติ, เอกํ จิตฺตํ โหติ, เอโก เวทนากฺขนฺโธ โหติ, เอโก สญฺญากฺขนฺโธ โหติ, เอโก สงฺขาร\-กฺขนฺโธ โหติ, เอโก วิญฺญาณ\-กฺขนฺโธ โหติ, เอกํ มนายตนํ โหติ, เอกํ มนินฺทฺริยํ โหติ, เอกา มโนวิญฺญาณ\-ธาตุ โหติ, เอกํ ธมฺมายตนํ โหติ, เอกา ธมฺมธาตุ โหติ. เย วา ปน ตสฺมิํ สมเย อญฺเญปิ อตฺถิ ปฏิจฺจ\-สมุปฺปนฺนา อรูปิโน ธมฺมา. อิเม ธมฺมา กุสลา ฯเปฯ}

\chapterheadpageread{1. จิตฺตุปฺปาท\-กณฺฑ}
\proseitem{338}{\csromanfolio{83}กตโม ตสฺมิํ สมเย สงฺขาร\-กฺขนฺโธ โหติ. ผสฺโส เจตนา วิตกฺโก วิจาโร ปีติ จิตฺตสฺเส\-กคฺคตา สทฺธินฺทฺริยํ วีริยินฺทฺริยํ สตินฺทฺริยํ สมาธินฺทฺริยํ ปญฺญินฺทฺริยํ ชีวิตินฺทฺริยํ อนญฺญาต\-ญฺญสฺสามี\-ตินฺทฺริยํ สมฺมาทิฏฺฐิ สมฺมาสงฺกปฺโป สมฺมาวาจา สมฺมากมฺมนฺโต สมฺมาอาชีโว สมฺมาวายาโม สมฺมาสติ สมฺมาสมาธิ สทฺธาพลํ วีริยพลํ สติพลํ สมาธิพลํ ปญฺญาพลํ หิริพลํ โอตฺตปฺปพลํ อโลโภ อโทโส อโมโห อนภิชฺฌา อพฺยาปาโท สมฺมาทิฏฺฐิ หิรี โอตฺตปฺปํ กายปสฺสทฺธิ จิตฺตปสฺสทฺธิ กายลหุตา จิตฺตลหุตา กายมุทุตา จิตฺตมุทุตา กายกมฺมญฺญตา จิตฺต\-กมฺมญฺญตา กายปาคุญฺญตา จิตฺตปาคุญฺญตา กายุชุกตา จิตฺตุชุกตา สติ สมฺปชญฺญํ สมโถ วิปสฺสนา ปคฺคาโห อวิกฺเขโป. เย วา ปน ตสฺมิํ สมเย อญฺเญปิ อตฺถิ ปฏิจฺจ\-สมุปฺปนฺนา อรูปิโน ธมฺมา ฐเปตฺวา เวทนา\-กฺขนฺธํ ฐเปตฺวา สญฺญากฺขนฺธํ ฐเปตฺวา วิญฺญาณ\-กฺขนฺธํ. อยํ ตสฺมิํ สมเย สงฺขาร\-กฺขนฺโธ โหติ ฯเปฯ. อิเม ธมฺมา กุสลา.}

\proseitem{339}{กตเม ธมฺมา กุสลา. ยสฺมิํ สมเย โลกุตฺตรํ ฌานํ ภาเวติ นิยฺยานิกํ อปจยคามิํ ทิฏฺฐิคตานํ ปหานาย ปฐมาย ภูมิยา ปตฺติยา วิวิจฺเจว กาเมหิ ฯเปฯ ปฐมํ ฌานํ อุปสมฺปชฺช วิหรติ ทุกฺข\-ปฏิปทํ ขิปฺปาภิญฺญํ. ตสฺมิํ สมเย ผสฺโส โหติ ฯเปฯ อวิกฺเขโป โหติ ฯเปฯ. อิเม ธมฺมา กุสลา.}

\proseitem{340}{กตเม ธมฺมา กุสลา. ยสฺมิํ สมเย โลกุตฺตรํ ฌานํ ภาเวติ นิยฺยานิกํ อปจยคามิํ ทิฏฺฐิคตานํ ปหานาย ปฐมาย ภูมิยา ปตฺติยา วิวิจฺเจว กาเมหิ ฯเปฯ ปฐมํ ฌานํ อุปสมฺปชฺช วิหรติ สุขปฏิปทํ ทนฺธาภิญฺญํ. ตสฺมิํ สมเย ผสฺโส โหติ ฯเปฯ อวิกฺเขโป โหติ ฯเปฯ. อิเม ธมฺมา กุสลา.}

\proseitem{341}{กตเม ธมฺมา กุสลา. ยสฺมิํ สมเย โลกุตฺตรํ ฌานํ ภาเวติ นิยฺยานิกํ อปจยคามิํ ทิฏฺฐิคตานํ ปหานาย ปฐมาย ภูมิยา ปตฺติยา วิวิจฺเจว กาเมหิ ฯเปฯ ปฐมํ ฌานํ อุปสมฺปชฺช วิหรติ สุขปฏิปทํ ขิปฺปาภิญฺญํ. ตสฺมิํ สมเย ผสฺโส โหติ ฯเปฯ อวิกฺเขโป โหติ ฯเปฯ. อิเม ธมฺมา กุสลา.}

\proseitem{342}{กตเม ธมฺมา กุสลา. ยสฺมิํ สมเย โลกุตฺตรํ ฌานํ ภาเวติ นิยฺยานิกํ อปจยคามิํ ทิฏฺฐิคตานํ ปหานาย ปฐมาย ภูมิยา \csromanfolio{84}ปตฺติยา วิตกฺก\-วิจารานํ วูปสมา ฯเปฯ ทุติยํ ฌานํ ฯเปฯ ตติยํ ฌานํ ฯเปฯ จตุตฺถํ ฌานํ ฯเปฯ ปฐมํ ฌานํ ฯเปฯ ปญฺจมํ ฌานํ อุปสมฺปชฺช วิหรติ ทุกฺข\-ปฏิปทํ ทนฺธาภิญฺญํ ฯเปฯ ทุกฺข\-ปฏิปทํ ขิปฺปาภิญฺญํ ฯเปฯ สุขปฏิปทํ ทนฺธาภิญฺญํ ฯเปฯ สุขปฏิปทํ ขิปฺปาภิญฺญํ. ตสฺมิํ สมเย ผสฺโส โหติ ฯเปฯ อวิกฺเขโป โหติ ฯเปฯ. อิเม ธมฺมา กุสลา.}

\vspace{\tipitakaparskip}
\csromancenter{สุทฺธิก\-ปฏิปทา.}
\csromanmatikaanchor{mtk.45}
\csromantocmarkref{subsection}{สุญฺญต}{mtk.45}
\bookmark[dest={mtk.45},level=2]{สุญฺญต}
\csromanheader{2}{\textbf{สุญฺญต}}
\proseitem{343}{กตเม ธมฺมา กุสลา. ยสฺมิํ สมเย โลกุตฺตรํ ฌานํ ภาเวติ นิยฺยานิกํ อปจยคามิํ ทิฏฺฐิคตานํ ปหานาย ปฐมาย ภูมิยา ปตฺติยา วิวิจฺเจว กาเมหิ ฯเปฯ ปฐมํ ฌานํ อุปสมฺปชฺช วิหรติ สุญฺญตํ. ตสฺมิํ สมเย ผสฺโส โหติ ฯเปฯ อวิกฺเขโป โหติ ฯเปฯ. อิเม ธมฺมา กุสลา.}

\proseitem{344}{กตเม ธมฺมา กุสลา. ยสฺมิํ สมเย โลกุตฺตรํ ฌานํ ภาเวติ นิยฺยานิกํ อปจยคามิํ ทิฏฺฐิคตานํ ปหานาย ปฐมาย ภูมิยา ปตฺติยา วิตกฺก\-วิจารานํ วูปสมา ฯเปฯ ทุติยํ ฌานํ ฯเปฯ ตติยํ ฌานํ ฯเปฯ จตุตฺถํ ฌานํ ฯเปฯ ปฐมํ ฌานํ ฯเปฯ ปญฺจมํ ฌานํ อุปสมฺปชฺช วิหรติ สุญฺญตํ. ตสฺมิํ สมเย ผสฺโส โหติ ฯเปฯ อวิกฺเขโป โหติ ฯเปฯ. อิเม ธมฺมา กุสลา.}

\vspace{\tipitakaparskip}
\csromancenter{\textbf{สุญฺญต}ํ.}
\csromanmatikaanchor{mtk.46}
\csromantocmarkref{subsection}{สุญฺญตมูลกปฏิปทา}{mtk.46}
\bookmark[dest={mtk.46},level=2]{สุญฺญตมูลกปฏิปทา}
\csromanheader{2}{สุญฺญต\-มูลก\-ปฏิปทา}
\proseitem{345}{กตเม ธมฺมา กุสลา. ยสฺมิํ สมเย โลกุตฺตรํ ฌานํ ภาเวติ นิยฺยานิกํ อปจยคามิํ ทิฏฺฐิคตานํ ปหานาย ปฐมาย ภูมิยา ปตฺติยา วิวิจฺเจว กาเมหิ ฯเปฯ ปฐมํ ฌานํ อุปสมฺปชฺช วิหรติ ทุกฺข\-ปฏิปทํ ทนฺธาภิญฺญํ สุญฺญตํ. ตสฺมิํ สมเย ผสฺโส โหติ ฯเปฯ อวิกฺเขโป โหติ ฯเปฯ. อิเม ธมฺมา กุสลา.}

\chapterheadpageread{1. จิตฺตุปฺปาท\-กณฺฑ}
\proseitem{346}{\csromanfolio{85}กตเม ธมฺมา กุสลา. ยสฺมิํ สมเย โลกุตฺตรํ ฌานํ ภาเวติ นิยฺยานิกํ อปจยคามิํ ทิฏฺฐิคตานํ ปหานาย ปฐมาย ภูมิยา ปตฺติยา วิวิจฺเจว กาเมหิ ฯเปฯ ปฐมํ ฌานํ อุปสมฺปชฺช วิหรติ ทุกฺข\-ปฏิปทํ ขิปฺปาภิญฺญํ สุญฺญตํ. ตสฺมิํ สมเย ผสฺโส โหติ ฯเปฯ อวิกฺเขโป โหติ ฯเปฯ. อิเม ธมฺมา กุสลา.}

\proseitem{347}{กตเม ธมฺมา กุสลา. ยสฺมิํ สมเย โลกุตฺตรํ ฌานํ ภาเวติ นิยฺยานิกํ อปจยคามิํ ทิฏฺฐิคตานํ ปหานาย ปฐมาย ภูมิยา ปตฺติยา วิวิจฺเจว กาเมหิ ฯเปฯ ปฐมํ ฌานํ อุปสมฺปชฺช วิหรติ สุขปฏิปทํ ทนฺธาภิญฺญํ สุญฺญตํ. ตสฺมิํ สมเย ผสฺโส โหติ ฯเปฯ อวิกฺเขโป โหติ ฯเปฯ. อิเม ธมฺมา กุสลา.}

\proseitem{348}{กตเม ธมฺมา กุสลา. ยสฺมิํ สมเย โลกุตฺตรํ ฌานํ ภาเวติ นิยฺยานิกํ อปจยคามิํ ทิฏฺฐิคตานํ ปหานาย ปฐมาย ภูมิยา ปตฺติยา วิวิจฺเจว กาเมหิ ฯเปฯ ปฐมํ ฌานํ อุปสมฺปชฺช วิหรติ สุขปฏิปทํ ขิปฺปาภิญฺญํ สุญฺญตํ. ตสฺมิํ สมเย ผสฺโส โหติ ฯเปฯ อวิกฺเขโป โหติ ฯเปฯ. อิเม ธมฺมา กุสลา.}

\proseitem{349}{กตเม ธมฺมา กุสลา. ยสฺมิํ สมเย โลกุตฺตรํ ฌานํ ภาเวติ นิยฺยานิกํ อปจยคามิํ ทิฏฺฐิคตานํ ปหานาย ปฐมาย ภูมิยา ปตฺติยา วิตกฺก\-วิจารานํ วูปสมา ฯเปฯ ทุติยํ ฌานํ ฯเปฯ ตติยํ ฌานํ ฯเปฯ จตุตฺถํ ฌนํ ฯเปฯ ปฐมํ ฌานํ ฯเปฯ ปญฺจมํ ฌานํ อุปสมฺปชฺช วิหรติ ทุกฺข\-ปฏิปทํ ทนฺธาภิญฺญํ สุญฺญตํ ฯเปฯ ทุกฺข\-ปฏิปทํ ขิปฺปาภิญฺญํ สุญฺญตํ ฯเปฯ สุขปฏิปทํ ทนฺธาภิญฺญํ สุญฺญตํ ฯเปฯ สุขปฏิปทํ ขิปฺปาภิญฺญํ สุญฺญตํ. ตสฺมิํ สมเย ผสฺโส โหติ ฯเปฯ อวิกฺเขโป โหติ ฯเปฯ. อิเม ธมฺมา กุสลา.}

\vspace{\tipitakaparskip}
\csromancenter{สุญฺญต\-มูลก\-ปฏิปทา.}
\csromanmatikaanchor{mtk.47}
\csromantocmarkref{subsection}{อปฺปณิหิต}{mtk.47}
\bookmark[dest={mtk.47},level=2]{อปฺปณิหิต}
\csromanheader{2}{\textbf{อปฺปณิหิต}}
\proseitem{350}{กตเม ธมฺมา กุสลา. ยสฺมิํ สมเย โลกุตฺตรํ ฌานํ ภาเวติ นิยฺยานิกํ อปจยคามิํ ทิฏฺฐิคตานํ ปหานาย ปฐมาย ภูมิยา ปตฺติยา วิวิจฺเจว กาเมหิ ฯเปฯ ปฐมํ ฌานํ อุปสมฺปชฺช วิหรติ \csromanfolio{86}อปฺปณิหิตํ. ตสฺมิํ สมเย ผสฺโส โหติ ฯเปฯ อวิกฺเขโป โหติ ฯเปฯ. อิเม ธมฺมา กุสลา.}

\proseitem{351}{กตเม ธมฺมา กุสลา. ยสฺมิํ สมเย โลกุตฺตรํ ฌานํ ภาเวติ นิยฺยานิกํ อปจยคามิํ ทิฏฺฐิคตานํ ปหานาย ปฐมาย ภูมิยา ปตฺติยา วิตกฺก\-วิจารานํ วูปสมา ฯเปฯ ทุติยํ ฌานํ ฯเปฯ ตติยํ ฌานํ ฯเปฯ จตุตฺถํ ฌานํ ฯเปฯ ปฐมํ ฌานํ ฯเปฯ ปญฺจมํ ฌานํ อุปสมฺปชฺช วิหรติ อปฺปณิหิตํ. ตสฺมิํ สมเย ผสฺโส โหติ ฯเปฯ อวิกฺเขโป โหติ ฯเปฯ. อิเม ธมฺมา กุสลา.}

\vspace{\tipitakaparskip}
\csromancenter{อปฺปณิหิตํ.}
\csromanmatikaanchor{mtk.48}
\csromantocmarkref{subsection}{อปฺปณิหิตมูลกปฏิปทา}{mtk.48}
\bookmark[dest={mtk.48},level=2]{อปฺปณิหิตมูลกปฏิปทา}
\csromanheader{2}{อปฺปณิหิต\-มูลก\-ปฏิปทา}
\proseitem{352}{กตเม ธมฺมา กุสลา. ยสฺมิํ สมเย โลกุตฺตรํ ฌานํ ภาเวติ นิยฺยานิกํ อปจยคามิํ ทิฏฺฐิคตานํ ปหานาย ปฐมาย ภูมิยา ปตฺติยา วิวิจฺเจว กาเมหิ ฯเปฯ ปฐมํ ฌานํ อุปสมฺปชฺช วิหรติ ทุกฺข\-ปฏิปทํ ทนฺธาภิญฺญํ อปฺปณิหิตํ. ตสฺมิํ สมเย ผสฺโส โหติ ฯเปฯ อวิกฺเขโป โหติ ฯเปฯ. อิเม ธมฺมา กุสลา.}

\proseitem{353}{กตเม ธมฺมา กุสลา. ยสฺมิํ สมเย โลกุตฺตรํ ฌานํ ภาเวติ นิยฺยานิกํ อปจยคามิํ ทิฏฺฐิคตานํ ปหานาย ปฐมาย ภูมิยา ปตฺติยา วิวิจฺเจว กาเมหิ ฯเปฯ ปฐมํ ฌานํ อุปสมฺปชฺช วิหรติ ทุกฺข\-ปฏิปทํ ขิปฺปาภิญฺญํ อปฺปณิหิตํ. ตสฺมิํ สมเย ผสฺโส โหติ ฯเปฯ อวิกฺเขโป โหติ ฯเปฯ. อิเม ธมฺมา กุสลา.}

\proseitem{354}{กตเม ธมฺมา กุสลา. ยสฺมิํ สมเย โลกุตฺตรํ ฌานํ ภาเวติ นิยฺยานิกํ อปจยคามิํ ทิฏฺฐิคตานํ ปหานาย ปฐมาย ภูมิยา ปตฺติยา วิวิจฺเจว กาเมหิ ฯเปฯ ปฐมํ ฌานํ อุปสมฺปชฺช วิหรติ สุขปฏิปทํ ทนฺธาภิญฺญํ อปฺปณิหิตํ. ตสฺมิํ สมเย ผสฺโส โหติ ฯเปฯ อวิกฺเขโป โหติ ฯเปฯ. อิเม ธมฺมา กุสลา.}

\chapterheadpageread{1. จิตฺตุปฺปาท\-กณฺฑ}
\proseitem{355}{\csromanfolio{87}กตเม ธมฺมา กุสลา. ยสฺมิํ สมเย โลกุตฺตรํ ฌานํ ภาเวติ นิยฺยานิกํ อปจยคามิํ ทิฏฺฐิคตานํ ปหานาย ปฐมาย ภูมิยา ปตฺติยา วิวิจฺเจว กาเมหิ ฯเปฯ ปฐมํ ฌานํ อุปสมฺปชฺช วิหรติ สุขปฏิปทํ ขิปฺปาภิญฺญํ อปฺปณิหิตํ. ตสฺมิํ สมเย ผสฺโส โหติ ฯเปฯ อวิกฺเขโป โหติ ฯเปฯ. อิเม ธมฺมา กุสลา.}

\proseitem{356}{กตเม ธมฺมา กุสลา. ยสฺมิํ สมเย โลกุตฺตรํ ฌานํ ภาเวติ นิยฺยานิกํ อปจยคามิํ ทิฏฺฐิคตานํ ปหานาย ปฐมาย ภูมิยา ปตฺติยา วิตกฺก\-วิจารานํ วูปสมา ฯเปฯ ทุติยํ ฌานํ ฯเปฯ ตติยํ ฌานํ ฯเปฯ จตุตฺถํ ฌานํ ฯเปฯ ปฐมํ ฌานํ ฯเปฯ ปญฺจมํ ฌานํ อุปสมฺปชฺช วิหรติ ทุกฺข\-ปฏิปทํ ทนฺธาภิญฺญํ อปฺปณิหิตํ ฯเปฯ ทุกฺข\-ปฏิปทํ ขิปฺปาภิญฺญํ อปฺปณิหิตํ ฯเปฯ สุขปฏิปทํ ทนฺธาภิญฺญํ อปฺปณิหิตํ ฯเปฯ สุขปฏิปทํ ขิปฺปาภิญฺญํ อปฺปณิหิตํ. ตสฺมิํ สมเย ผสฺโส โหติ ฯเปฯ อวิกฺเขโป โหติ ฯเปฯ. อิเม ธมฺมา กุสลา.}

\vspace{\tipitakaparskip}
\csromancenter{อปฺปณิหิต\-มูลก\-ปฏิปทา.}
\csromanmatikaanchor{mtk.49}
\csromantocmarkref{subsection}{วีสติ มหานย}{mtk.49}
\bookmark[dest={mtk.49},level=2]{วีสติ มหานย}
\csromanheader{2}{\textbf{วีสติ มหานย}}
\proseitem{357}{กตเม ธมฺมา กุสลา. ยสฺมิํ สมเย โลกุตฺตรํ มคฺคํ ภาเวติ ฯเปฯ โลกุตฺตรํ สติปฏฺฐานํ ภาเวติ ฯเปฯ โลกุตฺตรํ สมฺมปฺปธานํ ภาเวติ ฯเปฯ โลกุตฺตรํ อิทฺธิปาทํ ภาเวติ ฯเปฯ โลกุตฺตรํ อินฺทฺริยํ ภาเวติ ฯเปฯ โลกุตฺตรํ พลํ ภาเวติ ฯเปฯ โลกุตฺตรํ โพชฺฌงฺคํ ภาเวติ ฯเปฯ โลกุตฺตรํ สจฺจํ ภาเวติ ฯเปฯ โลกุตฺตรํ สมถํ ภาเวติ ฯเปฯ โลกุตฺตรํ ธมฺมํ ภาเวติ ฯเปฯ โลกุตฺตรํ ขนฺธํ ภาเวติ ฯเปฯ โลกุตฺตรํ อายตนํ ภาเวติ ฯเปฯ โลกุตฺตรํ ธาตุํ ภาเวติ ฯเปฯ โลกุตฺตรํ อาหารํ ภาเวติ ฯเปฯ โลกุตฺตรํ ผสฺสํ ภาเวติ ฯเปฯ โลกุตฺตรํ เวทนํ ภาเวติ ฯเปฯ โลกุตฺตรํ สญฺญํ ภาเวติ ฯเปฯ โลกุตฺตรํ เจตนํ ภาเวติ ฯเปฯ โลกุตฺตรํ จิตฺตํ ภาเวติ นิยฺยานิกํ อปจยคามิํ ทิฏฺฐิคตานํ ปหานาย ปฐมาย ภูมิยา ปตฺติยา วิวิจฺเจว กาเมหิ ฯเปฯ ปฐมํ ฌานํ อุปสมฺปชฺช วิหรติ ทุกฺข\-ปฏิปทํ ทนฺธาภิญฺญํ. \csromanfolio{88}ตสฺมิํ สมเย ผสฺโส โหติ ฯเปฯ อวิกฺเขโป โหติ ฯเปฯ. อิเม ธมฺมา กุสลา.}

\vspace{\tipitakaparskip}
\csromancenter{วีสติ มหานยา.}
\csromanmatikaanchor{mtk.50}
\csromantocmarkref{subsection}{อธิปติ}{mtk.50}
\bookmark[dest={mtk.50},level=2]{อธิปติ}
\csromanheader{2}{\textbf{อธิปติ}}
\proseitem{358}{กตเม ธมฺมา กุสลา. ยสฺมิํ สมเย โลกุตฺตรํ ฌานํ ภาเวติ นิยฺยานิกํ อปจยคามิํ ทิฏฺฐิคตานํ ปหานาย ปฐมาย ภูมิยา ปตฺติยา วิวิจฺเจว กาเมหิ ฯเปฯ ปฐมํ ฌานํ อุปสมฺปชฺช วิหรติ ทุกฺข\-ปฏิปทํ ทนฺธาภิญฺญํ ฉนฺทา\-ธิปเตยฺยํ ฯเปฯ วีริยา\-ธิปเตยฺยํ ฯเปฯ จิตฺตา\-ธิปเตยฺยํ ฯเปฯ วีมํสา\-ธิปเตยฺยํ. ตสฺมิํ สมเย ผสฺโส โหติ ฯเปฯ อวิกฺเขโป โหติ ฯเปฯ. อิเม ธมฺมา กุสลา.}

\proseitem{359}{กตเม ธมฺมา กุสลา. ยสฺมิํ สมเย โลกุตฺตรํ ฌานํ ภาเวติ นิยฺยานิกํ อปจยคามิํ ทิฏฺฐิคตานํ ปหานาย ปฐมาย ภูมิยา ปตฺติยา วิตกฺก\-วิจารานํ วูปสมา ฯเปฯ ทุติยํ ฌานํ ฯเปฯ ตติยํ ฌานํ ฯเปฯ จตุตฺถํ ฌานํ ฯเปฯ ปฐมํ ฌานํ ฯเปฯ ปญฺจมํ ฌานํ อุปสมฺปชฺช วิหรติ ทุกฺข\-ปฏิปทํ ทนฺธาภิญฺญํ ฉนฺทา\-ธิปเตยฺยํ ฯเปฯ วีริยา\-ธิปเตยฺยํ ฯเปฯ จิตฺตา\-ธิปเตยฺยํ ฯเปฯ วีมํสา\-ธิปเตยฺยํ. ตสฺมิํ สมเย ผสฺโส โหติ ฯเปฯ อวิกฺเขโป โหติ ฯเปฯ. อิเม ธมฺมา กุสลา.}

\proseitem{360}{กตเม ธมฺมา กุสลา. ยสฺมิํ สมเย โลกุตฺตรํ มคฺคํ ภาเวติ ฯเปฯ โลกุตฺตรํ สติปฏฺฐานํ ภาเวติ ฯเปฯ โลกุตฺตรํ สมฺมปฺปธานํ ภาเวติ ฯเปฯ โลกุตฺตรํ อิทฺธิปาทํ ภาเวติ ฯเปฯ โลกุตฺตรํ อินฺทฺริยํ ภาเวติ ฯเปฯ โลกุตฺตรํ พลํ ภาเวติ ฯเปฯ โลกุตฺตรํ โพชฺฌงฺคํ ภาเวติ ฯเปฯ โลกุตฺตรํ สจฺจํ ภาเวติ ฯเปฯ โลกุตฺตรํ สมถํ ภาเวติ ฯเปฯ โลกุตฺตรํ ธมฺมํ ภาเวติ ฯเปฯ โลกุตฺตรํ ขนฺธํ ภาเวติ ฯเปฯ โลกุตฺตรํ อายตนํ ภาเวติ ฯเปฯ โลกุตฺตรํ ธาตุํ ภาเวติ ฯเปฯ โลกุตฺตรํ อาหารํ ภาเวติ ฯเปฯ โลกุตฺตรํ ผสฺสํ ภาเวติ ฯเปฯ โลกุตฺตรํ เวทนํ ภาเวติ ฯเปฯ โลกุตฺตรํ สญฺญํ ภาเวติ ฯเปฯ โลกุตฺตรํ เจตนํ ภาเวติ ฯเปฯ โลกุตฺตรํ จิตฺตํ ภาเวติ นิยฺยานิกํ อปจยคามิํ ทิฏฺฐิคตานํ ปหานาย ปฐมาย ภูมิยา ปตฺติยา วิวิจฺเจว \csromanfolio{89}กาเมหิ ฯเปฯ ปฐมํ ฌานํ อุปสมฺปชฺช วิหรติ ทุกฺข\-ปฏิปทํ ทนฺธาภิญฺญํ ฉนฺทา\-ธิปเตยฺยํ ฯเปฯ วีริยา\-ธิปเตยฺยํ ฯเปฯ จิตฺตา\-ธิปเตยฺยํ ฯเปฯ วีมํสา\-ธิปเตยฺยํ. ตสฺมิํ สมเย ผสฺโส โหติ ฯเปฯ อวิกฺเขโป โหติ ฯเปฯ. อิเม ธมฺมา กุสลา.}

\vspace{\tipitakaparskip}
\csromancenter{อธิปติ.}
\csromancenter{ปฐโม มคฺโค.}
\proseitem{361}{กตเม ธมฺมา กุสลา. ยสฺมิํ สมเย โลกุตฺตรํ ฌานํ ภาเวติ นิยฺยานิกํ อปจยคามิํ กามราค\-พฺยาปาทานํ ตนุภาวาย ทุติยาย ภูมิยา ปตฺติยา วิวิจฺเจว กาเมหิ ฯเปฯ ปฐมํ ฌานํ อุปสมฺปชฺช วิหรติ ทุกฺข\-ปฏิปทํ ทนฺธาภิญฺญํ. ตสฺมิํ สมเย ผสฺโส โหติ ฯเปฯ อญฺญินฺทฺริยํ โหติ ฯเปฯ อวิกฺเขโป โหติ ฯเปฯ. อิเม ธมฺมา กุสลา.}

\vspace{\tipitakaparskip}
\csromancenter{ทุติโย มคฺโค.}
\proseitem{362}{กตเม ธมฺมา กุสลา. ยสฺมิํ สมเย โลกุตฺตรํ ฌานํ ภาเวติ นิยฺยานิกํ อปจยคามิํ กามราค\-พฺยาปาทานํ อนวเสส\-ปฺปหานาย ตติยาย ภูมิยา ปตฺติยา วิวิจฺเจว กาเมหิ ฯเปฯ ปฐมํ ฌานํ อุปสมฺปชฺช วิหรติ ทุกฺข\-ปฏิปทํ ทนฺธาภิญฺญํ. ตสฺมิํ สมเย ผสฺโส โหติ ฯเปฯ อญฺญินฺทฺริยํ โหติ ฯเปฯ อวิกฺเขโป โหติ ฯเปฯ. อิเม ธมฺมา กุสลา.}

\vspace{\tipitakaparskip}
\csromancenter{ตติโย มคฺโค.}
\proseitem{363}{กตเม ธมฺมา กุสลา. ยสฺมิํ สมเย โลกุตฺตรํ ฌานํ ภาเวติ นิยฺยานิกํ อปจยคามิํ รูปราคอรูปราคมานอุทฺธจฺจ- อวิชฺชาย อนวเสส\-ปฺปหานาย จตุตฺถาย ภูมิยา ปตฺติยา วิวิจฺเจว กาเมหิ ฯเปฯ ปฐมํ ฌานํ อุปสมฺปชฺช วิหรติ ทุกฺข\-ปฏิปทํ ทนฺธาภิญฺญํ. ตสฺมิํ สมเย \csromanfolio{90}ผสฺโส โหติ ฯเปฯ อญฺญินฺทฺริยํ โหติ ฯเปฯ อวิกฺเขโป โหติ ฯเปฯ. อิเม ธมฺมา กุสลา ฯเปฯ.}

\proseitem{364}{กตมํ ตสฺมิํ สมเย อญฺญินฺทฺริยํ โหติ. ยา เตสํ ธมฺมานํ ญาตานํ ทิฏฺฐานํ ปตฺตานํ วิทิตานํ สจฺฉิกตานํ สจฺฉิกิริยาย ปญฺญา ปชานนา วิจโย ปวิจโย ธมฺมวิจโย สลฺลกฺขณา อุปลกฺขณา ปจฺจุ\-ปลกฺขณา ปณฺฑิจฺจํ โกสลฺลํ เนปุญฺญํ เวภพฺยา จินฺตา อุปปริกฺขา ภูรี เมธา ปริณายิกา วิปสฺสนา สมฺปชญฺญํ ปโตโท ปญฺญา ปญฺญินฺทฺริยํ ปญฺญาพลํ ปญฺญาสตฺถํ ปญฺญาปาสาโท ปญฺญาอาโลโก ปญฺญาโอภาโส ปญฺญาปชฺโชโต ปญฺญารตนํ อโมโห ธมฺมวิจโย สมฺมาทิฏฺฐิ ธมฺมวิจย\-สมฺโพชฺฌงฺโค มคฺคงฺคํ มคฺค\-ปริยาปนฺนํ. อิทํ ตสฺมิํ สมเย อญฺญินฺทฺริยํ โหติ ฯเปฯ อวิกฺเขโป โหติ ฯเปฯ. เย วา ปน ตสฺมิํ สมเย อญฺเญปิ อตฺถิ ปฏิจฺจ\-สมุปฺปนฺนา อรูปิโน ธมฺมา. อิเม ธมฺมา กุสลา.}

\vspace{\tipitakaparskip}
\csromancenter{จตุตฺโถ มคฺโค.}
\csromancenter{โลกุตฺตรํ จิตฺตํ.}
\csromanmatikaanchor{mtk.51}
\csromantocmarkref{subsection}{ทฺวาทส อกุสล}{mtk.51}
\bookmark[dest={mtk.51},level=2]{ทฺวาทส อกุสล}
\csromanheader{2}{\textbf{ทฺวาทส อกุสล}}
\proseitem{365}{กตเม ธมฺมา อกุสลา. ยสฺมิํ สมเย อกุสลํ จิตฺตํ อุปฺปนฺนํ โหติ โสมนสฺส\-สหคตํ ทิฏฺฐิคต\-สมฺปยุตฺตํ รูปารมฺมณํ วา สทฺทารมฺมณํ วา คนฺธารมฺมณํ วา รสารมฺมณํ วา โผฏฺฐพฺพา\-รมฺมณํ วา ธมฺมารมฺมณํ วา ยํ ยํ วา ปนารพฺภ. ตสฺมิํ สมเย ผสฺโส โหติ, เวทนา โหติ, สญฺญา โหติ, เจตนา โหติ, จิตฺตํ โหติ, วิตกฺโก โหติ, วิจาโร โหติ, ปีติ โหติ, สุขํ โหติ, จิตฺตสฺเส\-กคฺคตา โหติ, วีริยินฺทฺริยํ โหติ, สมาธินฺทฺริยํ โหติ, มนินฺทฺริยํ โหติ, โสมนสฺสิ\-นฺทฺริยํ โหติ, ชีวิตินฺทฺริยํ โหติ, มิจฺฉาทิฏฺฐิ โหติ, มิจฺฉาสงฺกปฺโป โหติ, มิจฺฉาวายาโม โหติ, มิจฺฉาสมาธิ โหติ, วีริยพลํ โหติ, สมาธิพลํ โหติ, อหิริกพลํ โหติ, อโนตฺตปฺปพลํ โหติ, โลโภ โหติ, โมโห โหติ, อภิชฺฌา โหติ, มิจฺฉาทิฏฺฐิ โหติ, อหิริกํ \csromanfolio{91}โหติ, อโนตฺตปฺปํ โหติ, สมโถ โหติ, ปคฺคาโห โหติ, อวิกฺเขโป โหติ. เย วา ปน ตสฺมิํ สมเย อญฺเญปิ อตฺถิ ปฏิจฺจ\-สมุปฺปนฺนา อรูปิโน ธมฺมา. อิเม ธมฺมา อกุสลา.}

\proseitem{366}{กตโม ตสฺมิํ สมเย ผสฺโส โหติ. โย ตสฺมิํ สมเย ผสฺโส ผุสนา สมฺผุสนา สมฺผุสิตตฺตํ. อยํ ตสฺมิํ สมเย ผสฺโส โหติ.}

\proseitem{367}{กตมา ตสฺมิํ สมเย เวทนา โหติ. ยํ ตสฺมิํ สมเย ตชฺชา\-มโนวิญฺญาณธาตุ\-สมฺผสฺสชํ เจตสิกํ สาตํ เจตสิกํ สุขํ เจโต\-สมฺผสฺสชํ สาตํ สุขํ เวทยิตํ เจโต\-สมฺผสฺสชา สาตา สุขา เวทนา. อยํ ตสฺมิํ สมเย เวทนา โหติ.}

\proseitem{368}{กตมา ตสฺมิํ สมเย สญฺญา โหติ. ยา ตสฺมิํ สมเย ตชฺชา\-มโนวิญฺญาณธาตุ\-สมฺผสฺสชา สญฺญา สญฺชานนา สญฺชานิตตฺตํ. อยํ ตสฺมิํ สมเย สญฺญา โหติ.}

\proseitem{369}{กตมา ตสฺมิํ สมเย เจตนา โหติ. ยา ตสฺมิํ สมเย ตชฺชา\-มโนวิญฺญาณธาตุ\-สมฺผสฺสชา เจตนา สญฺเจตนา เจตยิตตฺตํ. อยํ ตสฺมิํ สมเย เจตนา โหติ.}

\proseitem{370}{กตมํ ตสฺมิํ สมเย จิตฺตํ โหติ. ยํ ตสฺมิํ สมเย จิตฺตํ มโน มานสํ หทยํ ปณฺฑรํ มโน มนายตนํ มนินฺทฺริยํ วิญฺญาณํ วิญฺญาณ\-กฺขนฺโธ ตชฺชา\-มโนวิญฺญาณ\-ธาตุ. อิทํ ตสฺมิํ สมเย จิตฺตํ โหติ.}

\proseitem{371}{กตโม ตสฺมิํ สมเย วิตกฺโก โหติ. โย ตสฺมิํ สมเย ตกฺโก วิตกฺโก สงฺกปฺโป อปฺปนา พฺยปฺปนา เจตโส อภินิโรปนา มิจฺฉาสงฺกปฺโป. อยํ ตสฺมิํ สมเย วิตกฺโก โหติ.}

\proseitem{372}{กตโม ตสฺมิํ สมเย วิจาโร โหติ. โย ตสฺมิํ สมเย จาโร วิจาโร อนุวิจาโร อุปวิจาโร จิตฺตสฺส อนุสนฺธานตา อนุเปกฺขนตา. อยํ ตสฺมิํ สมเย วิจาโร โหติ.}

\proseitem{373}{\csromanfolio{92}กตมา ตสฺมิํ สมเย ปีติ โหติ. ยา ตสฺมิํ สมเย ปีติ ปาโมชฺชํ อาโมทนา ปโมทนา หาโส ปหาโส วิตฺติ โอทคฺยํ อตฺตมนตา จิตฺตสฺส. อยํ ตสฺมิํ สมเย ปีติ โหติ.}

\proseitem{374}{กตมํ ตสฺมิํ สมเย สุขํ โหติ. ยํ ตสฺมิํ สมเย เจตสิกํ สาตํ เจตสิกํ สุขํ เจโต\-สมฺผสฺสชํ สาตํ สุขํ เวทยิตํ เจโต\-สมฺผสฺสชา สาตา สุขา เวทนา. อิทํ ตสฺมิํ สมเย สุขํ โหติ.}

\proseitem{375}{กตมา ตสฺมิํ สมเย จิตฺตสฺเส\-กคฺคตา โหติ. ยา ตสฺมิํ สมเย จิตฺตสฺส ฐิติ สณฺฐิติ อวฏฺฐิติ อวิสาหาโร อวิกฺเขโป อวิสาหฏ\-มานสตา สมโถ สมาธินฺทฺริยํ สมาธิพลํ มิจฺฉาสมาธิ. อยํ ตสฺมิํ สมเย จิตฺตสฺเส\-กคฺคตา โหติ.}

\proseitem{376}{กตมํ ตสฺมิํ สมเย วีริยินฺทฺริยํ โหติ. โย ตสฺมิํ สมเย เจตสิโก วีริยารมฺโภ นิกฺกโม ปรกฺกโม อุยฺยาโม วายาโม อุสฺสาโห อุสฺโสฬฺหี ถาโม ธิติ อสิถิล\-ปรกฺกมตา อนิกฺขิตฺต\-ฉนฺทตา อนิกฺขิตฺตธุรตา ธุรสมฺปคฺคาโห วีริยํ วีริยินฺทฺริยํ วีริยพลํ มิจฺฉาวายาโม. อิทํ ตสฺมิํ สมเย วีริยินฺทฺริยํ โหติ.}

\proseitem{377}{กตมํ ตสฺมิํ สมเย สมาธินฺทฺริยํ โหติ. ยา ตสฺมิํ สมเย จิตฺตสฺส ฐิติ สณฺฐิติ อวฏฺฐิติ อวิสาหาโร อวิกฺเขโป อวิสาหฏ\-มานสตา สมโถ สมาธินฺทฺริยํ สมาธิพลํ มิจฺฉาสมาธิ อิทํ ตสฺมิํ สมเย สมาธินฺทฺริยํ โหติ.}

\proseitem{378}{กตมํ ตสฺมิํ สมเย มนินฺทฺริยํ โหติ. ยํ ตสฺมิํ สมเย จิตฺตํ มโน มานสํ หทยํ ปณฺฑรํ มโน มนายตนํ มนินฺทฺริยํ วิญฺญาณํ วิญฺญาณ\-กฺขนฺโธ ตชฺชา\-มโนวิญฺญาณ\-ธาตุ. อิทํ ตสฺมิํ สมเย มนินฺทฺริยํ โหติ.}

\proseitem{379}{กตมํ ตสฺมิํ สมเย โสมนสฺสิ\-นฺทฺริยํ โหติ. ยํ ตสฺมิํ สมเย เจตสิกํ สาตํ เจตสิกํ สุขํ เจโต\-สมฺผสฺสชํ สาตํ สุขํ เวทยิตํ เจโต\-สมฺผสฺสชา สาตา สุขา เวทนา. อิทํ ตสฺมิํ สมเย โสมนสฺสิ\-นฺทฺริยํ โหติ.}

\chapterheadpageread{1. จิตฺตุปฺปาท\-กณฺฑ}
\proseitem{380}{\csromanfolio{93}กตมํ ตสฺมิํ สมเย ชีวิตินฺทฺริยํ โหติ. โย เตสํ อรูปีนํ ธมฺมานํ อายุ ฐิติ ยปนา ยาปนา อิริยนา วตฺตนา ปาลนา ชีวิตํ ชีวิตินฺทฺริยํ. อิทํ ตสฺมิํ สมเย ชีวิตินฺทฺริยํ โหติ.}

\proseitem{381}{กตมา ตสฺมิํ สมเย มิจฺฉาทิฏฺฐิ โหติ. ยา ตสฺมิํ สมเย ทิฏฺฐิ ทิฏฺฐิคตํ ทิฏฺฐิคหนํ ทิฏฺฐิกนฺตาโร ทิฏฺฐิ\-วิสูกายิกํ ทิฏฺฐิ\-วิปฺผนฺทิตํ ทิฏฺฐิ\-สญฺโญชนํ คาโห ปติฏฺฐาโห\csromansharedfootnote{adfc48d9bcacdc40}{ปฏิคฺคาโห (สี, สฺยา)} อภินิเวโส ปรามาโส กุมฺมคฺโค มิจฺฉาปโถ มิจฺฉตฺตํ ติตฺถายตนํ วิปริยาส\-คฺคาโห\csromansharedfootnote{4de7dc47ba037bea}{วิปริเยสคาโห (ก)}. อยํ ตสฺมิํ สมเย มิจฺฉาทิฏฺฐิ โหติ.}

\proseitem{382}{กตโม ตสฺมิํ สมเย มิจฺฉาสงฺกปฺโป โหติ. โย ตสฺมิํ สมเย ตกฺโก วิตฺตกฺโก สงฺกปฺโป อปฺปนา พฺยปฺปนา เจตโส อภินิโรปนา มิจฺฉาสงฺกปฺโป. อยํ ตสฺมิํ สมเย มิจฺฉาสงฺกปฺโป โหติ.}

\proseitem{383}{กตโม ตสฺมิํ สมเย มิจฺฉาวายาโม โหติ. โย ตสฺมิํ สมเย เจตสิโก วีริยารมฺโภ นิกฺกโม ปรกฺกโม อุยฺยาโม วายาโม อุสฺสาโห อุสฺโสฬฺหี ถาโม ธิติ อสิถิล\-ปรกฺกมตา อนิกฺขิตฺต\-ฉนฺทตา อนิกฺขิตฺตธุรตา ธุรสมฺปคฺคาโห วีริยํ วีริยินฺทฺริยํ วีริยพลํ มิจฺฉาวายาโม. อยํ ตสฺมิํ สมเย มิจฺฉาวายาโม โหติ.}

\proseitem{384}{กตโม ตสฺมิํ สมเย มิจฺฉาสมาธิ โหติ. ยา ตสฺมิํ สมเย จิตฺตสฺส ฐิติ สณฺฐิติ อวฏฺฐิติ อวิสาหาโร อวิกฺเขโป อวิสาหฏ\-มานสตา สมโถ สมาธินฺทฺริยํ สมาธิพลํ มิจฺฉาสมาธิ. อยํ ตสฺมิํ สมเย มิจฺฉาสมาธิ โหติ.}

\proseitem{385}{กตมํ ตสฺมิํ สมเย วีริยพลํ โหติ. โย ตสฺมิํ สมเย เจตสิโก วีริยารมฺโภ นิกฺกโม ปรกฺกโม อุยฺยาโม วายาโม อุสฺสาโห อุสฺโสฬฺหี ถาโม ธิติ อสิถิล\-ปรกฺกมตา อนิกฺขิตฺต\-ฉนฺทตา อนิกฺขิตฺตธุรตา ธุรสมฺปคฺคาโห วีริยํ วีริยินฺทฺริยํ วีริยพลํ มิจฺฉาวายาโม. อิทํ ตสฺมิํ สมเย วีริยพลํ โหติ.}

\proseitem{386}{\csromanfolio{94}กตมํ ตสฺมิํ สมเย สมาธิพลํ โหติ. ยา ตสฺมิํ สมเย จิตฺตสฺส ฐิติ สณฺฐิติ อวฏฺฐิติ อวิสาหาโร อวิกฺเขโป อวิสาหฏ\-มานสตา สมโถ สมาธินฺทฺริยํ สมาธิพลํ มิจฺฉาสมาธิ. อิทํ ตสฺมิํ สมเย สมาธิพลํ โหติ.}

\proseitem{387}{กตมํ ตสฺมิํ สมเย อหิริกพลํ โหติ. ยํ ตสฺมิํ สมเย น หิรียติ หิริยิตพฺเพน น หิรียติ ปาปกานํ อกุสลานํ ธมฺมานํ สมาปตฺติยา. อิทํ ตสฺมิํ สมเย อหิริกพลํ โหติ.}

\proseitem{388}{กตมํ ตสฺมิํ สมเย อโนตฺตปฺปพลํ โหติ. ยํ ตสฺมิํ สมเย น โอตฺตปฺปติ โอตฺตปฺปิตพฺเพน น โอตฺตปฺปติ ปาปกานํ อกุสลานํ ธมฺมานํ สมาปตฺติยา. อิทํ ตสฺมิํ สมเย อโนตฺตปฺปพลํ โหติ.}

\proseitem{389}{กตโม ตสฺมิํ สมเย โลโภ โหติ. โย ตสฺมิํ สมเย โลโภ ลุพฺภนา ลุพฺภิตตฺตํ สาราโค สารชฺชนา สารชฺชิตตฺตํ อภิชฺฌา โลโภ อกุสลมูลํ. อยํ ตสฺมิํ สมเย โลโภ โหติ.}

\proseitem{390}{กตโม ตสฺมิํ สมเย โมโห โหติ. ยํ ตสฺมิํ สมเย อญฺญาณํ อทสฺสนํ อนภิสมโย อนนุโพโธ อสมฺโพโธ อปฺปฏิเวโธ อสํคาหนา อปริโยคาหนา อสมเปกฺขนา อปจฺจเวกฺขนา อปจฺจกฺขกมฺมํ ทุมฺเมชฺฌํ พาลฺยํ อสมฺปชญฺญํ โมโห ปโมโห สมฺโมโห อวิชฺชา อวิชฺโชโฆ อวิชฺชาโยโค อวิชฺชานุสโย อวิชฺชา\-ปริยุฏฺฐานํ อวิชฺชาลงฺคี โมโห อกุสลมูลํ. อยํ ตสฺมิํ สมเย โมโห โหติ.}

\proseitem{391}{กตมา ตสฺมิํ สมเย อภิชฺฌา โหติ. โย ตสฺมิํ สมเย โลโภ ลุพฺภนา ลุพฺภิตตฺตํ สาราโค สารชฺชนา สารชฺชิตตฺตํ อภิชฺฌา โลโภ อกุสลมูลํ. อยํ ตสฺมิํ สมเย อภิชฺฌา โหติ.}

\proseitem{392}{กตมา ตสฺมิํ สมเย มิจฺฉาทิฏฺฐิ โหติ. ยา ตสฺมิํ สมเย ทิฏฺฐิ ทิฏฺฐิคตํ ทิฏฺฐิคหนํ ทิฏฺฐิกนฺตาโร ทิฏฺฐิ\-วิสูกายิกํ ทิฏฺฐิ\-วิปฺผนฺทิตํ ทิฏฺฐิ\-สญฺโญชนํ คาโห ปติฏฺฐาโห อภินิเวโส ปรามาโส กุมฺมคฺโค \csromanfolio{95}มิจฺฉาปโถ มิจฺฉตฺตํ ติตฺถายตนํ วิปริยาส\-คฺคาโห. อยํ ตสฺมิํ สมเย มิจฺฉาทิฏฺฐิ โหติ.}

\proseitem{393}{กตมํ ตสฺมิํ สมเย อหิริกํ โหติ. ยํ ตสฺมิํ สมเย น หิรียติ หิริยิตพฺเพน น หิรียติ ปาปกานํ อกุสลานํ ธมฺมานํ สมาปตฺติยา. อิทํ ตสฺมิํ สมเย อหิริกํ โหติ.}

\proseitem{394}{กตมํ ตสฺมิํ สมเย อโนตฺตปฺปํ โหติ. ยํ ตสฺมิํ สมเย น โอตฺตปฺปติ โอตฺตปฺปิตพฺเพน น โอตฺตปฺปติ ปาปกานํ อกุสลานํ ธมฺมานํ สมาปตฺติยา. อิทํ ตสฺมิํ สมเย อโนตฺตปฺปํ โหติ.}

\proseitem{395}{กตโม ตสฺมิํ สมเย สมโถ โหติ. ยา ตสฺมิํ สมเย จิตฺตสฺส ฐิติ สณฺฐิติ อวฏฺฐิติ อวิสาหาโร อวิกฺเขโป อวิสาหฏ\-มานสตา สมโถ สมาธินฺทฺริยํ สมาธิพลํ มิจฺฉาสมาธิ. อยํ ตสฺมิํ สมเย สมโถ โหติ.}

\proseitem{396}{กตโม ตสฺมิํ สมเย ปคฺคาโห โหติ. โย ตสฺมิํ สมเย เจตสิโก วีริยารมฺโภ นิกฺกโม ปรกฺกโม อุยฺยาโม วายาโม อุสฺสาโห อุสฺโสฬฺหี ถาโม ธิติ อสิถิล\-ปรกฺกมตา อนิกฺขิตฺต\-ฉนฺทตา อนิกฺขิตฺตธุรตา ธุรสมฺปคฺคาโห วีริยํ วีริยินฺทฺริยํ วีริยพลํ มิจฺฉาวายาโม. อยํ ตสฺมิํ สมเย ปคฺคาโห โหติ.}

\proseitem{397}{กตโม ตสฺมิํ สมเย อวิกฺเขโป โหติ. ยา ตสฺมิํ สมเย จิตฺตสฺส ฐิติ สณฺฐิติ อวฏฺฐิติ อวิสาหาโร อวิกฺเขโป อวิสาหฏ\-มานสตา สมโถ สมาธินฺทฺริยํ สมาธิพลํ มิจฺฉาสมาธิ. อยํ ตสฺมิํ สมเย อวิกฺเขโป โหติ. เย วา ปน ตสฺมิํ สมเย อญฺเญปิ อตฺถิ ปฏิจฺจ\-สมุปฺปนฺนา อรูปิโน ธมฺมา. อิเม ธมฺมา อกุสลา.}

\prose{ตสฺมิํ โข ปน สมเย จตฺตาโร ขนฺธา โหนฺติ, ทฺวายตนานิ โหนฺติ, เทฺว ธาตุโย โหนฺติ, ตโย อาหารา โหนฺติ, ปญฺจินฺทฺริยานิ โหนฺติ, ปญฺจงฺคิกํ ฌานํ โหติ, จตุรงฺคิโก มคฺโค โหติ, จตฺตาริ พลานิ โหนฺติ, เทฺว เหตู โหนฺติ, เอโก ผสฺโส โหติ ฯเปฯ เอกํ ธมฺมายตนํ โหติ, เอกา ธมฺมาธาตุ โหติ. เย วา ปน ตสฺมิํ สมเย อญฺเญปิ อตฺถิ ปฏิจฺจ\-สมุปฺปนฺนา อรูปิโน ธมฺมา. อิเม ธมฺมา อกุสลา ฯเปฯ.}

\proseitem{398}{\csromanfolio{96}กตโม ตสฺมิํ สมเย สงฺขาร\-กฺขนฺโธ โหติ. ผสฺโส เจตนา วิตกฺโก วิจาโร ปีติ จิตฺตสฺเส\-กคฺคตา วีริยินฺทฺริยํ สมาธินฺทฺริยํ ชีวิตินฺทฺริยํ มิจฺฉาทิฏฺฐิ มิจฺฉาสงฺกปฺโป มิจฺฉาวายาโม มิจฺฉาสมาธิ วีริยพลํ สมาธิพลํ อหิริกพลํ อโนตฺตปฺปพลํ โลโภ โมโห อภิชฺฌา มิจฺฉาทิฏฺฐิ อหิริกํ อโนตฺตปฺปํ สมโถ ปคฺคาโห อวิกฺเขโป. เย วา ปน ตสฺมิํ สมเย อญฺเญปิ อตฺถิ ปฏิจฺจ\-สมุปฺปนฺนา อรูปิโน ธมฺมา ฐเปตฺวา เวทนา\-กฺขนฺธํ ฐเปตฺวา สญฺญากฺขนฺธํ ฐเปตฺวา วิญฺญาณ\-กฺขนฺธํ. อยํ ตสฺมิํ สมเย สงฺขาร\-กฺขนฺโธ โหติ ฯเปฯ. อิเม ธมฺมา อกุสลา.}

\proseitem{399}{กตเม ธมฺมา อกุสลา. ยสฺมิํ สมเย อกุสลํ จิตฺตํ อุปฺปนฺนํ โหติ โสมนสฺส\-สหคตํ ทิฏฺฐิคต\-สมฺปยุตฺตํ สสงฺขาเรน รูปารมฺมณํ วา ฯเปฯ ธมฺมารมฺมณํ วา ยํ ยํ วา ปนารพฺภ. ตสฺมิํ สมเย ผสฺโส โหติ ฯเปฯ อวิกฺเขโป โหติ ฯเปฯ. อิเม ธมฺมา อกุสลา.}

\proseitem{400}{กตเม ธมฺมา อกุสลา. ยสฺมิํ สมเย อกุสลํ จิตฺตํ อุปฺปนฺนํ โหติ โสมนสฺส\-สหคตํ ทิฏฺฐิคต\-วิปฺปยุตฺตํ รูปารมฺมณํ วา สทฺทารมฺมณํ วา คนฺธารมฺมณํ วา รสารมฺมณํ วา โผฏฺฐพฺพา\-รมฺมณํ วา ธมฺมารมฺมณํ วา ยํ ยํ วา ปนารพฺภ. ตสฺมิํ สมเย ผสฺโส โหติ, เวทนา โหติ, สญฺญา โหติ, เจตนา โหติ, จิตฺตํ โหติ, วิตกฺโก โหติ, วิจาโร โหติ, ปีติ โหติ, สุขํ โหติ, จิตฺตสฺเส\-กคฺคตา โหติ, วีริยินฺทฺริยํ โหติ, สมาธินฺทฺริยํ โหติ, มนินฺทฺริยํ โหติ, โสมนสฺสิ\-นฺทฺริยํ โหติ, ชีวิตินฺทฺริยํ โหติ, มิจฺฉาสงฺกปฺโป โหติ, มิจฺฉาวายาโม โหติ, มิจฺฉาสมาธิ โหติ, วีริยพลํ โหติ, สมาธิพลํ โหติ, อหิริกพลํ โหติ, อโนตฺตปฺปพลํ โหติ, โลโภ โหติ, โมโห โหติ, อภิชฺฌา โหติ, อหิริกํ โหติ, อโนตฺตปฺปํ โหติ, สมโถ โหติ, ปคฺคาโห โหติ อวิกฺเขโป โหติ. เย วา ปน ตสฺมิํ สมเย อญฺเญปิ อตฺถิ ปฏิจฺจ\-สมุปฺปนฺนา อรูปิโน ธมฺมา. อิเม ธมฺมา อกุสลา ฯเปฯ.}

\prose{ตสฺมิํ โข ปน สมเย จตฺตาโร ขนฺธา โหนฺติ, ทฺวายตนานิ โหนฺติ, เทฺว ธาตุโย โหนฺติ, ตโย อาหารา โหนฺติ, ปญฺจินฺทฺริยานิ โหนฺติ, ปญฺจงฺคิกํ ฌานํ โหติ, ติวงฺคิโก มคฺโค โหติ, จตฺตาริ พลานิ โหนฺติ, เทฺว เหตู โหนฺติ, เอโก ผสฺโส โหติ ฯเปฯ เอกํ \csromanfolio{97}ธมฺมายตนํ โหติ. เอกา ธมฺมธาตุ โหติ. เย วา ปน ตสฺมิํ สมเย อญฺเญปิ อตฺถิ ปฏิจฺจ\-สมุปฺปนฺนา อรูปิโน ธมฺมา. อิเม ธมฺมา อกุสลา ฯเปฯ.}

\proseitem{401}{กตโม ตสฺมิํ สมเย สงฺขาร\-กฺขนฺโธ โหติ. ผสฺโส เจตนา วิตกฺโก วิจาโร ปีติ จิตฺตสฺเส\-กคฺคตา วีริยินฺทฺริยํ สมาธินฺทฺริยํ ชีวิตินฺทฺริยํ มิจฺฉาสงฺกปฺโป มิจฺฉาวายาโม มิจฺฉาสมาธิ วีริยพลํ สมาธิพลํ อหิริกพลํ อโนตฺตปฺปพลํ โลโภ โมโห อภิชฺฌา อหิริกํ อโนตฺตปฺปํ สมโถ ปคฺคาโห อวิกฺเขโป. เย วา ปน ตสฺมิํ สมเย อญฺเญปิ อตฺถิ ปฏิจฺจ\-สมุปฺปนฺนา อรูปิโน ธมฺมา ฐเปตฺวา เวทนา\-กฺขนฺธํ ฐเปตฺวา สญฺญากฺขนฺธํ ฐเปตฺวา วิญฺญาณ\-กฺขนฺธํ. อยํ ตสฺมิํ สมเย สงฺขาร\-กฺขนฺโธ โหติ ฯเปฯ. อิเม ธมฺมา อกุสลา.}

\proseitem{402}{กตเม ธมฺมา อกุสลา. ยสฺมิํ สมเย อกุสลํ จิตฺตํ อุปฺปนฺนํ โหติ โสมนสฺส\-สหคตํ ทิฏฺฐิคต\-วิปฺปยุตฺตํ สสงฺขาเรน รูปารมฺมณํ วา ฯเปฯ ธมฺมารมฺมณํ วา ยํ ยํ ปนารพฺภ. ตสฺมิํ สมเย ผสฺโส โหติ ฯเปฯ อวิกฺเขโป โหติ ฯเปฯ. อิเม ธมฺมา อกุสลา.}

\proseitem{403}{กตเม ธมฺมา อกุสลา. ยสฺมิํ สมเย อกุสลํ จิตฺตํ อุปฺปนฺนํ โหติ อุเปกฺขา\-สหคตํ ทิฏฺฐิคต\-สมฺปยุตฺตํ รูปารมฺมณํ วา สทฺทารมฺมณํ วา คนฺธารมฺมณํ วา รสารมฺมณํ วา โผฏฺฐพฺพา\-รมฺมณํ วา ธมฺมารมฺมณํ วา ยํ ยํ วา ปนารพฺภ. ตสฺมิํ สมเย ผสฺโส โหติ, เวทนา โหติ, สญฺญา โหติ, เจตนา โหติ, จิตฺตํ โหติ, วิตกฺโก โหติ, วิจาโร โหติ, อุเปกฺขา โหติ, จิตฺตสฺเส\-กคฺคตา โหติ, วีริยินฺทฺริยํ โหติ, สมาธินฺทฺริยํ โหติ, มนินฺทฺริยํ โหติ, อุเปกฺขินฺทฺริยํ โหติ, ชีวิตินฺทฺริยํ โหติ, มิจฺฉาทิฏฺฐิ โหติ, มิจฺฉาสงฺกปฺโป โหติ, มิจฺฉาวายาโม โหติ, มิจฺฉาสมาธิ โหติ, วีริยพลํ โหติ, สมาธิพลํ โหติ, อหิริกพลํ โหติ, อโนตฺตปฺปพลํ โหติ, โลโภ โหติ, โมโห โหติ, อภิชฺฌา โหติ, มิจฺฉาทิฏฺฐิ โหติ, อหิริกํ โหติ, อโนตฺตปฺปํ โหติ, สมโถ โหติ, ปคฺคาโห โหติ, อวิกฺเขโป โหติ. เย วา ปน ตสฺมิํ สมเย อญฺเญปิ อตฺถิ ปฏิจฺจ\-สมุปฺปนฺนา อรูปิโน ธมฺมา. อิเม ธมฺมา อกุสลา.}

\proseitem{404}{\csromanfolio{98}กตโม ตสฺมิํ สมเย ผสฺโส โหติ. โย ตสฺมิํ สมเย ผสฺโส ผุสนา สมฺผุสนา สมฺผุสิตตฺตํ. อยํ ตสฺมิํ สมเย ผสฺโส โหติ.}

\proseitem{405}{กตมา ตสฺมิํ สมเย เวทนา โหติ. ยํ ตสฺมิํ สมเย ตชฺชา\-มโนวิญฺญาณธาตุ\-สมฺผสฺสชํ เจตสิกํ เนว สาตํ นาสาตํ เจโต\-สมฺผสฺสชํ อทุกฺขมสุขํ เวทยิตํ เจโต\-สมฺผสฺสชา อทุกฺขมสุขา เวทนา. อยํ ตสฺมิํ สมเย เวทนา โหติ ฯเปฯ.}

\proseitem{406}{กตมา ตสฺมิํ สมเย อุเปกฺขา โหติ. ยํ ตสฺมิํ สมเย เจตสิกํ เนว สาตํ นาสาตํ เจโต\-สมฺผสฺสชํ อทุกฺขมสุขํ เวทยิตํ เจโต\-สมฺผสฺสชา อทุกฺขมสุขา เวทนา. อยํ ตสฺมิํ สมเย อุเปกฺขา โหติ ฯเปฯ.}

\proseitem{407}{กตมํ ตสฺมิํ สมเย อุเปกฺขินฺทฺริยํ โหติ. ยํ ตสฺมิํ สมเย เจตสิกํ เนว สาตํ นาสาตํ เจโต\-สมฺผสฺสชํ อทุกฺขมสุขํ เวทยิตํ เจโต\-สมฺผสฺสชา อทุกฺขมสุขา เวทนา. อิทํ ตสฺมิํ สมเย อุเปกฺขินฺทฺริยํ โหติ ฯเปฯ. เย วา ปน ตสฺมิํ สมเย อญฺเญปิ อตฺถิ ปฏิจฺจ\-สมุปฺปนฺนา อรูปิโน ธมฺมา. อิเม ธมฺมา อกุสลา.}

\prose{ตสฺมิํ โข ปน สมเย จตฺตาโร ขนฺธา โหนฺติ, ทฺวายตนานิ โหนฺติ, เทฺว ธาตุโย โหนฺติ, ตโย อาหารา โหนฺติ, ปญฺจินฺทฺริยานิ โหนฺติ, จตุรงฺคิกํ ฌานํ โหติ, จตุรงฺคิโก มคฺโค โหติ, จตฺตาริ พลานิ โหนฺติ, เทฺว เหตู โหนฺติ, เอโก ผสฺโส โหติ ฯเปฯ เอกํ ธมฺมายตนํ โหติ, เอกา ธมฺมธาตุ โหติ. เย วา ปน ตสฺมิํ สมเย อญฺเญปิ อตฺถิ ปฏิจฺจ\-สมุปฺปนฺนา อรูปิโน ธมฺมา. อิเม ธมฺมา อกุสลา ฯเปฯ.}

\proseitem{408}{กตโม ตสฺมิํ สมเย สงฺขาร\-กฺขนฺโธ โหติ. ผสฺโส เจตนา วิตกฺโก วิจาโร จิตฺตสฺเส\-กคฺคตา วีริยินฺทฺริยํ สมาธินฺทฺริยํ ชีวิตินฺทฺริยํ มิจฺฉาทิฏฺฐิ มิจฺฉาสงฺกปฺโป มิจฺฉาวายาโม มิจฺฉาสมาธิ วีริยพลํ สมาธิพลํ อหิริกพลํ อโนตฺตปฺปพลํ โลโภ โมโห อภิชฺฌา มิจฺฉาทิฏฺฐิ อหิริกํ อโนตฺตปฺปํ สมโถ ปคฺคาโห อวิกฺเขโป. เย วา ปน ตสฺมิํ สมเย อญฺเญปิ อตฺถิ ปฏิจฺจ\-สมุปฺปนฺนา อรูปิโน ธมฺมา ฐเปตฺวา เวทนา\-กฺขนฺธํ ฐเปตฺวา สญฺญากฺขนฺธํ ฐเปตฺวา วิญฺญาณ\-กฺขนฺธํ. อยํ ตสฺมิํ สมเย สงฺขาร\-กฺขนฺโธ โหติ ฯเปฯ. อิเม ธมฺมา อกุสลา.}

\chapterheadpageread{1. จิตฺตุปฺปาท\-กณฺฑ}
\proseitem{409}{\csromanfolio{99}กตเม ธมฺมา อกุสลา. ยสฺมิํ สมเย อกุสลํ จิตฺตํ อุปฺปนฺนํ โหติ อุเปกฺขา\-สหคตํ ทิฏฺฐิคต\-สมฺปยุตฺตํ สสงฺขาเรน รูปารมฺมณํ วา ฯเปฯ ธมฺมารมฺมณํ วา ยํ ยํ วา ปนารพฺภ. ตสฺมิํ สมเย ผสฺโส โหติ ฯเปฯ อวิกฺเขโป โหติ ฯเปฯ. อิเม ธมฺมา อกุสลา.}

\proseitem{410}{กตเม ธมฺมา อกุสลา. ยสฺมิํ สมเย อกุสลํ จิตฺตํ อุปฺปนฺนํ โหติ อุเปกฺขา\-สหคตํ ทิฏฺฐิคต\-วิปฺปยุตฺตํ รูปารมฺมณํ วา สทฺทารมฺมณํ วา คนฺธารมฺมณํ วา รสารมฺมณํ วา โผฏฺฐพฺพา\-รมฺมณํ วา ธมฺมารมฺมณํ วา ยํ ยํ วา ปนารพฺภ. ตสฺมิํ สมเย ผสฺโส โหติ, เวทนา โหติ, สญฺญา โหติ, เจตนา โหติ, จิตฺตํ โหติ, วิตกฺโก โหติ, วิจาโร โหติ, อุเปกฺขา โหติ, จิตฺตสฺเส\-กคฺคตา โหติ, วีริยินฺทฺริยํ โหติ, สมาธินฺทฺริยํ โหติ, มนินฺทฺริยํ โหติ, อุเปกฺขินฺทฺริยํ ชีวิตินฺทฺริยํ โหติ, มิจฺฉาสงฺกปฺโป โหติ, มิจฺฉาวายาโม โหติ, มิจฺฉาสมาธิ โหติ, วีริยพลํ โหติ, สมาธิพลํ โหติ, อหิริกพลํ โหติ, อโนตฺตปฺปพลํ โหติ, โลโภ โหติ, โมโห โหติ, อภิชฺฌา โหติ, อหิริกํ โหติ, อโนตฺตปฺปํ โหติ, สมโถ โหติ, ปคฺคาโห โหติ, อวิกฺเขโป โหติ. เย วา ปน ตสฺมิํ สมเย อญฺเญปิ อตฺถิ ปฏิจฺจ\-สมุปฺปนฺนา อรูปิโน ธมฺมา. อิเม ธมฺมา อกุสลา ฯเปฯ.}

\prose{ตสฺมิํ โข ปน สมเย จตฺตาโร ขนฺธา โหนฺติ, ทฺวายตนานิ โหนฺติ, เทฺว ธาตุโย โหนฺติ, ตโย อาหารา โหนฺติ, ปญฺจินฺทฺริยานิ โหนฺติ, จตุรงฺคิกํ ฌานํ โหติ, ติวงฺคิโก มคฺโค โหติ, จตฺตาริ พลานิ โหนฺติ, เทฺว เหตู โหนฺติ, เอโก ผสฺโส โหติ ฯเปฯ เอกํ ธมฺมายตนํ โหติ, เอกา ธมฺมาธาตุ โหติ. เย วา ปน ตสฺมิํ สมเย อญฺเญปิ อตฺถิ ปฏิจฺจ\-สมุปฺปนฺนา อรูปิโน ธมฺมา. อิเม ธมฺมา อกุสลา ฯเปฯ.}

\proseitem{411}{กตโม ตสฺมิํ สมเย สงฺขาร\-กฺขนฺโธ โหติ. ผสฺโส เจตนา วิตกฺโก วิจาโร จิตฺตสฺเส\-กคฺคตา วีริยินฺทฺริยํ สมาธินฺทฺริยํ ชีวิตินฺทฺริยํ มิจฺฉาสงฺกปฺโป มิจฺฉาวายาโม มิจฺฉาสมาธิ วีริยพลํ สมาธิพลํ อหิริกพลํ อโนตฺตปฺปพลํ โลโภ โมโห อภิชฺฌา อหิริกํ อโนตฺตปฺปํ สมโถ ปคฺคาโห อวิกฺเขโป. เย วา ปน ตสฺมิํ สมเย อญฺเญปิ อตฺถิ ปฏิจฺจ\-สมุปฺปนฺนา อรูปิโน ธมฺมา ฐเปตฺวา \csromanfolio{100}เวทนา\-กฺขนฺธํ ฐเปตฺวา สญฺญากฺขนฺธํ ฐเปตฺวา วิญฺญาณ\-กฺขนฺธํ. อยํ ตสฺมิํ สมเย สงฺขาร\-กฺขนฺโธ โหติ ฯเปฯ. อิเม ธมฺมา อกุสลา.}

\proseitem{412}{กตเม ธมฺมา อกุสลา. ยสฺมิํ สมเย อกุสลํ จิตฺตํ อุปฺปนฺนํ โหติ อุเปกฺขา\-สหคตํ ทิฏฺฐิคต\-วิปฺปยุตฺตํ สสงฺขาเรน รูปารมฺมณํ วา ฯเปฯ ธมฺมารมฺมณํ วา ยํ ยํ วา ปนารพฺภ. ตสฺมิํ สมเย ผสฺโส โหติ ฯเปฯ อวิกฺเขโป โหติ ฯเปฯ. อิเม ธมฺมา อกุสลา.}

\proseitem{413}{กตเม ธมฺมา อกุสลา. ยสฺมิํ สมเย อกุสลํ จิตฺตํ อุปฺปนฺนํ โหติ โทมนสฺส\-สหคตํ ปฏิฆ\-สมฺปยุตฺตํ รูปารมฺมณํ วา สทฺทารมฺมณํ วา คนฺธารมฺมณํ วา รสารมฺมณํ วา โผฏฺฐพฺพา\-รมฺมณํ วา ธมฺมารมฺมณํ วา ยํ ยํ วา ปนารพฺภ. ตสฺมิํ สมเย ผสฺโส โหติ, เวทนา โหติ, สญฺญา โหติ, เจตนา โหติ, จิตฺตํ โหติ, วิตกฺโก โหติ, วิจาโร โหติ, ทุกฺขํ โหติ, จิตฺตสฺเส\-กคฺคตา โหติ, วีริยินฺทฺริยํ โหติ, สมาธินฺทฺริยํ โหติ, มนินฺทฺริยํ โหติ, โทมนสฺสิ\-นฺทฺริยํ โหติ, ชีวิตินฺทฺริยํ โหติ, มิจฺฉาสงฺกปฺโป โหติ, มิจฺฉาวายาโม โหติ, มิจฺฉาสมาธิ โหติ, วีริยพลํ โหติ, สมาธิพลํ โหติ, อหิริกพลํ โหติ, อโนตฺตปฺปพลํ โหติ, โทโส โหติ, โมโห โหติ, พฺยาปาโท โหติ, อหิริกํ โหติ, อโนตฺตปฺปํ โหติ, สมโถ โหติ, ปคฺคาโห โหติ, อวิกฺเขโป โหติ. เย วา ปน ตสฺมิํ สมเย อญฺเญปิ อตฺถิ ปฏิจฺจ\-สมุปฺปนฺนา อรูปิโน ธมฺมา. อิเม ธมฺมา อกุสลา.}

\proseitem{414}{กตโม ตสฺมิํ สมเย ผสฺโส โหติ. โย ตสฺมิํ สมเย ผสฺโส ผุสนา สมฺผุสนา สมฺผุสิตตฺตํ. อยํ ตสฺมิํ สมเย ผสฺโส โหติ.}

\proseitem{415}{กตมา ตสฺมิํ สมเย เวทนา โหติ. ยํ ตสฺมิํ สมเย ตชฺชา\-มโนวิญฺญาณธาตุ\-สมฺผสฺสชํ เจตสิกํ อสาตํ เจตสิกํ ทุกฺขํ เจโต\-สมฺผสฺสชํ อสาตํ ทุกฺขํ เวทยิตํ เจโต\-สมฺผสฺสชา อสาตา ทุกฺขา เวทนา. อยํ ตสฺมิํ สมเย เวทนา โหติ ฯเปฯ.}

\proseitem{416}{กตมํ ตสฺมิํ สมเย ทุกฺขํ โหติ. ยํ ตสฺมิํ สมเย เจตสิกํ อสาตํ เจตสิกํ ทุกฺขํ เจโต\-สมฺผสฺสชํ อสาตํ ทุกฺขํ \csromanfolio{101}เวทยิตํ เจโต\-สมฺผสฺสชา อสาตา ทุกฺขา เวทนา. อิทํ ตสฺมิํ สมเย ทุกฺขํ โหติ ฯเปฯ.}

\proseitem{417}{กตมํ ตสฺมิํ สมเย โทมนสฺสิ\-นฺทฺริยํ โหติ. ยํ ตสฺมิํ สมเย เจตสิกํ อสาตํ เจตสิกํ ทุกฺขํ เจโต\-สมฺผสฺสชํ อสาตํ ทุกฺขํ เวทยิตํ เจโต\-สมฺผสฺสชา อสาตา ทุกฺขา เวทนา. อิทํ ตสฺมิํ สมเย โทมนสฺสิ\-นฺทฺริยํ โหติ ฯเปฯ.}

\proseitem{418}{กตโม ตสฺมิํ สมเย โทโส โหติ. โย ตสฺมิํ สมเย โทโส ทุสฺสนา ทุสฺสิตตฺตํ พฺยาปตฺติ พฺยาปชฺชนา พฺยาปชฺชิตตฺตํ วิโรโธ ปฏิวิโรโธ จณฺฑิกฺกํ อสุโรโป อนตฺตมนตา จิตฺตสฺส. อยํ ตสฺมิํ สมเย โทโส โหติ ฯเปฯ.}

\proseitem{419}{กตโม ตสฺมิํ สมเย พฺยาปาโท โหติ. โย ตสฺมิํ สมเย โทโส ทุสฺสนา ทุสฺสิตตฺตํ พฺยาปตฺติ พฺยาปชฺชนา พฺยาปชฺชิตตฺตํ วิโรโธ ปฏิวิโรโธ จณฺฑิกฺกํ อสุโรโป อนตฺตมนตา จิตฺตสฺส. อยํ ตสฺมิํ สมเย พฺยาปาโท โหติ ฯเปฯ. เย วา ปน ตสฺมิํ สมเย อญฺเญปิ อตฺถิ ปฏิจฺจ\-สมุปฺปนฺนา อรูปิโน ธมฺมา. อิเม ธมฺมา อกุสลา.}

\prose{ตสฺมิํ โข ปน สมเย จตฺตาโร ขนฺธา โหนฺติ, ทฺวายตนานิ โหนฺติ, เทฺว ธาตุโย โหนฺติ, ตโย อาหารา โหนฺติ, ปญฺจินฺทฺริยานิ โหนฺติ, จตุรงฺคิกํ ฌานํ โหติ, ติวงฺคิโก มคฺโค โหติ, จตฺตาริ พลานิ โหนฺติ, เทฺว เหตู โหนฺติ, เอโก ผสฺโส โหติ ฯเปฯ เอกํ ธมฺมายตนํ โหติ, เอกา ธมฺมาธาตุ โหติ. เย วา ปน ตสฺมิํ สมเย อญฺเญปิ อตฺถิ ปฏิจฺจ\-สมุปฺปนฺนา อรูปิโน ธมฺมา. อิเม ธมฺมา อกุสลา ฯเปฯ.}

\proseitem{420}{กตโม ตสฺมิํ สมเย สงฺขาร\-กฺขนฺโธ โหติ. ผสฺโส เจตนา วิตกฺโก วิจาโร จิตฺตสฺเส\-กคฺคตา วีริยินฺทฺริยํ สมาธินฺทฺริยํ ชีวิตินฺทฺริยํ มิจฺฉาสงฺกปฺโป มิจฺฉาวายาโม มิจฺฉาสมาธิ วีริยพลํ สมาธิพลํ อหิริกพลํ อโนตฺตปฺปพลํ โทโส โมโห พฺยาปาโท อหิริกํ อโนตฺตปฺปํ สมโถ ปคฺคาโห อวิกฺเขโป. เย วา ปน ตสฺมิํ สมเย อญฺเญปิ อตฺถิ ปฏิจฺจ\-สมุปฺปนฺนา อรูปิโน ธมฺมา ฐเปตฺวา เวทนา\-กฺขนฺธํ ฐเปตฺวา สญฺญากฺขนฺธํ ฐเปตฺวา วิญฺญาณ\-กฺขนฺธํ. อยํ ตสฺมิํ สมเย สงฺขาร\-กฺขนฺโธ โหติ ฯเปฯ. อิเม ธมฺมา อกุสลา.}

\proseitem{421}{\csromanfolio{102}กตเม ธมฺมา อกุสลา. ยสฺมิํ สมเย อกุสลํ จิตฺตํ อุปฺปนฺนํ โหติ โทมนสฺส\-สหคตํ ปฏิฆ\-สมฺปยุตฺตํ สสงฺขาเรน รูปารมฺมณํ วา ฯเปฯ ธมฺมารมฺมณํ วา ยํ ยํ วา ปนารพฺภ. ตสฺมิํ สมเย ผสฺโส โหติ ฯเปฯ อวิกฺเขโป โหติ ฯเปฯ. อิเม ธมฺมา อกุสลา.}

\proseitem{422}{กตเม ธมฺมา อกุสลา. ยสฺมิํ สมเย อกุสลํ จิตฺตํ อุปฺปนฺนํ โหติ อุเปกฺขา\-สหคตํ วิจิกิจฺฉา\-สมฺปยุตฺตํ รูปารมฺมณํ วา สทฺทารมฺมณํ วา คนฺธารมฺมณํ วา รสารมฺมณํ วา โผฏฺฐพฺพา\-รมฺมณํ วา ธมฺมารมฺมณํ วา ยํ ยํ วา ปนารพฺภ. ตสฺมิํ สมเย ผสฺโส โหติ, เวทนา โหติ, สญฺญา โหติ, เจตนา โหติ, จิตฺตํ โหติ, วิตกฺโก โหติ, วิจาโร โหติ, อุเปกฺขา โหติ, จิตฺตสฺเส\-กคฺคตา โหติ, วีริยินฺทฺริยํ โหติ, มนินฺทฺริยํ โหติ, อุเปกฺขินฺทฺริยํ โหติ, ชีวิตินฺทฺริยํ โหติ, มิจฺฉาสงฺกปฺโป โหติ, มิจฺฉาวายาโม โหติ, วีริยพลํ โหติ, อหิริกพลํ โหติ, อโนตฺตปฺปพลํ โหติ, วิจิกิจฺฉา โหติ, โมโห โหติ, อหิริกํ โหติ, อโนตฺตปฺปํ โหติ, ปคฺคาโห โหติ. เย วา ปน ตสฺมิํ สมเย อญฺเญปิ อตฺถิ ปฏิจฺจ\-สมุปฺปนฺนา อรูปิโน ธมฺมา. อิเม ธมฺมา อกุสลา.}

\proseitem{423}{กตโม ตสฺมิํ สมเย ผสฺโส โหติ. โย ตสฺมิํ สมเย ผสฺโส ผุสนา สมฺผุสนา สมฺผุสิตตฺตํ. อยํ ตสฺมิํ สมเย ผสฺโส โหติ ฯเปฯ.}

\proseitem{424}{กตมา ตสฺมิํ สมเย จิตฺตสฺเส\-กคฺคตา โหติ. ยา ตสฺมิํ สมเย จิตฺตสฺส ฐิติ. อยํ ตสฺมิํ สมเย จิตฺตสฺเส\-กคฺคตา โหติ ฯเปฯ.}

\proseitem{425}{กตมา ตสฺมิํ สมเย วิจิกิจฺฉา โหติ. ยา ตสฺมิํ สมเย กงฺขา กงฺขายนา กงฺขายิตตฺตํ วิมติ วิจิกิจฺฉา เทฺวฬฺหกํ เทฺวธาปโถ สํสโย อเนกํสคฺคาโห อาสปฺปนา ปริสปฺปนา อปริโยคาหนา ถมฺภิตตฺตํ จิตฺตสฺส มโนวิเลโข. อยํ ตสฺมิํ สมเย วิจิกิจฺฉา โหติ ฯเปฯ. เย วา ปน ตสฺมิํ สมเย อญฺเญปิ อตฺถิ ปฏิจฺจ\-สมุปฺปนฺนา อรูปิโน ธมฺมา. อิเม ธมฺมา อกุสลา.}

\prose{ตสฺมิํ โข ปน สมเย จิตฺตาโร ขนฺธา โหนฺติ, ทฺวายตนานิ โหนฺติ, เทฺว ธาตุโย โหนฺติ, ตโย อาหารา โหนฺติ, จตฺตาริ อินฺทฺริยานิ \csromanfolio{103}โหนฺติ, จตุรงฺคิกํ ฌานํ โหติ, ทุวงฺคิโก มคฺโค โหติ, ตีณิ พลานิ โหนฺติ, เอโก เหตุ โหติ, เอโก ผสฺโส โหติ ฯเปฯ เอกํ ธมฺมายตนํ โหติ, เอกา ธมฺมธาตุ โหติ. เย วา ปน ตสฺมิํ สมเย อญฺเญปิ อตฺถิ ปฏิจฺจ\-สมุปฺปนฺนา อรูปิโน ธมฺมา. อิเม ธมฺมา อกุสลา ฯเปฯ.}

\proseitem{426}{กตโม ตสฺมิํ สมเย สงฺขาร\-กฺขนฺโธ โหติ. ผสฺโส เจตนา วิตกฺโก วิจาโร จิตฺตสฺเส\-กคฺคตา วีริยินฺทฺริยํ ชีวิตินฺทฺริยํ มิจฺฉาสงฺกปฺโป มิจฺฉาวายาโม วีริยพลํ อหิริกพลํ อโนตฺตปฺปพลํ วิจิกิจฺฉา โมโห อหิริกํ อโนตฺตปฺปํ ปคฺคาโห. เย วา ปน ตสฺมิํ สมเย อญฺเญปิ อตฺถิ ปฏิจฺจ\-สมุปฺปนฺนา อรูปิโน ธมฺมา ฐเปตฺวา เวทนา\-กฺขนฺธํ ฐเปตฺวา สญฺญากฺขนฺธํ ฐเปตฺวา วิญฺญาณ\-กฺขนฺธํ. อยํ ตสฺมิํ สมเย สงฺขาร\-กฺขนฺโธ โหติ ฯเปฯ. อิเม ธมฺมา อกุสลา.}

\proseitem{427}{กตเม ธมฺมา อกุสลา. ยสฺมิํ สมเย อกุสลํ จิตฺตํ อุปฺปนฺนํ โหติ อุเปกฺขา\-สหคตํ อุทฺธจฺจ\-สมฺปยุตฺตํ รูปารมฺมณํ วา ฯเปฯ ยํ ยํ วา ปนารพฺภ. ตสฺมิํ สมเย ผสฺโส โหติ, เวทนา โหติ, สญฺญา โหติ, เจตนา โหติ, จิตฺตํ โหติ, วิตกฺโก โหติ, วิจาโร โหติ, อุเปกฺขา โหติ, จิตฺตสฺเส\-กคฺคตา โหติ, วีริยินฺทฺริยํ โหติ, สมาธินฺทฺริยํ โหติ, มนินฺทฺริยํ โหติ, อุเปกฺขินฺทฺริยํ โหติ, ชีวิตินฺทฺริยํ โหติ, มิจฺฉาสงฺกปฺโป โหติ, มิจฺฉาวายาโม โหติ, มิจฺฉาสมาธิ โหติ, วีริยพลํ โหติ, สมาธิพลํ โหติ, อหิริกพลํ โหติ, อโนตฺตปฺปพลํ โหติ, อุทฺธจฺจํ โหติ, โมโห โหติ, อหิริกํ โหติ, อโนตฺตปฺปํ โหติ, สมโถ โหติ, ปคฺคาโห โหติ, อวิกฺเขโป โหติ. เย วา ปน ตสฺมิํ สมเย อญฺเญปิ อตฺถิ ปฏิจฺจ\-สมุปฺปนฺนา อรูปิโน ธมฺมา. อิเม ธมฺมา อกุสลา.}

\proseitem{428}{กตโม ตสฺมิํ สมเย ผสฺโส โหติ. โย ตสฺมิํ สมเย ผสฺโส ผุสนา สมฺผุสนา สมฺผุสิตตฺตํ. อยํ ตสฺมิํ สมเย ผสฺโส โหติ ฯเปฯ.}

\proseitem{429}{กตมํ ตสฺมิํ สมเย อุทฺธจฺจํ โหติ. ยํ ตสฺมิํ สมเย จิตฺตสฺส อุทฺธจฺจํ อวูปสโม เจตโส วิกฺเขโป ภนฺตตฺตํ จิตฺตสฺส. อิทํ ตสฺมิํ \csromanfolio{104}สมเย อุทฺธจฺจํ โหติ ฯเปฯ. เย วา ปน ตสฺมิํ สมเย อญฺเญปิ อตฺถิ ปฏิจฺจ\-สมุปฺปนฺนา อรูปิโน ธมฺมา. อิเม ธมฺมา อกุสลา.}

\prose{ตสฺมิํ โข ปน สมเย จตฺตาโร ขนฺธา โหนฺติ, ทฺวายตนานิ โหนฺติ, เทฺว ธาตุโย โหนฺติ, ตโย อาหารา โหนฺติ, ปญฺจินฺทฺริยานิ โหนฺติ, จตุรงฺคิกํ ฌานํ โหติ, ติวงฺคิโก มคฺโค โหติ, จตฺตาริ พลานิ โหนฺติ, เอโก เหตุ โหติ, เอโก ผสฺโส โหติ ฯเปฯ เอกํ ธมฺมายตนํ โหติ, เอโก ธมฺมาธาตุ โหติ. เย วา ปน ตสฺมิํ สมเย อญฺเญปิ อตฺถิ ปฏิจฺจ\-สมุปฺปนฺนา อรูปิโน ธมฺมา. อิเม ธมฺมา อกุสลา ฯเปฯ.}

\proseitem{430}{กตโม ตสฺมิํ สมเย สงฺขาร\-กฺขนฺโธ โหติ. ผสฺโส เจตนา วิตกฺโก วิจาโร จิตฺตสฺเส\-กคฺคตา วีริยินฺทฺริยํ สมาธินฺทฺริยํ ชีวิตินฺทฺริยํ มิจฺฉาสงฺกปฺโป มิจฺฉาวายาโม มิจฺฉาสมาธิ วีริยพลํ สมาธิพลํ อหิริกพลํ อโนตฺตปฺปพลํ อุทฺธจฺจํ โมโห อหิริกํ อโนตฺตปฺปํ สมโถ ปคฺคาโห อวิกฺเขโป. เย วา ปน ตสฺมิํ สมเย อญฺเญปิ อตฺถิ ปฏิจฺจ\-สมุปฺปนฺนา อรูปิโน ธมฺมา ฐเปตฺวา เวทนา\-กฺขนฺธํ ฐเปตฺวา สญฺญากฺขนฺธํ ฐเปตฺวา วิญฺญาณ\-กฺขนฺธํ. อยํ ตสฺมิํ สมเย สงฺขาร\-กฺขนฺโธ โหติ ฯเปฯ. อิเม ธมฺมา อกุสลา.}

\vspace{\tipitakaparskip}
\csromancenter{ทฺวาทส อกุสลจิตฺตานิ.}
\csromanmatikaanchor{mtk.52}
\csromantocmarkref{subsection}{อพฺยากตวิปากกุสลวิปากปญฺจวิญฺญาณ}{mtk.52}
\bookmark[dest={mtk.52},level=2]{อพฺยากตวิปากกุสลวิปากปญฺจวิญฺญาณ}
\csromanheader{2}{\textbf{อพฺยากตวิปาก}}
\csromanheader{2}{กุสลวิปาก\-ปญฺจวิญฺญาณ}
\proseitem{431}{กตเม ธมฺมา อพฺยากตา. ยสฺมิํ สมเย กามาวจรสฺส กุสลสฺส กมฺมสฺส กตตฺตา อุปจิตตฺตา วิปากํ จกฺขุวิญฺญาณํ อุปฺปนฺนํ โหติ อุเปกฺขา\-สหคตํ รูปารมฺมณํ. ตสฺมิํ สมเย ผสฺโส โหติ, เวทนา โหติ, สญฺญา โหติ, เจตนา โหติ, จิตฺตํ โหติ, อุเปกฺขา โหติ, จิตฺตสฺเส\-กคฺคตา โหติ, มนินฺทฺริยํ โหติ, อุเปกฺขินฺทฺริยํ โหติ, ชีวิตินฺทฺริยํ โหติ. เย วา ปน ตสฺมิํ สมเย อญฺเญปิ อตฺถิ ปฏิจฺจ\-สมุปฺปนฺนา อรูปิโน ธมฺมา. อิเม ธมฺมา อพฺยากตา.}

\chapterheadpageread{1. จิตฺตุปฺปาท\-กณฺฑ}
\proseitem{432}{\csromanfolio{105}กตโม ตสฺมิํ สมเย ผสฺโส โหติ. โย ตสฺมิํ สมเย ผสฺโส ผุสนา สมฺผุสนา สมฺผุสิตตฺตํ. อยํ ตสฺมิํ สมเย ผสฺโส โหติ.}

\proseitem{433}{กตมา ตสฺมิํ สมเย เวทนา โหติ. ยํ ตสฺมิํ สมเย ตชฺชา\-จกฺขุวิญฺญาณธาตุ\-สมฺผสฺสชํ เจตสิกํ เนว สาตํ นาสาตํ เจโต\-สมฺผสฺสชํ อทุกฺขมสุขํ เวทยิตํ เจโต\-สมฺผสฺสชา อทุกฺขมสุขา เวทนา. อยํ ตสฺมิํ สมเย เวทนา โหติ.}

\proseitem{434}{กตมา ตสฺมิํ สมเย สญฺญา โหติ. ยา ตสฺมิํ สมเย ตชฺชา\-จกฺขุวิญฺญาณธาตุ\-สมฺผสฺสชา สญฺญา สญฺชานนา สญฺชานิตตฺตํ. อยํ ตสฺมิํ สมเย สญฺญา โหติ.}

\proseitem{435}{กตมา ตสฺมิํ สมเย เจตนา โหติ. ยา ตสฺมิํ สมเย ตชฺชา\-จกฺขุวิญฺญาณธาตุ\-สมฺผสฺสชา เจตนา สญฺเจตนา เจตยิตตฺตํ. อยํ ตสฺมิํ สมเย เจตนา โหติ.}

\proseitem{436}{กตมํ ตสฺมิํ สมเย จิตฺตํ โหติ. ยํ ตสฺมิํ สมเย จิตฺตํ มโน มานสํ หทยํ ปณฺฑรํ มโน มนายตนํ มนินฺทฺริยํ วิญฺญาณํ วิญฺญาณ\-กฺขนฺโธ ตชฺชา\-จกฺขุวิญฺญาณ\-ธาตุ. อิทํ ตสฺมิํ สมเย จิตฺตํ โหติ.}

\proseitem{437}{กตมา ตสฺมิํ สมเย อุเปกฺขา โหติ. ยํ ตสฺมิํ สมเย เจตสิกํ เนว สาตํ นาสาตํ เจโต\-สมฺผสฺสชํ อทุกฺขมสุขํ เวทยิตํ เจโต\-สมฺผสฺสชา อทุกฺขมสุขา เวทนา. อยํ ตสฺมิํ สมเย อุเปกฺขา โหติ.}

\proseitem{438}{กตมา ตสฺมิํ สมเย จิตฺตสฺเส\-กคฺคตา โหติ. ยา ตสฺมิํ สมเย จิตฺตสฺส ฐิติ. อยํ ตสฺมิํ สมเย จิตฺตสฺเส\-กคฺคตา โหติ.}

\proseitem{439}{กตมํ ตสฺมิํ สมเย มนินฺทฺริยํ โหติ. ยํ ตสฺมิํ สมเย จิตฺตํ มโน มานสํ หทยํ ปณฺฑรํ มโน มนายตนํ มนินฺทฺริยํ วิญฺญาณํ วิญฺญาณ\-กฺขนฺโธ ตชฺชา\-จกฺขุวิญฺญาณ\-ธาตุ. อิทํ ตสฺมิํ สมเย มนินฺทฺริยํ โหติ.}

\proseitem{440}{\csromanfolio{106}กตมํ ตสฺมิํ สมเย อุเปกฺขินฺทฺริยํ โหติ. ยํ ตสฺมิํ สมเย เจตสิกํ เนว สาตํ นาสาตํ เจโต\-สมฺผสฺสชํ อทุกฺขมสุขํ เวทยิตํ เจโต\-สมฺผสฺสชา อทุกฺขมสุขา เวทนา. อิทํ ตสฺมิํ สมเย อุเปกฺขินฺทฺริยํ โหติ.}

\proseitem{441}{กตมํ ตสฺมิํ สมเย ชีวิตินฺทฺริยํ โหติ. โย เตสํ อรูปีนํ ธมฺมานํ อายุ ฐิติ ยปนา ยาปนา อิริยนา วตฺตนา ปาลนา ชีวิตํ ชีวิตินฺทฺริยํ. อิทํ ตสฺมิํ สมเย ชีวิตินฺทฺริยํ โหติ. เย วา ปน ตสฺมิํ สมเย อญฺเญปิ อตฺถิ ปฏิจฺจ\-สมุปฺปนฺนา อรูปิโน ธมฺมา. อิเม ธมฺมา อพฺยากตา.}

\prose{ตสฺมิํ โข ปน สมเย จตฺตาโร ขนฺธา โหนฺติ, ทฺวายตนานิ โหนฺติ, เทฺว ธาตุโย โหนฺติ, ตโย อาหารา โหนฺติ, ตีณินฺทฺริยานิ โหนฺติ, เอโก ผสฺโส โหติ ฯเปฯ เอกา จกฺขุวิญฺญาณ\-ธาตุ โหติ, เอกํ ธมฺมายตนํ โหติ, เอกา ธมฺมธาตุ โหติ. เย วา ปน ตสฺมิํ สมเย อญฺเญปิ อตฺถิ ปฏิจฺจ\-สมุปฺปนฺนา อรูปิโน ธมฺมา. อิเม ธมฺมา อพฺยากตา ฯเปฯ.}

\proseitem{442}{กตโม ตสฺมิํ สมเย สงฺขาร\-กฺขนฺโธ โหติ. ผสฺโส เจตนา จิตฺตสฺเส\-กคฺคตา ชีวิตินฺทฺริยํ. เย วา ปน ตสฺมิํ สมเย อญฺเญปิ อตฺถิ ปฏิจฺจ\-สมุปฺปนฺนา อรูปิโน ธมฺมา ฐเปตฺวา เวทนา\-กฺขนฺธํ ฐเปตฺวา สญฺญากฺขนฺธํ ฐเปตฺวา วิญฺญาณ\-กฺขนฺธํ. อยํ ตสฺมิํ สมเย สงฺขาร\-กฺขนฺโธ โหติ ฯเปฯ. อิเม ธมฺมา อพฺยากตา.}

\proseitem{443}{กตเม ธมฺมา อพฺยากตา. ยสฺมิํ สมเย กามาวจรสฺส กุสลสฺส กมฺมสฺส กตตฺตา อุปจิตตฺตา วิปากํ โสตวิญฺญาณํ อุปฺปนฺนํ โหติ อุเปกฺขา\-สหคตํ สทฺทารมฺมณํ ฯเปฯ ฆานวิญฺญาณํ อุปฺปนฺนํ โหติ อุเปกฺขา\-สหคตํ คนฺธารมฺมณํ ฯเปฯ ชีวฺหาวิญฺญาณํ อุปฺปนฺนํ โหติ อุเปกฺขา\-สหคตํ รสารมฺมณํ ฯเปฯ กายวิญฺญาณํ อุปฺปนฺนํ โหติ สุขสหคตํ โผฏฺฐพฺพา\-รมฺมณํ. ตสฺมิํ สมเย ผสฺโส โหติ, เวทนา โหติ, สญฺญา โหติ, เจตนา โหติ, จิตฺตํ โหติ, สุขํ โหติ, จิตฺตสฺเส\-กคฺคตา โหติ, มนินฺทฺริยํ โหติ, สุขินฺทฺริยํ โหติ, ชีวิตินฺทฺริยํ โหติ. เย วา ปน ตสฺมิํ สมเย อญฺเญปิ อตฺถิ ปฏิจฺจ\-สมุปฺปนฺนา อรูปิโน ธมฺมา. อิเม ธมฺมา อพฺยากตา.}

\chapterheadpageread{1. จิตฺตุปฺปาท\-กณฺฑ}
\proseitem{444}{\csromanfolio{107}กตโม ตสฺมิํ สมเย ผสฺโส โหติ. โย ตสฺมิํ สมเย ผสฺโส ผุสนา สมฺผุสนา สมฺผุสิตตฺตํ. อยํ ตสฺมิํ สมเย ผสฺโส โหติ.}

\proseitem{445}{กตมา ตสฺมิํ สมเย เวทนา โหติ. ยํ ตสฺมิํ สมเย ตชฺชา\-กายวิญฺญาณธาตุ\-สมฺผสฺสชํ กายิกํ สาตํ กายิกํ สุขํ กาย\-สมฺผสฺสชํ สาตํ สุขํ เวทยิตํ กาย\-สมฺผสฺสชา สาตา สุขา เวทนา. อยํ ตสฺมิํ สมเย เวทนา โหติ.}

\proseitem{446}{กตมา ตสฺมิํ สมเย สญฺญา โหติ. ยา ตสฺมิํ สมเย ตชฺชา\-กายวิญฺญาณธาตุ\-สมฺผสฺสชา สญฺญา สญฺชานนา สญฺชานิตตฺตํ. อยํ ตสฺมิํ สมเย สญฺญา โหติ.}

\proseitem{447}{กตมา ตสฺมิํ สมเย เจตนา โหติ. ยา ตสฺมิํ สมเย ตชฺชา\-กายวิญฺญาณธาตุ\-สมฺผสฺสชา เจตนา สญฺเจตนา เจตยิตตฺตํ. อยํ ตสฺมิํ สมเย เจตนา โหติ.}

\proseitem{448}{กตมํ ตสฺมิํ สมเย จิตฺตํ โหติ. ยํ ตสฺมิํ สมเย จิตฺตํ มโน มานสํ หทยํ ปณฺฑรํ มโน มนายตนํ มนินฺทฺริยํ วิญฺญาณํ วิญฺญาณ\-กฺขนฺโธ ตชฺชา\-กายวิญฺญาณ\-ธาตุ. อิทํ ตสฺมิํ สมเย จิตฺตํ โหติ.}

\proseitem{449}{กตมํ ตสฺมิํ สมเย สุขํ โหติ. ยํ ตสฺมิํ สมเย กายิกํ สาตํ กายิกํ สุขํ กาย\-สมฺผสฺสชํ สาตํ สุขํ เวทยิตํ กาย\-สมฺผสฺสชา สาตา สุขา เวทนา. อิทํ ตสฺมิํ สมเย สุขํ โหติ.}

\proseitem{450}{กตมา ตสฺมิํ สมเย จิตฺตสฺเส\-กคฺคตา โหติ. ยา ตสฺมิํ สมเย จิตฺตสฺส ฐิติ. อยํ ตสฺมิํ สมเย จิตฺตสฺเส\-กคฺคตา โหติ.}

\proseitem{451}{กตมํ ตสฺมิํ สมเย มนินฺทฺริยํ โหติ. ยํ ตสฺมิํ สมเย จิตฺตํ มโน มานสํ หทยํ ปณฺฑรํ มโน มนายตนํ มนินฺทฺริยํ วิญฺญาณํ วิญฺญาณ\-กฺขนฺโธ ตชฺชา\-กายวิญฺญาณ\-ธาตุ. อิทํ ตสฺมิํ สมเย มนินฺทฺริยํ โหติ.}

\proseitem{452}{\csromanfolio{108}กตมํ ตสฺมิํ สมเย สุขินฺทฺริยํ โหติ. ยํ ตสฺมิํ สมเย กายิกํ สาตํ กายิกํ สุขํ กาย\-สมฺผสฺสชํ สาตํ สุขํ เวทยิตํ กาย\-สมฺผสฺสชา สาตา สุขา เวทนา. อิทํ ตสฺมิํ สมเย สุขินฺทฺริยํ โหติ.}

\proseitem{453}{กตมํ ตสฺมิํ สมเย ชีวิตินฺทฺริยํ โหติ. โย เตสํ อรูปีนํ ธมฺมานํ อายุ ฐิติ ยปนา ยาปนา อิริยนา วตฺตนา ปาลนา ชีวิตํ ชีวิตินฺทฺริยํ. อิทํ ตสฺมิํ สมเย ชีวิตินฺทฺริยํ โหติ. เย วา ปน ตสฺมิํ สมเย อญฺเญปิ อตฺถิ ปฏิจฺจ\-สมุปฺปนฺนา อรูปิโน ธมฺมา. อิเม ธมฺมา อพฺยากตา.}

\prose{ตสฺมิํ โข ปน สมเย จตฺตาโร ขนฺธา โหนฺติ, ทฺวายตนานิ โหนฺติ, เทฺว ธาตุโย โหนฺติ, ตโย อาหารา โหนฺติ, ตีณินฺทฺริยานิ โหนฺติ, เอโก ผสฺโส โหติ ฯเปฯ เอกา กายวิญฺญาณ\-ธาตุ โหติ, เอกํ ธมฺมายตนํ โหติ, เอกา ธมฺมธาตุ โหติ. เย วา ปน ตสฺมิํ สมเย อญฺเญปิ อตฺถิ ปฏิจฺจ\-สมุปฺปนฺนา อรูปิโน ธมฺมา. อิเม ธมฺมา อพฺยากตา ฯเปฯ.}

\proseitem{454}{กตโม ตสฺมิํ สมเย สงฺขาร\-กฺขนฺโธ โหติ. ผสฺโส เจตนา จิตฺตสฺเส\-กคฺคตา ชีวิตินฺทฺริยํ. เย วา ปน ตสฺมิํ สมเย อญฺเญปิ อตฺถิ ปฏิจฺจ\-สมุปฺปนฺนา อรูปิโน ธมฺมา ฐเปตฺวา เวทนา\-กฺขนฺธํ ฐเปตฺวา สญฺญากฺขนฺธํ ฐเปตฺวา วิญฺญาณ\-กฺขนฺธํ. อยํ ตสฺมิํ สมเย สงฺขาร\-กฺขนฺโธ โหติ ฯเปฯ. อิเม ธมฺมา อพฺยากตา.}

\vspace{\tipitakaparskip}
\csromancenter{กุสลวิปากานิ ปญฺจวิญฺญาณานิ.}
\csromanmatikaanchor{mtk.53}
\csromantocmarkref{subsection}{กุสลวิปากมโนธาตุ}{mtk.53}
\bookmark[dest={mtk.53},level=2]{กุสลวิปากมโนธาตุ}
\csromanheader{2}{กุสลวิปาก\-มโนธาตุ}
\proseitem{455}{กตเม ธมฺมา อพฺยากตา. ยสฺมิํ สมเย กามาวจรสฺส กุสลสฺส กมฺมสฺส กตตฺตา อุปจิตตฺตา วิปากา มโนธาตุ อุปฺปนฺนา โหติ อุเปกฺขา\-สหคตา รูปารมฺมณา วา ฯเปฯ โผฏฺฐพฺพา\-รมฺมณา วา ยํ ยํ วา ปนารพฺภ. ตสฺมิํ สมเย ผสฺโส โหติ, เวทนา โหติ, สญฺญา โหติ, เจตนา โหติ, จิตฺตํ โหติ, วิตฺตกฺโก โหติ, วิจาโร \csromanfolio{109}โหติ, อุเปกฺขา โหติ, จิตฺตสฺเส\-กคฺคตา โหติ, มนินฺทฺริยํ โหติ, อุเปกฺขินฺทฺริยํ โหติ, ชีวิตินฺทฺริยํ โหติ. เย วา ปน ตสฺมิํ สมเย อญฺเญปิ อตฺถิ ปฏิจฺจ\-สมุปฺปนฺนา อรูปิโน ธมฺมา. อิเม ธมฺมา อพฺยากตา.}

\proseitem{456}{กตโม ตสฺมิํ สมเย ผสฺโส โหติ. โย ตสฺมิํ สมเย ผสฺโส ผุสนา สมฺผุสนา สมฺผุสิตตฺตํ. อยํ ตสฺมิํ สมเย ผสฺโส โหติ.}

\proseitem{457}{กตมา ตสฺมิํ สมเย เวทนา โหติ. ยํ ตสฺมิํ สมเย ตชฺชา\-มโนธาตุ\-สมฺผสฺสชํ เจตสิกํ เนว สาตํ นาสาตํ เจโต\-สมฺผสฺสชํ อทุกฺขมสุขํ เวทยิตํ เจโต\-สมฺผสฺสชา อทุกฺขมสุขา เวทนา. อยํ ตสฺมิํ สมเย เวทนา โหติ.}

\proseitem{458}{กตมา ตสฺมิํ สมเย สญฺญา โหติ. ยา ตสฺมิํ สมเย ตชฺชา\-มโนธาตุ\-สมฺผสฺสชา สญฺญา สญฺชานนา สญฺชานิตตฺตํ. อยํ ตสฺมิํ สมเย สญฺญา โหติ.}

\proseitem{459}{กตมา ตสฺมิํ สมเย เจตนา โหติ. ยา ตสฺมิํ สมเย ตชฺชา\-มโนธาตุ\-สมฺผสฺสชา เจตนา สญฺเจตนา เจตยิตตฺตํ. อยํ ตสฺมิํ สมเย เจตนา โหติ.}

\proseitem{460}{กตมํ ตสฺมิํ สมเย จิตฺตํ โหติ. ยํ ตสฺมิํ สมเย จิตฺตํ มโน มานสํ หทยํ ปณฺฑรํ มโน มนายตนํ มนินฺทฺริยํ วิญฺญาณํ วิญฺญาณ\-กฺขนฺโธ ตชฺชามโนธาตุ. อิทํ ตสฺมิํ สมเย จิตฺตํ โหติ.}

\proseitem{461}{กตโม ตสฺมิํ สมเย วิตกฺโก โหติ. โย ตสฺมิํ สมเย ตกฺโก วิตกฺโก สงฺกปฺโป อปฺปนา พฺยปฺปนา เจตโส อภินิโรปนา. อยํ ตสฺมิํ สมเย วิตกฺโก โหติ.}

\proseitem{462}{กตโม ตสฺมิํ สมเย วิจาโร โหติ. โย ตสฺมิํ สมเย จาโร วิจาโร อนุวิจาโร อุปวิจาโร จิตฺตสฺส อนุสนฺธานตา อนุเปกฺขนตา. อยํ ตสฺมิํ สมเย วิจาโร โหติ.}

\proseitem{463}{\csromanfolio{110}กตมา ตสฺมิํ สมเย อุเปกฺขา โหติ. ยํ ตสฺมิํ สมเย เจตสิกํ เนว สาตํ นาสาตํ เจโต\-สมฺผสฺสชํ อทุกฺขมสุขํ เวทยิตํ เจโต\-สมฺผสฺสชา อทุกฺขมสุขา เวทนา. อยํ ตสฺมิํ สมเย อุเปกฺขา โหติ.}

\proseitem{464}{กตมา ตสฺมิํ สมเย จิตฺตสฺเส\-กคฺคตา โหติ. ยา ตสฺมิํ สมเย จิตฺตสฺส ฐิติ. อยํ ตสฺมิํ สมเย จิตฺตสฺเส\-กคฺคตา โหติ.}

\proseitem{465}{กตมํ ตสฺมิํ สมเย มนินฺทฺริยํ โหติ. ยํ ตสฺมิํ สมเย จิตฺตํ มโน มานสํ หทยํ ปณฺฑรํ มโน มนายตนํ มนินฺทฺริยํ วิญฺญาณํ วิญฺญาณ\-กฺขนฺโธ ตชฺชามโนธาตุ. อิทํ ตสฺมิํ สมเย มนินฺทฺริยํ โหติ.}

\proseitem{466}{กตมํ ตสฺมิํ สมเย อุเปกฺขินฺทฺริยํ โหติ. ยํ ตสฺมิํ สมเย เจตสิกํ เนว สาตํ นาสาตํ เจโต\-สมฺผสฺสชํ อทุกฺขมสุขํ เวทยิตํ เจโต\-สมฺผสฺสชา อทุกฺขมสุขา เวทนา. อิทํ ตสฺมิํ สมเย อุเปกฺขินฺทฺริยํ โหติ.}

\proseitem{467}{กตมํ ตสฺมิํ สมเย ชีวิตินฺทฺริยํ โหติ. โย เตสํ อรูปีนํ ธมฺมานํ อายุ ฐิติ ยปนา ยาปนา อิริยนา วตฺตนา ปาลนา ชีวิตํ ชีวิตินฺทฺริยํ. อิทํ ตสฺมิํ สมเย ชีวิตินฺทฺริยํ โหติ. เย วา ปน ตสฺมิํ สมเย อญฺเญปิ อตฺถิ ปฏิจฺจ\-สมุปฺปนฺนา อรูปิโน ธมฺมา. อิเม ธมฺมา อพฺยากตา.}

\prose{ตสฺมิํ โข ปน สมเย จตฺตาโร ขนฺธา โหนฺติ, ทฺวายตนานิ โหนฺติ, เทฺว ธาตุโย โหนฺติ, ตโย อาหารา โหนฺติ, ตีณินฺทฺริยานิ โหนฺติ, เอโก ผสฺโส โหติ ฯเปฯ เอกา มโนธาตุ โหติ, เอกํ ธมฺมายตนํ โหติ, เอกา ธมฺมธาตุ โหติ. เย วา ปน ตสฺมิํ สมเย อญฺเญปิ อตฺถิ ปฏิจฺจ\-สมุปฺปนฺนา อรูปิโน ธมฺมา. อิเม ธมฺมา อพฺยากตา ฯเปฯ.}

\proseitem{468}{กตโม ตสฺมิํ สมเย สงฺขาร\-กฺขนฺโธ โหติ. ผสฺโส เจตนา วิตกฺโก วิจาโร จิตฺตสฺเส\-กคฺคตา ชีวิตินฺทฺริยํ. เย วา ปน ตสฺมิํ สมเย อญฺเญปิ อตฺถิ ปฏิจฺจ\-สมุปฺปนฺนา อรูปิโน ธมฺมา ฐเปตฺวา \csromanfolio{111}เวทนา\-กฺขนฺธํ ฐเปตฺวา สญฺญากฺขนฺธํ ฐเปตฺวา วิญฺญาณ\-กฺขนฺธํ. อยํ ตสฺมิํ สมเย สงฺขาร\-กฺขนฺโธ โหติ ฯเปฯ. อิเม ธมฺมา อพฺยากตา.}

\vspace{\tipitakaparskip}
\csromancenter{กุสลวิปากา มโนธาตุ.}
\csromanmatikaanchor{mtk.54}
\csromantocmarkref{subsection}{กุสลวิปากมโนวิญฺญาณธาตุโสมนสฺสสหคต}{mtk.54}
\bookmark[dest={mtk.54},level=2]{กุสลวิปากมโนวิญฺญาณธาตุโสมนสฺสสหคต}
\csromanheader{2}{กุสลวิปาก\-มโนวิญฺญาณธาตุ\-โสมนสฺสสหคต}
\proseitem{469}{กตเม ธมฺมา อพฺยากตา. ยสฺมิํ สมเย กามาวจรสฺส กุสลสฺส กมฺมสฺส กตตฺตา อุปจิตตฺตา วิปากา มโนวิญฺญาณ\-ธาตุ อุปฺปนฺนา โหติ โสมนสฺส\-สหคตา รูปารมฺมณา วา ฯเปฯ ธมฺมารมฺมณา วา ยํ ยํ วา ปนารพฺภ. ตสฺมิํ สมเย ผสฺโส โหติ, เวทนา โหติ, สญฺญา โหติ, เจตนา โหติ, จิตฺตํ โหติ, วิตกฺโก โหติ, วิจาโร โหติ, ปีติ โหติ, สุขํ โหติ, จิตฺตสฺเส\-กคฺคตา โหติ, มนินฺทฺริยํ โหติ, โสมนสฺสิ\-นฺทฺริยํ โหติ, ชีวิตินฺทฺริยํ โหติ. เย วา ปน ตสฺมิํ สมเย อญฺเญปิ อตฺถิ ปฏิจฺจ\-สมุปฺปนฺนา อรูปิโน ธมฺมา. อิเม ธมฺมา อพฺยากตา.}

\proseitem{470}{กตโม ตสฺมิํ สมเย ผสฺโส โหติ. โย ตสฺมิํ สมเย ผสฺโส ผุสนา สมฺผุสนา สมฺผุสิตตฺตํ. อยํ ตสฺมิํ สมเย ผสฺโส โหติ.}

\proseitem{471}{กตมา ตสฺมิํ สมเย เวทนา โหติ. ยํ ตสฺมิํ สมเย ตชฺชา\-มโนวิญฺญาณธาตุ\-สมฺผสฺสชํ เจตสิกํ สาตํ เจตสิกํ สุขํ เจโต\-สมฺผสฺสชํ สาตํ สุขํ เวทยิตํ เจโต\-สมฺผสฺสชา สาตา สุขา เวทนา. อยํ ตสฺมิํ สมเย เวทนา โหติ.}

\proseitem{472}{กตมา ตสฺมิํ สมเย สญฺญา โหติ. ยา ตสฺมิํ สมเย ตชฺชามโนวิญฺญาณาธาตุสมฺผสฺสชา สญฺญา สญฺชานนา สญฺชานิตตฺตํ. อยํ ตสฺมิํ สมเย สญฺญา โหติ.}

\proseitem{473}{กตมา ตสฺมิํ สมเย เจตนา โหติ. ยา ตสฺมิํ สมเย ตชฺชา\-มโนวิญฺญาณธาตุ\-สมฺผสฺสชา เจตนา สญฺเจตนา เจตยิตตฺตํ. อยํ ตสฺมิํ สมเย เจตนา โหติ.}

\proseitem{474}{\csromanfolio{112}กตมํ ตสฺมิํ สมเย จิตฺตํ โหติ. ยํ ตสฺมิํ สมเย จิตฺตํ มโน มานสํ หทยํ ปณฺฑรํ มโน มนายตนํ มนินฺทฺริยํ วิญฺญาณํ วิญฺญาณ\-กฺขนฺโธ ตชฺชา\-มโนวิญฺญาณ\-ธาตุ. อิทํ ตสฺมิํ สมเย จิตฺตํ โหติ.}

\proseitem{475}{กตโม ตสฺมิํ สมเย วิตกฺโก โหติ. โย ตสฺมิํ สมเย ตกฺโก วิตกฺโก สงฺกปฺโป อปฺปนา พฺยปฺปนา เจตโส อภินิโรปนา. อยํ ตสฺมิํ สมเย วิตกฺโก โหติ.}

\proseitem{476}{กตโม ตสฺมิํ สมเย วิจาโร โหติ. โย ตสฺมิํ สมเย จาโร วิจาโร อนุวิจาโร อุปวิจาโร จิตฺตสฺส อนุสนฺธานตา อนุเปกฺขนตา. อยํ ตสฺมิํ สมเย วิจาโร โหติ.}

\proseitem{477}{กตมา ตสฺมิํ สมเย ปีติ โหติ. ยา ตสฺมิํ สมเย ปีติ ปาโมชฺชํ อาโมทนา ปโมทนา หาโส ปหาโส วิตฺติ โอทคฺยํ อตฺตมนตา จิตฺตสฺส. อยํ ตสฺมิํ สมเย ปีติ โหติ.}

\proseitem{478}{กตมํ ตสฺมิํ สมเย สุขํ โหติ. ยํ ตสฺมิํ สมเย เจตสิกํ สาตํ เจตสิกํ สุขํ เจโต\-สมฺผสฺสชํ สาตํ สุขํ เวทยิตํ เจโต\-สมฺผสฺสชา สาตา สุขา เวทนา. อิทํ ตสฺมิํ สมเย สุขํ โหติ.}

\proseitem{479}{กตมา ตสฺมิํ สมเย จิตฺตสฺเส\-กคฺคตา โหติ. ยา ตสฺมิํ สมเย จิตฺตสฺส ฐิติ. อยํ ตสฺมิํ สมเย จิตฺตสฺเส\-กคฺคตา โหติ.}

\proseitem{480}{กตมํ ตสฺมิํ สมเย มนินฺทฺริยํ โหติ. ยํ ตสฺมิํ สมเย จิตฺตํ มโน มานสํ หทยํ ปณฺฑรํ มโน มนายตนํ มนินฺทฺริยํ วิญฺญาณํ วิญฺญาณ\-กฺขนฺโธ ตชฺชา\-มโนวิญฺญาณ\-ธาตุ. อิทํ ตสฺมิํ สมเย มนินฺทฺริยํ โหติ.}

\proseitem{481}{กตมํ ตสฺมิํ สมเย โสมนสฺสิ\-นฺทฺริยํ โหติ. ยํ ตสฺมิํ สมเย เจตสิกํ สาตํ เจตสิกํ สุขํ เจโต\-สมฺผสฺสชํ สาตํ สุขํ เวทยิตํ เจโต\-สมฺผสฺสชา สาตา สุขา เวทนา. อิทํ ตสฺมิํ สมเย โสมนสฺสิ\-นฺทฺริยํ โหติ.}

\chapterheadpageread{1. จิตฺตุปฺปาท\-กณฺฑ}
\proseitem{482}{\csromanfolio{113}กตมํ ตสฺมิํ สมเย ชีวิตินฺทฺริยํ โหติ. โย เตสํ อรูปีนํ ธมฺมานํ อายุ ฐิติ ยปนา ยาปนา อิริยนา วตฺตนา ปาลนา ชีวิตํ ชีวิตินฺทฺริยํ. อิทํ ตสฺมิํ สมเย ชีวิตินฺทฺริยํ โหติ. เย วา ปน ตสฺมิํ สมเย อญฺเญปิ อตฺถิ ปฏิจฺจ\-สมุปฺปนฺนา อรูปิโน ธมฺมา. อิเม ธมฺมา อพฺยากตา.}

\prose{ตสฺมิํ โข ปน สมเย จตฺตาโร ขนฺธา โหนฺติ, ทฺวายตนานิ โหนฺติ, เทฺว ธาตุโย โหนฺติ, ตโย อาหารา โหนฺติ, ตีณินฺทฺริยานิ โหนฺติ, เอโก ผสฺโส โหติ ฯเปฯ เอกา มโนวิญฺญาณ\-ธาตุ โหติ, เอกํ ธมฺมายตนํ โหติ, เอกา ธมฺมธาตุ โหติ. เย วา ปน ตสฺมิํ สมเย อญฺเญปิ อตฺถิ ปฏิจฺจ\-สมุปฺปนฺนา อรูปิโน ธมฺมา. อิเม ธมฺมา อพฺยากตา ฯเปฯ.}

\proseitem{483}{กตโม ตสฺมิํ สมเย สงฺขาร\-กฺขนฺโธ โหติ. ผสฺโส เจตนา วิตกฺโก วิจาโร ปีติ จิตฺตสฺเส\-กคฺคตา ชีวิตินฺทฺริยํ. เย วา ปน ตสฺมิํ สมเย อญฺเญปิ อตฺถิ ปฏิจฺจ\-สมุปฺปนฺนา อรูปิโน ธมฺมา ฐเปตฺวา เวทนา\-กฺขนฺธํ ฐเปตฺวา สญฺญากฺขนฺธํ ฐเปตฺวา วิญฺญาณ\-กฺขนฺธํ. อยํ ตสฺมิํ สมเย สงฺขาร\-กฺขนฺโธ โหติ ฯเปฯ. อิเม ธมฺมา อพฺยากตา.}

\vspace{\tipitakaparskip}
\csromancenter{กุสลวิปากา มโนวิญฺญาณ\-ธาตุ โสมนสฺส\-สหคตา.}
\csromanmatikaanchor{mtk.55}
\csromantocmarkref{subsection}{กุสลวิปากมโนวิญฺญาณธาตุอุเปกฺขาสหคต}{mtk.55}
\bookmark[dest={mtk.55},level=2]{กุสลวิปากมโนวิญฺญาณธาตุอุเปกฺขาสหคต}
\csromanheader{2}{\textbf{กุสลวิปากมโนวฺ}อิญฺญาณธาตุอุเปกฺขาสหคต}
\proseitem{484}{กตเม ธมฺมา อพฺยากตา. ยสฺมิํ สมเย กามาวจรสฺส กุสลสฺส กมฺมสฺส กตตฺตา อุปจิตตฺตา วิปากา มโนวิญฺญาณ\-ธาตุ อุปฺปนฺนา โหติ อุเปกฺขา\-สหคตา รูปารมฺมณา วา ฯเปฯ ธมฺมารมฺมณา วา ยํ ยํ วา ปนารพฺภ. ตสฺมิํ สมเย ผสฺโส โหติ, เวทนา โหติ, สญฺญา โหติ, เจตนา โหติ, จิตฺตํ โหติ, วิตกฺโก โหติ, วิจาโร โหติ, อุเปกฺขา โหติ, จิตฺตสฺเส\-กคฺคตา โหติ, มนินฺทฺริยํ โหติ, อุเปกฺขินฺทฺริยํ โหติ, ชีวิตินฺทฺริยํ โหติ. เย วา ปน ตสฺมิํ สมเย อญฺเญปิ อตฺถิ ปฏิจฺจ\-สมุปฺปนฺนา อรูปิโน ธมฺมา. อิเม ธมฺมา อพฺยากตา.}

\proseitem{485}{\csromanfolio{114}กตโม ตสฺมิํ สมเย ผสฺโส โหติ. โย ตสฺมิํ สมเย ผสฺโส ผุสนา สมฺผุสนา สมฺผุสิตตฺตํ. อยํ ตสฺมิํ สมเย ผสฺโส โหติ.}

\proseitem{486}{กตมา ตสฺมิํ สมเย เวทนา โหติ. ยํ ตสฺมิํ สมเย ตชฺชา\-มโนวิญฺญาณธาตุ\-สมฺผสฺสชํ เจตสิกํ เนว สาตํ นาสาตํ เจโต\-สมฺผสฺสชํ อทุกฺขมสุขํ เวทยิตํ เจโต\-สมฺผสฺสชา อทุกฺขมสุขา เวทนา. อยํ ตสฺมิํ สมเย เวทนา โหติ.}

\proseitem{487}{กตมา ตสฺมิํ สมเย สญฺญา โหติ. ยา ตสฺมิํ สมเย ตชฺชา\-มโนวิญฺญาณธาตุ\-สมฺผสฺสชา สญฺญา สญฺชานนา สญฺชานิตตฺตํ. อยํ ตสฺมิํ สมเย สญฺญา โหติ.}

\proseitem{488}{กตมา ตสฺมิํ สมเย เจตนา โหติ. ยา ตสฺมิํ สมเย ตชฺชา\-มโนวิญฺญาณธาตุ\-สมฺผสฺสชา เจตนา สญฺเจตนา เจตยิตตฺตํ. อยํ ตสฺมิํ สมเย เจตนา โหติ.}

\proseitem{489}{กตมํ ตสฺมิํ สมเย จิตฺตํ โหติ. ยํ ตสฺมิํ สมเย จิตฺตํ มโน มานสํ หทยํ ปณฺฑรํ มโน มนายตนํ มนินฺทฺริยํ วิญฺญาณํ วิญฺญาณ\-กฺขนฺโธ ตชฺชา\-มโนวิญฺญาณ\-ธาตุ. อิทํ ตสฺมิํ สมเย จิตฺตํ โหติ.}

\proseitem{490}{กตโม ตสฺมิํ สมเย วิตกฺโก โหติ. โย ตสฺมิํ สมเย ตกฺโก วิตกฺโก สงฺกปฺโป อปฺปนา พฺยปฺปนา เจตโส อภินิโรปนา. อยํ ตสฺมิํ สมเย วิตกฺโก โหติ.}

\proseitem{491}{กตโม ตสฺมิํ สมเย วิจาโร โหติ. โย ตสฺมิํ สมเย จาโร วิจาโร อนุวิจาโร อุปวิจาโร จิตฺตสฺส อนุสนฺธานตา อนุเปกฺขนตา. อยํ ตสฺมิํ สมเย วิจาโร โหติ.}

\proseitem{492}{กตมา ตสฺมิํ สมเย อุเปกฺขา โหติ. ยํ ตสฺมิํ สมเย เจตสิกํ เนว สาตํ นาสาตํ เจโต\-สมฺผสฺสชํ อทุกฺขมสุขํ เวทยิตํ เจโต\-สมฺผสฺสชา อทุกฺขมสุขา เวทนา. อยํ ตสฺมิํ สมเย อุเปกฺขา โหติ.}

\chapterheadpageread{1. จิตฺตุปฺปาท\-กณฺฑ}
\proseitem{493}{\csromanfolio{115}กตมา ตสฺมิํ สมเย จิตฺตสฺเส\-กคฺคตา โหติ. ยา ตสฺมิํ สมเย จิตฺตสฺส ฐิติ. อยํ ตสฺมิํ สมเย จิตฺตสฺเส\-กคฺคตา โหติ.}

\proseitem{494}{กตํ ตสฺมิํ สมเย มนินฺทฺริยํ โหติ. ยํ ตสฺมิํ สมเย จิตฺตํ มโน มานสํ หทยํ ปณฺฑรํ มโน มนายตนํ มนินฺทฺริยํ วิญฺญาณํ วิญฺญาณ\-กฺขนฺโธ ตชฺชา\-มโนวิญฺญาณ\-ธาตุ. อิทํ ตสฺมิํ สมเย มนินฺทฺริยํ โหติ.}

\proseitem{495}{กตมํ ตสฺมิํ สมเย อุเปกฺขินฺทฺริยํ โหติ. ยํ ตสฺมิํ สมเย เจตสิกํ เนว สาตํ นาสาตํ เจโต\-สมฺผสฺสชํ อทุกฺขมสุขํ เวทยิตํ เจโต\-สมฺผสฺสชา อทุกฺขมสุขา เวทนา. อิทํ ตสฺมิํ สมเย อุเปกฺขินฺทฺริยํ โหติ.}

\proseitem{496}{กตมํ ตสฺมิํ สมเย ชีวิตินฺทฺริยํ โหติ. โย เตสํ อรูปีนํ ธมฺมานํ อายุ ฐิติ ยปนา ยาปนา อิริยนา วตฺตนา ปาลนา ชีวิตํ ชีวิตินฺทฺริยํ. อิทํ ตสฺมิํ สมเย ชีวิตินฺทฺริยํ โหติ. เย วา ปน ตสฺมิํ สมเย อญฺเญปิ อตฺถิ ปฏิจฺจ\-สมุปฺปนฺนา อรูปิโน ธมฺมา. อิเม ธมฺมา อพฺยากตา.}

\prose{ตสฺมิํ โข ปน สมเย จตฺตาโร ขนฺธา โหนฺติ, ทฺวายตนานิ โหนฺติ, เทฺว ธาตุโย โหนฺติ, ตโย อาหารา โหนฺติ, ตีณินฺทฺริยานิ โหนฺติ, เอโก ผสฺโส โหติ ฯเปฯ เอกา มโนวิญฺญาณ\-ธาตุ โหติ, เอกํ ธมฺมายตนํ โหติ, เอกา ธมฺมธาตุ โหติ. เย วา ปน ตสฺมิํ สมเย อญฺเญปิ อตฺถิ ปฏิจฺจ\-สมุปฺปนฺนา อรูปิโน ธมฺมา. อิเม ธมฺมา อพฺยากตา ฯเปฯ.}

\proseitem{497}{กตโม ตสฺมิํ สมเย สงฺขาร\-กฺขนฺโธ โหติ. ผสฺโส เจตนา วิตกฺโก วิจาโร จิตฺตสฺเส\-กคฺคตา ชีวิตินฺทฺริยํ. เย วา ปน ตสฺมิํ สมเย อญฺเญปิ อตฺถิ ปฏิจฺจ\-สมุปฺปนฺนา อรูปิโน ธมฺมา ฐเปตฺวา เวทนา\-กฺขนฺธํ ฐเปตฺวา สญฺญากฺขนฺธํ ฐเปตฺวา วิญฺญาณ\-กฺขนฺธํ. อยํ ตสฺมิํ สมเย สงฺขาร\-กฺขนฺโธ โหติ ฯเปฯ. อิเม ธมฺมา อพฺยากตา.}

\vspace{\tipitakaparskip}
\csromancenter{กุสลวิปากา อุเปกฺขา\-สหคตา มโนวิญฺญาณ\-ธาตุ.}
\csromanmatikaanchor{mtk.56}
\csromantocmarkref{subsection}{อฏฺฐมหาวิปาก}{mtk.56}
\bookmark[dest={mtk.56},level=2]{อฏฺฐมหาวิปาก}
\csromanheader{2}{อฏฺฐ\-มหา\-วิปาก}
\proseitem{498}{\csromanfolio{116}กตเม ธมฺมา อพฺยากตา. ยสฺมิํ สมเย กามาวจรสฺส กุสลสฺส กมฺมสฺส กตตฺตา อุปจิตตฺตา วิปากา มโนวิญฺญาณ\-ธาตุ อุปฺปนฺนา โหติ โสมนสฺส\-สหคตา ญาณสมฺปยุตฺตา ฯเปฯ โสมนสฺส สหคตา ญาณสมฺปยุตฺตา สสงฺขาเรน ฯเปฯ โสมนสฺส\-สหคตา ญาณวิปฺปยุตฺตา ฯเปฯ โสมนสฺส\-สหคตา ญาณวิปฺปยุตฺตา สสงฺขาเรน ฯเปฯ อุเปกฺขา\-สหคตา ญาณสมฺปยุตฺตา ฯเปฯ อุเปกฺขา\-สหคตา ญาณสมฺปยุตฺตา สสงฺขาเรน ฯเปฯ อุเปกฺขา\-สหคตา ญาณวิปฺปยุตฺตา ฯเปฯ อุเปกฺขา\-สหคตา ญาณวิปฺปยุตฺตา สสงฺขาเรน รูปารมฺมณา วา ฯเปฯ ธมฺมารมฺมณา วา ยํ ยํ วา ปนารพฺภ. ตสฺมิํ สมเย ผสฺโส โหติ ฯเปฯ อวิกฺเขโป โหติ ฯเปฯ. อิเม ธมฺมา อพฺยากตา ฯเปฯ อโลโภ อพฺยากตมูลํ ฯเปฯ อโทโส อพฺยากตมูลํ ฯเปฯ. อิเม ธมฺมา อพฺยากตา.}

\vspace{\tipitakaparskip}
\csromancenter{อฏฺฐ\-มหา\-วิปากา.}
\csromanmatikaanchor{mtk.57}
\csromantocmarkref{subsection}{รูปาวจรวิปาก}{mtk.57}
\bookmark[dest={mtk.57},level=2]{รูปาวจรวิปาก}
\csromanheader{2}{รูปาวจร\-วิปาก}
\proseitem{499}{กตเม ธมฺมา อพฺยากตา. ยสฺมิํ สมเย รูปูปปตฺติยา มคฺคํ ภาเวติ วิวิจฺเจว กาเมหิ ฯเปฯ ปฐมํ ฌานํ อุปสมฺปชฺช วิหรติ ปถวีกสิณํ. ตสฺมิํ สมเย ผสฺโส โหติ ฯเปฯ อวิกฺเขโป โหติ ฯเปฯ. อิเม ธมฺมา กุสลา. ตสฺเสว รูปาวจรสฺส กุสลสฺส กมฺมสฺส กตตฺตา อุปจิตฺตา วิปากํ วิวิจฺเจว กาเมหิ ฯเปฯ ปฐมํ ฌานํ อุปสมฺปชฺช วิหรติ ปถวีกสิณํ. ตสฺมิํ สมเย ผสฺโส โหติ ฯเปฯ อวิกฺเขโป โหติ ฯเปฯ. อิเม ธมฺมา อพฺยากตา.}

\proseitem{500}{กตเม ธมฺมา อพฺยากตา. ยสฺมิํ สมเย รูปูปปตฺติยา มคฺคํ ภาเวติ วิตกฺก\-วิจารานํ วูปสมา ฯเปฯ ทุติยํ ฌานํ ฯเปฯ ตติยํ ฌานํ ฯเปฯ จตุตฺถํ ฌานํ ฯเปฯ ปฐมํ ฌานํ ฯเปฯ ปญฺจมํ ฌานํ อุปสมฺปชฺช วิหรติ ปถวีกสิณํ. ตสฺมิํ สมเย ผสฺโส โหติ ฯเปฯ อวิกฺเขโป โหติ ฯเปฯ. อิเม ธมฺมา กุสลา. ตสฺเสว รูปาวจรสฺส กุสลสฺส กมฺมสฺส กตตฺตา อุปจิตตฺตา วิปากํ สุขสฺส จ ปหานา ฯเปฯ ปญฺจมํ ฌานํ อุปสมฺปชฺช วิหรติ ปถวีกสิณํ. \csromanfolio{117}ตสฺมิํ สมเย ผสฺโส โหติ ฯเปฯ อวิกฺเขโป โหติ ฯเปฯ. อิเม ธมฺมา อพฺยากตา.}

\vspace{\tipitakaparskip}
\csromancenter{รูปาวจร\-วิปากา.}
\csromanmatikaanchor{mtk.58}
\csromantocmarkref{subsection}{อรูปาวจรวิปาก}{mtk.58}
\bookmark[dest={mtk.58},level=2]{อรูปาวจรวิปาก}
\csromanheader{2}{อรูปาวจร\-วิปาก}
\proseitem{501}{กตเม ธมฺมา อพฺยากตา. ยสฺมิํ สมเย อรูปูปปตฺติยา มคฺคํ ภาเวติ สพฺพโส รูปสญฺญานํ สมติกฺกมา ปฏิฆ\-สญฺญานํ อตฺถงฺคมา นานตฺต\-สญฺญานํ อมนสิการา อากาสา\-นญฺจา\-ยตนสญฺญาสหคตํ สุขสฺส จ ปหานา ฯเปฯ จตุตฺถํ ฌานํ อุปสมฺปชฺช วิหรติ. ตสฺมิํ สมเย ผสฺโส โหติ ฯเปฯ อวิกฺเขโป โหติ ฯเปฯ. อิเม ธมฺมา กุสลา. ตสฺเสว อรูปาวจรสฺส กุสลสฺส กมฺมสฺส กตตฺตา อุปจิตตฺตา วิปากํ สพฺพโส รูปสญฺญานํ สมติกฺกมา ปฏิฆ\-สญฺญานํ อตฺถงฺคมา นานตฺต\-สญฺญานํ อมนสิการา อากาสา\-นญฺจา\-ยตนสญฺญาสหคตํ สุขสฺส จ ปหานา ฯเปฯ จตุตฺถํ ฌานํ อุปสมฺปชฺช วิหรติ. ตสฺมิํ สมเย ผสฺโส โหติ ฯเปฯ อวิกฺเขโป โหติ ฯเปฯ. อิเม ธมฺมา อพฺยากตา.}

\proseitem{502}{กตเม ธมฺมา อพฺยากตา. ยสฺมิํ สมเย อรูปูปปตฺติยา มคฺคํ ภาเวติ สพฺพโส อากาสา\-นญฺจา\-ยตนํ สมติกฺกมฺม วิญฺญาณญฺจายตน\-สญฺญาสหคตํ สุขสฺส จ ปหานา ฯเปฯ จตุตฺถํ ฌานํ อุปสมฺปชฺช วิหรติ. ตสฺมิํ สมเย ผสฺโส โหติ ฯเปฯ อวิกฺเขโป โหติ ฯเปฯ. อิเม ธมฺมา กุสลา. ตสฺเสว อรูปาวจรสฺส กุสลสฺส กมฺมสฺส กตตฺตา อุปจิตตฺตา วิปากํ สพฺพโส อากาสา\-นญฺจา\-ยตนํ สมติกฺกมฺม วิญฺญาณญฺจายตน\-สญฺญาสหคตํ สุขสฺส จ ปหานา ฯเปฯ จตุตฺถํ ฌานํ อุปสมฺปชฺช วิหรติ. ตสฺมิํ สมเย ผสฺโส โหติ ฯเปฯ อวิกฺเขโป โหติ ฯเปฯ. อิเม ธมฺมา อพฺยากตา.}

\proseitem{503}{กตเม ธมฺมา อพฺยากตา. ยสฺมิํ สมเย อรูปูปปตฺติยา มคฺคํ ภาเวติ สพฺพโส วิญฺญาณญฺจา\-ยตนํ สมติกฺกมฺม อากิญฺจญฺญายตน\-สญฺญาสหคตํ สุขสฺส จ ปหานา ฯเปฯ จตุตฺถํ ฌานํ อุปสมฺปชฺช วิหรติ. ตสฺมิํ สมเย ผสฺโส โหติ ฯเปฯ อวิกฺเขโป โหติ ฯเปฯ. อิเม ธมฺมา กุสลา. ตสฺเสว อรูปาวจรสฺส กุสลสฺส กมฺมสฺส กตตฺตา อุปจิตตฺตา \csromanfolio{118}วิปากํ สพฺพโส วิญฺญาณญฺจา\-ยตนํ สมติกฺกมฺม อากิญฺจญฺญายตน\-สญฺญาสหคตํ สุขสฺส จ ปหานา ฯเปฯ จตุตฺถํ ฌานํ อุปสมฺปชฺช วิหรติ. ตสฺมิํ สมเย ผสฺโส โหติ ฯเปฯ อวิกฺเขโป โหติ ฯเปฯ. อิเม ธมฺมา อพฺยากตา.}

\proseitem{504}{กตเม ธมฺมา อพฺยากตา. ยสฺมิํ สมเย อรูปูปปตฺติยา มคฺคํ ภาเวติ สพฺพโส อากิญฺจญฺญา\-ยตนํ สมติกฺกมฺม เนว\-สญฺญา\-นา\-สญฺญา\-ยตนสญฺญาสหคตํ สุขสฺส จ ปหานา ฯเปฯ จตุตฺถํ ฌานํ อุปสมฺปชฺช วิหรติ. ตสฺมิํ สมเย ผสฺโส โหติ ฯเปฯ อวิกฺเขโป โหติ ฯเปฯ. อิเม ธมฺมา กุสลา. ตสฺเสว อรูปาวจรสฺส กุสลสฺส กมฺมสฺส กตตฺตา อุปจิตตฺตา วิปากํ สพฺพโส อากิญฺจญฺญา\-ยตนํ สมติกฺกมฺม เนว\-สญฺญา\-นา\-สญฺญา\-ยตนสญฺญาสหคตํ สุขสฺส จ ปหานา ฯเปฯ จตุตฺถํ ฌานํ อุปสมฺปชฺช วิหรติ. ตสฺมิํ สมเย ผสฺโส โหติ ฯเปฯ อวิกฺเขโป โหติ ฯเปฯ. อิเม ธมฺมา อพฺยากตา.}

\vspace{\tipitakaparskip}
\csromancenter{อรูปาวจร\-วิปากา.}
\csromanheader{2}{โลกุตฺตร\-วิปาก\-ปฐมมคฺค\-วิปาก}
\csromanmatikaanchor{mtk.59}
\csromantocmarkref{subsection}{โลกุตฺตรวิปากปฐมมคฺควิปากสุทฺธิกปฏิปทา}{mtk.59}
\bookmark[dest={mtk.59},level=2]{โลกุตฺตรวิปากปฐมมคฺควิปากสุทฺธิกปฏิปทา}
\csromanheader{2}{สุทฺธิก\-ปฏิปทา}
\proseitem{505}{กตเม ธมฺมา อพฺยากตา. ยสฺมิํ สมเย โลกุตฺตรํ ฌานํ ภาเวติ นิยฺยานิกํ อปจยคามิํ ทิฏฺฐิคตานํ ปหานาย ปฐมาย ภูมิยา ปตฺติยา วิวิจฺเจว กาเมหิ ฯเปฯ ปฐมํ ฌานํ อุปสมฺปชฺช วิหรติ ทุกฺข\-ปฏิปทํ ทนฺธาภิญฺญํ. ตสฺมิํ สมเย ผสฺโส โหติ ฯเปฯ อวิกฺเขโป โหติ ฯเปฯ. อิเม ธมฺมา กุสลา. ตสฺเสว โลกุตฺตรสฺส กุสลสฺส ฌานสฺส กตตฺตา ภาวิตตฺตา วิปากํ วิวิจฺเจว กาเมหิ ฯเปฯ ปฐมํ ฌานํ อุปสมฺปชฺช วิรหติ ทุกฺข\-ปฏิปทํ ทนฺธาภิญฺญํ สุญฺญตํ. ตสฺมิํ สมเย ผสฺโส โหติ ฯเปฯ อญฺญินฺทฺริยํ โหติ ฯเปฯ อวิกฺเขโป โหติ ฯเปฯ. อิเม ธมฺมา อพฺยากตา.}

\proseitem{506}{กตเม ธมฺมา อพฺยากตา. ยสฺมิํ สมเย โลกุตฺตรํ ฌานํ ภาเวติ นิยฺยานิกํ อปจยคามิํ ทิฏฺฐิคตานํ ปหานาย ปฐมาย ภูมิยา ปตฺติยา วิวิจฺเจว กาเมหิ ฯเปฯ ปฐมํ ฌานํ อุปสมฺปชฺช วิหรติ ทุกฺข\-ปฏิปทํ ทนฺธาภิญฺญํ. ตสฺมิํ สมเย ผสฺโส โหติ ฯเปฯ อวิกฺเขโป โหติ ฯเปฯ. อิเม \csromanfolio{119}ธมฺมา กุสลา. ตสฺเสว โลกุตฺตรสฺส กุสลสฺส ฌานสฺส กตตฺตา ภาวิตตฺตา วิปากํ วิวิจฺเจว กาเมหิ ฯเปฯ ปฐมํ ฌานํ อุปสมฺปชฺช วิหรติ ทุกฺข\-ปฏิปทํ ทนฺธาภิญฺญํ อนิมิตฺตํ. ตสฺมิํ สมเย ผสฺโส โหติ ฯเปฯ อญฺญินฺทฺริยํ โหติ ฯเปฯ อวิกฺเขโป โหติ ฯเปฯ. อิเม ธมฺมา อพฺยากตา.}

\proseitem{507}{กตเม ธมฺมา อพฺยากตา. ยสฺมิํ สมเย โลกุตฺตรํ ฌานํ ภาเวติ นิยฺยานิกํ อปจยคามิํ ทิฏฺฐิคตานํ ปหานาย ปฐมาย ภูมิยา ปตฺติยา วิวิจฺเจว กาเมหิ ฯเปฯ ปฐมํ ฌานํ อุปสมฺปชฺช วิหรติ ทุกฺข\-ปฏิปทํ ทนฺธาภิญฺญํ. ตสฺมิํ สมเย ผสฺโส โหติ ฯเปฯ อวิกฺเขโป โหติ ฯเปฯ. อิเม ธมฺมา กุสลา. ตสฺเสว โลกุตฺตรสฺส กุสลสฺส ฌานสฺส กตตฺตา ภาวิตตฺตา วิปากํ วิวิจฺเจว กาเมหิ ฯเปฯ ปฐมํ ฌานํ อุปสมฺปชฺช วิหรติ ทุกฺข\-ปฏิปทํ ทนฺธาภิญฺญํ อปฺปณิหิตํ. ตสฺมิํ สมเย ผสฺโส โหติ ฯเปฯ อญฺญินฺทฺริยํ โหติ ฯเปฯ อวิกฺเขโป โหติ ฯเปฯ. อิเม ธมฺมา อพฺยากตา.}

\proseitem{508}{กตเม ธมฺมา อพฺยากตา. ยสฺมิํ สมเย โลกุตฺตรํ ฌานํ ภาเวติ นิยฺยานิกํ อปจยคามิํ ทิฏฺฐิคตานํ ปหานาย ปฐมาย ภูมิยา ปตฺติยา วิตกฺก\-วิจารานํ วูปสมา ฯเปฯ ทุติยํ ฌานํ ฯเปฯ ตติยํ ฌานํ ฯเปฯ จตุตฺถํ ฌานํ ฯเปฯ ปฐมํ ฌานํ ฯเปฯ ปญฺจมํ ฌานํ อุปสมฺปชฺช วิหรติ ทุกฺข\-ปฏิปทํ ทนฺธา\-ภิญฺญนฺติ กุสลํ ฯเปฯ ทุกฺข\-ปฏิปทํ ทนฺธาภิญฺญํ สุญฺญตนฺติ วิปาโก ฯเปฯ ทุกฺข\-ปฏิปทํ ทนฺธา\-ภิญฺญนฺติ กุสลํ ฯเปฯ ทุกฺข\-ปฏิปทํ ทนฺธาภิญฺญํ อนิมิตฺตนฺติ วิปาโก ฯเปฯ ทุกฺข\-ปฏิปทํ ทนฺธา\-ภิญฺญนฺติ กุสลํ ฯเปฯ ทุกฺข\-ปฏิปทํ ทนฺธาภิญฺญํ อปฺปณิหิตนฺติ วิปาโก. ตสฺมิํ สมเย ผสฺโส โหติ ฯเปฯ อวิกฺเขโป โหติ ฯเปฯ. อิเม ธมฺมา อพฺยากตา.}

\proseitem{509}{กตเม ธมฺมา อพฺยากตา. ยสฺมิํ สมเย โลกุตฺตรํ ฌานํ ภาเวติ นิยฺยานิกํ อปจยคามิํ ทิฏฺฐิคตานํ ปหานาย ปฐมาย ภูมิยา ปตฺติยา วิวิจฺเจว กาเมหิ ฯเปฯ ปฐมํ ฌานํ อุปสมฺปชฺช วิหรติ ทุกฺข\-ปฏิปทํ ขิปฺปาภิญฺญํ ฯเปฯ สุขปฏิปทํ ทนฺธาภิญฺญํ ฯเปฯ สุขปฏิปทํ ขิปฺปาภิญฺญํ ฯเปฯ ทุติยํ ฌานํ ฯเปฯ ตติยํ ฌานํ ฯเปฯ จตุตฺถํ ฌานํ ฯเปฯ ปฐมํ ฌานํ ฯเปฯ ปญฺจมํ ฌานํ อุปสมฺปชฺช วิหรติ สุขปฏิปทํ ขิปฺปา\-ภิญฺญนฺติ กุสลํ ฯเปฯ สุขปฏิปทํ ขิปฺปาภิญฺญํ สุญฺญตนฺติ วิปาโก ฯเปฯ สุขปฏิปทํ ขิปฺปา\-ภิญฺญนฺติ กุสลํ ฯเปฯ สุขปฏิปทํ ขิปฺปาภิญฺญํ อนิมิตฺตนฺติ วิปาโก ฯเปฯ สุขปฏิปทํ ขิปฺปา\-ภิญฺญนฺติ กุสลํ ฯเปฯ \csromanfolio{120}สุขปฏิปทํ ขิปฺปาภิญฺญํ อปฺปณิหิตนฺติ วิปาโก. ตสฺมิํ สมเย ผสฺโส โหติ ฯเปฯ อวิกฺเขโป โหติ ฯเปฯ. อิเม ธมฺมา อพฺยากตา.}

\vspace{\tipitakaparskip}
\csromancenter{สุทฺธิก\-ปฏิปทา.}
\csromanmatikaanchor{mtk.60}
\csromantocmarkref{subsection}{สุทฺธิกสุญฺญต}{mtk.60}
\bookmark[dest={mtk.60},level=2]{สุทฺธิกสุญฺญต}
\csromanheader{2}{สุทฺธิก\-สุญฺญต}
\proseitem{510}{กตเม ธมฺมา อพฺยากตา. ยสฺมิํ สมเย โลกุตฺตรํ ฌานํ ภาเวติ นิยฺยานิกํ อปจยคามิํ ทิฏฺฐิคตานํ ปหานาย ปฐมาย ภูมิยา ปตฺติยา วิวิจฺเจว กาเมหิ ฯเปฯ ปฐมํ ฌานํ อุปสมฺปชฺช วิหรติ สุญฺญตํ. ตสฺมิํ สมเย ผสฺโส โหติ ฯเปฯ อวิกฺเขโป โหติ ฯเปฯ. อิเม ธมฺมา กุสลา. ตสฺเสว โลกุตฺตรสฺส กุสลสฺส ฌานสฺส กตตฺตา ภาวิตตฺตา วิปากํ วิวิจฺเจว กาเมหิ ฯเปฯ ปฐมํ ฌานํ อุปสมฺปชฺช วิหรติ สุญฺญตํ. ตสฺมิํ สมเย ผสฺโส โหติ ฯเปฯ อวิกฺเขโป โหติ ฯเปฯ. อิเม ธมฺมา อพฺยากตา.}

\proseitem{511}{กตเม ธมฺมา อพฺยากตา. ยสฺมิํ สมเย โลกุตฺตรํ ฌานํ ภาเวติ นิยฺยานิกํ อปจยคามิํ ทิฏฺฐิคตานํ ปหานาย ปฐมาย ภูมิยา ปตฺติยา วิวิจฺเจว กาเมหิ ฯเปฯ ปฐมํ ฌานํ อุปสมฺปชฺช วิหรติ สุญฺญตํ. ตสฺมิํ สมเย ผสฺโส โหติ ฯเปฯ อวิกฺเขโป โหติ ฯเปฯ. อิเม ธมฺมา กุสลา ตสฺเสว โลกุตฺตรสฺส กุสลสฺส ฌานสฺส กตตฺตา ภาวิตตฺตา วิปากํ วิวิจฺเจว กาเมหิ ฯเปฯ ปฐมํ ฌานํ อุปสมฺปชฺช วิหรติ อนิมิตฺตํ. ตสฺมิํ สมเย ผสฺโส โหติ ฯเปฯ อวิกฺเขโป โหติ ฯเปฯ. อิเม ธมฺมา อพฺยากตา.}

\proseitem{512}{กตเม ธมฺมา อพฺยากตา. ยสฺมิํ สมเย โลกุตฺตรํ ฌานํ ภาเวติ นิยฺยานิกํ อปจยคามิํ ทิฏฺฐิคตานํ ปหานาย ปฐมาย ภูมิยา ปตฺติยา วิวิจฺเจว กาเมหิ ฯเปฯ ปฐมํ ฌานํ อุปสมฺปชฺช วิหรติ สุญฺญตํ. ตสฺมิํ สมเย ผสฺโส โหติ ฯเปฯ อวิกฺเขโป โหติ ฯเปฯ. อิเม ธมฺมา กุสลา. ตสฺเสว โลกุตฺตรสฺส กุสลสฺส ฌานสฺส กตตฺตา ภาวิตตฺตา วิปากํ วิวิจฺเจว กาเมหิ ฯเปฯ ปฐมํ ฌานํ อุปสมฺปชฺช วิหรติ อปฺปณิหิตํ. ตสฺมิํ สมเย ผสฺโส โหติ ฯเปฯ อวิกฺเขโป โหติ ฯเปฯ. อิเม ธมฺมา อพฺยากตา.}

\chapterheadpageread{1. จิตฺตุปฺปาท\-กณฺฑ}
\proseitem{513}{\csromanfolio{121}กตเม ธมฺมา อพฺยากตา. ยสฺมิํ สมเย โลกุตฺตรํ ฌานํ ภาเวติ นิยฺยานิกํ อปจยคามิํ ทิฏฺฐิคตานํ ปหานาย ปฐมาย ภูมิยา ปตฺติยา วิตกฺก\-วิจารานํ วูปสมา ฯเปฯ ทุติยํ ฌานํ ฯเปฯ ตติยํ ฌานํ ฯเปฯ จตุตฺถํ ฌานํ ฯเปฯ ปฐมํ ฌานํ ฯเปฯ ปญฺจมํ ฌานํ อุปสมฺปชฺช วิหรติ สุญฺญตนฺติ กุสลํ ฯเปฯ สุญฺญตนฺติ วิปาโก ฯเปฯ สุญฺญตนฺติ กุสลํ ฯเปฯ อนิมิตฺตนฺติ วิปาโก ฯเปฯ สุญฺญตนฺติ กุสลํ ฯเปฯ อปฺปณิหิตนฺติ วิปาโก. ตสฺมิํ สมเย ผสฺโส โหติ ฯเปฯ อวิกฺเขโป โหติ ฯเปฯ. อิเม ธมฺมา อพฺยากตา.}

\vspace{\tipitakaparskip}
\csromancenter{สุทฺธิก\-สุญฺญตํ.}
\csromanmatikaanchor{mtk.61}
\csromantocmarkref{subsection}{สุญฺญตปฏิปทา}{mtk.61}
\bookmark[dest={mtk.61},level=2]{สุญฺญตปฏิปทา}
\csromanheader{2}{สุญฺญต\-ปฏิปทา}
\proseitem{514}{กตเม ธมฺมา อพฺยากตา. ยสฺมิํ สมเย โลกุตฺตรํ ฌานํ ภาเวติ นิยฺยานิกํ อปจยคามิํ ทิฏฺฐิคตานํ ปหานาย ปฐมาย ภูมิยา ปตฺติยา วิวิจฺเจว กาเมหิ ฯเปฯ ปฐมํ ฌานํ อุปสมฺปชฺช วิหรติ ทุกฺข\-ปฏิปทํ ทนฺธาภิญฺญํ สุญฺญตํ. ตสฺมิํ สมเย ผสฺโส โหติ ฯเปฯ อวิกฺเขโป โหติ ฯเปฯ. อิเม ธมฺมา กุสลา. ตสฺเสว โลกุตฺตรสฺส กุสลสฺส ฌานสฺส กตตฺตา ภาวิตตฺตา วิปากํ วิวิจฺเจว กาเมหิ ฯเปฯ ปฐมํ ฌานํ อุปสมฺปชฺช วิหรติ ทุกฺข\-ปฏิปทํ ทนฺธาภิญฺญํ สุญฺญตํ. ตสฺมิํ สมเย ผสฺโส โหติ ฯเปฯ อวิกฺเขโป โหติ ฯเปฯ. อิเม ธมฺมา อพฺยากตา.}

\proseitem{515}{กตเม ธมฺมา อพฺยากตา. ยสฺมิํ สมเย โลกุตฺตรํ ฌานํ ภาเวติ นิยฺยานิกํ อปจยคามิํ ทิฏฺฐิคตานํ ปหานาย ปฐมาย ภูมิยา ปตฺติยา วิวิจฺเจว กาเมหิ ฯเปฯ ปฐมํ ฌานํ อุปสมฺปชฺช วิหรติ ทุกฺข\-ปฏิปทํ ทนฺธาภิญฺญํ สุญฺญตํ. ตสฺมิํ สมเย ผสฺโส โหติ ฯเปฯ อวิกฺเขโป โหติ ฯเปฯ. อิเม ธมฺมา กุสลา. ตสฺเสว โลกุตฺตรสฺส กุสลสฺส ฌานสฺส กตตฺตา ภาวิตตฺตา วิปากํ วิวิจฺเจว กาเมหิ ฯเปฯ ปฐมํ ฌานํ อุปสมฺปชฺช วิหรติ ทุกฺข\-ปฏิปทํ ทนฺธาภิญฺญํ อนิมิตฺตํ. ตสฺมิํ สมเย ผสฺโส โหติ ฯเปฯ อวิกฺเขโป โหติ ฯเปฯ. อิเม ธมฺมา อพฺยากตา.}

\proseitem{516}{กตเม ธมฺมา อพฺยากตา. ยสฺมิํ สมเย โลกุตฺตรํ ฌานํ ภาเวติ นิยฺยานิกํ อปจยคามิํ ทิฏฺฐิคตานํ ปหานาย ปฐมาย ภูมิยา \csromanfolio{122}ปตฺติยา วิวิจฺเจว กาเมหิ ฯเปฯ ปฐมํ ฌานํ อุปสมฺปชฺช วิหรติ ทุกฺข\-ปฏิปทํ ทนฺธาภิญฺญํ สุญฺญตํ. ตสฺมิํ สมเย ผสฺโส โหติ ฯเปฯ อวิกฺเขโป โหติ ฯเปฯ. อิเม ธมฺมา กุสลา. ตสฺเสว โลกุตฺตรสฺส กุสลสฺส ฌานสฺส กตตฺตา ภาวิตตฺตา วิปากํ วิวิจฺเจว กาเมหิ ฯเปฯ ปฐมํ ฌานํ อุปสมฺปชฺช วิหรติ ทุกฺข\-ปฏิปทํ ทนฺธาภิญฺญํ อปฺปณิหิตํ. ตสฺมิํ สมเย ผสฺโส โหติ ฯเปฯ อวิกฺเขโป โหติ ฯเปฯ. อิเม ธมฺมา อพฺยากตา.}

\proseitem{517}{กตเม ธมฺมา อพฺยากตา. ยสฺมิํ สมเย โลกุตฺตรํ ฌานํ ภาเวติ นิยฺยานิกํ อปจยคามิํ ทิฏฺฐิคตานํ ปหานาย ปฐมาย ภูมิยา ปตฺติยา วิตกฺก\-วิจารานํ วูปสมา ฯเปฯ ทุติยํ ฌานํ ฯเปฯ ตติยํ ฌานํ ฯเปฯ จตุตฺถํ ฌานํ ฯเปฯ ปฐมํ ฌานํ ฯเปฯ ปญฺจมํ ฌานํ อุปสมฺปชฺช วิหรติ ทุกฺข\-ปฏิปทํ ทนฺธาภิญฺญํ สุญฺญตนฺติ กุสลํ ฯเปฯ ทุกฺข\-ปฏิปทํ ทนฺธาภิญฺญํ สุญฺญตนฺติ วิปาโก ฯเปฯ ทุกฺข\-ปฏิปทํ ทนฺธาภิญฺญํ สุญฺญตนฺติ กุสลํ ฯเปฯ ทุกฺข\-ปฏิปทํ ทนฺธาภิญฺญํ อนิมิตฺตนฺติ วิปาโก ฯเปฯ ทุกฺข\-ปฏิปทํ ทนฺธาภิญฺญํ สุญฺญตนฺติ กุสลํ ฯเปฯ ทุกฺข\-ปฏิปทํ ทนฺธาภิญฺญํ อปฺปณิหิตนฺติ วิปาโก. ตสฺมิํ สมเย ผสฺโส โหติ ฯเปฯ อวิกฺเขโป โหติ ฯเปฯ. อิเม ธมฺมา อพฺยากตา.}

\proseitem{518}{กตเม ธมฺมา อพฺยากตา. ยสฺมิํ สมเย โลกุตฺตรํ ฌานํ ภาเวติ นิยฺยานิกํ อปจยคามิํ ทิฏฺฐิคตานํ ปหานาย ปฐมาย ภูมิยา ปตฺติยา วิวิจฺเจว กาเมหิ ฯเปฯ ปฐมํ ฌานํ อุปสมฺปชฺช วิหรติ ทุกฺข\-ปฏิปทํ ขิปฺปาภิญฺญํ สุญฺญตํ ฯเปฯ สุขปฏิปทํ ทนฺธาภิญฺญํ สุญฺญตํ ฯเปฯ สุขปฏิปทํ ขิปฺปาภิญฺญํ สุญฺญตํ ฯเปฯ ทุติยํ ฌานํ ฯเปฯ ตติยํ ฌานํ ฯเปฯ จตุตฺถํ ฌานํ ฯเปฯ ปฐมํ ฌานํ ฯเปฯ ปญฺจมํ ฌานํ อุปสมฺปชฺช วิหรติ สุขปฏิปทํ ขิปฺปาภิญฺญํ สุญฺญตนฺติ กุสลํ ฯเปฯ สุขปฏิปทํ ขิปฺปาภิญฺญํ สุญฺญตนฺติ วิปาโก ฯเปฯ สุขปฏิปทํ ขิปฺปาภิญฺญํ สุญฺญตนฺติ กุสลํ ฯเปฯ สุขปฏิปทํ ขิปฺปาภิญฺญํ อนิมิตฺตนฺติ วิปาโก ฯเปฯ สุขปฏิปทํ ขิปฺปาภิญฺญํ สุญฺญตนฺติ กุสลํ ฯเปฯ สุขปฏิปทํ ขิปฺปาภิญฺญํ อปฺปณิหิตนฺติ วิปาโก. ตสฺมิํ สมเย ผสฺโส โหติ ฯเปฯ อวิกฺเขโป โหติ ฯเปฯ. อิเม ธมฺมา อพฺยากตา.}

\vspace{\tipitakaparskip}
\csromancenter{สุญฺญต\-ปฏิปทา.}
\csromanmatikaanchor{mtk.62}
\csromantocmarkref{subsection}{สุทฺธิกอปฺปณิหิต}{mtk.62}
\bookmark[dest={mtk.62},level=2]{สุทฺธิกอปฺปณิหิต}
\csromanheader{2}{1. จิตฺตุปฺปาท\-กณฺฑ สุทฺธิก\-อปฺปณิหิต}
\proseitem{519}{\csromanfolio{123}กตเม ธมฺมา อพฺยากตา. ยสฺมิํ สมเย โลกุตฺตรํ ฌานํ ภาเวติ นิยฺยานิกํ อปจยคามิํ ทิฏฺฐิคตานํ ปหานาย ปฐมาย ภูมิยา ปตฺติยา วิวิจฺเจว กาเมหิ ฯเปฯ ปฐมํ ฌานํ อุปสมฺปชฺช วิหรติ อปฺปณิหิตํ. ตสฺมิํ สมเย ผสฺโส โหติ ฯเปฯ อวิกฺเขโป โหติ ฯเปฯ. อิเม ธมฺมา กุสลา. ตสฺเสว โลกุตฺตรสฺส กุสลสฺส ฌานสฺส กตตฺตา ภาวิตตฺตา วิปากํ วิวิจฺเจว กาเมหิ ฯเปฯ ปฐมํ ฌานํ อุปสมฺปชฺช วิหรติ อปฺปณิหิตํ. ตสฺมิํ สมเย ผสฺโส โหติ ฯเปฯ อวิกฺเขโป โหติ ฯเปฯ. อิเม ธมฺมา อพฺยากตา.}

\proseitem{520}{กตเม ธมฺมา อพฺยากตา. ยสฺมิํ สมเย โลกุตฺตรํ ฌานํ ภาเวติ นิยฺยานิกํ อปจยคามิํ ทิฏฺฐิคตานํ ปหานาย ปฐมาย ภูมิยา ปตฺติยา วิวิจฺเจว กาเมหิ ฯเปฯ ปฐมํ ฌานํ อุปสมฺปชฺช วิหรติ อปฺปณิหิตํ. ตสฺมิํ สมเย ผสฺโส โหติ ฯเปฯ อวิกฺเขโป โหติ ฯเปฯ. อิเม ธมฺมา กุสลา. ตสฺเสว โลกุตฺตรสฺส กุสลสฺส ฌานสฺส กตตฺตา ภาวิตตฺตา วิปากํ วิวิจฺเจว กาเมหิ ฯเปฯ ปฐมํ ฌานํ อุปสมฺปชฺช วิหรติ อนิมิตฺตํ. ตสฺมิํ สมเย ผสฺโส โหติ ฯเปฯ อวิกฺเขโป โหติ ฯเปฯ. อิเม ธมฺมา อพฺยากตา.}

\proseitem{521}{กตเม ธมฺมา อพฺยากตา. ยสฺมิํ สมเย โลกุตฺตรํ ฌานํ ภาเวติ นิยฺยานิกํ อปจยคามิํ ทิฏฺฐิคตานํ ปหานาย ปฐมาย ภูมิยา ปตฺติยา วิวิจฺเจว กาเมหิ ฯเปฯ ปฐมํ ฌานํ อุปสมฺปชฺช วิหรติ อปฺปณิหิตํ. ตสฺมิํ สมเย ผสฺโส โหติ ฯเปฯ อวิกฺเขโป โหติ ฯเปฯ. อิเม ธมฺมา กุสลา. ตสฺเสว โลกุตฺตรสฺส กุสลสฺส ฌานสฺส กตตฺตา ภาวิตตฺตา วิปากํ วิวิจฺเจว กาเมหิ ฯเปฯ ปฐมํ ฌานํ อุปสมฺปชฺช วิหรติ สุญฺญตํ. ตสฺมิํ สมเย ผสฺโส โหติ ฯเปฯ อวิกฺเขโป โหติ ฯเปฯ. อิเม ธมฺมา อพฺยากตา.}

\proseitem{522}{กตเม ธมฺมา อพฺยากตา. ยสฺมิํ สมเย โลกุตฺตรํ ฌานํ ภาเวติ นิยฺยานิกํ อปจยคามิํ ทิฏฺฐิคตานํ ปหานาย ปฐมาย ภูมิยา ปตฺติยา วิตกฺก\-วิจารานํ วูปสมา ฯเปฯ ทุติยํ ฌานํ ฯเปฯ ตติยํ ฌานํ ฯเปฯ จตุตฺถํ ฌานํ ฯเปฯ ปฐมํ ฌานํ ฯเปฯ ปญฺจมํ ฌานํ อุปสมฺปชฺช วิหรติ อปฺปณิหิตนฺติ กุสลํ ฯเปฯ อปฺปณิหิตนฺติ วิปาโก ฯเปฯ อปฺปณิหิตนฺติ \csromanfolio{124}กุสลํ ฯเปฯ อนิมิตฺตนฺติ วิปาโก ฯเปฯ อปฺปณิหิตนฺติ กุสลํ ฯเปฯ สุญฺญตนฺติ วิปาโก. ตสฺมิํ สมเย ผสฺโส โหติ ฯเปฯ อวิกฺเขโป โหติ ฯเปฯ. อิเม ธมฺมา อพฺยากตา.}

\prose{สุทฺธิก\-อปฺปณิหิตํ.}

\csromanmatikaanchor{mtk.63}
\csromantocmarkref{subsection}{อปฺปณิหิตปฏิปทา}{mtk.63}
\bookmark[dest={mtk.63},level=2]{อปฺปณิหิตปฏิปทา}
\csromanheader{2}{อปฺปณิหิต\-ปฏิปทา}
\proseitem{523}{กตเม ธมฺมา อพฺยากตา. ยสฺมิํ สมเย โลกุตฺตรํ ฌานํ ภาเวติ นิยฺยานิกํ อปจยคามิํ ทิฏฺฐิคตานํ ปหานาย ปฐมาย ภูมิยา ปตฺติยา วิวิจฺเจว กาเมหิ ฯเปฯ ปฐมํ ฌานํ อุปสมฺปชฺช วิหรติ ทุกฺข\-ปฏิปทํ ทนฺธาภิญฺญํ อปฺปณิหิตํ. ตสฺมิํ สมเย ผสฺโส โหติ ฯเปฯ อวิกฺเขโป โหติ ฯเปฯ. อิเม ธมฺมา กุสลา. ตสฺเสว โลกุตฺตรสฺส กุสลสฺส ฌานสฺส กตตฺตา ภาวิตตฺตา วิปากํ วิวิจฺเจว กาเมหิ ฯเปฯ ปฐมํ ฌานํ อุปสมฺปชฺช วิหรติ ทุกฺข\-ปฏิปทํ ทนฺธาภิญฺญํ อปฺปณิหิตํ. ตสฺมิํ สมเย ผสฺโส โหติ ฯเปฯ อวิกฺเขโป โหติ ฯเปฯ. อิเม ธมฺมา อพฺยากตา.}

\proseitem{524}{กตเม ธมฺมา อพฺยากตา. ยสฺมิํ สมเย โลกุตฺตรํ ฌานํ ภาเวติ นิยฺยานิกํ อปจยคามิํ ทิฏฺฐิคตานํ ปหานาย ปฐมาย ภูมิยา ปตฺติยา วิวิจฺเจว กาเมหิ ฯเปฯ ปฐมํ ฌานํ อุปสมฺปชฺช วิหรติ ทุกฺข\-ปฏิปทํ ทนฺธาภิญฺญํ อปฺปณิหิตํ. ตสฺมิํ สมเย ผสฺโส โหติ ฯเปฯ อวิกฺเขโป โหติ ฯเปฯ. อิเม ธมฺมา กุสลา. ตสฺเสว โลกุตฺตรสฺส กุสลสฺส ฌานสฺส กตตฺตา ภาวิตตฺตา วิปากํ วิวิจฺเจว กาเมหิ ฯเปฯ ปฐมํ ฌานํ อุปสมฺปชฺช วิหรติ ทุกฺข\-ปฏิปทํ ทนฺธาภิญฺญํ อนิมิตฺตํ. ตสฺมิํ สมเย ผสฺโส โหติ ฯเปฯ อวิกฺเขโป โหติ ฯเปฯ. อิเม ธมฺมา อพฺยากตา.}

\csromanmatikaanchor{mtk.64}
\csromantocmarkref{subsection}{วีสติมหานย}{mtk.64}
\bookmark[dest={mtk.64},level=2]{วีสติมหานย}
\proseitem{525}{กตเม ธมฺมา อพฺยากตา. ยสฺมิํ สมเย โลกุตฺตรํ ฌานํ ภาเวติ นิยฺยานิกํ อปจยคามิํ ทิฏฺฐิคตานํ ปหานาย ปฐมาย ภูมิยา ปตฺติยา วิวิจฺเจว กาเมหิ ฯเปฯ ปฐมํ ฌานํ อุปสมฺปชฺช วิหรติ ทุกฺข\-ปฏิปทํ ทนฺธาภิญฺญํ อปฺปณิหิตํ. ตสฺมิํ สมเย ผสฺโส โหติ ฯเปฯ อวิกฺเขโป โหติ ฯเปฯ. อิเม ธมฺมา กุสลา. ตสฺเสว โลกุตฺตรสฺส \csromanfolio{125}กุสลสฺส ฌานสฺส กตตฺตา ภาวิตตฺตา วิปากํ วิวิจฺเจว กาเมหิ ฯเปฯ ปฐมํ ฌานํ อุปสมฺปชฺช วิหรติ ทุกฺข\-ปฏิปทํ ทนฺธาภิญฺญํ สุญฺญตํ. ตสฺมิํ สมเย ผสฺโส โหติ ฯเปฯ อวิกฺเขโป โหติ ฯเปฯ. อิเม ธมฺมา อพฺยากตา.}

\proseitem{526}{กตเม ธมฺมา อพฺยากตา. ยสฺมิํ สมเย โลกุตฺตรํ ฌานํ ภาเวติ นิยฺยานิกํ อปจยคามิํ ทิฏฺฐิคตานํ ปหานาย ปฐมาย ภูมิยา ปตฺติยา วิตกฺก\-วิจารานํ วูปสมา ฯเปฯ ทุติยํ ฌานํ ฯเปฯ ตติยํ ฌานํ ฯเปฯ จตุตฺถํ ฌานํ ฯเปฯ ปฐมํ ฌานํ ฯเปฯ ปญฺจมํ ฌานํ อุปสมฺปชฺช วิหรติ ทุกฺข\-ปฏิปทํ ทนฺธาภิญฺญํ อปฺปณิหิตนฺติ กุสลํ ฯเปฯ ทุกฺข\-ปฏิปทํ ทนฺธาภิญฺญํ อปฺปณิหิตนฺติ วิปาโก ฯเปฯ ทุกฺข\-ปฏิปทํ ทนฺธาภิญฺญํ อปฺปณิหิตนฺติ กุสลํ ฯเปฯ ทุกฺข\-ปฏิปทํ ทนฺธาภิญฺญํ อนิมิตฺตนฺติ วิปาโก ฯเปฯ ทุกฺข\-ปฏิปทํ ทนฺธาภิญฺญํ อปฺปณิหิตนฺติ กุสลํ ฯเปฯ ทุกฺข\-ปฏิปทํ ทนฺธาภิญฺญํ สุญฺญตนฺติ วิปาโก. ตสฺมิํ สมเย ผสฺโส โหติ ฯเปฯ อวิกฺเขโป โหติ ฯเปฯ. อิเม ธมฺมา อพฺยากตา.}

\proseitem{527}{กตเม ธมฺมา อพฺยากตา. ยสฺมิํ สมเย โลกุตฺตรํ ฌานํ ภาเวติ นิยฺยานิกํ อปจยคามิํ ทิฏฺฐิคตานํ ปหานาย ปฐมาย ภูมิยา ปตฺติยา วิวิจฺเจว กาเมหิ ฯเปฯ ปฐมํ ฌานํ อุปสมฺปชฺช วิหรติ ทุกฺข\-ปฏิปทํ ขิปฺปาภิญฺญํ อปฺปณิหิตํ ฯเปฯ สุขปฏิปทํ ทนฺธาภิญฺญํ อปฺปณิหิตํ ฯเปฯ สุขปฏิปทํ ขิปฺปาภิญฺญํ อปฺปณิหิตํ ฯเปฯ ทุติยํ ฌานํ ฯเปฯ ตติยํ ฌานํ ฯเปฯ จตุตฺถํ ฌานํ ฯเปฯ ปฐมํ ฌานํ ฯเปฯ ปญฺจมํ ฌานํ อุปสมฺปชฺช วิหรติ สุขปฏิปทํ ขิปฺปาภิญฺญํ อปฺปณิหิตนฺติ กุสลํ ฯเปฯ สุขปฏิปทํ ขิปฺปาภิญฺญํ อปฺปณิหิตนฺติ วิปาโก ฯเปฯ สุขปฏิปทํ ขิปฺปาภิญฺญํ อปฺปณิหิตนฺติ กุสลํ ฯเปฯ สุขปฏิปทํ ขิปฺปาภิญฺญํ อนิมิตฺตนฺติ วิปาโก ฯเปฯ สุขปฏิปทํ ขิปฺปาภิญฺญํ อปฺปณิหิตนฺติ กุสลํ ฯเปฯ สุขปฏิปทํ ขิปฺปาภิญฺญํ สุญฺญตนฺติ วิปาโก. ตสฺมิํ สมเย ผสฺโส โหติ ฯเปฯ อวิกฺเขโป โหติ ฯเปฯ. อิเม ธมฺมา อพฺยกตา.}

\vspace{\tipitakaparskip}
\csromancenter{อปฺปณิหิต\-ปฏิปทา.}
\csromanheader{2}{\textbf{วีสติ มหานย}}
\proseitem{528}{กตเม ธมฺมา อพฺยากตา. ยสฺมิํ สมเย โลกุตฺตรํ มคฺคํ ภาเวติ ฯเปฯ โลกุตฺตรํ สติปฏฺฐานํ ภาเวติ ฯเปฯ โลกุตฺตรํ สมฺมปฺปธานํ \csromanfolio{126}ภาเวติ ฯเปฯ โลกุตฺตรํ อิทฺธิปาทํ ภาเวหิ ฯเปฯ โลกุตฺตรํ อินฺทฺริยํ ภาเวติ ฯเปฯ โลกุตฺตรํ พลํ ภาเวติ ฯเปฯ โลกุตฺตรํ โพชฺฌงฺคํ ภาเวติ ฯเปฯ โลกุตฺตรํ สจฺจํ ภาเวติ ฯเปฯ โลกุตฺตรํ สมถํ ภาเวติ ฯเปฯ โลกุตฺตรํ ธมฺมํ ภาเวติ ฯเปฯ โลกุตฺตรํ ขนฺธํ ภาเวติ ฯเปฯ โลกุตฺตรํ อายตนํ ภาเวติ ฯเปฯ โลกุตฺตรํ ธาตุํ ภาเวติ ฯเปฯ โลกุตฺตรํ อาหารํ ภาเวติ ฯเปฯ โลกุตฺตรํ ผสฺสํ ภาเวติ ฯเปฯ โลกุตฺตรํ เวทนํ ภาเวติ ฯเปฯ โลกุตฺตรํ สญฺญํ ภาเวติ ฯเปฯ โลกุตฺตรํ เจตนํ ภาเวติ ฯเปฯ โลกุตฺตรํ จิตฺตํ ภาเวติ นิยฺยานิกํ อปจยคามิํ ทิฏฺฐิคตานํ ปหานาย ปฐมาย ภูมิยา ปตฺติยา วิวิจฺเจว กาเมหิ ฯเปฯ ปฐมํ ฌานํ อุปสมฺปชฺช วิหรติ ทุกฺข\-ปฏิปทํ ทนฺธาภิญฺญํ. ตสฺมิํ สมเย ผสฺโส โหติ ฯเปฯ อวิกฺเขโป โหติ ฯเปฯ. อิเม ธมฺมา กุสลา. ตสฺเสว โลกุตฺตรสฺส กุสลสฺส ฌานสฺส กตตฺตา ภาวิตตฺตา วิปากํ วิวิจฺเจว กาเมหิ ฯเปฯ ปฐมํ ฌานํ อุปสมฺปชฺช วิหรติ ทุกฺข\-ปฏิปทํ ทนฺธาภิญฺญํ สุญฺญตํ ฯเปฯ อนิมิตฺตํ ฯเปฯ อปฺปณิหิตํ. ตสฺมิํ สมเย ผสฺโส โหติ ฯเปฯ อวิกฺเขโป โหติ ฯเปฯ. อิเม ธมฺมา อพฺยากตา.}

\vspace{\tipitakaparskip}
\csromancenter{วีสติ มหานยา.}
\csromanmatikaanchor{mtk.65}
\csromantocmarkref{subsection}{ฉนฺทาธิปเตยฺยสุทฺธิกปฏิปทา}{mtk.65}
\bookmark[dest={mtk.65},level=2]{ฉนฺทาธิปเตยฺยสุทฺธิกปฏิปทา}
\csromanheader{2}{ฉนฺทา\-ธิปเตยฺย\-สุทฺธิก\-ปฏิปทา}
\proseitem{529}{กตเม ธมฺมา อพฺยากตา. ยสฺมิํ สมเย โลกุตฺตรํ ฌานํ ภาเวติ นิยฺยานิกํ อปจยคามิํ ทิฏฺฐิคตานํ ปหานาย ปฐมาย ภูมิยา ปตฺติยา วิวิจฺเจว กาเมหิ ฯเปฯ ปฐมํ ฌานํ อุปสมฺปชฺช วิหรติ ทุกฺข\-ปฏิปทํ ทนฺธาภิญฺญํ ฉนฺทา\-ธิปเตยฺยํ. ตสฺมิํ สมเย ผสฺโส โหติ ฯเปฯ อวิกฺเขโป โหติ ฯเปฯ. อิเม ธมฺมา กุสลา. ตสฺเสว โลกุตฺตรสฺส กุสลสฺส ฌานสฺส กตตฺตา ภาวิตตฺตา วิปากํ วิวิจฺเจว กาเมหิ ฯเปฯ ปฐมํ ฌานํ อุปสมฺปชฺช วิหรติ ทุกฺข\-ปฏิปทํ ทนฺธาภิญฺญํ สุญฺญตํ ฉนฺทา\-ธิปเตยฺยํ. ตสฺมิํ สมเย ผสฺโส โหติ ฯเปฯ อวิกฺเขโป โหติ ฯเปฯ. อิเม ธมฺมา อพฺยากตา.}

\proseitem{530}{กตเม ธมฺมา อพฺยากตา. ยสฺมิํ สมเย โลกุตฺตรํ ฌานํ ภาเวติ นิยฺยานิกํ อปจยคามิํ ทิฏฺฐิคตานํ ปหานาย ปฐมาย \csromanfolio{127}ภูมิยา ปตฺติยา วิวิจฺเจว กาเมหิ ฯเปฯ ปฐมํ ฌานํ อุปสมฺปชฺช วิหรติ ทุกฺข\-ปฏิปทํ ทนฺธาภิญฺญํ ฉนฺทา\-ธิปเตยฺยํ. ตสฺมิํ สมเย ผสฺโส โหติ ฯเปฯ อวิกฺเขโป โหติ ฯเปฯ. อิเม ธมฺมา กุสลา. ตสฺเสว โลกุตฺตรสฺส กุสลสฺส ฌานสฺส กตตฺตา ภาวิตตฺตา วิปากํ วิวิจฺเจว กาเมหิ ฯเปฯ ปฐมํ ฌานํ อุปสมฺปชฺช วิหรติ ทุกฺข\-ปฏิปทํ ทนฺธาภิญฺญํ อนิมิตฺตํ ฉนฺทา\-ธิปเตยฺยํ. ตสฺมิํ สมเย ผสฺโส โหติ ฯเปฯ อวิกฺเขโป โหติ ฯเปฯ. อิเม ธมฺมา อพฺยากตา.}

\proseitem{531}{กตเม ธมฺมา อพฺยากตา. ยสฺมิํ สมเย โลกุตฺตรํ ฌานํ ภาเวติ นิยฺยานิกํ อปจยคามิํ ทิฏฺฐิคตานํ ปหานาย ปฐมาย ภูมิยา ปตฺติยา วิวิจฺเจว กาเมหิ ฯเปฯ ปฐมํ ฌานํ อุปสมฺปชฺช วิหรติ ทุกฺข\-ปฏิปทํ ทนฺธาภิญฺญํ ฉนฺทา\-ธิปเตยฺยํ. ตสฺมิํ สมเย ผสฺโส โหติ ฯเปฯ อวิกฺเขโป โหติ ฯเปฯ. อิเม ธมฺมา กุสลา. ตสฺเสว โลกุตฺตรสฺส กุสลสฺส ฌานสฺส กตตฺตา ภาวิตตฺตา วิปากํ วิวิจฺเจว กาเมหิ ฯเปฯ ปฐมํ ฌานํ อุปสมฺปชฺช วิหรติ ทุกฺข\-ปฏิปทํ ทนฺธาภิญฺญํ อปฺปณิหิตํ ฉนฺทา\-ธิปเตยฺยํ. ตสฺมิํ สมเย ผสฺโส โหติ ฯเปฯ อวิกฺเขโป โหติ ฯเปฯ. อิเม ธมฺมา อพฺยากตา.}

\proseitem{532}{กตเม ธมฺมา อพฺยากตา. ยสฺมิํ สมเย โลกุตฺตรํ ฌานํ ภาเวติ นิยฺยานิกํ อปจยคามิํ ทิฏฺฐิคตานํ ปหานาย ปฐมาย ภูมิยา ปตฺติยา วิตกฺก\-วิจารานํ วูปสมา ฯเปฯ ทุติยํ ฌานํ ฯเปฯ ตติยํ ฌานํ ฯเปฯ จตุตฺถํ ฌานํ ฯเปฯ ปฐมํ ฌานํ ฯเปฯ ปญฺจมํ ฌานํ อุปสมฺปชฺช วิหรติ ทุกฺข\-ปฏิปทํ ทนฺธาภิญฺญํ ฉนฺทาธิปเตยฺยนฺติ กุสลํ ฯเปฯ ทุกฺข\-ปฏิปทํ ทนฺธาภิญฺญํ สุญฺญตํ ฉนฺทาธิปเตยฺยนฺติ วิปาโก ฯเปฯ ทุกฺข\-ปฏิปทํ ทนฺธาภิญฺญํ ฉนฺทาธิปเตยฺยนฺติ กุสลํ ฯเปฯ ทุกฺข\-ปฏิปทํ ทนฺธาภิญฺญํ อนิมิตฺตํ ฉนฺทาธิปเตยฺยนฺติ วิปาโก ฯเปฯ ทุกฺข\-ปฏิปทํ ทนฺธาภิญฺญํ ฉนฺทาธิปเตยฺยนฺติ กุสลํ ฯเปฯ ทุกฺข\-ปฏิปทํ ทนฺธาภิญฺญํ อปฺปณิหิตํ ฉนฺทาธิปเตยฺยนฺติ วิปาโก. ตสฺมิํ สมเย ผสฺโส โหติ ฯเปฯ อวิกฺเขโป โหติ ฯเปฯ. อิเม ธมฺมา อพฺยากตา.}

\proseitem{533}{กตเม ธมฺมา อพฺยากตา. ยสฺมิํ สมเย โลกุตฺตรํ ฌานํ ภาเวติ นิยฺยานิกํ อปจยคามิํ ทิฏฺฐิคตานํ ปหานาย ปฐมาย ภูมิยา ปตฺติยา วิวิจฺเจว กาเมหิ ฯเปฯ ปฐมํ ฌานํ อุปสมฺปชฺช วิหรติ \csromanfolio{128}ทุกฺข\-ปฏิปทํ ขิปฺปาภิญฺญํ ฉนฺทา\-ธิปเตยฺยํ ฯเปฯ สุขปฏิปทํ ทนฺธาภิญฺญํ ฉนฺทา\-ธิปเตยฺยํ ฯเปฯ สุขปฏิปทํ ขิปฺปาภิญฺญํ ฉนฺทา\-ธิปเตยฺยํ ฯเปฯ ทุติยํ ฌานํ ฯเปฯ ตติยํ ฌานํ ฯเปฯ จตุตฺถํ ฌานํ ฯเปฯ ปฐมํ ฌานํ ฯเปฯ ปญฺจมํ ฌานํ อุปสมฺปชฺช วิหรติ สุขปฏิปทํ ขิปฺปาภิญฺญํ ฉนฺทาธิปเตยฺยนฺติ กุสลํ ฯเปฯ สุขปฏิปทํ ขิปฺปาภิญฺญํ สุญฺญตฺตํ ฉนฺทาธิปเตยฺยนฺติ วิปาโก ฯเปฯ สุขปฏิปทํ ขิปฺปาภิญฺญํ ฉนฺทาธิปเตยฺยนฺติ กุสลํ ฯเปฯ สุขปฏิปทํ ขิปฺปาภิญฺญํ อนิมิตฺตํ ฉนฺทาธิปเตยฺยนฺติ วิปาโก ฯเปฯ สุขปฏิปทํ ขิปฺปาภิญฺญํ ฉนฺทาธิปเตยฺยนฺติ กุสลํ ฯเปฯ สุขปฏิปทํ ขิปฺปาภิญฺญํ อปฺปณิหิตํ ฉนฺทาธิปเตยฺยนฺติ วิปาโก. ตสฺมิํ สมเย ผสฺโส โหติ ฯเปฯ อวิกฺเขโป โหติ ฯเปฯ. อิเม ธมฺมา อพฺยากตา.}

\vspace{\tipitakaparskip}
\csromancenter{ฉนฺทา\-ธิปเตยฺย\-สุทฺธิก\-ปฏิปทา.}
\csromanmatikaanchor{mtk.66}
\csromantocmarkref{subsection}{ฉนฺทาธิปเตยฺยสุทฺธิกสุญฺญต}{mtk.66}
\bookmark[dest={mtk.66},level=2]{ฉนฺทาธิปเตยฺยสุทฺธิกสุญฺญต}
\csromanheader{2}{ฉนฺทา\-ธิปเตยฺย\-สุทฺธิก\-สุญฺญต}
\proseitem{534}{กตเม ธมฺมา อพฺยากตา. ยสฺมิํ สมเย โลกุตฺตรํ ฌานํ ภาเวติ นิยฺยานิกํ อปจยคามิํ ทิฏฺฐิคตานํ ปหานาย ปฐมาย ภูมิยา ปตฺติยา วิวิจฺเจว กาเมหิ ฯเปฯ ปฐมํ ฌานํ อุปสมฺปชฺช วิหรติ สุญฺญตํ ฉนฺทา\-ธิปเตยฺยํ. ตสฺมิํ สมเย ผสฺโส โหติ ฯเปฯ อวิกฺเขโป โหติ ฯเปฯ. อิเม ธมฺมา กุสลา. ตสฺเสว โลกุตฺตรสฺส กุสลสฺส ฌานสฺส กตตฺตา ภาวิตตฺตา วิปากํ วิวิจฺเจว กาเมหิ ฯเปฯ ปฐมํ ฌานํ อุปสมฺปชฺช วิหรติ สุญฺญตํ ฉนฺทา\-ธิปเตยฺยํ. ตสฺมิํ สมเย ผสฺโส โหติ ฯเปฯ อวิกฺเขโป โหติ ฯเปฯ. อิเม ธมฺมา อพฺยากตา.}

\proseitem{535}{กตเม ธมฺมา อพฺยากตา. ยสฺมิํ สมเย โลกุตฺตรํ ฌานํ ภาเวติ นิยฺยานิกํ อปจยคามิํ ทิฏฺฐิคตานํ ปหานาย ปฐมาย ภูมิยา ปตฺติยา วิวิจฺเจว กาเมหิ ฯเปฯ ปฐมํ ฌานํ อุปสมฺปชฺช วิหรติ สุญฺญตํ ฉนฺทา\-ธิปเตยฺยํ. ตสฺมิํ สมเย ผสฺโส โหติ ฯเปฯ อวิกฺเขโป โหติ ฯเปฯ. อิเม ธมฺมา กุสลา. ตสฺเสว โลกุตฺตรสฺส กุสลสฺส ฌานสฺส กตตฺตา ภาวิตตฺตา วิปากํ วิวิจฺเจว กาเมหิ ฯเปฯ ปฐมํ ฌานํ อุปสมฺปชฺช วิหรติ อนิมิตฺตํ ฉนฺทา\-ธิปเตยฺยํ. ตสฺมิํ สมเย ผสฺโส โหติ ฯเปฯ อวิกฺเขโป โหติ ฯเปฯ. อิเม ธมฺมา อพฺยากตา.}

\chapterheadpageread{1. จิตฺตุปฺปาท\-กณฺฑ}
\proseitem{536}{\csromanfolio{129}กตเม ธมฺมา อพฺยากตา. ยสฺมิํ สมเย โลกุตฺตรํ ฌานํ ภาเวติ นิยฺยานิกํ อปจยคามิํ ทิฏฺฐิคตานํ ปหานาย ปฐมาย ภูมิยา ปตฺติยา วิวิจฺเจว กาเมหิ ฯเปฯ ปฐมํ ฌานํ อุปสมฺปชฺช วิหรติ สุญฺญตํ ฉนฺทา\-ธิปเตยฺยํ. ตสฺมิํ สมเย ผสฺโส โหติ ฯเปฯ อวิกฺเขโป โหติ ฯเปฯ. อิเม ธมฺมา กุสลา. ตสฺเสว โลกุตฺตรสฺส กุสลสฺส ฌานสฺส กตตฺตา ภาวิตตฺตา วิปากํ วิวิจฺเจว กาเมหิ ฯเปฯ ปฐมํ ฌานํ อุปสมฺปชฺช วิหรติ อปฺปณิหิตํ ฉนฺทา\-ธิปเตยฺยํ. ตสฺมิํ สมเย ผสฺโส โหติ ฯเปฯ อวิกฺเขโป โหติ ฯเปฯ. อิเม ธมฺมา อพฺยากตา.}

\proseitem{537}{กตเม ธมฺมา อพฺยากตา. ยสฺมิํ สมเย โลกุตฺตรํ ฌานํ ภาเวติ นิยฺยานิกํ อปจยคามิํ ทิฏฺฐิคตานํ ปหานาย ปฐมาย ภูมิยา ปตฺติยา วิตกฺก\-วิจารานํ วูปสมา ฯเปฯ ทุติยํ ฌานํ ฯเปฯ ตติยํ ฌานํ ฯเปฯ จตุตฺถํ ฌานํ ฯเปฯ ปฐมํ ฌานํ ฯเปฯ ปญฺจมํ ฌานํ อุปสมฺปชฺช วิหรติ สุญฺญตํ ฉนฺทาธิปเตยฺยนฺติ กุสลํ ฯเปฯ สุญฺญตํ ฉนฺทาธิปเตยฺยนฺติ วิปาโก ฯเปฯ สุญฺญตํ ฉนฺทาธิปเตยฺยนฺติ กุสลํ ฯเปฯ อนิมิตฺตํ ฉนฺทาธิปเตยฺยนฺติ วิปาโก ฯเปฯ สุญฺญตํ ฉนฺทาธิปเตยฺยนฺติ กุสลํ ฯเปฯ อปฺปณิหิตํ ฉนฺทาธิปเตยฺยนฺติ วิปาโก. ตสฺมิํ สมเย ผสฺโส โหติ ฯเปฯ อวิกฺเขโป โหติ ฯเปฯ. อิเม ธมฺมา อพฺยากตา.}

\vspace{\tipitakaparskip}
\csromancenter{ฉนฺทา\-ธิปเตยฺย\-สุทฺธิก\-สุญฺญตา.}
\proseitem{538}{กตเม ธมฺมา อพฺยากตา. ยสฺมิํ สมเย โลกุตฺตรํ ฌานํ ภาเวติ นิยฺยานิกํ อปจยคามิํ ทิฏฺฐิคตานํ ปหานาย ปฐมาย ภูมิยา ปตฺติยา วิวิจฺเจว กาเมหิ ฯเปฯ ปฐมํ ฌานํ อุปสมฺปชฺช วิหรติ ทุกฺข\-ปฏิปทํ ทนฺธาภิญฺญํ สุญฺญตํ ฉนฺทา\-ธิปเตยฺยํ. ตสฺมิํ สมเย ผสฺโส โหติ ฯเปฯ อวิกฺเขโป โหติ ฯเปฯ. อิเม ธมฺมา กุสลา. ตสฺเสว โลกุตฺตรสฺส กุสลสฺส ฌานสฺส กตตฺตา ภาวิตตฺตา วิปาก วิวิจฺเจว กาเมหิ ฯเปฯ ปฐมํ ฌานํ อุปสมฺปชฺช วิหรติ ทุกฺข\-ปฏิปทํ ทนฺธาภิญฺญํ สุญฺญตํ ฉนฺทา\-ธิปเตยฺยํ. ตสฺมิํ สมเย ผสฺโส โหติ ฯเปฯ อวิกฺเขโป โหติ ฯเปฯ. อิเม ธมฺมา อพฺยากตา.}

\proseitem{539}{\csromanfolio{130}กตเม ธมฺมา อพฺยากตา. ยสฺมิํ สมเย โลกุตฺตรํ ฌานํ ภาเวติ นิยฺยานิกํ อปจยคามิํ ทิฏฺฐิคตานํ ปหานาย ปฐมาย ภูมิยา ปตฺติยา วิวิจฺเจว กาเมหิ ฯเปฯ ปฐมํ ฌานํ อุปสมฺปชฺช วิหรติ ทุกฺข\-ปฏิปทํ ทนฺธาภิญฺญํ สุญฺญตํ ฉนฺทา\-ธิปเตยฺยํ. ตสฺมิํ สมเย ผสฺโส โหติ ฯเปฯ อวิกฺเขโป โหติ ฯเปฯ. อิเม ธมฺมา กุสลา. ตสฺเสว โลกุตฺตรสฺส กุสลสฺส ฌานสฺส กตตฺตา ภาวิตตฺตา วิปากํ วิวิจฺเจว กาเมหิ ฯเปฯ ปฐมํ ฌานํ อุปสมฺปชฺช วิหรติ ทุกฺข\-ปฏิปทํ ทนฺธาภิญฺญํ อนิมิตฺตํ ฉนฺทา\-ธิปเตยฺยํ. ตสฺมิํ สมเย ผสฺโส โหติ ฯเปฯ อวิกฺเขโป โหติ ฯเปฯ. อิเม ธมฺมา อพฺยากตา.}

\proseitem{540}{กตเม ธมฺมา อพฺยากตา. ยสฺมิํ สมเย โลกุตฺตรํ ฌานํ ภาเวติ นิยฺยานิกํ อปจยคามิํ ทิฏฺฐิคตานํ ปหานาย ปฐมาย ภูมิยา ปตฺติยา วิวิจฺเจว กาเมหิ ฯเปฯ ปฐมํ ฌานํ อุปสมฺปชฺช วิหรติ ทุกฺข\-ปฏิปทํ ทนฺธาภิญฺญํ สุญฺญตํ ฉนฺทา\-ธิปเตยฺยํ. ตสฺมิํ สมเย ผสฺโส โหติ ฯเปฯ อวิกฺเขโป โหติ ฯเปฯ. อิเม ธมฺมา กุสลา. ตสฺเสว โลกุตฺตรสฺส กุสลสฺส ฌานสฺส กตตฺตา ภาวิตตฺตา วิปากํ วิวิจฺเจว กาเมหิ ฯเปฯ ปฐมํ ฌานํ อุปสมฺปชฺช วิหรติ ทุกฺข\-ปฏิปทํ ทนฺธาภิญฺญํ อปฺปณิหิตํ ฉนฺทา\-ธิปเตยฺยํ. ตสฺมิํ สมเย ผสฺโส โหติ ฯเปฯ อวิกฺเขโป โหติ ฯเปฯ. อิเม ธมฺมา อพฺยากตา.}

\proseitem{541}{กตเม ธมฺมา อพฺยากตา. ยสฺมิํ สมเย โลกุตฺตรํ ฌานํ ภาเวติ นิยฺยานิกํ อปจยคามิํ ทิฏฺฐิคตานํ ปหานาย ปฐมาย ภูมิยา ปตฺติยา วิตกฺก\-วิจารานํ วูปสมา ฯเปฯ ทุติยํ ฌานํ ฯเปฯ ตติยํ ฌานํ ฯเปฯ จตุตฺถํ ฌานํ ฯเปฯ ปฐมํ ฌานํ ฯเปฯ ปญฺจมํ ฌานํ อุปสมฺปชฺช วิหรติ ทุกฺข\-ปฏิปทํ ทนฺธาภิญฺญํ สุญฺญตํ ฉนฺทาธิปเตยฺยนฺติ กุสลํ ฯเปฯ ทุกฺข\-ปฏิปทํ ทนฺธาภิญฺญํ สุญฺญตํ ฉนฺทาธิปเตยฺยนฺติ วิปาโก ฯเปฯ ทุกฺข\-ปฏิปทํ ทนฺธาภิญฺญํ สุญฺญตํ ฉนฺทาธิปเตยฺยนฺติ กุสลํ ฯเปฯ ทุกฺข\-ปฏิปทํ ทนฺธาภิญฺญํ อนิมิตฺตํ ฉนฺทาธิปเตยฺยนฺติ วิปาโก ฯเปฯ ทุกฺข\-ปฏิปทํ ทนฺธาภิญฺญํ สุญฺญตํ ฉนฺทาธิปเตยฺยนฺติ กุสลํ ฯเปฯ ทุกฺข\-ปฏิปทํ ทนฺธาภิญฺญํ อปฺปณิหิตํ ฉนฺทาธิปเตยฺยนฺติ วิปาโก. ตสฺมิํ สมเย ผสฺโส โหติ ฯเปฯ อวิกฺเขโป โหติ ฯเปฯ. อิเม ธมฺมา อพฺยากตา.}

\proseitem{542}{กตเม ธมฺมา อพฺยากตา. ยสฺมิํ สมเย โลกุตฺตรํ ฌานํ ภาเวติ นิยฺยานิกํ อปจยคามิํ ทิฏฺฐิคตานํ ปหานาย ปฐมาย \csromanfolio{131}ภูมิยา ปตฺติยา วิวิจฺเจว กาเมหิ ฯเปฯ ปฐมํ ฌานํ อุปสมฺปชฺช วิหรติ ทุกฺข\-ปฏิปทํ ขิปฺปาภิญฺญํ สุญฺญตํ ฉนฺทา\-ธิปเตยฺยํ ฯเปฯ สุขปฏิปทํ ทนฺธาภิญฺญํ สุญฺญตํ ฉนฺทา\-ธิปเตยฺยํ ฯเปฯ สุขปฏิปทํ ขิปฺปาภิญฺญํ สุญฺญตํ ฉนฺทา\-ธิปเตยฺยํ ฯเปฯ ทุติยํ ฌานํ ฯเปฯ ตติยํ ฌานํ ฯเปฯ จตุตฺถํ ฌานํ ฯเปฯ ปฐมํ ฌานํ ฯเปฯ ปญฺจมํ ฌานํ อุปสมฺปชฺช วิหรติ สุขปฏิปทํ ขิปฺปาภิญฺญํ สุญฺญตํ ฉนฺทาธิปเตยฺยนฺติ กุสลํ ฯเปฯ สุขปฏิปทํ ขิปฺปาภิญฺญํ สุญฺญตํ ฉนฺทาธิปเตยฺยนฺติ วิปาโก ฯเปฯ สุขปฏิปทํ ขิปฺปาภิญฺญํ สุญฺญตํ ฉนฺทาธิปเตยฺยนฺติ กุสลํ ฯเปฯ สุขปฏิปทํ ขิปฺปาภิญฺญํ อนิมิตฺตํ ฉนฺทาธิปเตยฺยนฺติ วิปาโก ฯเปฯ สุขปฏิปทํ ขิปฺปาภิญฺญํ สุญฺญตํ ฉนฺทาธิปเตยฺยนฺติ กุสลํ ฯเปฯ สุขปฏิปทํ ขิปฺปาภิญฺญํ อปฺปณิหิตํ ฉนฺทาธิปเตยฺยนฺติ วิปาโก. ตสฺมิํ สมเย ผสฺโส โหติ ฯเปฯ อวิกฺเขโป โหติ ฯเปฯ. อิเม ธมฺมา อพฺยากตา.}

\proseitem{543}{กตเม ธมฺมา อพฺยากตา. ยสฺมิํ สมเย โลกุตฺตรํ ฌานํ ภาเวติ นิยฺยานิกํ อปจยคามิํ ทิฏฺฐิคตานํ ปหานาย ปฐมาย ภูมิยา ปตฺติยา วิวิจฺเจว กาเมหิ ฯเปฯ ปฐมํ ฌานํ อุปสมฺปชฺช วิหรติ อปฺปณิหิตํ ฉนฺทา\-ธิปเตยฺยํ. ตสฺมิํ สมเย ผสฺโส โหติ ฯเปฯ อวิกฺเขโป โหติ ฯเปฯ. อิเม ธมฺมา กุสลา. ตสฺเสว โลกุตฺตรสฺส กุสลสฺส ฌานสฺส กตตฺตา ภาวิตตฺตา วิปากํ วิวิจฺเจว กาเมหิ ฯเปฯ ปฐมํ ฌานํ อุปสมฺปชฺช วิหรติ อปฺปณิหิตํ ฉนฺทา\-ธิปเตยฺยํ. ตสฺมิํ สมเย ผสฺโส โหติ ฯเปฯ อวิกฺเขโป โหติ ฯเปฯ. อิเม ธมฺมา อพฺยากตา.}

\proseitem{544}{กตเม ธมฺมา อพฺยากตา. ยสฺมิํ สมเย โลกุตฺตรํ ฌานํ ภาเวติ นิยฺยานิกํ อปจยคามิํ ทิฏฺฐิคตานํ ปหานาย ปฐมาย ภูมิยา ปตฺติยา วิวิจฺเจว กาเมหิ ฯเปฯ ปฐมํ ฌานํ อุปสมฺปชฺช วิหรติ อปฺปณิหิตํ ฉนฺทา\-ธิปเตยฺยํ. ตสฺมิํ สมเย ผสฺโส โหติ ฯเปฯ อวิกฺเขโป โหติ ฯเปฯ. อิเม ธมฺมา กุสลา. ตสฺเสว โลกุตฺตรสฺส กุสลสฺส ฌานสฺส กตตฺตา ภาวิตตฺตา วิปากํ วิวิจฺเจว กาเมหิ ฯเปฯ ปฐมํ ฌานํ อุปสมฺปชฺช วิหรติ อนิมิตฺตํ ฉนฺทา\-ธิปเตยฺยํ. ตสฺมิํ สมเย ผสฺโส โหติ ฯเปฯ อวิกฺเขโป โหติ ฯเปฯ. อิเม ธมฺมา อพฺยากตา.}

\proseitem{545}{กตเม ธมฺมา อพฺยากตา. ยสฺมิํ สมเย โลกุตฺตรํ ฌานํ ภาเวติ นิยฺยานิกํ อปจยคามิํ ทิฏฺฐิคตานํ ปหานาย ปฐมาย ภูมิยา ปตฺติยา วิวิจฺเจว กาเมหิ ฯเปฯ ปฐมํ ฌานํ อุปสมฺปชฺช วิหรติ \csromanfolio{132}อปฺปณิหิตํ ฉนฺทา\-ธิปเตยฺยํ. ตสฺมิํ สมเย ผสฺโส โหติ ฯเปฯ อวิกฺเขโป โหติ ฯเปฯ. อิเม ธมฺมา กุสลา. ตสฺเสว โลกุตฺตรสฺส กุสลสฺส ฌานสฺส กตตฺตา ภาวิตตฺตา วิปากํ วิวิจฺเจว กาเมหิ ฯเปฯ ปฐมํ ฌานํ อุปสมฺปชฺช วิหรติ สุญฺญตํ ฉนฺทา\-ธิปเตยฺยํ. ตสฺมิํ สมเย ผสฺโส โหติ ฯเปฯ อวิกฺเขโป โหติ ฯเปฯ. อิเม ธมฺมา อพฺยากตา.}

\proseitem{546}{กตเม ธมฺมา อพฺยากตา. ยสฺมิํ สมเย โลกุตฺตรํ ฌานํ ภาเวติ นิยฺยานิกํ อปจยคามิํ ทิฏฺฐิคตานํ ปหานาย ปฐมาย ภูมิยา ปตฺติยา วิตกฺก\-วิจารานํ วูปสมา ฯเปฯ ทุติยํ ฌานํ ฯเปฯ ตติยํ ฌานํ ฯเปฯ จตุตฺถํ ฌานํ ฯเปฯ ปฐมํ ฌานํ ฯเปฯ ปญฺจมํ ฌานํ อุปสมฺปชฺช วิหรติ อปฺปณิหิตํ ฉนฺทาธิปเตยฺยนฺติ กุสลํ ฯเปฯ อปฺปณิหิตํ ฉนฺทาธิปเตยฺยนฺติ วิปาโก ฯเปฯ อปฺปณิหิตํ ฉนฺทาธิปเตยฺยนฺติ กุสลํ ฯเปฯ อนิมิตฺตํ ฉนฺทาธิปเตยฺยนฺติ วิปาโก ฯเปฯ อปฺปณิหิตํ ฉนฺทาธิปเตยฺยนฺติ กุสลํ ฯเปฯ สุญฺญตํ ฉนฺทาธิปเตยฺยนฺติ วิปาโก. ตสฺมิํ สมเย ผสฺโส โหติ ฯเปฯ อวิกฺเขโป โหติ ฯเปฯ. อิเม ธมฺมา อพฺยากตา.}

\proseitem{547}{กตเม ธมฺมา อพฺยากตา. ยสฺมิํ สมเย โลกุตฺตรํ ฌานํ ภาเวติ นิยฺยานิกํ อปจยคามิํ ทิฏฺฐิคตานํ ปหานาย ปฐมาย ภูมิยา ปตฺติยา วิวิจฺเจว กาเมหิ ฯเปฯ ปฐมํ ฌานํ อุปสมฺปชฺช วิหรติ ทุกฺข\-ปฏิปทํ ทนฺธาภิญฺญํ อปฺปณิหิตํ ฉนฺทา\-ธิปเตยฺยํ. ตสฺมิํ สมเย ผสฺโส โหติ ฯเปฯ อวิกฺเขโป โหติ ฯเปฯ. อิเม ธมฺมา กุสลา. ตสฺเสว โลกุตฺตรสฺส กุสลสฺส ฌานสฺส กตตฺตา ภาวิตตฺตา วิปากํ วิวิจฺเจว กาเมหิ ฯเปฯ ปฐมํ ฌานํ อุปสมฺปชฺช วิหรติ ทุกฺข\-ปฏิปทํ ทนฺธาภิญฺญํ อปฺปณิหิตํ ฉนฺทา\-ธิปเตยฺยํ. ตสฺมิํ สมเย ผสฺโส โหติ ฯเปฯ อวิกฺเขโป โหติ ฯเปฯ. อิเม ธมฺมา อพฺยากตา.}

\proseitem{548}{กตเม ธมฺมา อพฺยากตา. ยสฺมิํ สมเย โลกุตฺตรํ ฌานํ ภาเวติ นิยฺยานิกํ อปจยคามิํ ทิฏฺฐิคตานํ ปหานาย ปฐมาย ภูมิยา ปตฺติยา วิวิจฺเจว กาเมหิ ฯเปฯ ปฐมํ ฌานํ อุปสมฺปชฺช วิหรติ ทุกฺข\-ปฏิปทํ ทนฺธาภิญฺญํ อปฺปณิหิตํ ฉนฺทา\-ธิปเตยฺยํ. ตสฺมิํ สมเย ผสฺโส โหติ ฯเปฯ อวิกฺเขโป โหติ ฯเปฯ. อิเม ธมฺมา กุสลา. ตสฺเสว โลกุตฺตรสฺส กุสลสฺส ฌานสฺส กตตฺตา ภาวิตตฺตา วิปากํ วิวิจฺเจว กาเมหิ ฯเปฯ ปฐมํ ฌานํ อุปสมฺปชฺช วิหรติ ทุกฺข\-ปฏิปทํ ทนฺธาภิญฺญํ อนิมิตฺตํ \csromanfolio{133}ฉนฺทา\-ธิปเตยฺยํ. ตสฺมิํ สมเย ผสฺโส โหติ ฯเปฯ อวิกฺเขโป โหติ ฯเปฯ. อิเม ธมฺมา อพฺยากตา.}

\proseitem{549}{กตเม ธมฺมา อพฺยากตา. ยสฺมิํ สมเย โลกุตฺตรํ ฌานํ ภาเวติ นิยฺยานิกํ อปจยคามิํ ทิฏฺฐิคตานํ ปหานาย ปฐมาย ภูมิยา ปตฺติยา วิวิจฺเจว กาเมหิ ฯเปฯ ปฐมํ ฌานํ อุปสมฺปชฺช วิหรติ ทุกฺข\-ปฏิปทํ ทนฺธาภิญฺญํ อปฺปณิหิตํ ฉนฺทา\-ธิปเตยฺยํ. ตสฺมิํ สมเย ผสฺโส โหติ ฯเปฯ อวิกฺเขโป โหติ ฯเปฯ. อิเม ธมฺมา กุสลา. ตสฺเสว โลกุตฺตรสฺส กุสลสฺส ฌานสฺส กตตฺตา ภาวิตตฺตา วิปากํ วิวิจฺเจว กาเมหิ ฯเปฯ ปฐมํ ฌานํ อุปสมฺปชฺช วิหรติ ทุกฺข\-ปฏิปทํ ทนฺธาภิญฺญํ สุญฺญตํ ฉนฺทา\-ธิปเตยฺยํ. ตสฺมิํ สมเย ผสฺโส โหติ ฯเปฯ อวิกฺเขโป โหติ ฯเปฯ. อิเม ธมฺมา อพฺยากตา.}

\proseitem{550}{กตเม ธมฺมา อพฺยากตา. ยสฺมิํ สมเย โลกุตฺตรํ ฌานํ ภาเวติ นิยฺยานิกํ อปจยคามิํ ทิฏฺฐิคตานํ ปหานาย ปฐมาย ภูมิยา ปตฺติยา วิตกฺก\-วิจารานํ วูปสมา ฯเปฯ ทุติยํ ฌานํ ฯเปฯ ตติยํ ฌานํ ฯเปฯ จตุตฺถํ ฌานํ ฯเปฯ ปฐมํ ฌานํ ฯเปฯ ปญฺจมํ ฌานํ อุปสมฺปชฺช วิหรติ ทุกฺข\-ปฏิปทํ ทนฺธาภิญฺญํ อปฺปณิหิตํ ฉนฺทาธิปเตยฺยนฺติ กุสลํ ฯเปฯ ทุกฺข\-ปฏิปทํ ทนฺธาภิญฺญํ อปฺปณิหิตํ ฉนฺทาธิปเตยฺยนฺติ วิปาโก ฯเปฯ ทุกฺข\-ปฏิปทํ ทนฺธาภิญฺญํ อปฺปณิหิตํ ฉนฺทาธิปเตยฺยนฺติ กุสลํ ฯเปฯ ทุกฺข\-ปฏิปทํ ทนฺธาภิญฺญํ อนิมิตฺตํ ฉนฺทาธิปเตยฺยนฺติ วิปาโก ฯเปฯ ทุกฺข\-ปฏิปทํ ทนฺธาภิญฺญํ อปฺปณิหิตํ ฉนฺทาธิปเตยฺยนฺติ กุสลํ ฯเปฯ ทุกฺข\-ปฏิปทํ ทนฺธาภิญฺญํ สุญฺญตํ ฉนฺทาธิปเตยฺยนฺติ วิปาโก. ตสฺมิํ สมเย ผสฺโส โหติ ฯเปฯ อวิกฺเขโป โหติ ฯเปฯ. อิเม ธมฺมา อพฺยากตา.}

\proseitem{551}{กตเม ธมฺมา อพฺยากตา. ยสฺมิํ สมเย โลกุตฺตรํ ฌานํ ภาเวติ นิยฺยานิกํ อปจยคามิํ ทิฏฺฐิคตานํ ปหานาย ปฐมาย ภูมิยา ปตฺติยา วิวิจฺเจว กาเมหิ ฯเปฯ ปฐมํ ฌานํ อุปสมฺปชฺช วิหรติ ทุกฺข\-ปฏิปทํ ขิปฺปาภิญฺญํ อปฺปณิหิตํ ฉนฺทา\-ธิปเตยฺยํ ฯเปฯ สุขปฏิปทํ ทนฺธาภิญฺญํ อปฺปณิหิตํ ฉนฺทา\-ธิปเตยฺยํ ฯเปฯ สุขปฏิปทํ ขิปฺปาภิญฺญํ อปฺปณิหิตํ ฉนฺทา\-ธิปเตยฺยํ ฯเปฯ ทุติยํ ฌานํ ฯเปฯ ตติยํ ฌานํ ฯเปฯ จตุตฺถํ ฌานํ ฯเปฯ ปฐมํ ฌานํ ฯเปฯ ปญฺจมํ ฌานํ อุปสมฺปชฺช วิหรติ สุขปฏิปทํ ขิปฺปาภิญฺญํ อปฺปณิหิตํ ฉนฺทาธิปเตยฺยนฺติ กุสลํ ฯเปฯ สุขปฏิปทํ ขิปฺปาภิญฺญํ อปฺปณิหิตํ ฉนฺทาธิปเตยฺยนฺติ วิปาโก ฯเปฯ \csromanfolio{134}สุขปฏิปทํ ขิปฺปาภิญฺญํ อปฺปณิหิตํ ฉนฺทาธิปเตยฺยนฺติ กุสลํ ฯเปฯ สุขปฏิปทํ ขิปฺปาภิญฺญํ อนิมิตฺตํ ฉนฺทาธิปเตยฺยนฺติ วิปาโก ฯเปฯ สุขปฏิปทํ ขิปฺปาภิญฺญํ อปฺปณิหิตํ ฉนฺทาธิปเตยฺยนฺติ กุสลํ ฯเปฯ สุขปฏิปทํ ขิปฺปาภิญฺญํ สุญฺญตํ ฉนฺทาธิปเตยฺยนฺติ วิปาโก. ตสฺมิํ สมเย ผสฺโส โหติ ฯเปฯ อวิกฺเขโป โหติ ฯเปฯ. อิเม ธมฺมา อพฺยากตา.}

\proseitem{552}{กตเม ธมฺมา อพฺยากตา. ยสฺมิํ สมเย โลกุตฺตรํ มคฺคํ ภาเวติ ฯเปฯ โลกุตฺตรํ สติปฏฺฐานํ ภาเวติ ฯเปฯ โลกุตฺตรํ สมฺมปฺปธานํ ภาเวติ ฯเปฯ โลกุตฺตรํ อิทฺธิปาทํ ภาเวติ ฯเปฯ โลกุตฺตรํ อินฺทฺริยํ ภาเวติ ฯเปฯ โลกุตฺตรํ พลํ ภาเวติ ฯเปฯ โลกุตฺตรํ โพชฺฌงฺคํ ภาเวติ ฯเปฯ โลกุตฺตรํ สจฺจํ ภาเวติ ฯเปฯ โลกุตฺตรํ สมถํ ภาเวติ ฯเปฯ โลกุตฺตรํ ธมฺมํ ภาเวติ ฯเปฯ โลกุตฺตรํ ขนฺธํ ภาเวติ ฯเปฯ โลกุตฺตรํ อายตนํ ภาเวติ ฯเปฯ โลกุตฺตรํ ธาตุํ ภาเวติ ฯเปฯ โลกุตฺตรํ อาหารํ ภาเวติ ฯเปฯ โลกุตฺตรํ ผสฺสํ ภาเวติ ฯเปฯ โลกุตฺตรํ เวทนํ ภาเวติ ฯเปฯ โลกุตฺตรํ สญฺญํ ภาเวติ ฯเปฯ โลกุตฺตรํ เจตนํ ภาเวติ ฯเปฯ โลกุตฺตรํ จิตฺตํ ภาเวติ นิยฺยานิกํ อปจยคามิํ ทิฏฺฐิคตานํ ปหานาย ปฐมาย ภูมิยา ปตฺติยา วิวิจฺเจว กาเมหิ ฯเปฯ ปฐมํ ฌานํ อุปสมฺปชฺช วิหรติ ทุกฺข\-ปฏิปทํ ทนฺธาภิญฺญํ ฉนฺทา\-ธิปเตยฺยํ. ตสฺมิํ สมเย ผสฺโส โหติ ฯเปฯ อวิกฺเขโป โหติ ฯเปฯ. อิเม ธมฺมา กุสลา. ตสฺเสว โลกุตฺตรสฺส กุสลสฺส ฌานสฺส กตตฺตา ภาวิตตฺตา วิปากํ วิวิจฺเจว กาเมหิ ฯเปฯ ปฐมํ ฌานํ อุปสมฺปชฺช วิหรติ ทุกฺข\-ปฏิปทํ ทนฺธาภิญฺญํ สุญฺญตํ ฯเปฯ อนิมิตฺตํ ฯเปฯ อปฺปณิหิตํ ฉนฺทา\-ธิปเตยฺยํ ฯเปฯ วีริยา\-ธิปเตยฺยํ ฯเปฯ จิตฺตา\-ธิปเตยฺยํ ฯเปฯ วีมํสา\-ธิปเตยฺยํ. ตสฺมิํ สมเย ผสฺโส โหติ ฯเปฯ อวิกฺเขโป โหติ ฯเปฯ. อิเม ธมฺมา อพฺยากตา.}

\vspace{\tipitakaparskip}
\csromancenter{ปฐมมคฺค\-วิปาโก.}
\csromanmatikaanchor{mtk.67}
\csromantocmarkref{subsection}{ทุติยาทิมคฺควิปาก}{mtk.67}
\bookmark[dest={mtk.67},level=2]{ทุติยาทิมคฺควิปาก}
\csromanheader{2}{ทุติยา\-ทิ\-มคฺค\-วิปาก}
\proseitem{553}{กตเม ธมฺมา อพฺยากตา. ยสฺมิํ สมเย โลกุตฺตรํ ฌานํ ภาเวติ นิยฺยานิกํ อปจยคามิํ กามราค\-พฺยาปาทานํ \csromanfolio{135}ตนุภาวาย ทุติยาย ภูมิยา ปตฺติยา ฯเปฯ กามราค\-พฺยาปาทานํ อนวเสส\-ปฺปหานาย ตติยาย ภูมิยา ปตฺติยา ฯเปฯ รูปราคอรูปราคมาน- อุทฺธจฺจอวิชฺชาย อนวเสส\-ปฺปหานาย จตุตฺถาย ภูมิยา ปตฺติยา วิวิจฺเจว กาเมหิ ฯเปฯ ปฐมํ ฌานํ อุปสมฺปชฺช วิหรติ ทุกฺข\-ปฏิปทํ ทนฺธาภิญฺญํ. ตสฺมิํ สมเย ผสฺโส โหติ ฯเปฯ อญฺญินฺทฺริยํ โหติ ฯเปฯ อวิกฺเขโป โหติ ฯเปฯ. อิเม ธมฺมา กุสลา. ตสฺเสว โลกุตฺตรสฺส กุสลสฺส ฌานสฺส กตตฺตา ภาวิตตฺตา วิปากํ วิวิจฺเจว กาเมหิ ฯเปฯ ปฐมํ ฌานํ อุปสมฺปชฺช วิหรติ ทุกฺข\-ปฏิปทํ ทนฺธาภิญฺญํ สุญฺญตํ. ตสฺมิํ สมเย ผสฺโส โหติ ฯเปฯ อญฺญาตาวิ\-นฺทฺริยํ โหติ ฯเปฯ อวิกฺเขโป โหติ ฯเปฯ. เย วา ปน ตสฺมิํ สมเย อญฺเญปิ อตฺถิ ปฏิจฺจ\-สมุปฺปนฺนา อรูปิโน ธมฺมา. อิเม ธมฺมา อพฺยากตา.}

\proseitem{554}{กตโม ตสฺมิํ สมเย ผสฺโส โหติ. โย ตสฺมิํ สมเย ผสฺโส ผุสนา สมฺผุสนา สมฺผุสิตตฺตํ. อยํ ตสฺมิํ สมเย ผสฺโส โหติ ฯเปฯ.}

\proseitem{555}{กตมํ ตสฺมิํ สมเย อญฺญตาวินฺทฺริยํ โหติ. ยา เตสํ อญฺญาตาวีนํ ธมฺมานํ อญฺญา ปญฺญา ปชานนา วิจโย ปวิจโย ธมฺมวิจโย สลฺลกฺขณา อุปลกฺขณา ปจฺจุ\-ปลกฺขณา ปณฺฑิจฺจํ โกสลฺลํ เนปุญฺญํ เวภพฺยา จินฺตา อุปปริกฺขา ภูรี เมธา ปริณายิกา วิปสฺสนา สมฺปชญฺญํ ปโตโท ปญฺญา ปญฺญินฺทฺริยํ ปญฺญาพลํ ปญฺญาสตฺถํ ปญฺญาปาสาโท ปญฺญาอาโลโก ปญฺญาโอภาโส ปญฺญาปชฺโชโต ปญฺญารตนํ อโมโห ธมฺมวิจโย สมฺมาทิฏฺฐิ ธมฺมวิจย\-สมฺโพชฺฌงฺโค มคฺคงฺคํ มคฺค\-ปริยาปนฺนํ. อิทํ ตสฺมิํ สมเย อญฺญาตาวิ\-นฺทฺริยํ โหติ ฯเปฯ. เย วา ปน ตสฺมิํ สมเย อญฺเญปิ อตฺถิ ปฏิจฺจ\-สมุปฺปนฺนา อรูปิโน ธมฺมา. อิเม ธมฺมา อพฺยากตา.}

\vspace{\tipitakaparskip}
\csromancenter{ทุติยา\-ทิ\-มคฺค\-วิปาโก.}
\csromancenter{โลกุตฺตร\-วิปาโก.}
\csromanheader{2}{อกุสลวิปาก\-อพฺยากต}
\csromanmatikaanchor{mtk.68}
\csromantocmarkref{subsection}{อกุสลวิปากอพฺยากตอกุสลวิปากปญฺจวิญฺญาณ ...}{mtk.68}
\bookmark[dest={mtk.68},level=2]{อกุสลวิปากอพฺยากตอกุสลวิปากปญฺจวิญฺญาณ ...}
\csromanheader{2}{อกุสลวิปาก\-ปญฺจวิญฺญาณ}
\proseitem{556}{\csromanfolio{136}กตเม ธมฺมา อพฺยากตา. ยสฺมิํ สมเย อกุสลสฺส กมฺมสฺส กตตฺตา อุปจิตตฺตา วิปากํ จกฺขุวิญฺญาณํ อุปฺปนฺนํ โหติ อุเปกฺขา\-สหคตํ รูปารมฺมณํ ฯเปฯ โสตวิญฺญาณํ อุปฺปนฺนํ โหติ อุเปกฺขา\-สหคตํ สทฺทารมฺมณํ ฯเปฯ ฆานวิญฺญาณํ อุปฺปนฺนํ โหติ อุเปกฺขา\-สหคตํ คนฺธารมฺมณํ ฯเปฯ ชิวฺหาวิญฺญาณํ อุปฺปนฺนํ โหติ อุเปกฺขา\-สหคตํ รสารมฺมณํ ฯเปฯ กายวิญฺญาณํ อุปฺปนฺนํ โหติ ทุกฺข\-สหคตํ โผฏฺฐพฺพา\-รมฺมณํ. ตสฺมิํ สมเย ผสฺโส โหติ, เวทนา โหติ, สญฺญา โหติ, เจตนา โหติ, จิตฺตํ โหติ, ทุกฺขํ โหติ, จิตฺตสฺเส\-กคฺคตา โหติ, มนินฺทฺริยํ โหติ, ทุกฺขินฺทฺริยํ โหติ, ชีวิตินฺทฺริยํ โหติ. เย วา ปน ตสฺมิํ สมเย อญฺเญปิ อตฺถิ ปฏิจฺจ\-สมุปฺปนฺนา อรูปิโน ธมฺมา. อิเม ธมฺมา อพฺยากตา.}

\proseitem{557}{กตโม ตสฺมิํ สมเย ผสฺโส โหติ. โย ตสฺมิํ สมเย ผสฺโส ผุสนา สมฺผุสนา สมฺผุสิตตฺตํ. อยํ ตสฺมิํ สมเย ผสฺโส โหติ.}

\proseitem{558}{กตมา ตสฺมิํ สมเย เวทนา โหติ. ยํ ตสฺมิํ สมเย ตชฺชา\-กายวิญฺญาณธาตุ\-สมฺผสฺสชํ กายิกํ อสาตํ กายิกํ ทุกฺขํ กาย\-สมฺผสฺสชํ อสาตํ ทุกฺขํ เวทยิตํ กาย\-สมฺผสฺสชา อสาตา ทุกฺขา เวทนา. อยํ ตสฺมิํ สมเย เวทนา โหติ ฯเปฯ.}

\proseitem{559}{กตมํ ตสฺมิํ สมเย ทุกฺขํ โหติ. ยํ ตสฺมิํ สมเย กายิกํ อสาตํ กายิกํ ทุกฺขํ กาย\-สมฺผสฺสชํ อสาตํ ทุกฺขํ เวทยิตํ กาย\-สมฺผสฺสชา อสาตา ทุกฺขา เวทนา. อิทํ ตสฺมิํ สมเย ทุกฺขํ โหติ ฯเปฯ.}

\proseitem{560}{กตมํ ตสฺมิํ สมเย ทุกฺขินฺทฺริยํ โหติ. ยํ ตสฺมิํ สมเย กายิกํ อสาตํ กายิกํ ทุกฺขํ กาย\-สมฺผสฺสชํ อสาตํ ทุกฺขํ เวทยิตํ กาย\-สมฺผสฺสชา อสาตา ทุกฺขา เวทนา. อิทํ ตสฺมิํ สมเย ทุกฺขินฺทฺริยํ โหติ ฯเปฯ. เย วา ปน ตสฺมิํ สมเย อญฺเญปิ อตฺถิ ปฏิจฺจ\-สมุปฺปนฺนา อรูปิโน ธมฺมา. อิเม ธมฺมา อพฺยากตา.}

\chapterheadpageread{1. จิตฺตุปฺปาท\-กณฺฑ}
\prose{\csromanfolio{137}ตสฺมิํ โข ปน สมเย จตฺตาโร ขนฺธา โหนฺติ, ทฺวายตนานิ โหนฺติ, เทฺว ธาตุโย โหนฺติ, ตโย อาหารา โหนฺติ, ตีณินฺทฺริยานิ โหนฺติ, เอโก ผสฺโส โหติ ฯเปฯ เอกา กายวิญฺญาณ\-ธาตุ โหติ, เอกํ ธมฺมายตนํ โหติ, เอกา ธมฺมธาตุ โหติ. เย วา ปน ตสฺมิํ สมเย อญฺเญปิ อตฺถิ ปฏิจฺจ\-สมุปฺปนฺนา อรูปิโน ธมฺมา. อิเม ธมฺมา อพฺยากตา ฯเปฯ.}

\proseitem{561}{กตโม ตสฺมิํ สมเย สงฺขาร\-กฺขนฺโธ โหติ. ผสฺโส เจตนา จิตฺตสฺเส\-กคฺคตา ชีวิตินฺทฺริยํ. เย วา ปน ตสฺมิํ สมเย อญฺเญปิ อตฺถิ ปฏิจฺจ\-สมุปฺปนฺนา อรูปิโน ธมฺมา ฐเปตฺวา เวทนา\-กฺขนฺธํ ฐเปตฺวา สญฺญากฺขนฺธํ ฐเปตฺวา วิญฺญาณ\-กฺขนฺธํ. อยํ ตสฺมิํ สมเย สงฺขาร\-กฺขนฺโธ โหติ ฯเปฯ. อิเม ธมฺมา อพฺยากตา.}

\vspace{\tipitakaparskip}
\csromancenter{อกุสลวิปาก\-ปญฺจวิญฺญาณานิ.}
\csromanmatikaanchor{mtk.69}
\csromantocmarkref{subsection}{อกุสลวิปากมโนธาตุ}{mtk.69}
\bookmark[dest={mtk.69},level=2]{อกุสลวิปากมโนธาตุ}
\csromanheader{2}{อกุสลวิปาก\-มโนธาตุ}
\proseitem{562}{กตเม ธมฺมา อพฺยากตา. ยสฺมิํ สมเย อกุสลสฺส กมฺมสฺส กตตฺตา อุปจิตตฺตา วิปากา มโนธาตุ อุปฺปนฺนา โหติ อุเปกฺขา\-สหคตา รูปารมฺมณา วา ฯเปฯ โผฏฺฐพฺพา\-รมฺมณา วา ยํ ยํ วา ปนารพฺภ. ตสฺมิํ สมเย ผสฺโส โหติ, เวทนา โหติ, สญฺญา โหติ, เจตนา โหติ, จิตฺตํ โหติ, วิตกฺโก โหติ, วิจาโร โหติ, อุเปกฺขา โหติ, จิตฺตสฺเส\-กคฺคตา โหติ, มนินฺทฺริยํ โหติ, อุเปกฺขินฺทฺริยํ โหติ, ชีวิตินฺทฺริยํ โหติ. เย วา ปน ตสฺมิํ สมเย อญฺเญปิ อตฺถิ ปฏิจฺจ\-สมุปฺปนฺนา อรูปิโน ธมฺมา. อิเม ธมฺมา อพฺยากตา ฯเปฯ.}

\prose{ตสฺมิํ โข ปน สมเย จตฺตาโร ขนฺธา โหนฺติ, ทฺวายตนานิ โหนฺติ, เทฺว ธาตุโย โหนฺติ, ตโย อาหารา โหนฺติ, ตีณินฺทฺริยานิ โหนฺติ, เอโก ผสฺโส โหติ ฯเปฯ รฺกา มโนธาตุ โหติ, เอกํ ธมฺมายตนํ โหติ, เอกา ธมฺมธาตุ โหติ. เย วา ปน ตสฺมิํ สมเย อญฺเญปิ อตฺถิ ปฏิจฺจ\-สมุปฺปนฺนา อรูปิโน ธมฺมา. อิเม ธมฺมา อพฺยากตา ฯเปฯ.}

\proseitem{563}{\csromanfolio{138}กตโม ตสฺมิํ สมเย สงฺขาร\-กฺขนฺโธ โหติ. ผสฺโส เจตนา วิตกฺโก วิจาโร จิตฺตสฺเส\-กคฺคตา ชีวิตินฺทฺริยํ. เย วา ปน ตสฺมิํ สมเย อญฺเญปิ อตฺถิ ปฏิจฺจ\-สมุปฺปนฺนา อรูปิโน ธมฺมา ฐเปตฺวา เวทนา\-กฺขนฺธํ ฐเปตฺวา สญฺญากฺขนฺธํ ฐเปตฺวา วิญฺญาณ\-กฺขนฺธํ. อยํ ตสฺมิํ สมเย สงฺขาร\-กฺขนฺโธ โหติ ฯเปฯ. อิเม ธมฺมา อพฺยากตา.}

\vspace{\tipitakaparskip}
\csromancenter{อกุสลวิปากา มโนธาตุ.}
\csromanmatikaanchor{mtk.70}
\csromantocmarkref{subsection}{อกุสลวิปากมโนวิญฺญาณธาตุ}{mtk.70}
\bookmark[dest={mtk.70},level=2]{อกุสลวิปากมโนวิญฺญาณธาตุ}
\csromanheader{2}{อกุสลวิปาก\-มโนวิญฺญาณ\-ธาตุ}
\proseitem{564}{กตเม ธมฺมา อพฺยากตา. ยสฺมิํ สมเย อกุสลสฺส กมฺมสฺส กตตฺตา อุปจิตตฺตา วิปากา มโนวิญฺญาณ\-ธาตุ อุปฺปนฺนา โหติ อุเปกฺขา\-สหคตา รูปารมฺมณา วา ฯเปฯ ธมฺมารมฺมณา วา ยํ ยํ วา ปนารพฺภ. ตสฺมิํ สมเย ผสฺโส โหติ, เวทนา โหติ, สญฺญา โหติ, เจตนา โหติ, จิตฺตํ โหติ, วิตกฺโก โหติ, วิจาโร โหติ, อุเปกฺขา โหติ, จิตฺตสฺเส\-กคฺคตา โหติ, มนินฺทฺริยํ โหติ, อุเปกฺขินฺทฺริยํ โหติ, ชีวิตินฺทฺริยํ โหติ. เย วา ปน ตสฺมิํ สมเย อญฺเญปิ อตฺถิ ปฏิจฺจ\-สมุปฺปนฺนา อรูปิโน ธมฺมา. อิเม ธมฺมา อพฺยากตา ฯเปฯ.}

\prose{ตสฺมิํ โข ปน สมเย จตฺตาโร ขนฺธา โหนฺติ, ทฺวายตนานิ โหนฺติ, เทฺว ธาตุโย โหนฺติ, ตโย อาหารา โหนฺติ, ตีณินฺทฺริยานิ โหนฺติ, เอโก ผสฺโส โหติ ฯเปฯ เอกา มโนวิญฺญาณ\-ธาตุ โหติ, เอกํ ธมฺมายตนํ โหติ, เอกา ธมฺมธาตุ โหติ. เย วา ปน ตสฺมิํ สมเย อญฺเญปิ อตฺถิ ปฏิจฺจ\-สมุปฺปนฺนา อรูปิโน ธมฺมา. อิเม ธมฺมา อพฺยากตา ฯเปฯ.}

\proseitem{565}{กตโม ตสฺมิํ สมเย สงฺขาร\-กฺขนฺโธ โหติ. ผสฺโส เจตนา วิตกฺโก วิจาโร จิตฺตสฺเส\-กคฺคตา ชีวิตินฺทฺริยํ. เย วา ปน ตสฺมิํ สมเย อญฺเญปิ อตฺถิ ปฏิจฺจ\-สมุปฺปนฺนา อรูปิโน ธมฺมา ฐเปตฺวา เวทนา\-กฺขนฺธํ \csromanfolio{139}ฐเปตฺวา สญฺญากฺขนฺธํ ฐเปตฺวา วิญฺญาณ\-กฺขนฺธํ. อยํ ตสฺมิํ สมเย สงฺขาร\-กฺขนฺโธ โหติ ฯเปฯ. อิเม ธมฺมา อพฺยากตา.}

\vspace{\tipitakaparskip}
\csromancenter{อกุสลวิปากา มโนวิญฺญาณ\-ธาตุ.}
\csromancenter{วิปากา อพฺยากตา.}
\csromanheader{2}{อเหตุก\-กิริยา\-อพฺยากต}
\csromanmatikaanchor{mtk.71}
\csromantocmarkref{subsection}{อเหตุกกิริยาอพฺยากตกิริยามโนธาตุ}{mtk.71}
\bookmark[dest={mtk.71},level=2]{อเหตุกกิริยาอพฺยากตกิริยามโนธาตุ}
\csromanheader{2}{กิริยา\-มโนธาตุ}
\proseitem{566}{กตเม ธมฺมา อพฺยากตา. ยสฺมิํ สมเย มโนธาตุ อุปฺปนฺนา โหติ กิริยา เนว กุสลา นากุสลา น จ กมฺมวิปากา อุเปกฺขา\-สหคตา รูปารมฺมณา วา ฯเปฯ โผฏฺฐพฺพา\-รมฺมณา วา ยํ ยํ วา ปนารพฺภ. ตสฺมิํ สมเย ผสฺโส โหติ, เวทนา โหติ, สญฺญา โหติ, เจตนา โหติ, จิตฺตํ โหติ, วิตกฺโก โหติ, วิจาโร โหติ, อุเปกฺขา โหติ, จิตฺตสฺเส\-กคฺคตา โหติ, มนินฺทฺริยํ โหติ, อุเปกฺขินฺทฺริยํ โหติ, ชีวิตินฺทฺริยํ โหติ. เย วา ปน ตสฺมิํ สมเย อญฺเญปิ อตฺถิ ปฏิจฺจ\-สมุปฺปนฺนา อรูปิโน ธมฺมา. อิเม ธมฺมา อพฺยากตา ฯเปฯ.}

\prose{ตสฺมิํ โข ปน สมเย จตฺตาโร ขนฺธา โหนฺติ, ทฺวายตนานิ โหนฺติ, เทฺว ธาตุโย โหนฺติ, ตโย อาหารา โหนฺติ, ตีณินฺทฺริยานิ โหนฺติ, เอโก ผสฺโส โหติ ฯเปฯ เอกา มโนธาตุ โหติ, เอกํ ธมฺมายตนํ โหติ, เอกา ธมฺมธาตุ โหติ. เย วา ปน ตสฺมิํ สมเย อญฺเญปิ อตฺถิ ปฏิจฺจ\-สมุปฺปนฺนา อรูปิโน ธมฺมา. อิเม ธมฺมา อพฺยากตา ฯเปฯ.}

\proseitem{567}{กตโม ตสฺมิํ สมเย สงฺขาร\-กฺขนฺโธ โหติ. ผสฺโส เจตนา วิตกฺโก วิจาโร จิตฺตสฺเส\-กคฺคตา ชีวิตินฺทฺริยํ. เย วา ปน ตสฺมิํ สมเย อญฺเญปิ อตฺถิ ปฏิจฺจ\-สมุปฺปนฺนา อรูปิโน ธมฺมา ฐเปตฺวา เวทนา\-กฺขนฺธํ ฐเปตฺวา สญฺญากฺขนฺธํ ฐเปตฺวา วิญฺญาณ\-กฺขนฺธํ. อยํ ตสฺมิํ สมเย สงฺขาร\-กฺขนฺโธ โหติ ฯเปฯ. อิเม ธมฺมา อพฺยากตา.}

\vspace{\tipitakaparskip}
\csromancenter{กิริยา มโนธาตุ.}
\csromanmatikaanchor{mtk.72}
\csromantocmarkref{subsection}{กิริยามโนวิญฺญาณธาตุโสมนสฺสสหคต}{mtk.72}
\bookmark[dest={mtk.72},level=2]{กิริยามโนวิญฺญาณธาตุโสมนสฺสสหคต}
\csromanheader{2}{กิริยา\-มโนวิญฺญาณธาตุ\-โสมนสฺสสหคต}
\proseitem{568}{\csromanfolio{140}กตเม ธมฺมา อพฺยากตา. ยสฺมิํ สมเย มโนวิญฺญาณ\-ธาตุ อุปฺปนฺนา โหติ กิริยา เนว กุสลา นากุสลา น จ กมฺมวิปากา โสมนสฺส\-สหคตา รูปารมฺมณา วา ฯเปฯ ธมฺมารมฺมณา วา ยํ ยํ วา ปนารพฺภ. ตสฺมิํ สมเย ผสฺโส โหติ, เวทนา โหติ, สญฺญา โหติ, เจตนา โหติ, จิตฺตํ โหติ, วิตกฺโก โหติ, วิจาโร โหติ, ปีติ โหติ, สุขํ โหติ, จิตฺตสฺเส\-กคฺคตา โหติ, วีริยินฺทฺริยํ โหติ, สมาธินฺทฺริยํ โหติ, มนินฺทฺริยํ โหติ, โสมนสฺสิ\-นฺทฺริยํ โหติ, ชีวิตินฺทฺริยํ โหติ. เย วา ปน ตสฺมิํ สมเย อญฺเญปิ อตฺถิ ปฏิจฺจ\-สมุปฺปนฺนา อรูปิโน ธมฺมา. อิเม ธมฺมา อพฺยากตา.}

\proseitem{569}{กตโม ตสฺมิํ สมเย ผสฺโส โหติ. โย ตสฺมิํ สมเย ผสฺโส ผุสนา สมฺผุสนา สมฺผุสิตตฺตํ. อยํ ตสฺมิํ สมเย ผสฺโส โหติ ฯเปฯ.}

\proseitem{570}{กตมา ตสฺมิํ สมเย จิตฺตสฺเส\-กคฺคตา โหติ. ยา ตสฺมิํ สมเย จิตฺตสฺส ฐิติ สณฺฐิติ อวฏฺฐิติ อวิสาหาโร อวิกฺเขโป อวิสาหฏ\-มานสตา สมโถ สมาธินฺทฺริยํ สมาธิพลํ. อยํ ตสฺมิํ สมเย จิตฺตสฺเส\-กคฺคตา โหติ.}

\proseitem{571}{กตมํ ตสฺมิํ สมเย วีริยินฺทฺริยํ โหติ. โย ตสฺมิํ สมเย เจตสิโก วีริยารมฺโภ นิกฺกโม ปรกฺกโม อุยฺยาโม วายาโม อุสฺสาโห อุสฺโสฬฺหี ถาโม ธิติ อสิถิล\-ปรกฺกมตา อนิกฺขิตฺต\-ฉนฺทตา อนิกฺขิตฺตธุรตา ธุรสมฺปคฺคาโห วีริยํ วีริยินฺทฺริยํ วีริยพลํ. อิทํ ตสฺมิํ สมเย วีริยินฺทฺริยํ โหติ.}

\proseitem{572}{กตมํ ตสฺมิํ สมเย สมาธินฺทฺริยํ โหติ. ยา ตสฺมิํ สมเย จิตฺตสฺส ฐิติ สณฺฐิติ อวฏฺฐิติ อวิสาหาโร อวิกฺเขโป อวิสาหฏ\-มานสตา สมโถ สมาธินฺทฺริยํ สมาธิพลํ. อิทํ ตสฺมิํ สมเย สมาธินฺทฺริยํ โหติ ฯเปฯ. เย วา ปน ตสฺมิํ สมเย อญฺเญปิ อตฺถิ ปฏิจฺจ\-สมุปฺปนฺนา อรูปิโน ธมฺมา. อิเม ธมฺมา อพฺยากตา.}

\prose{ตสฺมิํ โข ปน สมเย จตฺตาโร ขนฺธา โหนฺติ, ทฺวายตนานิ โหนฺติ, เทฺว ธาตุโย โหนฺติ, ตโย อาหารา โหนฺติ, ปญฺจินฺทฺริยานิ โหนฺติ, \csromanfolio{141}เอโก ผสฺโส โหติ ฯเปฯ เอกา มโนวิญฺญาณ\-ธาตุ โหติ, เอกํ ธมฺมายตนํ โหติ, เอกา ธมฺมธาตุ โหติ. เย วา ปน ตสฺมิํ สมเย อญฺเญปิ อตฺถิ ปฏิจฺจ\-สมุปฺปนฺนา อรูปิโน ธมฺมา. อิเม ธมฺมา อพฺยากตา ฯเปฯ.}

\proseitem{573}{กตโม ตสฺมิํ สมเย สงฺขาร\-กฺขนฺโธ โหติ. ผสฺโส เจตนา วิตกฺโก วิจาโร ปีติ จิตฺตสฺเส\-กคฺคตา วีริยินฺทฺริยํ สมาธินฺทฺริยํ ชีวิตินฺทฺริยํ. เย วา ปน ตสฺมิํ สมเย อญฺเญปิ อตฺถิ ปฏิจฺจ\-สมุปฺปนฺนา อรูปิโน ธมฺมา ฐเปตฺวา เวทนา\-กฺขนฺธํ ฐเปตฺวา สญฺญากฺขนฺธํ ฐเปตฺวา วิญฺญาณ\-กฺขนฺธํ. อยํ ตสฺมิํ สมเย สงฺขาร\-กฺขนฺโธ โหติ ฯเปฯ. อิเม ธมฺมา อพฺยากตา.}

\vspace{\tipitakaparskip}
\csromancenter{กิริยา มโนวิญฺญาณ\-ธาตุ โสมนสฺส\-สหคตา.}
\csromanmatikaanchor{mtk.73}
\csromantocmarkref{subsection}{กิริยามโนวิญฺญาณธาตุอุเปกฺขาสหคต}{mtk.73}
\bookmark[dest={mtk.73},level=2]{กิริยามโนวิญฺญาณธาตุอุเปกฺขาสหคต}
\csromanheader{2}{กิริยา\-มโนวิญฺญาณธาตุ\-อุเปกฺขาสหคต}
\proseitem{574}{กตเม ธมฺมา อพฺยากตา. ยสฺมิํ สมเย มโนวิญฺญาณ\-ธาตุ อุปฺปนฺนา โหติ กิริยา เนว กุสลา นากุสลา น จ กมฺมวิปากา อุเปกฺขา\-สหคตา รูปารมฺมณา วา ฯเปฯ ธมฺมารมฺมณา วา ยํ ยํ วา ปนารพฺภ. ตสฺมิํ สมเย ผสฺโส โหติ, เวทนา โหติ, สญฺญา โหติ, เจตนา โหติ, จิตฺตํ โหติ, วิตกฺโก โหติ, วิจาโร โหติ, อุเปกฺขา โหติ, จิตฺตสฺเส\-กคฺคตา โหติ, วีริยินฺทฺริยํ โหติ, สมาธินฺทฺริยํ โหติ, มนินฺทฺริยํ โหติ, อุเปกฺขินฺทฺริยํ โหติ, ชีวิตินฺทฺริยํ โหติ. เย วา ปน ตสฺมิํ สมเย อญฺเญปิ อตฺถิ ปฏิจฺจ\-สมุปฺปนฺนา อรูปิโน ธมฺมา. อิเม ธมฺมา อพฺยากตา ฯเปฯ.}

\prose{ตสฺมิํ โข ปน สมเย จตฺตาโร ขนฺธา โหนฺติ, ทฺวายตนานิ โหนฺติ, เทฺว ธาตุโย โหนฺติ, ตโย อาหารา โหนฺติ, ปญฺจินฺทฺริยานิ โหนฺติ, เอโก ผสฺโส โหติ ฯเปฯ เอกา มโนวิญฺญาณ\-ธาตุ โหติ, เอกํ ธมฺมายตนํ โหติ, เอกา ธมฺมธาตุ โหติ. เย วา ปน ตสฺมิํ สมเย อญฺเญปิ อตฺถิ ปฏิจฺจ\-สมุปฺปนฺนา อรูปิโน ธมฺมา. อิเม ธมฺมา อพฺยากตา ฯเปฯ.}

\proseitem{575}{\csromanfolio{142}กตโม ตสฺมิํ สมเย สงฺขาร\-กฺขนฺโธ โหติ. ผสฺโส เจตนา วิตกฺโก วิจาโร จิตฺตสฺเส\-กคฺคตา วีริยินฺทฺริยํ สมาธินฺทฺริยํ ชีวิตินฺทฺริยํ. เย วา ปน ตสฺมิํ สมเย อญฺเญปิ อตฺถิ ปฏิจฺจ\-สมุปฺปนฺนา อรูปิโน ธมฺมา ฐเปตฺวา เวทนา\-กฺขนฺธํ ฐเปตฺวา สญฺญากฺขนฺธํ ฐเปตฺวา วิญฺญาณ\-กฺขนฺธํ. อยํ ตสฺมิํ สมเย สงฺขาร\-กฺขนฺโธ โหติ ฯเปฯ. อิเม ธมฺมา อพฺยากตา.}

\vspace{\tipitakaparskip}
\csromancenter{กิริยา มโนวิญฺญาณ\-ธาตุ อุเปกฺขา\-สหคตา.}
\csromancenter{อเหตุกา กิริยา อพฺยากตา.}
\csromanmatikaanchor{mtk.74}
\csromantocmarkref{subsection}{สเหตุกกามาวจรกิริยา}{mtk.74}
\bookmark[dest={mtk.74},level=2]{สเหตุกกามาวจรกิริยา}
\csromanheader{2}{สเหตุก\-กามาวจร\-กิริยา}
\proseitem{576}{กตเม ธมฺมา อพฺยากตา. ยสฺมิํ สมเย มโนวิญฺญาณ\-ธาตุ อุปฺปนฺนา โหติ กิริยา เนว กุสลา นากุสลา น จ กมฺมวิปากา โสมนสฺส\-สหคตา ญาณสมฺปยุตฺตา ฯเปฯ โสมนสฺส\-สหคตา ญาณสมฺปยุตฺตา สสงฺขาเรน ฯเปฯ โสมนสฺส\-สหคตา ญาณวิปฺปยุตฺตา ฯเปฯ โสมนสฺส\-สหคตา ญาณวิปฺปยุตฺตา สสงฺขาเรน ฯเปฯ อุเปกฺขา\-สหคตา ญาณสมฺปยุตฺตา ฯเปฯ อุเปกฺขา\-สหคตา ญาณสมฺปยุตฺตา สสงฺขาเรน ฯเปฯ อุเปกฺขา\-สหคตา ญาณวิปฺปยุตฺตา ฯเปฯ อุเปกฺขา\-สหคตา ญาณวิปฺปยุตฺตา สสงฺขาเรน รูปารมฺมณา วา ฯเปฯ ธมฺมารมฺมณา วา ยํ ยํ วา ปนารพฺภ. ตสฺมิํ สมเย ผสฺโส โหติ ฯเปฯ อวิกฺเขโป โหติ ฯเปฯ. อิเม ธมฺมา อพฺยากตา ฯเปฯ อโลโภ อพฺยากตมูลํ ฯเปฯ อโทโส อพฺยากตมูลํ ฯเปฯ. อิเม ธมฺมา อพฺยากตา.}

\vspace{\tipitakaparskip}
\csromancenter{สเหตุกา กามาวจร\-กิริยา.}
\csromanmatikaanchor{mtk.75}
\csromantocmarkref{subsection}{รูปาวจรกิริยา}{mtk.75}
\bookmark[dest={mtk.75},level=2]{รูปาวจรกิริยา}
\csromanheader{2}{รูปาวจร\-กิริยา}
\proseitem{577}{กตเม ธมฺมา อพฺยากตา. ยสฺมิํ สมเย รูปาวจรํ ฌานํ ภาเวติ กิริยํ เนว กุสลํ นากุสลํ น จ กมฺมวิปากํ ทิฏฺฐ\-ธมฺม\-สุข\-วิหารํ วิวิจฺเจว กาเมหิ ฯเปฯ ปฐมํ ฌานํ อุปสมฺปชฺช วิหรติ ปถวีกสิณํ. \csromanfolio{143}ตสฺมิํ สมเย ผสฺโส โหติ ฯเปฯ อวิกฺเขโป โหติ ฯเปฯ. อิเม ธมฺมา อพฺยากตา.}

\proseitem{578}{กตเม ธมฺมา อพฺยากตา. ยสฺมิํ สมเย รูปาวจรํ ฌานํ ภาเวติ กิริยํ เนว กุสลํ นากุสลํ น จ กมฺมวิปากํ ทิฏฺฐ\-ธมฺม\-สุข\-วิหารํ วิตกฺก\-วิจารานํ วูปสมา ฯเปฯ ทุติยํ ฌานํ ฯเปฯ ตติยํ ฌานํ ฯเปฯ จตุตฺถํ ฌานํ ฯเปฯ ปฐมํ ฌานํ ฯเปฯ ปญฺจมํ ฌานํ อุปสมฺปชฺช วิหรติ ปถวีกสิณํ. ตสฺมิํ สมเย ผสฺโส โหติ ฯเปฯ อวิกฺเขปฺโป โหติ ฯเปฯ. อิเม ธมฺมา อพฺยากตา.}

\vspace{\tipitakaparskip}
\csromancenter{รูปาวจร\-กิริยา.}
\csromanmatikaanchor{mtk.76}
\csromantocmarkref{subsection}{อรูปาวจรกิริยา}{mtk.76}
\bookmark[dest={mtk.76},level=2]{อรูปาวจรกิริยา}
\csromanheader{2}{อรูปาวจร\-กิริยา}
\proseitem{579}{กตเม ธมฺมา อพฺยากตา. ยสฺมิํ สมเย อรูปาวจรํ ฌานํ ภาเวติ กิริยํ เนว กุสลํ นากุสลํ น จ กมฺมวิปากํ ทิฏฺฐ\-ธมฺม\-สุข\-วิหารํ สพฺพโส รูปสญฺญานํ สมติกฺกมา ปฏิฆ\-สญฺญานํ อตฺถงฺคมา นานตฺต\-สญฺญานํ อมนสิการา อากาสา\-นญฺจา\-ยตนสญฺญาสหคตํ สุขสฺส จ ปหานา ฯเปฯ จตุตฺถํ ฌานํ อุปสมฺปชฺช วิหรติ. ตสฺมิํ สมเย ผสฺโส โหติ ฯเปฯ อวิกฺเขโป โหติ ฯเปฯ. อิเม ธมฺมา อพฺยากตา.}

\proseitem{580}{กตเม ธมฺมา อพฺยากตา. ยสฺมิํ สมเย อรูปาวจรํ ฌานํ ภาเวติ กิริยํ เนว กุสลํ นากุสลํ น จ กมฺมวิปากํ ทิฏฺฐ\-ธมฺม\-สุข\-วิหารํ สพฺพโส อากาสา\-นญฺจา\-ยตนํ สมติกฺกมฺม วิญฺญาณญฺจายตน\-สญฺญาสหคตํ สุขสฺส จ ปหานา ฯเปฯ จตุตฺถํ ฌานํ อุปสมฺปชฺช วิหรติ. ตสฺมิํ สมเย ผสฺโส โหติ ฯเปฯ อวิกฺเขโป โหติ ฯเปฯ. อิเม ธมฺมา อพฺยากตา.}

\proseitem{581}{กตเม ธมฺมา อพฺยากตา. ยสฺมิํ สมเย อรูปาวจรํ ฌานํ ภาเวติ กิริยํ เนว กุสลํ นากุสลํ น จ กมฺมวิปากํ ทิฏฺฐ\-ธมฺม\-สุข\-วิหารํ สพฺพโส วิญฺญาณญฺจา\-ยตนํ สมติกฺกมฺม อากิญฺจญฺญายตน\-สญฺญาสหคตํ สุขสฺส จ ปหานา ฯเปฯ จตุตฺถํ ฌานํ อุปสมฺปชฺช \csromanfolio{144}วิหรติ. ตสฺมิํ สมเย ผสฺโส โหติ ฯเปฯ อวิกฺเขโป โหติ ฯเปฯ. อิเม ธมฺมา อพฺยากตา.}

\proseitem{582}{กตเม ธมฺมา อพฺยากตา. ยสฺมิํ สมเย อรูปาวจรํ ฌานํ ภาเวติ กิริยํ เนว กุสลํ นากุสลํ น จ กมฺมวิปากํ ทิฏฺฐ\-ธมฺม\-สุข\-วิหารํ สพฺพโส อากิญฺจญฺญา\-ยตนํ สมติกฺกมฺม เนว\-สญฺญา\-นา\-สญฺญา\-ยตนสญฺญาสหคตํ สุขสฺส จ ปหานา ฯเปฯ จตุตฺถํ ฌานํ อุปสมฺปชฺช วิหรติ. ตสฺมิํ สมเย ผสฺโส โหติ ฯเปฯ อวิกฺเขโป โหติ ฯเปฯ. อิเม ธมฺมา อพฺยากตา ฯเปฯ อโลโภ อพฺยากตมูลํ ฯเปฯ อโทโส อพฺยากตมูลํ ฯเปฯ อโมโห อพฺยากตมูลํ ฯเปฯ. อิเม ธมฺมา อพฺยากตา.}

\vspace{\tipitakaparskip}
\csromancenter{อรูปาวจร\-กิริยา.}
\csromancenter{กิริยา อพฺยากตา.}
\nitthitammajor{จิตฺตุปฺปาท\-กณฺฑํ นิฏฺฐิตํ.}
\csromanmatikaanchor{mtk.77}
\csromantocmarknopage{chapter}{2. รูปกณฺฑ}{mtk.77}
\bookmark[dest={mtk.77},level=0]{2. รูปกณฺฑ}
\chapterheadpageread{2. \textbf{รูปกณฺฑ}}
\csromanmatikaanchor{mtk.78}
\csromantocmarkref{section}{อุทฺเทส}{mtk.78}
\bookmark[dest={mtk.78},level=1]{อุทฺเทส}
\csromanheader{1}{\textbf{อุทฺเทส}}
\proseitem{583}{\csromanfolio{145}กตเม ธมฺมา อพฺยากตา. กุสลา\-กุสลานํ ธมฺมานํ วิปากา กามาวจรา\-รูปาวจรา อรูปาวจรา อปริยาปนฺนา เวทนากฺขนฺโธ สญฺญากฺขนฺโธ สงฺขาร\-กฺขนฺโธ วิญฺญาณ\-กฺขนฺโธ. เย จ ธมฺมา กิริยา เนว กุสลา นากุสลา น จ กมฺมวิปากา. สพฺพญฺจ รูปํ อสงฺขตา จ ธาตุ. อิเม ธมฺมา อพฺยากตา.}

\csromanmatikaanchor{mtk.79}
\csromantocmarkref{section}{เอกก}{mtk.79}
\bookmark[dest={mtk.79},level=1]{เอกก}
\csromanheader{1}{\textbf{เอกก}}
\proseitem{584}{ตตฺถ กตมํ สพฺพํ รูปํ. จตฺตาโร จ มหาภูตา จตุนฺนญฺจ มหาภูตานํ อุปาทาย รูปํ. อิทํ วุจฺจติ สพฺพํ รูปํ. สพฺพํ รูปํ น เหตุ อเหตุกํ เหตุ\-วิปฺปยุตฺตํ สปฺปจฺจยํ สงฺขตํ รูปํ\csromansharedfootnote{2be01ad793d94dea}{รูปิยํ (สี)} โลกิยํ สาสวํ สํโยชนิยํ คนฺถนิยํ โอฆนิยํ โยคนิยํ นีวรณิยํ ปรามฏฺฐํ อุปาทานิยํ สํกิเลสิกํ อพฺยากตํ อนารมฺมณํ อเจตสิกํ จิตฺต\-วิปฺปยุตฺตํ เนว\-วิปาก\-น\-วิปากธมฺม\-ธมฺมํ อสํกิลิฏฺฐ\-สํกิเลสิกํ น สวิตกฺก\-สวิจารํ น อวิตกฺก\-วิจารมตฺตํ อวิตกฺก\-อวิจารํ น ปีติสหคตํ น สุขสหคตํ น อุเปกฺขา\-สหคตํ เนว ทสฺสเนน น ภาวนาย ปหาตพฺพํ เนว ทสฺสเนน น ภาวนาย ปหาตพฺพ\-เหตุกํ เนว อาจยคามิ น อปจยคามิ เนว\-เสกฺข\-นาเสกฺขํ ปริตฺตํ กามาวจรํ น รูปาวจรํ น อรูปาวจรํ ปริยาปนฺนํ โน อปริยาปนฺนํ อนิยตํ อนิยฺยานิกํ อุปฺปนฺนํ ฉหิ วิญฺญาเณหิ วิญฺเญยฺยํ อนิจฺจํ ชราภิภูตํ.}

\vspace{\tipitakaparskip}
\csromancenter{เอวํ เอกวิเธน รูปสงฺคโห.}
\csromancenter{\textbf{เอกก}ํ.}
\csromancenter{\textbf{ทุก} ทุวิเธน รูปสงฺคโห\csromandash{}}
\csromanmatikaanchor{mtk.80}
\csromantocmarkref{section}{ทุก}{mtk.80}
\bookmark[dest={mtk.80},level=1]{ทุก}
\prose{\csromanfolio{146}อตฺถิ รูปํ อุปาทา, อตฺถิ รูปํ โน อุปาทา.}

\prose{อตฺถิ รูปํ อุปาทิณฺณํ, อตฺถิ รูปํ อนุปาทิณฺณํ.}

\prose{อตฺถิ รูปํ อุปาทิณฺณุ\-ปาทานิยํ, อตฺถิ รูปํ อนุปาทิณฺณุ\-ปาทานิยํ.}

\prose{อตฺถิ รูปํ สนิทสฺสนํ, อตฺถิ รูปํ อนิทสฺสนํ.}

\prose{อตฺถิ รูปํ สปฺปฏิฆํ, อตฺถิ รูปํ อปฺปฏิฆํ.}

\prose{อตฺถิ รูปํ อินฺทฺริยํ, อตฺถิ รูปํ น อินฺทฺริยํ.}

\prose{อตฺถิ รูปํ มหาภูตํ, อตฺถิ รูปํ น มหาภูตํ.}

\prose{อตฺถิ รูปํ วิญฺญตฺติ, อตฺถิ รูปํ น วิญฺญตฺติ.}

\prose{อตฺถิ รูปํ จิตฺต\-สมุฏฺฐานํ, อตฺถิ รูปํ น จิตฺต\-สมุฏฺฐานํ.}

\prose{อตฺถิ รูปํ จิตฺตสหภุ, อตฺถิ รูปํ น จิตฺตสหภุ.}

\prose{อตฺถิ รูปํ จิตฺตา\-นุปริวตฺติ, อตฺถิ รูปํ น จิตฺตา\-นุปริวตฺติ.}

\prose{อตฺถิ รูปํ อชฺฌตฺติกํ, อตฺถิ รูปํ พาหิรํ.}

\prose{อตฺถิ รูปํ โอฬาริกํ, อตฺถิ รูปํ สุขุมํ.}

\prose{อตฺถิ รูปํ ทูเร, อตฺถิ รูปํ สนฺติเก.}

\prose{อตฺถิ รูปํ จกฺขุ\-สมฺผสฺสสฺส วตฺถุ, อตฺถิ รูปํ จกฺขุ\-สมฺผสฺสสฺส น วตฺถุ. อตฺถิ รูปํ จกฺขุ\-สมฺผสฺสชาย เวทนาย ฯเปฯ สญฺญาย ฯเปฯ เจตนาย ฯเปฯ จกฺขุ\-วิญฺญาณสฺส วตฺถุ, อตฺถิ รูปํ จกฺขุ\-วิญฺญาณสฺส น วตฺถุ.}

\prose{อตฺถิ รูปํ โสต\-สมฺผสฺสสฺส ฯเปฯ ฆาน\-สมฺผสฺสสฺส ฯเปฯ ชิวฺหา\-สมฺผสฺสสฺส ฯเปฯ กาย\-สมฺผสฺสสฺส วตฺถุ, อตฺถิ รูปํ กาย\-สมฺผสฺสสฺส น วตฺถุ. อตฺถิ รูปํ กาย\-สมฺผสฺสชาย เวทนาย ฯเปฯ สญฺญาย ฯเปฯ เจตนาย ฯเปฯ กายวิญฺญาณสฺส วตฺตุ, อตฺถิ รูปํ กายวิญฺญาณสฺส น วตฺถุ.}

\prose{อตฺถิ รูปํ จกฺขุ\-สมฺผสฺสสฺส อารมฺมณํ, อตฺถิ รูปํ จกฺขุ\-สมฺผสฺสสฺส นารมฺมณํ. อตฺถิ รูปํ จกฺขุ\-สมฺผสฺสชาย เวทนาย ฯเปฯ สญฺญาย ฯเปฯ เจตนาย ฯเปฯ จกฺขุ\-วิญฺญาณสฺส อารมฺมณํ, อตฺถิ รูปํ จกฺขุ\-วิญฺญาณสฺส นารมฺมณํ.}

\chapterheadpageread{2. รูปกณฺฑ}
\prose{\csromanfolio{147}อตฺถิ รูปํ โสต\-สมฺผสฺสสฺส ฯเปฯ ฆาน\-สมฺผสฺสสฺส ฯเปฯ ชิวฺหา\-สมฺผสฺสสฺส ฯเปฯ กาย\-สมฺผสฺสสฺส อารมฺมณํ, อตฺถิ รูปํ กาย\-สมฺผสฺสสฺส นารมฺมณํ. อตฺถิ รูปํ กาย\-สมฺผสฺสชาย เวทนาย ฯเปฯ สญฺญาย ฯเปฯ เจตนาย ฯเปฯ กายวิญฺญาณสฺส อารมฺมณํ, อตฺถิ รูปํ กายวิญฺญาณสฺส นารมฺมณํ.}

\prose{อตฺถิ รูปํ จกฺขายตนํ, อตฺถิ รูปํ น จกฺขายตนํ. อตฺถิ รูปํ โสตายตนํ ฯเปฯ ฆานายตนํ ฯเปฯ ชิวฺหายตนํ ฯเปฯ กายายตนํ, อตฺถิ รูปํ น กายายตนํ.}

\prose{อตฺถิ รูปํ รูปายตนํ, อตฺถิ รูปํ น รูปายตนํ. อตฺถิ รูปํ สทฺทายตนํ ฯเปฯ คนฺธายตนํ ฯเปฯ รสายตนํ ฯเปฯ โผฏฺฐพฺพา\-ยตนํ, อตฺถิ รูปํ น โผฏฺฐพฺพา\-ยตนํ.}

\prose{อตฺถิ รูปํ จกฺขุธาตุ, อตฺถิ รูปํ น จกฺขุธาตุ. อตฺถิ รูปํ โสตธาตุ ฯเปฯ ฆานธาตุ ฯเปฯ ชิวฺหาธาตุ ฯเปฯ กายธาตุ, อตฺถิ รูปํ น กายธาตุ.}

\prose{อตฺถิ รูปํ รูปธาตุ, อตฺถิ รูปํ น รูปธาตุ. อตฺถิ รูปํ สทฺทธาตุ ฯเปฯ คนฺธธาตุ ฯเปฯ รสธาตุ ฯเปฯ โผฏฺฐพฺพ\-ธาตุ, อตฺถิ รูปํ น โผฏฺฐพฺพ\-ธาตุ.}

\prose{อตฺถิ รูปํ จกฺขุนฺทฺริยํ, อตฺถิ รูปํ น จกฺขุนฺทฺริยํ. อตฺถิ รูปํ โสตินฺทฺริยํ ฯเปฯ ฆานินฺทฺริยํ ฯเปฯ ชิวฺหินฺทฺริยํ ฯเปฯ กายินฺทฺริยํ, อตฺถิ รูปํ น กายินฺทฺริยํ.}

\prose{อตฺถิ รูปํ อิตฺถินฺทฺริยํ, อตฺถิ รูปํ น อิตฺถินฺทฺริยํ.}

\prose{อตฺถิ รูปํ ปุริสินฺทฺริยํ, อตฺถิ รูปํ น ปุริสินฺทฺริยํ.}

\prose{อตฺถิ รูปํ ชีวิตินฺทฺริยํ, อตฺถิ รูปํ น ชีวิตินฺทฺริยํ.}

\prose{อตฺถิ รูปํ กายวิญฺญตฺติ, อตฺถิ รูปํ น กายวิญฺญตฺติ.}

\prose{อตฺถิ รูปํ วจีวิญฺญตฺติ, อตฺถิ รูปํ น วจีวิญฺญตฺติ.}

\prose{อตฺถิ รูปํ อากาสธาตุ, อตฺถิ รูปํ น อากาสธาตุ.}

\prose{อตฺถิ รูปํ อาโปธาตุ, อตฺถิ รูปํ น อาโปธาตุ.}

\prose{อตฺถิ รูปํ รูปสฺส ลหุตา, อตฺถิ รูปํ รูปสฺส น ลหุตา.}

\prose{อตฺถิ รูปํ รูปสฺส มุทุตา, อตฺถิ รูปํ รูปสฺส น มุทุตา.}

\prose{อตฺถิ รูปํ รูปสฺส กมฺมญฺญตา, อตฺถิ รูปํ รูปสฺส น กมฺมญฺญตา.}

\prose{อตฺถิ รูปํ รูปสฺส อุปจโย, อตฺถิ รูปํ รูปสฺส น อุปจโย.}

\prose{อตฺถิ รูปํ รูปสฺส สนฺตติ, อตฺถิ รูปํ รูปสฺส น สนฺตติ.}

\csromanmatikaanchor{mtk.81}
\csromantocmarkref{section}{ติก}{mtk.81}
\bookmark[dest={mtk.81},level=1]{ติก}
\prose{\csromanfolio{148}อตฺถิ รูปํ รูปสฺส ชรตา, อตฺถิ รูปํ รูปสฺส น ชรตา.}

\prose{อตฺถิ รูปํ รูปสฺส อนิจฺจตา, อตฺถิ รูปํ รูปสฺส น อนิจฺจตา.}

\prose{อตฺถิ รูปํ กพฬีกาโร อาหาโร, อตฺถิ รูปํ น กพฬีกาโร อาหาโร.}

\vspace{\tipitakaparskip}
\csromancenter{เอวํ ทุวิเธน รูปสงฺคโห.}
\csromancenter{ทุกํ.}
\csromancenter{\textbf{ติก} ติวิเธน รูปสงฺคโห\csromandash{}}
\proseitem{585}{ยํ ตํ รูปํ อชฺฌตฺติกํ ตํ อุปาทา, ยํ ตํ รูปํ พาหิรํ ตํ อตฺถิ อุปาทา, อตฺถิ โน อุปาทา.}

\prose{ยํ ตํ รูปํ อชฺฌตฺติกํ ตํ อุปาทิณฺณํ, ยํ ตํ รูปํ พาหิรํ ตํ อตฺถิ อุปาทิณฺณํ, อตฺถิ อนุปาทิณฺณํ.}

\prose{ยํ ตํ รูปํ อชฺฌตฺติกํ ตํ อุปาทิณฺณุ\-ปาทานิยํ, ยํ ตํ รูปํ พาหิรํ ตํ อตฺถิ อุปาทิณฺณุ\-ปาทานิยํ, อตฺถิ อนุปาทิณฺณุ\-ปาทานิยํ.}

\prose{ยํ ตํ รูปํ อชฺฌตฺติกํ ตํ อนิทสฺสนํ, ยํ ตํ รูปํ พาหิรํ ตํ อตฺถิ สนิทสฺสนํ, อตฺถิ อนิทสฺสนํ.}

\prose{ยํ ตํ รูปํ อชฺฌตฺติกํ ตํ สปฺปฏิฆํ, ยํ ตํ รูปํ พาหิรํ ตํ อตฺถิ สปฺปฏิฆํ, อตฺถิ อปฺปฏิฆํ.}

\prose{ยํ ตํ รูปํ อชฺฌตฺติกํ ตํ อินฺทฺริยํ, ยํ ตํ รูปํ พาหิรํ ตํ อตฺถิ อินฺทฺริยํ, อตฺถิ น อินฺทฺริยํ.}

\prose{ยํ ตํ รูปํ อชฺฌตฺติกํ ตํ น มหาภูตํ, ยํ ตํ รูปํ พาหิรํ ตํ อตฺถิ มหาภูตํ, อตฺถิ น มหาภูตํ.}

\prose{ยํ ตํ รูปํ อชฺฌตฺติกํ ตํ น วิญฺญตฺติ, ยํ ตํ รูปํ พาหิรํ ตํ อตฺถิ วิญฺญตฺติ, อตฺถิ น วิญฺญตฺติ.}

\chapterheadpageread{2. รูปกณฺฑ}
\prose{\csromanfolio{149}ยํ ตํ รูปํ อชฺฌตฺติกํ ตํ น จิตฺต\-สมุฏฺฐานํ, ยํ ตํ รูปํ พาหิรํ ตํ อตฺถิ จิตฺต\-สมุฏฺฐานํ, อตฺถิ น จิตฺต\-สมุฏฺฐานํ.}

\prose{ยํ ตํ รูปํ อชฺฌตฺติกํ ตํ น จิตฺตสหภุ, ยํ ตํ รูปํ พาหิรํ ตํ อตฺถิ จิตฺตสหภุ, อตฺถิ น จิตฺตสหภุ.}

\prose{ยํ ตํ รูปํ อชฺฌตฺติกํ ตํ น จิตฺตา\-นุปริวตฺติ, ยํ ตํ รูปํ พาหิรํ ตํ อตฺถิ จิตฺตา\-นุปริวตฺติ, อตฺถิ น จิตฺตา\-นุปริวตฺติ.}

\prose{ยํ ตํ รูปํ อชฺฌตฺติกํ ตํ โอฬาริกํ, ยํ ตํ รูปํ พาหิรํ ตํ อตฺถิ โอฬาริกํ, อตฺถิ สุขุมํ.}

\prose{ยํ ตํ รูปํ อชฺฌตฺติกํ ตํ สนฺติเก, ยํ ตํ รูปํ พาหิรํ ตํ อตฺถิ ทูเร, อตฺถิ สนฺติเก.}

\prose{ยํ ตํ รูปํ พาหิรํ ตํ จกฺขุ\-สมฺผสฺสสฺส น วตฺถุ, ยํ ตํ รูปํ อชฺฌตฺติกํ ตํ อตฺถิ จกฺขุ\-สมฺผสฺสสฺส วตฺถุ, อตฺถิ จกฺขุ\-สมฺผสฺสสฺส น วตฺถุ.}

\prose{ยํ ตํ รูปํ พาหิรํ ตํ จกฺขุ\-สมฺผสฺสชาย เวทนาย ฯเปฯ สญฺญาย ฯเปฯ เจตนาย ฯเปฯ จกฺขุ\-วิญฺญาณสฺส น วตฺถุ, ยํ ตํ รูปํ อชฺฌตฺติกํ ตํ อตฺถิ จกฺขุ\-วิญฺญาณสฺส วตฺถุ, อตฺถิ จกฺขุ\-วิญฺญาณสฺส น วตฺถุ.}

\prose{ยํ ตํ รูปํ พาหิรํ ตํ โสต\-สมฺผสฺสสฺส ฯเปฯ ฆาน\-สมฺผสฺสสฺส ฯเปฯ ชิวฺหา\-สมฺผสฺสสฺส ฯเปฯ กาย\-สมฺผสฺสสฺส น วตฺถุ, ยํ ตํ รูปํ อชฺฌตฺติกํ ตํ อตฺถิ กาย\-สมฺผสฺสสฺส วตฺถุ, อตฺถิ กาย\-สมฺผสฺสสฺส น วตฺถุ. ยํ ตํ รูปํ พาหิรํ ตํ กาย\-สมฺผสฺสชาย เวทนาย ฯเปฯ สญฺญาย ฯเปฯ เจตนาย ฯเปฯ กายวิญฺญาณสฺส น วตฺถุ, ยํ ตํ รูปํ อชฺฌตฺติกํ ตํ อตฺถิ กายวิญฺญาณสฺส วตฺถุ, อตฺถิ กายวิญฺญาณสฺส น วตฺถุ.}

\prose{ยํ ตํ รูปํ อชฺฌตฺติกํ ตํ จกฺขุ\-สมฺผสฺสสฺส นารมฺมณํ, ยํ ตํ รูปํ พาหิรํ ตํ อตฺถิ จกฺขุ\-สมฺผสฺสสฺส อารมฺมณํ, อตฺถิ จกฺขุ\-สมฺผสฺสสฺส นารมฺมณํ. ยํ ตํ รูปํ อชฺฌตฺติกํ ตํ จกฺขุ\-สมฺผสฺสชาย เวทนาย ฯเปฯ สญฺญาย ฯเปฯ เจตนาย ฯเปฯ จกฺขุ\-วิญฺญาณสฺส นารมฺมณํ, ยํ ตํ รูปํ พาหิรํ ตํ อตฺถิ จกฺขุ\-วิญฺญาณสฺส อารมฺมณํ, อตฺถิ จกฺขุ\-วิญฺญาณสฺส นารมฺมณํ.}

\prose{ยํ ตํ รูปํ อชฺฌตฺติกํ ตํ โสต\-สมฺผสฺสสฺส ฯเปฯ ฆาน\-สมฺผสฺสสฺส ฯเปฯ ชิวฺหา\-สมฺผสฺสสฺส ฯเปฯ กาย\-สมฺผสฺสสฺส นารมฺมณํ, ยํ ตํ รูปํ พาหิรํ ตํ อตฺถิ กาย\-สมฺผสฺสสฺส อารมฺมณํ, อตฺถิ กาย\-สมฺผสฺสสฺส นารมฺมณํ. ยํ ตํ รูปํ}

\prose{\csromanfolio{150}อชฺฌตฺติกํ ตํ กาย\-สมฺผสฺสชาย เวทนาย ฯเปฯ สญฺญาย ฯเปฯ เจตนาย ฯเปฯ กายวิญฺญาณสฺส นารมฺมณํ, ยํ ตํ รูปํ พาหิรํ ตํ อตฺถิ กายวิญฺญาณสฺส อารมฺมณํ, อตฺถิ กายวิญฺญาณสฺส นารมฺมณํ.}

\prose{ยํ ตํ รูปํ พาหิรํ ตํ น จกฺขายตนํ, ยํ ตํ รูปํ อชฺฌตฺติกํ ตํ อตฺถิ จกฺขายตนํ, อตฺถิ น จกฺขายตนํ. ยํ ตํ รูปํ พาหิรํ ตํ น โสตายตนํ ฯเปฯ น ฆานายตนํ ฯเปฯ น ชิวฺหายตนํ ฯเปฯ น กายายตนํ, ยํ ตํ รูปํ อชฺฌตฺติกํ ตํ อตฺถิ กายายตนํ, อตฺถิ น กายายตนํ.}

\prose{ยํ ตํ รูปํ อชฺฌตฺติกํ ตํ น รูปายตนํ, ยํ ตํ รูปํ พาหิรํ ตํ อตฺถิ รูปายตนํ, อตฺถิ น รูปายตนํ. ยํ ตํ รูปํ อชฺฌตฺติกํ ตํ น สทฺทายตนํ ฯเปฯ น คนฺธายตนํ ฯเปฯ น รสายตนํ ฯเปฯ น โผฏฺฐพฺพา\-ยตนํ, ยํ ตํ รูปํ พาหิรํ ตํ อตฺถิ โผฏฺฐพฺพา\-ยตนํ, อตฺถิ น โผฏฺฐพฺพา\-ยตนํ.}

\prose{ยํ ตํ รูปํ พาหิรํ ตํ น จกฺขุธาตุ, ยํ ตํ รูปํ อชฺฌตฺติกํ ตํ อตฺถิ จกฺขุธาตุ, อตฺถิ น จกฺขุธาตุ. ยํ ตํ รูปํ พาหิรํ ตํ น โสตธาตุ ฯเปฯ น ฆานธาตุ ฯเปฯ น ชิวฺหาธาตุ ฯเปฯ น กายธาตุ, ยํ ตํ รูปํ อชฺฌตฺติกํ ตํ อตฺถิ กายธาตุ, อตฺถิ น กายธาตุ.}

\prose{ยํ ตํ รูปํ อชฺฌตฺติกํ ตํ น รูปธาตุ, ยํ ตํ รูปํ พาหิรํ ตํ อตฺถิ รูปธาตุ, อตฺถิ น รูปธาตุ. ยํ ตํ รูปํ อชฺฌตฺติกํ ตํ น สทฺทธาตุ ฯเปฯ น คนฺธธาตุ ฯเปฯ น รสธาตุ ฯเปฯ น โผฏฺฐพฺพ\-ธาตุ, ยํ ตํ รูปํ พาหิรํ ตํ อตฺถิ โผฏฺฐพฺพ\-ธาตุ, อตฺถิ น โผฏฺฐพฺพ\-ธาตุ.}

\prose{ยํ ตํ รูปํ พาหิรํ ตํ น จกฺขุนฺทฺริยํ, ยํ ตํ รูปํ อชฺฌตฺติกํ ตํ อตฺถิ จกฺขุนฺทฺริยํ, อตฺถิ น จกฺขุนฺทฺริยํ. ยํ ตํ รูปํ พาหิรํ ตํ น โสตินฺทฺริยํ ฯเปฯ น ฆานินฺทฺริยํ ฯเปฯ น ชิวฺหินฺทฺริยํ ฯเปฯ น กายินฺทฺริยํ, ยํ ตํ รูปํ อชฺฌตฺติกํ ตํ อตฺถิ กายินฺทฺริยํ, อตฺถิ น กายินฺทฺริย.}

\prose{ยํ ตํ รูปํ อชฺฌตฺติกํ ตํ น อิตฺถินฺทฺริยํ, ยํ ตํ รูปํ พาหิรํ ตํ อตฺถิ อิตฺถินฺทฺริยํ, อตฺถิ น อิตฺถินฺทฺริยํ.}

\prose{ยํ ตํ รูปํ อชฺฌตฺติกํ ตํ น ปุริสินฺทฺริยํ, ยํ ตํ รูปํ พาหิรํ ตํ อตฺถิ ปุริสินฺทฺริยํ, อตฺถิ น ปุริสินฺทฺริยํ.}

\prose{ยํ ตํ รูปํ อชฺฌตฺติกํ ตํ น ชีวิตินฺทฺริยํ, ยํ ตํ รูปํ พาหิรํ ตํ อตฺถิ ชีวิตินฺทฺริยํ, อตฺถิ น ชีวิตินฺทฺริยํ.}

\chapterheadpageread{2. รูปกณฺฑ}
\prose{\csromanfolio{151}ยํ ตํ รูปํ อชฺฌตฺติกํ ตํ น กายวิญฺญตฺติ, ยํ ตํ รูปํ พาหิรํ ตํ อตฺถิ กายวิญฺญตฺติ, อตฺถิ น กายวิญฺญตฺติ.}

\prose{ยํ ตํ รูปํ อชฺฌตฺติกํ ตํ น วจีวิญฺญตฺติ, ยํ ตํ รูปํ พาหิรํ ตํ อตฺถิ วจีวิญฺญตฺติ, อตฺถิ น วจีวิญฺญตฺติ.}

\prose{ยํ ตํ รูปํ อชฺฌตฺติกํ ตํ น อากาสธาตุ, ยํ ตํ รูปํ พาหิรํ ตํ อตฺถิ อากาสธาตุ, อตฺถิ น อากาสธาตุ.}

\prose{ยํ ตํ รูปํ อชฺฌตฺติกํ ตํ น อาโปธาตุ, ยํ ตํ รูปํ พาหิรํ ตํ อตฺถิ อาโปธาตุ, อตฺถิ น อาโปธาตุ.}

\prose{ยํ ตํ รูปํ อชฺฌตฺติกํ ตํ รูปสฺส น ลหุตา, ยํ ตํ รูปํ พาหิรํ ตํ อตฺถิ รูปสฺส ลหุตา, อตฺถิ รูปสฺส น ลหุตา.}

\prose{ยํ ตํ รูปํ อชฺฌตฺติกํ ตํ รูปสฺส น มุทุตา, ยํ ตํ รูปํ พาหิรํ ตํ อตฺถิ รูปสฺส มุทุตา, อตฺถิ รูปสฺส น มุทุตา.}

\prose{ยํ ตํ รูปํ อชฺฌตฺติกํ ตํ รูปสฺส น กมฺมญฺญตา, ยํ ตํ รูปํ พาหิรํ ตํ อตฺถิ รูปสฺส กมฺมญฺญตา, อตฺถิ รูปสฺส น กมฺมญฺญตา.}

\prose{ยํ ตํ รูปํ อชฺฌตฺติกํ ตํ รูปสฺส น อุปจโย, ยํ ตํ รูปํ พาหิรํ ตํ อตฺถิ รูปสฺส อุปจโย, อตฺถิ รูปสฺส น อุปจโย.}

\prose{ยํ ตํ รูปํ อชฺฌตฺติกํ ตํ รูปสฺส น สนฺตติ, ยํ ตํ รูปํ พาหิรํ ตํ อตฺถิ รูปสฺส สนฺตติ, อตฺถิ รูปสฺส น สนฺตติ.}

\prose{ยํ ตํ รูปํ อชฺฌตฺติกํ ตํ รูปสฺส น ชรตา, ยํ ตํ รูปํ พาหิรํ ตํ อตฺถิ รูปสฺส ชรตา, อตฺถิ รูปสฺส น ชรตา.}

\prose{ยํ ตํ รูปํ อชฺฌตฺติกํ ตํ รูปสฺส น อนิจฺจตา, ยํ ตํ รูปํ พาหิรํ ตํ อตฺถิ รูปสฺส อนิจฺจตา, อตฺถิ รูปสฺส น อนิจฺจตา.}

\prose{ยํ ตํ รูปํ อชฺฌตฺติกํ ตํ น กพฬีกาโร อาหาโร, ยํ ตํ รูปํ พาหิรํ ตํ อตฺถิ กพฬีกาโร อาหาโร, อตฺถิ น กพฬีกาโร อาหาโร.}

\vspace{\tipitakaparskip}
\csromancenter{เอวํ ติวิเธน รูปสงฺคโห.}
\csromancenter{ติกํ.}
\csromancenter{\textbf{จตุกฺก} จตุพฺพิเธน รูปสงฺคโห\csromandash{}}
\csromanmatikaanchor{mtk.82}
\csromantocmarkref{section}{จตุกฺก}{mtk.82}
\bookmark[dest={mtk.82},level=1]{จตุกฺก}
\proseitem{586}{\csromanfolio{152}ยํ ตํ รูปํ อุปาทา ตํ อตฺถิ อุปาทิณฺณํ, อตฺถิ อนุปาทิณฺณํ. ยํ ตํ รูปํ โน อุปาทา ตํ อตฺถิ อุปาทิณฺณํ, อตฺถิ อนุปาทิณฺณํ.}

\prose{ยํ ตํ รูปํ อุปาทา ตํ อตฺถิ อุปาทิณฺณุ\-ปาทานิยํ, อตฺถิ อนุปาทิณฺณุ\-ปาทานิยํ. ยํ ตํ รูปํ โน อุปาทา ตํ อตฺถิ อุปาทิณฺณุ\-ปาทานิยํ, อตฺถิ อนุปาทิณฺณุ\-ปาทานิยํ.}

\prose{ยํ ตํ รูปํ อุปาทา ตํ อตฺถิ สปฺปฏิฆํ, อตฺถิ อปฺปฏิฆํ. ยํ ตํ รูปํ โน อุปาทา ตํ อตฺถิ สปฺปฏิฆํ, อตฺถิ อปฺปฏิฆํ.}

\prose{ยํ ตํ รูปํ อุปาทา ตํ อตฺถิ โอฬาริกํ, อตฺถิ สุขุมํ. ยํ ตํ รูปํ โน อุปาทา ตํ อตฺถิ โอฬาริกํ, อตฺถิ สุขุมํ.}

\prose{ยํ ตํ รูปํ อุปาทา ตํ อตฺถิ ทูเร, อตฺถิ สนฺติเก. ยํ ตํ รูปํ โน อุปาทา ตํ อตฺถิ ทูเร, อตฺถิ สนฺติเก.}

\prose{ยํ ตํ รูปํ อุปาทิณฺณํ ตํ อตฺถิ สนิทสฺสนํ, อตฺถิ อนิทสฺสนํ. ยํ ตํ รูปํ อนุปาทิณฺณํ ตํ อตฺถิ สนิทสฺสนํ, อตฺถิ อนิทสฺสนํ.}

\prose{ยํ ตํ รูปํ อุปาทิณฺณํ ตํ อตฺถิ สปฺปฏิฆํ, อตฺถิ อปฺปฏิฆํ. ยํ ตํ รูปํ อนุปาทิณฺณํ ตํ อตฺถิ สปฺปฏิฆํ, อตฺถิ อปฺปฏิฆํ.}

\prose{ยํ ตํ รูปํ อุปาทิณฺณํ ตํ อตฺถิ มหาภูตํ, อตฺถิ น มหาภูตํ. ยํ ตํ รูปํ อนุปาทิณฺณํ ตํ อตฺถิ มหาภูตํ, อตฺถิ น มหาภูตํ.}

\prose{ยํ ตํ รูปํ อุปาทิณฺณํ ตํ อตฺถิ โอฬาริกํ, อตฺถิ สุขุมํ. ยํ ตํ รูปํ อนุปาทิณฺณํ ตํ อตฺถิ โอฬาริกํ, อตฺถิ สุขุมํ.}

\prose{ยํ ตํ รูปํ อุปาทิณฺณํ ตํ อตฺถิ ทูเร, อตฺถิ สนฺติเก. ยํ ตํ รูปํ อนุปาทิณฺณํ ตํ อตฺถิ ทูเร, อตฺถิ สนฺติเก.}

\chapterheadpageread{2. รูปกณฺฑ}
\prose{\csromanfolio{153}ยํ ตํ รูปํ อุปาทิณฺณุ\-ปาทานิยํ ตํ อตฺถิ สนิทสฺสนํ, อตฺถิ อนิทสฺสนํ. ยํ ตํ รูปํ อนุปาทิณฺณุ\-ปาทานิยํ ตํ อตฺถิ สนิทสฺสนํ, อตฺถิ อนิทสฺสนํ.}

\prose{ยํ ตํ รูปํ อุปาทิณฺณุ\-ปาทานิยํ ตํ อตฺถิ สปฺปฏิฆํ, อตฺถิ อปฺปฏิฆํ. ยํ ตํ รูปํ อนุปาทิณฺณุ\-ปาทานิยํ ตํ อตฺถิ สปฺปฏิฆํ, อตฺถิ อปฺปฏิฆํ.}

\prose{ยํ ตํ รูปํ อุปาทิณฺณุ\-ปาทานิยํ ตํ อตฺถิ มหาภูตํ, อตฺถิ น มหาภูตํ. ยํ ตํ รูปํ อนุปาทิณฺณุ\-ปาทานิยํ ตํ อตฺถิ มหาภูตํ, อตฺถิ น มหาภูตํ.}

\prose{ยํ ตํ รูปํ อุปาทิณฺณุ\-ปาทานิยํ ตํ อตฺถิ โอฬาริกํ, อตฺถิ สุขุมํ. ยํ ตํ รูปํ อนุปาทิณฺณุ\-ปาทานิยํ ตํ อตฺถิ โอฬาริกํ, อตฺถิ สุขุมํ.}

\prose{ยํ ตํ รูปํ อุปาทิณฺณุ\-ปาทานิยํ ตํ อตฺถิ ทูเร, อตฺถิ สนฺติเก. ยํ ตํ รูปํ อนุปาทิณฺณุ\-ปาทานิยํ ตํ อตฺถิ ทูเร, อตฺถิ สนฺติเก.}

\prose{ยํ ตํ รูปํ สปฺปฏิฆํ ตํ อตฺถิ อินฺทฺริยํ, อตฺถิ น อินฺทฺริยํ. ยํ ตํ รูปํ อปฺปฏิฆํ ตํ อตฺถิ อินฺทฺริยํ, อตฺถิ น อินฺทฺริยํ.}

\prose{ยํ ตํ รูปํ สปฺปฏิฆํ ตํ อตฺถิ มหาภูตํ, อตฺถิ น มหาภูตํ. ยํ ตํ รูปํ อปฺปฏิฆํ ตํ อตฺถิ มหาภูตํ, อตฺถิ น มหาภูตํ.}

\prose{ยํ ตํ รูปํ อินฺทฺริยํ ตํ อตฺถิ โอฬาริกํ, อตฺถิ สุขุมํ. ยํ ตํ รูปํ น อินฺทฺริยํ ตํ อตฺถิ โอฬาริกํ, อตฺถิ สุขุมํ.}

\prose{ยํ ตํ รูปํ อินฺทฺริยํ ตํ อตฺถิ ทูเร, อตฺถิ สนฺติเก. ยํ ตํ รูปํ น อินฺทฺริยํ ตํ อตฺถิ ทูเร, อตฺถิ สนฺติเก.}

\prose{ยํ ตํ รูปํ มหาภูตํ ตํ อตฺถิ โอฬาริกํ, อตฺถิ สุขุมํ. ยํ ตํ รูปํ น มหาภูตํ ตํ อตฺถิ โอฬาริกํ, อตฺถิ สุขุมํ.}

\prose{ยํ ตํ รูปํ มหาภูตํ ตํ อตฺถิ ทูเร, อตฺถิ สนฺติเก. ยํ ตํ รูปํ น มหาภูตํ ตํ อตฺถิ ทูเร, อตฺถิ สนฺติเก.}

\prose{ทิฏฺฐํ สุตํ มุตํ วิญฺญาตํ รูปํ. เอวํ จตุพฺพิเธน รูปสงฺคโห. จตุกฺกํ.}

\vspace{\tipitakaparskip}
\csromancenter{\textbf{ปญฺจก} ปญฺจวิเธน รูปสงฺคโห\csromandash{}}
\csromanmatikaanchor{mtk.83}
\csromanmatikaanchor{mtk.84}
\csromanmatikaanchor{mtk.85}
\csromanmatikaanchor{mtk.86}
\csromantocmarkref{section}{ปญฺจก}{mtk.83}
\bookmark[dest={mtk.83},level=1]{ปญฺจก}
\csromantocmarkref{section}{ฉกฺก}{mtk.84}
\bookmark[dest={mtk.84},level=1]{ฉกฺก}
\csromantocmarkref{section}{สตฺตก}{mtk.85}
\bookmark[dest={mtk.85},level=1]{สตฺตก}
\csromantocmarkref{section}{อฏฺฐก}{mtk.86}
\bookmark[dest={mtk.86},level=1]{อฏฺฐก}
\proseitem{587}{\csromanfolio{154}ปถวีธาตุ อาโปธาตุ เตโชธาตุ วาโยธาตุ, ยญฺจ รูปํ อุปาทา. เอวํ ปญฺจวิเธน รูปสงฺคโห. ปญฺจกํ.}

\vspace{\tipitakaparskip}
\csromancenter{\textbf{ฉกฺก} ฉพฺพิเธน รูปสงฺคโห\csromandash{}}
\proseitem{588}{จกฺขุ\-วิญฺเญยฺยํ รูปํ โสตวิญฺเญยฺยํ รูปํ ฆานวิญฺเญยฺยํ รูปํ ชิวฺหาวิญฺเญยฺยํ รูปํ กายวิญฺเญยฺยํ รูปํ มโนวิญฺเญยฺยํ รูปํ. เอวํ ฉพฺพิเธน รูปสงฺคโห. ฉกฺกํ.}

\vspace{\tipitakaparskip}
\csromancenter{\textbf{สตฺตก} สตฺตวิเธน รูปสงฺคโห\csromandash{}}
\proseitem{589}{จกฺขุ\-วิญฺเญยฺยํ รูปํ โสตวิญฺเญยฺยํ รูปํ ฆานวิญฺเญยฺยํ รูปํ ชิวฺหาวิญฺเญยฺยํ รูปํ กายวิญฺเญยฺยํ รูปํ มโนธาตุ\-วิญฺเญยฺยํ รูปํ มโนวิญฺญาณ\-ธาตุวิญฺเญยฺยํ รูปํ. เอวํ สตฺตวิเธน รูปสงฺคโห. สตฺตกํ.}

\vspace{\tipitakaparskip}
\csromancenter{\textbf{อฏฺฐก} อฏฺฐวิเธน รูปสงฺคโห\csromandash{}}
\proseitem{590}{จกฺขุ\-วิญฺเญยฺยํ รูปํ โสตวิญฺเญยฺยํ รูปํ ฆานวิญฺเญยฺยํ รูปํ ชิวฺหาวิญฺเญยฺยํ รูปํ กายวิญฺเญยฺยํ รูปํ อตฺถิ สุขสมฺผสฺสํ อตฺถิ ทุกฺข\-สมฺผสฺสํ มโนธาตุ\-วิญฺเญยฺยํ รูปํ มโนวิญฺญาณ\-ธาตุวิญฺเญยฺยํ รูปํ. เอวํ อฏฺฐวิเธน รูปสงฺคโห. อฏฺฐกํ.}

\csromanmatikaanchor{mtk.87}
\csromantocmarkref{section}{นวก}{mtk.87}
\bookmark[dest={mtk.87},level=1]{นวก}
\csromanheader{1}{2. รูปกณฺฑ \textbf{นวก} นววิเธน รูปสงฺคโห\csromandash{}}
\csromanmatikaanchor{mtk.88}
\csromanmatikaanchor{mtk.89}
\csromantocmarkref{section}{ทสก}{mtk.88}
\bookmark[dest={mtk.88},level=1]{ทสก}
\csromantocmarkref{section}{เอกาทสก}{mtk.89}
\bookmark[dest={mtk.89},level=1]{เอกาทสก}
\proseitem{591}{\csromanfolio{155}จกฺขุนฺทฺริยํ โสตินฺทฺริยํ ฆานินฺทฺริยํ ชิวฺหินฺทฺริยํ กายินฺทฺริยํ อิตฺถินฺทฺริยํ ปุริสินฺทฺริยํ ชีวิตินฺทฺริยํ, ยญฺจ รูปํ น อินฺทฺริยํ. เอวํ นววิเธน รูปสงฺคโห. นวกํ.}

\vspace{\tipitakaparskip}
\csromancenter{\textbf{ทสก} ทสวิเธน รูปสงฺคโห\csromandash{}}
\proseitem{592}{จกฺขุนฺทฺริยํ โสตินฺทฺริยํ ฆานินฺทฺริยํ ชิวฺหินฺทฺริยํ กายินฺทฺริยํ อิตฺถินฺทฺริยํ ปุริสินฺทฺริยํ ชีวิตินฺทฺริยํ น อินฺทฺริยํ รูปํ อตฺถิ สปฺปฏิฆํ อตฺถิ อปฺปฏิฆํ. เอวํ ทสวิเธน รูปสงฺคโห. ทสกํ.}

\vspace{\tipitakaparskip}
\csromancenter{\textbf{เอกาทสก} เอกาทสวิเธน รูปสงฺคโห\csromandash{}}
\proseitem{593}{จกฺขายตนํ โสตายตนํ ฆานายตนํ ชิวฺหายตนํ กายายตนํ รูปายตนํ สทฺทายตนํ คนฺธายตนํ รสายตนํ โผฏฺฐพฺพา\-ยตนํ, ยญฺจ รูปํ อนิทสฺสน\-อปฺปฏิฆํ\csromansharedfootnote{51bc355ad04b300e}{อนิทสฺสนํ อปฺปฏิฆํ (สี, สฺยา)} ธมฺมายตน\-ปริยาปนฺนํ. เอวํ เอกาทสวิเธน รูปสงฺคโห. เอกาทสกํ. มาติกา.}

\csromanmatikaanchor{mtk.90}
\csromantocmarknopage{section}{รูปวิภตฺติ}{mtk.90}
\bookmark[dest={mtk.90},level=1]{รูปวิภตฺติ}
\csromanheader{1}{\textbf{รูปวิภตฺติ}}
\csromanmatikaanchor{mtk.91}
\csromantocmarkref{subsection}{เอกกนิทฺเทส}{mtk.91}
\bookmark[dest={mtk.91},level=2]{เอกกนิทฺเทส}
\csromanheader{2}{\textbf{เอกกนิทฺเทส}}
\proseitem{594}{\csromanfolio{156}สพฺพํ รูปํ น เหตุเมว อเหตุกเมว เหตุ\-วิปฺปยุตฺตเมว สปฺปจฺจยเมว สงฺขตเมว รูปเมว โลกิยเมว สาสวเมว สญฺโญชนิยเมว คนฺถนิยเมว โอฆนิยเมว โยคนิยเมว นีวรณิยเมว ปรามฏฺฐเมว อุปาทานิยเมว สํกิเลสิกเมว อพฺยากตเมว อนารมฺมณเมว อเจตสิกเมว จิตฺต\-วิปฺปยุตฺตเมว เนว\-วิปาก\-น\-วิปากธมฺม\-ธมฺมเมว อสํกิลิฏฺฐ\-สํกิเลสิกเมว น สวิตกฺกสวิจารเมว น อวิตกฺกวิจารมตฺตเมว อวิตกฺกอวิจารเมว น ปีติสหคตเมว น สุขสหคตเมว น อุเปกฺขาสหคตเมว เนว ทสฺสเนน น ภาวนาย ปหาตพฺพ\-เหตุกเมว เนว อาจยคามิ น อปจยคามิเมว เนว\-เสกฺข\-นาเสกฺขเมว ปริตฺตเมว กามาวจรเมว น รูปาวจรเมว น อรูปาวจรเมว ปริยาปนฺนเมว โน อปริยาปนฺนเมว อนิยตเมว อนิยฺยานิกเมว อุปฺปนฺนํ ฉหิ วิญฺญาเณหิ วิญฺเญยฺยเมว อนิจฺจเมว ชราภิ\-ภูตเมว.}

\vspace{\tipitakaparskip}
\csromancenter{เอวํ เอกวิเธน รูปสงฺคโห.}
\csromancenter{เอกกนิทฺเทโส.}
\csromanmatikaanchor{mtk.92}
\csromantocmarkref{subsection}{ทุกนิทฺเทสอุปาทาภาชนีย}{mtk.92}
\bookmark[dest={mtk.92},level=2]{ทุกนิทฺเทสอุปาทาภาชนีย}
\csromanheader{2}{\textbf{ทุกนิทฺเทส}}
\csromanheader{2}{\textbf{อุปาทาภาชนีย}}
\proseitem{595}{กตมํ ตํ รูปํ อุปาทา. จกฺขายตนํ โสตายตนํ ฆานายตนํ ชิวฺหายตนํ กายายตนํ รูปายตนํ สทฺทายตนํ คนฺธายตนํ รสายตนํ อิตฺถินฺทฺริยํ ปุริสินฺทฺริยํ ชีวิตินฺทฺริยํ กายวิญฺญตฺติ วจีวิญฺญตฺติ อากาสธาตุ รูปสฺส ลหุตา รูปสฺส มุทุตา รูปสฺส กมฺมญฺญตา รูปสฺส อุปจโย รูปสฺส สนฺตติ รูปสฺส ชรตา รูปสฺส อนิจฺจตา กพฬีกาโร อาหาโร.}

\chapterheadpageread{2. รูปกณฺฑ}
\proseitem{596}{\csromanfolio{157}กตมํ ตํ รูปํ จกฺขายตนํ. ยํ จกฺขุ\csromansharedfootnote{aa096c50224701ff}{จกฺขุํ (สี, สฺยา)} จตุนฺนํ มหาภูตานํ อุปาทาย ปสาโท อตฺตภาว\-ปริยาปนฺโน อนิทสฺสโน สปฺปฏิโฆ, เยน จกฺขุนา อนิทสฺสเนน สปฺปฏิเฆน รูปํ สนิทสฺสนํ สปฺปฏิฆํ ปสฺสิ วา ปสฺสติ วา ปสฺสิสฺสติ วา ปสฺเส วา, จกฺขุํเปตํ จกฺขายตนํเปตํ จกฺขุธาตุ\-เปสา จกฺขุนฺทฺริยํเปตํ โลโกเปโส ทฺวาราเปสา สมุทฺโทเปโส ปณฺฑรํเปตํ เขตฺตํเปตํ วตฺถุํเปตํ เนตฺตํเปตฺตํ นยนํเปตํ โอริมํ ตีรํเปตํ สุญฺโญ คาโมเปโส. อิทํ ตํ รูปํ จกฺขายตนํ.}

\proseitem{597}{กตมํ ตํ รูปํ จกฺขายตนํ. ยํ จกฺขุ จตุนฺนํ มหาภูตานํ อุปาทาย ปสาโท อตฺตภาว\-ปริยาปนฺโน อนิทสฺสโน สปฺปฏิโฆ, ยมฺหิ จกฺขุมฺหิ อนิทสฺสนมฺหิ สปฺปฏิฆมฺหิ รูปํ สนิทสฺสนํ สปฺปฏิฆํ ปฏิหญฺญิ วา ปฏิหญฺญติ วา ปฏิหญฺญิสฺสติ วา ปฏิหญฺเญ วา, จกฺขุํเปตํ จกฺขายตนํเปตํ จกฺขุธาตุ\-เปสา จกฺขุนฺทฺริยํเปตํ โลโกเปโส ทฺวาราเปสา สมุทฺโทเปโส ปณฺฑรํเปตํ เขตฺตํเปตํ วตฺถุํเปตํ เนตฺตํเปตํ นยนํเปตํ โอริมํ ตีรํเปตํ สุญฺโญ คาโมเปโส. อิทํ ตํ รูปํ จกฺขายตนํ.}

\proseitem{598}{กตมํ ตํ รูปํ จกฺขายตนํ. ยํ จกฺขุ จตุนฺนํ มหาภูตานํ อุปาทาย ปสาโท อตฺตภาว\-ปริยาปนฺโน อนิทสฺสโน สปฺปฏิโฆ, ยํ จกฺขุ อนิทสฺสนํ สปฺปฏิฆํ รูปมฺหิ สนิทสฺสนมฺหิ สปฺปฏิฆมฺหิ ปฏิหญฺญิ วา ปฏิหญฺญติ วา ปฏิหญฺญิสฺสติ วา ปฏิหญฺเญ วา, จกฺขุํเปตํ จกฺขายตนํเปตํ จกฺขุธาตุ\-เปสา จกฺขุนฺทฺริยํเปตํ โลโกเปโส ทฺวาราเปสา สมุทฺโทเปโส ปณฺฑรํเปตํ เขตฺตํเปตํ วตฺถุํเปตํ เนตฺตํเปตํ นยนํเปตํ โอริมํ ตีรํเปตํ สุญฺโญ คาโมเปโส. อิทํ ตํ รูปํ จกฺขายตนํ.}

\proseitem{599}{กตมํ ตํ รูปํ จกฺขายตนํ. ยํ จกฺขุ จตุนฺนํ มหาภูตานํ อุปาทาย ปสาโท อตฺตภาว\-ปริยาปนฺโน อนิทสฺสโน สปฺปฏิโฆ, ยํ จกฺขุํ นิสฺสาย รูปํ อารพฺภ จกฺขุ\-สมฺผสฺโส อุปฺปชฺชิ วา อุปฺปชฺชติ วา อุปฺปชฺชิสฺสติ วา อุปฺปชฺเช วา ฯเปฯ. ยํ จกฺขุํ นิสฺสาย รูปํ อารพฺภ จกฺขุ\-สมฺผสฺสชา เวทนา ฯเปฯ สญฺญา ฯเปฯ เจตนา ฯเปฯ จกฺขุวิญฺญาณํ อุปฺปชฺชิ วา อุปฺปชฺชติ วา อุปฺปชฺชิสฺสติ วา อุปฺปชฺเช วา ฯเปฯ. ยํ จกฺขุํ นิสฺสาย รูปารมฺมโณ จกฺขุ\-สมฺผสฺโส อุปฺปชฺชิ วา อุปฺปชฺชติ วา อุปฺปชฺชิสฺสติ วา อุปฺปชฺเช วา ฯเปฯ. ยํ จกฺขุํ นิสฺสาย รูปารมฺมณา จกฺขุ\-สมฺผสฺสชา เวทนา ฯเปฯ สญฺญา ฯเปฯ เจตนา ฯเปฯ \csromanfolio{158}จกฺขุวิญฺญาณํ อุปฺปชฺชิ วา อุปฺปชฺชติ วา อุปฺปชฺชิสฺสติ วา อุปฺปชฺเช วา, จกฺขุํเปตํ จกฺขายตนํเปตํ จกฺขุธาตุ\-เปสา จกฺขุนฺทฺริยํเปตํ โลโกเปโส ทฺวาราเปสา สมุทฺโทเปโส ปณฺฑรํเปตํ เขตฺตํเปตํ วตฺถุํเปตํ เนตฺตํเปตํ นยนํเปตํ โอริมํ ตีรํเปตํ สุญฺโญ คาโมเปโส. อิทํ ตํ รูปํ จกฺขายตนํ.}

\proseitem{600}{กตมํ ตํ รูปํ โสตายตนํ. ยํ โสตํ จตุนฺนํ มหาภูตานํ อุปาทาย ปสาโท อตฺตภาว\-ปริยาปนฺโน อนิทสฺสโน สปฺปฏิโฆ, เยน โสเตน อนิทสฺสเนน สปฺปฏิเฆน สทฺทํ อนิทสฺสนํ สปฺปฏิฆํ สุณิ วา สุณาติ วา สุณิสฺสติ วา สุเณ วา, โสตํเปตํ โสตายตนํเปตํ โสตธาตุเปสา โสตินฺทฺริยํเปตํ โลโกเปโส ทฺวาราเปสา สมุทฺโทเปโส ปณฺฑรํเปตํ เขตฺตํเปตํ วตฺถุํเปตํ โอริมํ ตีรํเปตํ สุญฺโญ คาโมเปโส. อิทํ ตํ รูปํ โสตายตนํ.}

\proseitem{601}{กตมํ ตํ รูปํ โสตายตนํ. ยํ โสตํ จตุนฺนํ มหาภูตานํ อุปาทาย ปสาโท อตฺตภาว\-ปริยาปนฺโน อนิทสฺสโน สปฺปฏิโฆ, ยมฺหิ โสตมฺหิ อนิทสฺสนมฺหิ สปฺปฏิฆมฺหิ สทฺโท อนิทสฺสโน สปฺปฏิโฆ ปฏิหญฺญิ วา ปฏิหญฺญติ วา ปฏิหญฺญิสฺสติ วา ปฏิหญฺเญ วา, โสตํเปตํ โสตายตนํเปตํ โสตธาตุเปสา โสตินฺทฺริยํเปตํ โลโกเปโส ทฺวาราเปสา สมุทฺโทเปโส ปณฺฑรํเปตํ เขตฺตํเปตํ วตฺถุํเปตํ โอริมํ ตีรํเปตํ สุญฺโญ คาโมเปโส. อิทํ ตํ รูปํ โสตายตนํ.}

\proseitem{602}{กตมํ ตํ รูปํ โสตายตนํ. ยํ โสตํ จตุนฺนํ มหาภูตานํ อุปาทาย ปสาโท อตฺตภาว\-ปริยาปนฺโน อนิทสฺสโน สปฺปฏิโฆ, ยํ โสตํ อนิทสฺสนํ สปฺปฏิฆํ สทฺทมฺหิ อนิทสฺสนมฺหิ สปฺปฏิฆมฺหิ ปฏิหญฺญิ วา ปฏิหญฺญติ วา ปฏิหญฺญิสฺสติ วา ปฏิหญฺเญ วา, โสตํเปตํ โสตายตนํเปตํ โสตธาตุเปสา โสตินฺทฺริยํเปตํ โลโกเปโส ทฺวาราเปสา สมุทฺโทเปโส ปณฺฑรํเปตํ เขตฺตํเปตํ วตฺถุํเปตํ โอริมํ ตีรํเปตํ สุญฺโญ คาโมเปโส. อิทํ ตํ รูปํ โสตายตนํ.}

\proseitem{603}{กตมํ ตํ รูปํ โสตายตนํ. ยํ โสตํ จตุนฺนํ มหาภูตานํ อุปาทาย ปสาโท อตฺตภาว\-ปริยาปนฺโน อนิทสฺสโน สปฺปฏิโฆ, ยํ โสตํ นิสฺสาย สทฺทํ อารพฺภ โสตสมฺผสฺโส อุปฺปชฺชิ วา อุปฺปชฺชติ วา \csromanfolio{159}อุปฺปชฺชิสฺสติ วา อุปฺปชฺเช วา ฯเปฯ. ยํ โสตํ นิสฺสาย สทฺทํ อารพฺภ โสต\-สมฺผสฺสชา เวทนา ฯเปฯ สญฺญา ฯเปฯ เจตนา ฯเปฯ โสตวิญฺญาณํ อุปฺปชฺชิ วา อุปฺปชฺชติ วา อุปฺปชฺชิสฺสติ วา อุปฺปชฺเช วา ฯเปฯ. ยํ โสตํ นิสฺสาย สทฺทารมฺมโณ โสตสมฺผสฺโส อุปฺปชฺชิ วา อุปฺปชฺชติ วา อุปฺปชฺชิสฺสติ วา อุปฺปชฺเช วา ฯเปฯ. ยํ โสตํ นิสฺสาย สทฺทารมฺมณา โสต\-สมฺผสฺสชา เวทนา ฯเปฯ สญฺญา ฯเปฯ เจตนา ฯเปฯ โสตวิญฺญาณํ อุปฺปชฺชิ วา อุปฺปชฺชติ วา อุปฺปชฺชิสฺสติ วา อุปฺปชฺเช วา, โสตํเปตํ โสตายตนํเปตํ โสตธาตุเปสา โสตินฺทฺริยํเปตํ โลโกเปโส ทฺวาราเปสา สมุทฺโทเปโส ปณฺฑรํเปตํ เขตฺตํเปตํ วตฺถุํเปตํ โอริมํ ตีรํเปตํ สุญฺโญ คาโมเปโส. อิทํ ตํ รูปํ โสตายตนํ.}

\proseitem{604}{กตมํ ตํ รูปํ ฆานายตนํ. ยํ ฆานํ จตุนฺนํ มหาภูตานํ อุปาทาย ปสาโท อตฺตภาว\-ปริยาปนฺโน อนิทสฺสโน สปฺปฏิโฆ, เยน ฆาเนน อนิทสฺสเนน สปฺปฏิเฆน คนฺธํ อนิทสฺสนํ สปฺปฏิฆํ ฆายิ วา ฆายติ วา ฆายิสฺสติ วา ฆาเย วา, ฆานํเปตํ ฆานายตนํเปตํ ฆานธาตุเปสา ฆานินฺทฺริยํเปตํ โลโกเปโส ทฺวาราเปสา สมุทฺโทเปโส ปณฺฑรํเปตํ เขตฺตํเปตํ วตฺถุํเปตํ โอริมํ ตีรํเปตํ สุญฺโญ คาโมเปโส. อิทํ ตํ รูปํ ฆานายตนํ.}

\proseitem{605}{กตมํ ตํ รูปํ ฆานายตนํ. ยํ ฆานํ จตุนฺนํ มหาภูตานํ อุปาทาย ปสาโท อตฺตภาว\-ปริยาปนฺโน อนิทสฺสโน สปฺปฏิโฆ, ยมฺหิ ฆานมฺหิ อนิทสฺสนมฺหิ สปฺปฏิฆมฺหิ คนฺโธ อนิทสฺสโน สปฺปฏิโฆ ปฏิหญฺญิ วา ปฏิหญฺญติ วา ปฏิหญฺญิสฺสติ วา ปฏิหญฺเญ วา, ฆานํเปตํ ฆานายตนํเปตํ ฆานธาตุเปสา ฆานินฺทฺริยํเปตํ โลโกเปโส ทฺวาราเปสา สมุทฺโทเปโส ปณฺฑรํเปตํ เขตฺตํเปตํ วตฺถุํเปตํ โอริมํ ตีรํเปตํ สุญฺโญ คาโมเปโส. อิทํ ตํ รูปํ ฆานายตนํ.}

\proseitem{606}{กตมํ ตํ รูปํ ฆานายตนํ. ยํ ฆานํ จตุนฺนํ มหาภูตานํ อุปาทาย ปสาโท อตฺตภาว\-ปริยาปนฺโน อนิทสฺสโน สปฺปฏิโฆ, ยํ ฆานํ อนิทสฺสนํ สปฺปฏิฆํ คนฺธมฺหิ อนิทสฺสนมฺหิ สปฺปฏิฆมฺหิ ปฏิหญฺญิ วา ปฏิหญฺญติ วา ปฏิหญฺญิสฺสติ วา ปฏิหญฺเญ วา, ฆานํเปตํ ฆานายตนํเปตํ ฆานธาตุเปสา ฆานินฺทฺริยํเปตํ โลโกเปโส ทฺวาราเปสา \csromanfolio{160}สมุทฺโทเปโส ปณฺฑรํเปตํ เขตฺตํเปตํ วตฺถุํเปตํ โอริมํ ตีรํเปตํ สุญฺโญ คาโมเปโส. อิทํ ตํ รูปํ ฆานายตนํ.}

\proseitem{607}{กตมํ ตํ รูปํ ฆานายตนํ. ยํ ฆานํ จตุนฺนํ มหาภูตานํ อุปาทาย ปสาโท อตฺตภาว\-ปริยาปนฺโน อนิทสฺสโน สปฺปฏิโฆ, ยํ ฆานํ นิสฺสาย คนฺธํ อารพฺภ ฆานสมฺผสฺโส อุปฺปชฺชิ วา อุปฺปชฺชติ วา อุปฺปชฺชิสฺสติ วา อุปฺปชฺเช วา ฯเปฯ. ยํ ฆานํ นิสฺสาย คนฺธํ อารพฺภ ฆาน\-สมฺผสฺสชา เวทนา ฯเปฯ สญฺญา ฯเปฯ เจตนา ฯเปฯ ฆานวิญฺญาณํ อุปฺปชฺชิ วา อุปฺปชฺชติ วา อุปฺปชฺชิสฺสติ วา อุปฺปชฺเช วา ฯเปฯ. ยํ ฆานํ นิสฺสาย คนฺธารมฺมโณ ฆานสมฺผสฺโส อุปฺปชฺชิ วา อุปฺปชฺชติ วา อุปฺปชฺชิสฺสติ วา อุปฺปชฺเช วา ฯเปฯ. ยํ ฆานํ นิสฺสาย คนฺธารมฺมณา ฆาน\-สมฺผสฺสชา เวทนา ฯเปฯ สญฺญา ฯเปฯ เจตนา ฯเปฯ ฆานวิญฺญาณํ อุปฺปชฺชิ วา อุปฺปชฺชติ วา อุปฺปชฺชิสฺสติ วา อุปฺปชฺเช วา, ฆานํเปตํ ฆานายตนํเปตํ ฆานธาตุเปสา ฆานินฺทฺริยํเปตํ โลโกเปโส ทฺวาราเปสา สมุทฺโทเปโส ปณฺฑรํเปตํ เขตฺตํเปตํ วตฺถุํเปตํ โอริมํ ตีรํเปตํ สุญฺโญ คาโมเปโส. อิทํ ตํ รูปํ ฆานายตนํ.}

\proseitem{608}{กตมํ ตํ รูปํ ชิวฺหายตนํ. ยา ชิวฺหา จตุนฺนํ มหาภูตานํ อุปาทาย ปสาโท อตฺตภาว\-ปริยาปนฺโน อนิทสฺสโน สปฺปฏิโฆ, ยาย ชิวฺหาย อนิทสฺสนาย สปฺปฏิฆาย รสํ อนิทสฺสนํ สปฺปฏิฆํ สายิ วา สายติ วา สายิสฺสติ วา สาเย วา, ชิวฺหาเปสา ชิวฺหายตนํเปตํ ชิวฺหาธาตุเปสา ชิวฺหินฺทฺริยํเปตํ โลโกเปโส ทฺวาราเปสา สมุทฺโทเปโส ปณฺฑรํเปตํ เขตฺตํเปตํ วตฺถุํเปตํ โอริมํ ตีรํเปตํ สุญฺโญ คาโมเปโส. อิทํ ตํ รูปํ ชิวฺหายตนํ.}

\proseitem{609}{กตมํ ตํ รูปํ ชิวฺหายตนํ. ยา ชิวฺหา จตุนฺนํ มหาภูตานํ อุปาทาย ปสาโท อตฺตภาว\-ปริยาปนฺโน อนิทสฺสโน สปฺปฏิโฆ, ยาย ชิวฺหาย อนิทสฺสนาย สปฺปฏิฆาย รโส อนิทสฺสโน สปฺปฏิโฆ ปฏิหญฺญิ วา ปฏิหญฺญติ วา ปฏิหญฺญิสฺสติ วา ปฏิหญฺเญ วา, ชิวฺหาเปสา ชิวฺหายตนํเปตํ ชิวฺหาธาตุเปสา ชิวฺหินฺทฺริยํเปตํ โลโกเปโส ทฺวาราเปสา สมุทฺโทเปโส ปณฺฑรํเปตํ เขตฺตํเปตํ วตฺถุํเปตํ โอริมํ ตีรํเปตํ สุญฺโญ คาโมเปโส. อิทํ ตํ รูปํ ชิวฺหายตนํ.}

\proseitem{610}{กตมํ ตํ รูปํ ชิวฺหายตนํ. ยา ชิวฺหา จตุนฺนํ มหาภูตานํ อุปาทาย ปสาโท อตฺตภาว\-ปริยาปนฺโน อนิทสฺสโน สปฺปฏิโฆ, \csromanfolio{161}ยา ชิวฺหา อนิทสฺสนา สปฺปฏิฆา รสมฺหิ อนิทสฺสนมฺหิ สปฺปฏิฆมฺหิ ปฏิหญฺญิ วา ปฏิหญฺญติ วา ปฏิหญฺญิสฺสติ วา ปฏิหญฺเญ วา, ชิวฺหาเปสา ชิวฺหายตนํเปตํ ชิวฺหาธาตุเปสา ชิวฺหินฺทฺริยํเปตํ โลโกเปโส ทฺวาราเปสา สมุทฺโทเปโส ปณฺฑรํเปตํ เขตฺตํเปตํ วตฺถุํเปตํ โอริมํ ตีรํเปตํ สุญฺโญ คาโมเปโส. อิทํ ตํ รูปํ ชิวฺหายตนํ.}

\proseitem{611}{กตมํ ตํ รูปํ ชิวฺหายตนํ. ยา ชิวฺหา จตุนฺนํ มหาภูตานํ อุปาทาย ปสาโท อตฺตภาว\-ปริยาปนฺโน อนิทสฺสโน สปฺปฏิโฆ, ยํ ชิวฺหํ นิสฺสาย รสํ อารพฺภ ชิวฺหาสมฺผสฺโส อุปฺปชฺชิ วา อุปฺปชฺชติ วา อุปฺปชฺชิสฺสติ วา อุปฺปชฺเช วา ฯเปฯ. ยํ ชิวฺหํ นิสฺสาย รสํ อารพฺภ ชิวฺหา\-สมฺผสฺสชา เวทนา ฯเปฯ สญฺญา ฯเปฯ เจตนา ฯเปฯ ชิวฺหาวิญฺญาณํ อุปฺปชฺชิ วา อุปฺปชฺชติ วา อุปฺปชฺชิสฺสติ วา อุปฺปชฺเช วา ฯเปฯ. ยํ ชิวฺหํ นิสฺสาย รสารมฺมโณ ชิวฺหาสมฺผสฺโส อุปฺปชฺชิ วา อุปฺปชฺชติ วา อุปฺปชฺชิสฺสติ วา อุปฺปชฺเช วา ฯเปฯ. ยํ ชิวฺหํ นิสฺสาย รสารมฺมณา ชิวฺหาสมฺปสฺสชา เวทนา ฯเปฯ สญฺญา ฯเปฯ เจตนา ฯเปฯ ชิวฺหาวิญฺญาณํ อุปฺปชฺชิ วา อุปฺปชฺชติ วา อุปฺปชฺชิสฺสติ วา อุปฺปชฺเช วา, ชิวฺหาเปสา ชิวฺหายตนํเปตํ ชิวฺหาธาตุเปสา ชิวฺหินฺทฺริยํเปตํ โลโกเปโส ทฺวาราเปสา สมุทฺโทเปโส ปณฺฑรํเปตํ เขตฺตํเปตํ วตฺถุํเปตํ โอริมํ ตีรํเปตํ สุญฺโญ คาโมเปโส. อิทํ ตํ รูปํ ชิวฺหายตนํ.}

\proseitem{612}{กตมํ ตํ รูปํ กายายตนํ. โย กาโย จตุนฺนํ มหาภูตานํ อุปาทาย ปสาโท อตฺตภาว\-ปริยาปนฺโน อนิทสฺสโน สปฺปฏิโฆ, เยน กาเยน อนิทสฺสเนน สปฺปฏิเฆน โผฏฺฐพฺพํ อนิทสฺสน\-สปฺปฏิฆํ ผุสิ วา ผุสติ วา ผุสิสฺสติ วา ผุเส วา, กาโยเปโส กายายตนํเปตํ กายธาตุเปสา กายินฺทฺริยํเปตํ โลโกเปโส ทฺวาราเปสา สมุทฺโทเปโส ปณฺฑรํเปตํ เขตฺตํเปตํ วตฺถุํเปตํ โอริมํ ตีรํเปตํ สุญฺโญ คาโมเปโส. อิทํ ตํ รูปํ กายายตนํ.}

\proseitem{613}{กตมํ ตํ รูปํ กายายตนํ. โย กาโย จตุนฺนํ มหาภูตานํ อุปาทาย ปสาโท อตฺตภาว\-ปริยาปนฺโน อนิทสฺสโน สปฺปฏิโฆ, ยมฺหิ กายมฺหิ อนิทสฺสนมฺหิ สปฺปฏิฆมฺหิ โผฏฺฐพฺโพ อนิทสฺสโน สปฺปฏิโฆ ปฏิหญฺญิ วา ปฏิหญฺญติ วา ปฏิหญฺญิสฺสติ วา ปฏิหญฺเญ วา, กาโยเปโส กายายตนํเปตํ กายธาตุเปสา กายินฺทฺริยํเปตํ โลโกเปโส ทฺวาราเปสา สมุทฺโทเปโส ปณฺฑรํเปตํ \csromanfolio{162}เขตฺตํเปตํ วตฺถุํเปตํ โอริมํ ตีรํเปตํ สุญฺโญ คาโมเปโส. อิทํ ตํ รูปํ กายายตนํ.}

\proseitem{614}{กตมํ ตํ รูปํ กายายตนํ. โย กาโย จตุนฺนํ มหาภูตานํ อุปาทาย ปสาโท อตฺตภาว\-ปริยาปนฺโน อนิทสฺสโน สปฺปฏิโฆ, โย กาโย อนิทสฺสโน สปฺปฏิโฆ โผฏฺฐพฺพมฺหิ อนิทสฺสนมฺหิ สปฺปฏิฆมฺหิ ปฏิหญฺญิ วา ปฏิหญฺญติ วา ปฏิหญฺญิสฺสติ วา ปฏิหญฺเญ วา, กาโยเปโส กายายตนํเปตํ กายธาตุเปสา กายินฺทฺริยํเปตํ โลโกเปโส ทฺวาราเปสา สมุทฺโทเปโส ปณฺฑรํเปตํ เขตฺตํเปตํ วตฺถุํเปตํ โอริมํ ตีรํเปตํ สุญฺโญ คาโมเปโส. อิทํ ตํ รูปํ กายายตนํ.}

\proseitem{615}{กตมํ ตํ รูปํ กายายตนํ. โย กาโย จตุนฺนํ มหาภูตานํ อุปาทาย ปสาโท อตฺตภาว\-ปริยาปนฺโน อนิทสฺสโน สปฺปฏิโฆ, ยํ กายํ นิสฺสาย โผฏฺฐพฺพํ อารพฺภ กายสมฺผสฺโส อุปฺปชฺชิ วา อุปฺปชฺชติ วา อุปฺปชฺชิสฺสติ วา อุปฺปชฺเช วา ฯเปฯ. ยํ กายํ นิสฺสาย โผฏฺฐพฺพํ อารพฺภ กาย\-สมฺผสฺสชา เวทนา ฯเปฯ สญฺญา ฯเปฯ เจตนา ฯเปฯ กายวิญฺญาณํ อุปฺปชฺชิ วา อุปฺปชฺชติ วา อุปฺปชฺชิสฺสติ วา อุปฺปชฺเช วา ฯเปฯ. ยํ กายํ นิสฺสาย โผฏฺฐพฺพา\-รมฺมโณ กายสมฺผสฺโส อุปฺปชฺชิ วา อุปฺปชฺชติ วา อุปฺปชฺชิสฺสติ วา อุปฺปชฺเช วา ฯเปฯ. ยํ กายํ นิสฺสาย โผฏฺฐพฺพา\-รมฺมณา กาย\-สมฺผสฺสชา เวทนา ฯเปฯ สญฺญา ฯเปฯ เจตนา ฯเปฯ กายวิญฺญาณํ อุปฺปชฺชิ วา อุปฺปชฺชติ วา อุปฺปชฺชิสฺสติ วา อุปฺปชฺเช วา, กาโยเปโส กายายตนํเปตํ กายธาตุเปสา กายินฺทฺริยํเปตํ โลโกเปโส ทฺวาราเปสา สมุทฺโทเปโส ปณฺฑรํเปตํ เขตฺตํเปตํ วตฺถุํเปตํ โอริมํ ตีรํเปตํ สุญฺโญ คาโมเปโส. อิทํ ตํ รูปํ กายายตนํ.}

\proseitem{616}{กตมํ ตํ รูปํ รูปายตนํ. ยํ รูปํ จตุนฺนํ มหาภูตานํ อุปาทาย วณฺณนิภา สนิทสฺสนํ สปฺปฏิฆํ นีลํ ปีตกํ โลหิตกํ โอทาตํ กาฬกํ มญฺชิฏฺฐกํ\csromansharedfootnote{1b3052498abae4da}{มญฺเชฏฺฐกํ (สี, สฺยา)} หริ หริวณฺณํ อมฺพงฺกุร\-วณฺณํ ทีฆํ รสฺสํ อณุํ ถูลํ วฏฺฏํ ปริมณฺฑลํ จตุรํสํ\csromansharedfootnote{7647424e47729ef9}{จตุรสฺสํ (สี, ก)} ฉฬํสํ อฏฺฐํสํ โสฬสํสํ นินฺนํ ถลํ ฉายา อาตโป อาโลโก อนฺธกาโร อพฺภา มหิกา ธูโม รโช จนฺท\-มณฺฑลสฺส วณฺณนิภา สูริย\-มณฺฑลสฺส\csromansharedfootnote{21dad02c13a60165}{สุริยมณฺฑลสฺส (สี, สฺยา)} วณฺณนิภา \csromanfolio{163}ตารกรูปานํ วณฺณนิภา อาทาส\-มณฺฑลสฺส วณฺณนิภา มณิสงฺขมุตฺตา\-เวฬุริยสฺส วณฺณนิภา ชาตรูป\-รชตสฺส วณฺณนิภา, ยํ วา ปนญฺญมฺปิ อตฺถิ รูปํ จตุนฺนํ มหาภูตานํ อุปาทาย วณฺณนิภา สนิทสฺสนํ สปฺปฏิฆํ, ยํ รูปํ สนิทสฺสนํ สปฺปฏิฆํ จกฺขุนา อนิทสฺสเนน สปฺปฏิเฆน ปสฺสิ วา ปสฺสติ วา ปสฺสิสฺสติ วา ปสฺเส วา, รูปํเปตํ รูปายตนํเปตํ รูปธาตุเปสา. อิทํ ตํ รูปํ รูปายตนํ.}

\proseitem{617}{กตมํ ตํ รูปํ รูปายตนํ. ยํ รูปํ จตุนฺนํ มหาภูตานํ อุปาทาย วณฺณนิภา สนิทสฺสนํ สปฺปฏิฆํ นีลํ ปีตกํ โลหิตกํ โอทาตํ กาฬกํ มญฺชิฏฺฐกํ หริ หริวณฺณํ อมฺพงฺกุร\-วณฺณํ ทีฆํ รสฺสํ อณุํ ถูลํ วฏฺฏํ ปริมณฺฑลํ จตุรํสํ ฉฬํสํ อฏฺฐํสํ โสฬสํสํ นินฺนํ ถลํ ฉายา อาตโป อาโลโก อนฺธกาโร อพฺภา มหิกา ธูโม รโช จนฺท\-มณฺฑลสฺส วณฺณนิภา สูริย\-มณฺฑลสฺส วณฺณนิภา ตารกรูปานํ วณฺณนิภา อาทาส\-มณฺฑลสฺส วณฺณนิภา มณิสงฺขมุตฺตา\-เวฬุริยสฺส วณฺณนิภา ชาตรูป\-รชตสฺส วณฺณนิภา, ยํ วา ปนญฺญมฺปิ อตฺถิ รูปํ จตุนฺนํ มหาภูตานํ อุปาทาย วณฺณนิภา สนิทสฺสนํ สปฺปฏิฆํ, ยมฺหิ รูปมฺหิ สนิทสฺสนมฺหิ สปฺปฏิฆมฺหิ จกฺขุํ อนิทสฺสนํ สปฺปฏิฆํ ปฏิหญฺญิ วา ปฏิหญฺญติ วา ปฏิหญฺญิสฺสติ วา ปฏิหญฺเญ วา, รูปํเปตํ รูปายตนํเปตํ รูปธาตุเปสา. อิทํ ตํ รูปํ รูปายตนํ.}

\proseitem{618}{กตมํ ตํ รูปํ รูปายตนํ. ยํ รูปํ จตุนฺนํ มหาภูตานํ อุปาทาย วณฺณนิภา สนิทสฺสนํ สปฺปฏิฆํ นีลํ ปีตกํ โลหิตกํ โอทาตํ กาฬกํ มญฺชิฏฺฐกํ หริ หริวณฺณํ อมฺพงฺกุร\-วณฺณํ ทีฆํ รสฺสํ อณุํ ถูลํ วฏฺฏํ ปริมณฺฑลํ จตุรํสํ ฉฬํสํ อฏฺฐํสํ โสฬสํสํ นินฺนํ ถลํ ฉายา อาตโป อาโลโก อนฺธกาโร อพฺภา มหิกา ธูโม รโช จนฺท\-มณฺฑลสฺส วณฺณนิภา สูริย\-มณฺฑลสฺส วณฺณนิภา ตารกรูปานํ วณฺณนิภา อาทาส\-มณฺฑลสฺส วณฺณนิภา มณิสงฺขามุตฺตาเวฬุริยสฺส วณฺณนิภา ชาตรูป\-รชตสฺส วณฺณนิภา, ยํ วา ปนญฺญมฺปิ อตฺถิ รูปํ จตุนฺนํ มหาภูตานํ อุปาทาย วณฺณนิภา สนิทสฺสนํ สปฺปฏิฆํ, ยํ รูปํ สนิทสฺสนํ สปฺปฏิฆํ จกฺขุมฺหิ อนิทสฺสนมฺหิ สปฺปติฆมฺหิ ปฏิหญฺญิ วา ปฏิหญฺญติ วา ปฏิหญฺญิสฺสติ วา ปฏิหญฺเญ วา, รูปํเปตํ รูปายตนํเปตํ รูปธาตุเปสา. อิทํ ตํ รูปํ รูปายตนํ.}

\proseitem{619}{\csromanfolio{164}กตมํ ตํ รูปํ รูปายตนํ. ยํ รูปํ จตุนฺนํ มหาภูตานํ อุปาทาย วณฺณนิภา สนิทสฺสนํ สปฺปฏิฆํ นีลํ ปีตกํ โลหิตกํ โอทาตํ กาฬกํ มญฺชิฏฺฐกํ หริ หริวณฺณํ อมฺพงฺกุร\-วณฺณํ ทีฆํ รสฺสํ อณุํ ถูลํ วฏฺฏํ ปริมณฺฑลํ จตุรํสํ ฉฬํสํ อฏฺฐํสํ โสฬสํสํ นินฺนํ ถลํ ฉายา อาตโป อาโลโก อนฺธกาโร อพฺภา มหิกา ธูโม รโช จนฺท\-มณฺฑลสฺส วณฺณนิภา สูริย\-มณฺฑลสฺส วณฺณนิภา ตารกรูปานํ วณฺณนิภา อาทาส\-มณฺฑลสฺส วณฺณนิภา มณิสงฺขมุตฺตา\-เวฬุริยสฺส วณฺณนิภา ชาตรูป\-รชตสฺส วณฺณนิภา, ยํ วา ปนญฺญมฺปิ อตฺถิ รูปํ จตุนฺนํ มหาภูตานํ อุปาทาย วณฺณนิภา สนิทสฺสนํ สปฺปฏิฆํ, ยํ รูปํ อารพฺภ จกฺขุํ นิสฺสาย จกฺขุ\-สมฺผสฺโส อุปฺปชฺชิ วา อุปฺปชฺชติ วา อุปฺปชฺชิสฺสติ วา อุปฺปชฺเช วา ฯเปฯ. ยํ รูปํ อารพฺภ จกฺขุํ นิสฺสาย จกฺขุ\-สมฺผสฺสชา เวทนา ฯเปฯ สญฺญา ฯเปฯ เจตนา ฯเปฯ จกฺขุวิญฺญาณํ อุปฺปชฺชิ วา อุปฺปชฺชติ วา อุปฺปชฺชิสฺสติ วา อุปฺปชฺเช วา ฯเปฯ. ยํ รูปารมฺมโณ จกฺขุํ นิสฺสย จกฺขุ\-สมฺผสฺโส อุปฺปชฺชิ วา อุปฺปชฺชติ วา อุปฺปชฺชิสฺสติ วา อุปฺปชฺเช วา ฯเปฯ. ยํ รูปารมฺมณา จกฺขุํ นิสฺสาย จกฺขุ\-สมฺผสฺสชา เวทนา ฯเปฯ สญฺญา ฯเปฯ เจตนา ฯเปฯ จกฺขุวิญฺญาณํ อุปฺปชฺชิ วา อุปฺปชฺชติ วา อุปฺปชฺชิสฺสติ วา อุปฺปชฺเช วา, รูปํเปตํ รูปายตนํเปตํ รูปธาตุเปสา. อิทํ ตํ รูปํ รูปายตนํ.}

\proseitem{620}{กตมํ ตํ รูปํ สทฺทายตนํ. โย สทฺโท จตุนฺนํ มหาภูตานํ อุปาทาย อนิทสฺสโน สปฺปฏิโฆ เภริสทฺโท มุทิงฺคสทฺโท สงฺขสทฺโท ปณวสทฺโท คีตสทฺโท วาทิตสทฺโท สมฺมสทฺโท ปาณิสทฺโท สตฺตานํ นิคฺโฆสสทฺโท ธาตูนํ สนฺนิฆาต\-สทฺโท วาตสทฺโท อุทกสทฺโท มนุสฺสสทฺโท อมนุสฺสสทฺโท, โย วา ปนญฺโญปิ อตฺถิ สทฺโท จตุนฺนํ มหาภูตานํ อุปาทาย อนิทสฺสโน สปฺปฏิโฆ, ยํ สทฺทํ อนิทสฺสนํ สปฺปฏิฆํ โสเตน อนิทสฺสเนน สปฺปฏิเฆน สุณิ วา สุณาติ วา สุณิสฺสติ วา สุเณ วา, สทฺโทเปโส สทฺทายตนํเปตํ สทฺทธาตุเปสา. อิทํ ตํ รูปํ สทฺทายตนํ.}

\proseitem{621}{กตมํ ตํ รูปํ สทฺทายตนํ. โย สทฺโท จตุนฺนํ มหาภูตานํ อุปาทาย อนิทสฺสโน สปฺปฏิโฆ เภริสทฺโท มุทิงฺคสทฺโท สงฺขสทฺโท ปณวสทฺโท คีตสทฺโท วาทิตสทฺโท สมฺมสทฺโท ปาณิสทฺโท สตฺตานํ นิคฺโฆสสทฺโท ธาตูนํ สนฺนิฆาต\-สทฺโท วาตสทฺโท อุทกสทฺโท มนุสฺสสทฺโท อมนุสฺสสทฺโท, โย วา ปนญฺโญปิ อตฺถิ สทฺโท จตุนฺนํ \csromanfolio{165}มหาภูตานํ อุปาทาย อนิทสฺสโน สปฺปฏิโฆ, ยมฺหิ สทฺทมฺหิ อนิทสฺสนมฺหิ สปฺปฏิฆมฺหิ โสตํ อนิทสฺสนํ สปฺปฏิฆํ ปฏิหญฺญิ วา ปฏิหญฺญติ วา ปฏิหญฺญิสฺสติ วา ปฏิหญฺเญ วา, สทฺโทเปโส สทฺทายตนํเปตํ สทฺทธาตุเปสา. อิทํ ตํ รูปํ สทฺทายตนํ.}

\proseitem{622}{กตมํ ตํ รูปํ สทฺทายตนํ. โย สทฺโท จตุนฺนํ มหาภูตานํ อุปาทาย อนิทสฺสโน สปฺปฏิโฆ เภริสทฺโท มุทิงฺคสทฺโท สงฺขสทฺโท ปณวสทฺโท คีตสทฺโท วาทิตสทฺโท สมฺมสทฺโท ปาณิสทฺโท สตฺตานํ นิคฺโฆสสทฺโท ธาตูนํ สนฺนิฆาต\-สทฺโท วาตสทฺโท อุทกสทฺโท มนุสฺสสทฺโท อมนุสฺสสทฺโท, โย วา ปนญฺโญปิ อตฺถิ สทฺโท จตุนฺนํ มหาภูตานํ อุปาทาย อนิทสฺสโน สปฺปฏิโฆ, โย สทฺโท อนิทสฺสโน สปฺปฏิโฆ โสตมฺหิ อนิทสฺสนมฺหิ สปฺปฏิฆมฺหิ ปฏิหญฺญิ วา ปฏิหญฺญติ วา ปฏิหญฺญิสฺสติ วา ปฏิหญฺเญ วา, สทฺโทเปโส สทฺทายตนํเปตํ สทฺทธาตุเปสา. อิทํ ตํ รูปํ สทฺทายตนํ.}

\proseitem{623}{กตมํ ตํ รูปํ สทฺทายตนํ. โย สทฺโท จตุนฺนํ มหาภูตานํ อุปาทาย อนิทสฺสโน สปฺปฏิโฆ เภริสทฺโท มุทิงฺคสทฺโท สงฺขสทฺโท ปณวสทฺโท คีตสทฺโท วาทิตสทฺโท สมฺมสทฺโท ปาณิสทฺโท สตฺตานํ นิคฺโฆสสทฺโท ธาตูนํ สนฺนิฆาต\-สทฺโท วาตสทฺโท อุทกสทฺโท มนุสฺสสทฺโท อมนุสฺสสทฺโท, โย วา ปนญฺโญปิ อตฺถิ สทฺโท จตุนฺนํ มหาภูตานํ อุปาทาย อนิทสฺสโน สปฺปฏิโฆ, ยํ สทฺทํ อารพฺภ โสตํ นิสฺสาย โสตสมฺผสฺโส อุปฺปชฺชิ วา อุปฺปชฺชติ วา อุปฺปชฺชิสฺสติ วา อุปฺปชฺเช วา ฯเปฯ. ยํ สทฺทํ อารพฺภ โสตํ นิสฺสาย โสต\-สมฺผสฺสชา เวทนา ฯเปฯ สญฺญา ฯเปฯ เจตนา ฯเปฯ โสตวิญฺญาณํ อุปฺปชฺชิ วา อุปฺปชฺชติ วา อุปฺปชฺชิสฺสติ วา อุปฺปชฺเช วา ฯเปฯ. ยํ สทฺทารมฺมโณ โสตํ นิสฺสาย โสตสมฺผสฺโส อุปฺปชฺชิ วา อุปฺปชฺชติ วา อุปฺปชฺชิสฺสติ วา อุปฺปชฺเช วา ฯเปฯ. ยํ สทฺทารมฺมณา โสตํ นิสฺสาย โสต\-สมฺผสฺสชา เวทนา ฯเปฯ สญฺญา ฯเปฯ เจตนา ฯเปฯ โสตวิญฺญาณํ อุปฺปชฺชิ วา อุปฺปชฺชติ วา อุปฺปชฺชิสฺสติ วา อุปฺปชฺเช วา, สทฺโทเปโส สทฺทายตนํเปตํ สทฺทธาตุเปสา. อิทํ ตํ รูปํ สทฺทายตนํ.}

\proseitem{624}{กตมํ ตํ รูปํ คนฺธายตนํ. โย คนฺโธ จตุนฺนํ มหาภูตานํ อุปาทาย อนิทสฺสโน สปฺปฏิโฆ มูลคนฺโธ สารคนฺโธ ตจคนฺโธ ปตฺตคนฺโธ ปุปฺผคนฺโธ ผลคนฺโธ อามกคนฺโธ วิสฺสคนฺโธ สุคนฺโธ \csromanfolio{166}ทุคฺคนฺโธ, โย วา ปนญฺโญปิ อตฺถิ คนฺโธ จตุนฺนํ มหาภูตานํ อุปาทาย อนิทสฺสโน สปฺปฏิโฆ, ยํ คนฺธํ อนิทสฺสนํ สปฺปฏิฆํ ฆาเนน อนิทสฺสเนน สปฺปฏิเฆน ฆายิ วา ฆายติ วา ฆายิสฺสติ วา ฆาเย วา, คนฺโธเปโส คนฺธายตนํเปตํ คนฺธธาตุ\-เปสา. อิทํ ตํ รูปํ คนฺธายตนํ.}

\proseitem{625}{กตมํ ตํ รูปํ คนฺธายตนํ. โย คนฺโธ จตุนฺนํ มหาภูตานํ อุปาทาย อนิทสฺสโน สปฺปฏิโฆ มูลคนฺโธ สารคนฺโธ ตจคนฺโธ ปตฺตคนฺโธ ปุปฺผคนฺโธ ผลคนฺโธ อามกคนฺโธ วิสฺสคนฺโธ สุคนฺโธ ทุคฺคนฺโธ, โย วา ปนญฺโญปิ อตฺถิ คนฺโธ จตุนฺนํ มหาภูตานํ อุปาทาย อนิทสฺสโน สปฺปฏิโฆ, ยมฺหิ คนฺธมฺหิ อนิทสฺสนมฺหิ สปฺปฏิฆมฺหิ ฆานํ อนิทสฺสนํ สปฺปฏิฆํ ปฏิหญฺญิ วา ปฏิหญฺญติ วา ปฏิหญฺญิสฺสติ วา ปฏิหญฺเญ วา, คนฺโธเปโส คนฺธายตนํเปตํ คนฺธธาตุ\-เปสา. อิทํ ตํ รูปํ คนฺธายตนํ.}

\proseitem{626}{กตมํ ตํ รูปํ คนฺธายตนํ. โย คนฺโธ จตุนฺนํ มหาภูตานํ อุปาทาย อนิทสฺสโน สปฺปฏิโฆ มูลคนฺโธ สารคนฺโธ ตจคนฺโธ ปตฺตคนฺโธ ปุปฺผคนฺโธ ผลคนฺโธ อามกคนฺโธ วิสฺสคนฺโธ สุคนฺโธ ทุคฺคนฺโธ, โย วา ปนญฺโญปิ อตฺถิ คนฺโธ จตุนฺนํ มหาภูตานํ อุปาทาย อนิทสฺสโน สปฺปฏิโฆ, โย คนฺโธ อนิทสฺสโน สปฺปฏิโฆ ฆานมฺหิ อนิทสฺสนมฺหิ สปฺปฏิฆมฺหิ ปฏิหญฺญิ วา ปฏิหญฺญติ วา ปฏิหญฺญิสฺสติ วา ปฏิหญฺเญ วา, คนฺโธเปโส คนฺธายตนํเปตํ คนฺธธาตุ\-เปสา. อิทํ ตํ รูปํ คนฺธายตนํ.}

\proseitem{627}{กตมํ ตํ รูปํ คนฺธายตนํ. โย คนฺโธ จตุนฺนํ มหาภูตานํ อุปาทาย อนิทสฺสโน สปฺปฏิโฆ มูลคนฺโธ สารคนฺโธ ตจคนฺโธ ปตฺตคนฺโธ ปุปฺผคนฺโธ ผลคนฺโธ อามกคนฺโธ วิสฺสคนฺโธ สุคนฺโธ ทุคฺคนฺโธ, โย วา ปนญฺโญปิ อตฺถิ คนฺโธ จตุนฺนํ มหาภูตานํ อุปาทาย อนิทสฺสโน สปฺปฏิโฆ, ยํ คนฺธํ อารพฺภ ฆานํ นิสฺสาย ฆานสมฺผสฺโส อุปฺปชฺชิ วา อุปฺปชฺชติ วา อุปฺปชฺชิสฺสติ วา อุปฺปชฺเช วา ฯเปฯ. ยํ คนฺธํ อารพฺภ ฆานํ นิสฺสาย ฆาน\-สมฺผสฺสชา เวทนา ฯเปฯ สญฺญา ฯเปฯ เจตนา ฯเปฯ ฆานวิญฺญาณํ อุปฺปชฺชิ วา อุปฺปชฺชติ วา อุปฺปชฺชิสฺสติ วา อุปฺปชฺเช วา ฯเปฯ. ยํ คนฺธารมฺมโณ ฆานํ นิสฺสาย ฆานสมฺผสฺโส อุปฺปชฺชิ วา อุปฺปชฺชติ วา \csromanfolio{167}อุปฺปชฺชิสฺสติ วา อุปฺปชฺเช วา ฯเปฯ. ยํ คนฺธารมฺมณา ฆานํ นิสฺสาย ฆาน\-สมฺผสฺสชา เวทนา ฯเปฯ สญฺญา ฯเปฯ เจตนา ฯเปฯ ฆานวิญฺญาณํ อุปฺปชฺชิ วา อุปฺปชฺชติ วา อุปฺปชฺชิสฺสติ วา อุปฺปชฺเช วา, คนฺโธเปโส คนฺธายตนํเปตํ คนฺธธาตุ\-เปสา. อิทํ ตํ รูปํ คนฺธายตนํ.}

\proseitem{628}{กตมํ ตํ รูปํ รสายตนํ. โย รโส จตุนฺนํ มหาภูตานํ อุปาทาย อนิทสฺสโน สปฺปฏิโฆ มูลรโส ขนฺธรโส ตจรโส ปตฺตรโส ปุปฺผรโส ผลรโส อมฺพิลํ มธุรํ ติตฺตกํ กฏุกํ โลณิกํ ขาริกํ ลมฺพิลํ\csromansharedfootnote{0662900a086b6b6c}{ลปิลกํ (สี)} กสาโว สาทุ อสาทุ, โย วา ปนญฺโญปิ อตฺถิ รโส จตุนฺนํ มหาภูตานํ อุปาทาย อนิทสฺสโน สปฺปฏิโฆ, ยํ รสํ อนิทสฺสนํ สปฺปฏิฆํ ชิวฺหาย อนิทสฺสนาย สปฺปฏิฆาย สายิ วา สายติ วา สายิสฺสติ วา สาเย วา, รโสเปโส รสายตนํเปตํ รสธาตุเปสา. อิทํ ตํ รูปํ รสายตนํ.}

\proseitem{629}{กตมํ ตํ รูปํ รสายตนํ. โย รโส จตุนฺนํ มหาภูตานํ อุปาทาย อนิทสฺสโน สปฺปฏิโฆ มูลรโส ขนฺธรโส ตจรโส ปตฺตรโส ปุปฺผรโส ผลรโส อมฺพิลํ มธุรํ ติตฺตกํ กฏุกํ โลณิกํ ขาริกํ ลมฺพิลํ กสาโว สาทุ อสาทุ, โย วา ปนญฺโญปิ อตฺถิ รโส จตุนฺนํ มหาภูตานํ อุปาทาย อนิทสฺสโน สปฺปฏิโฆ, ยมฺหิ รสมฺหิ อนิทสฺสนมฺหิ สปฺปฏิฆมฺหิ ชิวฺหา อนิทสฺสนา สปฺปฏิฆา ปฏิหญฺญิ วา ปฏิหญฺญติ วา ปฏิหญฺญิสฺสติ วา ปฏิหญฺเญ วา, รโสเปโส รสายตนํเปตํ รสธาตุเปสา. อิทํ ตํ รูปํ รสายตนํ.}

\proseitem{630}{กตมํ ตํ รูปํ รสายตนํ. โย รโส จตุนฺนํ มหาภูตานํ อุปาทาย อนิทสฺสโน สปฺปฏิโฆ มูลรโส ขนฺธรโส ตจรโส ปตฺตรโส ปุปฺผรโส ผลรโส อมฺพิลํ มธุรํ ติตฺตกํ กฏุกํ โลณิกํ ขาริกํ ลมฺพิลํ กสาโว สาทุ อสาทุ, โย วา ปนญฺโญปิ อตฺถิ รโส จตุนฺนํ มหาภูตานํ อุปาทาย อนิทสฺสโน สปฺปฏิโฆ, โย รโส อนิทสฺสโน สปฺปฏิโฆ ชิวฺหาย อนิทสฺสนาย สปฺปฏิฆาย ปฏิหญฺญิ วา ปฏิหญฺญติ วา ปฏิหญฺญิสฺสติ วา ปฏิหญฺเญ วา, รโสเปโส รสายตนํเปตํ รสธาตุเปสา. อิทํ ตํ รูปํ รสายตนํ.}

\proseitem{631}{\csromanfolio{168}กตมํ ตํ รูปํ รสายตนํ. โย รโส จตุนฺนํ มหาภูตานํ อุปาทาย อนิทสฺสโน สปฺปฏิโฆ มูลรโส ขนฺธรโส ตจรโส ปตฺตรโส ปุปฺผรโส ผลรโส อมฺพิลํ มธุรํ ติตฺตกํ กฏุกํ โลณิกํ ขาริกํ ลมฺพิลํ กสาโว สาทุ อสาทุ, โย วา ปนญฺโญปิ อตฺถิ รโส จตุนฺนํ มหาภูตานํ อุปาทาย อนิทสฺสโน สปฺปฏิโฆ, ยํ รสํ อารพฺภ ชิวฺหํ นิสฺสาย ชิวฺหาสมฺผสฺโส อุปฺปชฺชิ วา อุปฺปชฺชติ วา อุปฺปชฺชิสฺสติ วา อุปฺปชฺเช วา ฯเปฯ. ยํ รสํ อารพฺภ ชิวฺหํ นิสฺสาย ชิวฺหา\-สมฺผสฺสชา เวทนา ฯเปฯ สญฺญา ฯเปฯ เจตนา ฯเปฯ ชิวฺหาวิญฺญาณํ อุปฺปชฺชิ วา อุปฺปชฺชติ วา อุปฺปชฺชิสฺสติ วา อุปฺปชฺเช วา ฯเปฯ. ยํ รสารมฺมโณ ชิวฺหํ นิสฺสาย ชิวฺหาสมฺผสฺโส อุปฺปชฺชิ วา อุปฺปชฺชติ วา อุปฺปชฺชิสฺสติ วา อุปฺปชฺเช วา ฯเปฯ. ยํ รสารมฺมณา ชิวฺหํ นิสฺสาย ชิวฺหา\-สมฺผสฺสชา เวทนา ฯเปฯ สญฺญา ฯเปฯ เจตนา ฯเปฯ ชิวฺหาวิญฺญาณํ อุปฺปชฺชิ วา อุปฺปชฺชติ วา อุปฺปชฺชิสฺสติ วา อุปฺปชฺเช วา, รโสเปโส รสายตนํเปตํ รสธาตุเปสา. อิทํ ตํ รูปํ รสายตนํ.}

\proseitem{632}{กตมํ ตํ รูปํ อิตฺถินฺทฺริยํ. ยํ อิตฺถิยา อิตฺถิลิงฺคํ อิตฺถินิมิตฺตํ อิตฺถิกุตฺตํ อิตฺถากปฺโป อิตฺถตฺตํ อิตฺถิภาโว\csromansharedfootnote{1b053dcb70e220df}{อิตฺถิตฺตํ อิตฺถีภาโว (สฺยา)}. อิทํ ตํ รูปํ อิตฺถินฺทฺริยํ.}

\proseitem{633}{กตมํ ตํ รูปํ ปุริสินฺทฺริยํ. ยํ ปุริสสฺส ปุริสลิงฺคํ ปุริสนิมิตฺตํ ปุริสกุตฺตํ ปุริสากปฺโป ปุริสตฺตํ ปุริสภาโว. อิทํ ตํ รูปํ ปุริสินฺทฺริยํ.}

\proseitem{634}{กตมํ ตํ รูปํ ชีวิตินฺทฺริยํ. โย เตสํ รูปีนํ ธมฺมานํ อายุ ฐิติ ยปนา ยาปนา อิริยนา วตฺตนา ปาลนา ชีวิตํ ชีวิตินฺทฺริยํ. อิทํ ตํ รูปํ ชีวิตินฺทฺริยํ.}

\proseitem{635}{กตมํ ตํ รูปํ กายวิญฺญตฺติ. ยา กุสลจิตฺตสฺส วา อกุสลจิตฺตสฺส วา อพฺยากต\-จิตฺตสฺส วา อภิกฺกม\-นฺตสฺส วา ปฏิกฺกมนฺตสฺส วา อาโลเกนฺตสฺส วา วิโลเกนฺตสฺส วา สมิญฺเชนฺตสฺส วา ปสาเรนฺตสฺส วา กายสฺส ถมฺภนา สนฺถมฺภนา สนฺถมฺภิตตฺตํ วิญฺญตฺติ วิญฺญาปนา วิญฺญาปิตตฺตํ. อิทํ ตํ รูปํ กายวิญฺญตฺติ.}

\proseitem{636}{กตมํ ตํ รูปํ วจีวิญฺญตฺติ. ยา กุสลจิตฺตสฺส วา อกุสลจิตฺตสฺส วา อพฺยากต\-จิตฺตสฺส วา วาจา คิรา พฺยปฺปโถ อุทีรณํ โฆโส โฆสกมฺมํ วาจา วจีเภโท, อยํ วุจฺจติ วาจา. ยา ตาย วาจาย วิญฺญตฺติ วิญฺญาปนา วิญฺญาปิตตฺตํ. อิทํ ตํ รูปํ วจีวิญฺญตฺติ.}

\chapterheadpageread{2. รูปกณฺฑ}
\proseitem{637}{\csromanfolio{169}กตมํ ตํ รูปํ อากาสธาตุ. โย อากาโส อากาสคตํ อฆํ อฆคตํ วิวโร วิวรคตํ อสมฺผุฏฺฐํ จตูหิ มหาภูเตหิ. อิทํ ตํ รูปํ อากาสธาตุ.}

\proseitem{638}{กตมํ ตํ รูปํ รูปสฺส ลหุตา. ยา รูปสฺส ลหุตา ลหุปริณามตา อทนฺธนตา อวิตฺถนตา. อิทํ ตํ รูปํ รูปสฺส ลหุตา.}

\proseitem{639}{กตมํ ตํ รูปํ รูปสฺส มุทุตา. ยา รูปสฺส มุทุตา มทฺทวตา อกกฺขฬตา อกถินตา. อิทํ ตํ รูปํ รูปสฺส มุทุตา.}

\proseitem{640}{กตมํ ตํ รูปํ รูปสฺส กมฺมญฺญตา. ยา รูปสฺส กมฺมญฺญตา กมฺมญฺญตฺตํ กมฺมญฺญภาโว. อิทํ ตํ รูปํ รูปสฺส กมฺมญฺญตา.}

\proseitem{641}{กตมํ ตํ รูปํ รูปสฺส อุปจโย. โย อายตนานํ อาจโย, โส รูปสฺส อุปจโย. อิทํ ตํ รูปํ รูปสฺส อุปจโย.}

\proseitem{642}{กตมํ ตํ รูปํ รูปสฺส สนฺตติ. โย รูปสฺส อุปจโย, สา รูปสฺส สนฺตติ. อิทํ ตํ รูปํ รูปสฺส สนฺตติ.}

\proseitem{643}{กตมํ ตํ รูปํ รูปสฺส ชรตา. ยา รูปสฺส ชรา ชีรณตา ขณฺฑิจฺจํ ปาลิจฺจํ วลิตฺตจตา อายุโน สํหานิ อินฺทฺริยานํ ปริปาโก. อิทํ ตํ รูปํ รูปสฺส ชรตา.}

\proseitem{644}{กตมํ ตํ รูปํ รูปสฺส อนิจฺจตา. โย รูปสฺส ขโย วโย เภโท ปริเภโท อนิจฺจตา อนฺตรธานํ. อิทํ ตํ รูปํ รูปสฺส อนิจฺจตา.}

\proseitem{645}{กตมํ ตํ รูปํ กพฬีกาโร อาหาโร. โอทโน กุมฺมาโส สตฺตุ มจฺโฉ มํสํ ขีรํ ทธิ สปฺปิ นวนีตํ เตลํ มธุ ผาณิตํ, ยํ วา ปนญฺญมฺปิ อตฺถิ รูปํ ยมฺหิ ยมฺหิ ชนปเท เตสํ เตสํ สตฺตานํ มุขาสิยํ ทนฺตวิขาทนํ คลชฺโฌหรณียํ กุจฺฉิ\-วิตฺถมฺภนํ ยาย โอชาย สตฺตา ยาเปนฺติ. อิทํ ตํ รูปํ กพฬีกาโร อาหาโร. อิทํ ตํ รูปํ อุปาทา. อุปาทา\-ภาชนียํ. รูปกณฺเฑ ปฐม\-ภาณวาโร.}

\proseitem{646}{\csromanfolio{170}กตมํ ตํ รูปํ โน อุปาทา. โผฏฺฐพฺพา\-ยตนํ อาโปธาตุ.}

\proseitem{647}{กตมํ ตํ รูปํ โผฏฺฐพฺพา\-ยตนํ. ปถวีธาตุ เตโชธาตุ วาโยธาตุ กกฺขฬํ มุทุกํ สณฺหํ ผรุสํ สุขสมฺผสฺสํ ทุกฺข\-สมฺผสฺสํ ครุกํ ลหุกํ, ยํ โผฏฺฐพฺพํ อนิทสฺสนํ สปฺปฏิฆํ กาเยน อนิทสฺสเนน สปฺปฏิเฆน ผุสิ วา ผุสติ วา ผุสิสฺสติ วา ผุเส วา, โผฏฺฐพฺโพเปโส โผฏฺฐพฺพา\-ยตนํเปตํ โผฏฺฐพฺพ\-ธาตุเปสา. อิทํ ตํ รูปํ โผฏฺฐพฺพา\-ยตนํ.}

\proseitem{648}{กตมํ ตํ รูปํ โผฏฺฐพฺพา\-ยตนํ. ปถวีธาตุ เตโชธาตุ วาโยธาตุ กกฺขฬํ มุทุกํ สณฺหํ ผรุสํ สุขสมฺผสฺสํ ทุกฺข\-สมฺผสฺสํ ครุกํ ลหุกํ, ยมฺหิ โผฏฺฐพฺพมฺหิ อนิทสฺสนมฺหิ สปฺปฏิฆมฺหิ กาโย อนิทสฺสโน สปฺปฏิโฆ ปฏิหญฺญิ วา ปฏิหญฺญติ วา ปฏิหญฺญิสฺสติ วา ปฏิหญฺเญ วา, โผฏฺฐพฺโพเปโส โผฏฺฐพฺพา\-ยตนํเปตํ โผฏฺฐพฺพ\-ธาตุเปสา. อิทํ ตํ รูปํ โผฏฺฐพฺพา\-ยตนํ.}

\proseitem{649}{กตมํ ตํ รูปํ โผฏฺฐพฺพา\-ยตนํ. ปถวีธาตุ เตโชธาตุ วาโยธาตุ กกฺขฬํ มุทุกํ สณฺหํ ผรุสํ สุขสมฺผสฺสํ ทุกฺข\-สมฺผสฺสํ ครุกํ ลหุกํ, โย โผฏฺฐพฺโพ อนิทสฺสโน สปฺปฏิโฆ กายมฺหิ อนิทสฺสนมฺหิ สปฺปฏิฆมฺหิ ปฏิหญฺญิ วา ปฏิหญฺญติ วา ปฏิหญฺญิสฺสติ วา ปฏิหญฺเญ วา, โผฏฺฐพฺโพเปโส โผฏฺฐพฺพา\-ยตนํเปตํ โผฏฺฐพฺพ\-ธาตุเปสา. อิทํ ตํ รูปํ โผฏฺฐพฺพา\-ยตนํ.}

\proseitem{650}{กตมํ ตํ รูปํ โผฏฺฐพฺพา\-ยตนํ. ปถวีธาตุ เตโชธาตุ วาโยธาตุ กกฺขฬํ มุทุกํ สณฺหํ ผรุสํ สุขสมฺผสฺสํ ทุกฺข\-สมฺผสฺสํ ครุกํ ลหุกํ, ยํ โผฏฺฐพฺพํ อารพฺภ กายํ นิสฺสาย กายสมฺผสฺโส อุปฺปชฺชิ วา อุปฺปชฺชติ วา อุปฺปชฺชิสฺสติ วา อุปฺปชฺเช วา ฯเปฯ. ยํ โผฏฺฐพฺพํ อารพฺภ กายํ นิสฺสาย กาย\-สมฺผสฺสชา เวทนา ฯเปฯ สญฺญา ฯเปฯ เจตนา ฯเปฯ กายวิญฺญาณํ อุปฺปชฺชิ วา อุปฺปชฺชติ วา อุปฺปชฺชิสฺสติ วา อุปฺปชฺเช วา ฯเปฯ. ยํ โผฏฺฐพฺพา\-รมฺมโณ กายํ นิสฺสาย กายสมฺผสฺโส อุปฺปชฺชิ วา อุปฺปชฺชติ วา อุปฺปชฺชิสฺสติ วา อุปฺปชฺเช วา ฯเปฯ. ยํ โผฏฺฐพฺพา\-รมฺมณา กายํ นิสฺสาย กาย\-สมฺผสฺสชา เวทนา ฯเปฯ สญฺญา ฯเปฯ เจตนา ฯเปฯ กายวิญฺญาณํ อุปฺปชฺชิ วา อุปฺปชฺชติ วา อุปฺปชฺชิสฺสติ วา อุปฺปชฺเช วา, โผฏฺฐพฺโพเปโส โผฏฺฐพฺพา\-ยตนํเปตํ โผฏฺฐพฺพ\-ธาตุเปสา. อิทํ ตํ รูปํ โผฏฺฐพฺพา\-ยตนํ.}

\chapterheadpageread{2. รูปกณฺฑ}
\proseitem{651}{\csromanfolio{171}กตมํ ตํ รูปํ อาโปธาตุ. ยํ อาโป อาโปคตํ สิเนโห สิเนหคตํ พนฺธนตฺตํ รูปสฺส. อิทํ ตํ รูปํ อาโปธาตุ.}

\vspace{\tipitakaparskip}
\csromancenter{อิทํ ตํ รูปํ โน อุปาทา.}
\proseitem{652}{กตมํ ตํ รูปํ อุปาทิณฺณํ. จกฺขายตนํ โสตายตนํ ฆานายตนํ ชิวฺหายตนํ กายายตนํ อิตฺถินฺทฺริยํ ปุริสินฺทฺริยํ ชีวิตินฺทฺริยํ, ยํ วา ปนญฺญมฺปิ อตฺถิ รูปํ กมฺมสฺส กตตฺตา รูปายตนํ คนฺธายตนํ รสายตนํ โผฏฺฐพฺพา\-ยตนํ อากาสธาตุ อาโปธาตุ รูปสฺส อุปจโย รูปสฺส สนฺตติ กพฬีกาโร อาหาโร. อิทํ ตํ รูปํ อุปาทิณฺณํ.}

\proseitem{653}{กตมํ ตํ รูปํ อนุปาทิณฺณํ. สทฺทายตนํ กายวิญฺญตฺติ วจีวิญฺญตฺติ รูปสฺส ลหุตา รูปสฺส มุทุตา รูปสฺส กมฺมญฺญตา รูปสฺส ชรตา รูปสฺส อนิจฺจตา, ยํ วา ปนญฺญมฺปิ อตฺถิ รูปํ น กมฺมสฺส กตตฺตา รูปายตนํ คนฺธายตนํ รสายตนํ โผฏฺฐพฺพา\-ยตนํ อากาสธาตุ อาโปธาตุ รูปสฺส อุปจโย รูปสฺส สนฺตติ กพฬีกาโร อาหาโร. อิทํ ตํ รูปํ อนุปาทิณฺณํ.}

\proseitem{654}{กตมํ ตํ รูปํ อุปาทิณฺณุ\-ปาทานิยํ. จกฺขายตนํ ฯเปฯ กายายตนํ อิตฺถินฺทฺริยํ ปุริสินฺทฺริยํ ชีวิตินฺทฺริยํ, ยํ วา ปนญฺญมฺปิ อตฺถิ รูปํ กมฺมสฺส กตตฺตา รูปายตนํ คนฺธายตนํ รสายตนํ โผฏฺฐพฺพา\-ยตนํ อากาสธาตุ อาโปธาตุ รูปสฺส อุปจโย รูปสฺส สนฺตติ กพฬีกาโร อาหาโร. อิทํ ตํ รูปํ อุปาทิณฺณุ\-ปาทานิยํ.}

\proseitem{655}{กตมํ ตํ รูปํ อนุปาทิณฺณุ\-ปาทานิยํ. สทฺทายตนํ กายวิญฺญตฺติ วจีวิญฺญตฺติ รูปสฺส ลหุตา รูปสฺส มุทุตา รูปสฺส กมฺมญฺญตา รูปสฺส ชรตา รูปสฺส อนิจฺจตา, ยํ วา ปนญฺญมฺปิ อตฺถิ รูปํ น กมฺมสฺส กตตฺตา รูปายตนํ คนฺธายตนํ รสายตนํ โผฏฺฐพฺพา\-ยตนํ อากาสธาตุ อาโปธาตุ รูปสฺส อุปจโย รูปสฺส สนฺตติ กพฬีกาโร อาหาโร. อิทํ ตํ รูปํ อนุปาทิณฺณุ\-ปาทานิยํ.}

\proseitem{656}{กตมํ ตํ รูปํ สนิทสฺสนํ. รูปายตนํ. อิทํ ตํ รูปํ สนิทสฺสนํ.}

\proseitem{657}{\csromanfolio{172}กตมํ ตํ รูปํ อนิทสฺสนํ. จกฺขายตนํ ฯเปฯ กพฬีกาโร อาหาโร. อิทํ ตํ รูปํ อนิทสฺสนํ.}

\proseitem{658}{กตมํ ตํ รูปํ สปฺปฏิฆํ. จกฺขายตนํ โสตายตนํ ฆานายตนํ ชิวฺหายตนํ กายายตนํ รูปายตนํ สทฺทายตนํ คนฺธายตนํ รสายตนํ โผฏฺฐพฺพา\-ยตนํ. อิทํ ตํ รูปํ สปฺปฏิฆํ.}

\proseitem{659}{กตมํ ตํ รูปํ อปฺปฏิฆํ. อิตฺถินฺทฺริยํ ฯเปฯ กพฬีกาโร อาหาโร. อิทํ ตํ รูปํ อปฺปฏิฆํ.}

\proseitem{660}{กตมํ ตํ รูปํ อินฺทฺริยํ. จกฺขุนฺทฺริยํ โสตินฺทฺริยํ ฆานินฺทฺริยํ ชิวฺหินฺทฺริยํ กายินฺทฺริยํ อิตฺถินฺทฺริยํ ปุริสินฺทฺริยํ ชีวิตินฺทฺริยํ. อิทํ ตํ รูปํ อินฺทฺริยํ.}

\proseitem{661}{กตมํ ตํ รูปํ น อินฺทฺริยํ. รูปายตนํ ฯเปฯ กพฬีกาโร อาหาโร. อิทํ ตํ รูปํ น อินฺทฺริยํ.}

\proseitem{662}{กตมํ ตํ รูปํ มหาภูตํ. โผฏฺฐพฺพา\-ยตนํ อาโปธาตุ. อิทํ ตํ รูปํ มหาภูตํ.}

\proseitem{663}{กตมํ ตํ รูปํ น มหาภูตํ. จกฺขายตนํ ฯเปฯ กพฬีกาโร อาหาโร. อิทํ ตํ รูปํ น มหาภูตํ.}

\proseitem{664}{กตมํ ตํ รูปํ วิญฺญตฺติ. กายวิญฺญตฺติ วจีวิญฺญตฺติ. อิทํ ตํ รูปํ วิญฺญตฺติ.}

\proseitem{665}{กตมํ ตํ รูปํ น วิญฺญตฺติ. จกฺขายตนํ ฯเปฯ กพฬีกาโร อาหาโร. อิทํ ตํ รูปํ น วิญฺญตฺติ.}

\proseitem{666}{กตมํ ตํ รูปํ จิตฺต\-สมุฏฺฐานํ. กายวิญฺญตฺติ วจีวิญฺญตฺติ, ยํ วา ปนญฺญมฺปิ อตฺถิ รูปํ จิตฺตชํ จิตฺตเหตุกํ จิตฺต\-สมุฏฺฐานํ รูปายตนํ สทฺทายตนํ คนฺธายตนํ รสายตนํ โผฏฺฐพฺพา\-ยตนํ อากาสธาตุ อาโปธาตุ รูปสฺส ลหุตา รูปสฺส มุทุตา รูปสฺส กมฺมญฺญตา รูปสฺส อุปจโย รูปสฺส สนฺตติ กพฬีกาโร อาหาโร. อิทํ ตํ รูปํ จิตฺต\-สมุฏฺฐานํ.}

\proseitem{667}{กตมํ ตํ รูปํ น จิตฺต\-สมุฏฺฐานํ. จกฺขายตนํ ฯเปฯ กายายตนํ อิตฺถินฺทฺริยํ ปุริสินฺทฺริยํ ชีวิตินฺทฺริยํ รูปสฺส ชรตา รูปสฺส อนิจฺจตา, ยํ วา ปนญฺญมฺปิ อตฺถิ รูปํ น จิตฺตชํ น จิตฺตเหตุกํ น จิตฺต\-สมุฏฺฐานํ รูปายตนํ \csromanfolio{173}สทฺทายตนํ คนฺธายตนํ รสายตนํ โผฏฺฐพฺพา\-ยตนํ อากาสธาตุ อาโปธาตุ รูปสฺส ลหุตา รูปสฺส มุทุตา รูปสฺส กมฺมญฺญตา รูปสฺส อุปจโย รูปสฺส สนฺตติ กพฬีกาโร อาหาโร. อิทํ ตํ รูปํ น จิตฺต\-สมุฏฺฐานํ.}

\proseitem{668}{กตมํ ตํ รูปํ จิตฺตสหภุ. กายวิญฺญตฺติ วจีวิญฺญตฺติ. อิทํ ตํ รูปํ จิตฺตสหภุ.}

\proseitem{669}{กตมํ ตํ รูปํ น จิตฺตสหภุ. จกฺขายตนํ ฯเปฯ กพฬีกาโร อาหาโร. อิทํ ตํ รูปํ น จิตฺตสหภุ.}

\proseitem{670}{กตมํ ตํ รูปํ จิตฺตา\-นุปริวตฺติ. กายวิญฺญตฺติ วจีวิญฺญตฺติ. อิทํ ตํ รูปํ จิตฺตา\-นุปริวตฺติ.}

\proseitem{671}{กตมํ ตํ รูปํ น จิตฺตา\-นุปริวตฺติ. จกฺขายตนํ ฯเปฯ กพฬีกาโร อาหาโร. อิทํ ตํ รูปํ น จิตฺตา\-นุปริวตฺติ.}

\proseitem{672}{กตมํ ตํ รูปํ อชฺฌตฺติกํ. จกฺขายตนํ ฯเปฯ กายายตนํ. อิทํ ตํ รูปํ อชฺฌตฺติกํ.}

\proseitem{673}{กตมํ ตํ รูปํ พาหิรํ. รูปายตนํ ฯเปฯ กพฬีกาโร อาหาโร. อิทํ ตํ รูปํ พาหิรํ.}

\proseitem{674}{กตมํ ตํ รูปํ โอฬาริกํ. จกฺขายตนํ ฯเปฯ. โผฏฺฐพฺพา\-ยตนํ. อิทํ ตํ รูปํ โอฬาริกํ.}

\proseitem{675}{กตมํ ตํ รูปํ สุขุมํ. อิตฺถินฺทฺริยํ ฯเปฯ กพฬีกาโร อาหาโร. อิทํ ตํ รูปํ สุขุมํ.}

\proseitem{676}{กตมํ ตํ รูปํ ทูเร. อิตฺถินฺทฺริยํ ฯเปฯ กพฬีกาโร อาหาโร. อิทํ ตํ รูปํ ทูเร.}

\proseitem{677}{กตมํ ตํ รูปํ สนฺติเก. จกฺขายตนํ ฯเปฯ. โผฏฺฐพฺพา\-ยตนํ. อิทํ ตํ รูปํ สนฺติเก.}

\proseitem{678}{กตมํ ตํ รูปํ จกฺขุ\-สมฺผสฺสสฺส วตฺถุ. จกฺขายตนํ. อิทํ ตํ รูปํ จกฺขุ\-สมฺผสฺสสฺส วตฺถุ.}

\proseitem{679}{\csromanfolio{174}กตมํ ตํ รูปํ จกฺขุ\-สมฺผสฺสสฺส น วตฺถุ. โสตายตนํ ฯเปฯ กพฬีกาโร อาหาโร. อิทํ ตํ รูปํ จกฺขุ\-สมฺผสฺสสฺส น วตฺถุ.}

\proseitem{680}{กตมํ ตํ รูปํ จกฺขุ\-สมฺผสฺสชาย เวทนาย ฯเปฯ สญฺญาย ฯเปฯ เจตนาย ฯเปฯ จกฺขุ\-วิญฺญาณสฺส วตฺถุ. จกฺขายตนํ. อิทํ ตํ รูปํ จกฺขุ\-วิญฺญาณสฺส วตฺถุ.}

\proseitem{681}{กตมํ ตํ รูปํ จกฺขุ\-วิญฺญาณสฺส น วตฺถุ. โสตายตนํ ฯเปฯ กพฬีกาโร อาหาโร. อิทํ ตํ รูปํ จกฺขุ\-วิญฺญาณสฺส น วตฺถุ.}

\proseitem{682}{กตมํ ตํ รูปํ โสต\-สมฺผสฺสสฺส ฯเปฯ ฆาน\-สมฺผสฺสสฺส ฯเปฯ ชิวฺหา\-สมฺผสฺสสฺส ฯเปฯ กาย\-สมฺผสฺสสฺส วตฺถุ. กายายตนํ. อิทํ ตํ รูปํ กาย\-สมฺผสฺสสฺส วตฺถุ.}

\proseitem{683}{กตมํ ตํ รูปํ กาย\-สมฺผสฺสสฺส น วตฺถุ. จกฺขายตนํ ฯเปฯ กพฬีกาโร อาหาโร. อิทํ ตํ รูปํ กาย\-สมฺผสฺสสฺส น วตฺถุ.}

\proseitem{684}{กตมํ ตํ รูปํ กาย\-สมฺผสฺสชาย เวทนาย ฯเปฯ สญฺญาย ฯเปฯ เจตนาย ฯเปฯ กายวิญฺญาณสฺส วตฺถุ. กายายตนํ. อิทํ ตํ รูปํ กายวิญฺญาณสฺส วตฺถุ.}

\proseitem{685}{กตมํ ตํ รูปํ กายวิญฺญาณสฺส น วตฺถุ. จกฺขายตนํ ฯเปฯ กพฬีกาโร อาหาโร. อิทํ ตํ รูปํ กายวิญฺญาณสฺส น วตฺถุ.}

\proseitem{686}{กตมํ ตํ รูปํ จกฺขุ\-สมฺผสฺสสฺส อารมฺมณํ. รูปายตนํ. อิทํ ตํ รูปํ จกฺขุ\-สมฺผสฺสสฺส อารมฺมณํ.}

\proseitem{687}{กตมํ ตํ รูปํ จกฺขุ\-สมฺผสฺสสฺส น อารมฺมณํ. จกฺขายตนํ ฯเปฯ กพฬีกาโร อาหาโร. อิทํ ตํ รูปํ จกฺขุ\-สมฺผสฺสสฺส น อารมฺมณํ.}

\proseitem{688}{กตมํ ตํ รูปํ จกฺขุ\-สมฺผสฺสชาย เวทนาย ฯเปฯ สญฺญาย ฯเปฯ เจตนาย ฯเปฯ จกฺขุ\-วิญฺญาณสฺส อารมฺมณํ. รูปายตนํ. อิทํ ตํ รูปํ จกฺขุ\-วิญฺญาณสฺส อารมฺมณํ.}

\proseitem{689}{กตมํ ตํ รูปํ จกฺขุ\-วิญฺญาณสฺส น อารมฺมณํ. จกฺขายตนํ ฯเปฯ กพฬีกาโร อาหาโร. อิทํ ตํ รูปํ จกฺขุ\-วิญฺญาณสฺส น อารมฺมณํ.}

\chapterheadpageread{2. รูปกณฺฑ}
\proseitem{690}{\csromanfolio{175}กตมํ ตํ รูปํ โสต\-สมฺผสฺสสฺส ฯเปฯ ฆาน\-สมฺผสฺสสฺส ฯเปฯ ชิวฺหา\-สมฺผสฺสสฺส ฯเปฯ กาย\-สมฺผสฺสสฺส อารมฺมณํ. โผฏฺฐพฺพา\-ยตนํ. อิทํ ตํ รูปํ กาย\-สมฺผสฺสสฺส อารมฺมณํ.}

\proseitem{691}{กตมํ ตํ รูปํ กาย\-สมฺผสฺสสฺส น อารมฺมณํ. จกฺขายตนํ ฯเปฯ กพฬีกาโร อาหาโร. อิทํ ตํ รูปํ กาย\-สมฺผสฺสสฺส น อารมฺมณํ.}

\proseitem{692}{กตมํ ตํ รูปํ กาย\-สมฺผสฺสชาย เวทนาย ฯเปฯ สญฺญาย ฯเปฯ เจตนาย ฯเปฯ กายวิญฺญาณสฺส อารมฺมณํ. โผฏฺฐพฺพา\-ยตนํ. อิทํ ตํ รูปํ กายวิญฺญาณสฺส อารมฺมณํ.}

\proseitem{693}{กตมํ ตํ รูปํ กายวิญฺญาณสฺส น อารมฺมณํ. จกฺขายตนํ ฯเปฯ กพฬีกาโร อาหาโร. อิทํ ตํ รูปํ กายวิญฺญาณสฺส น อารมฺมณํ.}

\proseitem{694}{กตมํ ตํ รูปํ จกฺขายตนํ. ยํ จกฺขุ จตุนฺนํ มหาภูตานํ อุปาทาย ปสาโท ฯเปฯ สุญฺโญ คาโมเปโส. อิทํ ตํ รูปํ จกฺขายตนํ.}

\proseitem{695}{กตมํ ตํ รูปํ น จกฺขายตนํ. โสตายตนํ ฯเปฯ กพฬีกาโร อาหาโร. อิทํ ตํ รูปํ น จกฺขายตนํ.}

\proseitem{696}{กตมํ ตํ รูปํ โสตายตนํ ฯเปฯ ฆานายตนํ ฯเปฯ ชิวฺหายตนํ ฯเปฯ กายายตนํ. โย กาโย จตุนฺนํ มหาภูตานํ อุปาทาย ปสาโท ฯเปฯ สุญฺโญ คาโมเปโส. อิทํ ตํ รูปํ กายายตนํ.}

\proseitem{697}{กตมํ ตํ รูปํ น กายายตนํ. จกฺขายตนํ ฯเปฯ กพฬีกาโร อาหาโร. อิทํ ตํ รูปํ น กายายตนํ.}

\proseitem{698}{กตมํ ตํ รูปํ รูปายตนํ. ยํ รูปํ จตุนฺนํ มหาภูตานํ อุปาทาย วณฺณนิภา ฯเปฯ รูปธาตุเปสา. อิทํ ตํ รูปํ รูปายตนํ.}

\proseitem{699}{กตมํ ตํ รูปํ น รูปายตนํ. จกฺขายตนํ ฯเปฯ กพฬีกาโร อาหาโร. อิทํ ตํ รูปํ น รูปายตนํ.}

\proseitem{700}{กตมํ ตํ รูปํ สทฺทายตนํ ฯเปฯ คนฺธายตนํ ฯเปฯ รสายตนํ ฯเปฯ โผฏฺฐพฺพา\-ยตนํ. ปถวีธาตุ ฯเปฯ โผฏฺฐพฺพ\-ธาตุเปสา. อิทํ ตํ รูปํ โผฏฺฐพฺพา\-ยตนํ.}

\proseitem{701}{กตมํ ตํ รูปํ น โผฏฺฐพฺพา\-ยตนํ. จกฺขายตนํ ฯเปฯ กพฬีกาโร อาหาโร. อิทํ ตํ รูปํ น โผฏฺฐพฺพา\-ยตนํ.}

\proseitem{702}{\csromanfolio{176}กตมํ ตํ รุปํ จกฺขุธาตุ. จกฺขายตนํ. อิทํ ตํ รูปํ จกฺขุธาตุ.}

\proseitem{703}{กตมํ ตํ รูปํ น จกฺขุธาตุ. โสตายตนํ ฯเปฯ กพฬีกาโร อาหาโร. อิทํ ตํ รูปํ น จกฺขุธาตุ.}

\proseitem{704}{กตมํ ตํ รูปํ โสตธาตุ ฯเปฯ ฆานธาตุ ฯเปฯ ชิวฺหาธาตุ ฯเปฯ กายธาตุ. กายายตนํ. อิทํ ตํ รูปํ กายธาตุ.}

\proseitem{705}{กตมํ ตํ รูปํ น กายธาตุ. จกฺขายตนํ ฯเปฯ กพฬีกาโร อาหาโร. อิทํ ตํ รูปํ น กายธาตุ.}

\proseitem{706}{กตมํ ตํ รูปํ รูปธาตุ. รูปายตนํ. อิทํ ตํ รูปํ รูปธาตุ.}

\proseitem{707}{กตมํ ตํ รูปํ น รูปธาตุ. จกฺขายตนํ ฯเปฯ กพฬีกาโร อาหาโร. อิทํ ตํ รูปํ น รูปธาตุ.}

\proseitem{708}{กตมํ ตํ รูปํ สทฺทธาตุ ฯเปฯ คนฺธธาตุ ฯเปฯ รสธาตุ ฯเปฯ โผฏฺฐพฺพ\-ธาตุ. โผฏฺฐพฺพา\-ยตนํ. อิทํ ตํ รูปํ โผฏฺฐพฺพ\-ธาตุ.}

\proseitem{709}{กตมํ ตํ รูปํ น โผฏฺฐพฺพ\-ธาตุ. จกฺขายตนํ ฯเปฯ กพฬีกาโร อาราโห. อิทํ ตํ รูปํ น โผฏฺฐพฺพ\-ธาตุ.}

\proseitem{710}{กตมํ ตํ รูปํ จกฺขุนฺทฺริยํ. ยํ จกฺขุ จตุนฺนํ มหาภูตานํ อุปาทาย ปสาโท ฯเปฯ สุญฺโญ คาโมเปโส. อิทํ ตํ รูปํ จกฺขุนฺทฺริยํ.}

\proseitem{711}{กตมํ ตํ รูปํ น จกฺขุนฺทฺริยํ. โสตายตนํ ฯเปฯ กพฬีกาโร อาหาโร. อิทํ ตํ รูปํ น จกฺขุนฺทฺริยํ.}

\proseitem{712}{กตมํ ตํ รูปํ โสตินฺทฺริยํ ฯเปฯ ฆานินฺทฺริยํ ฯเปฯ ชิวฺหินฺทฺริยํ ฯเปฯ กายินฺทฺริยํ. โย กาโย จตุนฺนํ มหาภูตานํ อุปาทาย ปสาโท ฯเปฯ สุญฺโญ คาโมเปโส. อิทํ ตํ รูปํ กายินฺทฺริยํ.}

\proseitem{713}{กตมํ ตํ รูปํ น กายินฺทฺริยํ. จกฺขายตนํ. กพฬีกาโร อาหาโร. อิทํ ตํ รูปํ น กายินฺทฺริยํ.}

\proseitem{714}{กตมํ ตํ รูปํ อิตฺถินฺทฺริยํ. ยํ อิตฺถิยา อิตฺถิลิงฺคํ อิตฺถินิมิตฺตํ อิตฺถิกุตฺตํ อิตฺถากปฺโป อิตฺถตฺตํ อิตฺถิภาโว. อิทํ ตํ รูปํ อิตฺถินฺทฺริยํ.}

\chapterheadpageread{2. รูปกณฺฑ}
\proseitem{715}{\csromanfolio{177}กตมํ ตํ รูปํ น อิตฺถินฺทฺริยํ. จกฺขายตนํ ฯเปฯ กพฬีกาโร อาหาโร. อิทํ ตํ รูปํ น อิตฺถินฺทฺริยํ.}

\proseitem{716}{กตมํ ตํ รูปํ ปุริสินฺทฺริยํ. ยํ ปุริสสฺส ปุริสลิงฺคํ ปุริสนิมิตฺตํ ปุริสกุตฺตํ ปุริสากปฺโป ปุริสตฺตํ ปุริสภาโว. อิทํ ตํ รูปํ ปุริสินฺทฺริยํ.}

\proseitem{717}{กตมํ ตํ รูปํ น ปุริสินฺทฺริยํ. จกฺขายตนํ ฯเปฯ กพฬีกาโร อาหาโร. อิทํ ตํ รูปํ น ปุริสินฺทฺริยํ.}

\proseitem{718}{กตมํ ตํ รูปํ ชีวิตินฺทฺริยํ. โย เตสํ รูปีนํ ธมฺมานํ อายุ ฐิติ ยปนา ยาปนา อิริยนา วตฺตนา ปาลนา ชีวิตํ ชีวิตินฺทฺริยํ. อิทํ ตํ รูปํ ชีวิตินฺทฺริยํ.}

\proseitem{719}{กตมํ ตํ รูปํ น ชีวิตินฺทฺริยํ. จกฺขายตนํ ฯเปฯ กพฬีกาโร อาหาโร. อิทํ ตํ รูปํ น ชีวิตินฺทฺริยํ.}

\proseitem{720}{กตมํ ตํ รูปํ กายวิญฺญตฺติ. ยา กุสลจิตฺตสฺส วา อกุสลจิตฺตสฺส วา อพฺยากต\-จิตฺตสฺส วา อภิกฺกม\-นฺตสฺส วา ปฏิกฺกมนฺตสฺส วา อาโลเกนฺตสฺส วา วิโลเกนฺตสฺส วา สมิญฺเชนฺตสฺส วา ปสาเรนฺตสฺส วา กายสฺส ถมฺภนา สนฺถมฺภนา สนฺถมฺภิตตฺตํ วิญฺญตฺติ วิญฺญาปนา วิญฺญาปิตตฺตํ. อิทํ ตํ รูปํ กายวิญฺญตฺติ.}

\proseitem{721}{กตมํ ตํ รูปํ น กายวิญฺญตฺติ. จกฺขายตนํ ฯเปฯ กพฬีกาโร อาหาโร. อิทํ ตํ รูปํ น กายวิญฺญตฺติ.}

\proseitem{722}{กตมํ ตํ รูปํ วจีวิญฺญตฺติ. ยา กุสลจิตฺตสฺส วา อกุสลจิตฺตสฺส วา อพฺยากต\-จิตฺตสฺส วา วาจา คิรา พฺยปฺปโถ อุทีรณํ โฆโส โฆสกมฺมํ วาจา วจีเภโท, อยํ วุจฺจติ วาจา. ยา ตาย วาจาย วิญฺญตฺติ วิญฺญาปนา วิญฺญาปิตตฺตํ. อิทํ ตํ รูปํ วจีวิญฺญตฺติ.}

\proseitem{723}{กตมํ ตํ รูปํ น วจีวิญฺญตฺติ. จกฺขายตนํ ฯเปฯ กพฬีกาโร อาหาโร. อิทํ ตํ รูปํ น วจีวิญฺญตฺติ.}

\proseitem{724}{กตมํ ตํ รูปํ อากาสธาตุ. โย อากาโส อากาสคตํ อฆํ อฆคตํ วิวโร วิวรคตํ อสมฺผุฏฺฐํ จตูหิ มหาภูเตหิ. อิทํ ตํ รูปํ อากาสธาตุ.}

\proseitem{725}{\csromanfolio{178}กตมํ ตํ รูปํ น อากาสธาตุ. จกฺขายตนํ ฯเปฯ กพฬีกาโร อาหาโร. อิทํ ตํ รูปํ น อากาสธาตุ.}

\proseitem{726}{กตมํ ตํ รูปํ อาโปธาตุ. ยํ อาโป อาโปคตํ สิเนโห สิเนหคตํ พนฺธนตฺตํ รูปสฺส. อิทํ ตํ รูปํ อาโปธาตุ.}

\proseitem{727}{กตมํ ตํ รูปํ น อาโปธาตุ. จกฺขายตนํ ฯเปฯ กพฬีกาโร อาหาโร. อิทํ ตํ รูปํ น อาโปธาตุ.}

\proseitem{728}{กตมํ ตํ รูปํ รูปสฺส ลหุตา. ยา รูปสฺส ลหุตา ลหุปริณามตา อทนฺธนตา อวิตฺถนตา. อิทํ ตํ รูปํ รูปสฺส ลหุตา.}

\proseitem{729}{กตมํ ตํ รูปํ รูปสฺส น ลหุตา. จกฺขายตนํ ฯเปฯ กพฬีกาโร อาหาโร. อิทํ ตํ รูปํ รูปสฺส น ลหุตา.}

\proseitem{730}{กตมํ ตํ รูปํ รูปสฺส มุทุตา. ยา รูปสฺส มุทุตา มทฺทวตา อกกฺขฬตา อกถินตา. อิทํ ตํ รูปํ รูปสฺส มุทุตา.}

\proseitem{731}{กตมํ ตํ รูปํ รูปสฺส น มุทุตา. จกฺขายตนํ ฯเปฯ กพฬีกาโร อาหาโร. อิทํ ตํ รูปํ รูปสฺส น มุทุตา.}

\proseitem{732}{กตมํ ตํ รูปํ รูปสฺส กมฺมญฺญตา. ยา รูปสฺส กมฺมญฺญตา กมฺมญฺญตฺตํ กมฺมญฺญภาโว. อิทํ ตํ รูปํ รูปสฺส กมฺมญฺญตา.}

\proseitem{733}{กตมํ ตํ รูปํ รูปสฺส น กมฺมญฺญตา. จกฺขายตนํ ฯเปฯ กพฬีกาโร อาหาโร. อิทํ ตํ รูปํ รูปสฺส น กมฺมญฺญตา.}

\proseitem{734}{กตมํ ตํ รูปํ รูปสฺส อุปจโย. โย อายตนานํ อาจโย โส รูปสฺส อุปจโย. อิทํ ตํ รูปํ รูปสฺส อุปจโย.}

\proseitem{735}{กตมํ ตํ รูปํ รูปสฺส น อุปจโย. จกฺขายตนํ ฯเปฯ กพฬีกาโร อาหาโร. อิทํ ตํ รูปํ รูปสฺส น อุปจโย.}

\proseitem{736}{กตมํ ตํ รูปํ รูปสฺส สนฺตติ. โย รูปสฺส อุปจโย, สา รูปสฺส สนฺตติ. อิทํ ตํ รูปํ รูปสฺส สนฺตติ.}

\proseitem{737}{กตมํ ตํ รูปํ รูปสฺส น สนฺตติ. จกฺขายตนํ ฯเปฯ กพฬีกาโร อาหาโร. อิทํ ตํ รูปํ รูปสฺส น สนฺตติ.}

\chapterheadpageread{2. รูปกณฺฑ}
\proseitem{738}{\csromanfolio{179}กตมํ ตํ รูปํ รูปสฺส ชรตา. ยา รูปสฺส ชรา ชีรณตา ขณฺฑิจฺจํ ปาลิจฺจํ วลิตฺตจตา อายุโน สํหานิ อินฺทฺริยานํ ปริปาโก. อิทํ ตํ รูปํ รูปสฺส ชรตา.}

\proseitem{739}{กตมํ ตํ รูปํ รูปสฺส น ชรตา. จกฺขายตนํ ฯเปฯ กพฬีกาโร อาหาโร. อิทํ ตํ รูปํ รูปสฺส น ชรตา.}

\proseitem{740}{กตมํ ตํ รูปํ รูปสฺส อนิจฺจตา. โย รูปสฺส ขโย วโย เภโท ปริเภโท อนิจฺจตา อนฺตรธานํ. อิทํ ตํ รูปํ รูปสฺส อนิจฺจตา.}

\proseitem{741}{กตมํ ตํ รูปํ รูปสฺส น อนิจฺจตา. จกฺขายตนํ ฯเปฯ กพฬีกาโร อาหาโร. อิทํ ตํ รูปํ รูปสฺส น อนิจฺจตา.}

\proseitem{742}{กตมํ ตํ รูปํ กพฬีกาโร อาหาโร. โอทโน กุมฺมาโส สตฺตุ มจฺโฉ มํสํ ขีรํ ทธิ สปฺปิ นวนีตํ เตลํ มธุ ผาณิตํ, ยํ วา ปนญฺญมฺปิ อตฺถิ รูปํ ยมฺหิ ยมฺหิ ชนปเท เตสํ เตสํ สตฺตานํ มุขาสิยํ ทนฺตวิขาทนํ คลชฺโฌหรณียํ กุจฺฉิ\-วิตฺถมฺภนํ ยาย โอชาย สตฺตา ยาเปนฺติ. อิทํ ตํ รูปํ กพฬีกาโร อาหาโร.}

\proseitem{743}{กตมํ ตํ รูปํ น กพฬีกาโร อาหาโร. จกฺขายตนํ ฯเปฯ รูปสฺส อนิจฺจตา. อิทํ ตํ รูปํ น กพฬีกาโร อาหาโร.เอวํ ทุวิเธน รูปสงฺคโห.ทุกนิทฺเทโส.}

\csromanmatikaanchor{mtk.93}
\csromantocmarkref{subsection}{ติกนิทฺเทส}{mtk.93}
\bookmark[dest={mtk.93},level=2]{ติกนิทฺเทส}
\csromanheader{2}{\textbf{ติกนิทฺเทส}}
\proseitem{744}{กตมํ ตํ รูปํ อชฺฌตฺติกํ อุปาทา. จกฺขายตนํ ฯเปฯ กายายตนํ. อิทํ ตํ รูปํ อชฺฌตฺติกํ อุปาทา.}

\proseitem{745}{กตมํ ตํ รูปํ พาหิรํ อุปาทา. รูปายตนํ ฯเปฯ กพฬีกาโร อาหาโร. อิทํ ตํ รูปํ พาหิรํ อุปาทา.}

\proseitem{746}{กตมํ ตํ รูปํ พาหิรํ โน อุปาทา. โผฏฺฐพฺพา\-ยตนํ อาโปธาตุ. อิทํ ตํ รูปํ พาหิรํ โน อุปาทา.}

\proseitem{747}{\csromanfolio{180}กตมํ ตํ รูปํ อชฺฌตฺติกํ อุปาทิณฺณํ. จกฺขายตนํ ฯเปฯ กายายตนํ. อิทํ ตํ รูปํ อชฺฌตฺติกํ อุปาทิณฺณํ.}

\proseitem{748}{กตมํ ตํ รูปํ พาหิรํ อุปาทิณฺณํ. อิตฺถินฺทฺริยํ ปุริสินฺทฺริยํ ชีวิตินฺทฺริยํ, ยํ วา ปนญฺญมฺปิ อตฺถิ รูปํ กมฺมสฺส กตตฺตา รูปายตนํ คนฺธายตนํ รสายตนํ โผฏฺฐพฺพา\-ยตนํ อากาสธาตุ อาโปธาตุ รูปสฺส อุปจโย รูปสฺส สนฺตติ กพฬีกาโร อาหาโร. อิทํ ตํ รูปํ พาหิรํ อุปาทิณฺณํ.}

\proseitem{749}{กตมํ ตํ รูปํ พาหิรํ อนุปาทิณฺณํ. สทฺทายตนํ กายวิญฺญตฺติ วจีวิญฺญตฺติ รูปสฺส ลหุตา รูปสฺส มุทุตา รูปสฺส กมฺมญฺญตา รูปสฺส ชรตา รูปสฺส อนิจฺจตา, ยํ วา ปนญฺญมฺปิ อตฺถิ รูปํ น กมฺมสฺส กตตฺตา รูปายตนํ คนฺธายตนํ รสายตนํ โผฏฺฐพฺพา\-ยตนํ อากาสธาตุ อาโปธาตุ รูปสฺส อุปจโย รูปสฺส สนฺตติ กพฬีกาโร อาหาโร. อิทํ ตํ รูปํ พาหิรํ อนุปาทิณฺณํ.}

\proseitem{750}{กตมํ ตํ รูปํ อชฺฌตฺติกํ อุปาทิณฺณุ\-ปาทานิยํ. จกฺขายตนํ ฯเปฯ กายายตนํ. อิทํ ตํ รูปํ อชฺฌตฺติกํ อุปาทิณฺณุ\-ปาทานิยํ.}

\proseitem{751}{กตมํ ตํ รูปํ พาหิรํ อุปาทิณฺณุ\-ปาทานิยํ. อิตฺถินฺทฺริยํ ปุริสินฺทฺริยํ ชีวิตินฺทฺริยํ, ยํ วา ปนญฺญมฺปิ อตฺถิ รูปํ กมฺมสฺส กตตฺตา รูปายตนํ คนฺธายตนํ รสายตนํ โผฏฺฐพฺพา\-ยตนํ อากาสธาตุ อาโปธาตุ รูปสฺส อุปจโย รูปสฺส สนฺตติ กพฬีกาโร อาหาโร. อิทํ ตํ รูปํ พาหิรํ อุปาทิณฺณุ\-ปาทานิยํ.}

\proseitem{752}{กตมํ ตํ รูปํ พาหิรํ อนุปาทิณฺณุ\-ปาทานิยํ. สทฺทายตนํ กายวิญฺญตฺติ วจีวิญฺญตฺติ รูปสฺส ลหุตา รูปสฺส มุทุตา รูปสฺส กมฺมญฺญตา รูปสฺส ชรตา รูปสฺส อนิจฺจตา, ยํ วา ปนญฺญมฺปิ อตฺถิ รูปํ น กมฺมสฺส กตตฺตา รูปายตนํ คนฺธายตนํ รสายตนํ โผฏฺฐพฺพา\-ยตนํ อากาสธาตุ อาโปธาตุ รูปสฺส อุปจโย รูปสฺส สนฺตติ กพฬีกาโร อาหาโร. อิทํ ตํ รูปํ พาหิรํ อนุปาทิณฺณุ\-ปาทานิยํ.}

\proseitem{753}{กตมํ ตํ รูปํ อชฺฌตฺติกํ อนิทสฺสนํ. จกฺขายตนํ ฯเปฯ กายายตนํ. อิทํ ตํ รูปํ อชฺฌตฺติกํ อนิทสฺสนํ.}

\proseitem{754}{กตมํ ตํ รูปํ พาหิรํ สนิทสฺสนํ. รูปายตนํ. อิทํ ตํ รูปํ พาหิรํ สนิทสฺสนํ.}

\chapterheadpageread{2. รูปกณฺฑ}
\proseitem{755}{\csromanfolio{181}กตมํ ตํ รูปํ พาหิรํ อนิทสฺสนํ. สทฺทายตนํ ฯเปฯ กพฬีกาโร อาหาโร. อิทํ ตํ รูปํ พาหิรํ อนิทสฺสนํ.}

\proseitem{756}{กตมํ ตํ รูปํ อชฺฌตฺติกํ สปฺปฏิฆํ. จกฺขายตนํ ฯเปฯ กายายตนํ. อิทํ ตํ รูปํ อชฺฌตฺติกํ สปฺปฏิฆํ.}

\proseitem{757}{กตมํ ตํ รูปํ พาหิรํ สปฺปฏิฆํ. รูปายตนํ ฯเปฯ โผฏฺฐพฺพา\-ยตนํ. อิทํ ตํ รูปํ พาหิรํ สปฺปฏิฆํ.}

\proseitem{758}{กตมํ ตํ รูปํ พาหิรํ อปฺปฏิฆํ. อิตฺถินฺทฺริยํ ฯเปฯ กพฬีกาโร อาหาโร. อิทํ ตํ รูปํ พาหิรํ อปฺปฏิฆํ.}

\proseitem{759}{กตมํ ตํ รูปํ อชฺฌตฺติกํ อินฺทฺริยํ. จกฺขุนฺทฺริยํ ฯเปฯ กายินฺทฺริยํ. อิทํ ตํ รูปํ อชฺฌตฺติกํ อินฺทฺริยํ.}

\proseitem{760}{กตมํ ตํ รูปํ พาหิรํ อินฺทฺริยํ. อิตฺถินฺทฺริยํ ปุริสินฺทฺริยํ ชีวิตินฺทฺริยํ. อิทํ ตํ รูปํ พาหิรํ อินฺทฺริยํ.}

\proseitem{761}{กตมํ ตํ รูปํ พาหิรํ น อินฺทฺริยํ. รูปายตนํ ฯเปฯ กพฬีกาโร อาหาโร. อิทํ ตํ รูปํ พาหิรํ น อินฺทฺริยํ.}

\proseitem{762}{กตมํ ตํ รูปํ อชฺฌตฺติกํ น มหาภูตํ. จกฺขายตนํ ฯเปฯ กายายตนํ. อิทํ ตํ รูปํ อชฺฌตฺติกํ น มหาภูตํ.}

\proseitem{763}{กตมํ ตํ รูปํ พาหิรํ มหาภูตํ. โผฏฺฐพฺพา\-ยตนํ อาโปธาตุ. อิทํ ตํ รูปํ พาหิรํ มหาภูตํ.}

\proseitem{764}{กตมํ ตํ รูปํ พาหิรํ น มหาภูตํ. รูปายตนํ ฯเปฯ กพฬีกาโร อาหาโร. อิทํ ตํ รูปํ พาหิรํ น มหาภูตํ.}

\proseitem{765}{กตมํ ตํ รูปํ อชฺฌตฺติกํ น วิญฺญตฺติ. จกฺขายตนํ ฯเปฯ กายายตนํ. อิทํ ตํ รูปํ อชฺฌตฺติกํ น วิญฺญตฺติ.}

\proseitem{766}{กตมํ ตํ รูปํ พาหิรํ วิญฺญตฺติ. กายวิญฺญตฺติ วจีวิญฺญตฺติ. อิทํ ตํ รูปํ พาหิรํ วิญฺญตฺติ.}

\proseitem{767}{กตมํ ตํ รูปํ พาหิรํ น วิญฺญตฺติ. รูปายตนํ ฯเปฯ กพฬีกาโร อาหาโร. อิทํ ตํ รูปํ พาหิรํ น วิญฺญตฺติ.}

\proseitem{768}{\csromanfolio{182}กตมํ ตํ รูปํ อชฺฌตฺติกํ น จิตฺต\-สมุฏฺฐานํ. จกฺขายตนํ ฯเปฯ กายายตนํ. อิทํ ตํ รูปํ อชฺฌตฺติกํ น จิตฺต\-สมุฏฺฐานํ.}

\proseitem{769}{กตมํ ตํ รูปํ พาหิรํ จิตฺต\-สมุฏฺฐานํ. กายวิญฺญตฺติ วจีวิญฺญตฺติ, ยํ วา ปนญฺญมฺปิ อตฺถิ รูปํ จิตฺตชํ จิตฺตเหตุกํ จิตฺต\-สมุฏฺฐานํ รูปายตนํ สทฺทายตนํ คนฺธายตนํ รสายตนํ โผฏฺฐพฺพา\-ยตนํ อากาสธาตุ อาโปธาตุ รูปสฺส ลหุตา รูปสฺส มุทุตา รูปสฺส กมฺมญฺญตา รูปสฺส อุปจโย รูปสฺส สนฺตติ กพฬีกาโร อาหาโร. อิทํ ตํ รูปํ พาหิรํ จิตฺต\-สมุฏฺฐานํ.}

\proseitem{770}{กตมํ ตํ รูปํ พหิรํ น จิตฺต\-สมุฏฺฐานํ. อิตฺถินฺทฺริยํ ปุริสินฺทฺริยํ ชีวิตินฺทฺริยํ รูปสฺส ชรตา รูปสฺส อนิจฺจตา, ยํ วา ปนญฺญมฺปิ อตฺถิ รูปํ น จิตฺตชํ น จิตฺตเหตุกํ น จิตฺต\-สมุฏฺฐานํ รูปายตนํ สทฺทายตนํ คนฺธายตนํ รสายตนํ โผฏฺฐพฺพา\-ยตนํ อากาสธาตุ อาโปธาตุ รูปสฺส ลหุตา รูปสฺส มุทุตา รูปสฺส กมฺมญฺญตา รูปสฺส อุปจโย รูปสฺส สนฺตติ กพฬีกาโร อาหาโร. อิทํ ตํ รูปํ พาหิรํ น จิตฺต\-สมุฏฺฐานํ.}

\proseitem{771}{กตมํ ตํ รูปํ อชฺฌตฺติกํ น จิตฺตสหภุ. จกฺขายตนํ ฯเปฯ กายายตนํ. อิทํ ตํ รูปํ อชฺฌตฺติกํ น จิตฺตสหภุ.}

\proseitem{772}{กตมํ ตํ รูปํ พาหิรํ จิตฺตสหภุ. กายวิญฺญตฺติ วจีวิญฺญตฺติ. อิทํ ตํ รูปํ พาหิรํ จิตฺตสหภุ.}

\proseitem{773}{กตมํ ตํ รูปํ พาหิรํ น จิตฺตสหภุ. รูปายตนํ ฯเปฯ กพฬีกาโร อาหาโร. อิทํ ตํ รูปํ พาหิรํ น จิตฺตสหภุ.}

\proseitem{774}{กตมํ ตํ รูปํ อชฺฌตฺติกํ น จิตฺตา\-นุปริวตฺติ. จกฺขายตนํ ฯเปฯ กายายตนํ. อิทํ ตํ รูปํ อชฺฌตฺติกํ น จิตฺตา\-นุปริวตฺติ.}

\proseitem{775}{กตมํ ตํ รูปํ พาหิรํ จิตฺตา\-นุปริวตฺติ. กายวิญฺญตฺติ วจีวิญฺญตฺติ. อิทํ ตํ รูปํ พาหิรํ จิตฺตา\-นุปริวตฺติ.}

\proseitem{776}{กตมํ ตํ รูปํ พาหิรํ น จิตฺตา\-นุปริวตฺติ. รูปายตนํ ฯเปฯ กพฬีกาโร อาหาโร. อิทํ ตํ รูปํ พาหิรํ น จิตฺตา\-นุปริวตฺติ.}

\proseitem{777}{กตมํ ตํ รูปํ อชฺฌตฺติกํ โอฬาริกํ. จกฺขายตนํ ฯเปฯ กายายตนํ. อิทํ ตํ รูปํ อชฺฌตฺติกํ โอฬาริกํ.}

\chapterheadpageread{2. รูปกณฺฑ}
\proseitem{778}{\csromanfolio{183}กตมํ ตํ รูปํ พาหิรํ โอฬาริกํ. รูปายตนํ ฯเปฯ โผฏฺฐพฺพา\-ยตนํ. อิทํ ตํ รูปํ พาหิรํ โอฬาริกํ.}

\proseitem{779}{กตมํ ตํ รูปํ พาหิรํ สุขุมํ. อิตฺถินฺทฺริยํ ฯเปฯ กพฬีกาโร อาหาโร. อิทํ ตํ รูปํ พาหิรํ สุขุมํ.}

\proseitem{780}{กตมํ ตํ รูปํ อชฺฌตฺติกํ สนฺติเก. จกฺขายตนํ ฯเปฯ กายายตนํ. อิทํ ตํ รูปํ อชฺฌตฺติกํ สนฺติเก.}

\proseitem{781}{กตมํ ตํ รูปํ พาหิรํ ทูเร. อิตฺถินฺทฺริยํ ฯเปฯ กพฬีกาโร อาหาโร. อิทํ ตํ รูปํ พาหิรํ ทูเร.}

\proseitem{782}{กตมํ ตํ รูปํ พาหิรํ สนฺติเก. รูปายตนํ ฯเปฯ โผฏฺฐพฺพา\-ยตนํ. อิทํ ตํ รูปํ พาหิรํ สนฺติเก.}

\proseitem{783}{กตมํ ตํ รูปํ พาหิรํ จกฺขุ\-สมฺผสฺสสฺส น วตฺถุ. รูปายตนํ ฯเปฯ กพฬีกาโร อาหาโร. อิทํ ตํ รูปํ พาหิรํ จกฺขุ\-สมฺผสฺสสฺส น วตฺถุ.}

\proseitem{784}{กตมํ ตํ รูปํ อชฺฌตฺติกํ จกฺขุ\-สมฺผสฺสสฺส วตฺถุ. จกฺขายตนํ. อิทํ ตํ รูปํ อชฺฌตฺติกํ จกฺขุ\-สมฺผสฺสสฺส วตฺถุ.}

\proseitem{785}{กตมํ ตํ รูปํ อชฺฌตฺติกํ จกฺขุ\-สมฺผสฺสสฺส น วตฺถุ. โสตายตนํ ฯเปฯ กายายตนํ. อิทํ ตํ รูปํ อชฺฌตฺติกํ จกฺขุ\-สมฺผสฺสสฺส น วตฺถุ.}

\proseitem{786}{กตมํ ตํ รูปํ พาหิรํ จกฺขุ\-สมฺผสฺสชาย เวทนาย ฯเปฯ สญฺญาย ฯเปฯ เจตนาย ฯเปฯ จกฺขุ\-วิญฺญาณสฺส น วตฺถุ. รูปายตนํ ฯเปฯ กพฬีกาโร อาหาโร. อิทํ ตํ รูปํ พาหิรํ จกฺขุ\-วิญฺญาณสฺส น วตฺถุ.}

\proseitem{787}{กตมํ ตํ รูปํ อชฺฌตฺติกํ จกฺขุ\-วิญฺญาณสฺส วตฺถุ. จกฺขายตนํ. อิทํ ตํ รูปํ อชฺฌตฺติกํ จกฺขุ\-วิญฺญาณสฺส วตฺถุ.}

\proseitem{788}{กตมํ ตํ รูปํ อชฺฌตฺติกํ จกฺขุ\-วิญฺญาณสฺส น วตฺถุ. โสตายตนํ ฯเปฯ กายายตนํ. อิทํ ตํ รูปํ อชฺฌตฺติกํ จกฺขุ\-วิญฺญาณสฺส น วตฺถุ.}

\proseitem{789}{กตมํ ตํ รูปํ พาหิรํ โสต\-สมฺผสฺสสฺส ฯเปฯ ฆาน\-สมฺผสฺสสฺส ฯเปฯ ชิวฺหา\-สมฺผสฺสสฺส ฯเปฯ กาย\-สมฺผสฺสสฺส น วตฺถุ. รูปายตนํ ฯเปฯ กพฬีกาโร อาหาโร. อิทํ ตํ รูปํ พาหิรํ กาย\-สมฺผสฺสสฺส น วตฺถุ.}

\proseitem{790}{กตมํ ตํ รูปํ อชฺฌตฺติกํ กาย\-สมฺผสฺสสฺส วตฺถุ. กายายตนํ. อิทํ ตํ รูปํ อชฺฌตฺติกํ กาย\-สมฺผสฺสสฺส วตฺถุ.}

\proseitem{791}{\csromanfolio{184}กตมํ ตํ รูปํ อชฺฌตฺติกํ กาย\-สมฺผสฺสสฺส น วตฺถุ. จกฺขายตนํ ฯเปฯ ชิวฺหายตนํ. อิทํ ตํ รูปํ อชฺฌตฺติกํ กาย\-สมฺผสฺสสฺส น วตฺถุ.}

\proseitem{792}{กตมํ ตํ รูปํ พาหิรํ กาย\-สมฺผสฺสชาย เวทนาย ฯเปฯ สญฺญาย ฯเปฯ เจตนาย ฯเปฯ กายวิญฺญาณสฺส น วตฺถุ. รูปายตนํ ฯเปฯ กพฬีกาโร อาหาโร. อิทํ ตํ รูปํ พาหิรํ กายวิญฺญาณสฺส น วตฺถุ.}

\proseitem{793}{กตมํ ตํ รูปํ อชฺฌตฺติกํ กายวิญฺญาณสฺส วตฺถุ. กายายตนํ. อิทํ ตํ รูปํ อชฺฌตฺติกํ กายวิญฺญาณสฺส วตฺถุ.}

\proseitem{794}{กตมํ ตํ รูปํ อชฺฌตฺติกํ กายวิญฺญาณสฺส น วตฺถุ. จกฺขายตนํ ฯเปฯ ชิวฺหายตนํ. อิทํ ตํ รูปํ อชฺฌตฺติกํ กายวิญฺญาณสฺส น วตฺถุ.}

\proseitem{795}{กตมํ ตํ รูปํ อชฺฌตฺติกํ จกฺขุ\-สมฺผสฺสสฺส น อารมฺมณํ จกฺขายตนํ ฯเปฯ กายายตน. อิทํ ตํ รูปํ อชฺฌตฺติกํ จกฺขุ\-สมฺผสฺสสฺส น อารมฺมณํ.}

\proseitem{796}{กตมํ ตํ รูปํ พาหิรํ จกฺขุ\-สมฺผสฺสสฺส อารมฺมณํ. รูปายตนํ. อิทํ ตํ รูปํ พาหิรํ จกฺขุ\-สมฺผสฺสสฺส อารมฺมณํ.}

\proseitem{797}{กตมํ ตํ รูปํ พาหิรํ จกฺขุ\-สมฺผสฺสสฺส น อารมฺมณํ. สทฺทายตนํ ฯเปฯ กพฬีโร อาหาโร. อิทํ ตํ รูปํ พาหิรํ จกฺขุ\-สมฺผสฺสสฺส น อารมฺมณํ.}

\proseitem{798}{กตมํ ตํ รูปํ อชฺฌตฺติกํ จกฺขุ\-สมฺผสฺสชาย เวทนายา ฯเปฯ สญฺญาย ฯเปฯ เจตนาย ฯเปฯ จกฺขุ\-วิญฺญาณสฺส น อารมฺมณํ. จกฺขายตนํ ฯเปฯ กายายตนํ. อิทํ ตํ รูปํ อชฺฌตฺติกํ จกฺขุ\-วิญฺญาณสฺส น อารมฺมณํ.}

\proseitem{799}{กตมํ ตํ รูปํ พาหิรํ จกฺขุ\-วิญฺญาณสฺส อารมฺมณํ. รูปายตนํ. อิทํ ตํ รูปํ พาหิรํ จกฺขุ\-วิญฺญาณสฺส อารมฺมณํ.}

\proseitem{800}{กตมํ ตํ รูปํ พาหิรํ จกฺขุ\-วิญฺญาณสฺส น อารมฺมณํ. สทฺทายตนํ ฯเปฯ กพฬีกาโร อาหาโร. อิทํ ตํ รูปํ พาหิรํ จกฺขุ\-วิญฺญาณสฺส น อารมฺมณํ.}

\proseitem{801}{กตมํ ตํ รูปํ อชฺฌตฺติกํ โสต\-สมฺผสฺสสฺส ฯเปฯ ฆาน\-สมฺผสฺสสฺส ฯเปฯ ชิวฺหา\-สมฺผสฺสสฺส ฯเปฯ กาย\-สมฺผสฺสสฺส น อารมฺมณํ. \csromanfolio{185}จกฺขายตนํ ฯเปฯ กายายตนํ. อิทํ ตํ รูปํ อชฺฌตฺติกํ กาย\-สมฺผสฺสสฺส น อารมฺมณํ.}

\proseitem{802}{กตมํ ตํ รูปํ พาหิรํ กาย\-สมฺผสฺสสฺส อารมฺมณํ. โผฏฺฐพฺพา\-ยตนํ. อิทํ ตํ รูปํ พาหิรํ กาย\-สมฺผสฺสสฺส อารมฺมณํ.}

\proseitem{803}{กตมํ ตํ รูปํ พาหิรํ กาย\-สมฺผสฺสสฺส น อารมฺมณํ. รูปายตนํ ฯเปฯ กพฬีกาโร อาหาโร. อิทํ ตํ รูปํ พาหิรํ กาย\-สมฺผสฺสสฺส น อารมฺมณํ.}

\proseitem{804}{กตมํ ตํ รูปํ อชฺฌตฺติกํ กาย\-สมฺผสฺสชาย เวทนาย ฯเปฯ สญฺญาย ฯเปฯ เจตนาย ฯเปฯ กายวิญฺญาณสฺส น อารมฺมณํ. จกฺขายตนํ ฯเปฯ กายายตนํ. อิทํ ตํ รูปํ อชฺฌตฺติกํ กายวิญฺญาณสฺส น อารมฺมณํ.}

\proseitem{805}{กตมํ ตํ รูปํ พาหิรํ กายวิญฺญาณสฺส อารมฺมณํ. โผฏฺฐพฺพา\-ยตนํ. อิทํ ตํ รูปํ พาหิรํ กายวิญฺญาณสฺส อารมฺมณํ.}

\proseitem{806}{กตมํ ตํ รูปํ พาหิรํ กายวิญฺญาณสฺส น อารมฺมณํ. รูปายตนํ ฯเปฯ กพฬีกาโร อาหาโร. อิทํ ตํ รูปํ พาหิรํ กายวิญฺญาณสฺส น อารมฺมณํ.}

\proseitem{807}{กตมํ ตํ รูปํ พาหิรํ น จกฺขายตนํ. รูปายตนํ ฯเปฯ กพฬีกาโร อาหาโร. อิทํ ตํ รูปํ พาหิรํ น จกฺขายตนํ.}

\proseitem{808}{กตมํ ตํ รูปํ อชฺฌตฺติกํ จกฺขายตนํ. ยํ จกฺขุ จตุนฺนํ มหาภูตานํ อุปาทาย ปสาโท ฯเปฯ สุญฺโญ คาโมเปโส. อิทํ ตํ รูปํ อชฺฌตฺติกํ จกฺขายตนํ.}

\proseitem{809}{กตมํ ตํ รูปํ อชฺฌตฺติกํ น จกฺขายตนํ. โสตายตนํ ฯเปฯ กายายตนํ. อิทํ ตํ รูปํ อชฺฌตฺติกํ น จกฺขายตนํ.}

\proseitem{810}{กตมํ ตํ รูปํ พาหิรํ น โสตายตนํ ฯเปฯ น ฆานายตนํ ฯเปฯ น ชิวฺหายตนํ ฯเปฯ น กายายตนํ. รูปายตนํ ฯเปฯ กพฬีกาโร อาหาโร. อิทํ ตํ รูปํ พาหิรํ น กายายตนํ.}

\proseitem{811}{\csromanfolio{186}กตมํ ตํ รูปํ อชฺฌตฺติกํ กายายตนํ. โย กาโย จตุนฺนํ มหาภูตานํ อุปาทาย ปสาโท ฯเปฯ สุญฺโญ คาโมเปโส. อิทํ ตํ รูปํ อชฺฌตฺติกํ กายายตนํ.}

\proseitem{812}{กตมํ ตํ รูปํ อชฺฌตฺติกํ น กายายตนํ. จกฺขายตนํ ฯเปฯ ชิวฺหายตนํ. อิทํ ตํ รูปํ อชฺฌตฺติกํ น กายายตนํ.}

\proseitem{813}{กตมํ ตํ รูปํ อชฺฌตฺติกํ น รูปายตนํ. จกฺขายตนํ ฯเปฯ กายายตนํ. อิทํ ตํ รูปํ อชฺฌตฺติกํ น รูปายตนํ.}

\proseitem{814}{กตมํ ตํ รูปํ พาหิรํ รูปายตนํ. ยํ รูปํ จตุนฺนํ มหาภูตานํ อุปาทาย วณฺณนิภา ฯเปฯ รูปธาตุเปสา. อิทํ ตํ รูปํ พาหิรํ รูปายตนํ.}

\proseitem{815}{กตมํ ตํ รูปํ พาหิรํ น รูปายตนํ. สทฺทายตนํ ฯเปฯ กพฬีกาโร อาหาโร. อิทํ ตํ รูปํ พาหิรํ น รูปายตนํ.}

\proseitem{816}{กตมํ ตํ รูปํ อชฺฌตฺติกํ น สทฺทายตนํ ฯเปฯ น คนฺธายตนํ ฯเปฯ น รสายตนํ ฯเปฯ น โผฏฺฐพฺพา\-ยตนํ. จกฺขายตนํ ฯเปฯ กายายตนํ. อิทํ ตํ รูปํ อชฺฌตฺติกํ น โผฏฺฐพฺพา\-ยตนํ.}

\proseitem{817}{กตมํ ตํ รูปํ พาหิรํ โผฏฺฐพฺพา\-ยตนํ. ปถวีธาตุ ฯเปฯ โผฏฺฐพฺพ\-ธาตุเปสา. อิทํ ตํ รูปํ พาหิรํ โผฏฺฐพฺพา\-ยตนํ.}

\proseitem{818}{กตมํ ตํ รูปํ พาหิรํ น โผฏฺฐพฺพา\-ยตนํ. รูปายตนํ ฯเปฯ กพฬีกาโร อาหาโร. อิทํ ตํ รูปํ พาหิรํ น โผฏฺฐพฺพา\-ยตนํ.}

\proseitem{819}{กตมํ ตํ รูปํ พาหิรํ น จกฺขุธาตุ. รูปายตนํ ฯเปฯ กพฬีกาโร อาหาโร. อิทํ ตํ รูปํ พาหิรํ น จกฺขุธาตุ.}

\proseitem{820}{กตมํ ตํ รูปํ อชฺฌตฺติกํ จกฺขุธาตุ. จกฺขายตนํ. อิทํ ตํ รูปํ อชฺฌตฺติกํ จกฺขุธาตุ.}

\proseitem{821}{กตมํ ตํ รูปํ อชฺฌตฺติกํ น จกฺขุธาตุ. โสตายตนํ ฯเปฯ กายายตนํ. อิทํ ตํ รูปํ อชฺฌตฺติกํ น จกฺขุธาตุ.}

\proseitem{822}{กตมํ ตํ รูปํ พาหิรํ น โสตาธาตุ ฯเปฯ น ฆานธาตุ ฯเปฯ น ชิวฺหาธาตุ ฯเปฯ น กายธาตุ. รูปายตนํ ฯเปฯ กพฬีกาโร อาหาโร. อิทํ ตํ รูปํ พาหิรํ น กายธาตุ.}

\chapterheadpageread{2. รูปกณฺฑ}
\proseitem{823}{\csromanfolio{187}กตมํ ตํ รูปํ อชฺฌตฺติกํ กายธาตุ. กายายตนํ. อิทํ ตํ รูปํ อชฺฌตฺติกํ กายธาตุ.}

\proseitem{824}{กตมํ ตํ รูปํ อชฺฌตฺติกํ น กายธาตุ. จกฺขายตนํ ฯเปฯ ชิวฺหายตนํ. อิทํ ตํ รูปํ อชฺฌตฺติกํ น กายธาตุ.}

\proseitem{825}{กตมํ ตํ รูปํ อชฺฌตฺติกํ น รูปธาตุ. จกฺขายตนํ ฯเปฯ กายายตนํ. อิทํ ตํ รูปํ อชฺฌตฺติกํ น รูปธาตุ.}

\proseitem{826}{กตมํ ตํ รูปํ พาหิรํ รูปธาตุ. รูปายตนํ. อิทํ ตํ รูปํ พาหิรํ รูปธาตุ.}

\proseitem{827}{กตมํ ตํ รูปํ พาหิรํ น รูปธาตุ. สทฺทายตนํ ฯเปฯ กพฬีกาโร อาหาโร. อิทํ ตํ รูปํ พาหิรํ น รูปธาตุ.}

\proseitem{828}{กตมํ ตํ รูปํ อชฺฌตฺติกํ น สทฺทธาตุ ฯเปฯ น คนฺธธาตุ ฯเปฯ น รสธาตุ. น โผฏฺฐพฺพ\-ธาตุ. จกฺขายตนํ ฯเปฯ กายายตนํ. อิทํ ตํ รูปํ อชฺฌตฺติกํ น โผฏฺฐพฺพาธาตุ.}

\proseitem{829}{กตมํ ตํ รูปํ พาหิรํ โผฏฺฐพฺพ\-ธาตุ. โผฏฺฐพฺพา\-ยตนํ. อิทํ ตํ รูปํ พาหิรํ โผฏฺฐพฺพ\-ธาตุ.}

\proseitem{830}{กตมํ ตํ รูปํ พาหิรํ น โผฏฺฐพฺพ\-ธาตุ ฯเปฯ รูปายตนํ ฯเปฯ กพฬีกาโร อาหาโร. อิทํ ตํ รูปํ พาหิรํ น โผฏฺฐพฺพ\-ธาตุ.}

\proseitem{831}{กตมํ ตํ รูปํ พาหิรํ น จกฺขุนฺทฺริยํ. รูปายตนํ ฯเปฯ กพฬีกาโร อาหาโร. อิทํ ตํ รูปํ พาหิรํ น จกฺขุนฺทฺริยํ.}

\proseitem{832}{กตมํ ตํ รูปํ อชฺฌตฺติกํ จกฺขุนฺทฺริยํ. ยํ จกฺขุ จตุนฺนํ มหาภูตานํ อุปาทาย ปสาโท ฯเปฯ สุญฺโญ คาโมเปโส. อิทํ ตํ รูปํ อชฺฌตฺติกํ จกฺขุนฺทฺริยํ.}

\proseitem{833}{กตมํ ตํ รูปํ อชฺฌตฺติกํ น จกฺขุนฺทฺริยํ. โสตายตนํ ฯเปฯ กายายตนํ. อิทํ ตํ รูปํ อชฺฌตฺติกํ น จกฺขุนฺทฺริยํ.}

\proseitem{834}{กตมํ ตํ รูปํ พาหิรํ น โสตินฺทฺริยํ ฯเปฯ น ฆานินฺทฺริยํ ฯเปฯ น ชิวฺหินฺทฺริยํ ฯเปฯ น กายินฺทฺริยํ. รูปายตนํ ฯเปฯ กพฬีกาโร อาหาโร. อิทํ ตํ รูปํ พาหิรํ น กายินฺทฺริยํ.}

\proseitem{835}{\csromanfolio{188}กตมํ ตํ รูปํ อชฺฌตฺติกํ กายินฺทฺริยํ. โย กาโย จตุนฺนํ มหาภูตานํ อุปาทาย ปสาโท ฯเปฯ สุญฺโญ คาโมเปโส. อิทํ ตํ รูปํ อชฺฌตฺติกํ กายินฺทฺริยํ.}

\proseitem{836}{กตมํ ตํ รูปํ อชฺฌตฺติกํ น กายินฺทฺริยํ. จกฺขายตนํ ฯเปฯ ชิวฺหายตนํ. อิทํ ตํ รูปํ อชฺฌตฺติกํ น กายินฺทฺริยํ.}

\proseitem{837}{กตมํ ตํ รูปํ อชฺฌตฺติกํ น อิตฺถินฺทฺริยํ. จกฺขายตนํ ฯเปฯ กายายตนํ. อิทํ ตํ รูปํ อชฺฌตฺติกํ น อิตฺถินฺทฺริยํ.}

\proseitem{838}{กตมํ ตํ รูปํ พาหิรํ อิตฺถินฺทฺริยํ. ยํ อิตฺถิยา อิตฺถิลิงฺคํ อิตฺถินิมิตฺตํ อิตฺถิกุตฺตํ อิตฺถากปฺโป อิตฺถตฺตํ อิตฺถิภาโว. อิทํ ตํ รูปํ พาหิรํ อิตฺถินฺทฺริยํ.}

\proseitem{839}{กตมํ ตํ รูปํ พาหิรํ น อิตฺถินฺทฺริยํ. รูปายตนํ ฯเปฯ กพฬีกาโร อาหาโร. อิทํ ตํ รูปํ พาหิรํ น อิตฺถินฺทฺริยํ.}

\proseitem{840}{กตมํ ตํ รูปํ อชฺฌตฺติกํ น ปุริสินฺทฺริยํ. จกฺขายตนํ ฯเปฯ กายายตนํ. อิทํ ตํ รูปํ อชฺฌตฺติกํ น ปุริสินฺทฺริยํ.}

\proseitem{841}{กตมํ ตํ รูปํ พาหิรํ ปุริสินฺทฺริยํ. ยํ ปุริสสฺส ปุริสลิงฺคํ ปุริสนิมิตฺตํ ปุริสกุตฺตํ ปุริสากปฺโป ปุริสตฺตํ ปุริสภาโว. อิทํ ตํ รูปํ พาหิรํ ปุริสินฺทฺริยํ.}

\proseitem{842}{กตมํ ตํ รูปํ พาหิรํ น ปุริสินฺทฺริยํ. รูปายตนํ ฯเปฯ กพฬีกาโร อาหาโร. อิทํ ตํ รูปํ พาหิรํ น ปุริสินฺทฺริยํ.}

\proseitem{843}{กตมํ ตํ รูปํ อชฺฌตฺติกํ น ชีวิตินฺทฺริยํ. จกฺขายตนํ ฯเปฯ กายายตนํ. อิทํ ตํ รูปํ อชฺฌตฺติกํ น ชีวิตินฺทฺริยํ.}

\proseitem{844}{กตมํ ตํ รูปํ พาหิรํ ชีวิตินฺทฺริยํ. โย เตสํ รูปีนํ ธมฺมานํ อายุ ฐิติ ยปนา ยาปนา อิริยนา วตฺตนา ปาลนา ชีวิตํ ชีวิตินฺทฺริยํ. อิทํ ตํ รูปํ พาหิรํ ชีวิตินฺทฺริยํ.}

\proseitem{845}{กตมํ ตํ รูปํ พาหิรํ น ชีวิตินฺทฺริยํ. รูปายตนํ ฯเปฯ กพฬีกาโร อาหาโร. อิทํ ตํ รูปํ พาหิรํ น ชีวิตินฺทฺริยํ.}

\proseitem{846}{กตมํ ตํ รูปํ อชฺฌตฺติกํ น กายวิญฺญตฺติ. จกฺขายตนํ ฯเปฯ กายายตนํ. อิทํ ตํ รูปํ อชฺฌตฺติกํ น กายวิญฺญตฺติ.}

\chapterheadpageread{2. รูปกณฺฑ}
\proseitem{847}{\csromanfolio{189}กตมํ ตํ รูปํ พาหิรํ กายวิญฺญตฺติ. ยา กุสลจิตฺตสฺส วา อกุสลจิตฺตสฺส วา อพฺยากต\-จิตฺตสฺส วา อภิกฺกม\-นฺตสฺส วา ปฏิกฺกมนฺตสฺส วา อาโลเกนฺตสฺส วา วิโลเกนฺตสฺส วา สมิญฺเชนฺตสฺส วา ปสาเรนฺตสฺส วา กายสฺส ถมฺภนา สนฺถมฺภนา สนฺถมฺภิตตฺตํ วิญฺญตฺติ วิญฺญาปนา วิญฺญาปิตตฺตํ. อิทํ ตํ รูปํ พาหิรํ กายวิญฺญตฺติ.}

\proseitem{848}{กตมํ ตํ รูปํ พาหิรํ น กายวิญฺญตฺติ. รูปายตนํ ฯเปฯ กพฬีกาโร อาหาโร. อิทํ ตํ รูปํ พาหิรํ น กายวิญฺญตฺติ.}

\proseitem{849}{กตมํ ตํ รูปํ อชฺฌตฺติกํ น วจีวิญฺญตฺติ. จกฺขายตนํ ฯเปฯ กายายตนํ. อิทํ ตํ รูปํ อชฺฌตฺติกํ น วจีวิญฺญตฺติ.}

\proseitem{850}{กตมํ ตํ รูปํ พาหิรํ วจีวิญฺญตฺติ. ยา กุสลจิตฺตสฺส วา อกุสลจิตฺตสฺส วา อพฺยากต\-จิตฺตสฺส วา วาจา คิรา พฺยปฺปโถ อุทีรณํ โฆโส โฆสกมฺมํ วาจา วจีเภโท, อยํ วุจฺจติ วาจา. ยา ตาย วาจาย วิญฺญตฺติ วิญฺญาปนา วิญฺญาปิตตฺตํ. อิทํ ตํ รูปํ พาหิรํ วจีวิญฺญตฺติ.}

\proseitem{851}{กตมํ ตํ รูปํ พาหิรํ น วจีวิญฺญตฺติ. รูปายตนํ ฯเปฯ กพฬีกาโร อาหาโร. อิทํ ตํ รูปํ พาหิรํ น วจีวิญฺญตฺติ.}

\proseitem{852}{กตมํ ตํ รูปํ อชฺฌตฺติกํ น อากาสธาตุ. จกฺขายตนํ ฯเปฯ กายายตนํ. อิทํ ตํ รูปํ อชฺฌตฺติกํ น อากาสธาตุ.}

\proseitem{853}{กตมํ ตํ รูปํ พาหิรํ อากาสธาตุ. โย อากาโส อากาสคตํ อฆํ อฆคตํ วิวโร วิวรคตํ อสมฺผุฏฺฐํ จตูหิ มหาภูเตหิ. อิทํ ตํ รูปํ พาหิรํ อากาสธาตุ.}

\proseitem{854}{กตมํ ตํ รูปํ พาหิรํ น อากาสธาตุ. รูปายตนํ ฯเปฯ กพฬีกาโร อาหาโร. อิทํ ตํ รูปํ พาหิรํ น อากาสธาตุ.}

\proseitem{855}{กตมํ ตํ รูปํ อชฺฌตฺติกํ น อาโปธาตุ. จกฺขายตนํ ฯเปฯ กายายตนํ. อิทํ ตํ รูปํ อชฺฌตฺติกํ น อาโปธาตุ.}

\proseitem{856}{กตมํ ตํ รูปํ พาหิรํ อาโปธาตุ. ยํ อาโป อาโปคตํ สิเนโห สิเนหคตํ พนฺธนตฺตํ รูปสฺส. อิทํ ตํ รูปํ พาหิรํ อาโปธาตุ.}

\proseitem{857}{กตมํ ตํ รูปํ พาหิรํ น อาโปธาตุ. รูปายตนํ ฯเปฯ กพฬีกาโร อาหาโร. อิทํ ตํ รูปํ พาหิรํ น อาโปธาตุ.}

\proseitem{858}{\csromanfolio{190}กตมํ ตํ รูปํ อชฺฌตฺติกํ รูปสฺส น ลหุตา. จกฺขายตนํ ฯเปฯ กายายตนํ. อิทํ ตํ รูปํ อชฺฌตฺติกํ รูปสฺส น ลหุตา.}

\proseitem{859}{กตมํ ตํ รูปํ พาหิรํ รูปสฺส ลหุตา. ยา รูปสฺส ลหุตา ลหุปริณามตา อทนฺธนตา อวิตฺถนตา. อิทํ ตํ รูปํ พาหิรํ รูปสฺส ลหุตา.}

\proseitem{860}{กตมํ ตํ รูปํ พาหิรํ รูปสฺส น ลหุตา. รูปายตนํ ฯเปฯ กพฬีกาโร อาหาโร. อิทํ ตํ รูปํ พาหิรํ รูปสฺส น ลหุตา.}

\proseitem{861}{กตมํ ตํ รูปํ อชฺฌตฺติกํ รูปสฺส น มุทุตา. จกฺขายตนํ ฯเปฯ กายายตนํ. อิทํ ตํ รูปํ อชฺฌตฺติกํ รูปสฺส น มุทุตา.}

\proseitem{862}{กตมํ ตํ รูปํ พาหิรํ รูปสฺส มุทุตา. ยา รูปสฺส มุทุตา มทฺทวตา อกกฺขฬตา อกถินตา. อิทํ ตํ รูปํ พาหิรํ รูปสฺส มุทุตา.}

\proseitem{863}{กตมํ ตํ รูปํ พาหิรํ รูปสฺส น มุทุตา. รูปายตนํ ฯเปฯ กพฬีกาโร อาหาโร. อิทํ ตํ รูปํ พาหิรํ รูปสฺส น มุทุตา.}

\proseitem{864}{กตมํ ตํ รูปํ อชฺฌตฺติกํ รูปสฺส น กมฺมญฺญตา. จกฺขายตนํ ฯเปฯ กายายตนํ. อิทํ ตํ รูปํ อชฺฌตฺติกํ รูปสฺส น กมฺมญฺญตา.}

\proseitem{865}{กตมํ ตํ รูปํ พาหิรํ รูปสฺส กมฺมญฺญตา. ยา รูปสฺส กมฺมญฺญตา กมฺมญฺญตฺตํ กมฺมญฺญภาโว. อิทํ ตํ รูปํ พาหิรํ รูปสฺส กมฺมญฺญตา.}

\proseitem{866}{กตมํ ตํ รูปํ พาหิรํ รูปสฺส น กมฺมญฺญตา. รูปายตนํ ฯเปฯ กพฬีกาโร อาหาโร. อิทํ ตํ รูปํ พาหิรํ รูปสฺส น กมฺมญฺญตา.}

\proseitem{867}{กตมํ ตํ รูปํ อชฺฌตฺติกํ รูปสฺส น อุปจโย. จกฺขายตนํ ฯเปฯ กายายตนํ. อิทํ ตํ รูปํ อชฺฌตฺติกํ รูปสฺส น อุปจโย.}

\proseitem{868}{กตมํ ตํ รูปํ พาหิรํ รูปสฺส อุปจโย. โย อายตนานํ อาจโย, โส รูปสฺส อุปจโย. อิทํ ตํ รูปํ พาหิรํ รูปสฺส อุปจโย.}

\proseitem{869}{กตมํ ตํ รูปํ พาหิรํ รูปสฺส น อุปจโย. รูปายตนํ ฯเปฯ กพฬีกาโร อาหาโร. อิทํ ตํ รูปํ พาหิรํ รูปสฺส น อุปจโย.}

\proseitem{870}{กตมํ ตํ รูปํ อชฺฌตฺติกํ รูปสฺส น สนฺตติ. จกฺขายตนํ ฯเปฯ กายายตนํ. อิทํ ตํ รูปํ อชฺฌตฺติกํ รูปสฺส น สนฺตติ.}

\chapterheadpageread{2. รูปกณฺฑ}
\proseitem{871}{\csromanfolio{191}กตมํ ตํ รูปํ พาหิรํ รูปสฺส สนฺตติ. โย รูปสฺส อุปจโย, สา รูปสฺส สนฺตติ. อิทํ ตํ รูปํ พาหิรํ รูปสฺส สนฺตติ.}

\proseitem{872}{กตมํ ตํ รูปํ พาหิรํ รูปสฺส น สนฺตติ. รูปายตนํ ฯเปฯ กพฬีกาโร อาหาโร. อิทํ ตํ รูปํ พาหิรํ รูปสฺส น สนฺตติ.}

\proseitem{873}{กตมํ ตํ รูปํ อชฺฌตฺติกํ รูปสฺส น ชรตา. จกฺขายตนํ ฯเปฯ กายายตนํ. อิทํ ตํ รูปํ อชฺฌตฺติกํ รูปสฺส น ชรตา.}

\proseitem{874}{กตมํ ตํ รูปํ พาหิรํ รูปสฺส ชรตา. ยา รูปสฺส ชรา ชีรณตา ขณฺฑิจฺจํ ปาลิจฺจํ วลิตฺตจตา อายุโน สํหานิ อินฺทฺริยานํ ปริปาโก. อิทํ ตํ รูปํ พาหิรํ รูปสฺส ชรตา.}

\proseitem{875}{กตมํ ตํ รูปํ พาหิรํ รูปสฺส น ชรตา. รูปายตนํ ฯเปฯ กพฬีกาโร อาหาโร. อิทํ ตํ รูปํ พาหิรํ รูปสฺส น ชรตา.}

\proseitem{876}{กตมํ ตํ รูปํ อชฺฌตฺติกํ รูปสฺส น อนิจฺจตา. จกฺขายตนํ ฯเปฯ กายายตนํ. อิทํ ตํ รูปํ อชฺฌตฺติกํ รูปสฺส น อนิจฺจตา.}

\proseitem{877}{กตมํ ตํ รูปํ พาหิรํ รูปสฺส อนิจฺจตา. โย รูปสฺส ขโย วโย เภโท ปริเภโท อนิจฺจตา อนฺตรธานํ. อิทํ ตํ รูปํ พาหิรํ รูปสฺส อนิจฺจตา.}

\proseitem{878}{กตมํ ตํ รูปํ พาหิรํ รูปสฺส น อนิจฺจตา. รูปายตนํ ฯเปฯ กพฬีกาโร อาหาโร. อิทํ ตํ รูปํ พาหิรํ รูปสฺส น อนิจฺจตา.}

\proseitem{879}{กตมํ ตํ รูปํ อชฺฌตฺติกํ น กพฬีกาโร อาหาโร. จกฺขายตนํ ฯเปฯ กายายตนํ. อิทํ ตํ รูปํ อชฺฌตฺติกํ น กพฬีกาโร อาหาโร.}

\proseitem{880}{กตมํ ตํ รูปํ พาหิรํ กพฬีกาโร อาหาโร. โอทโน กุมฺมาโส สตฺตุ มจฺโฉ มํสํ ขีรํ ทธิ สปฺปิ นวนีตํ เตลํ มธุ ผาณิตํ, ยํ วา ปนญฺญมฺปิ อตฺถิ รูปํ ยมฺหิ ยมฺหิ ชนปเท เตสํ เตสํ สตฺตานํ มุขาสิยํ ทนฺตวิขาทนํ คลชฺโฌหรณียํ กุจฺฉิ\-วิตฺถมฺภนํ ยาย โอชาย สตฺตา ยาเปนฺติ. อิทํ ตํ รูปํ พาหิรํ กพฬีกาโร อาหาโร.}

\proseitem{881}{\csromanfolio{192}กตมํ ตํ รูปํ พาหิรํ น กพฬีกาโร อาหาโร. รูปายตนํ ฯเปฯ รูปสฺส อนิจฺจตา. อิทํ ตํ รูปํ พาหิรํ น กพฬีกาโร อาหาโร. เอวํ ติวิเธน รูปสงฺคโห. ติกนิทฺเทโส.}

\csromanmatikaanchor{mtk.94}
\csromantocmarkref{subsection}{จตุกฺก}{mtk.94}
\bookmark[dest={mtk.94},level=2]{จตุกฺก}
\csromanheader{2}{\textbf{จตุกฺก}}
\proseitem{882}{กตมํ ตํ รูปํ อุปาทา อุปาทิณฺณํ. จกฺขายตนํ ฯเปฯ กายายตนํ อิตฺถินฺทฺริยํ ปุริสินฺทฺริยํ ชีวิตินฺทฺริยํ, ยํ วา ปนญฺญมฺปิ อตฺถิ รูปํ กมฺมสฺส กตตฺตา รูปายตนํ คนฺธายตนํ รสายตนํ อากาสธาตุ รูปสฺส อุปจโย รูปสฺส สนฺตติ กพฬีกาโร อาหาโร. อิทํ ตํ รูปํ อุปาทา อุปาทิณฺณํ.}

\proseitem{883}{กตมํ ตํ รูปํ อุปาทา อนุปาทิณฺณํ. สทฺทายตนํ กายวิญฺญตฺติ วจีวิญฺญตฺติ รูปสฺส ลหุตา รูปสฺส มุทุตา รูปสฺส กมฺมญฺญตา รูปสฺส ชรตา รูปสฺส อนิจฺจตา, ยํ วา ปนญฺญมฺปิ อตฺถิ รูปํ น กมฺมสฺส กตตฺตา รูปายตนํ คนฺธายตนํ รสายตนํ อากาสธาตุ รูปสฺส อุปจโย รูปสฺส สนฺตติ กพฬีกาโร อาหาโร. อิทํ ตํ รูปํ อุปาทา อนุปาทิณฺณํ.}

\proseitem{884}{กตมํ ตํ รูปํ โน อุปาทา อุปาทิณฺณํ. กมฺมสฺส กตตฺตา โผฏฺฐพฺพา\-ยตนํ อาโปธาตุ. อิทํ ตํ รูปํ โน อุปาทา อุปาทิณฺณํ.}

\proseitem{885}{กตมํ ตํ รูปํ โน อุปาทา อนุปาทิณฺณํ. น กมฺมสฺส กตตฺตา โผฏฺฐพฺพา\-ยตนํ อาโปธาตุ. อิทํ ตํ รูปํ โน อุปาทา อนุปาทิณฺณํ.}

\proseitem{886}{กตมํ ตํ รูปํ อุปาทา อุปาทิณฺณุ\-ปาทานิยํ. จกฺขายตนํ ฯเปฯ กายายตนํ อิตฺถินฺทฺริยํ ปุริสินฺทฺริยํ ชีวิตินฺทฺริยํ, ยํ วา ปนญฺญมฺปิ อตฺถิ รูปํ กมฺมสฺส กตตฺตา รูปายตนํ คนฺธายตนํ รสายตนํ อากาสธาตุ รูปสฺส อุปจโย รูปสฺส สนฺตติ กพฬีกาโร อาหาโร. อิทํ ตํ รูปํ อุปาทา อุปาทิณฺณุ\-ปาทานิยํ.}

\proseitem{887}{กตมํ ตํ รูปํ อุปาทา อนุปาทิณฺณุ\-ปาทานิยํ. สทฺทายตนํ กายวิญฺญตฺติ วจีวิญฺญตฺติ รูปสฺส ลหุตา รูปสฺส มุทุตา รูปสฺส กมฺมญฺญตา \csromanfolio{193}รูปสฺส ชรตา รูปสฺส อนิจฺจตา, ยํ วา ปนญฺญมฺปิ อตฺถิ รูปํ น กมฺมสฺส กตตฺตา รูปายตนํ คนฺธายตนํ รสายตนํ อากาสธาตุ รูปสฺส อุปจโย รูปสฺส สนฺตติ กพฬีกาโร อาหาโร. อิทํ ตํ รูปํ อุปาทา อนุปาทิณฺณุ\-ปาทานิยํ.}

\proseitem{888}{กตมํ ตํ รูปํ โน อุปาทา อุปาทิณฺณุ\-ปาทานิยํ. กมฺมสฺส กตตฺตา โผฏฺฐพฺพา\-ยตนํ อาโปธาตุ. อิทํ ตํ รูปํ โน อุปาทา อุปาทิณฺณุ\-ปาทานิยํ.}

\proseitem{889}{กตมํ ตํ รูปํ โน อุปาทา อนุปาทิณฺณุ\-ปาทานิยํ. น กมฺมสฺส กตตฺตา โผฏฺฐพฺพา\-ยตนํ อาโปธาตุ. อิทํ ตํ รูปํ โน อุปาทา อนุปาทิณฺณุ\-ปาทานิยํ.}

\proseitem{890}{กตมํ ตํ รูปํ อุปาทา สปฺปฏิฆํ. จกฺขายตนํ ฯเปฯ รสายตนํ. อิทํ ตํ รูปํ อุปาทา สปฺปฏิฆํ.}

\proseitem{891}{กตมํ ตํ รูปํ อุปาทา อปฺปฏิฆํ. อิตฺถินฺทฺริยํ ฯเปฯ กพฬีกาโร อาหาโร. อิทํ ตํ รูปํ อุปาทา อปฺปฏิฆํ.}

\proseitem{892}{กตมํ ตํ รูปํ โน อุปาทา สปฺปฏิฆํ. โผฏฺฐพฺพา\-ยตนํ. อิทํ ตํ รูปํ โน อุปาทา สปฺปฏิฆํ.}

\proseitem{893}{กตมํ ตํ รูปํ โน อุปาทา อปฺปฏิฆํ. อาโปธาตุ. อิทํ ตํ รูปํ โน อุปาทา อปฺปฏิฆํ.}

\proseitem{894}{กตมํ ตํ รูปํ อุปาทา โอฬาริกํ. จกฺขายตนํ ฯเปฯ รสายตนํ. อิทํ ตํ รูปํ อุปาทา โอฬาริกํ.}

\proseitem{895}{กตมํ ตํ รูปํ อุปาทา สุขุมํ. อิตฺถินฺทฺริยํ ฯเปฯ กพฬีกาโร อาหาโร. อิทํ ตํ รูปํ อุปาทา สุขุมํ.}

\proseitem{896}{กตมํ ตํ รูปํ โน อุปาทา โอฬาริกํ. โผฏฺฐพฺพา\-ยตนํ. อิทํ ตํ รูปํ โน อุปาทา โอฬาริกํ.}

\proseitem{897}{กตมํ ตํ รูปํ โน อุปาทา สุขุมํ. อาโปธาตุ. อิทํ ตํ รูปํ โน อุปาทา สุขุมํ.}

\proseitem{898}{กตมํ ตํ รูปํ อุปาทา ทูเร. อิตฺถินฺทฺริยํ ฯเปฯ กพฬีกาโร อาหาโร. อิทํ ตํ รูปํ อุปาทา ทูเร.}

\proseitem{899}{\csromanfolio{194}กตมํ ตํ รูปํ อุปาทา สนฺติเก. จกฺขายตนํ ฯเปฯ รสายตนํ. อิทํ ตํ รูปํ อุปาทา สนฺติเก.}

\proseitem{900}{กตมํ ตํ รูปํ โน อุปาทา ทูเร. อาโปธาตุ. อิทํ ตํ รูปํ โน อุปาทา ทูเร.}

\proseitem{901}{กตมํ ตํ รูปํ โน อุปาทา สนฺติเก. โผฏฺฐพฺพา\-ยตนํ. อิทํ ตํ รูปํ โน อุปาทา สนฺติเก.}

\proseitem{902}{กตมํ ตํ รูปํ อุปาทิณฺณํ สนิทสฺสนํ. กมฺมสฺส กตตฺตา รูปายตนํ. อิทํ ตํ รูปํ อุปาทิณฺณํ สนิทสฺสนํ.}

\proseitem{903}{กตมํ ตํ รูปํ อุปาทิณฺณํ อนิทสฺสนํ. จกฺขายตนํ ฯเปฯ กายายตนํ อิตฺถินฺทฺริยํ ปุริสินฺทฺริยํ ชีวิตินฺทฺริยํ, ยํ วา ปนญฺญมฺปิ อตฺถิ รูปํ กมฺมสฺส กตตฺตา คนฺธายตนํ รสายตนํ โผฏฺฐพฺพา\-ยตนํ อากาสธาตุ อาโปธาตุ รูปสฺส อุปจโย รูปสฺส สนฺตติ กพฬีกาโร อาหาโร. อิทํ ตํ รูปํ อุปาทิณฺณํ อนิทสฺสนํ.}

\proseitem{904}{กตมํ ตํ รูปํ อนุปาทิณฺณํ สนิทสฺสนํ. น กมฺมสฺส กตตฺตา รูปายตนํ. อิทํ ตํ รูปํ อนุปาทิณฺณํ สนิทสฺสนํ.}

\proseitem{905}{กตมํ ตํ รูปํ อนุปาทิณฺณํ อนิทสฺสนํ. สทฺทายตนํ กายวิญฺญตฺติ วจีวิญฺญตฺติ รูปสฺส ลหุตา รูปสฺส มุทุตา รูปสฺส กมฺมญฺญตา รูปสฺส ชรตา รูปสฺส อนิจฺจตา, ยํ วา ปนญฺญมฺปิ อตฺถิ รูปํ น กมฺมสฺส กตตฺตา คนฺธายตนํ รสายตนํ โผฏฺฐพฺพา\-ยตนํ อากาสธาตุ อาโปธาตุ รูปสฺส อุปจโย รูปสฺส สนฺตติ กพฬีกาโร อาหาโร. อิทํ ตํ รูปํ อนุปาทิณฺณํ อนิทสฺสนํ.}

\proseitem{906}{กตมํ ตํ รูปํ อุปาทิณฺณํ สปฺปฏิฆํ. จกฺขายตนํ ฯเปฯ กายายตนํ, ยํ วา ปนญฺญมฺปิ รูปํ กมฺมสฺส กตตฺตา รูปายตนํ คนฺธายตนํ รสายตนํ โผฏฺฐพฺพา\-ยตนํ. อิทํ ตํ รูปํ อุปาทิณฺณํ สปฺปฏิฆํ.}

\proseitem{907}{กตมํ ตํ รูปํ อุปาทิณฺณํ อปฺปฏิฆํ. อิตฺถินฺทฺริยํ ปุริสินฺทฺริยํ ชีวิตินฺทฺริยํ, ยํ วา ปนญฺญมฺปิ อตฺถิ รูปํ กมฺมสฺส กตตฺตา อากาสธาตุ อาโปธาตุ รูปสฺส อุปจโย รูปสฺส สนฺตติ กพฬีกาโร อาหาโร. อิทํ ตํ รูปํ อุปาทิณฺณํ อปฺปฏิฆํ.}

\chapterheadpageread{2. รูปกณฺฑ}
\proseitem{908}{\csromanfolio{195}กตมํ ตํ รูปํ อนุปาทิณฺณํ สปฺปฏิฆํ. สทฺทายตนํ, ยํ วา ปนญฺญมฺปิ อตฺถิ รูปํ น กมฺมสฺส กตตฺตา รูปายตนํ คนฺธายตนํ รสายตนํ โผฏฺฐพฺพา\-ยตนํ. อิทํ ตํ รูปํ อนุปาทิณฺณํ สปฺปฏิฆํ.}

\proseitem{909}{กตมํ ตํ รูปํ อนุปาทิณฺณํ อปฺปฏิฆํ. กายวิญฺญตฺติ วจีวิญฺญตฺติ รูปสฺส ลหุตา รูปสฺส มุทุตา รูปสฺส กมฺมญฺญตา รูปสฺส ชรตา รูปสฺส อนิจฺจตา, ยํ วา ปนญฺญมฺปิ อตฺถิ รูปํ น กมฺมสฺส กตตฺตา อากาสธาตุ อาโปธาตุ รูปสฺส อุปจโย รูปสฺส สนฺตติ กพฬีกาโร อาหาโร. อิทํ ตํ รูปํ อนุปาทิณฺณํ อปฺปฏิฆํ.}

\proseitem{910}{กตมํ ตํ รูปํ อุปาทิณฺณํ มหาภูตํ. กมฺมสฺส กตตฺตา โผฏฺฐพฺพา\-ยตนํ อาโปธาตุ. อิทํ ตํ รูปํ อุปาทิณฺณํ มหาภูตํ.}

\proseitem{911}{กตมํ ตํ รูปํ อุปาทิณฺณํ น มหาภูตํ. จกฺขายตนํ ฯเปฯ กายายตนํ อิตฺถินฺทฺริยํ ปุริสินฺทฺริยํ ชีวิตินฺทฺริยํ, ยํ วา ปนญฺญมฺปิ อตฺถิ รูปํ กมฺมสฺส กตตฺตา รูปายตนํ คนฺธายตนํ รสายตนํ อากาสธาตุ รูปสฺส อุปจโย รูปสฺส สนฺตติ กพฬีกาโร อาหาโร. อิทํ ตํ รูปํ อุปาทิณฺณํ น มหาภูตํ.}

\proseitem{912}{กตมํ ตํ รูปํ อนุปาทิณฺณํ มหาภูตํ. น กมฺมสฺส กตตฺตา โผฏฺฐพฺพา\-ยตนํ อาโปธาตุ. อิทํ ตํ รูปํ อนุปาทิณฺณํ มหาภูตํ.}

\proseitem{913}{กตมํ ตํ รูปํ อนุปาทิณฺณํ น มหาภูตํ. สทฺทายตนํ กายวิญฺญตฺติ วจีวิญฺญตฺติ รูปาสฺส ลหุตา รูปสฺส มุทุตา รูปสฺส กมฺมญฺญตา รูปสฺส ชรตา รูปสฺส อนิจฺจตา, ยํ วา ปนญฺญมฺปิ อตฺถิ รูปํ น กมฺมสฺส กตตฺตา รูปายตนํ คนฺธายตนํ รสายตนํ อากาสธาตุ รูปสฺส อุปจโย รูปสฺส สนฺตติ กพฬีกาโร อาหาโร. อิทํ ตํ รูปํ อนุปาทิณฺณํ น มหาภูตํ.}

\proseitem{914}{กตมํ ตํ รูปํ อุปาทิณฺณํ โอฬาริกํ. จกฺขายตนํ ฯเปฯ กายายตนํ, ยํ วา ปนญฺญมฺปิ อตฺถิ รูปํ กมฺมสฺส กตตฺตา รูปายตนํ คนฺธายตนํ รสายตนํ โผฏฺฐพฺพา\-ยตนํ. อิทํ ตํ รูปํ อุปาทิณฺณํ โอฬาริกํ.}

\proseitem{915}{กตมํ ตํ รูปํ อุปาทิณฺณํ สุขุมํ. อิตฺถินฺทฺริยํ ปุริสินฺทฺริยํ ชีวิตินฺทฺริยํ, ยํ วา ปนญฺญมฺปิ อตฺถิ รูปํ กมฺมสฺส กตตฺตา อากาสธาตุ อาโปธาตุ \csromanfolio{196}รูปสฺส อุปจโย รูปสฺส สนฺตติ กพฬีกาโร อาหาโร. อิทํ ตํ รูปํ อุปาทิณฺณํ สุขุมํ.}

\proseitem{916}{กตมํ ตํ รูปํ อนุปาทิณฺณํ โอฬาริกํ. สทฺทายตนํ, ยํ วา ปนญฺญมฺปิ อตฺถิ รูปํ น กมฺมสฺส กตตฺตา รูปายตนํ คนฺธายตนํ รสายตนํ โผฏฺฐพฺพา\-ยตนํ. อิทํ ตํ รูปํ อนุปาทิณฺณํ โอฬาริกํ.}

\proseitem{917}{กตมํ ตํ รูปํ อนุปาทิณฺณํ สุขุมํ. กายวิญฺญตฺติ วจีวิญฺญตฺติ รูปสฺส ลหุตา รูปสฺส มุทุตา รูปสฺส กมฺมญฺญตา รูปสฺส ชรตา รูปสฺส อนิจฺจตา, ยํ วา ปนญฺญมฺปิ อตฺถิ รูปํ น กมฺมสฺส กตตฺตา อากาสธาตุ อาโปธาตุ รูปสฺส อุปจโย รูปสฺส สนฺตติ กพฬีกาโร อาหาโร. อิทํ ตํ รูปํ อนุปาทิณฺณํ สุขุมํ.}

\proseitem{918}{กตมํ ตํ รูปํ อุปาทิณฺณํ ทูเร. อิตฺถินฺทฺริยํ ปุริสินฺทฺริยํ ชีวิตินฺทฺริยํ, ยํ วา ปนญฺญมฺปิ อตฺถิ รูปํ กมฺมสฺส กตตฺตา อากาสธาตุ อาโปธาตุ รูปสฺส อุปจโย รูปสฺส สนฺตติ กพฬีกาโร อาหาโร. อิทํ ตํ รูปํ อุปาทิณฺณํ ทูเร.}

\proseitem{919}{กตมํ ตํ รูปํ อุปาทิณฺณํ สนฺติเก. จกฺขายตนํ ฯเปฯ กายายตนํ, ยํ วา ปนญฺญมฺปิ อตฺถิ รูปํ กมฺมสฺส กตตฺตา รูปายตนํ คนฺธายตนํ รสายตนํ โผฏฺฐพฺพา\-ยตนํ. อิทํ ตํ รูปํ อุปาทิณฺณํ สนฺติเก.}

\proseitem{920}{กตมํ ตํ รูปํ อนุปาทิณฺณํ ทูเร. กายวิญฺญตฺติ วจีวิญฺญตฺติ รูปสฺส ลหุตา รูปสฺส มุทุตา รูปสฺส กมฺมญฺญตา รูปสฺส ชรตา รูปสฺส อนิจฺจตา, ยํ วา ปนญฺญมฺปิ อตฺถิ รูปํ น กมฺมสฺส กตตฺตา อากาสธาตุ อาโปธาตุ รูปสฺส อุปจโย รูปสฺส สนฺตติ กพฬีกาโร อาหาโร. อิทํ ตํ รูปํ อนุปาทิณฺณํ ทูเร.}

\proseitem{921}{กตมํ ตํ รูปํ อนุปาทิณฺณํ สนฺติเก. สทฺทายตนํ, ยํ วา ปนญฺญมฺปิ อตฺถิ รูปํ น กมฺมสฺส กตตฺตา รูปายตนํ คนฺธายตนํ รสายตนํ โผฏฺฐพฺพา\-ยตนํ. อิทํ ตํ รูปํ อนุปาทิณฺณํ สนฺติเก.}

\proseitem{922}{กตมํ ตํ รูปํ อุปาทิณฺณุ\-ปาทานิยํ สนิทสฺสนํ. กมฺมสฺส กตตฺตา รูปายตนํ. อิทํ ตํ รูปํ อุปาทิณฺณุ\-ปาทานิยํ สนิทสฺสนํ.}

\chapterheadpageread{2. รูปกณฺฑ}
\proseitem{923}{\csromanfolio{197}กตมํ ตํ รูปํ อุปทิณฺณุปาทานิยํ อนิทสฺสนํ. จกฺขายตนํ ฯเปฯ กายายตนํ อิตฺถินฺทฺริยํ ปุริสินฺทฺริยํ ชีวิตินฺทฺริยํ, ยํ วา ปนญฺญมฺปิ อตฺถิ รูปํ กมฺมสฺส กตตฺตา คนฺธายตนํ รสายตนํ โผฏฺฐพฺพา\-ยตนํ อากาสธาตุ อาโปธาตุ รูปสฺส อุปจโย รูปสฺส สนฺตติ กพฬีกาโร อาหาโร. อิทํ ตํ รูปํ อุปาทิณฺณุ\-ปาทานิยํ อนิทสฺสนํ.}

\proseitem{924}{กตมํ ตํ รูปํ อนุปาทิณฺณุ\-ปาทานิยํ สนิทสฺสนํ. น กมฺมสฺส กตตฺตา รูปายตนํ. อิทํ ตํ รูปํ อนุปาทิณฺณุ\-ปาทานิยํ สนิทสฺสนํ.}

\proseitem{925}{กตมํ ตํ รูปํ อนุปาทิณฺณุ\-ปาทานิยํ อนิทสฺสนํ. สทฺทายตนํ กายวิญฺญตฺติ วจีวิญฺญตฺติ รูปสฺส ลหุตา รูปสฺส มุทุตา รูปสฺส กมฺมญฺญตา รูปสฺส ชรตา รูปสฺส อนิจฺจตา, ยํ วา ปนญฺญมฺปิ อตฺถิ รูปํ น กมฺมสฺส กตตฺตา คนฺธายตนํ รสายตนํ โผฏฺฐพฺพา\-ยตนํ อากาสธาตุ อาโปธาตุ รูปสฺส อุปจโย รูปสฺส สนฺตติ กาพฬีกาโร อาหาโร. อิทํ ตํ รูปํ อนุปาทิณฺณุ\-ปาทานิยํ อนิทสฺสนํ.}

\proseitem{926}{กตมํ ตํ รูปํ อุปาทิณฺณุ\-ปาทานิยํ สปฺปฏิฆํ. จกฺขายตนํ ฯเปฯ กายายตนํ, ยํ วา ปนญฺญมฺปิ อตฺถิ รูปํ กมฺมสฺส กตตฺตา รูปายตนํ คนฺธายตนํ รสายตนํ โผฏฺฐพฺพา\-ยตนํ. อิทํ ตํ รูปํ อุปาทิณฺณุ\-ปาทานิยํ สปฺปฏิฆํ.}

\proseitem{927}{กตมํ ตํ รูปํ อุปาทิณฺณุ\-ปาทานิยํ อปฺปฏิฆํ. อิตฺถินฺทฺริยํ ปุริสินฺทฺริยํ ชีวิตินฺทฺริยํ, ยํ วา ปนญฺญมฺปิ อตฺถิ รูปํ กมฺมสฺส กตตฺตา อากาสธาตุ อาโปธาตุ รูปสฺส อุปจโย รูปสฺส สนฺตติ กพฬีกาโร อาหาโร. อิทํ ตํ รูปํ อุปาทิณฺณุ\-ปาทานิยํ อปฺปฏิฆํ.}

\proseitem{928}{กตมํ ตํ รูปํ อนุปาทิณฺณุ\-ปาทานิยํ สปฺปฏิฆํ. สทฺทายตนํ, ยํ วา ปนญฺญมฺปิ อตฺถิ รูปํ น กมฺมสฺส กตตฺตา รูปายตนํ คนฺธายตนํ รสายตนํ โผฏฺฐพฺพา\-ยตนํ. อิทํ ตํ รูปํ อนุปาทิณฺณุ\-ปาทานิยํ สปฺปฏิฆํ.}

\proseitem{929}{กตมํ ตํ รูปํ อนุปาทิณฺณุ\-ปาทานิยํ อปฺปฏิฆํ. กายวิญฺญตฺติ วจีวิญฺญตฺติ รูปสฺส ลหุตา รูปสฺส มุทุตา รูปสฺส กมฺมญฺญตา รูปสฺส ชรตา รูปสฺส อนิจฺจตา, ยํ วา ปนญฺญมฺปิ อตฺถิ รูปํ น กมฺมสฺส กตตฺตา อากาสธาตุ อาโปธาตุ รูปสฺส อุปจโย รูปสฺส สนฺตติ กพฬีกาโร อาหาโร. อิทํ ตํ รูปํ อนุปาทิณฺณุ\-ปาทานิยํ อปฺปฏิฆํ.}

\proseitem{930}{\csromanfolio{198}กตมํ ตํ รูปํ อุปาทิณฺณุ\-ปาทานิยํ มหาภูตํ. กมฺมสฺส กตตฺตา โผฏฺฐพฺพา\-ยตนํ อาโปธาตุ. อิทํ ตํ รูปํ อุปาทิณฺณุ\-ปาทานิยํ มหาภูตํ.}

\proseitem{931}{กตมํ ตํ รูปํ อุปาทิณฺณุ\-ปาทานิยํ น มหาภูตํ. จกฺขายตนํ ฯเปฯ กายายตนํ อิตฺถินฺทฺริยํ ปุริสินฺทฺริยํ ชีวิตินฺทฺริยํ, ยํ วา ปนญฺญมฺปิ อตฺถิ รูปํ กมฺมสฺส กตตฺตา รูปายตนํ คนฺธายตนํ รสายตนํ อากาสธาตุ รูปสฺส อุปจโย รูปสฺส สนฺตติ กพฬีกาโร อาหาโร. อิทํ ตํ รูปํ อุปาทิณฺณุ\-ปาทานิยํ น มหาภูตํ.}

\proseitem{932}{กตมํ ตํ รูปํ อนุปาทิณฺณุ\-ปาทานิยํ มหาภูตํ. น กมฺมสฺส กตตฺตา โผฏฺฐพฺพา\-ยตนํ อาโปธาตุ. อิทํ ตํ รูปํ อนุปาทิณฺณุ\-ปาทานิยํ มหาภูตํ.}

\proseitem{933}{กตมํ ตํ รูปํ อนุปาทิณฺณุ\-ปาทานิยํ น มหาภูตํ. สทฺทายตนํ กายวิญฺญตฺติ วจีวิญฺญตฺติ รูปสฺส ลหุตา รูปสฺส มุทุตา รูปสฺส กมฺมญฺญตา รูปสฺส ชรตา รูปสฺส อนิจฺจตา, ยํ วา ปนญฺญมฺปิ อตฺถิ รูปํ น กมฺมสฺส อุปจโย รูปสฺส สนฺตติ กพฬีกาโร อาหาโร. อิทํ ตํ รูปํ อนุปาทิณฺณุ\-ปาทานิยํ น มหาภูตํ.}

\proseitem{934}{กตมํ ตํ รูปํ อุปาทิณฺณุ\-ปาทานิยํ โอฬาริกํ. จกฺขายตนํ ฯเปฯ กายายตนํ, ยํ วา ปนญฺญมฺปิ อตฺถิ รูปํ กมฺมสฺส กตตฺตา รูปายตนํ คนฺธายตนํ รสายตนํ โผฏฺฐพฺพา\-ยตนํ. อิทํ ตํ รูปํ อุปาทิณฺณุ\-ปาทานิยํ โอฬาริกํ.}

\proseitem{935}{กตมํ ตํ รูปํ อุปาทิณฺณุ\-ปาทานิยํ สุขุมํ. อิตฺถินฺทฺริยํ ปุริสินฺทฺริยํ ชีวิตินฺทฺริยํ, ยํ วา ปนญฺญมฺปิ อตฺถิ รูปํ กมฺมสฺส กตตฺตา อากาสธาตุ อาโปธาตุ รูปสฺส อุปจโย รูปสฺส สนฺตติ กพฬีกาโร อาหาโร. อิทํ ตํ รูปํ อุปาทิณฺณุ\-ปาทานิยํ สุขุมํ.}

\proseitem{936}{กตมํ ตํ รูปํ อนุปาทิณฺณุ\-ปาทานิยํ โอฬาริกํ. สทฺทายตนํ, ยํ วา ปนญฺญมฺปิ อตฺถิ รูปํ น กมฺมสฺส กตตฺตา รูปายตนํ คนฺธายตนํ รสายตนํ โผฏฺฐพฺพา\-ยตนํ. อิทํ ตํ รูปํ อนุปาทิณฺณุ\-ปาทานิยํ โอฬาริกํ.}

\chapterheadpageread{2. รูปกณฺฑ}
\proseitem{937}{\csromanfolio{199}กตมํ ตํ รูปํ อนุปาทิณฺณุ\-ปาทานิยํ สุขุมํ. กายวิญฺญตฺติ วจีวิญฺญตฺติ รูปสฺส ลหุตา รูปสฺส มุทุตา รูปสฺส กมฺมญฺญตา รูปสฺส ชรตา รูปสฺส อนิจฺจตา, ยํ วา ปนญฺญมฺปิ อตฺถิ รูปํ น กมฺมสฺส กตตฺตา อากาสธาตุ อาโปธาตุ รูปสฺส อุปจโย รูปสฺส สนฺตติ กพฬีกาโร อาหาโร. อิทํ ตํ รูปํ อนุปาทิณฺณุ\-ปาทานิยํ สุขุมํ.}

\proseitem{938}{กตมํ ตํ รูปํ อุปาทิณฺณุ\-ปาทานิยํ ทูเร. อิตฺถินฺทฺริยํ ปุริสินฺทฺริยํ ชีวิตินฺทฺริยํ, ยํ วา ปนญฺญมฺปิ อตฺถิ รูปํ กมฺมสฺส กตตฺตา อากาสธาตุ อาโปธาตุ รูปสฺส อุปจโย รูปสฺส สนฺตติ กพฬีกาโร อาหาโร. อิทํ ตํ รูปํ อุปาทิณฺณุ\-ปาทานิยํ ทูเร.}

\proseitem{939}{กตมํ ตํ รูปํ อุปาทิณฺณุ\-ปาทานิยํ สนฺติเก. จกฺขายตนํ ฯเปฯ กายายตนํ, ยํ วา ปนญฺญมฺปิ อตฺถิ รูปํ กมฺมสฺส กตตฺตา รูปายตนํ คนฺธายตนํ รสายตนํ โผฏฺฐพฺพา\-ยตนํ. อิทํ ตํ รูปํ อุปาทิณฺณุ\-ปาทานิยํ สนฺติเก.}

\proseitem{940}{กตมํ ตํ รูปํ อนุปาทิณฺณุ\-ปาทานิยํ ทูเร. กายวิญฺญตฺติ วจีวิญฺญตฺติ รูปสฺส ลหุตา รูปสฺส มุทุตา รูปสฺส กมฺมญฺญตา รูปสฺส ชรตา รูปสฺส อนิจฺจตา, ยํ วา ปนญฺญมฺปิ อตฺถิ รูปํ น กมฺมสฺส กตตฺตา อากาสธาตุ อาโปธาตุ รูปสฺส อุปจโย รูปสฺส สนฺตติ กพฬีกาโร อาหาโร. อิทํ ตํ รูปํ อนุปาทิณฺณุ\-ปาทานิยํ ทูเร.}

\proseitem{941}{กตมํ ตํ รูปํ อนุปาทิณฺณุ\-ปาทานิยํ สนฺติเก. สทฺทายตนํ, ยํ วา ปนญฺญมฺปิ อตฺถิ รูปํ น กมฺมสฺส กตตฺตา รูปายตนํ คนฺธายตนํ รสายตนํ โผฏฺฐพฺพา\-ยตนํ. อิทํ ตํ รูปํ อนุปาทิณฺณุ\-ปาทานิยํ สนฺติเก.}

\proseitem{942}{กตมํ ตํ รูปํ สปฺปฏิฆํ อินฺทฺริยํ. จกฺขุนฺทฺริยํ ฯเปฯ กายินฺทฺริยํ. อิทํ ตํ รูปํ สปฺปฏิฆํ อินฺทฺริยํ.}

\proseitem{943}{กตมํ ตํ รูปํ สปฺปฏิฆํ น อินฺทฺริยํ. รูปายตนํ ฯเปฯ โผฏฺฐพฺพา\-ยตนํ. อิทํ ตํ รูปํ สปฺปฏิฆํ น อินฺทฺริยํ.}

\proseitem{944}{กตมํ ตํ รูปํ อปฺปฏิฆํ อินฺทฺริยํ. อิตฺถินฺทฺริยํ ปุริสินฺทฺริยํ ชีวิตินฺทฺริยํ. อิทํ ตํ รูปํ อปฺปฏิฆํ อินฺทฺริยํ.}

\proseitem{945}{\csromanfolio{200}กตมํ ตํ รูปํ อปฺปฏิฆํ น อินฺทฺริยํ. กายวิญฺญตฺติ วจีวิญฺญตฺติ ฯเปฯ กพฬีกาโร อาหาโร. อิทํ ตํ รูปํ อปฺปฏิฆํ น อินฺทฺริยํ.}

\proseitem{946}{กตมํ ตํ รูปํ สปฺปฏิฆํ มหาภูตํ. โผฏฺฐพฺพา\-ยตนํ. อิทํ ตํ รูปํ สปฺปฏิฆํ มหาภูตํ.}

\proseitem{947}{กตมํ ตํ รูปํ สปฺปฏิฆํ น มหาภูตํ. จกฺขายตนํ ฯเปฯ รสายตนํ. อิทํ ตํ รูปํ สปฺปฏิฆํ น มหาภูตํ.}

\proseitem{948}{กตมํ ตํ รูปํ อปฺปฏิฆํ มหาภูตํ. อาโปธาตุ. อิทํ ตํ รูปํ อปฺปฏิฆํ มหาภูตํ.}

\proseitem{949}{กตมํ ตํ รูปํ อปฺปฏิฆํ น มหาภูตํ. อิตฺถินฺทฺริยํ ฯเปฯ กพฬีกาโร อาหาโร. อิทํ ตํ รูปํ อปฺปฏิฆํ น มหาภูตํ.}

\proseitem{950}{กตมํ ตํ รูปํ อินฺทฺริยํ โอฬาริกํ. จกฺขุนฺทฺริยํ ฯเปฯ กายินฺทฺริยํ. อิทํ ตํ รูปํ อินฺทฺริยํ โอฬาริกํ.}

\proseitem{951}{กตมํ ตํ รูปํ อินฺทฺริยํ สุขุมํ. อิตฺถินฺทฺริยํ ปุริสินฺทฺริยํ ชีวิตินฺทฺริยํ. อิทํ ตํ รูปํ อินฺทฺริยํ สุขุมํ.}

\proseitem{952}{กตมํ ตํ รูปํ น อินฺทฺริยํ โอฬาริกํ. รูปายตนํ ฯเปฯ โผฏฺฐพฺพา\-ยตนํ. อิทํ ตํ รูปํ น อินฺทฺริยํ โอฬาริกํ.}

\proseitem{953}{กตมํ ตํ รูปํ น อินฺทฺริยํ สุขุมํ. กายวิญฺญตฺติ วจีวิญฺญตฺติ ฯเปฯ กพฬีกาโร อาหาโร. อิทํ ตํ รูปํ น อินฺทฺริยํ สุขุมํ.}

\proseitem{954}{กตมํ ตํ รูปํ อินฺทฺริยํ ทูเร. อิตฺถินฺทฺริยํ ปุริสินฺทฺริยํ ชีวิตินฺทฺริยํ. อิทํ ตํ รูปํ อินฺทฺริยํ ทูเร.}

\proseitem{955}{กตมํ ตํ รูปํ อินฺทฺริยํ สนฺติเก. จกฺขุนฺทฺริยํ ฯเปฯ กายินฺทฺริยํ. อิทํ ตํ รูปํ อินฺทฺริยํ สนฺติเก.}

\proseitem{956}{กตมํ ตํ รูปํ น อินฺทฺริยํ ทูเร. กายวิญฺญตฺติ วจีวิญฺญตฺติ ฯเปฯ กพฬีกาโร อาหาโร. อิทํ ตํ รูปํ น อินฺทฺริยํ ทูเร.}

\proseitem{957}{กตมํ ตํ รูปํ น อินฺทฺริยํ สนฺติเก. รูปายตนํ ฯเปฯ โผฏฺฐพฺพา\-ยตนํ. อิทํ ตํ รูปํ น อินฺทฺริยํ สนฺติเก.}

\proseitem{958}{กตมํ ตํ รูปํ มหาภูตํ โอฬาริกํ. โผฏฺฐพฺพา\-ยตนํ. อิทํ ตํ รูปํ มหาภูตํ โอฬาริกํ.}

\chapterheadpageread{2. รูปกณฺฑ}
\proseitem{959}{\csromanfolio{201}กตมํ ตํ รูปํ มหาภูตํ สุขุมํ. อาโปธาตุ. อิทํ ตํ รูปํ มหาภูตํ สุขุมํ.}

\proseitem{960}{กตมํ ตํ รูปํ น มหาภูตํ โอฬาริกํ. จกฺขายตนํ ฯเปฯ รสายตนํ. อิทํ ตํ รูปํ น มหาภูตํ โอฬาริกํ.}

\proseitem{961}{กตมํ ตํ รูปํ น มหาภูตํ สุขุมํ. อิตฺถินฺทฺริยํ ฯเปฯ กพฬีกาโร อาหาโร. อิทํ ตํ รูปํ น มหาภูตํ สุขุมํ.}

\proseitem{962}{กตมํ ตํ รูปํ มหาภูตํ ทูเร. อาโปธาตุ. อิทํ ตํ รูปํ มหาภูตํ ทูเร.}

\proseitem{963}{กตมํ ตํ รูปํ น มหาภูตํ สนฺติเก. โผฏฺฐพฺพา\-ยตนํ. อิทํ ตํ รูปํ มหาภูตํ สนฺติเก.}

\proseitem{964}{กตมํ ตํ รูปํ น มหาภูตํ ทูเร. อิตฺถินฺทฺริยํ ฯเปฯ กพฬีกาโร อาหาโร. อิทํ ตํ รูปํ น มหาภูตํ ทูเร.}

\proseitem{965}{กตมํ ตํ รูปํ น มหาภูตํ สนฺติเก. จกฺขายตนํ ฯเปฯ รสายตนํ. อิทํ ตํ รูปํ น มหาภูตํ สนฺติเก.}

\proseitem{966}{รูปายตนํ ทิฏฺฐํ สทฺทายตนํ สุตํ คนฺธายตนํ รสายตนํ โผฏฺฐพฺพา\-ยตนํ มุตํ สพฺพํ รูปํ มนสา วิญฺญาตํ รูปํ. เอวํ จตุพฺพิเธน รูปสงฺคโห. จตุกฺกํ.}

\csromanmatikaanchor{mtk.95}
\csromantocmarkref{subsection}{ปญฺจก}{mtk.95}
\bookmark[dest={mtk.95},level=2]{ปญฺจก}
\csromanheader{2}{\textbf{ปญฺจก}}
\proseitem{967}{กตมํ ตํ รูปํ ปถวีธาตุ. ยํ กกฺขฬํ ขรคตํ\csromansharedfootnote{6e77dcf9ff9751fb}{ขริคตํ (ก)} กกฺขฬตฺตํ กกฺขฬภาโว อชฺฌตฺตํ วา พหิทฺธา วา อุปาทิณฺณํ วา อนุปาทิณฺณํ วา. อิทํ ตํ รูปํ ปถวีธาตุ.}

\proseitem{968}{กตมํ ตํ รูปํ อาโปธาตุ. ยํ อาโป อาโปคตํ สิเนโห สิเนหคตํ พนฺธนตฺตํ รูปสฺส อชฺฌตฺตํ วา พหิทฺธา วา อุปาทิณฺณํ วา อนุปาทิณฺณํ วา. อิทํ ตํ รูปํ อาโปธาตุ.}

\proseitem{969}{\csromanfolio{202}กตมํ ตํ รูปํ เตโชธาตุ. ยํ เตโช เตโชคตํ อุสฺมา อุสฺมาคตํ อุสุมํ อุสุมคตํ อชฺฌตฺตํ วา พหิทฺธา วา อุปาทิณฺณํ วา อนุปาทิณฺณํ วา. อิทํ ตํ รูปํ เตโชธาตุ.}

\proseitem{970}{กตมํ ตํ รูปํ วาโยธาตุ. ยํ วาโย วาโยคตํ ถมฺภิตตฺตํ รูปสฺส อชฺฌตฺตํ วา พหิทฺธา วา อุปาทิณฺณํ วา อนุปาทิณฺณํ วา. อิทํ ตํ รูปํ วาโยธาตุ.}

\proseitem{971}{กตมํ ตํ รูปํ อุปาทา. จกฺขายตนํ ฯเปฯ กพฬีกาโร อาหาโร. อิทํ ตํ รูปํ อุปาทา. เอวํ ปญฺจวิเธน รูปสงฺคโห. ปญฺจกํ.}

\csromanmatikaanchor{mtk.96}
\csromantocmarkref{subsection}{ฉกฺก}{mtk.96}
\bookmark[dest={mtk.96},level=2]{ฉกฺก}
\csromanheader{2}{\textbf{ฉกฺก}}
\proseitem{972}{รูปายตนํ จกฺขุ\-วิญฺเญยฺยํ รูปํ สทฺทายตนํ โสตวิญฺเญยฺยํ รูปํ คนฺธายตนํ ฆานวิญฺเญยฺยํ รูปํ รสายตนํ ชิวฺหาวิญฺเญยฺยํ รูปํ โผฏฺฐพฺพา\-ยตนํ กายวิญฺเญยฺยํ รูปํ สพฺพํ รูปํ มโนวิญฺเญยฺยํ รูปํ.เอวํ ฉพฺพิเธน รูปสงฺคโห.ฉกฺกํ.}

\csromanmatikaanchor{mtk.97}
\csromantocmarkref{subsection}{สตฺตก}{mtk.97}
\bookmark[dest={mtk.97},level=2]{สตฺตก}
\csromanheader{2}{\textbf{สตฺตก}}
\proseitem{973}{รูปายตนํ จกฺขุ\-วิญฺเญยฺยํ รูปํ สทฺทายตนํ โสตวิญฺเญยฺยํ รูปํ คนฺธายตนํ ฆานวิญฺเญยฺยํ รูปํ รสายตนํ ชิวฺหาวิญฺเญยฺยํ รูปํ โผฏฺฐพฺพา\-ยตนํ กายวิญฺเญยฺยํ รูปํ รูปายตนํ สทฺทายตนํ คนฺธายตนํ รสายตนํ โผฏฺฐพฺพา\-ยตนํ มโนธาตุ\-วิญฺเญยฺยํ รูปํ สพฺพํ รูปํ มโนวิญฺญาณ\-ธาตุวิญฺเญยฺยํ รูปํ. เอวํ สตฺตวิเธน รูปสงฺคโห. สตฺตกํ.}

\csromanmatikaanchor{mtk.98}
\csromantocmarkref{subsection}{อฏฺฐก}{mtk.98}
\bookmark[dest={mtk.98},level=2]{อฏฺฐก}
\csromanheader{2}{2. รูปกณฺฑ \textbf{อฏฺฐก}}
\proseitem{974}{\csromanfolio{203}รูปายตนํ จกฺขุ\-วิญฺเญยฺยํ รูปํ สทฺทายตนํ โสตวิญฺเญยฺยํ รูปํ คนฺธายตนํ ฆานวิญฺเญยฺยํ รูปํ รสายตนํ ชิวฺหาวิญฺเญยฺยํ รูปํ มนาปิโย โผฏฺฐพฺโพ สุขสมฺผสฺโส กายวิญฺเญยฺยํ รูปํ อมนาปิโย โผฏฺฐพฺโพ ทุกฺข\-สมฺผสฺโส กายวิญฺเญยฺยํ รูปํ รูปายตนํ สทฺทายตนํ คนฺธายตนํ รสายตนํ โผฏฺฐพฺพา\-ยตนํ มโนธาตุ\-วิญฺเญยฺยํ รูปํ สพฺพํ รูปํ มโนวิญฺญาณ\-ธาตุวิญฺเญยฺยํ รูปํ. เอวํ อฏฺฐวิเธน รูปสงฺคโห. อฏฺฐกํ.}

\csromanmatikaanchor{mtk.99}
\csromantocmarkref{subsection}{นวก}{mtk.99}
\bookmark[dest={mtk.99},level=2]{นวก}
\csromanheader{2}{\textbf{นวก}}
\proseitem{975}{กตมํ ตํ รูปํ จกฺขุนฺทฺริยํ. ยํ จกฺขุ จตุนฺนํ มหาภูตานํ อุปาทาย ปสาโท ฯเปฯ สุญฺโญ คาโมเปโส. อิทํ ตํ รูปํ จกฺขุนฺทฺริยํ.}

\proseitem{976}{กตมํ ตํ รูปํ โสตินฺทฺริยํ ฯเปฯ ฆานินฺทฺริยํ ฯเปฯ ชิวฺหินฺทฺริยํ ฯเปฯ กายินฺทฺริยํ ฯเปฯ อิตฺถินฺทฺริยํ ฯเปฯ ปุริสินฺทฺริยํ ฯเปฯ ชีวิตินฺทฺริยํ. โย เตสํ รูปีนํ ธมฺมานํ อายุ ฐิติ ยปนา ยาปนา อิริยนา วตฺตนา ปาลนา ชีวิตํ ชีวิตินฺทฺริยํ. อิทํ ตํ รูปํ ชีวิตินฺทฺริยํ.}

\proseitem{977}{กตมํ ตํ รูปํ น อินฺทฺริยํ. รูปายตนํ ฯเปฯ กพฬีกาโร อาหาโร. อิทํ ตํ รูปํ น อินฺทฺริยํ. เอวํ นววิเธน รูปสงฺคโห. นวกํ.}

\csromanmatikaanchor{mtk.100}
\csromantocmarkref{subsection}{ทสก}{mtk.100}
\bookmark[dest={mtk.100},level=2]{ทสก}
\csromanheader{2}{\textbf{ทสก}}
\proseitem{978}{กตมํ ตํ รูปํ จกฺขุนฺทฺริยํ. ยํ จกฺขุ จตุนฺนํ มหาภูตานํ อุปาทาย ปสาโท ฯเปฯ สุญฺโญ คาโมเปโส. อิทํ ตํ รูปํ จกฺขุนฺทฺริยํ.}

\proseitem{979}{กตมํ ตํ รูปํ โสตินฺทฺริยํ ฯเปฯ ฆานินฺทฺริยํ ฯเปฯ ชิวฺหินฺทฺริยํ ฯเปฯ กายินฺทฺริยํ ฯเปฯ อิตฺถินฺทฺริยํ ฯเปฯ ปุริสินฺทฺริยํ ฯเปฯ ชีวิตินฺทฺริยํ. โย เตสํ รูปีนํ \csromanfolio{204}ธมฺมานํ อายุ ฐิติ ยปนา ยาปนา อิริยนา วตฺตนา ปาลนา ชีวิตํ ชีวิตินฺทฺริยํ. อิทํ ตํ รูปํ ชีวิตินฺทฺริยํ.}

\proseitem{980}{กตมํ ตํ รูปํ น อินฺทฺริยํ สปฺปฏิฆํ. รูปายตนํ ฯเปฯ โผฏฺฐพฺพา\-ยตนํ. อิทํ ตํ รูปํ น อินฺทฺริยํ สปฺปฏิฆํ.}

\proseitem{981}{กตมํ ตํ รูปํ น อินฺทฺริยํ อปฺปฏิฆํ. กายวิญฺญตฺติ ฯเปฯ กพฬีกาโร อาหาโร. อิทํ ตํ รูปํ น อินฺทฺริยํ อปฺปฏิฆํ. เอวํ ทสวิเธน รูปสงฺคโห. ทสกํ.}

\csromanmatikaanchor{mtk.101}
\csromantocmarkref{subsection}{เอกาทสก}{mtk.101}
\bookmark[dest={mtk.101},level=2]{เอกาทสก}
\csromanheader{2}{\textbf{เอกาทสก}}
\proseitem{982}{กตมํ ตํ รูปํ จกฺขายตนํ. ยํ จกฺขุ จตุนฺนํ มหาภูตานํ อุปาทาย ปสาโท ฯเปฯ สุญฺโญ คาโมเปโส. อิทํ ตํ รูปํ จกฺขายตนํ.}

\proseitem{983}{กตมํ ตํ รูปํ โสตายตนํ ฯเปฯ ฆานายตนํ ฯเปฯ ชิวฺหายตนํ ฯเปฯ กายายตนํ ฯเปฯ รูปายตนํ ฯเปฯ สทฺทายตนํ ฯเปฯ คนฺธายตนํ ฯเปฯ รสายตนํ ฯเปฯ โผฏฺฐพฺพา\-ยตนํ. ปถวีธาตุ ฯเปฯ โผฏฺฐพฺพ\-ธาตุเปสา. อิทํ ตํ รูปํ โผฏฺฐพฺพา\-ยตนํ.}

\proseitemwithclosermajor{984}{กตมํ ตํ รูปํ อนิทสฺสนํ อปฺปฏิฆํ ธมฺมายตน\-ปริยาปนฺนํ. อิตฺถินฺทฺริยํ ฯเปฯ กพฬีกาโร อาหาโร. อิทํ ตํ รูปํ อนิทสฺสนํ อปฺปฏิฆํ ธมฺมายตน\-ปริยาปนฺนํ. เอวํ เอกาทสวิเธน รูปสงฺคโห. เอกาทสกํ. อฏฺฐโม ภาณวาโร. รูปวิภตฺติ.}{\textbf{รูปกณฺฑํ นิฏฺฐิตํ.}}
\csromanmatikaanchor{mtk.102}
\csromantocmarknopage{chapter}{3. นิกฺเขปกณฺฑ}{mtk.102}
\bookmark[dest={mtk.102},level=0]{3. นิกฺเขปกณฺฑ}
\chapterheadpageread{3. \textbf{นิกฺเขปกณฺฑ}}
\csromanmatikaanchor{mtk.103}
\csromantocmarkref{section}{ติกนิกฺเขป}{mtk.103}
\bookmark[dest={mtk.103},level=1]{ติกนิกฺเขป}
\csromanheader{1}{\textbf{ติกนิกฺเขป}}
\proseitem{985}{\csromanfolio{205}กตเม ธมฺมา กุสลา. ตีณิ กุสลมูลานิ อโลโภ อโทโส อโมโห, ตํสมฺปยุตฺโต เวทนากฺขนฺโธ สญฺญากฺขนฺโธ สงฺขาร\-กฺขนฺโธ วิญฺญาณ\-กฺขนฺโธ, ตํสมุฏฺฐานํ กายกมฺมํ วจีกมฺมํ มโนกมฺมํ. อิเม ธมฺมา กุสลา.}

\proseitem{986}{กตเม ธมฺมา อกุสลา. ตีณิ อกุสลมูลานิ โลโภ โทโส โมโห, ตเทกฏฺฐา จ กิเลสา, ตํสมฺปยุตฺโต เวทนากฺขนฺโธ สญฺญากฺขนฺโธ สงฺขาร\-กฺขนฺโธ วิญฺญาณ\-กฺขนฺโธ, ตํสมุฏฺฐานํ กายกมฺมํ วจีกมฺมํ มโนกมฺมํ. อิเม ธมฺมา อกุสลา.}

\proseitem{987}{กตเม ธมฺมา อพฺยากตา. กุสลา\-กุสลานํ ธมฺมานํ วิปากา กามาวจรา รูปาวจรา อรูปาวจรา อปริยาปนฺนา เวทนากฺขนฺโธ สญฺญากฺขนฺโธ สงฺขาร\-กฺขนฺโธ วิญฺญาณ\-กฺขนฺโธ, เย จ ธมฺมา กิริยา เนว กุสลา นากุสลา น จ กมฺมวิปากา สพฺพญฺจ รูปํ อสงฺขตา จ ธาตุ. อิเม ธมฺมา อพฺยากตา.}

\proseitem{988}{กตเม ธมฺมา สุขาย เวทนาย สมฺปยุตฺตา. สุขภูมิยํ กามาวจเร รูปาวจเร อปริยาปนฺเน สุขํ เวทนํ ฐเปตฺวา ตํสมฺปยุตฺโต สญฺญากฺขนฺโธ สงฺขาร\-กฺขนฺโธ วิญฺญาณ\-กฺขนฺโธ. อิเม ธมฺมา สุขาย เวทนาย สมฺปยุตฺตา.}

\proseitem{989}{กตเม ธมฺมา ทุกฺขาย เวทนาย สมฺปยุตฺตา. ทุกฺขภูมิยํ กามาวจเร ทุกฺขํ เวทนํ ฐเปตฺวา ตํสมฺปยุตฺโต สญฺญากฺขนฺโธ สงฺขาร\-กฺขนฺโธ วิญฺญาณ\-กฺขนฺโธ. อิเม ธมฺมา ทุกฺขาย เวทนาย สมฺปยุตฺตา.}

\proseitem{990}{กตเม ธมฺมา อทุกฺขม\-สุขาย เวทนาย สมฺปยุตฺตา. อทุกฺขมสุข\-ภูมิยํ กามาวจเร รูปาวจเร อรูปาวจเร อปริยาปนฺเน อทุกฺขมสุขํ เวทนํ ฐเปตฺวา ตํสมฺปยุตฺโต สญฺญากฺขนฺโธ สงฺขาร\-กฺขนฺโธ วิญฺญาณ\-กฺขนฺโธ. อิเม ธมฺมา อทุกฺขม\-สุขาย เวทนาย สมฺปยุตฺตา.}

\proseitem{991}{\csromanfolio{206}กตเม ธมฺมา วิปากา. กุสลา\-กุสลานํ ธมฺมานํ วิปากา กามาวจรา รูปาวจรา อรูปาวจรา อปริยาปนฺนา เวทนากฺขนฺโธ ฯเปฯ วิญฺญาณ\-กฺขนฺโธ. อิเม ธมฺมา วิปากา.}

\proseitem{992}{กตเม ธมฺมา วิปาก\-ธมฺม\-ธมฺมา. กุสลากุสลา ธมฺมา กามาวจรา รูปาวจรา อรูปาวจรา อปริยาปนฺนา เวทนากฺขนฺโธ ฯเปฯ วิญฺญาณ\-กฺขนฺโธ. อิเม ธมฺมา วิปาก\-ธมฺม\-ธมฺมา.}

\proseitem{993}{กตเม ธมฺมา เนว\-วิปาก\-น\-วิปากธมฺม\-ธมฺมา. เย จ ธมฺมา กิริยา เนว กุสลา นากุสลา น จ กมฺมวิปากา สพฺพญฺจ รูปํ อสงฺขตา จ ธาตุ. อิเม ธมฺมา เนว\-วิปาก\-น\-วิปากธมฺม\-ธมฺมา.}

\proseitem{994}{กตเม ธมฺมา อุปาทิณฺณุ\-ปาทานิยา. สาสวา กุสลา\-กุสลานํ ธมฺมานํ วิปากา กามาวจรา รูปาวจรา อรูปาวจรา เวทนากฺขนฺโธ ฯเปฯ วิญฺญาณ\-กฺขนฺโธ, ยญฺจ รูปํ กมฺมสฺส กตตฺตา. อิเม ธมฺมา อุปาทิณฺณุ\-ปาทานิยา.}

\proseitem{995}{กตเม ธมฺมา อนุปาทิณฺณุ\-ปาทานิยา. สาสวา กุสลากุสลา ธมฺมา กามาวจรา รูปาวจรา อรูปาวจรา เวทนากฺขนฺโธ ฯเปฯ วิญฺญาณ\-กฺขนฺโธ, เย จ ธมฺมา กิริยา เนว กุสลา นากุสลา น จ กมฺมวิปากา, ยญฺจ รูปํ น กมฺมสฺส กตตฺตา. อิเม ธมฺมา อนุปาทิณฺณุ\-ปาทานิยา.}

\proseitem{996}{กตเม ธมฺมา อนุปาทิณฺณ\-อนุปาทานิยา. อปริยาปนฺนา มคฺคา จ มคฺคผลานิ จ อสงฺขตา จ ธาตุ. อิเม ธมฺมา อนุปาทิณฺณ- อนุปาทานิยา.}

\proseitem{997}{กตเม ธมฺมา สํกิลิฏฺฐ\-สํกิเลสิกา. ตีณิ อกุสลมูลานิ โลโภ โทโส โมโห, ตเทกฏฺฐา จ กิเลสา, ตํสมฺปยุตฺโต เวทนากฺขนฺโธ ฯเปฯ วิญฺญาณ\-กฺขนฺโธ, ตํสมุฏฺฐานํ กายกมฺมํ วจีกมฺมํ มโนกมฺมํ. อิเม ธมฺมา สํกิลิฏฺฐ\-สํกิเลสิกา.}

\proseitem{998}{กตเม ธมฺมา อสํกิลิฏฺฐ\-สํกิเลสิกา. สาสวา กุสลาพฺยากตา ธมฺมา กามาวจรา รูปาวจรา อรูปาวจรา รูปกฺขนฺโธ \csromanfolio{207}เวทนากฺขนฺโธ สญฺญากฺขนฺโธ สงฺขาร\-กฺขนฺโธ วิญฺญาณ\-กฺขนฺโธ. อิเม ธมฺมา อสํกิลิฏฺฐ\-สํกิเลสิกา.}

\proseitem{999}{กตเม ธมฺมา อสํกิลิฏฺฐ\-อสํกิเลสิกา. อปริยาปนฺนา มคฺคา จ มคฺคผลานิ จ อสงฺขตา จ ธาตุ. อิเม ธมฺมา อสํกิลิฏฺฐ- อสํกิเลสิกา.}

\proseitem{1000}{กตเม ธมฺมา สวิตกฺก\-สวิจารา. สวิตกฺกสวิจาร\-ภูมิยํ กามาวจเร รูปาวจเร อปริยาปนฺเน วิตกฺกวิจาเร ฐเปตฺวา ตํสมฺปยุตฺโต เวทนากฺขนฺโธ ฯเปฯ วิญฺญาณ\-กฺขนฺโธ. อิเม ธมฺมา สวิตกฺก\-สวิจารา.}

\proseitem{1001}{กตเม ธมฺมา อวิตกฺก\-วิจารมตฺตา. อวิตกฺกวิจารมตฺต\-ภูมิยํ รูปาวจเร อปริยาปนฺเน วิจารํ ฐเปตฺวา ตํสมฺปยุตฺโต เวทนากฺขนฺโธ ฯเปฯ วิญฺญาณ\-กฺขนฺโธ. อิเม ธมฺมา อวิตกฺก\-วิจารมตฺตา.}

\proseitem{1002}{กตเม ธมฺมา อวิตกฺก\-อวิจารา. อวิตกฺก\-อวิจารภูมิยํ กามาวจเร รูปาวจเร อรูปาวจเร อปริยาปนฺเน เวทนากฺขนฺโธ ฯเปฯ วิญฺญาณ\-กฺขนฺโธ สพฺพญฺจ รูปํ อสงฺขตา จ ธาตุ. อิเม ธมฺมา อวิตกฺก\-อวิจารา.}

\proseitem{1003}{กตเม ธมฺมา ปีติสหคตา. ปีติภูมิยํ กามาวจเร รูปาวจเร อปริยาปนฺเน ปีติํ ฐเปตฺวา ตํสมฺปยุตฺโต เวทนากฺขนฺโธ ฯเปฯ วิญฺญาณ\-กฺขนฺโธ. อิเม ธมฺมา ปีติสหคตา.}

\proseitem{1004}{กตเม ธมฺมา สุขสหคตา. สุขภูมิยํ กามาวจเร รูปาวจเร อปริยาปนฺเน สุขํ ฐเปตฺวา ตํสมฺปยุตฺโต สญฺญากฺขนฺโธ สงฺขาร\-กฺขนฺโธ วิญฺญาณ\-กฺขนฺโธ. อิเม ธมฺมา สุขสหคตา.}

\proseitem{1005}{กตเม ธมฺมา อุเปกฺขา\-สหคตา. อุเปกฺขา\-ภูมิยํ กามาวจเร รูปาวจเร อรูปาวจเร อปริยาปนฺเน อุเปกฺขํ ฐเปตฺวา ตํสมฺปยุตฺโต สญฺญากฺขนฺโธ สงฺขาร\-กฺขนฺโธ วิญฺญาณ\-กฺขนฺโธ. อิเม ธมฺมา อุเปกฺขา\-สหคตา.}

\proseitem{1006}{กตเม ธมฺมา ทสฺสเนน ปหาตพฺพา. ตีณิ สญฺโญชนานิ สกฺกายทิฏฺฐิ วิจิกิจฺฉา สีลพฺพต\-ปรามาโส.}

\proseitem{1007}{\csromanfolio{208}ตตฺถ กตมา สกฺกายทิฏฺฐิ. อิธ อสฺสุตวา ปุถุชฺชโน อริยานํ อทสฺสาวี อริยธมฺมสฺส อโกวิโท อริยธมฺเม อวินีโต สปฺปุริสานํ อทสฺสาวี สปฺปุริส\-ธมฺมสฺส อโกวิโท สปฺปุริส\-ธมฺเม อวินีโต รูปํ อตฺตโต สมนุปสฺสติ, รูปวนฺตํ วา อตฺตานํ, อตฺตนิ วา รูปํ, รูปสฺมิํ วา อตฺตานํ. เวทนํ อตฺตโต สมนุปสฺสติ, เวทนาวนฺตํ วา อตฺตานํ, อตฺตนิ วา เวทนํ, เวทนาย วา อตฺตานํ. สญฺญํ อตฺตโต สมนุปสฺสติ, สญฺญาวนฺตํ วา อตฺตานํ, อตฺตนิ วา สญฺญํ, สญฺญาย วา อตฺตานํ. สงฺขาเร อตฺตโต สมนุปสฺสติ, สงฺขารวนฺตํ วา อตฺตานํ, อตฺตนิ วา สงฺขาเร, สงฺขาเรสุ วา อตฺตานํ. วิญฺญาณํ อตฺตโต สมนุปสฺสติ, วิญฺญาณวนฺตํ วา อตฺตานํ, อตฺตนิ วา วิญฺญาณํ, วิญฺญาณสฺมิํ วา อตฺตานํ. ยา เอวรูปา ทิฏฺฐิ ทิฏฺฐิคตํ ทิฏฺฐิคหนํ ทิฏฺฐิกนฺตาโร ทิฏฺฐิ\-วิสูกายิกํ\csromansharedfootnote{677280a54adcb188}{ทิฏฺฐิวิสูกายิตํ (สี)} ทิฏฺฐิ\-วิปฺผนฺทิตํ ทิฏฺฐิ\-สญฺโญชนํ คาโห ปฏิคฺคาโห อภินิเวโส ปรามาโส กุมฺมคฺโค มิจฺฉาปโถ มิจฺฉตฺตํ ติตฺถายตนํ วิปริยาส\-คฺคาโห. อยํ วุจฺจติ สกฺกายทิฏฺฐิ.}

\proseitem{1008}{ตตฺถ กตมา วิจิกิจฺฉา. สตฺถริ กงฺขติ วิจิกิจฺฉติ, ธมฺเม กงฺขติ วิจิกิจฺฉติ, สํเฆ กงฺขติ วิจิกิจฺฉติ, สิกฺขาย กงฺขติ วิจิกิจฺฉติ, ปุพฺพนฺเต กงฺขติ วิจิกิจฺฉติ, อปรนฺเต กงฺขติ วิจิกิจฺฉติ, ปุพฺพนฺตา\-ปรนฺเต กงฺขติ วิจิกิจฺฉติ, อิทปฺปจฺจยตา ปฏิจฺจ\-สมุปฺปนฺเนสุ ธมฺเมสุ กงฺขติ วิจิกิจฺฉติ. ยา เอวรูปา กงฺขา กงฺขายนา กงฺขายิตตฺตํ วิมติ วิจิกิจฺฉา เทฺวฬฺหกํ เทฺวธาปโถ สํสโย อเนกํสคฺคาโห อาสปฺปนา ปริสปฺปนา อปริโยคาหนา ถมฺภิตตฺตํ จิตฺตสฺส มโนวิเลโข. อยํ วุจฺจติ วิจิกิจฺฉา.}

\proseitem{1009}{ตตฺถ กตโม สีลพฺพต\-ปรามาโส. อิโต พหิทฺธา สมณ\-พฺราหฺมณานํ สีเลน สุทฺธิ วเตน สุทฺธิ สีลพฺพเตน สุทฺธีติ. ยา เอวรูปา ทิฏฺฐิ ทิฏฺฐิคตํ ทิฏฺฐิคหนํ ทิฏฺฐิกนฺตาโร ทิฏฺฐิ\-วิสูกายิกํ ทิฏฺฐิ\-วิปฺผนฺทิตํ ทิฏฺฐิ\-สญฺโญชนํ คาโห ปฏิคฺคาโห อภินิเวโส ปรามาโส กุมฺมคฺโค มิจฺฉาปโถ มิจฺฉตฺตํ ติตฺถายตนํ วิปริยาส\-คฺคาโห. อยํ วุจฺจติ สีลพฺพต\-ปรามาโส.}

\chapterheadpageread{3. นิกฺเขปกณฺฑ}
\proseitem{1010}{\csromanfolio{209}อิมานิ ตีณิ สํโยชนานิ, ตเทกฏฺฐา จ กิเลสา ตํสมฺปยุตฺโต เวทนากฺขนฺโธ ฯเปฯ วิญฺญาณ\-กฺขนฺโธ, ตํสมุฏฺฐานํ กายกมฺมํ วจีกมฺมํ มโนกมฺมํ. อิเม ธมฺมา ทสฺสเนน ปหาตพฺพา.}

\proseitem{1011}{กตเม ธมฺมา ภาวนาย ปหาตพฺพา. อวเสโส โลโภ โทโส โมโห, ตเทกฏฺฐา จ กิเลสา, ตํสมฺปยุตฺโต เวทนากฺขนฺโธ ฯเปฯ วิญฺญาณ\-กฺขนฺโธ, ตํสมุฏฺฐานํ กายกมฺมํ วจีกมฺมํ มโนกมฺมํ. อิเม ธมฺมา ภาวนาย ปหาตพฺพา.}

\proseitem{1012}{กตเม ธมฺมา เนว ทสฺสเนน น ภาวนาย ปหาตพฺพา. กุสลาพฺยากตา ธมฺมา กามาวจรา รูปาวจรา อรูปาวจรา อปริยาปนฺนา เวทนากฺขนฺโธ ฯเปฯ วิญฺญาณ\-กฺขนฺโธ สพฺพญฺจ รูปํ อสงฺขตา จ ธาตุ. อิเม ธมฺมา เนว ทสฺสเนน น ภาวนาย ปหาตพฺพา.}

\proseitem{1013}{กตเม ธมฺมา ทสฺสเนน ปหาตพฺพ\-เหตุกา. ตีณิ สญฺโญชนานิ สกฺกายทิฏฺฐิ วิจิกิจฺฉา สีลพฺพต\-ปรามาโส.}

\proseitem{1014}{ตตฺถ กตมา สกฺกายทิฏฺฐิ ฯเปฯ. อยํ วุจฺจติ สกฺกายทิฏฺฐิ.}

\proseitem{1015}{ตตฺถ กตมา วิจิกิจฺฉา ฯเปฯ. อยํ วุจฺจติ วิจิกิจฺฉา.}

\proseitem{1016}{ตตฺถ กตโม สีลพฺพต\-ปรามาโส ฯเปฯ. อยํ วุจฺจติ สีลพฺพต\-ปรามาโส.}

\proseitem{1017}{อิมานิ ตีณิ สญฺโญชนานิ, ตเทกฏฺฐา จ กิเลสา ตํสมฺปยุตฺโต เวทนากฺขนฺโธ ฯเปฯ วิญฺญาณ\-กฺขนฺโธ, ตํสมุฏฺฐานํ กายกมฺมํ วจีกมฺมํ มโนกมฺมํ, อิเม ธมฺมา ทสฺสเนน ปหาตพฺพ\-เหตุกา. ตีณิ สญฺโญชนานิ สกฺกายทิฏฺฐิ วิจิกิจฺฉา สีลพฺพต\-ปรามาโส, อิเม ธมฺมา ทสฺสเนน ปหาตพฺพา. ตเทกฏฺโฐ โลโภ โทโส โมโห, อิเม ธมฺมา ทสฺสเนน ปหาตพฺพเหตู. ตเทกฏฺฐา จ กิเลสา ตํสมฺปยุตฺโต เวทนากฺขนฺโธ ฯเปฯ วิญฺญาณ\-กฺขนฺโธ, ตํสมุฏฺฐานํ กายกมฺมํ วจีกมฺมํ มโนกมฺมํ. อิเม ธมฺมา ทสฺสเนน ปหาตพฺพ\-เหตุกา.}

\proseitem{1018}{กตเม ธมฺมา ภาวนาย ปหาตพฺพ\-เหตุกา. อวเสโส โลโภ โทโส โมโห, อิเม ธมฺมา ภาวนาย ปหาตพฺพเหตู. \csromanfolio{210}ตเทกฏฺฐา จ กิเลสา ตํสมฺปยุตฺโต เวทนากฺขนฺโธ ฯเปฯ วิญฺญาณ\-กฺขนฺโธ, ตํสมุฏฺฐานํ กายกมฺมํ วจีกมฺมํ มโนกมฺมํ. อิเม ธมฺมา ภาวนาย ปหาตพฺพ\-เหตุกา.}

\proseitem{1019}{กตเม ธมฺมา เนว ทสฺสเนน น ภาวนาย ปหาตพฺพ\-เหตุกา. เต ธมฺเม ฐเปตฺวา อวเสสา กุสลากุสลา\-พฺยากตา ธมฺมา กามาวจรา รูปาวจรา อรูปาวจรา อปริยาปนฺนา เวทนากฺขนฺโธ สญฺญากฺขนฺโธ สงฺขาร\-กฺขนฺโธ วิญฺญาณ\-กฺขนฺโธ สพฺพญฺจ รูปํ อสงฺขตา จ ธาตุ. อิเม ธมฺมา เนว ทสฺสเนน น ภาวนาย ปหาตพฺพ\-เหตุกา.}

\proseitem{1020}{กตเม ธมฺมา อาจยคามิโน. สาสวา กุสลากุสลา ธมฺมา กามาวจรา รูปาวจรา อรูปาวจรา เวทนากฺขนฺโธ ฯเปฯ วิญฺญาณ\-กฺขนฺโธ. อิเม ธมฺมา อาจยคามิโน.}

\proseitem{1021}{กตเม ธมฺมา อปจยคามิโน. จตฺตาโร มคฺคา อปริยาปนฺนา. อิเม ธมฺมา อปจยคามิโน.}

\proseitem{1022}{กตเม ธมฺมา เนว อาจยคามิ น อปจยคามิโน. กุสลา\-กุสลานํ ธมฺมานํ วิปากา กามาวจรา รูปาวจรา อรูปาวจรา อปริยาปนฺนา เวทนากฺขนฺโธ ฯเปฯ วิญฺญาณ\-กฺขนฺโธ, เย จ ธมฺมา กิริยา เนว กุสลา นากุสลา น จ กมฺมวิปากา สพฺพญฺจ รูปํ อสงฺขตา จ ธาตุ. อิเม ธมฺมา เนว อาจยคามิ น อปจยคามิโน.}

\proseitem{1023}{กตเม ธมฺมา เสกฺขา. จตฺตาโร มคฺคา อปริยาปนฺนา เหฏฺฐิมานิ จ ตีณิ สามญฺญผลานิ. อิเม ธมฺมา เสกฺขา.}

\proseitem{1024}{กตเม ธมฺมา อเสกฺขา. อุปริฏฺฐิมํ\csromansharedfootnote{21a280921f1b264a}{อุปริมํ (สฺยา)} อรหตฺตผลํ. อิเม ธมฺมา อเสกฺขา.}

\proseitem{1025}{กตเม ธมฺมา เนว\-เสกฺข\-นาเสกฺขา. เต ธมฺเม ฐเปตฺวา อวเสสา กุสลากุสลา\-พฺยากตา ธมฺมา กามาวจรา รูปาวจรา อรูปาวจรา เวทนากฺขนฺโธ ฯเปฯ วิญฺญาณ\-กฺขนฺโธ สพฺพญฺจ รูปํ อสงฺขตา จ ธาตุ. อิเม ธมฺมา เนว\-เสกฺข\-นาเสกฺขา.}

\chapterheadpageread{3. นิกฺเขปกณฺฑ}
\proseitem{1026}{\csromanfolio{211}กตเม ธมฺมา ปริตฺตา. สพฺเพว กามาวจรา กุสลากุสลา\-พฺยากตา ธมฺมา รูปกฺขนฺโธ ฯเปฯ วิญฺญาณ\-กฺขนฺโธ. อิเม ธมฺมา ปริตฺตา.}

\proseitem{1027}{กตเม ธมฺมา มหคฺคตา. รูปาวจรา อรูปาวจรา กุสลาพฺยากตา ธมฺมา เวทนากฺขนฺโธ ฯเปฯ วิญฺญาณ\-กฺขนฺโธ. อิเม ธมฺมา มหคฺคตา.}

\proseitem{1028}{กตเม ธมฺมา อปฺปมาณา. อปริยาปนฺนา มคฺคา จ มคฺคผลานิ จ อสงฺขตา จ ธาตุ. อิเม ธมฺมา อปฺปมาณา.}

\proseitem{1029}{กตเม ธมฺมา ปริตฺตารมฺมณา. ปริตฺเต ธมฺเม อารพฺภ เย อุปฺปชฺชนฺติ จิตฺตเจตสิกา ธมฺมา. อิเม ธมฺมา ปริตฺตารมฺมณา.}

\proseitem{1030}{กตเม ธมฺมา มหคฺคตา\-รมฺมณา. มหคฺคเต ธมฺเม อารพฺภ เย อุปฺปชฺชนฺติ จิตฺตเจตสิกา ธมฺมา. อิเม ธมฺมา มหคฺคตา\-รมฺมณา.}

\proseitem{1031}{กตเม ธมฺมา อปฺปมาณา\-รมฺมณา. อปฺปมาเณ ธมฺเม อารพฺภ เย อุปฺปชฺชนฺติ จิตฺตเจตสิกา ธมฺมา. อิเม ธมฺมา อปฺปมาณา\-รมฺมณา.}

\proseitem{1032}{กตเม ธมฺมา หีนา. ตีณิ อกุสลมูลานิ โลโภ โทโส โมโห, ตเทกฏฺฐา จ กิเลสา ตํสมฺปยุตฺโต เวทนากฺขนฺโธ ฯเปฯ วิญฺญาณ\-กฺขนฺโธ, ตํสมุฏฺฐานํ กายกมฺมํ วจีกมฺมํ มโนกมฺมํ. อิเม ธมฺมา หีนา.}

\proseitem{1033}{กตเม ธมฺมา มชฺฌิมา. สาสวา กุสลาพฺยากตา ธมฺมา กามาวจรา รูปาวจรา อรูปาวจรา รูปกฺขนฺโธ ฯเปฯ วิญฺญาณ\-กฺขนฺโธ. อิเม ธมฺมา มชฺฌิมา.}

\proseitem{1034}{กตเม ธมฺมา ปณีตา. อปริยาปนฺนา มคฺคา จ มคฺคผลานิ จ อสงฺขตา จ ธาตุ. อิเม ธมฺมา ปณีตา.}

\proseitem{1035}{กตเม ธมฺมา มิจฺฉตฺต\-นิยตา. ปญฺจ กมฺมานิ อานนฺตริกานิ ยา จ มิจฺฉาทิฏฺฐิ นิยตา. อิเม ธมฺมา มิจฺฉตฺต\-นิยตา.}

\proseitem{1036}{กตเม ธมฺมา สมฺมตฺตนิยตา. จตฺตาโร มคฺคา อปริยาปนฺนา. อิเม ธมฺมา สมฺมตฺตนิยตา.}

\proseitem{1037}{\csromanfolio{212}กตเม ธมฺมา อนิยตา. เต ธมฺเม ฐเปตฺวา อวเสสา กุสลากุสลา\-พฺยากตา ธมฺมา กามาวจรา รูปาวจรา อรูปาวจรา อปริยาปนฺนา เวทนากฺขนฺโธ ฯเปฯ วิญฺญาณ\-กฺขนฺโธ สพฺพญฺจ รูปํ อสงฺขตา จ ธาตุ. อิเม ธมฺมา อนิยตา.}

\proseitem{1038}{กตเม ธมฺมา มคฺคารมฺมณา. อริยมคฺคํ อารพฺภ เย อุปฺปชฺชนฺติ จิตฺตเจตสิกา ธมฺมา. อิเม ธมฺมา มคฺคารมฺมณา.}

\proseitem{1039}{กตเม ธมฺมา มคฺคเหตุกา. อริยมคฺค\-สมงฺคิสฺส มคฺคงฺคานิ ฐเปตา ตํสมฺปยุตฺโต เวทนากฺขนฺโธ ฯเปฯ วิญฺญาณ\-กฺขนฺโธ, อิเม ธมฺมา มคฺคเหตุกา. อริยมคฺค\-สมงฺคิสฺส สมฺมาทิฏฺฐิ มคฺโค เจว เหตุ จ, สมฺมาทิฏฺฐิํ ฐเปตฺวา ตํสมฺปยุตฺโต เวทนากฺขนฺโธ ฯเปฯ วิญฺญาณ\-กฺขนฺโธ, อิเม ธมฺมา มคฺคเหตุกา. อริยมคฺค\-สมงฺคิสฺส อโลโภ อโทโส อโมโห, อิเม ธมฺมา มคฺคเหตู. ตํสมฺปยุตฺโต เวทนากฺขนฺโธ ฯเปฯ วิญฺญาณ\-กฺขนฺโธ. อิเม ธมฺมา มคฺคเหตุกา.}

\proseitem{1040}{กตเม ธมฺมา มคฺคาธิปติโน. อริยมคฺคํ อธิปติํ กริตฺวา เย อุปฺปชฺชนฺติ จิตฺตเจตสิกา ธมฺมา, อิเม ธมฺมา มคฺคาธิปติโน. อริยมคฺค\-สมงฺคิสฺส วีมํสา\-ธิปเตยฺยํ มคฺคํ ภาวยนฺตสฺส วีมํสํ ฐเปตฺวา ตํสมฺปยุตฺโต เวทนากฺขนฺโธ ฯเปฯ วิญฺญาณ\-กฺขนฺโธ. อิเม ธมฺมา มคฺคาธิปติโน.}

\proseitem{1041}{กตเม ธมฺมา อุปฺปนฺนา. เย ธมฺมา ชาตา ภูตา สญฺชาตา นิพฺพตฺตา อภินิพฺพตฺตา ปาตุภูตา อุปฺปนฺนา สมุปฺปนฺนา อุฏฺฐิตา สมุฏฺฐิตา อุปฺปนฺนา อุปฺปนฺนํเสน สงฺคหิตา รูปํ\csromansharedfootnote{36636428f9dafb9a}{รูปา (พหูสุ)} เวทนา สญฺญา สงฺขารา วิญฺญาณํ. อิเม ธมฺมา อุปฺปนฺนา.}

\proseitem{1042}{กตเม ธมฺมา อนุปฺปนฺนา. เย ธมฺมา อชาตา อภูตา อสญฺชาตา อนิพฺพตฺตา อนภินิพฺพตฺตา อปาตุภูตา อนุปฺปนฺนา อสมุปฺปนฺนา อนุฏฺฐิตา อสมุฏฺฐิตา อนุปฺปนฺนา อนุปฺปนฺนํเสน สงฺคหิตา รูปํ เวทนา สญฺญา สงฺขารา วิญฺญาณํ. อิเม ธมฺมา อนุปฺปนฺนา.}

\proseitem{1043}{กตเม ธมฺมา อุปฺปาทิโน. กุสลา\-กุสลานํ ธมฺมานํ อวิปกฺก\-วิปากานํ วิปากา กามาวจรา รูปาวจรา อรูปาวจรา อปริยาปนฺนา เวทนากฺขนฺโธ ฯเปฯ วิญฺญาณ\-กฺขนฺโธ, ยญฺจ รูปํ กมฺมสฺส กตตฺตา อุปฺปชฺชิสฺสติ. อิเม ธมฺมา อุปฺปาทิโน.}

\chapterheadpageread{3. นิกฺเขปกณฺฑ}
\proseitem{1044}{\csromanfolio{213}กตเม ธมฺมา อตีตา. เย ธมฺมา อตีตา นิรุทฺธา วิคตา วิปริณตา อตฺถงฺคตา อพฺภตฺถงฺคตา อุปฺปชฺชิตฺวา วิคตา อตีตา อตีตํเสน สงฺคหิตา รูปํ เวทนา สญฺญา สงฺขารา วิญฺญาณํ. อิเม ธมฺมา อตีตา.}

\proseitem{1045}{กตเม ธมฺมา อนาคตา. เย ธมฺมา อชาตา อภูตา อสญฺชาตา อนิพฺพตฺตา อนภินิพฺพตฺตา อปาตุภูตา อนุปฺปนฺนา อสมุปฺปนฺนา อนุฏฺฐิตา อสมุฏฺฐิตา อนาคตา อนาคตํเสน สงฺคหิตา รูปํ เวทนา สญฺญา สงฺขารา วิญฺญาณํ. อิเม ธมฺมา อนาคตา.}

\proseitem{1046}{กตเม ธมฺมา ปจฺจุปฺปนฺนา. เย ธมฺมา ชาตา ภูตา สญฺชาตา นิพฺพตฺตา อภินิพฺพตฺตา ปาตุภูตา อุปฺปนฺนา สมุปฺปนฺนา อุฏฺฐิตา สมุฏฺฐิตา ปจฺจุปฺปนฺนา ปจฺจุปฺปนฺนํ\-เสน สงฺคหิตา รูปํ เวทนา สญฺญา สงฺขารา วิญฺญาณํ. อิเม ธมฺมา ปจฺจุปฺปนฺนา.}

\proseitem{1047}{กตเม ธมฺมา อตีตารมฺมณา. อตีเต ธมฺเม อารพฺภ เย อุปฺปชฺชนฺติ จิตฺตเจตสิกา ธมฺมา. อิเม ธมฺมา อตีตารมฺมณา.}

\proseitem{1048}{กตเม ธมฺมา อนาคตารมฺมณา. อนาคเต ธมฺเม อารพฺภ เย อุปฺปชฺชนฺติ จิตฺตเจตสิกา ธมฺมา. อิเม ธมฺมา อนาคตารมฺมณา.}

\proseitem{1049}{กตเม ธมฺมา ปจฺจุปฺปนฺนา\-รมฺมณา. ปจฺจุปฺปนฺเน ธมฺเม อารพฺภ เย อุปฺปชฺชนฺติ จิตฺตเจตสิกา ธมฺมา. อิเม ธมฺมา ปจฺจุปฺปนฺนา\-รมฺมณา.}

\proseitem{1050}{กตเม ธมฺมา อชฺฌตฺตา. เย ธมฺมา เตสํ เตสํ สตฺตานํ อชฺฌตฺตํ ปจฺจตฺตํ นิยตา ปาฏิปุคฺคลิกา อุปาทิณฺณา รูปํ เวทนา สญฺญา สงฺขารา วิญฺญาณํ. อิเม ธมฺมา อชฺฌตฺตา.}

\proseitem{1051}{กตเม ธมฺมา พหิทฺธา. เย ธมฺมา เตสํ เตสํ ปรสตฺตานํ ปรปุคฺคลานํ อชฺฌตฺตํ ปจฺจตฺตํ นิยตา ปาฏิปุคฺคลิกา อุปาทิณฺณา รูปํ เวทนา สญฺญา สงฺขารา วิญฺญาณํ. อิเม ธมฺมา พหิทฺธา.}

\proseitem{1052}{กตเม ธมฺมา อชฺฌตฺต\-พหิทฺธา. ตทุภยํ. อิเม ธมฺมา อชฺฌตฺต\-พหิทฺธา.}

\proseitem{1053}{กตเม ธมฺมา อชฺฌตฺตา\-รมฺมณา. อชฺฌตฺเต ธมฺเม อารพฺภ เย อุปฺปชฺชนฺติ จิตฺตเจตสิกา ธมฺมา. อิเม ธมฺมา อชฺฌตฺตา\-รมฺมณา.}

\proseitem{1054}{\csromanfolio{214}กตเม ธมฺมา พหิทฺธา\-รมฺมณา. พหิทฺธา ธมฺเม อารพฺภ เย อุปฺปชฺชนฺติ จิตฺตเจตสิกา ธมฺมา. อิเม ธมฺมา พหิทฺธา\-รมฺมณา.}

\proseitem{1055}{กตเม ธมฺมา อชฺฌตฺต\-พหิทฺธา\-รมฺมณา. อชฺฌตฺต\-พหิทฺธา ธมฺเม อารพฺภ เย อุปฺปชฺชนฺติ จิตฺตเจตสิกา ธมฺมา. อิเม ธมฺมา อชฺฌตฺต\-พหิทฺธา\-รมฺมณา.}

\proseitem{1056}{กตเม ธมฺมา สนิทสฺสน\-สปฺปฏิฆา. รูปายตนํ. อิเม ธมฺมา สนิทสฺสน\-สปฺปฏิฆา.}

\proseitem{1057}{กตเม ธมฺมา อนิทสฺสน\-สปฺปฏิฆา. จกฺขายตนํ โสตายตนํ ฆานายตนํ ชิวฺหายตนํ กายายตนํ สทฺทายตนํ คนฺธายตนํ รสายตนํ โผฏฺฐพฺพา\-ยตนํ. อิเม ธมฺมา อนิทสฺสน\-สปฺปฏิฆา.}

\proseitem{1058}{กตเม ธมฺมา อนิทสฺสน\-อปฺปฏิฆา. เวทนากฺขนฺโธ สญฺญากฺขนฺโธ สงฺขาร\-กฺขนฺโธ วิญฺญาณ\-กฺขนฺโธ, ยญฺจ รูปํ อนิทสฺสนํ อปฺปฏิฆํ ธมฺมายตน\-ปริยาปนฺนํ อสงฺขตา จ ธาตุ. อิเม ธมฺมา อนิทสฺสน\-อปฺปฏิฆา.}

\vspace{\tipitakaparskip}
\csromancenter{ติกํ.}
\csromanmatikaanchor{mtk.104}
\csromantocmarkref{section}{ทุกนิกฺเขปเหตุโคจฺฉก}{mtk.104}
\bookmark[dest={mtk.104},level=1]{ทุกนิกฺเขปเหตุโคจฺฉก}
\csromanheader{1}{\textbf{ทุกนิกฺเขป}}
\csromanheader{2}{\textbf{เหตุโคจฺฉก}}
\proseitem{1059}{กตเม ธมฺมา เหตู. ตโย กุสลเหตู, ตโย อกุสลเหตู, ตโย อพฺยากตเหตู, นว กามาวจรเหตู, ฉ รูปาวจรเหตู, ฉ อรูปาวจร\-เหตู, ฉ อปริยาปนฺน\-เหตู.}

\proseitem{1060}{ตตฺถ กตเม ตโย กุสลเหตู. อโลโภ อโทโส อโมโห.}

\proseitem{1061}{ตตฺถ กตโม อโลโภ. โย อโลโภ อลุพฺภนา อลุพฺภิตตฺตํ อสาราโค อสารชฺชนา อสารชฺชิตตฺตํ อนภิชฺฌา อโลโภ กุสลมูลํ. อยํ วุจฺจติ อโลโภ.}

\chapterheadpageread{3. นิกฺเขปกณฺฑ}
\proseitem{1062}{\csromanfolio{215}ตตฺถ กตโม อโทโส. โย อโทโส อทุสฺสนา อทุสฺสิตตฺตํ เมตฺติ เมตฺตายนา เมตฺตายิตตฺตํ อนุทฺทา อนุทฺทายนา อนุทายิตตฺตํ หิเตสิตา อนุกมฺปา อพฺยาปาโท อพฺยาปชฺโช อโทโส กุสลมูลํ. อยํ วุจฺจติ อโทโส.}

\proseitem{1063}{ตตฺถ กตโม อโมโห. ทุกฺเข ญาณํ, ทุกฺขสมุทเย ญาณํ, ทุกฺขนิโรเธ ญาณํ, ทุกฺขนิโรธ\-คามินิยา ปฏิปทาย ญาณํ, ปุพฺพนฺเตญาณํ, อปรนฺเต ญาณํ, ปุพฺพนฺตา\-ปรนฺเต ญาณํ, อิทปฺปจฺจยตา ปฏิจฺจ\-สมุปฺปนฺเนสุ ธมฺเมสุ ญาณํ. ยา เอวรูปา ปญฺญา ปชานนา วิจโย ปวิจโย ธมฺมวิจโย สลฺลกฺขณา อุปลกฺขณา ปจฺจุ\-ปลกฺขณา ปณฺฑิจฺจํ โกสลฺลํ เนปุญฺญํ เวภพฺยา จินฺตา อุปปริกฺขา ภูรี เมธา ปริณายิกา วิปสฺสนา สมฺปชญฺญํ ปโตโท ปญฺญา ปญฺญินฺทฺริยํ ปญฺญาพลํ ปญฺญาสตฺถํ ปญฺญาปาสาโท ปญฺญาอาโลโก ปญฺญาโอภาโส ปญฺญาปชฺโชโต ปญฺญารตนํ อโมโห ธมฺมวิจโย สมฺมาทิฏฺฐิ. อยํ วุจฺจติ อโมโห. อิเม ตโย กุสลเหตู.}

\proseitem{1064}{ตตฺถ กตเม ตโย อกุสลเหตู. โลโภ โทโส โมโห.}

\proseitem{1065}{ตตฺถ กตโม โลโภ. โย ราโค สาราโค อนุนโย อนุโรโธ นนฺที นนฺทีราโค\csromansharedfootnote{2a08338d936ea781}{นนฺทิราโค (สี)} จิตฺตสฺส สาราโค อิจฺฉา มุจฺฉา อชฺโฌสานํ เคโธ ปลิเคโธ สงฺโค ปงฺโค เอชา มายา ชนิกา สญฺชนนี สิพฺพินี\csromansharedfootnote{38799d9af8877ecf}{สิพฺพนี (สี)} ชาลินี สริตา วิสตฺติกา สุตฺตํ วิสฏา อายูหินี\csromansharedfootnote{b49c8d7eb1025482}{อายูหนี (สี, สฺยา)} ทุติยา ปณิธิ ภวเนตฺติ วนํ วนโถ สนฺถโว สิเนโห อเปกฺขา ปฏิพนฺธุ อาสา อาสิสนา อาสิสิตตฺตํ\csromansharedfootnote{276379439a218fd2}{อาสิํสนา อาสิํสิตตฺตํ (สี, สฺยา)} รูปาสา สทฺทาสา คนฺธาสา รสาสา โผฏฺฐพฺพาสา ลาภาสา ธนาสา ปุตฺตาสา ชีวิตาสา ชปฺปา ปชปฺปา อภิชปฺปา ชปฺปา ชปฺปนา ชปฺปิตตฺตํ โลลุปฺปํ โลลุปฺปายนา โลลุปฺปา\-ยิตตฺตํ ปุจฺฉญฺชิกตา\csromansharedfootnote{2c4ea6fdb7fd1562}{ปุญฺจิกตา (สฺยา) ปุจฺฉิกตา (สี)} สาธุกมฺยตา อธมฺมราโค วิสมโลโภ นิกนฺติ นิกามนา ปตฺถนา ปิหนา สมฺปตฺถนา กามตณฺหา ภวตณฺหา วิภวตณฺหา รูปตณฺหา อรูปตณฺหา นิโรธตณฺหา รูปตณฺหา สทฺทตณฺหา \csromanfolio{216}คนฺธตณฺหา รสตณฺหา โผฏฺฐพฺพ\-ตณฺหา ธมฺมตณฺหา โอโฆ โยโค คนฺโถ อุปาทานํ อาวรณํ นีวรณํ ฉาทนํ พนฺธนํ อุปกฺกิเลโส อนุสโย ปริยุฏฺฐานํ ลตา เววิจฺฉํ ทุกฺขมูลํ ทุกฺขนิทานํ ทุกฺขปฺปภโว มารปาโส มารพฬิสํ มารวิสโย ตณฺหานที ตณฺหาชาลํ ตณฺหาคทฺทุลํ ตณฺหาสมุทฺโท อภิชฺฌา โลโภ อกุสลมูลํ. อยํ วุจฺจติ โลโภ.}

\proseitem{1066}{ตตฺถ กตโม โทโส. “อนตฺถํ เม อจรี”ติ อาฆาโต ชายติ, “อนตฺถํ เม จรตี”ติ อาฆาโต ชายติ, “อนตฺถํ เม จริสฺสตี”ติ อาฆาโต ชายติ, “ปิยสฺส เม มนาปสฺส อนตฺถํ อจริ ฯเปฯ อนตฺถํ จรติ ฯเปฯ อนตฺถํ จริสฺสตี”ติ อาฆาโต ชายติ, “อปฺปิยสฺส เม อมนาปสฺส อตฺถํ อจริ ฯเปฯ อตฺถํ จรติ ฯเปฯ อตฺถํ จริสฺสตี”ติ อาฆาโต ชายติ, อฏฺฐาเน วา ปน อาฆาโต ชายติ. โย เอวรูโป จิตฺตสฺส อาฆาโต ปฏิฆาโต ปฏิฆํ ปฏิวิโรโธ โกโป ปโกโป สมฺปโกโป โทโส ปโทโส สมฺปโทโส จิตฺตสฺส พฺยาปตฺติ มโนปโทโส โกโธ กุชฺฌนา กุชฺฌิตตฺตํ โทโส ทุสฺสนา ทุสฺสิตตฺตํ พฺยาปตฺติ พฺยาปชฺชนา พฺยาปชฺชิตตฺตํ วิโรโธ ปฏิวิโรโธ จณฺฑิกฺกํ อสุโรโป อนตฺตมนตา จิตฺตสฺส. อยํ วุจฺจติ โทโส.}

\proseitem{1067}{ตตฺถ กตโม โมโห. ทุกฺเข อญฺญาณํ, ทุกฺขสมุทเย อญฺญาณํ, ทุกฺขนิโรเธ อญฺญาณํ, ทุกฺขนิโรธ\-คามินิยา ปฏิปทาย อญฺญาณํ, ปุพฺพนฺเต อญฺญาณํ, อปรนฺเต อญฺญาณํ, ปุพฺพนฺตา\-ปรนฺเต อญฺญาณํ, อิทปฺปจฺจยตา ปฏิจฺจ\-สมุปฺปนฺเนสุ ธมฺเมสุ อญฺญาณํ. ยํ เอวรูปํ อญฺญาณํ อทสฺสนํ อนภิสมโย อนนุโพโธ อสมฺโพโธ อปฺปฏิเวโธ อสํคาหนา อปริโยคาหนา อสมเปกฺขนา อปจฺจเวกฺขณา อปจฺจกฺขกมฺมํ ทุมฺเมชฺฌํ พาลฺยํ อสมฺปชญฺญํ โมโห ปโมโห สมฺโมโห อวิชฺชา อวิชฺโชโฆ อวิชฺชาโยโค อวิชฺชานุสโย อวิชฺชา\-ปริยุฏฺฐานํ อวิชฺชาลงฺคี โมโห อกุสลมูลํ. อยํ วุจฺจติ โมโห. อิเม ตโย อกุสลเหตู.}

\proseitem{1068}{ตตฺถ กตเม ตโย อพฺยากตเหตู. กุสลานํ วา ธมฺมานํ วิปากโต กิริยา\-พฺยากเตสุ วา ธมฺเมสุ อโลโภ อโทโส อโมโห. อิเม ตโย อพฺยากตเหตู.}

\chapterheadpageread{3. นิกฺเขปกณฺฑ}
\proseitem{1069}{\csromanfolio{217}ตตฺถ กตเม นว กามาวจรเหตู. ตโย กุสลเหตู, ตโย อกุสลเหตู, ตโย อพฺยากตเหตู. อิเม นว กามาวจรเหตู.}

\proseitem{1070}{ตตฺถ กตเม ฉ รูปาวจรเหตู. ตโย กุสลเหตู, ตโย อพฺยากตเหตู. อิเม ฉ รูปาวจรเหตู.}

\proseitem{1071}{ตตฺถ กตเม ฉ อรูปาวจร\-เหตู. ตโย กุสลเหตู, ตโย อพฺยากตเหตู. อิเม ฉ อรูปาวจร\-เหตู.}

\proseitem{1072}{ตตฺถ กตเม ฉ อปริยาปนฺน\-เหตู. ตโย กุสลเหตู, ตโย อพฺยากตเหตู. อิเม ฉ อปริยาปนฺน\-เหตู.}

\proseitem{1073}{ตตฺถ กตเม ตโย กุสลเหตู. อโลโภ อโทโส อโมโห.}

\proseitem{1074}{ตตฺถ กตโม อโลโภ. โย อโลโภ อลุพฺภนา อลุพฺภิตตฺตํ อสาราโค อสารชฺชนา อสารชฺชิตตฺตํ อนภิชฺฌา อโลโภ กุสลมูลํ. อยํ วุจฺจติ อโลโภ.}

\proseitem{1075}{ตตฺถ กตโม อโทโส. โย อโทโส อทุสฺสนา อทุสฺสิตตฺตํ ฯเปฯ อพฺยาปาโท อพฺยาปชฺโช อโทโส กุสลมูลํ. อยํ วุจฺจติ อโทโส.}

\proseitem{1076}{ตตฺถ กตโม อโมโห. ทุกฺเข ญาณํ, ทุกฺขสมุทเย ญาณํ, ทุกฺขนิโรเธ ญาณํ, ทุกฺขนิโรธ\-คามินิยา ปฏิปทาย ญาณํ, ปุพฺพนฺเต ญาณํ, อปรนฺเต ญาณํ, ปุพฺพนฺตา\-ปรนฺเต ญาณํ, อิทปฺปจฺจยตา ปฏิจฺจ\-สมุปฺปนฺเนสุ ธมฺเมสุ ญาณํ. ยา เอวรูปา ปญฺญา ปชานนา วิจโย ปวิจโย ธมฺมวิจโย สลฺลกฺขณา อุปลกฺขณา ปจฺจุ\-ปลกฺขณา ปณฺฑิจฺจํ โกสลฺลํ เนปุญฺญํ เวภพฺยา จินฺตา อุปปริกฺขา ภูรี เมธา ปริณายิกา วิปสฺสนา สมฺปชญฺญํ ปโตโท ปญฺญา ปญฺญินฺทฺริยํ ปญฺญาพลํ ปญฺญาสตฺถํ ปญฺญาปาสาโท ปญฺญาอาโลโก ปญฺญาโอภาโส ปญฺญาปชฺโชโต ปญฺญารตนํ อโมโห ธมฺมวิจโย สมฺมาทิฏฺฐิ ธมฺมวิจย\-สมฺโพชฺฌงฺโค มคฺคงฺคํ มคฺค\-ปริยาปนฺนํ. อยํ วุจฺจติ อโมโห.}

\vspace{\tipitakaparskip}
\csromancenter{อิเม ตโย กุสลเหตู.}
\proseitem{1077}{\csromanfolio{218}ตตฺถ กตเม ตโย อพฺยากตเหตู. กุสลานํ ธมฺมานํ วิปากโต อโลโภ อโทโส อโมโห. อิเม ตโย อพฺยากตเหตู. อิเม ฉ อปริยาปนฺน\-เหตู. อิเม ธมฺมา เหตู.}

\proseitem{1078}{กตเม ธมฺมา น เหตู. เต ธมฺเม ฐเปตฺวา อวเสสา กุสลากุสลา\-พฺยากตา ธมฺมา กามาวจรา รูปาวจรา อรูปวจรา อปริยาปนฺนา เวทนากฺขนฺโธ ฯเปฯ วิญฺญาณ\-กฺขนฺโธ สพฺพญฺจ รูปํ อสงฺขตา จ ธาตุ. อิเม ธมฺมา น เหตู.}

\proseitem{1079}{กตเม ธมฺมา สเหตูกา. เตหิ ธมฺเมหิ เย ธมฺมา สเหตุกา เวทนากฺขนฺโธ ฯเปฯ วิญฺญาณ\-กฺขนฺโธ. อิเม ธมฺมา สเหตุกา.}

\proseitem{1080}{กตเม ธมฺมา อเหตุกา. เตหิ ธมฺเมหิ เย ธมฺมา อเหตุกา เวทนากฺขนฺโธ ฯเปฯ วิญฺญาณ\-กฺขนฺโธ สพฺพญฺจ รูปํ อสงฺขตา จ ธาตุ. อิเม ธมฺมา อเหตุกา.}

\proseitem{1081}{กตเม ธมฺมา เหตุสมฺปยุตฺตา. เตหิ ธมฺเมหิ เย ธมฺมา สมฺปยุตฺตา เวทนากฺขนฺโธ ฯเปฯ วิญฺญาณ\-กฺขนฺโธ. อิเม ธมฺมา เหตุสมฺปยุตฺตา.}

\proseitem{1082}{กตเม ธมฺมา เหตุวิปฺปยุตฺตา. เตหิ ธมฺเมหิ เย ธมฺมา วิปฺปยุตฺตา เวทนากฺขนฺโธ ฯเปฯ วิญฺญาณ\-กฺขนฺโธ สพฺพญฺจ รูปํ อสงฺขตา จ ธาตุ. อิเม ธมฺมา เหตุวิปฺปยุตฺตา.}

\proseitem{1083}{กตเม ธมฺมา เหตู เจว สเหตุกา จ. โลโภ โมเหน เหตุ เจว สเหตุโก จ, โมโห โลเภน เหตุ เจว สเหตุโก จ, โทโส โมเหน เหตุ เจว สเหตุโก จ, โมโห โทเสน เหตุ เจว สเหตุโก จ, อโลโภ อโทโส อโมโห, เต อญฺญมญฺญํ เหตู เจว สเหตุกา จ. อิเม ธมฺมา เหตู เจว สหตุกา จ.}

\proseitem{1084}{กตเม ธมฺมา สเหตุกา เจว น จ เหตู. เตหิ ธมฺเมหิ เย ธมฺมา สเหตุกา, เต ธมฺเม ฐเปตฺวา เวทนากฺขนฺโธ ฯเปฯ วิญฺญาณ\-กฺขนฺโธ. อิเม ธมฺมา สเหตุกา เจว น จ เหตู.}

\chapterheadpageread{3. นิกฺเขปกณฺฑ}
\proseitem{1085}{\csromanfolio{219}กตเม ธมฺมา เหตู เจว เหตุสมฺปยุตฺตา จ. โลโภ โมเหน เหตุ เจว เหตุสมฺปยุตฺโต จ, โมโห โลเภน เหตุ เจว เหตุสมฺปยุตฺโต จ, โทโส โมเหน เหตุ เจว เหตุสมฺปยุตฺโต จ, โมโห โทเสน เหตุ เจว เหตุสมฺปยุตฺโต จ, อโลโภ อโทโส อโมโห, เต อญฺญมญฺญํ เหตู เจว เหตุสมฺปยุตฺตา จ. อิเม ธมฺมา เหตู เจว เหตุสมฺปยุตฺตา จ.}

\proseitem{1086}{กตเม ธมฺมา เหตุสมฺปยุตฺตา เจว น จ เหตู. เตหิ ธมฺเมหิ เย ธมฺมา สมฺปยุตฺตา, เต ธมฺเม ฐเปตฺวา เวทนากฺขนฺโธ ฯเปฯ วิญฺญาณ\-กฺขนฺโธ. อิเม ธมฺมา เหตุสมฺปยุตฺตา เจว น จ เหตู.}

\proseitem{1087}{กตเม ธมฺมา น เหตู สเหตุกา. เตหิ ธมฺเมหิ เย ธมฺมา น เหตู สเหตุกา เวทนากฺขนฺโธ ฯเปฯ วิญฺญาณ\-กฺขนฺโธ. อิเม ธมฺมา น เหตู สเหตุกา.}

\proseitem{1088}{กตเม ธมฺมา น เหตู อเหตุกา. เตหิ ธมฺเมหิ เย ธมฺมา น เหตู อเหตุกา เวทนากฺขนฺโธ ฯเปฯ วิญฺญาณ\-กฺขนฺโธ สพฺพญฺจ รูปํ อสงฺขตา จ ธาตุ. อิเม ธมฺมา น เหตู อเหตุกา.}

\csromanmatikaanchor{mtk.105}
\csromantocmarkref{section}{จูฬนฺตรทุก}{mtk.105}
\bookmark[dest={mtk.105},level=1]{จูฬนฺตรทุก}
\csromanheader{1}{\textbf{จูฬนฺตรทุก}}
\proseitem{1089}{กตเม ธมฺมา สปฺปจฺจยา. ปญฺจกฺขนฺธา รูปกฺขนฺโธ เวทนากฺขนฺโธ สญฺญากฺขนฺโธ สงฺขาร\-กฺขนฺโธ วิญฺญาณ\-กฺขนฺโธ. อิเม ธมฺมา สปฺปจฺจยา.}

\proseitem{1090}{กตเม ธมฺมา อปฺปจฺจยา. อสงฺขตา ธาตุ. อิเม ธมฺมา อปฺปจฺจยา.}

\proseitem{1091}{กตเม ธมฺมา สงฺขตา. เยว เต ธมฺมา สปฺปจฺจยา, เตว เต ธมฺมา สงฺขตา.}

\proseitem{1092}{กตเม ธมฺมา อสงฺขตา. โย เอว โส ธมฺโม อปฺปจฺจโย, โส เอว โส ธมฺโม อสงฺขโต.}

\proseitem{1093}{กตเม ธมฺมา สนิทสฺสนา. รูปายตนํ. อิเม ธมฺมา สนิทสฺสนา.}

\proseitem{1094}{\csromanfolio{220}กตเม ธมฺมา อนิทสฺสนา. จกฺขายตนํ ฯเปฯ โผฏฺฐพฺพา\-ยตนํ เวทนากฺขนฺโธ ฯเปฯ วิญฺญาณ\-กฺขนฺโธ, ยญฺจ รูปํ อนิทสฺสนํ อปฺปฏิฆํ ธมฺมายตน\-ปริยาปนฺนํ อสงฺขตา จ ธาตุ. อิเม ธมฺมา อนิทสฺสนา.}

\proseitem{1095}{กตเม ธมฺมา สปฺปฏิฆา. จกฺขายตนํ ฯเปฯ โผฏฺฐพฺพา\-ยตนํ. อิเม ธมฺมา สปฺปฏิฆา.}

\proseitem{1096}{กตเม ธมฺมา อปฺปฏิฆา. เวทนากฺขนฺโธ ฯเปฯ วิญฺญาณ\-กฺขนฺโธ, ยญฺจ รูปํ อนิทสฺสนํ อปฺปฏิฆํ ธมฺมายตน\-ปริยาปนฺนํ อสงฺขตา จ ธาตุ. อิเม ธมฺมา อปฺปฏิฆา.}

\proseitem{1097}{กตเม ธมฺมา รูปิโน. จตฺตาโร จ มหาภูตา จตุนฺนญฺจ มหาภูตานํ อุปาทาย รูปํ. อิเม ธมฺมา รูปิโน.}

\proseitem{1098}{กตเม ธมฺมา อรูปิโน. เวทนากฺขนฺโธ ฯเปฯ วิญฺญาณ\-กฺขนฺโธ อสงฺขตา จ ธาตุ. อิเม ธมฺมา อรูปิโน.}

\proseitem{1099}{กตเม ธมฺมา โลกิยา. สาสวา กุสลากุสลา\-พฺยากตา ธมฺมา กามาวจรา รูปาวจรา อรูปาวจรา รูปกฺขนฺโธ ฯเปฯ วิญฺญาณ\-กฺขนฺโธ. อิเม ธมฺมา โลกิยา.}

\proseitem{1100}{กตเม ธมฺมา โลกุตฺตรา. อปริยาปนฺนา มคฺคา จ มคฺคผลานิ จ อสงฺขตา จ ธาตุ. อิเม ธมฺมา โลกุตฺตรา.}

\proseitem{1101}{กตเม ธมฺมา เกนจิ วิญฺเญยฺยา เกนจิ น วิญฺเญยฺยา. เย เต ธมฺมา จกฺขุวิญฺเญยฺยา น เต ธมฺมา โสตวิญฺเญยฺยา, เย วา ปน เต ธมฺมา โสตวิญฺเญยฺยา น เต ธมฺมา จกฺขุวิญฺเญยฺยา. เย เต ธมฺมา จกฺขุวิญฺเญยฺยา น เต ธมฺมา ฆานวิญฺเญยฺยา, เย วา ปน เต ธมฺมา ฆานวิญฺเญยฺยา น เต ธมฺมา จกฺขุวิญฺเญยฺยา. เย เต ธมฺมา จกฺขุวิญฺเญยฺยา น เต ธมฺมา ชิวฺหาวิญฺเญยฺยา, เย วา ปน เต ธมฺมา ชิวฺหาวิญฺเญยฺยา น เต ธมฺมา จกฺขุวิญฺเญยฺยา. เย เต ธมฺมา จกฺขุวิญฺเญยฺยา น เต ธมฺมา กายวิญฺเญยฺยา, เย วา ปน เต ธมฺมา กายวิญฺเญยฺยา น เต ธมฺมา จกฺขุวิญฺเญยฺยา. เย เต ธมฺมา โสตวิญฺเญยฺยา น เต ธมฺมา ฆานวิญฺเญยฺยา, เย วา ปน เต ธมฺมา ฆานวิญฺเญยฺยา น เต ธมฺมา โสตวิญฺเญยฺยา. เย เต ธมฺมา โสตวิญฺเญยฺยา น เต ธมฺมา ชิวฺหาวิญฺเญยฺยา, เย วา ปน เต \csromanfolio{221}ธมฺมา ชิวฺหาวิญฺเญยฺยา น เต ธมฺมา โสตวิญฺเญยฺยา. เย เต ธมฺมา โสตวิญฺเญยฺยา น เต ธมฺมา กายวิญฺเญยฺยา, เย วา ปน เต ธมฺมา กายวิญฺเญยฺยา น เต ธมฺมา โสตวิญฺเญยฺยา. เย เต ธมฺมา โสตวิญฺเญยฺยา น เต ธมฺมา จกฺขุวิญฺเญยฺยา, เย วา ปน เต ธมฺมา จกฺขุวิญฺเญยฺยา น เต ธมฺมา โสตวิญฺเญยฺยา. เย เต ธมฺมา ฆานวิญฺเญยฺยา น เต ธมฺมา ชิวฺหาวิญฺเญยฺยา, เย วา ปน เต ธมฺมา ชิวฺหาวิญฺเญยฺยา น เต ธมฺมา ฆานวิญฺเญยฺยา. เย เต ธมฺมา ฆานวิญฺเญยฺยา น เต ธมฺมา กายวิญฺเญยฺยา, เย วา ปน เต ธมฺมา กายวิญฺเญยฺยา น เต ธมฺมา ฆานวิญฺเญยฺยา. เย เต ธมฺมา ฆานวิญฺเญยฺยา น เต ธมฺมา จกฺขุวิญฺเญยฺยา, เย วา ปน เต ธมฺมา จกฺขุวิญฺเญยฺยา น เต ธมฺมา ฆานวิญฺเญยฺยา. เย เต ธมฺมา ฆานวิญฺเญยฺยา น เต ธมฺมา โสตวิญฺเญยฺยา, เย วา ปน เต ธมฺมา โสตวิญฺเญยฺยา น เต ธมฺมา ฆานวิญฺเญยฺยา. เย เต ธมฺมา ชิวฺหาวิญฺเญยฺยา น เต ธมฺมา กายวิญฺเญยฺยา, เย วา ปน เต ธมฺมา กายวิญฺเญยฺยา น เต ธมฺมา ชิวฺหาวิญฺเญยฺยา. เย เต ธมฺมา ชิวฺหาวิญฺเญยฺยา น เต ธมฺมา จกฺขุวิญฺเญยฺยา, เย วา ปน เต ธมฺมา จกฺขุวิญฺเญยฺยา น เต ธมฺมา ชิวฺหาวิญฺเญยฺยา. เย เต ธมฺมา ชิวฺหาวิญฺเญยฺยา น เต ธมฺมา โสตวิญฺเญยฺยา, เย วา ปน เต ธมฺมา โสตวิญฺเญยฺยา น เต ธมฺมา ชิวฺหาวิญฺเญยฺยา. เย เต ธมฺมา ชิวฺหาวิญฺเญยฺยา น เต ธมฺมา ฆานวิญฺเญยฺยา, เย วา ปน เต ธมฺมา ฆานวิญฺเญยฺยา น เต ธมฺมา ชิวฺหาวิญฺเญยฺยา. เย เต ธมฺมา กายวิญฺเญยฺยา น เต ธมฺมา จกฺขุวิญฺเญยฺยา, เย วา ปน เต ธมฺมา จกฺขุวิญฺเญยฺยา น เต ธมฺมา กายวิญฺเญยฺยา. เย เต ธมฺมา กายวิญฺเญยฺยา น เต ธมฺมา โสตวิญฺเญยฺยา, เย วา ปน เต ธมฺมา โสตวิญฺเญยฺยา น เต ธมฺมา กายวิญฺเญยฺยา. เย เต ธมฺมา กายวิญฺเญยฺยา น เต ธมฺมา ฆานวิญฺเญยฺยา, เย วา ปน เต ธมฺมา ฆานวิญฺเญยฺยา น เต ธมฺมา กายวิญฺเญยฺยา. เย เต ธมฺมา กายวิญฺเญยฺยา น เต ธมฺมา ชิวฺหาวิญฺเญยฺยา, เย วา ปน เต ธมฺมา ชิวฺหาวิญฺเญยฺยา น เต ธมฺมา กายวิญฺเญยฺยา. อิเม ธมฺมา เกนจิ วิญฺเญยฺยา เกนจิ น วิญฺเญยฺยา.}

\csromanmatikaanchor{mtk.106}
\csromantocmarkref{section}{อาสวโคจฺฉก}{mtk.106}
\bookmark[dest={mtk.106},level=1]{อาสวโคจฺฉก}
\csromanheader{1}{\textbf{อาสวโคจฺฉก}}
\proseitem{1102}{กตเม ธมฺมา อาสวา. จตฺตาโร อาสวา กามาสโว ภวาสโว ทิฏฺฐาสโว อวิชฺชาสโว.}

\proseitem{1103}{\csromanfolio{222}ตตฺถ กตโม กามาสโว. โย กาเมสุ กามจฺฉนฺโท กามราโค กามนนฺที กามตณฺหา กามสิเนโห กามปริฬาโห กามมุจฺฉา กามชฺโฌสานํ. อยํ วุจฺจติ กามาสโว.}

\proseitem{1104}{ตตฺถ กตโม ภวาสโว. โย ภเวสุ ภวฉนฺโท\csromansharedfootnote{c17ba5d596b5b391}{ภวจฺฉนฺโท (สี, สฺยา)} ภวราโค ภวนนฺที ภวตณฺหา ภวสิเนโห ภวปริฬาโห ภวมุจฺฉา ภวชฺโฌสานํ. อยํ วุจฺจติ ภวาสโว.}

\proseitem{1105}{ตตฺถ กตโม ทิฏฺฐาสโว. “สสฺสโต โลโก”ติ วา, “อสสฺสโต โลโก”ติ วา, “อนฺตวา โลโก”ติ วา, “อนนฺตวา โลโก”ติ วา, “ตํ ชีวํ ตํ สรีรนฺ”ติ วา, “อญฺญํ ชีวํ อญฺญํ สรีรนฺ”ติ วา, “โหติ ตถาคโต ปรํ มรณา”ติ วา, “น โหติ ตถาคโต ปรํ มรณา”ติ วา, “โหติ จ น จ โหติ ตถาคโต ปรํ มรณา”ติ วา, “เนว โหติ น น โหติ ตถาคโต ปรํ มรณา”ติ วา. ยา เอวรูปา ทิฏฺฐิ ทิฏฺฐิคตํ ทิฏฺฐิคหนํ ทิฏฺฐิกนฺตาโร ทิฏฺฐิ\-วิสูกายิกํ ทิฏฺฐิ\-วิปฺผนฺทิตํ ทิฏฺฐิ\-สญฺโญชนํ คาโห ปฏิคฺคาโห อภินิเวโส ปรามาโส กุมฺมคฺโค มิจฺฉาปโถ มิจฺฉตฺตํ ติตฺถายตนํ วิปริยาส\-คฺคาโห. อยํ วุจฺจติ ทิฏฺฐาสโว. สพฺพาปิ มิจฺฉาทิฏฺฐิ ทิฏฺฐาสโว.}

\proseitem{1106}{ตตฺถ กตโม อวิชฺชาสโว. ทุกฺเข อญฺญาณํ, ทุกฺขสมุทเย อญฺญาณํ, ทุกฺขนิโรเธ อญฺญาณํ, ทุกฺขนิโรธ\-คามินิยา ปฏิปทาย อญฺญาณํ, ปุพฺพนฺเต อญฺญาณํ, อปรนฺเต อญฺญาณํ, ปุพฺพนฺตา\-ปรนฺเต อญฺญาณํ, อิทปฺปจฺจยตา ปฏิจฺจ\-สมุปฺปนฺเนสุ ธมฺเมสุ อญฺญาณํ. ยํ เอวรูปํ อญฺญาณํ อทสฺสนํ อนภิสมโย อนนุโพโธ อสมฺโพโธ อปฺปฏิเวโธ อสํคาหนา อปริโยคาหนา อสมเปกฺขนา อปจฺจเวกฺขณา อปจฺจกฺขกมฺมํ ทุมฺเมชฺฌํ พาลฺยํ อสมฺปชญฺญํ โมโห ปโมโห สมฺโมโห อวิชฺชา อวิชฺโชโฆ อวิชฺชาโยโค อวิชฺชานุสโย อวิชฺชา\-ปริยุฏฺฐานํ อวิชฺชาลงฺคี โมโห อกุสลมูลํ. อยํ วุจฺจติ อวิชฺชาสโว. อิเม ธมฺมา อาสวา.}

\proseitem{1107}{กตเม ธมฺมา โน อาสวา. เต ธมฺเม ฐเปตฺวา อวเสสา กุสลากุสลา\-พฺยากตา ธมฺมา กามาวจรา รูปาวจรา อรูปวจรา \csromanfolio{223}อปริยาปนฺนา เวทนากฺขนฺโธ ฯเปฯ วิญฺญาณ\-กฺขนฺโธ สพฺพญฺจ รูปํ อสงฺขตา จ ธาตุ. อิเม ธมฺมา โน อาสวา.}

\proseitem{1108}{กตเม ธมฺมา สาสวา. กุสลากุสลา\-พฺยากตา ธมฺมา กามาวจรา รูปาวจรา อรูปาวจรา รูปกฺขนฺโธ ฯเปฯ วิญฺญาณ\-กฺขนฺโธ. อิเม ธมฺมา สาสวา.}

\proseitem{1109}{กตเม ธมฺมา อนาสวา. อปริยาปนฺนา มคฺคา จ มคฺคผลานิ จ อสงฺขตา จ ธาตุ. อิเม ธมฺมา อนาสวา.}

\proseitem{1110}{กตเม ธมฺมา อาสว\-สมฺปยุตฺตา. เตหิ ธมฺเมหิ เย ธมฺมา สมฺปยุตฺตา เวทนากฺขนฺโธ ฯเปฯ วิญฺญาณ\-กฺขนฺโธ. อิเม ธมฺมา อาสว\-สมฺปยุตฺตา.}

\proseitem{1111}{กตเม ธมฺมา อาสว\-วิปฺปยุตฺตา. เตหิ ธมฺเมหิ เย ธมฺมา วิปฺปยุตฺตา เวทนากฺขนฺโธ ฯเปฯ วิญฺญาณ\-กฺขนฺโธ สพฺพญฺจ รูปํ อสงฺขตา จ ธาตุ. อิเม ธมฺมา อาสว\-วิปฺปยุตฺตา.}

\proseitem{1112}{กตเม ธมฺมา อาสวา เจว สาสวา จ. เต เยว อาสวา อาสวา เจว สาสวา จ.}

\proseitem{1113}{กตเม ธมฺมา สาสวา เจว โน จ อาสวา. เตหิ ธมฺเมหิ เย ธมฺมา สาสวา, เต ธมฺเม ฐเปตฺวา อวเสสา สาสวา กุสลากุสลา\-พฺยากตา ธมฺมา กามาวจรา รูปาวจรา อรูปาวจรา รูปกฺขนฺโธ ฯเปฯ วิญฺญาณ\-กฺขนฺโธ. อิเม ธมฺมา สาสวา เจว โน จ อาสวา.}

\proseitem{1114}{กตเม ธมฺมา อาสวา เจว อาสว\-สมฺปยุตฺตา จ. กามาสโว อวิชฺชาสเวน อาสโว เจว อาสว\-สมฺปยุตฺโต จ, อวิชฺชาสโว กามาสเวน อาสโว เจว อาสว\-สมฺปยุตฺโต จ, ภวาสโว อวิชฺชาสเวน อาสโว เจว อาสว\-สมฺปยุตฺโต จ, อวิชฺชาสโว ภวาสเวน อาสโว เจว อาสว\-สมฺปยุตฺโต จ, ทิฏฺฐาสโว อวิชฺชาสเวน อาสโว เจว อาสว\-สมฺปยุตฺโต จ, อวิชฺชาสโว ทิฏฺฐาสเวน อาสโว เจว อาสว\-สมฺปยุตฺโต จ. อิเม ธมฺมา อาสวา เจว อาสว\-สมฺปยุตฺตา จ.}

\proseitem{1115}{\csromanfolio{224}กตเม ธมฺมา อาสว\-สมฺปยุตฺตา เจว โน จ อาสวา. เตหิ ธมฺเมหิ เย ธมฺมา สมฺปยุตฺตา, เต ธมฺเม ฐเปตฺวา เวทนากฺขนฺโธ ฯเปฯ วิญฺญาณ\-กฺขนฺโธ. อิเม ธมฺมา อาสว\-สมฺปยุตฺตา เจว โน จ อาสวา.}

\proseitem{1116}{กตเม ธมฺมา อาสว\-วิปฺปยุตฺตา สาสวา. เตหิ ธมฺเมหิ เย ธมฺมา วิปฺปยุตฺตา สาสวา กุสลากุสลา\-พฺยากตา ธมฺมา กามาวจรา รูปาวจรา อรูปาวจรา รูปกฺขนฺโธ ฯเปฯ วิญฺญาณ\-กฺขนฺโธ. อิเม ธมฺมา อาสว\-วิปฺปยุตฺตา สาสวา.}

\proseitem{1117}{กตเม ธมฺมา อาสว\-วิปฺปยุตฺตา อนาสวา. อปริยาปนฺนา มคฺคา จ มคฺคผลานิ จ อสงฺขตา จ ธาตุ. อิเม ธมฺมา อาสว\-วิปฺปยุตฺตา อนาสวา.}

\vspace{\tipitakaparskip}
\csromancenter{นิกฺเขปกณฺเฑ ปฐม\-ภาณวาโร.}
\csromanmatikaanchor{mtk.107}
\csromantocmarkref{section}{สญฺโญชนโคจฺฉก}{mtk.107}
\bookmark[dest={mtk.107},level=1]{สญฺโญชนโคจฺฉก}
\csromanheader{1}{สญฺโญชน\-โคจฺฉก}
\proseitem{1118}{กตเม ธมฺมา สญฺโญชานา. ทส สญฺโญชนานิ กามราค\-สญฺโญชนํ ปฏิฆ\-สญฺโญชนํ มานสญฺโญชนํ ทิฏฺฐิ\-สญฺโญชนํ วิจิกิจฺฉา\-สญฺโญชนํ สีลพฺพตปรามาส\-สญฺโญชนํ ภวราค\-สญฺโญชนํ อิสฺสาสญฺโญชนํ มจฺฉริย\-สญฺโญชนํ อวิชฺชา\-สญฺโญชนํ.}

\proseitem{1119}{ตตฺถ กตมํ กามราค\-สญฺโญชนํ. โย กาเมสุ กามจฺฉนฺโท กามราโค กามนนฺที กามตณฺหา กามสิเนโห กามปริฬาโห กามมุจฺฉา กามชฺโฌสานํ. อิทํ วุจฺจติ กามราค\-สญฺโญชนํ.}

\proseitem{1120}{ตตฺถ กตมํ ปฏิฆ\-สญฺโญชนํ. “อนตฺถํ เม อจรี”ติ อาฆาโต ชายติ, “อนตฺถํ เม จรตี”ติ อาฆาโต ชายติ, “อนตฺถํ เม จริสฺสตี”ติ อาฆาโต ชายติ, “ปิยสฺส เม มนาปสฺส อนตฺถํ อจริ ฯเปฯ อนตฺถํ จรติ ฯเปฯ อนตฺถํ จริสฺสตี”ติ อาฆาโต ชายติ, “อปฺปิยสฺส เม อมนาปสฺส อตฺถํ อจริ ฯเปฯ อตฺถํ จรติ ฯเปฯ อตฺถํ จริสฺสตี”ติ \csromanfolio{225}อาฆาโต ชายติ, อฏฺฐาเน วา ปน อาฆาโต ชายติ. โย เอวรูโป จิตฺตสฺส อาฆาโต ปฏิฆาโต ปฏิฆํ ปฏิวิโรโธ โกโป ปโกโป สมฺปโกโป โทโส ปโทโส สมฺปโทโส จิตฺตสฺส พฺยาปตฺติ มโนปโทโส โกโท กุชฺฌนา กุชฺฌิตตฺตํ โทโส ทุสฺสนา ทุสฺสิตตฺตํ พฺยาปตฺติ พฺยาปชฺชนา พฺยาปชฺชิตตฺตํ วิโรโธ ปฏิวิโรโธ จณฺฑิกฺกํ อสุโรโป อนตฺตมนตา จิตฺตสฺส. อิทํ วุจฺจติ ปฏิฆ\-สญฺโญชนํ.}

\proseitem{1121}{ตตฺถ กตมํ มานสญฺโญชนํ. “เสยฺโยหมสฺมี”ติ มาโน, “สทิโสหมสฺมี”ติ มาโน, “หีโนหมสฺมี”ติ มาโน. โย เอวรูโป มาโน มญฺญนา มญฺญิตตฺตํ อุนฺนติ อุนฺนโม\csromansharedfootnote{00810625007bb0d1}{อุณฺณติ อุณฺณาโม (สฺยา)} ธโช สมฺปคฺคาโห เกตุกมฺยตา จิตฺตสฺส. อิทํ วุจฺจติ มานสญฺโญชนํ.}

\proseitem{1122}{ตตฺถ กตมํ ทิฏฺฐิ\-สญฺโญชนํ. “สสฺสโต โลโก”ติ วา, “อสสฺสโต โลโก”ติ วา, “อนฺตวา โลโก”ติ วา, “อนนฺตวา โลโก”ติ วา, “ตํ ชีวํ ตํ สรีรนฺ”ติ วา, “อญฺญํ ชีวํ อญฺญํ สรีรนฺ”ติ วา, “โหติ ตถาคโต ปรํ มรณา”ติ วา, “น โหติ ตถาคโต ปรํ มรณา”ติ วา, “โหติ จ น จ โหติ ตถาคโต ปรํ มรณา”ติ วา, “เนว โหติ น น โหติ ตถาคโต ปรํ มรณา”ติ วา. ยา เอวรูปา ทิฏฺฐิ ทิฏฺฐิคตํ ทิฏฺฐิคหนํ ทิฏฺฐิกนฺตาโร ทิฏฺฐิ\-วิสูกายิกํ ทิฏฺฐิ\-วิปฺผนฺทิตํ ทิฏฺฐิ\-สญฺโญชนํ คาโห ปฏิคฺคาโห อภินิเวโส ปรามาโส กุมฺมคฺโค มิจฺฉาปโถ มิจฺฉตฺตํ ติตฺถายตนํ วิปริยาส\-คฺคาโห. อิทํ วุจฺจติ ทิฏฺฐิ\-สญฺโญชนํ. ฐเปตฺวา สีลพฺพตปรามาส\-สญฺโญชนํ สพฺพาปิ มิจฺฉาทิฏฺฐิ ทิฏฺฐิ\-สญฺโญชนํ.}

\proseitem{1123}{ตตฺถา กตมํ วิจิกิจฺฉา\-สญฺโญชนํ. สตฺถริ กงฺขติ วิจิกิจฺฉติ, ธมฺเม กงฺขติ วิจิกิจฺฉติ, สํเฆ กงฺขติ วิจิกิจฺฉติ, สิกฺขาย กงฺขติ วิจิกิจฺฉาติ, ปุพฺพนฺเต กงฺขติ วิจิกิจฺฉติ, อปรนฺเต กงฺขติ วิจิกิจฺฉติ, ปุพฺพนฺตา\-ปรนฺเต กงฺขติ วิจิกิจฺฉาติ, อิทปฺปจฺจยตา ปฏิจฺจ\-สมุปฺปนฺเนสุ ธมฺเมสุ กงฺขติ วิจิกิจฺฉติ. ยา เอวรูปา กงฺขา กงฺขายนา กงฺขายิตตฺตํ วิมติ วิจิกิจฺฉา เทฺวฬฺหกํ เทฺวธาปโถ สํสโย อเนกํสคฺคาโห อาสปฺปนา ปริสปฺปนา อปริโยคาหนา ถมฺภิตตฺตํ จิตฺตสฺส มโนวิเลโข. อิทํ วุจฺจติ วิจิกิจฺฉา\-สญฺโญชนํ.}

\proseitem{1124}{\csromanfolio{226}ตตฺถ กตมํ สีลพฺพตปรามาส\-สญฺโญชนํ. “อิโต พหิทฺธา สมณ\-พฺราหฺมณานํ สีเลน สุทฺธิ วเตน สุทฺธิ สีลพฺพเตน สุทฺธี”ติ ยา เอวรูปา ทิฏฺฐิ ทิฏฺฐิคตํ ทิฏฺฐิคหนํ ทิฏฺฐิกนฺตาโร ทิฏฺฐิ\-วิสูกายิกํ ทิฏฺฐิ\-วิปฺผนฺทิตํ ทิฏฺฐิ\-สญฺโญชนํ คาโห ปติฏฺฐาโห อภินิเวโส ปรามาโส กุมฺมคฺเค มิจฺฉาปโถ มิจฺฉตฺตํ ติตฺถายตนํ วิปริยาส\-คฺคาโห. อิทํ วุจฺจติ สีลพฺพตปรามาส\-สญฺโญชนํ.}

\proseitem{1125}{ตตฺถ กตมํ ภวราค\-สญฺโญชนํ. โย ภเวสุ ภวฉนฺโท ภวราโค ภวนนฺที ภวตณฺหา ภวสิเนโห ภวปริฬาโห ภวมุจฺฉา ภวชฺโฌสานํ. อิทํ วุจฺจติ ภวราค\-สญฺโญชนํ.}

\proseitem{1126}{ตตฺถ กตมํ อิสฺสาสญฺโญชนํ. ยา ปรลาภ\-สกฺการ\-ครุการ\-มานน\-วนฺทน\-ปูชนาสุ อิสฺสา อิสฺสายนา อิสฺสายิตตฺตํ อุสูยา อุสูยนา อุสูยิตตฺตํ\csromansharedfootnote{d1a1d0afd746682e}{อุสฺสุยา อุสฺสุยนา อุสฺสุยิตตฺตํ (ก)}. อิทํ วุจฺจติ อิสฺสาสญฺโญชนํ.}

\proseitem{1127}{ตตฺถ กตมํ มจฺฉริย\-สญฺโญชนํ. ปญฺจ มจฺฉริยานิ อาวาส มจฺฉริยํ กุลมจฺฉริยํ ลาภ\-มจฺฉริยํ วณฺณ\-มจฺฉริยํ ธมฺม\-มจฺฉริยํ. ยํ เอวรูปํ มจฺเฉรํ มจฺฉรายนา มจฺฉรายิ\-ตตฺตํ เววิจฺฉํ กทริยํ กฏุกญฺจุกตา อคฺคหิตตฺตํ จิตฺตสฺส. อิทํ วุจฺจติ มจฺฉริย\-สญฺโญชนํ.}

\proseitem{1128}{ตตฺถ กตมํ อวิชฺชา\-สญฺโญชนํ. ทุกฺเข อญฺญาณํ, ทุกฺขสมุทเย อญฺญาณํ, ทุกฺขนิโรเธ อญฺญาณํ, ทุกฺขนิโรธ\-คามินิยา ปฏิปทาย อญฺญาณํ, ปุพฺพนฺเต อญฺญาณํ, อปรนฺเต อญฺญาณํ, ปุพฺพนฺตา\-ปรนฺเต อญฺญาณํ, อิทปฺปจฺจยตา ปฏิจฺจ\-สมุปฺปนฺเนสุ ธมฺเมสุ อญฺญาณํ. ยํ เอวรูปํ อญฺญาณํ อทสฺสนํ อนภิสมโย อนนุโพโธ อสมฺโพโธ อปฺปฏิเวโธ อสํคาหนา อปริโยคาหนา อสมเปกฺขนา อปจฺจเวกฺขณา อปจฺจกฺขกมฺมํ ทุมฺเมชฺฌํ พาลฺยํ อสมฺปชญฺญํ โมโห ปโมโห สมฺโมโห อวิชฺชา อวิชฺโชโฆ อวิชฺชาโยโค อวิชฺชานุสโย อวิชฺชา\-ปริยุฏฺฐานํ อวิชฺชาลงฺคี โมโห อกุสลมูลํ. อิทํ วุจฺจติ อวิชฺชา\-สญฺโญชนํ.}

\vspace{\tipitakaparskip}
\csromancenter{อิเม ธมฺมา สญฺโญชนา.}
\chapterheadpageread{3. นิกฺเขปกณฺฑ}
\proseitem{1129}{\csromanfolio{227}กตเม ธมฺมา โน สญฺโญชนา. เต ธมฺเม ฐเปตฺวา อวเสสา กุสลากุสลา\-พฺยากตา ธมฺมา กามาวจรา รูปาวจรา อรูปาวจรา อปริยาปนฺนา เวทนากฺขนฺโธ ฯเปฯ วิญฺญาณ\-กฺขนฺโธ สพฺพญฺจ รูปํ อสงฺขตา จ ธาตุ. อิเม ธมฺมา โน สญฺโญชนา.}

\proseitem{1130}{กตเม ธมฺมา สญฺโญชนิยา. สาสวา กุสลากุสลา\-พฺยากตา ธมฺมา กามาวจรา รูปาวจรา อรูปาวจรา รูปกฺขนฺโธ ฯเปฯ วิญฺญาณ\-กฺขนฺโธ. อิเม ธมฺมา สญฺโญชนิยา.}

\proseitem{1131}{กตเม ธมฺมา อสญฺโญชนิยา. อปริยาปนฺนา มคฺคา จ มคฺคผลานิ จ อสงฺขตา จ ธาตุ. อิเม ธมฺมา อสญฺโญชนิยา.}

\proseitem{1132}{กตเม ธมฺมา สญฺโญชน\-สมฺปยุตฺตา. เตหิ ธมฺเมหิ เย ธมฺมา สมฺปยุตฺตา เวทนากฺขนฺโธ ฯเปฯ วิญฺญาณ\-กฺขนฺโธ. อิเม ธมฺมา สญฺโญชน\-สมฺปยุตฺตา.}

\proseitem{1133}{กตเม ธมฺมา สญฺโญชน\-วิปฺปยุตฺตา. เตหิ ธมฺเมหิ เย ธมฺมา วิปฺปยุตฺตา เวทนากฺขนฺโธ ฯเปฯ วิญฺญาณ\-กฺขนฺโธ สพฺพญฺจ รูปํ อสงฺขตา จ ธาตุ. อิเม ธมฺมา สญฺโญชน\-วิปฺปยุตฺตา.}

\proseitem{1134}{กตเม ธมฺมา สญฺโญชนา เจว สญฺโญชนิยา จ. ตาเนว สญฺโญชนานิ สญฺโญชนา เจว สญฺโญชนิยา จ.}

\proseitem{1135}{กตเม ธมฺมา สญฺโญชนิยา เจว โน จ สญฺโญชนา. เตหิ ธมฺเมหิ เย ธมฺมา สญฺโญชนิยา, เต ธมฺเม ฐเปตฺวา อวเสสา สาสวา กุสลากุสลา\-พฺยากตา ธมฺมา กามาวจรา รูปาวจรา อรูปาวจรา รูปกฺขนฺโธ ฯเปฯ วิญฺญาณ\-กฺขนฺโธ. อิเม ธมฺมา สญฺโญชนิยา เจว โน จ สญฺโญชนา.}

\proseitem{1136}{กตเม ธมฺมา สญฺโญชนา เจว สญฺโญชน\-สมฺปยุตฺตา จ. กามราค\-สญฺโญชนํ อวิชฺชา\-สญฺโญชเนน สญฺโญชนญฺเจว สญฺโญชน\-สมฺปยุตฺตญฺจ, อวิชฺชา\-สญฺโญชนํ กามราค\-สญฺโญชเนน สญฺโญชนญฺเจว สญฺโญชน\-สมฺปยุตฺตญฺจ, ปฏิฆ\-สญฺโญชนํ อวิชฺชา\-สญฺโญชเนน สญฺโญชนญฺเจว สญฺโญชน\-สมฺปยุตฺตญฺจ, อวิชฺชา\-สญฺโญชนํ ปฏิฆ\-สญฺโญชเนน \csromanfolio{228}สญฺโญชนญฺเจว สญฺโญชน\-สมฺปยุตฺตญฺจ, มานสญฺโญชนํ อวิชฺชา\-สญฺโญชเนน สญฺโญชนญฺเจว สญฺโญชน\-สมฺปยุตฺตญฺจ, อวิชฺชา\-สญฺโญชนํ มาน\-สญฺโญชเนน สญฺโญชนญฺเจว สญฺโญชน\-สมฺปยุตฺตญฺจ, ทิฏฺฐิ\-สญฺโญชนํ อวิชฺชา\-สญฺโญชเนน สญฺโญชนญฺเจว สญฺโญชน\-สมฺปยุตฺตญฺจ, อวิชฺชา\-สญฺโญชนํ ทิฏฺฐิ\-สญฺโญชเนน สญฺโญชนญฺเจว สญฺโญชน\-สมฺปยุตฺตญฺจ, วิจิกิจฺฉา\-สญฺโญชนํ อวิชฺชา\-สญฺโญชเนน สญฺโญชนญฺเจว สญฺโญชน\-สมฺปยุตฺตญฺจ, อวิชฺชา\-สญฺโญชนํ วิจิกิจฺฉา\-สญฺโญชเนน สญฺโญชนญฺเจว สญฺโญชน\-สมฺปยุตฺตญฺจ, สีลพฺพตปรามาส\-สญฺโญชนํ อวิชฺชา\-สญฺโญชเนน สญฺโญชนญฺเจว สญฺโญชน\-สมฺปยุตฺตญฺจ, อวิชฺชา\-สญฺโญชนํ สีลพฺพตปรามาส\-สญฺโญชเนน สญฺโญชนญฺเจว สญฺโญชน\-สมฺปยุตฺตญฺจ, ภวราค\-สญฺโญชนํ อวิชฺชา\-สญฺโญชเนน สญฺโญชนญฺเจว สญฺโญชน\-สมฺปยุตฺตญฺจ, อวิชฺชา\-สญฺโญชนํ ภวราค\-สญฺโญชเนน สญฺโญชนญฺเจว สญฺโญชน\-สมฺปยุตฺตญฺจ, อิสฺสาสญฺโญชนํ อวิชฺชา\-สญฺโญชเนน สญฺโญชนญฺเจว สญฺโญชน\-สมฺปยุตฺตญฺจ, อวิชฺชา\-สญฺโญชนํ อิสฺสา\-สญฺโญชเนน สญฺโญชนญฺเจว สญฺโญชน\-สมฺปยุตฺตญฺจ, มจฺฉริย\-สญฺโญชนํ อวิชฺชา\-สญฺโญชเนน สญฺโญชนญฺเจว สญฺโญชน\-สมฺปยุตฺตญฺจ, อวิชฺชา\-สญฺโญชนํ มจฺฉริย\-สญฺโญชเนน สญฺโญชนญฺเจว สญฺโญชน\-สมฺปยุตฺตญฺจ. อิเม ธมฺมา สญฺโญชนา เจว สญฺโญชน\-สมฺปยุตฺตา จ.}

\proseitem{1137}{กตเม ธมฺมา สญฺโญชน\-สมฺปยุตฺตา เจว โน จ สญฺโญชนา. เตหิ ธมฺเมหิ เย ธมฺมา สมฺปยุตฺตา, เต ธมฺเม ฐเปตฺวา เวทนากฺขนฺโธ ฯเปฯ วิญฺญาณ\-กฺขนฺโธ. อิเม ธมฺมา สญฺโญชน\-สมฺปยุตฺตา เจว โน จ สญฺโญชนา.}

\proseitem{1138}{กตเม ธมฺมา สญฺโญชน\-วิปฺปยุตฺตา สญฺโญชนิยา. เตหิ ธมฺเมหิ เย ธมฺมา วิปฺปยุตฺตา สาสวา กุสลากุสลา\-พฺยากตา ธมฺมา กามาวจรา รูปาวจรา อรูปาวจรา รูปกฺขนฺโธ ฯเปฯ วิญฺญาณ\-กฺขนฺโธ. อิเม ธมฺมา สญฺโญชน\-วิปฺปยุตฺตา สญฺโญชนิยา.}

\proseitem{1139}{กตเม ธมฺมา สญฺโญชน\-วิปฺปยุตฺตา อสญฺโญชนิยา. อปริยาปนฺนา มคฺคา จ มคฺคผลานิ จ อสงฺขตา จ ธาตุ. อิเม ธมฺมา สญฺโญชน\-วิปฺปยุตฺตา อสญฺโญชนิยา.}

\csromanmatikaanchor{mtk.108}
\csromantocmarkref{section}{คนฺถโคจฺฉก}{mtk.108}
\bookmark[dest={mtk.108},level=1]{คนฺถโคจฺฉก}
\csromanheader{1}{3. นิกฺเขปกณฺฑ \textbf{คนฺถโคจฺฉก}}
\proseitem{1140}{\csromanfolio{229}กตเม ธมฺมา คนฺถา. จตฺตาโร คนฺถา อภิชฺฌา\-กายคนฺโถ พฺยาปาโท กายคนฺโถ สีลพฺพต\-ปรามาโส กายคนฺโถ อิทํสจฺจา\-ภินิเวโส กายคนฺโถ.}

\proseitem{1141}{ตตฺถ กตโม อภิชฺฌา\-กายคนฺโถ. โย ราโค สาราโค อนุนโย อนุโรโธ นนฺที นนฺทีราโค จิตฺตสฺส สาราโค อิจฺฉา มุจฺฉา อชฺโฌสานํ เคโธ ปลิเคโธ สงฺโค ปงฺโก เอชา มายา ชนิกา สญฺชนนี สิพฺพินี ชาลินี สริตา วิสตฺติกา สุตฺตํ วิสฏา อายูหินี ทุติยา ปณิธิ ภวเนตฺติ วนํ วนโถ สนฺถโว สิเนโห อเปกฺขา ปฏิพนฺธุ อาสา อาสิสนา อาสิสิตตฺตํ รูปาสา สทฺทาสา คนฺธาสา รสาสา โผฏฺฐพฺพาสา ลาภาสา ธนาสา ปุตฺตาสา ชีวิตาสา ชปฺปา ปชปฺปา อภิชปฺปา ชปฺปา ชปฺปนา ชปฺปิตตฺตํ โลลุปฺปํ โลลุปฺปายนา โลลุปฺปา\-ยิตตฺตํ ปุจฺฉญฺชิกตา สาธุกมฺยตา อธมฺมราโค วิสมโลโภ นิกนฺติ นิกามนา ปตฺถนา ปิหนา สมฺปตฺถนา กามตณฺหา ภวตณฺหา วิภวตณฺหา รูปตณฺหา อรูปตณฺหา นิโรธตณฺหา รูปตณฺหา สทฺทตณฺหา คนฺธตณฺหา รสตณฺหา โผฏฺฐพฺพ\-ตณฺหา ธมฺมตณฺหา โอโฆ โยโค คนฺโถ อุปาทานํ อาวรณํ นีวรณํ ฉาทนํ พนฺธนํ อุปกฺกิเลโส อนุสโย ปริยุฏฺฐานํ ลตา เววิจฺฉํ ทุกฺขมูลํ ทุกฺขนิทานํ ทุกฺขปฺปภโว มารปาโส มารพฬิสํ มารวิสโย ตณฺหานที ตณฺหาชาลํ ตณฺหาคทฺทุลํ ตณฺหาสมุทฺโท อภิชฺฌา โลโภ อกุสลมูลํ. อยํ วุจฺจติ อภิชฺฌา\-กายคนฺโถ.}

\proseitem{1142}{ตตฺถ กตโม พฺยาปาโท กายคนฺโถ. “อนตฺถํ เม อจรี”ติ อาฆาโต ชายติ, “อนตฺถํ เม จรตี”ติ อาฆาโต ชายติ, “อนตฺถํ เม จริสฺสตี”ติ อาฆาโต ชายติ, “ปิยสฺส เม มนาปสฺส อนตฺถํ อจริ ฯเปฯ อนตฺถํ จรติ ฯเปฯ อนตฺถํ จริสฺสตี”ติ อาฆาโต ชายติ, “อปฺปิยสฺส เม อมนาปสฺส อตฺถํ อจริ ฯเปฯ อตฺถํ จรติ ฯเปฯ อตฺถํ จริสฺสตี”ติ อาฆาโต ชายติ, อฏฺฐาเน วา ปน อาฆาโต ชายติ. โย เอวรูโป จิตฺตสฺส อาฆาโต ปฏิฆาโต ปฏิฆํ ปฏิวิโรโธ โกโป ปโกโป สมฺปโกโป โทโส ปโทโส สมฺปโทโส จิตฺตสฺส พฺยาปตฺติ มโนปโทโส โกโธ กุชฺฌนา กุชฺฌิตตฺตํ โทโส \csromanfolio{230}ทุสฺสนา ทุสฺสิตตฺตํ พฺยาปตฺติ พฺยาปชฺชนา พฺยาปชฺชิตตฺตํ วิโรโธ ปฏิวิโรโธ จณฺฑิกฺกํ อสุโรโป อนตฺตมนตา จิตฺตสฺส. อยํ วุจฺจติ พฺยาปาโท กายคนฺโถ.}

\proseitem{1143}{ตตฺถ กตโม สีลพฺพต\-ปรามาโส กายคนฺโถ. “อิโต พหิทฺธา สมณ\-พฺราหฺมณานํ สีเลน สุทฺธิ วเตน สุทฺธิ สีลพฺพเตน สุทฺธี”ติ ยา เอวรูปา ทิฏฺฐิ ทิฏฺฐิคตํ ทิฏฺฐิคหนํ ทิฏฺฐิกนฺตาโร ทิฏฺฐิ\-วิสูกายิกํ ทิฏฺฐิ\-วิปฺผนฺทิตํ ทิฏฺฐิ\-สญฺโญชนํ คาโห ปติฏฺฐาโห อภินิเวโส ปรามาโส กุมฺมคฺโค มิจฺฉาปโถ มิจฺฉตฺตํ หิตฺถายตนํ วิปริยาส\-คฺคาโห. อยํ วุจฺจติ สีลพฺพต\-ปรามาโส กายคนฺโถ.}

\proseitem{1144}{ตตฺถ กตโม อิทํสจฺจา\-ภินิเวโส กายคนฺโถ. “สสฺสโต โลโก อิทเมว สจฺจํ โมฆมญฺญนฺ”ติ วา, “อสสฺสโต โลโก อิทเมว สจฺจํ โมฆมญฺญนฺ”ติ วา, “อนฺตวา โลโก อิทเมว สจฺจํ โมฆมญฺญนฺ”ติ วา, “อนนฺตวา โลโก อิทเมว สจฺจํ โมฆมญฺญนฺ”ติ วา, “ตํ ชีวํ ตํ สรีรํ อิทเมว สจฺจํ โมฆมญฺญนฺ”ติ วา, “อญฺญํ ชีวํ อญฺญํ สรีรํ อิทเมว สจฺจํ โมฆมญฺญนฺ”ติ วา, “โหติ ตถาคโต ปรํ มรณา อิทเมว สจฺจํ โมฆมญฺญนฺ”ติ วา, “น โหติ ตถาคโต ปรํ มรณา อิทเมว สจฺจํ โมฆมญฺญนฺ”ติ วา, “โหติ จ น จ โหติ ตถาคโต ปรํ มรณา อิทเมว สจฺจํ โมฆมญฺญนฺ”ติ วา, “เนว โหติ น น โหติ ตถาคโต ปรํ มรณา อิทเมว สจฺจํ โมฆมญฺญนฺ”ติ วา. ยา เอวรูปา ทิฏฺฐิ ทิฏฺฐิคตํ ทิฏฺฐิคหนํ ทิฏฺฐิกนฺตาโร ทิฏฺฐิ\-วิสูกายิกํ ทิฏฺฐิ\-วิปฺผนฺทิตํ ทิฏฺฐิ\-สญฺโญชนํ คาโห ปติฏฺฐาโห อภินิเวโส ปรามาโส กุมฺมคฺโค มิจฺฉาปโถ มิจฺฉตฺตํ ติตฺถายตนํ วิปริยาส\-คฺคาโห. อยํ วุจฺจติ อิทํสจฺจา\-ภินิเวโส กายคนฺโถ. ฐเปตฺวา สีลพฺพต\-ปรามาสํ กายคนฺถํ สพฺพาปิ มิจฺฉาทิฏฺฐิ อิทํสจฺจา\-ภินิเวโส กายคนฺโถ. อิเม ธมฺมา คนฺถา.}

\proseitem{1145}{กตเม ธมฺมา โน คนฺถา. เต ธมฺเม ฐเปตฺวา อวเสสา กุสลากุสลา\-พฺยากตา ธมฺมา กามาวจรา รูปาวจรา อรูปาวจรา อปริยาปนฺนา เวทนากฺขนฺโธ ฯเปฯ วิญฺญาณ\-กฺขนฺโธ สพฺพญฺจ รูปํ อสงฺขตา จ ธาตุ. อิเม ธมฺมา โน คนฺถา.}

\chapterheadpageread{3. นิกฺเขปกณฺฑ}
\proseitem{1146}{\csromanfolio{231}กตเม ธมฺมา คนฺถนิยา. สาสวา กุสลากุสลา\-พฺยากตา ธมฺมา กามาวจรา รูปาวจรา อรูปาวจรา รูปกฺขนฺโธ ฯเปฯ วิญฺญาณ\-กฺขนฺโธ. อิเม ธมฺมา คนฺถนิยา.}

\proseitem{1147}{กตเม ธมฺมา อคนฺถนิยา. อปริยาปนฺนา มคฺคา จ มคฺคผลานิ จ อสงฺขตา จ ธาตุ. อิเม ธมฺมา อคนฺถนิยา.}

\proseitem{1148}{กตเม ธมฺมา คนฺถ\-สมฺปยุตฺตา. เตหิ ธมฺเมหิ เย ธมฺมา สมฺปยุตฺตา เวทนากฺขนฺโธ ฯเปฯ วิญฺญาณ\-กฺขนฺโธ. อิเม ธมฺมา คนฺธสมฺปยุตฺตา.}

\proseitem{1149}{กตเม ธมฺมา คนฺถ\-วิปฺปยุตฺตา. เตหิ ธมฺเมหิ เย ธมฺมา วิปฺปยุตฺตา เวทนากฺขนฺโธ ฯเปฯ วิญฺญาณ\-กฺขนฺโธ สพฺพญฺจ รูปํ อสงฺขตา จ ธาตุ. อิเม ธมฺมา คนฺถ\-วิปฺปยุตฺตา.}

\proseitem{1150}{กตเม ธมฺมา คนฺถา เจว คนฺถนิยา จ. เตว คนฺถา คนฺถา เจว คนฺถนิยา จ.}

\proseitem{1151}{กตเม ธมฺมา คนฺถนิยา เจว โน จ คนฺถา. เตหิ ธมฺเมหิ เย ธมฺมา คนฺถนิยา, เต ธมฺเม ฐเปตฺวา อวเสสา สาสวา กุสลากุสลา\-พฺยากตา ธมฺมา กามาวจรา รูปาวจรา อรูปาวจรา รูปกฺขนฺโธ ฯเปฯ วิญฺญาณ\-กฺขนฺโธ. อิเม ธมฺมา คนฺถนิยา เจว โน จ คนฺถา.}

\proseitem{1152}{กตเม ธมฺมา คนฺถา เจว คนฺถ\-สมฺปยุตฺตา จ. สีลพฺพต\-ปรามาโส กายคนฺโถ อภิชฺฌา\-กายคนฺเถน คนฺโถ เจว คนฺถ\-สมฺปยุตฺโต จ, อภิชฺฌา\-กายคนฺโถ สีลพฺพต\-ปรามาเสน กายคนฺเถน คนฺโถ เจว คนฺถ\-สมฺปยุตฺโต จ, อิทํสจฺจา\-ภินิเวโส กายคนฺโถ อภิชฺฌา\-กายคนฺเถน คนฺโถ เจว คนฺถ\-สมฺปยุตฺโต จ, อภิชฺฌา\-กายคนฺโถ อิทํสจฺจา\-ภินิเวเสน กายคนฺเถน คนฺโถ เจว คนฺถ\-สมฺปยุตฺโต จ. อิเม ธมฺมา คนฺถา เจว คนฺถ\-สมฺปยุตฺตา จ.}

\proseitem{1153}{กตเม ธมฺมา คนฺถ\-สมฺปยุตฺตา เจว โน จ คนฺถา. เตหิ ธมฺเมหิ เย ธมฺมา สมฺปยุตฺตา, เต ธมฺเม ฐเปตฺวา เวทนากฺขนฺโธ ฯเปฯ วิญฺญาณ\-กฺขนฺโธ. อิเม ธมฺมา คนฺถ\-สมฺปยุตฺตา เจว โน จ คนฺถา.}

\proseitem{1154}{\csromanfolio{232}กตเม ธมฺมา คนฺถ\-วิปฺปยุตฺตา คนฺถนิยา. เตหิ ธมฺเมหิ เย ธมฺมา วิปฺปยุตฺตา สาสวา กุสลากุสลา\-พฺยากตา ธมฺมา กามาวจรา รูปาวจรา อรูปาวจรา รูปกฺขนฺโธ ฯเปฯ วิญฺญาณ\-กฺขนฺโธ. อิเม ธมฺมา คนฺถ\-วิปฺปยุตฺตา คนฺถนิยา.}

\proseitem{1155}{กตเม ธมฺมา คนฺถ\-วิปฺปยุตฺตา อคนฺถนิยา. อปริยาปนฺนา มคฺคา จ มคฺคาผลานิ จ อสงฺขตา จ ธาตุ. อิเม ธมฺมา คนฺถ\-วิปฺปยุตฺตา อคนฺถนิยา.}

\csromanmatikaanchor{mtk.109}
\csromantocmarkref{section}{โอฆโคจฺฉก}{mtk.109}
\bookmark[dest={mtk.109},level=1]{โอฆโคจฺฉก}
\csromanheader{1}{\textbf{โอฆโคจฺฉก}}
\vspace{\tipitakaparskip}
\csromancenter{กตเม ธมฺมา โอฆา ฯเปฯ.}
\csromanmatikaanchor{mtk.110}
\csromantocmarkref{section}{โยคโคจฺฉก}{mtk.110}
\bookmark[dest={mtk.110},level=1]{โยคโคจฺฉก}
\csromanheader{1}{\textbf{โยคโคจฺฉก}}
\vspace{\tipitakaparskip}
\csromancenter{กตเม ธมฺมา โยคา ฯเปฯ.}
\csromanmatikaanchor{mtk.111}
\csromantocmarkref{section}{นีวรณโคจฺฉก}{mtk.111}
\bookmark[dest={mtk.111},level=1]{นีวรณโคจฺฉก}
\csromanheader{1}{นีวรณ\-โคจฺฉก}
\proseitem{1158}{กตเม ธมฺมา นีวรณา. ฉ นีวรณา\csromansharedfootnote{b5865fe4ac202808}{นีวรณานิ (สฺยา)} กามจฺฉนฺท\-นีวรณํ พฺยาปาท\-นีวรณํ ถินมิทฺธ\-นีวรณํ อุทฺธจฺจ\-กุกฺกุจฺจ\-นีวรณํ วิจิกิจฺฉา\-นีวรณํ อวิชฺชา\-นีวรณํ.}

\proseitem{1159}{ตตฺถ กตมํ กามจฺฉนฺท\-นีวรณํ. โย กาเมสุ กามจฺฉนฺโท กามราโค กามนนฺที กามตณฺหา กามสิเนโห กามปริฬาโห กามมุจฺฉา กามชฺโฌสานํ. อิทํ วุจฺจติ กามจฺฉนฺท\-นีวรณํ.}

\proseitem{1160}{ตตฺถ กตมํ พฺยาปาท\-นีวรณํ. “อนตฺถํ เม อจรี”ติ อาฆาโต ชายติ, “อนตฺถํ เม จรตี”ติ อาฆาโต ชายติ, “อนตฺถํ เม จริสฺสตี”ติ อาฆาโต ชายติ, “ปิยสฺส เม มนาปสฺส อนตฺถํ อจริ ฯเปฯ อนตฺถํ จรติ ฯเปฯ อนตฺถํ จริสฺสตี”ติ อาฆาโต ชายติ, \csromanfolio{233}“อปฺปิยสฺส เม อมนาปสฺส อตฺถํ อจริ ฯเปฯ อตฺถํ จรติ ฯเปฯ อตฺถํ จริสฺสตี”ติ อาฆาโต ชายติ, อฏฺฐาเน วา ปน อาฆาโต ชายติ. โย เอวรูโป จิตฺตสฺส อาฆาโต ปฏิฆาโต ปฏิฆํ ปฏิวิโรโธ โกโป ปโกโป สมฺปโกโป โทโส ปโทโส สมฺปโทโส จิตฺตสฺส พฺยาปตฺติ มโนปโทโส โกโธ กุชฺฌนา กุชฺฌิตตฺตํ โทโส ทุสฺสนา ทุสฺสิตตฺตํ พฺยาปตฺติ พฺยาปชฺชนา พฺยาปชฺชิตตฺตํ วิโรโธ ปฏิวิโรโธ จณฺฑิกฺกํ อสุโรโป อนตฺตมนตา จิตฺตสฺส. อิทํ วุจฺจติ พฺยาปาท\-นีวรณํ.}

\proseitem{1161}{ตตฺถ กตมํ ถินมิทฺธ\-นีวรณํ. อตฺถิ ถินํ, อตฺถิ มิทฺธํ.}

\proseitem{1162}{ตตฺถ กตมํ ถินํ. ยา จิตฺตสฺส อกลฺลตา อกมฺมญฺญตา โอลียนา สลฺลียนา ลีนํ ลียนา ลียิตตฺตํ ถินํ ถิยนา ถิยิตตฺตํ จิตฺตสฺส. อิทํ วุจฺจติ ถินํ.}

\proseitem{1163}{ตตฺถ กตมํ มิทฺธํ. ยา กายสฺส อกลฺลตา อกมฺมญฺญตา โอนาโห ปริโยนาโห อนฺโตสโมโรโธ มิทฺธํ โสปฺปํ ปจลายิกา โสปฺปํ สุปนา สุปิตตฺตํ. อิทํ วุจฺจติ มิทฺธํ. อิติ อิทญฺจ ถินํ อิทญฺจ มิทฺธํ, อิทํ วุจฺจติ ถินมิทฺธ\-นีวรณํ.}

\proseitem{1164}{ตตฺถ กตมํ อุทฺธจฺจ\-กุกฺกุจฺจ\-นีวรณํ. อตฺถิ อุทฺธจฺจํ, อตฺถิ กุกฺกุจฺจํ.}

\proseitem{1165}{ตตฺถ กตมํ อุทฺธจฺจํ. ยํ จิตฺตสฺส อุทฺธจฺจํ อวูปสโม เจตโส วิกฺเขโป ภนฺตตฺตํ จิตฺตสฺส. อิทํ วุจฺจติ อุทฺธจฺจํ.}

\proseitem{1166}{ตตฺถ กตมํ กุกฺกุจฺจํ. อกปฺปิเย กปฺปิยสญฺญิตา, กปฺปิเย อกปฺปิย\-สญฺญิตา, อวชฺเช วชฺชสญฺญิตา, วชฺเช อวชฺชสญฺญิตา. ยํ เอวรูปํ กุกฺกุจฺจํ กุกฺกุจฺจายนา กุกฺกุจฺจายิ\-ตตฺตํ เจตโส วิปฺปฏิสาโร มโนวิเลโข. อิทํ วุจฺจติ กุกฺกุจฺจํ. อิติ อิทญฺจ อุทฺธจฺจํ อิทญฺจ กุกฺกุจฺจํ, อิทํ วุจฺจติ อุทฺธจฺจ\-กุกฺกุจฺจ\-นีวรณํ.}

\proseitem{1167}{ตตฺถ กตมํ วิจิกิจฺฉา\-นีวรณํ. สตฺถริ กงฺขติ วิจิกิจฺฉติ, ธมฺเม กงฺขติ วิจิกิจฺฉติ, สํเฆ กงฺขติ วิจิกิจฺฉติ, สิกฺขาย กงฺขติ วิจิกิจฺฉติ, ปุพฺพนฺเต กงฺขติ วิจิกิจฺฉติ, อปรนฺเต กงฺขติ วิจิกิจฺฉติ, \csromanfolio{234}ปุพฺพนฺตา\-ปรนฺเต กงฺขติ วิจิกิจฺฉติ, อิทปฺปจฺจยตา ปฏิจฺจ\-สมุปฺปนฺเนสุ ธมฺเมสุ กงฺขติ วิจิกิจฺฉติ. ยา เอวรูปา กงฺขา กงฺขายนา กงฺขายิตตฺตํ วิมติ วิจิกิจฺฉา เทฺวฬฺหกํ เทฺวธาปโถ สํสโย อเนกํสคฺคาโห อาสปฺปนา ปริสปฺปนา อปริโยคาหนา ถมฺภิตตฺตํ จิตฺตสฺส มโนวิเลโข. อิทํ วุจฺจติ วิจิกิจฺฉา\-นีวรณํ.}

\proseitem{1168}{ตตฺถ กตมํ อวิชฺชา\-นีวรณํ. ทุกฺเข อญฺญาณํ, ทุกฺขสมุทเย อญฺญาณํ, ทุกฺขนิโรเธ อญฺญาณํ, ทุกฺขนิโรธ\-คามินิยา ปฏิปทาย อญฺญาณํ, ปุพฺพนฺเต อญฺญาณํ, อปรนฺเต อญฺญาณํ, ปุพฺพนฺตา\-ปรนฺเต อญฺญาณํ, อิทปฺปจฺจยตา ปฏิจฺจ\-สมุปฺปนฺเนสุ ธมฺเมสุ อญฺญาณํ. ยํ เอวรูปํ อญฺญาณํ อทสฺสนํ อนภิสมโย อนนุโพโธ อสมฺโพโธ อปฺปฏิเวโธ อสํคาหนา อปริโยคาหนา อสมเปกฺขนา อปจฺจเวกฺขณา อปจฺจกฺขกมฺมํ ทุมฺเมชฺฌํ พาลฺยํ อสมฺปชญฺญํ โมโห ปโมโห สมฺโมโห อวิชฺชา อวิชฺโชโฆ อวิชฺชาโยโค อวิชฺชานุสโย อวิชฺชา\-ปริยุฏฺฐานํ อวิชฺชาลงฺคี โมโห อกุสลมูลํ. อิทํ วุจฺจติ อวิชฺชา\-นีวรณํ. อิเม ธมฺมา นีวรณา.}

\proseitem{1169}{กตเม ธมฺมา โน นีวรณา. เต ธมฺเม ฐเปตฺวา อวเสสา กุสลากุสลา\-พฺยากตา ธมฺมา กามาวจรา รูปาวจรา อรูปาวจรา อปริยาปนฺนา เวทนากฺขนฺโธ ฯเปฯ วิญฺญาณ\-กฺขนฺโธ สพฺพญฺจ รูปํ อสงฺขตา จ ธาตุ. อิเม ธมฺมา โน นีวรณา.}

\proseitem{1170}{กตเม ธมฺมา นีวรณิยา. สาสวา กุสลากุสลา\-พฺยากตา ธมฺมา กามาวจรา รูปาวจรา อรูปาวจรา รูปกฺขนฺโธ ฯเปฯ วิญฺญาณ\-กฺขนฺโธ. อิเม ธมฺมา นีวรณิยา.}

\proseitem{1171}{กตเม ธมฺมา อนีวรณิยา. อปริยาปนฺนา มคฺคา จ มคฺคผลานิ จ อสงฺขตา จ ธาตุ. อิเม ธมฺมา อนีวรณิยา.}

\proseitem{1172}{กตเม ธมฺมา นีวรณ\-สมฺปยุตฺตา. เตหิ ธมฺเมหิ เย ธมฺมา สมฺปยุตฺตา เวทนากฺขนฺโธ ฯเปฯ วิญฺญาณ\-กฺขนฺโธ. อิเม ธมฺมา นีวรณ\-สมฺปยุตฺตา.}

\chapterheadpageread{3. นิกฺเขปกณฺฑ}
\proseitem{1173}{\csromanfolio{235}กตเม ธมฺมา นีวรณ\-วิปฺปยุตฺตา. เตหิ ธมฺเมหิ เย ธมฺมา วิปฺปยุตฺตา เวทนากฺขนฺโธ ฯเปฯ วิญฺญาณ\-กฺขนฺโธ สพฺพญฺจ รูปํ อสงฺขตา จ ธาตุ. อิเม ธมฺมา นีวรณ\-วิปฺปยุตฺตา.}

\proseitem{1174}{กตเม ธมฺมา นีวรณา เจว นีวรณิยา จ. ตาเนว นีวรณานิ นีวรณา เจว นีวรณิยา จ.}

\proseitem{1175}{กตเม ธมฺมา นีวรณิยา เจว โน จ นีวรณา. เตหิ ธมฺเมหิ เย ธมฺมา นีวรณิยา, เต ธมฺเม ฐเปตฺวา อวเสสา สาสวา กุสลากุสลา\-พฺยากตา ธมฺมา กามาวจรา รูปาวจรา อรูปาวจรา รูปกฺขนฺโธ ฯเปฯ วิญฺญาณ\-กฺขนฺโธ. อิเม ธมฺมา นีวรณิยา เจว โน จ นีวรณา.}

\proseitem{1176}{กตเม ธมฺมา นีวรณา เจว นีวรณ\-สมฺปยุตฺตา จ. กามจฺฉนฺท\-นีวรณํ อวิชฺชา\-นีวรเณน นีวรณญฺเจว นีวรณ\-สมฺปยุตฺตญฺจ, อวิชฺชา\-นีวรณํ กามจฺฉนฺท\-นีวรเณน นีวรณญฺเจว นีวรณ\-สมฺปยุตฺตญฺจ, พฺยาปาท\-นีวรณํ อวิชฺชา\-นีวรเณน นีวรณญฺเจว นีวรณาสมฺปยุตฺตญฺจ, อวิชฺชา\-นีวรณํ พฺยาปาท\-นีวรเณน นีวรณญฺเจว นีวรณ\-สมฺปยุตฺตญฺจ, ถินมิทฺธ\-นีวรณํ อวิชฺชา\-นีวรเณน นีวรณญฺเจว นีวรณ\-สมฺปยุตฺตญฺจ, อวิชฺชา\-นีวรณํ ถินมิทฺธ\-นีวรเณน นีวรณญฺเจว นีวรณ\-สมฺปยุตฺตญฺจ, อุทฺธจฺจ\-นีวรณํ อวิชฺชา\-นีวรเณน นีวรณญฺเจว นีวรณาสมฺปยุตฺตญฺจ, อวิชฺชา\-นีวรณํ อุทฺธจฺจ\-นีวรเณน นีวรณญฺเจว นีวรณ\-สมฺปยุตฺตญฺจ, กุกฺกุจฺจ\-นีวรณํ อวิชฺชา\-นีวรเณน นีวรณญฺเจว นีวรณ\-สมฺปยุตฺตญฺจ, อวิชฺชา\-นีวรณํ กุกฺกุจฺจ\-นีวรเณน นีวรณญฺเจว นีวรณ\-สมฺปยุตฺตญฺจ, วิจิกิจฺฉา\-นีวรณํ อวิชฺชา\-นีวรเณน นีวรณญฺเจว นีวรณ\-สมฺปยุตฺตญฺจ, อวิชฺชา\-นีวรณํ วิจิกิจฺฉา\-นีวรเณน นีวรณญฺเจว นีวรณ\-สมฺปยุตฺตญฺจ, กามจฺฉนฺท\-นีวรณํ อุทฺธจฺจ\-นีวรเณน นีวรณญฺเจว นีวรณ\-สมฺปยุตฺตญฺจ, อุทฺธจฺจ\-นีวรณํ กามจฺฉนฺท\-นีวรเณน นีวรณญฺเจว นีวรณ\-สมฺปยุตฺตญฺจ, พฺยาปาท\-นีวรณํ อุทฺธจฺจ\-นีวรเณน นีวรณญฺเจว นีวรณ\-สมฺปยุตฺตญฺจ, อุทฺธจฺจ\-นีวรณํ พฺยาปาท\-นีวรเณน นีวรณญฺเจว นีวรณ\-สมฺปยุตฺตญฺจ, ถินมิทฺธ\-นีวรณํ อุทฺธจฺจ\-นีวรเณน นีวรณญฺเจว นีวรณ\-สมฺปยุตฺตญฺจ, อุทฺธจฺจ\-นีวรณํ ถินมิทฺธ\-นีวรเณน นีวรณญฺเจว นีวรณ\-สมฺปยุตฺตญฺจ, กุกฺกุจฺจ\-นีวรณํ อุทฺธจฺจ\-นีวรเณน นีวรณญฺเจว นีวรณ\-สมฺปยุตฺตญฺจ, อุทฺธจฺจ\-นีวรณํ กุกฺกุจฺจ\-นีวรเณน นีวรณญฺเจว นีวรณ\-สมฺปยุตฺตญฺจ, วิจิกิจฺฉา\-นีวรณํ อุทฺธจฺจ\-นีวรเณน \csromanfolio{236}นีวรณญฺเจว นีวรณ\-สมฺปยุตฺตญฺจ, อุทฺธจฺจ\-นีวรณํ วิจิกิจฺฉา\-นีวรเณน นีวรณญฺเจว นีวรณ\-สมฺปยุตฺตญฺจ, อวิชฺชา\-นีวรณํ อุทฺธจฺจ\-นีวรเณน นีวรณญฺเจว นีวรณ\-สมฺปยุตฺตญฺจ, อุทฺธจฺจ\-นีวรณํ อวิชฺชา\-นีวรเณน นีวรณญฺเจว นีวรณ\-สมฺปยุตฺตญฺจ. อิเม ธมฺมา นีวรณา เจว นีวรณ\-สมฺปยุตฺตา จ.}

\proseitem{1177}{กตเม ธมฺมา นีวรณ\-สมฺปยุตฺตา เจว โน จ นีวรณา. เตหิ ธมฺเมหิ เย ธมฺมา สมฺปยุตฺตา, เต ธมฺเม ฐเปตฺวา เวทนากฺขนฺโธ ฯเปฯ วิญฺญาณ\-กฺขนฺโธ. อิเม ธมฺมา นีวรณ\-สมฺปยุตฺตา เจว โน จ นีวรณา.}

\proseitem{1178}{กตเม ธมฺมา นีวรณ\-วิปฺปยุตฺตา นีวรณิยา. เตหิ ธมฺเมหิ เย ธมฺมา วิปฺปยุตฺตา สาสวา กุสลาพฺยากตา ธมฺมา กามาวจรา รูปาวจรา อรูปาวจรา รูปกฺขนฺโธ ฯเปฯ วิญฺญาณ\-กฺขนฺโธ. อิเม ธมฺมา นีวรณ\-วิปฺปยุตฺตา นีวรณิยา.}

\proseitem{1179}{กตเม ธมฺมา นีวรณ\-วิปฺปยุตฺตา อนีวรณิยา. อปริยาปนฺนา มคฺคา จ มคฺคผลานิ จ อสงฺขตา จ ธาตุ. อิเม ธมฺมา นีวรณ\-วิปฺปยุตฺตา อนีวรณิยา.}

\csromanmatikaanchor{mtk.112}
\csromantocmarkref{section}{ปรามาสโคจฺฉก}{mtk.112}
\bookmark[dest={mtk.112},level=1]{ปรามาสโคจฺฉก}
\csromanheader{1}{ปรามาส\-โคจฺฉก}
\proseitem{1180}{กตเม ธมฺมา ปรามาสา. ทิฏฺฐิปรามาโส.}

\proseitem{1181}{ตตฺถ กตโม ทิฏฺฐิปรามาโส. “สสฺสโต โลโก”ติ วา, “อสสฺสโต โลโก”ติ วา, “อนฺตวา โลโก”ติ วา, “อนนฺตวา โลโก”ติ วา, “ตํ ชีวํ ตํ สรีรนฺ”ติ วา, “อญฺญํ ชีวํ อญฺญํ สรีรนฺ”ติ วา, “โหติ ตถาคโต ปรํ มรณา”ติ วา, “น โหติ ตถาคโต ปรํ มรณา”ติ วา, โหติ จ น จ โหติ ตถาคโต ปรํ มรณา”ติ วา, “เนว โหติ น น โหติ ตถาคโต ปรํ มรณา”ติ วา. ยา เอวรูปา ทิฏฺฐิ ทิฏฺฐิคตํ ทิฏฺฐิคหนํ ทิฏฺฐิกนฺตาโร ทิฏฺฐิ\-วิสูกายิกํ ทิฏฺฐิ\-วิปฺผนฺทิตํ ทิฏฺฐิ\-สญฺโญชนํ คาโห ปฏิคฺคาโห อภินิเวโส ปรามาโส กุมฺมคฺโค มิจฺฉาปโถ มิจฺฉตฺตํ ติตฺถายตนํ วิปริยาส\-คฺคาโห. อยํ วุจฺจติ ทิฏฺฐิปรามาโส. สพฺพาปิ มิจฺฉาทิฏฺฐิ ทิฏฺฐิปรามาโส.}

\vspace{\tipitakaparskip}
\csromancenter{อิเม ธมฺมา ปรามาสา.}
\chapterheadpageread{3. นิกฺเขปกณฺฑ}
\proseitem{1182}{\csromanfolio{237}กตเม ธมฺมา โน ปรามาสา. เต ธมฺเม ฐเปตฺวา อวเสสา กุสลากุสลา\-พฺยากตา ธมฺมา กามาวจรา รูปาวจรา อรูปาวจรา อปริยาปนฺนา เวทนากฺขนฺโธ ฯเปฯ วิญฺญาณ\-กฺขนฺโธ สพฺพญฺจ รูปํ อสงฺขตา จ ธาตุ. อิเม ธมฺมา โน ปรามาสา.}

\proseitem{1183}{กตเม ธมฺมา ปรามฏฺฐา. สาสวา กุสลากุสลา\-พฺยากตา ธมฺมา กามาวจรา รูปาวจรา อรูปาวจรา รูปกฺขนฺโธ ฯเปฯ วิญฺญาณ\-กฺขนฺโธ. อิเม ธมฺมา ปรามฏฺฐา.}

\proseitem{1184}{กตเม ธมฺมา อปรามฏฺฐา. อปริยาปนฺนา มคฺคา จ มคฺคผลานิ จ อสงฺขตา จ ธาตุ. อิเม ธมฺมา อปรามฏฺฐา.}

\proseitem{1185}{กตเม ธมฺมา ปรามาส\-สมฺปยุตฺตา. เตหิ ธมฺเมหิ เย ธมฺมา สมฺปยุตฺตา เวทนากฺขนฺโธ ฯเปฯ วิญฺญาณ\-กฺขนฺโธ. อิเม ธมฺมา ปรามาส\-สมฺปยุตฺตา.}

\proseitem{1186}{กตเม ธมฺมา ปรามาส\-วิปฺปยุตฺตา. เตหิ ธมฺเมหิ เย ธมฺมา วิปฺปยุตฺตา เวทนากฺขนฺโธ ฯเปฯ วิญฺญาณ\-กฺขนฺโธ สพฺพญฺจ รูปํ อสงฺขตา จ ธาตุ. อิเม ธมฺมา ปรามาส\-วิปฺปยุตฺตา.}

\proseitem{1187}{กตเม ธมฺมา ปรามาสา เจว ปรามฏฺฐา จ. เสฺวว ปรามาโส ปรามาโส เจว ปรามฏฺโฐ จ.}

\proseitem{1188}{กตเม ธมฺมา ปรามฏฺฐา เจว โน จ ปรามาสา. เตหิ ธมฺเมหิ เย ธมฺมา ปรามฏฺฐา, เต ธมฺเม ฐเปตฺวา อวเสสา สาสวา กุสลากุสลา\-พฺยากตา ธมฺมา กามาวจรา รูปาวจรา อรูปาวจรา รูปกฺขนฺโธ ฯเปฯ วิญฺญาณ\-กฺขนฺโธ. อิเม ธมฺมา ปรามฏฺฐา เจว โน จ ปรามาสา.}

\proseitem{1189}{กตเม ธมฺมา ปรามาส\-วิปฺปยุตฺตา ปรามฏฺฐา. เตหิ ธมฺเมหิ เย ธมฺมา วิปฺปยุตฺตา สาสวา กุสลากุสลา\-พฺยากตา ธมฺมา กามาวจรา รูปาวจรา อรูปาวจรา รูปกฺขนฺโธ ฯเปฯ วิญฺญาณ\-กฺขนฺโธ. อิเม ธมฺมา ปรามาส\-วิปฺปยุตฺตา ปรามฏฺฐา.}

\proseitem{1190}{กตเม ธมฺมา ปรามาส\-วิปฺปยุตฺตา อปรามฏฺฐา. อปริยาปนฺนา มคฺคา จ มคฺคผลานิ จ อสงฺขตา จ ธาตุ. อิเม ธมฺมา ปรามาส\-วิปฺปยุตฺตา อปรามฏฺฐา.}

\csromanmatikaanchor{mtk.113}
\csromantocmarkref{section}{มหนฺตรทุก}{mtk.113}
\bookmark[dest={mtk.113},level=1]{มหนฺตรทุก}
\csromanheader{1}{\textbf{มหนฺตรทุก}}
\proseitem{1191}{\csromanfolio{238}กตเม ธมฺมา สารมฺมณา. เวทนากฺขนฺโธ สญฺญากฺขนฺโธ สงฺขาร\-กฺขนฺโธ วิญฺญาณ\-กฺขนฺโธ. อิเม ธมฺมา สารมฺมณา.}

\proseitem{1192}{กตเม ธมฺมา อนารมฺมณา. สพฺพญฺจ รูปํ อสงฺขตา จ ธาตุ. อิเม ธมฺมา อนารมฺมณา.}

\proseitem{1193}{กตเม ธมฺมา จิตฺตา. จกฺขุวิญฺญาณํ โสตวิญฺญาณํ ฆานวิญฺญาณํ ชิวฺหาวิญฺญาณํ กายวิญฺญาณํ มโนธาตุ มโนวิญฺญาณ\-ธาตุ. อิเม ธมฺมา จิตฺตา.}

\proseitem{1194}{กตเม ธมฺมา โน จิตฺตา. เวทนากฺขนฺโธ สญฺญากฺขนฺโธ สงฺขรกฺขนฺโธ สพฺพญฺจ รูปํ อสงฺขตา จ ธาตุ. อิเม ธมฺมา โน จิตฺตา.}

\proseitem{1195}{กตเม ธมฺมา เจตสิกา. เวทนากฺขนฺโธ สญฺญากฺขนฺโธ สงฺขาร\-กฺขนฺโธ. อิเม ธมฺมา เจตสิกา.}

\proseitem{1196}{กตเม ธมฺมา อเจตสิกา. จิตฺตญฺจ สพฺพญฺจ รูปํ อสงฺขตา จ ธาตุ. อิเม ธมฺมา อเจตสิกา.}

\proseitem{1197}{กตเม ธมฺมา จิตฺต\-สมฺปยุตฺตา. เวทนากฺขนฺโธ สญฺญากฺขนฺโธ สงฺขาร\-กฺขนฺโธ. อิเม ธมฺมา จิตฺต\-สมฺปยุตฺตา.}

\proseitem{1198}{กตเม ธมฺมา จิตฺต\-วิปฺปยุตฺตา. สพฺพญฺจ รูปํ อสงฺขตา จ ธาตุ. อิเม ธมฺมา จิตฺต\-วิปฺปยุตฺตา. จิตฺตํ น วตฺตพฺพํ จิตฺเตน สมฺปยุตฺตนฺติปิ จิตฺเตน วิปฺปยุตฺตนฺติปิ.}

\proseitem{1199}{กตเม ธมฺมา จิตฺตสํสฏฺฐา. เวทนากฺขนฺโธ สญฺญากฺขนฺโธ สงฺขรกฺขนฺโธ. อิเม ธมฺมา จิตฺตสํสฏฺฐา.}

\proseitem{1200}{กตเม ธมฺมา จิตฺต\-วิสํสฏฺฐา. สพฺพญฺจ รูปํ อสงฺขตา จ ธาตุ. อิเม ธมฺมา จิตฺต\-วิสํสฏฺฐา. จิตฺตํ น วตฺตพฺพํ จิตฺเตน สํสฏฺฐนฺติปิ จิตฺเตน วิสํสฏฺฐนฺติปิ.}

\chapterheadpageread{3. นิกฺเขปกณฺฑ}
\proseitem{1201}{\csromanfolio{239}กตเม ธมฺมา จิตฺต\-สมุฏฺฐานา. เวทนากฺขนฺโธ สญฺญากฺขนฺโธ สงฺขาร\-กฺขนฺโธ กายวิญฺญตฺติ วจีวิญฺญตฺติ, ยํ วา ปนญฺญมฺปิ อตฺถิ รูปํ จิตฺตชํ จิตฺตเหตุกํ จิตฺต\-สมุฏฺฐานํ รูปายตนํ สทฺทายตนํ คนฺธายตนํ รสายตนํ โผฏฺฐพฺพา\-ยตนํ อากาสธาตุ อาโปธาตุ รูปสฺส ลหุตา รูปสฺส มุทุตา รูปสฺส กมฺมญฺญตา รูปสฺส อุปจโย รูปสฺส สนฺตติ กพฬีกาโร อาหาโร. อิเม ธมฺมา จิตฺต\-สมุฏฺฐานา.}

\proseitem{1202}{กตเม ธมฺมา โน จิตฺต\-สมุฏฺฐานา. จิตฺตญฺจ อวเสสญฺจ รูปํ อสงฺขตา จ ธาตุ. อิเม ธมฺมา โน จิตฺต\-สมุฏฺฐานา.}

\proseitem{1203}{กตเม ธมฺมา จิตฺตสหภุโน. เวทนากฺขนฺโธ สญฺญากฺขนฺโธ สงฺขาร\-กฺขนฺโธ กายวิญฺญตฺติ วจีวิญฺญตฺติ. อิเม ธมฺมา จิตฺตสหภุโน.}

\proseitem{1204}{กตเม ธมฺมา โน จิตฺตสหภุโน. จิตฺตญฺจ อวเสสญฺจ รูปํ อสงฺขตา จ ธาตุ. อิเม ธมฺมา โน จิตฺตสหภุโน.}

\proseitem{1205}{กตเม ธมฺมา จิตฺตา\-นุปริวตฺติโน. เวทนากฺขนฺโธ สญฺญากฺขนฺโธ สงฺขาร\-กฺขนฺโธ กายวิญฺญตฺติ วจีวิญฺญตฺติ. อิเม ธมฺมา จิตฺตา\-นุปริวตฺติโน.}

\proseitem{1206}{กตเม ธมฺมา โน จิตฺตา\-นุปริวตฺติโน. จิตฺตญฺจ อวเสสญฺจ รูปํ อสงฺขตา จ ธาตุ. อิเม ธมฺมา โน จิตฺตา\-นุปริวตฺติโน.}

\proseitem{1207}{กตเม ธมฺมา จิตฺต\-สํสฏฺฐ\-สมุฏฺฐานา. เวทนากฺขนฺโธ สญฺญากฺขนฺโธ สงฺขาร\-กฺขนฺโธ. อิเม ธมฺมา จิตฺต\-สํสฏฺฐ\-สมุฏฺฐานา.}

\proseitem{1208}{กตเม ธมฺมา โน จิตฺต\-สํสฏฺฐ\-สมุฏฺฐานา. จิตฺตญฺจ สพฺพญฺจ รูปํ อสงฺขตา จ ธาตุ. อิเม ธมฺมา โน จิตฺต\-สํสฏฺฐ\-สมุฏฺฐานา.}

\proseitem{1209}{กตเม ธมฺมา จิตฺต\-สํสฏฺฐ\-สมุฏฺฐาน\-สหภุโน. เวทนากฺขนฺโธ สญฺญากฺขนฺโธ สงฺขาร\-กฺขนฺโธ. อิเม ธมฺมา จิตฺต\-สํสฏฺฐ\-สมุฏฺฐาน\-สหภุโน.}

\proseitem{1210}{กตเม ธมฺมา โน จิตฺต\-สํสฏฺฐ\-สมุฏฺฐาน\-สหภุโน. จิตฺตญฺจ สพฺพญฺจ รูปํ อสงฺขตา จ ธาตุ. อิเม ธมฺมา โน จิตฺต\-สํสฏฺฐ\-สมุฏฺฐาน\-สหภุโน.}

\proseitem{1211}{กตเม ธมฺมา จิตฺต\-สํสฏฺฐ\-สมุฏฺฐานา\-นุปริวตฺติโน. เวทนากฺขนฺโธ สญฺญากฺขนฺโธ สงฺขาร\-กฺขนฺโธ. อิเม ธมฺมา จิตฺต\-สํสฏฺฐ\-สมุฏฺฐานา\-นุปริวตฺติโน.}

\proseitem{1212}{\csromanfolio{240}กตเม ธมฺมา โน จิตฺตสํสฏฺฐสมุฏฺฐนานุปริวตฺติโน. จิตฺตญฺจ สพฺพญฺจ รูปํ อสงฺขตา จ ธาตุ. อิเม ธมฺมา โน จิตฺต\-สํสฏฺฐ\-สมุฏฺฐานา\-นุปริวตฺติโน.}

\proseitem{1213}{กตเม ธมฺมา อชฺฌตฺติกา. จกฺขายตนํ ฯเปฯ มนายตนํ. อิเม ธมฺมา อชฺฌตฺติกา.}

\proseitem{1214}{กตเม ธมฺมา พาหิรา. รูปายตนํ ฯเปฯ ธมฺมายตนํ. อิเม ธมฺมา พาหิรา.}

\proseitem{1215}{กตเม ธมฺมา อุปาทา. จกฺขายตนํ ฯเปฯ กพฬีกาโร อาหาโร. อิเม ธมฺมา อุปาทา.}

\proseitem{1216}{กตเม ธมฺมา โน อุปาทา. เวทนากฺขนฺโธ สญฺญากฺขนฺโธ สงฺขาร\-กฺขนฺโธ วิญฺญาณ\-กฺขนฺโธ จตฺตาโร จ มหาภูตา อสงฺขตา จ ธาตุ. อิเม ธมฺมา โน อุปาทา.}

\proseitem{1217}{กตเม ธมฺมา อุปาทิณฺณา. สาสวา กุสลา\-กุสลานํ ธมฺมานํ วิปากา กามาวจรา รูปาวจรา อรูปาวจรา เวทนากฺขนฺโธ ฯเปฯ วิญฺญาณ\-กฺขนฺโธ, ยญฺจ รูปํ น กมฺมสฺส กตตฺตา. อิเม ธมฺมา อุปาทิณฺณา.}

\proseitem{1218}{กตเม ธมฺมา อนุปาทิณฺณา. สาสวา กุสลากุสลา ธมฺมา กามาวจรา รูปาวจรา อรูปาวจรา เวทนากฺขนฺโธ ฯเปฯ วิญฺญาณ\-กฺขนฺโธ, เย จ ธมฺมา กิริยา เนว กุสลา นากุสลา น จ กมฺมวิปากา, ยญฺจ รูปํ น กมฺมสฺส กตตฺตา อปริยาปนฺนา มคฺคา จ มคฺคผลานิ จ อสงฺขตา จ ธาตุ. อิเม ธมฺมา อนุปาทิณฺณา.}

\csromanmatikaanchor{mtk.114}
\csromantocmarkref{section}{อุปาทานโคจฺฉก}{mtk.114}
\bookmark[dest={mtk.114},level=1]{อุปาทานโคจฺฉก}
\csromanheader{1}{อุปาทาน\-โคจฺฉก}
\proseitem{1219}{กตเม ธมฺมา อุปาทานา. จตฺตาริ อุปาทานานิ กามุปาทานํ ทิฏฺฐุปาทานํ สีลพฺพตุ\-ปาทานํ อตฺตวาทุ\-ปาทานํ\csromansharedfootnote{5ca29d2afd46c948}{กามูปาทานํ ทิฏฺฐูปาทานํ สีลพฺพตูปาทานํ อตฺตวาทูปาทานํ (สี)}.}

\chapterheadpageread{3. นิกฺเขปกณฺฑ}
\proseitem{1220}{\csromanfolio{241}ตตฺถ กตมํ กามุปาทานํ. โย กาเมสุ กามจฺฉนฺโท กามราโค กามนนฺที กามตณฺหา กามสิเนโห กามปริฬาโห กามมุจฺฉา กามชฺโฌสานํ. อิทํ วุจฺจติ กามุปาทานํ.}

\proseitem{1221}{ตตฺถ กตมํ ทิฏฺฐุปาทานํ. “นตฺถิ ทินฺนํ, นตฺถิ ยิฏฺฐํ, นตฺถิ หุตํ, นตฺถิ สุกต\-ทุกฺกฏานํ กมฺมานํ ผลํ วิปาโก, นตฺถิ อยํ โลโก, นตฺถิ ปโร โลโก, นตฺถิ มาตา, นตฺถิ ปิตา, นตฺถิ สตฺตา โอปปาติกา, นตฺถิ โลเก สมณพฺราหฺมณา สมฺมคฺคตา\csromansharedfootnote{44029fb9de358ac7}{สมคฺคตา (ก)} สมฺมาปฏิปนฺนา, เย อิมญฺจ โลกํ ปรญฺจ โลกํ สยํ อภิญฺญา สจฺฉิกตฺวา ปเวเทนฺตี”ติ ยา เอวรูปา ทิฏฺฐิ ทิฏฺฐิคตํ ทิฏฺฐิคหนํ ทิฏฺฐิกนฺตาโร ทิฏฺฐิ\-วิสูกายิกํ ทิฏฺฐิ\-วิปฺผนฺทิตํ ทิฏฺฐิ\-สญฺโญชนํ คาโห ปติฏฺฐาโห อภินิเวโส ปรามาโส กุมฺมคฺโค มิจฺฉาปโถ มิจฺฉตฺตํ ติตฺถายตนํ วิปริยาส\-คฺคาโห. อิทํ วุจฺจติ ทิฏฺฐุปาทานํ. ฐเปตฺวา สีลพฺพตุปาทานญฺจ อตฺตวาทุปาทานญฺจ สพฺพาปิ มิจฺฉาทิฏฺฐิ ทิฏฺฐุปาทานํ.}

\proseitem{1222}{ตตฺถ กตมํ สีลพฺพตุ\-ปาทานํ. “อิโต พหิทฺธา สมณ\-พฺราหฺมณานํ สีเลน สุทฺธิ วเตน สุทฺธิ สีลพฺพเตน สุทฺธี”ติ ยา เอวรูปา ทิฏฺฐิ ทิฏฺฐิคตํ ทิฏฺฐิคหนํ ทิฏฺฐิกนฺตาโร ทิฏฺฐิ\-วิสูกายิกํ ทิฏฺฐิ\-วิปฺผนฺทิตํ ทิฏฺฐิ\-สญฺโญชนํ คาโห ปติฏฺฐาโห อภินิเวโส ปรามาโส กุมฺมคฺโค มิจฺฉาปโถ มิจฺฉตฺตํ ติตฺถายตนํ วิปริยาส\-คฺคาโห. อิทํ วุจฺจติ สีลพฺพตุ\-ปาทานํ.}

\proseitem{1223}{ตตฺถ กตมํ อตฺตวาทุ\-ปาทานํ. อิธ อสฺสุตวา ปุถุชฺชโน อริยานํ อทสฺสาวี อริยธมฺมสฺส อโกวิโท อริยธมฺเม อวินีโต สปฺปุริสานํ อทสฺสาวี สปฺปุริส\-ธมฺมสฺส อโกวิโท สปฺปุริส\-ธมฺเม อวินีโต รูปํ อตฺตโต สมนุปสฺสติ, รูปวนฺตํ วา อตฺตานํ, อตฺตนิ วา รูปํ, รูปสฺมิํ วา อตฺตานํ. เวทนํ ฯเปฯ. สญฺญํ ฯเปฯ. สงฺขาเร ฯเปฯ. วิญฺญาณํ อตฺตโต สมนุปสฺสติ, วิญฺญาณวนฺตํ วา อตฺตานํ, อตฺตนิ วา วิญฺญาณํ, วิญฺญาณสฺมิํ วา อตฺตานํ. ยา เอวรูปา ทิฏฺฐิ ทิฏฺฐิคตํ ทิฏฺฐิคหนํ ทิฏฺฐิกนฺตาโร ทิฏฺฐิ\-วิสูกายิกํ ทิฏฺฐิ\-วิปฺผนฺทิตํ ทิฏฺฐิ\-สญฺโญชนํ คาโห ปติฏฺฐาโห อภินิเวโส ปรามาโส กุมฺมคฺโค มิจฺฉาปโถ มิจฺฉตฺตํ ติตฺถายตนํ วิปริยาส\-คฺคาโห. อิทํ วุจฺจติ อตฺตวาทุ\-ปาทานํ.}

\vspace{\tipitakaparskip}
\csromancenter{อิเม ธมฺมา อุปาทานา.}
\proseitem{1224}{\csromanfolio{242}กตเม ธมฺมา โน อุปาทานา. เต ธมฺเม ฐเปตฺวา อวเสสา กุสลากุสลา\-พฺยากตา ธมฺมา กามาวจรา รูปาวจรา อรูปาวจรา อปริยาปนฺนา เวทนากฺขนฺโธ ฯเปฯ วิญฺญาณ\-กฺขนฺโธ สพฺพญฺจ รูปํ อสงฺขตา จ ธาตุ. อิเม ธมฺมา โน อุปาทานา.}

\proseitem{1225}{กตเม ธมฺมา อุปาทานิยา. สาสวา กุสลากุสลา\-พฺยากตา ธมฺมา กามาวจรา รูปาวจรา อรูปวจรา รูปกฺขนฺโธ ฯเปฯ วิญฺญาณ\-กฺขนฺโธ. อิเม ธมฺมา อุปาทานิยา.}

\proseitem{1226}{กตเม ธมฺมา อนุปาทานิยา. อปริยาปนฺนา มคฺคา จ มคฺคผลานิ จ อสงฺขตา จ ธาตุ. อิเม ธมฺมา อนุปาทานิยา.}

\proseitem{1227}{กตเม ธมฺมา อุปาทาน\-สมฺปยุตฺตา. เตหิ ธมฺเมหิ เย ธมฺมา สมฺปยุตฺตา เวทนากฺขนฺโธ. วิญฺญาณ\-กฺขนฺโธ. อิเม ธมฺมา อุปาทาน\-สมฺปยุตฺตา.}

\proseitem{1228}{กตเม ธมฺมา อุปาทาน\-วิปฺปยุตฺตา. เตหิ ธมฺเมหิ เย ธมฺมา วิปฺปยุตฺตา เวทนากฺขนฺโธ ฯเปฯ วิญฺญาณ\-กฺขนฺโธ สพฺพญฺจ รูปํ อสงฺขตา จ ธาตุ. อิเม ธมฺมา อุปาทาน\-วิปฺปยุตฺตา.}

\proseitem{1229}{กตเม ธมฺมา อุปาทานา เจว อุปาทานิยา จ. ตาเนว อุปาทานานิ อุปาทานา เจว อุปาทานิยา จ.}

\proseitem{1230}{กตเม ธมฺมา อุปาทานิยา เจว โน จ อุปาทานา. เตหิ ธมฺเมหิ เย ธมฺมา อุปาทานิยา, เต ธมฺเม ฐเปตฺวา อวเสสา สาสวา กุสลากุสลา\-พฺยากตา ธมฺมา กามาวจรา รูปาวจรา อรูปาวจรา รูปกฺขนฺโธ ฯเปฯ วิญฺญาณ\-กฺขนฺโธ. อิเม ธมฺมา อุปาทานิยา เจว โน จ อุปาทานา.}

\proseitem{1231}{กตเม ธมฺมา อุปาทานา เจว อุปาทาน\-สมฺปยุตฺตา จ. ทิฏฺฐุปาทานํ กามุปาทาเนน อุปาทานญฺเจว อุปาทาน\-สมฺปยุตฺตญฺจ, กามุปาทานํ ทิฏฺฐุปาทาเนน อุปาทานญฺเจว อุปาทาน\-สมฺปยุตฺตญฺจ, สีลพฺพตุ\-ปาทานํ กามุปาทาเนน อุปาทานญฺเจว อุปาทาน\-สมฺปยุตฺตญฺจ, กามุปาทานํ สีลพฺพตุ\-ปาทาเนน อุปาทานญฺเจว \csromanfolio{243}อุปาทาน\-สมฺปยุตฺตญฺจ, อตฺตวาทุ\-ปาทานํ กามุปาทาเนน อุปาทานญฺเจว อุปาทาน\-สมฺปยุตฺตญฺจ, กามุปาทานํ อตฺตวาทุ\-ปาทาเนน อุปาทานญฺเจว อุปาทาน\-สมฺปยุตฺตญฺจ. อิเม ธมฺมา อุปาทานา เจว อุปาทาน\-สมฺปยุตฺตา จ.}

\proseitem{1232}{กตเม ธมฺมา อุปาทาน\-สมฺปยุตฺตา เจว โน จ อุปาทานา. เตหิ ธมฺเมหิ เย ธมฺมา สมฺปยุตฺตา, เต ธมฺเม ฐเปตฺวา เวทนากฺขนฺโธ ฯเปฯ วิญฺญาณ\-กฺขนฺโธ. อิเม ธมฺมา อุปาทาน\-สมฺปยุตฺตา เจว โน จ อุปาทานา.}

\proseitem{1233}{กตเม ธมฺมา อุปาทาน\-วิปฺปยุตฺตา อุปาทานิยา. เตหิ ธมฺเมหิ เย ธมฺมา วิปฺปยุตฺตา สาสวา กุสลากุสลา\-พฺยากตา ธมฺมา กามาวจรา รูปาวจรา อรูปาวจรา รูปกฺขนฺโธ ฯเปฯ วิญฺญาณ\-กฺขนฺโธ. อิเม ธมฺมา อุปาทาน\-วิปฺปยุตฺตา อุปาทานิยา.}

\proseitem{1234}{กตเม ธมฺมา อุปาทาน\-วิปฺปยุตฺตา อนุปาทานิยา. อปริยาปนฺนา มคฺคา จ มคฺคผลานิ จ อสงฺขตา จ ธาตุ. อิเม ธมฺมา อุปาทาน\-วิปฺปยุตฺตา อนุปาทานิยา.}

\vspace{\tipitakaparskip}
\csromancenter{นิกฺเขปกณฺเฑ ทุติย\-ภาณวาโร.}
\csromanmatikaanchor{mtk.115}
\csromantocmarkref{section}{กิเลสโคจฺฉก}{mtk.115}
\bookmark[dest={mtk.115},level=1]{กิเลสโคจฺฉก}
\csromanheader{1}{\textbf{กิเลสโคจฺฉก}}
\proseitem{1235}{กตเม ธมฺมา กิเลสา. ทส กิเลสวตฺถูนิ โลโภ โทโส โมโห มาโน ทิฏฺฐิ วิจิกิจฺฉา ถินํ อุทฺธจฺจํ อหิรีกํ อโนตฺตปฺปํ.}

\proseitem{1236}{ตตฺถ กตโม โลโภ. โย ราโค สาราโค อนุนโย อนุโรโธ นนฺที นนฺทีราโค จิตฺตสฺส สาราโค อิจฺฉา มุจฺฉา อชฺโฌสานํ เคโธ ปลิเคโธ สงฺโค ปงฺโก เอชา มายา ชนิกา สญฺชนนี สิพฺพินี ชาลินี สริตา วิสตฺติกา สุตฺตํ วิสฏา อายูหินี ทุติยา ปณิธิ ภวเนตฺติ วนํ วนโถ สนฺถโว สิเนโห อเปกฺขา ปฏิพนฺธุ อาสา อาสิสนา อาสิสิตตฺตํ รูปาสา สทฺทาสา คนฺธาสา \csromanfolio{244}รสาสา โผฏฺฐพฺพาสา ลาภาสา ธนาสา ปุตฺตาสา ชีวิตาสา ชปฺปา ปชปฺปา อภิชปฺปา ชปฺปา ชปฺปนา ชปฺปิตตฺตํ โลลุปฺปํ โลลุปฺปายนา โลลุปฺปา\-ยิตตฺตํ ปุจฺฉญฺชิกตา สาธุกมฺยตา อธมฺมราโค วิสมโลโภ นิกนฺติ นิกามนา ปตฺถนา ปิหนา สมฺปตฺถานา กามตณฺหา ภวตณฺหา วิภวตณฺหา รูปตณฺหา อรูปตณฺหา นิโรธตณฺหา รูปตณฺหา สทฺทตณฺหา คนฺธตณฺหา รสตณฺหา โผฏฺฐพฺพ\-ตณฺหา ธมฺมตณฺหา โอโฆ โยโค คนฺโถ อุปาทานํ อาวรณํ นีวรณํ ฉาทนํ พนฺธนํ อุปกฺกิเลโส อนุสโย ปริยุฏฺฐานํ ลตา เววิจฺฉํ ทุกฺขมูลํ ทุกฺขนิทานํ ทุกฺขปฺปภโว มารปาโส มารพฬิสํ มารวิสโย ตณฺหานที ตณฺหาชาลํ ตณฺหาคทฺทุลํ ตณฺหาสมุทฺโท อภิชฺฌา โลโภ อกุสลมูลํ. อยํ วุจฺจติ โลโภ.}

\proseitem{1237}{ตตฺถ กตโม โทโส. “อนตฺถํ เม อจรี”ติ อาฆาโต ชายติ, “อนตฺถํ เม จรตี”ติ อาฆาโต ชายติ, “อนตฺถํ เม จริสฺสตี”ติ อาฆาโต ชายติ, “ปิยสฺส เม มนาปสฺส อนตฺถํ อจริ ฯเปฯ อนตฺถํ จรติ ฯเปฯ อนตฺถํ จริสฺสตี”ติ อาฆาโต ชายติ, “อปฺปิยสฺส เม อมนาปสฺส อตฺถํ อจริ ฯเปฯ อตฺถํ จรติ ฯเปฯ อตฺถํ จริสฺสตี”ติ อาฆาโต ชายติ, อฏฺฐาเน วา ปน อาฆาโต ชายติ. โย เอวรูโป จิตฺตสฺส อาฆาโต ปฏิฆาโต ปฏิฆํ ปฏิวิโรโธ โกโป ปโกโป สมฺปโกโป โทโส ปโทโส สมฺปโทโส จิตฺตสฺส พฺยาปตฺติ มโนปโทโส โกโธ กุชฺฌนา กุชฺฌิตตฺตํ โทโส ทุสฺสนา ทุสฺสิตตฺตํ พฺยาปตฺติ พฺยาปชฺชนา พฺยาปชฺชิตตฺตํ วิโรโธ ปฏิวิโรโธ จณฺฑิกฺกํ อสุโรโป อนตฺตมนตา จิตฺตสฺส. อยํ วุจฺจติ โทโส.}

\proseitem{1238}{ตตฺถ กตโม โมโห. ทุกฺเข อญฺญาณํ, ทุกฺขสมุทเย อญฺญาณํ, ทุกฺขนิโรเธ อญฺญาณํ, ทุกฺขนิโรธ\-คามินิยา ปฏิปทาย อญฺญาณํ ปุพฺพนฺเต อญฺญาณํ, อปรนฺเต อญฺญาณํ, ปุพฺพนฺตา\-ปรนฺเต อญฺญาณํ, อิทปฺปจฺจยตา ปฏิจฺจ\-สมุปฺปนฺเนสุ ธมฺเมสุ อญฺญาณํ. ยํ เอวรูปํ อญฺญาณํ อทสฺสนํ อนภิสมโย อนนุโพโธ อสมฺโพโธ อปฺปฏิเวโธ อสํคาหนา อปริโยคาหนา อสมเปกฺขนา อปจฺจเวกฺขณา อปจฺจกฺขกมฺมํ ทุมฺเมชฺฌํ พาลฺยํ อสมฺปชญฺญํ โมโห ปโมโห สมฺโมโห อวิชฺชา อวิชฺโชโฆ อวิชฺชาโยโค อวิชฺชานุสโย อวิชฺชา\-ปริยุฏฺฐานํ อวิชฺชาลงฺคี โมโห อกุสลมูลํ. อยํ วุจฺจติ โมโห.}

\chapterheadpageread{3. นิกฺเขปกณฺฑ}
\proseitem{1239}{\csromanfolio{245}ตตฺถ กตโม มาโน. “เสยฺโยหมสฺมี”ติ มาโน, “สทิโสหมสฺมี”ติ มาโน, “หีโนหมสฺมี”ติ มาโน. โย เอวรูโป มาโน มญฺญนา มญฺญิตตฺตํ อุนฺนติ อุนฺนโม ธโช สมฺปคฺคาโห เกตุกมฺยตา จิตฺตสฺส. อยํ วุจฺจติ มาโน.}

\proseitem{1240}{ตตฺถ กตมา ทิฏฺฐิ. “สสฺสโต โลโก”ติ วา, “อสสฺสโต โลโก”ติ วา, “อนฺตวา โลโก”ติ วา, “อนนฺตวา โลโก”ติ วา, “ตํ ชีวํ ตํ สรีรนฺ”ติ วา, “อญฺญํ ชีวํ อญฺญํ สรีรนฺ”ติ วา, “โหติ ตถาคโต ปรํ มรณา”ติ วา, “น โหติ ตถาคโต ปรํ มรณา”ติ วา, “โหติ จ น จ โหติ ตถาคโต ปรํ มรณา”ติ วา, “เนว โหติ น น โหติ ตถาคโต ปรํ มรณา”ติ วา. ยา เอวรูปา ทิฏฺฐิ ทิฏฺฐิคตํ ทิฏฺฐิคหนํ ทิฏฺฐิกนฺตาโร ทิฏฺฐิ\-วิสูกายิกํ ทิฏฺฐิ\-วิปฺผนฺทิตํ ทิฏฺฐิ\-สญฺโญชนํ คาโห ปติฏฺฐาโห อภินิเวโส ปรามาโส กุมฺมคฺโค มิจฺฉาปโถ มิจฺฉตฺตํ ติตฺถายตนํ วิปริยาส\-คฺคาโห. อยํ วุจฺจติ ทิฏฺฐิ. สพฺพาปิ มิจฺฉาทิฏฺฐิ ทิฏฺฐิ.}

\proseitem{1241}{ตตฺถ กตมา วิจิกิจฺฉา. สตฺถริ กงฺขติ วิจิกิจฺฉติ, ธมฺเม กงฺขติ วิจิกิจฺฉติ, สํเฆ กงฺขติ วิจิกิจฺฉติ, สิกฺขาย กงฺขติ วิจิกิจฺฉติ, ปุพฺพนฺเต กงฺขติ วิจิกิจฺฉติ, อปรนฺเต กงฺขติ วิจิกิจฺฉติ, ปุพฺพนฺตา\-ปรนฺเต กงฺขติ วิจิกิจฺฉติ, อิทปฺปจฺจยตา ปฏิจฺจ\-สมุปฺปนฺเนสุ ธมฺเมสุ กงฺขติ วิจิกิจฺฉติ. ยา เอวรูปา กงฺขา กงฺขายนา กงฺขายิตตฺตํ วิมติ วิจิกิจฺฉา เทฺวฬฺหกํ เทฺวธาปโถ สํสโย อเนกํสคฺคาโห อาสปฺปนา ปริสปฺปนา อปริโยคาหนา ถมฺภิตตฺตํ จิตฺตสฺส มโนวิเลโข. อยํ วุจฺจติ วิจิกิจฺฉา.}

\proseitem{1242}{ตตฺถ กตมํ ถินํ. ยา จิตฺตสฺส อกลฺลตา อกมฺมญฺญตา โอลียนา สลฺลียนา ลีนํ ลียนา ลียิตตฺตํ ถินํ ถิยนา ถิยิตตฺตํ จิตฺตสฺส. อิทํ วุจฺจติ ถินํ.}

\proseitem{1243}{ตตฺถ กตมํ อุทฺธจฺจํ. ยํ จิตฺตสฺส อุทฺธจฺจํ อวูปสโม เจตโส วิกฺเขโป ภนฺตตฺตํ จิตฺตสฺส. อิทํ วุจฺจติ อุทฺธจฺจํ.}

\proseitem{1244}{ตตฺถ กตมํ อหิริกํ. ยํ น หิรียติ หิริยิตพฺเพน น หิรียติ ปาปกานํ อกุสลานํ ธมฺมานํ สมาปตฺติยา. อิทํ วุจฺจติ อหิริกํ.}

\proseitem{1245}{\csromanfolio{246}ตตฺถ กตมํ อโนตฺตปฺปํ. ยํ น โอตฺตปฺปติ โอตฺตปฺปิตพฺเพน น โอตฺตปฺปติ ปาปกานํ อกุสลานํ ธมฺมานํ สมาปตฺติยา. อิทํ วุจฺจติ อโนตฺตปฺปํ. อิเม ธมฺมา กิเลสา.}

\proseitem{1246}{กตเม ธมฺมา โน กิเลสา. เต ธมฺเม ฐเปตฺวา อวเสสา กุสลากุสลา\-พฺยากตา ธมฺมา กามาวจรา รูปาวจรา อรูปาวจรา อปริยาปนฺนา เวทนากฺขนฺโธ ฯเปฯ วิญฺญาณ\-กฺขนฺโธ สพฺพญฺจ รูปํ อสงฺขตา จ ธาตุ. อิเม ธมฺมา โน กิเลสา.}

\proseitem{1247}{กตเม ธมฺมา สํกิเลสิกา. สาสวา กุสลากุสลา\-พฺยากตา ธมฺมา. กามาวจรา รูปาวจรา อรูปาวจรา รูปกฺขนฺโธ ฯเปฯ วิญฺญาณ\-กฺขนฺโธ. อิเม ธมฺมา สํกิเลสิกา.}

\proseitem{1248}{กตเม ธมฺมา อสํกิเลสิกา อปริยาปนฺนา มคฺคา จ มคฺคผลานิ จ อสงฺขตา จ ธาตุ. อิเม ธมฺมา อสํกิเลสิกา.}

\proseitem{1249}{กตเม ธมฺมา สํกิลิฏฺฐา. ตีณิ อกุสลมูลานิ โลโภ โทโส โมโห, ตเทกฏฺฐา จ กิเลสา ตํสมฺปยุตฺโต เวทนากฺขนฺโธ ฯเปฯ วิญฺญาณ\-กฺขนฺโธ, ตํสมุฏฺฐานํ กายกมฺมํ วจีกมฺมํ มโนกมฺมํ. อิเม ธมฺมา สํกิลิฏฺฐา.}

\proseitem{1250}{กตเม ธมฺมา อสํกิลิฏฺฐา. กุสลาพฺยากตา ธมฺมา กามาวจรา รูปาวจรา อรูปาวจรา อปริยาปนฺนา เวทนากฺขนฺโธ ฯเปฯ วิญฺญาณ\-กฺขนฺโธ สพฺพญฺจ รูปํ อสงฺขตา จ ธาตุ. อิเม ธมฺมา อสํกิลิฏฺฐา.}

\proseitem{1251}{กตเม ธมฺมา กิเลส\-สมฺปยุตฺตา. เตหิ ธมฺเมหิ เย ธมฺมา สมฺปยุตฺตา เวทนากฺขนฺโธ ฯเปฯ วิญฺญาณ\-กฺขนฺโธ. อิเม ธมฺมา กิเลส\-สมฺปยุตฺตา.}

\proseitem{1252}{กตเม ธมฺมา กิเลส\-วิปฺปยุตฺตา. เตหิ ธมฺเมหิ เย ธมฺมา วิปฺปยุตฺตา เวทนากฺขนฺโธ ฯเปฯ วิญฺญาณ\-กฺขนฺโธ สพฺพญฺจ รูปํ อสงฺขตา จ ธาตุ. อิเม ธมฺมา กิเลส\-วิปฺปยุตฺตา.}

\chapterheadpageread{3. นิกฺเขปกณฺฑ}
\proseitem{1253}{\csromanfolio{247}กตเม ธมฺมา กิเลสา เจว สํกิเลสิกา จ. เตว กิเลสา กิเลสา เจว สํกิเลสิกา จ.}

\proseitem{1254}{กตเม ธมฺมา สํกิเลสิกา เจว โน จ กิเลสา. เตหิ ธมฺเมหิ เย ธมฺมา สํกิเลสิกา, เต ธมฺเม ฐเปตฺวา อวเสสา สาสวา กุสลากุสลา\-พฺยากตา ธมฺมา กามาวจรา รูปาวจรา อรูปาวจรา รูปกฺขนฺโธ ฯเปฯ วิญฺญาณ\-กฺขนฺโธ. อิเม ธมฺมา สํกิเลสิกา เจว โน จ กิเลสา.}

\proseitem{1255}{กตเม ธมฺมา กิเลสา เจว สํกิลิฏฺฐา จ. เตว กิเลสา กิเลสา เจว สํกิลิฏฺฐา จ.}

\proseitem{1256}{กตเม ธมฺมา สํกิลิฏฺฐา เจว โน จ กิเลสา. เตหิ ธมฺเมหิ เย ธมฺมา สํกิลิฏฺฐา, เต ธมฺเม ฐเปตฺวา เวทนากฺขนฺโธ ฯเปฯ วิญฺญาณ\-กฺขนฺโธ. อิเม ธมฺมา สํกิลิฏฺฐา เจว โน จ กิเลสา.}

\proseitem{1257}{กตเม ธมฺมา กิเลสา เจว กิเลส\-สมฺปยุตฺตา จ. โลโภ โมเหน กิเลโส เจว กิเลส\-สมฺปยุตฺโต จ, โมโห โลเภน กิเลโส เจว กิเลส\-สมฺปยุตฺโต จ, โทโส โมเหน กิเลโส เจว กิเลส\-สมฺปยุตฺโต จ, โมโห โทเสน กิเลโส เจว กิเลส\-สมฺปยุตฺโต จ, มาโน โมเหน กิเลโส เจว กิเลส\-สมฺปยุตฺโต จ, โมโห มาเนน กิเลโส เจว กิเลส\-สมฺปยุตฺโต จ, ทิฏฺฐิ โมเหน กิเลโส เจว กิเลส\-สมฺปยุตฺตา จ, โมโห ทิฏฺฐิยา กิเลโส เจว กิเลส\-สมฺปยุตฺโต จ, วิจิกิจฺฉา โมเหน กิเลโส เจว กิเลส\-สมฺปยุตฺตา จ, โมโห วิจิกิจฺฉาย กิเลโส เจว กิเลส\-สมฺปยุตฺโต จ, ถินํ โมเหน กิเลโส เจว กิเลส\-สมฺปยุตฺตญฺจ, โมโห ถิเนน กิเลโส เจว กิเลส\-สมฺปยุตฺโต จ, อุทฺธจฺจํ โมเหน กิเลโส เจว กิเลส\-สมฺปยุตฺตญฺจ, โมโห อุทฺธจฺเจน กิเลโส เจว กิเลส\-สมฺปยุตฺโต จ, อหิริกํ โมเหน กิเลโส เจว กิเลส\-สมฺปยุตฺตญฺจ, โมโห อหิริเกน กิเลโส เจว กิเลส\-สมฺปยุตฺโต จ, อโนตฺตปฺปํ โมเหน กิเลโส เจว กิเลส\-สมฺปยุตฺตญฺจ, โมโห อโนตฺตปฺเปน กิเลโส เจว กิเลส\-สมฺปยุตฺโต จ, โลโภ \csromanfolio{248}อุทฺธจฺเจน กิเลโส เจว กิเลส\-สมฺปยุตฺโต จ, อุทฺธจฺจํ โลเภน กิเลโส เจว กิเลส\-สมฺปยุตฺตญฺจ, โทโส อุทฺธจฺเจน กิเลโส เจว กิเลส\-สมฺปยุตฺโต จ, อุทฺธจฺจํ โทเสน กิเลโส เจว กิเลส\-สมฺปยุตฺตญฺจ, โมโห อุทฺธจฺเจน กิเลโส เจว กิเลส\-สมฺปยุตฺโต จ, อุทฺธจฺจํ โมเหน กิเลโส เจว กิเลส\-สมฺปยุตฺตญฺจ, มาโน อุทฺธจฺเจน กิเลโส เจว กิเลส\-สมฺปยุตฺโต จ, อุทฺธจฺจํ มาเนน กิเลโส เจว กิเลส\-สมฺปยุตฺตญฺจ, ทิฏฺฐิ อุทฺธจฺเจน กิเลโส เจว กิเลส\-สมฺปยุตฺตา จ, อุทฺธจฺจํ ทิฏฺฐิยา กิเลโส เจว กิเลส\-สมฺปยุตฺตญฺจ, วิจิกิจฺฉา อุทฺธจฺเจน กิเลโส เจว กิเลส\-สมฺปยุตฺตา จ, อุทฺธจฺจํ วิจิกิจฺฉาย กิเลโส เจว กิเลส\-สมฺปยุตฺตญฺจ, ถินํ อุทฺธจฺเจน กิเลโส เจว กิเลส\-สมฺปยุตฺตญฺจ, อุทฺธจฺจํ ถิเนน กิเลโส เจว กิเลส\-สมฺปยุตฺตญฺจ, อหิริกํ อุทฺธจฺเจน กิเลโส เจว กิเลส\-สมฺปยุตฺตญฺจ, อุทฺธจฺจํ อหิริเกน กิเลโส เจว กิเลส\-สมฺปยุตฺตญฺจ, อโนตฺตปฺปํ อุทฺธจฺเจน กิเลโส เจว กิเลส\-สมฺปยุตฺตญฺจ, อุทฺธจฺจํ อโนตฺตปฺเปน กิเลโส เจว กิเลส\-สมฺปยุตฺตญฺจ, โลโภ อหิริเกน กิเลโส เจว กิเลส\-สมฺปยุตฺโต จ, อหิริกํ โลเภน กิเลโส เจว กิเลส\-สมฺปยุตฺตญฺจ, โทโส อหิริเกน กิเลโส เจว กิเลส\-สมฺปยุตฺโต จ, อหิริกํ โทเสน กิเลโส เจว กิเลส\-สมฺปยุตฺตญฺจ, โมโห อหิริเกน กิเลโส เจว กิเลส\-สมฺปยุตฺโต จ, อหิริกํ โมเหน กิเลโส เจว กิเลส\-สมฺปยุตฺตญฺจ, มาโน อหิริเกน กิเลโส เจว กิเลส\-สมฺปยุตฺโต จ, อหิริกํ มาเนน กิเลโส เจว กิเลส\-สมฺปยุตฺตญฺจ, ทิฏฺฐิ อหิริเกน กิเลโส เจว กิเลส\-สมฺปยุตฺตา จ, อหิริกํ ทิฏฺฐิยา กิเลโส เจว กิเลส\-สมฺปยุตฺตญฺจ, วิจิกิจฺฉา อหิริเกน กิเลโส เจว กิเลส\-สมฺปยุตฺตา จ, อหิริกํ วิจิกิจฺฉาย กิเลโส เจว กิเลส\-สมฺปยุตฺตญฺจ, ถินํ อหิริเกน กิเลโส เจว กิเลส\-สมฺปยุตฺตญฺจ, อหิริกํ ถิเนน กิเลโส เจว กิเลส\-สมฺปยุตฺตญฺจ, อุทฺธจฺจํ อหิริเกน กิเลโส เจว กิเลส\-สมฺปยุตฺตญฺจ, อหิริกํ อุทฺธจฺเจน กิเลโส เจว กิเลส\-สมฺปยุตฺตญฺจ, อโนตฺตปฺปํ อหิริเกน กิเลโส เจว กิเลส\-สมฺปยุตฺตญฺจ, อหิริกํ อโนตฺตปฺเปน กิเลโส เจว กิเลส\-สมฺปยุตฺตญฺจ, โลโภ อโนตฺตปฺเปน กิเลโส เจว กิเลส\-สมฺปยุตฺโต จ, อโนตฺตปฺปํ โลเภน กิเลโส เจว \csromanfolio{249}กิเลส\-สมฺปยุตฺตญฺจ, โทโส อโนตฺตปฺเปน กิเลโส เจว กิเลส\-สมฺปยุตฺโต จ, อโนตฺตปฺปํ โทเสน กิเลโส เจว กิเลส\-สมฺปยุตฺตญฺจ, โมโห อโนตฺตปฺเปน กิเลโส เจว กิเลส\-สมฺปยุตฺโต จ, อโนตฺตปฺปํ โมเหน กิเลโส เจว กิเลส\-สมฺปยุตฺตญฺจ, มาโน อโนตฺตปฺเปน กิเลโส เจว กิเลส\-สมฺปยุตฺโต จ, อโนตฺตปฺปํ มาเนน กิเลโส เจว กิเลส\-สมฺปยุตฺตญฺจ, ทิฏฺฐิ อโนตฺตปฺเปน กิเลโส เจว กิเลส\-สมฺปยุตฺตา จ, อโนตฺตปฺปํ ทิฏฺฐิยา กิเลโส เจว กิเลส\-สมฺปยุตฺตญฺจ, วิจิกิจฺฉา อโนตฺตปฺเปน กิเลโส เจว กิเลส\-สมฺปยุตฺตา จ, อโนตฺตปฺปํ วิจิกิจฺฉาย กิเลโส เจว กิเลส\-สมฺปยุตฺตญฺจ, ถินํ อโนตฺตปฺเปน กิเลโส เจว กิเลส\-สมฺปยุตฺตญฺจ, อโนตฺตปฺปํ ถิเนน กิเลโส เจว กิเลส\-สมฺปยุตฺตญฺจ, อุทฺธจฺจํ อโนตฺตปฺเปน กิเลโส เจว กิเลส\-สมฺปยุตฺตญฺจ, อโนตฺตปฺปํ อุทฺธจฺเจน กิเลโส เจว กิเลส\-สมฺปยุตฺตญฺจ, อหิริกํ อโนตฺตปฺเปน กิเลโส เจว กิเลส\-สมฺปยุตฺตญฺจ, อโนตฺตปฺปํ อหิริเกน กิเลโส เจว กิเลส\-สมฺปยุตฺตญฺจ. อิเม ธมฺมา กิเลสา เจว กิเลส\-สมฺปยุตฺตา จ.}

\proseitem{1258}{กตเม ธมฺมา กิเลส\-สมฺปยุตฺตา เจว โน จ กิเลสา. เตหิ ธมฺเมหิ เย ธมฺมา สมฺปยุตฺตา, เต ธมฺเม ฐเปตฺวา เวทนากฺขนฺโธ ฯเปฯ วิญฺญาณ\-กฺขนฺโธ. อิเม ธมฺมา กิเลส\-สมฺปยุตฺตา เจว โน จ กิเลสา.}

\proseitem{1259}{กตเม ธมฺมา กิเลส\-วิปฺปยุตฺตา สํกิเลสิกา. เตหิ ธมฺเมหิ เย ธมฺมา วิปฺปยุตฺตา สาสวา กุสลาพฺยากตา ธมฺมา กามาวจรา รูปาวจรา อรูปาวจรา รูปกฺขนฺโธ ฯเปฯ วิญฺญาณ\-กฺขนฺโธ. อิเม ธมฺมา กิเลส\-วิปฺปยุตฺตา สํกิเลสิกา.}

\proseitem{1260}{กตเม ธมฺมา กิเลส\-วิปฺปยุตฺตา อสํกิเลสิกา. อปริยาปนฺนา มคฺคา จ มคฺคผลานิ จ อสงฺขตา จ ธาตุ. อิเม ธมฺมา กิเลส\-วิปฺปยุตฺตา อสํกิเลสิกา.}

\csromanmatikaanchor{mtk.116}
\csromantocmarkref{section}{ปิฏฺฐิทุก}{mtk.116}
\bookmark[dest={mtk.116},level=1]{ปิฏฺฐิทุก}
\csromanheader{1}{\textbf{ปิฏฺฐิทุก}}
\proseitem{1261}{กตเม ธมฺมา ทสฺสเนน ปหาตพฺพา. ตีณิ สญฺโญชนานิ สกฺกายทิฏฺฐิ วิจิกิจฺฉา สีลพฺพต\-ปรามาโส.}

\proseitem{1262}{\csromanfolio{250}ตตฺถ กตมา สกฺกายทิฏฺฐิ. อิธ อสฺสุตวา ปุถุชฺชโน อริยานํ อทสฺสาวี อริยธมฺมสฺส อโกวิโท อริยธมฺเม อวินีโต สปฺปุริสานํ อทสฺสาวี สปฺปุริส\-ธมฺมสฺส อโกวิโท สปฺปุริส\-ธมฺเม อวินีโต รูปํ อตฺตโต สมนุปสฺสติ, รูปวนฺตํ วา อตฺตานํ, อตฺตนิ วา รูปํ, รูปสฺมิํ วา อตฺตานํ. เวทนํ ฯเปฯ. สญฺญํ ฯเปฯ. สงฺขาเร ฯเปฯ. วิญฺญาณํ อตฺตโต สมนุปสฺสติ, วิญฺญาณวนฺตํ วา อตฺตานํ, อตฺตนิ วา วิญฺญาณํ, วิญฺญาณสฺมิํ วา อตฺตานํ. ยา เอวรูปา ทิฏฺฐิ ทิฏฺฐิคตํ ฯเปฯ วิปริยาส\-คฺคาโห. อยํ วุจฺจติ สกฺกายทิฏฺฐิ.}

\proseitem{1263}{ตตฺถ กตมา วิจิกิจฺฉา. สตฺถริ กงฺขติ วิจิกิจฺฉติ ฯเปฯ ถมฺภิตตฺตํ จิตฺตสฺส มโนวิเลโข. อยํ วุจฺจติ วิจิกิจฺฉา.}

\proseitem{1264}{ตตฺถ กตโม สีลพฺพต\-ปรามาโส. “อิโต พหิทฺธา สมณ\-พฺราหฺมณานํ สีเลน สุทฺธิ วเตน สุทฺธิ สีลพฺพเตน สุทฺธี”ติ ยา เอวรูปา ทิฏฺฐิ ทิฏฺฐิคตํ ฯเปฯ วิปริยาส\-คฺคาโห. อยํ วุจฺจติ สีลพฺพต\-ปรามาโส. อิมานิ ตีณิ สญฺโญชนานิ, ตเทกฏฺฐา จ กิเลสา ตํสมฺปยุตฺโต เวทนากฺขนฺโธ ฯเปฯ วิญฺญาณ\-กฺขนฺโธ, ตํสมุฏฺฐานํ กายกมฺมํ วจีกมฺมํ มโนกมฺมํ. อิเม ธมฺมา ทสฺสเนน ปหาตพฺพา.}

\proseitem{1265}{กตเม ธมฺมา น ทสฺสเนน ปหาตพฺพา. เต ธมฺเม ฐเปตฺวา อวเสสา กุสลากุสลา\-พฺยากตา ธมฺมา กามาวจรา รูปาวจรา อรูปาวจรา อปริยาปนฺนา เวทนากฺขนฺโธ ฯเปฯ วิญฺญาณ\-กฺขนฺโธ สพฺพญฺจ รูปํ อสงฺขตา จ ธาตุ. อิเม ธมฺมา น ทสฺสเนน ปหาตพฺพา.}

\proseitem{1266}{กตเม ธมฺมา ภาวนาย ปหาตพฺพา. อวเสโส โลโภ โทโส โมโห, ตเทกฏฺฐา จ กิเลสา ตํสมฺปยุตฺโต เวทนากฺขนฺโธ ฯเปฯ วิญฺญาณ\-กฺขนฺโธ, ตํสมุฏฺฐานํ กายกมฺมํ วจีกมฺมํ มโนกมฺมํ. อิเม ธมฺมา ภาวนาย ปหาตพฺพา.}

\proseitem{1267}{กตเม ธมฺมา น ภาวนาย ปหาตพฺพา. เต ธมฺเม ฐเปตฺวา อวเสสา กุสลากุสลา\-พฺยากตา ธมฺมา กามาวจร รูปวจรา อรูปาวจรา อปริยาปนฺนา เวทนากฺขนฺโธ ฯเปฯ วิญฺญาณ\-กฺขนฺโธ สพฺพญฺจ รูปํ อสงฺขตา จ ธาตุ. อิเม ธมฺมา น ภาวนาย ปหาตพฺพา.}

\chapterheadpageread{3. นิกฺเขปกณฺฑ}
\proseitem{1268}{\csromanfolio{251}กตเม ธมฺมา ทสฺสเนน ปหาตพฺพ\-เหตุกา. ตีณิ สญฺโญชนานิ สกฺกายทิฏฺฐิ วิจิกิจฺฉา สีลพฺพต\-ปรามาโส.}

\proseitem{1269}{ตตฺถ กตมา สกฺกายทิฏฺฐิ. อิธ อสฺสุตวา ปุถุชฺชโน อริยานํ อทสฺสาวี อริยธมฺมสฺส อโกวิโท อริยธมฺเม อวินีโต สปฺปุริสานํ อทสฺสาวี สปฺปุริส\-ธมฺมสฺส อโกวิโท สปฺปุริส\-ธมฺเม อวินีโต รูปํ อตฺตโต สมนุปสฺสติ, รูปวนฺตํ วา อตฺตานํ, อตฺตนิ วา รูปํ, รูปสฺมิํ วา อตฺตานํ. เวทนํ ฯเปฯ. สญฺญํ ฯเปฯ. สงฺขาเร ฯเปฯ. วิญฺญาณํ อตฺตโต สมนุปสฺสติ, วิญฺญาณวนฺตํ วา อตฺตานํ, อตฺตนิ วา วิญฺญาณํ, วิญฺญาณสฺมิํ วา อตฺตานํ. ยา เอวรูปา ทิฏฺฐิ ทิฏฺฐิคตํ ฯเปฯ วิปริยาสคฺคโห. อยํ วุจฺจติ สกฺกายทิฏฺฐิ.}

\proseitem{1270}{ตตฺถ กตมา วิจิกิจฺฉา. สตฺถริ กงฺขติ วิจิกิจฺฉติ ฯเปฯ ถมฺภิตตฺตํ จิตฺตสฺส มโนวิเลโข. อยํ วุจฺจติ วิจิกิจฺฉา.}

\proseitem{1271}{ตตฺถ กตโม สีลพฺพต\-ปรามาโส. “อิโต พหิทฺธา สมณ\-พฺราหฺมณานํ สีเลน สุทฺธิ วเตน สุทฺธิ สีลพฺพเตน สุทฺธี”ติ ยา เอวรูปา ทิฏฺฐิ ทิฏฺฐิคตํ ฯเปฯ วิปริยาส\-คฺคาโห. อยํ วุจฺจติ สีลพฺพต\-ปรามาโส. อิมานิ ตีณิ สํโยชนานิ, ตเทกฏฺฐา จ กิเลสา ตํสมฺปยุตฺโต เวทนากฺขนฺโธ ฯเปฯ วิญฺญาณ\-กฺขนฺโธ, ตํสมุฏฺฐานํ กายกมฺมํ วจีกมฺมํ มโนกมฺมํ, อิเม ธมฺมา ทสฺสเนน ปหาตพฺพ\-เหตุกา. ตีณิ สญฺโญชนานิ สกฺกายทิฏฺฐิ วิจิกิจฺฉา สีลพฺพต\-ปรามาโส, อิเม ธมฺมา ทสฺสเนน ปหาตพฺพา. ตเทกฏฺโฐ โลโภ โทโส โมโห, อิเม ธมฺมา ทสฺสเนน ปหาตพฺพเหตู. ตเทกฏฺฐา จ กิเลสา ตํสมฺปยุตฺโต เวทนากฺขนฺโธ ฯเปฯ วิญฺญาณ\-กฺขนฺโธ, ตํสมุฏฺฐานํ กายกมฺมํ วจีกมฺมํ มโนกมฺมํ. อิเม ธมฺมา ทสฺสเนน ปหาตพฺพ\-เหตุกา.}

\proseitem{1272}{กตเม ธมฺมา น ทสฺสเนน ปหาตพฺพ\-เหตุกา. เต ธมฺเม ฐเปตฺวา อวเสสา กุสลากุสลา\-พฺยากตา ธมฺมา กามาวจรา รูปาวจรา อรูปาวจรา อปริยาปนฺนา เวทนากฺขนฺโธ ฯเปฯ วิญฺญาณ\-กฺขนฺโธ สพฺพญฺจ รูปํ อสงฺขตา จ ธาตุ. อิเม ธมฺมา น ทสฺสเนน ปหาตพฺพ\-เหตุกา.}

\proseitem{1273}{กตเม ธมฺมา ภาวนาย ปหาตพฺพ\-เหตุกา. อวเสโส โลโภ โทโส โมโห, อิเม ธมฺมา ภาวนาย ปหาตพฺพเหตู. \csromanfolio{252}ตเทกฏฺฐา จ กิเลสา ตํสมฺปยุตฺโต เวทนากฺขนฺโธ ฯเปฯ วิญฺญาณ\-กฺขนฺโธ, ตํสมุฏฺฐานํ กายกมฺมํ วจีกมฺมํ มโนกมฺมํ. อิเม ธมฺมา ภาวนาย ปหาตพฺพ\-เหตุกา.}

\proseitem{1274}{กตเม ธมฺมา น ภาวนาย ปหาตพฺพ\-เหตุกา. เต ธมฺเม ฐเปตฺวา อวเสสา กุสลากุสลา\-พฺยากตา ธมฺมา กามาวจรา รูปาวจรา อรูปาวจรา อปริยาปนฺนา เวทนากฺขนฺโธ ฯเปฯ วิญฺญาณ\-กฺขนฺโธ สพฺพญฺจ รูปํ อสงฺขตา จ ธาตุ. อิเม ธมฺมา น ภาวนาย ปหาตพฺพ\-เหตุกา.}

\proseitem{1275}{กตเม ธมฺมา สวิตกฺกา. สวิตกฺก\-ภูมิยํ กามาวจเร รูปาวจเร อปริยาปนฺเน วิตกฺกํ ฐเปตฺวา ตํสมฺปยุตฺโต เวทนากฺขนฺโธ ฯเปฯ วิญฺญาณ\-กฺขนฺโธ. อิเม ธมฺมา สวิตกฺกา.}

\proseitem{1276}{กตเม ธมฺมา อวิตกฺกา. อวิตกฺก\-ภูมิยํ กามาวจเร รูปาวจเร อรูปาวจเร อปริยาปนฺเน เวทนากฺขนฺโธ ฯเปฯ วิญฺญาณ\-กฺขนฺโธ วิตกฺโก จ สพฺพญฺจ รูปํ อสงฺขตา จ ธาตุ. อิเม ธมฺมา อวิตกฺกา.}

\proseitem{1277}{กตเม ธมฺมา สวิจารา. สวิจาร\-ภูมิยํ กามาวจเร รูปาวจเร อปริยาปนฺเน วิจารํ ฐเปตฺวา ตํสมฺปยุตฺโต เวทนากฺขนฺโธ สญฺญากฺขนฺโธ สงฺขาร\-กฺขนฺโธ วิญฺญาณ\-กฺขนฺโธ. อิเม ธมฺมา สวิจารา.}

\proseitem{1278}{กตเม ธมฺมา อวิจารา. อวิจารภูมิยํ กามาวจเร รูปาวจเร อรูปาวจเร อปริยาปนฺเน เวทนากฺขนฺโธ ฯเปฯ วิญฺญาณ\-กฺขนฺโธ วิจาโร จ สพฺพญฺจ รูปํ อสงฺขตา จ ธาตุ. อิเม ธมฺมา อวิจารา.}

\proseitem{1279}{กตเม ธมฺมา สปฺปีติกา. สปฺปีติก\-ภูมิยํ กามาวจเร รูปาวจเร อปริยาปนฺเน ปีติํ ฐเปตฺวา ตํสมฺปยุตฺโต เวทนากฺขนฺโธ ฯเปฯ วิญฺญาณ\-กฺขนฺโธ. อิเม ธมฺมา สปฺปีติกา.}

\proseitem{1280}{กตเม ธมฺมา อปฺปีติกา. อปฺปีติกภูมิยํ กามาวจเร รูปาวจเร อรูปาวจเร อปริยาปนฺเน เวทนากฺขนฺโธ ฯเปฯ วิญฺญาณ\-กฺขนฺโธ ปีติ จ สพฺพญฺจ รูปํ อสงฺขตา จ ธาตุ. อิเม ธมฺมา อปฺปีติกา.}

\proseitem{1281}{กตเม ธมฺมา ปีติสหคตา. ปีติภูมิยํ กามาวจเร รูปาวจเร อปริยาปนฺเน ปีติํ ฐเปตฺวา ตํสมฺปยุตฺโต เวทนากฺขนฺโธ ฯเปฯ วิญฺญาณ\-กฺขนฺโธ. อิเม ธมฺมา ปีติสหคตา.}

\chapterheadpageread{3. นิกฺเขปกณฺฑ}
\proseitem{1282}{\csromanfolio{253}กตเม ธมฺมา น ปีติสหคตา. น ปีติภูมิยํ กามาวจเร รูปาวจเร อรูปาวจเร อปริยาปนฺเน เวทนากฺขนฺโธ ฯเปฯ วิญฺญาณ\-กฺขนฺโธ ปีติ จ สพฺพญฺจ รูปํ อสงฺขตา จ ธาตุ. อิเม ธมฺมา น ปีติสหคตา.}

\proseitem{1283}{กตเม ธมฺมา สุขสหคตา. สุขภูมิยํ กามาวจเร รูปาวจเร อปริยาปนฺเน สุขํ ฐเปตฺวา ตํสมฺปยุตฺโต สญฺญากฺขนฺโธ สงฺขาร\-กฺขนฺโธ วิญฺญาณ\-กฺขนฺโธ. อิเม ธมฺมา สุขสหคตา.}

\proseitem{1284}{กตเม ธมฺมา น สุขสหคตา. น สุขภูมิยํ กามาวจเร รูปาวจเร อรูปาวจเร อปริยาปนฺเน เวทนากฺขนฺโธ ฯเปฯ วิญฺญาณ\-กฺขนฺโธ สุขญฺจ สพฺพญฺจ รูปํ อสงฺขตา จ ธาตุ. อิเม ธมฺมา น สุขสหคตา.}

\proseitem{1285}{กตเม ธมฺมา อุเปกฺขา\-สหคตา. อุเปกฺขา\-ภูมิยํ กามาวจเร รูปาวจเร อรูปาวจเร อปริยาปนฺเน อุเปกฺขํ ฐเปตฺวา ตํสมฺปยุตฺโต สญฺญากฺขนฺโธ สงฺขาร\-กฺขนฺโธ วิญฺญาณ\-กฺขนฺโธ. อิเม ธมฺมา อุเปกฺขา\-สหคตา.}

\proseitem{1286}{กตเม ธมฺมา น อุเปกฺขา\-สหคตา. น อุเปกฺขา\-ภูมิยํ กามาวจเร รูปาวจเร อปริยาปนฺเน เวทนากฺขนฺโธ ฯเปฯ วิญฺญาณ\-กฺขนฺโธ อุเปกฺขา จ สพฺพญฺจ รูปํ อสงฺขตา จ ธาตุ. อิเม ธมฺมา น อุเปกฺขา\-สหคตา.}

\proseitem{1287}{กตเม ธมฺมา กามาวจรา. เหฏฺฐโต อวิจินิรยํ ปริยนฺตํ กริตฺวา อุปริโต ปรนิมฺมิต\-วส\-วตฺตี เทเว\csromansharedfootnote{262a37cf8af809ef}{ปรนิมฺมิตวสวตฺติเทเว (สี, ก)} อนฺโต กริตฺวา ยํ เอตสฺมิํ อนฺตเร เอตฺถาวจรา เอตฺถ ปริยาปนฺนา ขนฺธธาตุ อายตนา รูปํ เวทนา สญฺญา สงฺขารา วิญฺญาณํ. อิเม ธมฺมา กามาวจรา.}

\proseitem{1288}{กตเม ธมฺมา น กามาวจรา. รูปาวจรา อรูปาวจรา อปริยาปนฺนา. อิเม ธมฺมา น กามาวจรา.}

\proseitem{1289}{กตเม ธมฺมา รูปาวจรา. เหฏฺฐโต พฺรหฺมโลกํ ปริยนฺตํ กริตฺวา อุปริโต อกนิฏฺเฐ เทเว\csromansharedfootnote{bc6122d71f0103bb}{อกนิฏฺฐเทเว (สี, ก)} อนฺโต กริตฺวา ยํ เอตสฺมิํ อนฺตเร เอตฺถาวจรา เอตฺถ ปริยาปนฺนา สมาปนฺนสฺส วา อุปปนฺนสฺส วา ทิฏฺฐิธมฺมสุขวิหาริสฺส\csromansharedfootnote{25e8f3677728d90f}{ทิฏฺฐธมฺมสุขวิหารสฺส (ก)} วา จิตฺตเจตสิกา ธมฺมา. อิเม ธมฺมา รูปาวจรา.}

\proseitem{1290}{\csromanfolio{254}กตเม ธมฺมา น รูปาวจรา. กามาวจรา อรูปาวจรา อปริยาปนฺนา. อิเม ธมฺมา น รูปาวจรา.}

\proseitem{1291}{กตเม ธมฺมา อรูปาวจรา. เหฏฺฐโต อากาสานญฺจายตนุ\-ปเค เทเว ปริยนฺตํ กริตฺวา อุปริโต เนวสญฺญานาสญฺญายตนุ\-ปเค เทเว อนฺโต กริตฺวา ยํ เอตสฺมิํ อนฺตเร เอตฺถาวจรา เอตฺถ ปริยาปนฺนา สมาปนฺนสฺส วา อุปปนฺนสฺส วา ทิฏฺฐธมฺม\-สุขวิหาริสฺส วา จิตฺตเจตสิกา ธมฺมา. อิเม ธมฺมา อรูปาวจรา.}

\proseitem{1292}{กตเม ธมฺมา น อรูปาวจรา. กามาวจรา รูปาวจรา อปริยาปนฺนา. อิเม ธมฺมา น อรูปาวจรา.}

\proseitem{1293}{กตเม ธมฺมา ปริยาปนฺนา. สาสวา กุสลากุสลา\-พฺยากตา ธมฺมา กามาวจรา รูปาวจรา อรูปาวจรา รูปกฺขนฺโธ ฯเปฯ วิญฺญาณ\-กฺขนฺโธ. อิเม ธมฺมา ปริยาปนฺนา.}

\proseitem{1294}{กตเม ธมฺมา อปริยาปนฺนา. มคฺคา จ มคฺคผลานิ จ อสงฺขตา จ ธาตุ. อิเม ธมฺมา อปริยาปนฺนา.}

\proseitem{1295}{กตเม ธมฺมา นิยฺยานิกา. จตฺตาโร มคฺคา อปริยาปนฺนา. อิเม ธมฺมา นิยฺยานิกา.}

\proseitem{1296}{กตเม ธมฺมา อนิยฺยานิกา. เต ธมฺเม ฐเปตฺวา อวเสสา กุสลากุสลา\-พฺยากตา ธมฺมา กามาวจรา รูปาวจรา อรูปาวจรา อปริยาปนฺนา เวทนากฺขนฺโธ ฯเปฯ วิญฺญาณ\-กฺขนฺโธ สพฺพญฺจ รูปํ อสงฺขตา จ ธาตุ. อิเม ธมฺมา อนิยฺยานิกา.}

\proseitem{1297}{กตเม ธมฺมา นิยตา. ปญฺจ กมฺมานิ อานนฺตริกานิ ยา จ มิจฺฉาทิฏฺฐิ นิยตา จตฺตาโร มคฺคา อปริยาปนฺนา. อิเม ธมฺมา นิยตา.}

\proseitem{1298}{กตเม ธมฺมา อนิยตา. เต ธมฺเม ฐเปตฺวา อวเสสา กุสลากุสลา\-พฺยากตา ธมฺมา กามาวจรา รูปาวจรา อรูปาวจรา อปริยาปนฺนา เวทนากฺขนฺโธ ฯเปฯ วิญฺญาณ\-กฺขนฺโธ สพฺพญฺจ รูปํ อสงฺขตา จ ธาตุ. อิเม ธมฺมา อนิยตา.}

\chapterheadpageread{3. นิกฺเขปกณฺฑ}
\proseitem{1299}{\csromanfolio{255}กตเม ธมฺมา สอุตฺตรา. สาสวา กุสลากุสลา\-พฺยากตา ธมฺมา กามาวจรา รูปาวจรา อรูปาวจรา รูปกฺขนฺโธ ฯเปฯ วิญฺญาณ\-กฺขนฺโธ. อิเม ธมฺมา สอุตฺตรา.}

\proseitem{1300}{กตเม ธมฺมา อนุตฺตรา. อปริยาปนฺนา มคฺคา จ มคฺคผลานิ จ อสงฺขตา จ ธาตุ. อิเม ธมฺมา อนุตฺตรา.}

\proseitem{1301}{กตเม ธมฺมา สรณา. ตีณิ อกุสลมูลานิ โลโภ โทโส โมโห, ตเทกฏฺฐา จ กิเลสา ตํสมฺปยุตฺโต เวทนากฺขนฺโธ ฯเปฯ วิญฺญาณ\-กฺขนฺโธ, ตํสมุฏฺฐานํ กายกมฺมํ วจีกมฺมํ มโนกมฺมํ. อิเม ธมฺมา สรณา.}

\proseitem{1302}{กตเม ธมฺมา อรณา. กุสลาพฺยากตา ธมฺมา กามาวจรา รูปาวจรา อรูปาวจรา อปริยาปนฺนา เวทนากฺขนฺโธ ฯเปฯ วิญฺญาณ\-กฺขนฺโธ สพฺพญฺจ รูปํ อสงฺขตา จ ธาตุ. อิเม ธมฺมา อรณา.}

\vspace{\tipitakaparskip}
\csromancenter{อภิธมฺมทุกํ.}
\csromanmatikaanchor{mtk.117}
\csromantocmarkref{section}{สุตฺตนฺติกทุกนิกฺเขป}{mtk.117}
\bookmark[dest={mtk.117},level=1]{สุตฺตนฺติกทุกนิกฺเขป}
\csromanheader{1}{สุตฺตนฺติก\-ทุก\-นิกฺเขป}
\proseitem{1303}{กตเม ธมฺมา วิชฺชาภาคิโน. วิชฺชาย สมฺปยุตฺตกา ธมฺมา. อิเม ธมฺมา วิชฺชาภาคิโน.}

\proseitem{1304}{กตเม ธมฺมา อวิชฺชาภาคิโน. อวิชฺชาย สมฺปยุตฺตกา ธมฺมา. อิเม ธมฺมา อวิชฺชาภาคิโน.}

\proseitem{1305}{กตเม ธมฺมา วิชฺชูปมา. เหฏฺฐิเมสุ ตีสุ อริยมคฺเคสุ ปญฺญา. อิเม ธมฺมา วิชฺชูปมา.}

\proseitem{1306}{กตเม ธมฺมา วชิรูปมา. อุปริฏฺฐิเม อรหตฺตมคฺเค ปญฺญา. อิเม ธมฺมา วชิรูปมา.}

\proseitem{1307}{กตเม ธมฺมา พาลา. อหิรีกญฺจ อโนตฺตปฺปญฺจ. อิเม ธมฺมา พาลา. สพฺเพปิ อกุสลา ธมฺมา พาลา.}

\proseitem{1308}{\csromanfolio{256}กตเม ธมฺมา ปณฺฑิตา. หิรี จ โอตฺตปฺปญฺจ. อิเม ธมฺมา ปณฺฑิตา. สพฺเพปิ กุสลา ธมฺมา ปณฺฑิตา.}

\proseitem{1309}{กตเม ธมฺมา กณฺหา. อหิรีกญฺจ อโนตฺตปฺปญฺจ. อิเม ธมฺมา กณฺหา. สพฺเพปิ อกุสลา ธมฺมา กณฺหา.}

\proseitem{1310}{กตเม ธมฺมา สุกฺกา. หิรี จ โอตฺตปฺปญฺจ. อิเม ธมฺมา สุกฺกา. สพฺเพปิ กุสลา ธมฺมา สุกฺกา.}

\proseitem{1311}{กตเม ธมฺมา ตปนียา. กายทุจฺจริตํ วจีทุจฺจริตํ มโนทุจฺจริตํ. อิเม ธมฺมา ตปนียา. สพฺเพปิ อกุสลา ธมฺมา ตปนียา.}

\proseitem{1312}{กตเม ธมฺมา อตปนียา. กายสุจริตํ วจีสุจริตํ มโนสุจริตํ. อิเม ธมฺมา อตปนียา. สพฺเพปิ กุสลา ธมฺมา อตปนียา.}

\proseitem{1313}{กตเม ธมฺมา อธิวจนา. ยา เตสํ เตสํ ธมฺมานํ สงฺขา สมญฺญา ปญฺญตฺติ โวหาโร นามํ นามกมฺมํ นามเธยฺยํ นิรุตฺติ พฺยญฺชนํ อภิลาโป. อิเม ธมฺมา อธิวจนา. สพฺเพว ธมฺมา อธิวจน\-ปถา.}

\proseitem{1314}{กตเม ธมฺมา นิรุตฺติ. ยา เตสํ เตสํ ธมฺมานํ สงฺขา สมญฺญา ปญฺญตฺติ โวหาโร นามํ นามกมฺมํ นามเธยฺยํ นิรุตฺติ พฺยญฺชนํ อภิลาโป. อิเม ธมฺมา นิรุตฺติ. สพฺเพว ธมฺมา นิรุตฺติปถา.}

\proseitem{1315}{กตเม ธมฺมา ปญฺญตฺติ. ยา เตสํ เตสํ ธมฺมานํ สงฺขา สมญฺญา ปญฺญตฺติ โวหาโร นามํ นามกมฺมํ นามเธยฺยํ นิรุตฺติ พฺยญฺชนํ อภิลาโป. อิเม ธมฺมา ปญฺญตฺติ. สพฺเพว ธมฺมา ปญฺญตฺติปถา.}

\proseitem{1316}{ตตฺถ กตมํ นามํ. เวทนากฺขนฺโธ สญฺญากฺขนฺโธ สงฺขาร\-กฺขนฺโธ วิญฺญาณ\-กฺขนฺโธ อสงฺขตา จ ธาตุ. อิทํ วุจฺจติ นามํ.}

\proseitem{1317}{ตตฺถ กตมํ รูปํ. จตฺตาโร จ มหาภูตา จตุนฺนญฺจ มหาภูตานํ อุปาทาย รูปํ. อิทํ วุจฺจติ รูปํ.}

\proseitem{1318}{ตตฺถ กตมา อวิชฺชา. ยํ อญฺญาณํ อทสฺสนํ ฯเปฯ อวิชฺชาลงฺคี โมโห อกุสลมูลํ. อยํ วุจฺจติ อวิชฺชา.}

\chapterheadpageread{3. นิกฺเขปกณฺฑ}
\proseitem{1319}{\csromanfolio{257}ตตฺถ กตมา ภวตณฺหา. โย ภเวสุ ภวฉนฺโท ฯเปฯ ภวชฺโฌสานํ. อยํ วุจฺจติ ภวตณฺหา.}

\proseitem{1320}{ตตฺถ กตมา ภวทิฏฺฐิ. “ภวิสฺสติ อตฺตา จ โลโก จา”ติ ยา เอวรูปา ทิฏฺฐิ ทิฏฺฐิคตํ ฯเปฯ วิปริยาส\-คฺคาโห. อยํ วุจฺจติ ภวทิฏฺฐิ.}

\proseitem{1321}{ตตฺถ กตมา วิภวทิฏฺฐิ. “น ภวิสฺสติ อตฺตา จ โลโก จา”ติ ยา เอวรูปา ทิฏฺฐิ ทิฏฺฐิคตํ ฯเปฯ วิปริยาส\-คฺคาโห. อยํ วุจฺจติ วิภวทิฏฺฐิ.}

\proseitem{1322}{ตตฺถ กตมา สสฺสตทิฏฺฐิ. “สสฺสโต อตฺตา จ โลโก จา”ติ ยา เอวรูปา ทิฏฺฐิ ทิฏฺฐิคตํ ฯเปฯ วิปริยาส\-คฺคาโห. อยํ วุจฺจติ สสฺสตทิฏฺฐิ.}

\proseitem{1323}{ตตฺถ กตมา อุจฺเฉททิฏฺฐิ. “อุจฺฉิชฺชิสฺสติ อตฺตา จ โลโก จา”ติ ยา เอวรูปา ทิฏฺฐิ ทิฏฺฐิคตํ ฯเปฯ วิปริยาส\-คฺคาโห. อยํ วุจฺจติ อุจฺเฉททิฏฺฐิ.}

\proseitem{1324}{ตตฺถ กตมา อนฺตวา ทิฏฺฐิ. “อนฺตวา อตฺตา จ โลโก จา”ติ ยา เอวรูปา ทิฏฺฐิ ทิฏฺฐิคตํ ฯเปฯ วิปริยาส\-คฺคาโห. อยํ วุจฺจติ อนฺตวา ทิฏฺฐิ.}

\proseitem{1325}{ตตฺถ กตมา อนนฺตวา ทิฏฺฐิ. “อนนฺตวา อตฺตา จ โลโก จา”ติ ยา เอวรูปา ทิฏฺฐิ ทิฏฺฐิคตํ ฯเปฯ วิปริยาส\-คฺคาโห. อยํ วุจฺจติ อนนฺตวา ทิฏฺฐิ.}

\proseitem{1326}{ตตฺถ กตมา ปุพฺพนฺตา\-นุทิฏฺฐิ. ปุพฺพนฺตํ อารพฺภ ยา อุปฺปชฺชติ ทิฏฺฐิ ทิฏฺฐิคตํ ฯเปฯ วิปริยาส\-คฺคาโห. อยํ วุจฺจติ ปุพฺพนฺตา\-นุทิฏฺฐิ.}

\proseitem{1327}{ตตฺถ กตมา อปรนฺตา\-นุทิฏฺฐิ. อปรนฺตํ อารพฺภ ยา อุปฺปชฺชติ ทิฏฺฐิ ทิฏฺฐิคตํ ฯเปฯ วิปริยาส\-คฺคาโห. อยํ วุจฺจติ อปรนฺตา\-นุทิฏฺฐิ.}

\proseitem{1328}{ตตฺถ กตมํ อหิริกํ. ยํ น หิรียติ หิริยิตพฺเพน น หิรียติ ปาปกานํ อกุสลานํ ธมฺมานํ สมาปตฺติยา. อิทํ วุจฺจติ อหิริกํ.}

\proseitem{1329}{ตตฺถ กตมํ อโนตฺตปฺปํ. ยํ น โอตฺตปฺปติ โอตฺตปฺปิตพฺเพน น โอตฺตปฺปติ ปาปกานํ อกุสลานํ ธมฺมานํ สมาปตฺติยา. อิทํ วุจฺจติ อโนตฺตปฺปํ.}

\proseitem{1330}{\csromanfolio{258}ตตฺถ กตมา หิรี. ยํ หิรียติ หิริยิตพฺเพน หิรียติ ปาปกานํ อกุสลานํ ธมฺมานํ สมาปตฺติยา. อยํ วุจฺจติ หิรี.}

\proseitem{1331}{ตตฺถ กตมํ โอตฺตปฺปํ. ยํ โอตฺตปฺปติ โอตฺตปฺปิตพฺเพน โอตฺตปฺปติ ปาปกานํ อกุสลานํ ธมฺมานํ สมาปตฺติยา. อิทํ วุจฺจติ โอตฺตปฺปํ.}

\proseitem{1332}{ตตฺถ กตมา โทวจสฺสตา. สหธมฺมิเก วุจฺจมาเน โทวจสฺสายํ โทวจสฺสิยํ โทวจสฺสตา วิปฺปฏิกูล\-คฺคาหิตา วิปจฺจนีก\-สาตตา อนาทริยํ อนาทรตา อคารวตา อปฺปฏิสฺสวตา. อยํ วุจฺจติ โทวจสฺสตา.}

\proseitem{1333}{ตตฺถ กตมา ปาปมิตฺตตา. เย เต ปุคฺคลา อสฺสทฺธา ทุสฺสีลา อปฺปสฺสุตา มจฺฉริโน ทุปฺปญฺญา. ยา เตสํ เสวนา นิเสวนา สํเสวนา ภชนา สมฺภชนา ภตฺติ สมฺภตฺติ ตํสมฺปวงฺกตา. อยํ วุจฺจติ ปาปมิตฺตตา.}

\proseitem{1334}{ตตฺถ กตมา โสวจสฺสตา. สหธมฺมิเก วุจฺจมาเน โสวจสฺสายํ โสวจสฺสิยํ โสวจสฺสตา อปฺปฏิกูล\-คฺคาหิตา อวิปจฺจนีกสาตตา สคารวตา สาทริยํ สาทรตา สปฺปฏิสฺสวตา. อยํ วุจฺจติ โสวจสฺสตา.}

\proseitem{1335}{ตตฺถ กตมา กลฺยาณมิตฺตตา. เย เต ปุคฺคลา สทฺธา สีลวนฺโต พหุสฺสุตา จาควนฺโต ปญฺญวนฺโต. ยา เตสํ เสวนา นิเสวนา สํเสวนา ภชนา สมฺภชนา ภตฺติ สมฺภตฺติ ตํสมฺปวงฺกตา. อยํ วุจฺจติ กลฺยาณมิตฺตตา.}

\proseitem{1336}{ตตฺถ กตมา อาปตฺติกุสลตา. ปญฺจปิ อาปตฺติกฺขนฺธา อาปตฺติโย, สตฺตปิ อาปตฺติกฺขนฺธา อาปตฺติโย. ยา ตาสํ อาปตฺตีนํ อาปตฺติกุสลตา ปญฺญา ปชานนา ฯเปฯ อโมโห ธมฺมวิจโย สมฺมาทิฏฺฐิ. อยํ วุจฺจติ อาปตฺติกุสลตา.}

\proseitem{1337}{ตตฺถ กตมา อาปตฺติวุฏฺฐาน\-กุสลตา. ยา ตาหิ อาปตฺตีหิ วุฏฺฐาน\-กุสลตา ปญฺญา ปชานนา ฯเปฯ อโมโห ธมฺมวิจโย สมฺมาทิฏฺฐิ. อยํ วุจฺจติ อาปตฺติวุฏฺฐาน\-กุสลตา.}

\chapterheadpageread{3. นิกฺเขปกณฺฑ}
\proseitem{1338}{\csromanfolio{259}ตตฺถ กตมา สมาปตฺติ\-กุสลตา. อตฺถิ สวิตกฺก\-สวิจารา สมาปตฺติ, อตฺถิ อวิตกฺก\-วิจารมตฺตา สมาปตฺติ, อตฺถิ อวิตกฺก\-อวิจารา สมาปตฺติ. ยา ตาสํ สมาปตฺตีนํ สมาปตฺติ\-กุสลตา ปญฺญา ปชานนา ฯเปฯ อโมโห ธมฺมวิจโย สมฺมาทิฏฺฐิ. อยํ วุจฺจติ สมาปตฺติ\-กุสลตา.}

\proseitem{1339}{ตตฺถ กตมา สมาปตฺติ\-วุฏฺฐาน\-กุสลตา. ยา ตาหิ สมาปตฺตีหิ วุฏฺฐาน\-กุสลตา ปญฺญา ปชานนา ฯเปฯ อโมโห ธมฺมวิจโย สมฺมาทิฏฺฐิ. อยํ วุจฺจติ สมาปตฺติ\-วุฏฺฐาน\-กุสลตา.}

\proseitem{1340}{ตตฺถ กตมา ธาตุกุสลตา. อฏฺฐารส ธาตุโย จกฺขุธาตุ รูปธาตุ จกฺขุวิญฺญาณ\-ธาตุ, โสตธาตุ สทฺทธาตุ โสตวิญฺญาณ\-ธาตุ, ฆานธาตุ คนฺธธาตุ ฆานวิญฺญาณ\-ธาตุ, ชิวฺหาธาตุ รสธาตุ ชิวฺหาวิญฺญาณ\-ธาตุ, กายธาตุ โผฏฺฐพฺพ\-ธาตุ กายวิญฺญาณ\-ธาตุ, มโนธาตุ ธมฺมธาตุ มโนวิญฺญาณ\-ธาตุ. ยา ตาสํ ธาตูนํ ธาตุกุสลตา ปญฺญา ปชานนา ฯเปฯ อโมโห ธมฺมวิจโย สมฺมาทิฏฺฐิ. อยํ วุจฺจติ ธาตุกุสลตา.}

\proseitem{1341}{ตตฺถ กตมา มนสิการ\-กุสลตา. ยา ตาสํ ธาตูนํ มนสิการ\-กุสลตา ปญฺญา ปชานนา ฯเปฯ อโมโห ธมฺมวิจโย สมฺมาทิฏฺฐิ. อยํ วุจฺจติ มนสิการ\-กุสลตา.}

\proseitem{1342}{ตตฺถ กตมา อายตน\-กุสลตา. ทฺวาทสา\-ยตนานิ จกฺขายตนํ รูปายตนํ, โสตายตนํ สทฺทายตนํ, ฆานายตนํ คนฺธายตนํ, ชิวฺหายตนํ รสายตนํ, กายายตนํ โผฏฺฐพฺพา\-ยตนํ, มนายตนํ ธมฺมายตนํ. ยา เตสํ อายตนานํ อายตน\-กุสลตา ปญฺญา ปชานนา ฯเปฯ อโมโห ธมฺมวิจโย สมฺมาทิฏฺฐิ. อยํ วุจฺจติ อายตน\-กุสลตา.}

\proseitem{1343}{ตตฺถ กตมา ปฏิจฺจสมุปฺปาท\-กุสลตา. “อวิชฺชาปจฺจยา สงฺขารา, สงฺขาร\-ปจฺจยา วิญฺญาณํ, วิญฺญาณปจฺจยา นามรูปํ, นามรูป\-ปจฺจยา สฬายตนํ, สฬายตน\-ปจฺจยา ผสฺโส, ผสฺสปจฺจยา เวทนา, เวทนาปจฺจยา ตณฺหา, ตณฺหาปจฺจยา อุปาทานํ, อุปาทานปจฺจยา ภโว, ภวปจฺจยา ชาติ, ชาติปจฺจยา ชรามรณํ, โสก\-ปริเทว\-ทุกฺข\-โทมนสฺสุ\-ปายาสา สมฺภวนฺติ, \csromanfolio{260}เอวเมตสฺส เกวลสฺส ทุกฺข\-กฺขนฺธสฺส สมุทโย โหตี”ติ ยา ตตฺถ ปญฺญา ปชานนา ฯเปฯ อโมโห ธมฺมวิจโย สมฺมาทิฏฺฐิ. อยํ วุจฺจติ ปฏิจฺจสมุปฺปาท\-กุสลตา.}

\proseitem{1344}{ตตฺถ กตมา ฐานกุสลตา. “เย เย ธมฺมา เยสํ เยสํ ธมฺมานํ เหตู ปจฺจยา อุปฺปาทาย\csromansharedfootnote{9739128e5ddb6077}{อุปาทาย (สี, ก)} ตํ ตํ ฐานนฺ”ติ ยา ตตฺถ ปญฺญา ปชานนา ฯเปฯ อโมโห ธมฺมวิจโย สมฺมาทิฏฺฐิ. อยํ วุจฺจติ ฐานกุสลตา.}

\proseitem{1345}{ตตฺถ กตมา อฏฺฐาน\-กุสลตา. “เย เย ธมฺมา เยสํ เยสํ ธมฺมานํ น เหตู น ปจฺจยา อุปฺปาทาย ตํ ตํ อฏฺฐานนฺ”ติ ยา ตตฺถ ปญฺญา ปชานนา ฯเปฯ อโมโห ธมฺมวิจโย สมฺมาทิฏฺฐิ. อยํ วุจฺจติ อฏฺฐาน\-กุสลตา.}

\proseitem{1346}{ตตฺถ กตโม อชฺชโว. ยา อชฺชวตา อชิมฺหตา อวงฺกตา อกุฏิลตา. อยํ วุจฺจติ อชฺชโว.}

\proseitem{1347}{ตตฺถ กตโม มทฺทโว. ยา มุทุตา มทฺทวตา อกกฺขฬตา อกถินตา นีจจิตฺตตา. อยํ วุจฺจติ มทฺทโว.}

\proseitem{1348}{ตตฺถ กตมา ขนฺติ. ยา ขนฺติ ขมนตา อธิวาสนตา อจณฺฑิกฺกํ อนสุโรโป อตฺตมนตา จิตฺตสฺส. อยํ วุจฺจติ ขนฺติ.}

\proseitem{1349}{ตตฺถ กตมํ โสรจฺจํ. โย กายิโก อวีติกฺกโม, วาจสิโก อวีติกฺกโม, กายิกวาจสิโก อวีติกฺกโม. อิทํ วุจฺจติ โสรจฺจํ. สพฺโพปิ สีลสํวโร โสรจฺจํ.}

\proseitem{1350}{ตตฺถ กตมํ สาขลฺยํ. ยา สา วาจา อณฺฑกา กกฺกสา ปรกฏุกา ปราภิสชฺชนี โกธสามนฺตา อสมาธิสํวตฺตนิกา ตถารูปิํ วาจํ ปหาย ยา สา วาจา เนฬา กณฺณสุขา เปมนียา หทยงฺคมา โปรี พหุชนกนฺตา พหุชนมนาปา ตถารูปิํ วาจํ ภาสิตา โหติ. ยา ตตฺถ สณฺหวาจตา สขิลวาจตา อผรุสวาจตา. อิทํ วุจฺจติ สาขลฺยํ.}

\chapterheadpageread{3. นิกฺเขปกณฺฑ}
\proseitem{1351}{\csromanfolio{261}ตตฺถ กตโม ปฏิสนฺถาโร. เทฺว ปฏิสนฺถารา อามิส\-ปฏิสนฺถาโร จ ธมฺม\-ปฏิสนฺถาโร จ. อิเธกจฺโจ ปฏิสนฺถารโก โหติ อามิส\-ปฏิสนฺถาเรน วา ธมฺม\-ปฏิสนฺถาเรน วา. อยํ วุจฺจติ ปฏิสนฺถาโร.}

\proseitem{1352}{ตตฺถ กตมา อินฺทฺริเยสุ อคุตฺตทฺวารตา. อิเธกจฺโจ จกฺขุนา รูปํ ทิสฺวา นิมิตฺตคฺคาหี โหติ อนุพฺยญฺชนคฺคาหี ยตฺวาธิกรณเมนํ จกฺขุนฺทฺริยํ อสํวุตํ วิหรนฺตํ อภิชฺฌา\-โทมนสฺสา ปาปกา อกุสลา ธมฺมา อนฺวาสฺสเวยฺยุํ ตสฺส สํวราย น ปฏิปชฺชติ น รกฺขติ จกฺขุนฺทฺริยํ จกฺขุนฺทฺริเย น สํวรํ อาปชฺชติ, โสเตน สทฺทํ สุตฺวา ฯเปฯ ฆาเนน คนฺธํ ฆายิตฺวา ฯเปฯ ชิวฺหาย รสํ สายิตฺวา ฯเปฯ กาเยน โผฏฺฐพฺพํ ผุสิตฺวา ฯเปฯ มนสา ธมฺมํ วิญฺญาย นิมิตฺตคฺคาหี โหติ อนุพฺยญฺชนคฺคาหี ยตฺวาธิกรณเมนํ มนินฺทฺริยํ อสํวุตํ วิหรนฺตํ อภิชฺฌา\-โทมนสฺสา ปาปกา อกุสลา ธมฺมา อนฺวาสฺสเวยฺยุํ ตสฺส สํวราย น ปฏิปชฺชติ น รกฺขติ มนินฺทฺริยํ มนินฺทฺริเย น สํวรํ อาปชฺชติ. ยา อิเมสํ ฉนฺนํ อินฺทฺริยานํ อคุตฺติ อโคปนา อนารกฺโข อสํวโร. อยํ วุจฺจติ อินฺทฺริเยสุ อคุตฺตทฺวารตา.}

\proseitem{1353}{ตตฺถ กตมา โภชเน อมตฺตญฺญุตา. อิเธกจฺโจ อปฺปฏิสงฺขา อโยนิโส อาหารํ อาหาเรติ ทวาย มทาย มณฺฑนาย วิภูสนาย. ยา ตตฺถ อสนฺตุฏฺฐิตา อมตฺตญฺญุตา อปฺปฏิสงฺขา โภชเน. อยํ วุจฺจติ โภชเน อมตฺตญฺญุตา.}

\proseitem{1354}{ตตฺถ กตมา อินฺทฺริเยสุ คุตฺตทฺวารตา. อิเธกจฺโจ จกฺขุนา รูปํ ทิสฺวา น นิมิตฺตคฺคาหี โหติ นานุพฺยญฺชน\-คฺคาหี ยตฺวาธิกรณเมนํ จกฺขุนฺทฺริยํ อสํวุตํ วิหรนฺตํ อภิชฺฌา\-โทมนสฺสา ปาปกา อกุสลา ธมฺมา อนฺวาสฺสเวยฺยุํ ตสฺส สํวราย ปฏิปชฺชติ รกฺขติ จกฺขุนฺทฺริยํ จกฺขุนฺทฺริเย สํวรํ อาปชฺชติ, โสเตน สทฺทํ สุตฺวา ฯเปฯ ฆาเนน คนฺธํ ฆายิตฺวา ฯเปฯ ชิวฺหาย รสํ สายิตฺวา ฯเปฯ กาเยน โผฏฺฐพฺพํ ผุสิตฺวา ฯเปฯ มนสา ธมฺมา วิญฺญาย น นิมิตฺตคฺคาหี โหติ นานุพฺยญฺชน\-คฺคาหี ยตฺวาธิกรณเมนํ มนินฺทฺริยํ อสํวุตํ วิหรนฺตํ อภิชฺฌา\-โทมนสฺสา ปาปกา อกุสลา ธมฺมา อนฺวาสฺสเวยฺยุํ ตสฺส สํวราย ปฏิปชฺชติ รกฺขติ มนินฺทฺริยํ มนินฺทฺริเย สํวรํ อาปชฺชติ. ยา อิเมสํ ฉนฺนํ อินฺทฺริยานํ คุตฺติ โคปนา อารกฺโข สํวโร. อยํ วุจฺจติ อินฺทฺริเยสุ คุตฺตทฺวารตา.}

\proseitem{1355}{\csromanfolio{262}ตตฺถ กตมา โภชเน มตฺตญฺญุตา. อิเธกจฺโจ ปฏิสงฺขา โยนิโส อาหารํ อาหาเรติ “เนว ทวาย น มทาย น มณฺฑนาย น วิภูสนาย ยาวเทว อิมสฺส กายสฺส ฐิติยา ยาปนาย วิหิํสู\-ปรติยา พฺรหฺมจริยา\-นุคฺคหาย อิติ ปุราณญฺจ เวทนํ ปฏิหงฺขามิ, นวญฺจ เวทนํ น อุปฺปาเทสฺสามิ, ยาตฺรา จ เม ภวิสฺสติ, อนวชฺชตา จ ผาสุวิหาโร จา”ติ ยา ตตฺถ สนฺตุฏฺฐิตา มตฺตญฺญุตา ปฏิสงฺขา โภชเน. อยํ วุจฺจติ โภชเน มตฺตญฺญุตา.}

\proseitem{1356}{ตตฺถ กตมํ มุฏฺฐสจฺจํ. ยา อสติ อนนุสฺสติ อปฺปฏิสฺสติ อสติ อสรณตา อธารณตา ปิลาปนตา สมฺมุสนตา. อิทํ วุจฺจติ มุฏฺฐสจฺจํ.}

\proseitem{1357}{ตตฺถ กตมํ อสมฺปชญฺญํ. ยํ อญฺญาณํ อทสฺสนํ ฯเปฯ อวิชฺชาลงฺคี โมโห อกุสลมูลํ. อิทํ วุจฺจติ อสมฺปชญฺญํ.}

\proseitem{1358}{ตตฺถ กตมา สติ. ยา สติ อนุสฺสติ ปฏิสฺสติ สติ สรณตา ธารณตา อปิลาปนตา อสมฺมุสนตา สติ สตินฺทฺริยํ สติพลํ สมฺมาสติ. อยํ วุจฺจติ สติ.}

\proseitem{1359}{ตตฺถ กตมํ สมฺปชญฺญํ. ยา ปญฺญา ปชานนา ฯเปฯ อโมโห ธมฺมวิจโย สมฺมาทิฏฺฐิ. อิทํ วุจฺจติ สมฺปชญฺญํ.}

\proseitem{1360}{ตตฺถ กตมํ ปฏิสงฺขานพลํ. ยา ปญฺญา ปชานนา ฯเปฯ อโมโห ธมฺมวิจโย สมฺมาทิฏฺฐิ. อิทํ วุจฺจติ ปฏิสงฺขานพลํ.}

\proseitem{1361}{ตตฺถ กตมํ ภาวนาพลํ. ยา กุสลานํ ธมฺมานํ อาเสวนา ภาวนา พหุลีกมฺมํ. อิทํ วุจฺจติ ภาวนาพลํ. สตฺตปิ โภชฺฌงฺคา ภาวนาพลํ.}

\proseitem{1362}{ตตฺถ กตโม สมโถ. ยา จิตฺตสฺส ฐิติ ฯเปฯ สมฺมาสมาธิ. อยํ วุจฺจติ สมโถ.}

\proseitem{1363}{ตตฺถ กตมา วิปสฺสนา. ยา ปญฺญา ปชานนา ฯเปฯ อโมโห ธมฺมวิจโย สมฺมาทิฏฺฐิ. อยํ วุจฺจติ วิปสฺสนา.}

\chapterheadpageread{3. นิกฺเขปกณฺฑ}
\proseitem{1364}{\csromanfolio{263}ตตฺถ กตมํ สมถ\-นิมิตฺตํ. ยา จิตฺตสฺส ฐิติ ฯเปฯ สมฺมาสมาธิ. อิทํ วุจฺจติ สมถ\-นิมิตฺตํ.}

\proseitem{1365}{ตตฺถ กตมํ ปคฺคาห\-นิมิตฺตํ. โย เจตสิโก วีริยารมฺโภ ฯเปฯ สมฺมาวายาโม. อิทํ วุจฺจติ ปคฺคาห\-นิมิตฺตํ.}

\proseitem{1366}{ตตฺถ กตโม ปคฺคาโห. โย เจตสิโก วีริยารมฺโภ ฯเปฯ สมฺมาวายาโม. อยํ วุจฺจติ ปคฺคาโห.}

\proseitem{1367}{ตตฺถ กตโม อวิกฺเขโป. ยา จิตฺตสฺส ฐิติ ฯเปฯ สมฺมาสมาธิ. อยํ วุจฺจติ อวิกฺเขโป.}

\proseitem{1368}{ตตฺถ กตมา สีลวิปตฺติ. โย กายิโก วีติกฺกโม, วาจสิโก วีติกฺกโม, กายิกวาจสิโก วีติกฺกโม. อยํ วุจฺจติ สีลวิปตฺติ. สพฺพมฺปิ ทุสฺสิลฺยํ สีลวิปตฺติ.}

\proseitem{1369}{ตตฺถ กตมา ทิฏฺฐิวิปตฺติ. “นตฺถิ ทินฺนํ, นตฺถิ ยิฏฺฐํ, นตฺถิ หุตํ, นตฺถิ สุกต\-ทุกฺกฏานํ กมฺมานํ ผลํ วิปาโก, นตฺถิ อยํ โลโก, นตฺถิ ปโร โลโก, นตฺถิ มาตา, นตฺถิ ปิตา, นตฺถิ สตฺตา โอปปาติกา, นตฺถิ โลเก สมณพฺราหฺมณา สมฺมคฺคตา สมฺมาปฏิปนฺนา, เย อิมญฺจ โลกํ ปรญฺจ โลกํ สยํ อภิญฺญา สจฺฉิกตฺวา ปเวเทนฺตี”ติ ยา เอวรูปา ทิฏฺฐิ ทิฏฺฐิคตํ ฯเปฯ วิปริยาส\-คฺคาโห. อยํ อุจฺจติ ทิฏฺฐิวิปตฺติ. สพฺพาปิ มิจฺฉาทิฏฺฐิ ทิฏฺฐิวิปตฺติ.}

\proseitem{1370}{ตตฺถ กตมา สีลสมฺปทา. โย กายิโก อวีติกฺกโม, วาจสิโก อวีติกฺกโม, กายิกวาจสิโก อวีติกฺกโม. อยํ วุจฺจติ สีลสมฺปทา. สพฺโพปิ สีลสํวโร สีลสมฺปทา.}

\proseitem{1371}{ตตฺถ กตมา ทิฏฺฐิสมฺปทา. “อตฺถิ ทินฺนํ, อตฺถิ ยิฏฺฐํ, อตฺถิ หุตํ, อตฺถิ สุกต\-ทุกฺกฏานํ กมฺมานํ ผลํ วิปาโก, อตฺถิ อยํ โลโก, อตฺถิ ปโร โลโก, อตฺถิ มาตา, อตฺถิ ปิตา, อตฺถิ สตฺตา โอปปาติกา, อตฺถิ โลเก สมณพฺราหฺมณา สมฺมคฺคตา สมฺมาปฏิปนฺนา, เย อิมญฺจ โลกํ ปรญฺจ โลกํ สยํ อภิญฺญา สจฺฉิกตฺวา ปเวเทนฺตี”ติ ยา เอวรูปา ปญฺญา ปชานนา ฯเปฯ อโมโห ธมฺมวิจโย สมฺมาทิฏฺฐิ. อยํ วุจฺจติ ทิฏฺฐิสมฺปทา. สพฺพาปิ สมฺมาทิฏฺฐิ ทิฏฺฐิสมฺปทา.}

\proseitem{1372}{\csromanfolio{264}ตตฺถ กตมา สีลวิสุทฺธิ. โย กายิโก อวีติกฺกโม, วาจสิโก อวีติกฺกโม, กายิกวาจสิโก อวีติกฺกโม. อยํ วุจฺจติ สีลวิสุทฺธิ. สพฺโพปิ สีลสํวโร สีลวิสุทฺธิ.}

\proseitem{1373}{ตตฺถ กตมา ทิฏฺฐิวิสุทฺธิ. กมฺมสฺสกต\-ญาณํ สจฺจานุโลมิกญาณํ มคฺค\-สมงฺคิสฺส ญาณํ ผลสมงฺคิสฺส ญาณํ.}

\proseitem{1374}{\textbf{ทิฏฺฐิวิสุทฺธิ โข ปนา}ติ ยา ปญฺญา ปชานนา ฯเปฯ อโมโห ธมฺมวิจโย สมฺมาทิฏฺฐิ.}

\proseitem{1375}{\textbf{ยถาทิฏฺฐิสฺส จ ปธาน}นฺติ โย เจตสิโก วีริยารมฺโภ ฯเปฯ สมฺมาวายาโม.}

\proseitem{1376}{\textbf{สํเวโค}ติ ชาติภยํ ชราภยํ พฺยาธิภยํ มรณาภยํ, สํเวชนิยํ\csromansharedfootnote{a6192dbd7df7e1ee}{สํเวชนียํ (สี)} ฐานนฺติ ชาติ ชรา พฺยาธิ มรณํ.}

\proseitem{1377}{\textbf{สํวิคฺคสฺส จ โยนิโส ปธาน}นฺติ อิธ ภิกฺขุ อนุปฺปนฺนานํ ปาปกานํ อกุสลานํ ธมฺมานํ อนุปฺปาทาย ฉนฺทํ ชเนติ วายมติ วีริยํ อารภติ จิตฺตํ ปคฺคณฺหาติ ปทหติ, อุปฺปนฺนานํ ปาปกานํ อกุสลานํ ธมฺมานํ ปหานาย ฉนฺทํ ชเนติ วายมติ วีริยํ อารภติ จิตฺตํ ปคฺคณฺหาติ ปทหติ, อนุปฺปนฺนานํ กุสลานํ ธมฺมานํ อุปฺปาทาย ฉนฺทํ ชเนติ วายมติ วีริยํ อารภติ จิตฺตํ ปคฺคณฺหาติ ปทหติ, อุปฺปนฺนานํ กุสลานํ ธมฺมานํ ฐิติยา อสมฺโมสาย ภิยฺโยภาวาย เวปุลฺลาย ภาวนาย ปาริปูริยา ฉนฺทํ ชเนติ วายมติ วีริยํ อารภติ จิตฺตํ ปคฺคณฺหาติ ปทหติ.}

\proseitem{1378}{\textbf{อสนฺตุฏฺฐิตา จ กุสเลสุ ธมฺเมสู}ติ ยา กุสลานํ ธมฺมานํ ภาวนาย อสนฺตุฏฺฐสฺส ภิยฺโยกมฺยตา.}

\proseitem{1379}{อปฺปฏิวานิตา จ ปธานสฺมินฺติ ยา กุสลานํ ธมฺมานํ ภาวนาย สกฺกจฺจ\-กิริยตา สาตจฺจกิริยตา อฏฺฐิต\-กิริยตา อโนลีนวุตฺติตา อนิกฺขิตฺต\-ฉนฺทตา อนิกฺขิตฺตธุรตา อาเสวนา ภาวนา พหุลีกมฺมํ.}

\chapterheadpageread{3. นิกฺเขปกณฺฑ}
\proseitem{1380}{\csromanfolio{265}วิชฺชาติ ติสฺโส วิชฺชา. ปุพฺเพ\-นิวาสา\-นุสฺสติ\-ญาณํ วิชฺชา, สตฺตานํ จุตูปปาเต ญาณํ วิชฺชา, อาสวานํ ขเย ญาณํ วิชฺชา.}

\proseitem{1381}{\textbf{วิมุตฺตี}ติ เทฺว วิมุตฺติโย. จิตฺตสฺส\csromansharedfootnote{50c47fac685b37a1}{จิตฺตสฺส จ (สี, สฺยา)} อธิมุตฺติ นิพฺพานญฺจ.}

\proseitem{1382}{ขเย ญาณนฺติ มคฺค\-สมงฺคิสฺส ญาณํ.}

\proseitemwithcloser{1383}{\textbf{อนุปฺปาเท ญาณนฺ}ติ ผลสมงฺคิสฺส ญาณํ.}{\textbf{นิกฺเขปกณฺโฑ นิฏฺฐิโต.}}
\csromanmatikaanchor{mtk.118}
\csromantocmarknopage{chapter}{4. อฏฺฐกถากณฺฑ}{mtk.118}
\bookmark[dest={mtk.118},level=0]{4. อฏฺฐกถากณฺฑ}
\chapterheadpageread{4. อฏฺฐกถา\-กณฺฑ}
\csromanmatikaanchor{mtk.119}
\csromantocmarkref{section}{ติกอตฺถุทฺธาร}{mtk.119}
\bookmark[dest={mtk.119},level=1]{ติกอตฺถุทฺธาร}
\prose{\csromanfolio{266}ติกอตฺถุทฺธาร}

\proseitem{1384}{กตเม ธมฺมา กุสลา. จตูสุ ภูมีสุ กุสลํ. อิเม ธมฺมา กุสลา.}

\proseitem{1385}{กตเม ธมฺมา อกุสลา. ทฺวาทส อกุสล\-จิตฺตุปฺปาทา. อิเม ธมฺมา อกุสลา.}

\proseitem{1386}{กตเม ธมฺมา อพฺยากตา. จตูสุ ภูมีสุ วิปาโก ตีสุ ภูมีสุ กิริยาพฺยากตํ รูปญฺจ นิพฺพานญฺจ. อิเม ธมฺมา อพฺยากตา.}

\proseitem{1387}{กตเม ธมฺมา สุขาย เวทนาย สมฺปยุตฺตา. กามาวจร\-กุสลโต จตฺตาโร โสมนสฺสสหคต\-จิตฺตุปฺปาทา อกุสลโต จตฺตาโร กามาวจร\-กุสลสฺส วิปากโต จ กิริยโต จ ปญฺจ รูปาวจร\-ติกจตุกฺกชฺฌานา กุสลโต จ วิปากโต จ กิริยโต จ โลกุตฺตร\-ติกจตุกฺกชฺฌานา กุสลโต จ วิปากโต จ เอตฺถุปฺปนฺนํ สุขํ เวทนํ ฐเปตฺวา. อิเม ธมฺมา สุขาย เวทนาย สมฺปยุตฺตา.}

\proseitem{1388}{กตเม ธมฺมา ทุกฺขาย เวทนาย สมฺปยุตฺตา. เทฺว โทมนสฺสสหคต\-จิตฺตุปฺปาทา ทุกฺข\-สหคตํ กายวิญฺญาณํ เอตฺถุปฺปนฺนํ ทุกฺขํ เวทนํ ฐเปตฺวา. อิเม ธมฺมา ทุกฺขาย เวทนาย สมฺปยุตฺตา.}

\proseitem{1389}{กตเม ธมฺมา อทุกฺขม\-สุขาย เวทนาย สมฺปยุตฺตา. กามาวจร\-กุสลโต จตฺตาโร อุเปกฺขาสหคต\-จิตฺตุปฺปาทา อกุสลโต ฉ กามาวจร\-กุสลสฺส วิปากโต ทส อกุสลสฺส วิปากโต ฉ กิริยโต ฉ รูปาวจรํ จตุตฺถํ ฌานํ กุสลโต จ วิปากโต จ กิริยโต จ จตฺตาโร อรูปาวจรา กุสลโต จ วิปากโต จ เอตฺถุปฺปนฺนํ อทุกฺขมสุขํ เวทนํ ฐเปตฺวา. อิเม ธมฺมา อทุกฺขม\-สุขาย เวทนาย สมฺปยุตฺตา. ติสฺโส จ เวทนา รูปญฺจ นิพฺพานญฺจ. \csromanfolio{267}อิเม ธมฺมา น วตฺตพฺพา สุขาย เวทนาย สมฺปยุตฺตาติปิ ทุกฺขาย เวทนาย สมฺปยุตฺตาติปิ อทุกฺขม\-สุขาย เวทนาย สมฺปยุตฺตาติปิ.}

\proseitem{1390}{กตเม ธมฺมา วิปากา. จตูสุ ภูมีสุ วิปาโก. อิเม ธมฺมา วิปากา.}

\proseitem{1391}{กตเม ธมฺมา วิปาก\-ธมฺม\-ธมฺมา. จตูสุ ภูมีสุ กุสลํ อกุสลํ. อิเม ธมฺมา วิปาก\-ธมฺม\-ธมฺมา.}

\proseitem{1392}{กตเม ธมฺมา เนว\-วิปาก\-น\-วิปากธมฺม\-ธมฺมา. ตีสุ ภูมีสุ กิริยาพฺยากตํ รูปญฺจ นิพฺพานญฺจ. อิเม ธมฺมา เนว\-วิปาก\-น\-วิปากธมฺม\-ธมฺมา.}

\proseitem{1393}{กตเม ธมฺมา อุปาทิณฺณุ\-ปาทานิยา. ตีสุ ภูมีสุ วิปาโก, ยญฺจ รูปํ กมฺมสฺส กตตฺตา. อิเม ธมฺมา อุปาทิณฺณุ\-ปาทานิยา.}

\proseitem{1394}{กตเม ธมฺมา อนุปาทิณฺณุ\-ปาทานิยา. ตีสุ ภูมีสุ กุสลํ อกุสลํ ตีสุ ภูมีสุ กิริยาพฺยากตํ, ยญฺจ รูปํ น กมฺมสฺส กตตฺตา. อิเม ธมฺมา อนุปาทิณฺณุ\-ปาทานิยา.}

\proseitem{1395}{กตเม ธมฺมา อนุปาทิณฺณ\-อนุปาทานิยา. จตฺตาโร มคฺคา อปริยาปนฺนา จตฺตาริ จ สามญฺญาผลานิ นิพฺพานญฺจ. อิเม ธมฺมา อนุปาทิณฺณ\-อนุปาทานิยา.}

\proseitem{1396}{กตเม ธมฺมา สํกิลิฏฺฐ\-สํกิเลสิกา. ทฺวาทสา\-กุสล\-จิตฺตุปฺปาทา. อิเม ธมฺมา สํกิลิฏฺฐ\-สํกิเลสิกา.}

\proseitem{1397}{กตเม ธมฺมา อสํกิลิฏฺฐ\-สํกิเลสิกา. ตีสุ ภูมีสุ กุสลํ ตีสุ ภูมีสุ วิปาโก ตีสุ ภูมีสุ กิริยาพฺยากตํ สพฺพญฺจ รูปํ. อิเม ธมฺมา อสํกิลิฏฺฐ\-สํกิเลสิกา.}

\proseitem{1398}{กตเม ธมฺมา อสํกิลิฏฺฐ\-อสํกิเลสิกา. จตฺตาโร มคฺคา อปริยาปนฺนา จตฺตาริ จ สามญฺญาผลานิ นิพฺพานญฺจ. อิเม ธมฺมา อสํกิลิฏฺฐ\-อสํกิเลสิกา.}

\proseitem{1399}{กตเม ธมฺมา สวิตกฺก\-สวิจารา. กามาวจรํ กุสลํ อกุสลํ กามาวจร\-กุสลสฺส วิปากโต เอกาทส จิตฺตุปฺปาทา อกุสลสฺส วิปากโต เทฺว กิริยโต เอกาทส รูปาวจรํ ปฐมํ \csromanfolio{268}ฌานํ กุสลโต จ วิปากโต จ กิริยโต จ โลกุตฺตรํ ปฐมํ ฌานํ กุสลโต จ วิปากโต จ เอตฺถุปฺปนฺเน วิตกฺกวิจาเร ฐเปตฺวา. อิเม ธมฺมา สวิตกฺก\-สวิจารา.}

\proseitem{1400}{กตเม ธมฺมา อวิตกฺก\-วิจารมตฺตา. รูปาวจร\-ปญฺจกนเย ทุติยํ ฌานํ กุสลโต จ วิปากโต จ กิริยโต จ โลกุตฺตร\-ปญฺจกนเย ทุติยํ ฌานํ กุสลโต จ วิปากโต จ เอตฺถุปฺปนฺนํ วิจารํ ฐเปตฺวา วิตกฺโก จ. อิเม ธมฺมา อวิตกฺก\-วิจารมตฺตา.}

\proseitem{1401}{กตเม ธมฺมา อวิตกฺก\-อวิจารา. เทฺว\-ปญฺจวิญฺญาณานิ รูปาวจร\-ติก\-ติก\-ชฺฌานา กุสลโต จ วิปากโต จ กิริยโต จ จตฺตาโร อารุปฺปา กุสลโต จ วิปากโต จ กิริยโต จ โลกุตฺตร\-ติก\-ติก\-ชฺฌานา กุสลโต จ วิปากโต จ ปญฺจกนเย ทุติเย ฌาเน\csromansharedfootnote{aace4c5353719a0d}{ทุติยชฺฌาเน (สี)} อุปฺปนฺโน จ วิจาโร รูปญฺจ นิพฺพานญฺจ. อิเม ธมฺมา อวิตกฺก\-อวิจารา. วิตกฺก\-สหชาโต วิจาโร น วตฺตพฺโพ สวิตกฺกสวิจาโรติปิ อวิตกฺกวิจารมตฺโตติปิ อวิตกฺกอวิจาโรติปิ.}

\proseitem{1402}{กตเม ธมฺมา ปีติสหคตา. กามาวจร\-กุสลโต จตฺตาโร โสมนสฺสสหคต\-จิตฺตุปฺปาทา อกุสลโต จตฺตาโร กามาวจร\-กุสลสฺส วิปากโต ปญฺจ กิริยโต ปญฺจ รูปาวจร\-ทุ\-กติก\-ชฺฌานา กุสลโต จ วิปากโต จ กิริยโต จ โลกุตฺตร\-ทุ\-กติก\-ชฺฌานา กุสลโต จ วิปากโต จ เอตฺถุปฺปนฺนํ ปีติํ ฐเปตฺวา. อิเม ธมฺมา ปีติสหคตา.}

\proseitem{1403}{กตเม ธมฺมา สุขสหคตา. กามาวจร\-กุสลโต จตฺตาโร โสมนสฺสสหคต\-จิตฺตุปฺปาทา อกุสลโต จตฺตาโร กามาวจร\-กุสลสฺส วิปากโต ฉ กิริยโต ปญฺจ รูปาวจร\-ติกจตุกฺกชฺฌานา กุสลโต จ วิปากโต จ กิริยโต จ โลกุตฺตร\-ติกจตุกฺกชฺฌานา กุสลโต จ วิปากโต จ เอตฺถุปฺปนฺนํ สุขํ ฐเปตฺวา. อิเม ธมฺมา สุขสหคตา.}

\proseitem{1404}{กตเม ธมฺมา อุเปกฺขา\-สหคตา. กามาวจร\-กุสลโต จตฺตาโร อุเปกฺขาสหคต\-จิตฺตุปฺปาทา อกุสลโต ฉ กามาวจร- \csromanfolio{269}กุสลสฺส วิปากโต ทส อกุสลสฺส วิปากโต ฉ กิริยโต ฉ รูปาวจรํ จตุตฺถํ ฌานํ กุสลโต จ วิปากโต จ กิริยโต จ จตฺตาโร อารุปฺปา กุสลโต จ วิปากโต จ กิริยโต จ โลกุตฺตรํ จตุตฺถํ ฌานํ กุสลโต จ วิปากโต จ เอตฺถุปฺปนฺนํ อุเปกฺขํ ฐเปตฺวา. อิเม ธมฺมา อุเปกฺขา\-สหคตา. ปีติ น ปีติสหคตา สุขสหคตา น อุเปกฺขา\-สหคตา, สุขํ น สุขสหคตํ สิยา ปีติสหคตํ น อุเปกฺขา\-สหคตํ สิยา น วตฺตพฺพํ ปีติสหคตนฺติ. เทฺว โทมนสฺสสหคต\-จิตฺตุปฺปาทา ทุกฺข\-สหคต\-กายวิญฺญาณํ, ยา จ เวทนา อุเปกฺขา รูปญฺจ นิพฺพานญฺจ. อิเม ธมฺมา น วตฺตพฺพา ปีติสหคตาติปิ สุขสหคตาติปิ อุเปกฺขา\-สหคตาติปิ.}

\proseitem{1405}{กตเม ธมฺมา ทสฺสเนน ปหาตพฺพา. จตฺตาโร ทิฏฺฐิคต\-สมฺปยุตฺต\-จิตฺตุปฺปาทา วิจิกิจฺฉา\-สหคโต จิตฺตุปฺปาโท. อิเม ธมฺมา ทสฺสเนน\-ปหาตพฺพา.}

\proseitem{1406}{กตเม ธมฺมา ภาวนาย ปหาตพฺพา. อุทฺธจฺจ\-สหคโต จิตฺตุปฺปาโท. อิเม ธมฺมา ภาวนาย ปหาตพฺพา. จตฺตาโร ทิฏฺฐิคต\-วิปฺปยุตฺตา โลภ\-สหคต\-จิตฺตุปฺปาทา เทฺว โทมนสฺสสหคต\-จิตฺตุปฺปาทา. อิเม ธมฺมา สิยา ทสฺสเนน ปหาตพฺพา, สิยา ภาวนาย ปหาตพฺพา.}

\proseitem{1407}{กตเม ธมฺมา เนว ทสฺสเนน น ภาวนาย ปหาตพฺพา. จตูสุ ภูมีสุ กุสลํ จตูสุ ภูมีสุ วิปาโก ตีสุ ภูมีสุ กิริยาพฺยากตํ รูปญฺจ นิพฺพานญฺจ. อิเม ธมฺมา เนว ทสฺสเนน น ภาวนาย ปหาตพฺพา.}

\proseitem{1408}{กตเม ธมฺมา ทสฺสเนน ปหาตพฺพ\-เหตุกา. จตฺตาโร ทิฏฺฐิคต\-สมฺปยุตฺต\-จิตฺตุปฺปาทา วิจิกิจฺฉา\-สหคโต จิตฺตุปฺปาโท เอตฺถุปฺปนฺนํ โมหํ ฐเปตฺวา. อิเม ธมฺมา ทสฺสเนน ปหาตพฺพ\-เหตุกา.}

\proseitem{1409}{กตเม ธมฺมา ภาวนาย ปหาตพฺพ\-เหตุกา. อุทฺธจฺจ\-สหคโต จิตฺตุปฺปาโท เอตฺถุปฺปนฺนํ โมหํ ฐเปตฺวา. อิเม ธมฺมา ภาวนาย ภาวนาย ปหาตพฺพ\-เหตุกา. จตฺตาโร ทิฏฺฐิคต\-วิปฺปยุตฺตา โลภ\-สหคต\-จิตฺตุปฺปาทา เทฺว โทมนสฺสสหคต\-จิตฺตุปฺปาทา. อิเม ธมฺมา สิยา ทสฺสเนน ปหาตพฺพ\-เหตุกา, สิยา ภาวนาย ปหาตพฺพ\-เหตุกา.}

\proseitem{1410}{\csromanfolio{270}กตเม ธมฺมา เนว ทสฺสเนน น ภาวนาย ปหาตพฺพ\-เหตุกา. วิจิกิจฺฉา\-สหคโต โมโห อุทฺธจฺจ\-สหคโต โมโห จตูสุ ภูมีสุ กุสลํ จตูสุ ภูมีสุ วิปาโก ตีสุ ภูมีสุ กิริยาพฺยากตํ รูปญฺจ นิพฺพานญฺจ. อิเม ธมฺมา เนว ทสฺสเนน น ภาวนาย ปหาตพฺพ\-เหตุกา.}

\proseitem{1411}{กตเม ธมฺมา อาจยคามิโน. ตีสุ ภูมีสุ กุสลํ อกุสลํ. อิเม ธมฺมา อาจยคามิโน.}

\proseitem{1412}{กตเม ธมฺมา อปจยคามิโน. จตฺตาโร มคฺคา อปริยาปนฺนา. อิเม ธมฺมา อปจยคามิโน.}

\proseitem{1413}{กตเม ธมฺมา เนวา\-จยคามิ\-นา\-ปจยคามิโน. จตูสุ ภูมีสุ วิปาโก ตีสุ ภูมีสุ กิริยาพฺยากตํ รูปญฺจ นิพฺพานญฺจ. อิเม ธมฺมา เนวา\-จยคามิ\-นา\-ปจยคามิโน.}

\proseitem{1414}{กตเม ธมฺมา เสกฺขา. จตฺตาโร มคฺคา อปริยาปนฺนา เหฏฺฐิมานิ จ ตีณิ สามญฺญผลานิ. อิเม ธมฺมา เสกฺขา.}

\proseitem{1415}{กตเม ธมฺมา อเสกฺขา. อุปริฏฺฐิมํ อรหตฺตผลํ. อิเม ธมฺมา อเสกฺขา.}

\proseitem{1416}{กตเม ธมฺมา เนว\-เสกฺข\-นาเสกฺขา. ตีสุ ภูมีสุ กุสลํ อกุสลํ ตีสุ ภูมีสุ วิปาโก ตีสุ ภูมีสุ กิริยาพฺยากตํ รูปญฺจ นิพฺพานญฺจ. อิเม ธมฺมา เนว\-เสกฺข\-นาเสกฺขา.}

\proseitem{1417}{กตเม ธมฺมา ปริตฺตา. กามาวจร\-กุสลํ อกุสลํ สพฺโพ กามาวจรสฺส วิปาโก กามาวจร\-กิริยา\-พฺยากตํ สพฺพญฺจ รูปํ. อิเม ธมฺมา ปริตฺตา.}

\proseitem{1418}{กตเม ธมฺมา มหคฺคตา. รูปาวจรา อรูปาวจรา กุสลาพฺยากตา. อิเม ธมฺมา มหคฺคตา.}

\proseitem{1419}{กตเม ธมฺมา อปฺปมาณา. จตฺตาโร มคฺคา อปริยาปนฺนา จตฺตาริ จ สามญฺญาผลานิ นิพฺพานญฺจ. อิเม ธมฺมา อปฺปมาณา.}

\chapterheadpageread{4. อฏฺฐกถา\-กณฺฑ}
\proseitem{1420}{\csromanfolio{271}กตเม ธมฺมา ปริตฺตารมฺมณา. สพฺโพ กามาวจรสฺส วิปาโก กิริยา\-มโนธาตุ กิริยา\-เหตุก\-มโนวิญฺญาณธาตุ โสมนสฺส\-สหคตา. อิเม ธมฺมา ปริตฺตารมฺมณา.}

\proseitem{1421}{กตเม ธมฺมา มหคฺคตา\-รมฺมณา. วิญฺญาณญฺจา\-ยตนํ เนว\-สญฺญา\-นา\-สญฺญา\-ยตนํ. อิเม ธมฺมา มหคฺคตา\-รมฺมณา.}

\proseitem{1422}{กตเม ธมฺมา อปฺปมาณา\-รมฺมณา. จตฺตาโร มคฺคา อปริยาปนฺนา จตฺตาริ จ สามญฺญผลานิ. อิเม ธมฺมา อปฺปมาณา\-รมฺมณา. กามาวจร\-กุสลโต จตฺตาโร ญาณวิปฺปยุตฺต\-จิตฺตุปฺปาทา กิริยโต จตฺตาโร ญาณวิปฺปยุตฺต\-จิตฺตุปฺปาทา สพฺพํ อกุสลํ. อิเม ธมฺมา สิยา ปริตฺตารมฺมณา, สิยา มหคฺคตา\-รมฺมณา, น อปฺปมาณา\-รมฺมณา. สิยา น วตฺตพฺพา ปริตฺตารมฺมณาติปิ มหคฺคตา\-รมฺมณาติปิ. กามาวจร\-กุสลโต จตฺตาโร ญาณสมฺปยุตฺต\-จิตฺตุปฺปาทา กิริยโต จตฺตาโร ญาณสมฺปยุตฺต\-จิตฺตุปฺปาทา รูปาวจรํ จตุตฺถํ ฌานํ กุสลโต จ กิริยโต จ กิริยา\-เหตุก\-มโนวิญฺญาณธาตุ อุเปกฺขา\-สหคตา. อิเม ธมฺมา สิยา ปริตฺตารมฺมณา, สิยา มหคฺคตา\-รมฺมณา, สิยา อปฺปมาณา\-รมฺมณา. สิยา น วตฺตพฺพา ปริตฺตารมฺมณาติปิ มหคฺคตา\-รมฺมณาติปิ อปฺปมาณา\-รมฺมณาติปิ. รูปาวจร\-ติกจตุกฺกชฺฌานา กุสลโต จ วิปากโต จ กิริยโต จ จตุตฺถสฺส ฌานสฺส วิปาโก อากาสา\-นญฺจา\-ยตนํ อากิญฺจญฺญา\-ยตนํ. อิเม ธมฺมา น วตฺถพฺพา ปริตฺตารมฺมณาติปิ มหคฺคตา\-รมฺมณาติปิ อปฺปมาณา\-รมฺมณาติปิ. รูปญฺจ นิพฺพานญฺจ อนารมฺมณา.}

\proseitem{1423}{กตเม ธมฺมา หีนา. ทฺวาทส อกุสล\-จิตฺตุปฺปาทา. อิเม ธมฺมา หีนา.}

\proseitem{1424}{กตเม ธมฺมา มชฺฌิมา. ตีสุ ภูมีสุ กุสลํ ตีสุ ภูมีสุ วิปาโก ตีสุ ภูมีสุ กิริยาพฺยากตํ สพฺพญฺจ รูปํ. อิเม ธมฺมา มชฺฌิมา.}

\proseitem{1425}{กตเม ธมฺมา ปณีตา. จตฺตาโร มคฺคา อปริยาปนฺนา จตฺตาริ จ สามญฺญผลานิ นิพฺพานญฺจ. อิเม ธมฺมา ปณีตา.}

\proseitem{1426}{กตเม ธมฺมา มิจฺฉตฺต\-นิยตา. จตฺตาโร ทิฏฺฐิคต\-สมฺปยุตฺต\-จิตฺตุปฺปาทา เทฺว โทมนสฺสสหคต\-จิตฺตุปฺปาทา. อิเม ธมฺมา สิยา มิจฺฉตฺต\-นิยตา, สิยา อนิยตา.}

\proseitem{1427}{\csromanfolio{272}กตเม ธมฺมา สมฺมตฺตนิยตา. จตฺตาโร มคฺคา อปริยาปนฺนา. อิเม ธมฺมา สมฺมตฺตนิยตา.}

\proseitem{1428}{กตเม ธมฺมา อนิยตา. จตฺตาโร ทิฏฺฐิคต\-วิปฺปยุตฺต\-โลภ\-สหคต\-จิตฺตุปฺปาทา วิจิกิจฺฉา\-สหคโต จิตฺตุปฺปาโท อุทฺธจฺจ\-สหคโต จิตฺตุปฺปาโท ตีสุ ภูมีสุ กุสลํ จตูสุ ภูมีสุ วิปาโก ตีสุ ภูมีสุ กิริยาพฺยากตํ รูปญฺจ นิพฺพานญฺจ. อิเม ธมฺมา อนิยตา.}

\proseitem{1429}{กตเม ธมฺมา มคฺคารมฺมณา. กามาวจร\-กุสลโต จตฺตาโร ญาณสมฺปยุตฺต\-จิตฺตุปฺปาทา กิริยโต จตฺตาโร ญาณสมฺปยุตฺต\-จิตฺตุปฺปาทา. อิเม ธมฺมา สิยา มคฺคารมฺมณา, น มคฺคเหตุกา, สิยา มคฺคาธิปติโน. สิยา น วตฺตพฺพา มคฺคารมฺมณาติปิ มคฺคาธิปติโนติปิ. จตฺตาโร อริยมคฺคา น มคฺคารมฺมณา มคฺคเหตุกา, สิยา มคฺคาธิปติโน. สิยา น วตฺตพฺพา มคฺคาธิปติโนติ. รูปาวจร\-จตุตฺถํ ฌานํ กุสลโต จ กิริยโต จ กิริยา\-เหตุก\-มโนวิญฺญาณธาตุ อุเปกฺขา\-สหคตา. อิเม ธมฺมา สิยา มคฺคารมฺมณา, น มคฺคเหตุกา, น มคฺคาธิปติโน. สิยา น วตฺตพฺพา มคฺคารมฺมณาติ. กามาวจร\-กุสลโต จตฺตาโร ญาณวิปฺปยุตฺต\-จิตฺตุปฺปาทา สพฺพํ อกุสลํ สพฺโพ กามาวจรสฺส วิปาโก กิริยโต ฉ จิตฺตุปฺปาทา รูปาวจร\-ติกจตุกฺกชฺฌานา กุสลโต จ วิปากโต จ กิริยโต จ จตุตฺถสฺส ฌานสฺส วิปาโก จตฺตาโร อารุปฺปา กุสลโต จ วิปากโต จ กิริยโต จ จตฺตาริ จ สามญฺญาผลานิ. อิเม ธมฺมา น วตฺตพฺพา มคฺคารมฺมณาติปิ มคฺคเหตุกาติปิ มคฺคาธิปติโนติปิ. รูปญฺจ นิพฺพานญฺจ อนารมฺมณา.}

\proseitem{1430}{กตเม ธมฺมา อุปฺปานฺนา. จตูสุ ภูมีสุ วิปาโก, ยญฺจ รูปํ กมฺมสฺส กตตฺตา. อิเม ธมฺมา สิยา อุปฺปนฺนา สิยา อุปฺปาทิโน. น วตฺตพฺพา อนุปฺปนฺนาติ. จตูสุ ภูมีสุ กุสลํ อกุสลํ ตีสุ ภูมีสุ กิริยาพฺยากตํ, ยญฺจ รูปํ น กมฺมสฺส กตตฺตา. อิเม ธมฺมา สิยา อุปฺปานฺนา สิยา อนุปฺปนฺนา. น วตฺตพฺพา อุปฺปาทิโนติ. นิพฺพานํ น วตฺตพฺพํ อุปฺปนฺนนฺติปิ อนุปฺปนฺนนฺติปิ อุปฺปาทิโนติปิ.}

\proseitem{1431}{นิพฺพานํ ฐเปตฺวา สพฺเพ ธมฺมา สิยา อตีตา สิยา อนาคตา สิยา ปจฺจุปฺปนฺนา. นิพฺพานํ น วตฺตพฺพํ อตีตนฺติปิ อนาคตนฺติปิ ปจฺจุปฺปนฺนนฺติปิ.}

\chapterheadpageread{4. อฏฺฐกถา\-กณฺฑ}
\proseitem{1432}{\csromanfolio{273}กตเม ธมฺมา อตีตารมฺมณา. วิญฺญาณญฺจา\-ยตนํ เนว\-สญฺญา\-นา\-สญฺญา\-ยตนํ. อิเม ธมฺมา อตีตารมฺมณา.}

\proseitem{1433}{นิโยคา อนาคตารมฺมณา นตฺถิ.}

\proseitem{1434}{กตเม ธมฺมา ปจฺจุปฺปนฺนา\-รมฺมณา. เทฺว\-ปญฺจวิญฺญาณานิ ติสฺโส จ มโนธาตุโย. อิเม ธมฺมา ปจฺจุปฺปนฺนา\-รมฺมณา. กามาวจร\-กุสลสฺส วิปากโต ทส จิตฺตุปฺปาทา อกุสลสฺส วิปากโต มโนวิญฺญาณ\-ธาตุ อุเปกฺขา\-สหคตา กิริยา\-เหตุก\-มโนวิญฺญาณ\-ธาตุ โสมนสฺส\-สหคตา, อิเม ธมฺมา สิยา อตีตารมฺมณา สิยา อนาคตารมฺมณา สิยา ปจฺจุปฺปนฺนา\-รมฺมณา. กามาวจร\-กุสลํ อกุสลํ กิริยโต นว จิตฺตุปฺปาทา รูปาวจรํ จตุตฺถํ ฌานํ กุสลโต จ กิริยโต จ, อิเม ธมฺมา สิยา อตีตารมฺมณา สิยา อนาคตารมฺมณา สิยา ปจฺจุปฺปนฺนา\-รมฺมณา. สิยา น วตฺตพฺพา อตีตารมฺมณาติปิ อนาคตารมฺมณาติปิ ปจฺจุปฺปนฺนา\-รมฺมณาติปิ. รูปาวจร\-ติกจตุกฺกชฺฌานา กุสลโต จ วิปากโต จ กิริยโต จ จตุตฺถสฺส ฌานสฺส วิปาโก อากาสา\-นญฺจา\-ยตนํ อากิญฺจญฺญา\-ยตนํ จตฺตาโร มคฺคา อปริยาปนฺนา จตฺตาริ จ สามญฺญาผลานิ, อิเม ธมฺมา น วตฺตพฺพา อตีตารมฺมณาติปิ อนาคตารมฺมณาติปิ ปจฺจุปฺปนฺนา\-รมฺมณาติปิ. รูปญฺจ นิพฺพานญฺจ อนารมฺมณา.}

\proseitem{1435}{อนินฺทฺริยพทฺธ\-รู\-ปญฺจ นิพฺพานญฺจ ฐเปตฺวา สพฺเพ ธมฺมา สิยา อชฺฌตฺตา สิยา พหิทฺธา สิยา อชฺฌตฺต\-พหิทฺธา. อนินฺทฺริยพทฺธ\-รู\-ปญฺจ นิพฺพานญฺจ พหิทฺธา.}

\proseitem{1436}{กตเม ธมฺมา อชฺฌตฺตา\-รมฺมณา. วิญฺญาณญฺจา\-ยตนํ เนว\-สญฺญา\-นา\-สญฺญา\-ยตนํ. อิเม ธมฺมา อชฺฌตฺตา\-รมฺมณา.}

\proseitem{1437}{กตเม ธมฺมา พหิทฺธา\-รมฺมณา. รูปาวจร\-ติกจตุกฺกชฺฌานา กุสลโต จ วิปากโต จ กิริยโต จ จตุตฺถสฺส ฌานสฺส วิปาโก อากาสา\-นญฺจา\-ยตนํ จตฺตาโร มคฺคา อปริยาปนฺนา จตฺตาริ จ สามญฺญผลานิ. อิเม ธมฺมา พหิทฺธา\-รมฺมณา. รูปํ ฐเปตฺวา สพฺเพว กามาวจรา กุสลากุสลา\-พฺยากตา ธมฺมา รูปาวจรํ จตุตฺถํ ฌานํ กุสลโต จ กิริยโต จ, อิเม ธมฺมา สิยา อชฺฌตฺตา\-รมฺมณา \csromanfolio{274}สิยา พหิทฺธา\-รมฺมณา สิยา อชฺฌตฺตพหิทฺธา\-รมฺมณา. อากิญฺจญฺญา\-ยตนํ น วตฺตพฺพํ อชฺฌตฺตารมฺมณานฺติปิ พหิทฺธารมฺมณนฺติปิ อชฺฌตฺตพหิทฺธารมฺมณนฺติปิ. รูปญฺจ นิพฺพานญฺจ อนารมฺมณา.}

\proseitem{1438}{กตเม ธมฺมา สนิทสฺสน\-สปฺปฏิฆา. รูปายตนํ. อิเม ธมฺมา สนิทสฺสน\-สปฺปฏิฆา.}

\proseitem{1439}{กตเม ธมฺมา อนิทสฺสน\-สปฺปฏิฆา. จกฺขายตนํ ฯเปฯ โผฏฺฐพฺพา\-ยตนํ. อิเม ธมฺมา อนิทสฺสน\-สปฺปฏิฆา.}

\proseitem{1440}{กตเม ธมฺมา อนิทสฺสน\-อปฺปฏิฆา. จตูสุ ภูมีสุ กุสลํ อกุสลํ จตูสุ ภูมีสุ วิปาโก ตีสุ ภูมีสุ กิริยาพฺยากตํ, ยญฺจ รูปํ อนิทสฺสนํ อปฺปฏิฆํ ธมฺมายตน\-ปริยาปนฺนํ นิพฺพานญฺจ. อิเม ธมฺมา อนิทสฺสน\-อปฺปฏิฆา.}

\vspace{\tipitakaparskip}
\csromancenter{ติกํ.}
\csromanheader{2}{ทุกอตฺถุทฺธาร}
\csromanmatikaanchor{mtk.120}
\csromantocmarkref{section}{ทุกอตฺถุทฺธารเหตุโคจฺฉก}{mtk.120}
\bookmark[dest={mtk.120},level=1]{ทุกอตฺถุทฺธารเหตุโคจฺฉก}
\csromanheader{1}{\textbf{เหตุโคจฺฉก}}
\proseitem{1441}{กตเม ธมฺมา เหตู. ตโย กุสลเหตู ตโย อกุสลเหตู ตโย อพฺยากตเหตู, อโลโภ กุสลเหตุ อโทโส กุสลเหตุ จตูสุ ภูมีสุ กุสเลสุ อุปฺปชฺชนฺติ, อโมโห กุสลเหตุ กามาวจร\-กุสลโต จตฺตาโร ญาณวิปฺปยุตฺเต จิตฺตุปฺปาเท ฐเปตฺวา จตูสุ ภูมีสุ กุสเลสุ อุปฺปชฺชติ.}

\prose{โลโภ อฏฺฐสุ โลภ\-สหคเตสุ จิตฺตุปฺปาเทสุ อุปฺปชฺชติ, โทโส ทฺวีสุ โทมนสฺส\-สหคเตสุ จิตฺตุปฺปาเทสุ อุปฺปชฺชติ, โมโห สพฺพากุสเลสุ อุปฺปชฺชติ.}

\prose{อโลโภ วิปากเหตุ อโทโส วิปากเหตุ กามาวจรสฺส วิปากโต อเหตุเก จิตฺตุปฺปาเท ฐเปตฺวา จตูสุ ภูมีสุ วิปาเกสุ อุปฺปชฺชนฺติ, อโมโห วิปากเหตุ กามาวจรสฺส วิปากโต \csromanfolio{275}อเหตุเก จิตฺตุปฺปาเท ฐเปตฺวา จตฺตาโร ญาณวิปฺปยุตฺเต จิตฺตุปฺปาเท ฐเปตฺวา จตูสุ ภูมีสุ วิปาเกสุ อุปฺปชฺชติ.}

\prose{อโลโภ กิริยเหตุ อโทโส กิริยเหตุ กามาวจร\-กิริยโต อเหตุเก จิตฺตุปฺปาเท ฐเปตฺวา ตีสุ ภูมีสุ กิริเยสุ อุปฺปชฺชนฺติ, อโมโห กิริยเหตุ กามาวจร\-กิริยโต อเหตุเก จิตฺตุปฺปาเท ฐเปตฺวา จตฺตาโร ญาณวิปฺปยุตฺเต จิตฺตุปฺปาเท ฐเปตฺวา ตีสุ ภูมีสุ กิริเยสุ อุปฺปชฺชติ. อิเม ธมฺมา เหตู.}

\proseitem{1442}{กตเม ธมฺมา น เหตู. ฐเปตฺวา เหตู จตูสุ ภูมีสุ กุสลํ อกุสลํ จตูสุ ภูมีสุ วิปาโก ตีสุ ภูมีสุ กิริยาพฺยากตํ รูปญฺจ นิพฺพานญฺจ. อิเม ธมฺมา น เหตู.}

\proseitem{1443}{กตเม ธมฺมา สเหตุกา. วิจิกิจฺฉา\-สหคตํ อุทฺธจฺจ\-สหคตํ โมหํ ฐเปตฺวา อวเสสํ อกุสลํ จตูสุ ภูมีสุ กุสลํ กามาวจรสฺส วิปากโต อเหตุเก จิตฺตุปฺปาเท ฐเปตฺวา จตูสุ ภูมีสุ วิปาโก กามาวจร\-กิริยโต อเหตุเก จิตฺตุปฺปาเท ฐเปตฺวา ตีสุ ภูมีสุ กิริยาพฺยากตํ. อิเม ธมฺมา สเหตุกา.}

\proseitem{1444}{กตเม ธมฺมา อเหตุกา. วิจิกิจฺฉา\-สหคโต โมโห อุทฺธจฺจ\-สหคโต โมโห เทฺว\-ปญฺจวิญฺญาณานิ ติสฺโส จ มโนธาตุโย ปญฺจ จ อเหตุก\-มโนวิญฺญาณธาตุโย รูปญฺจ นิพฺพานญฺจ. อิเม ธมฺมา อเหตุกา.}

\proseitem{1445}{กตเม ธมฺมา เหตุสมฺปยุตฺตา. วิจิกิจฺฉา\-สหคตํ อุทฺธจฺจ\-สหคตํ โมหํ ฐเปตฺวา อวเสสํ อกุสลํ จตูสุ ภูมีสุ กุสลํ กามาวจรสฺส วิปากโต อเหตุเก จิตฺตุปฺปาเท ฐเปตฺวา จตูสุ ภูมีสุ วิปาโก กามาวจร\-กิริยโต อเหตุเก จิตฺตุปฺปาเท ฐเปตฺวา ตีสุ ภูมีสุ กิริยาพฺยากตํ. อิเม ธมฺมา เหตุสมฺปยุตฺตา.}

\proseitem{1446}{กตเม ธมฺมา เหตุวิปฺปยุตฺตา. วิจิกิจฺฉา\-สหคโต โมโห อุทฺธจฺจ\-สหคโต โมโห เทฺว\-ปญฺจวิญฺญาณานิ ติสฺโส จ มโนธาตุโย ปญฺจ จ อเหตุก\-มโนวิญฺญาณธาตุโย รูปญฺจ นิพฺพานญฺจ. อิเม ธมฺมา เหตุวิปฺปยุตฺตา.}

\proseitem{1447}{\csromanfolio{276}กตเม ธมฺมา เหตู เจว สเหตุกา จ. ยตฺถ เทฺว ตโย เหตู เอกโต อุปฺปชฺชนฺติ. อิเม ธมฺมา เหตู เจว สเหตุกา จ.}

\proseitem{1448}{กตเม ธมฺมา สเหตุกา เจว น จ เหตู. จตูสุ ภูมีสุ กุสลํ อกุสลํ กามาวจรสฺส วิปากโต อเหตุเก จิตฺตุปฺปาเท ฐเปตฺวา จตูสุ ภูมีสุ วิปาโก กามาวจร\-กิริยโต อเหตุเก จิตฺตุปฺปาเท ฐเปตฺวา ตีสุ ภูมีสุ กิริยาพฺยากตํ เอตฺถุปฺปนฺเน เหตู ฐเปตฺวา. อิเม ธมฺมา สเหตุกา เจว น จ เหตู. อเหตุกา ธมฺมา น วตฺตพฺพา เหตู เจว สเหตุกา จาติปิ สเหตุกา เจว น จ เหตูหิปิ.}

\proseitem{1449}{กตเม ธมฺมา เหตู เจว เหตุสมฺปยุตฺตา จ. ยตฺถ เทฺว ตโย เหตู เอกโต อุปฺปชฺชนฺติ. อิเม ธมฺมา เหตู เจว เหตุสมฺปยุตฺตา จ.}

\proseitem{1450}{กตเม ธมฺมา เหตุสมฺปยุตฺตา เจว น จ เหตู. จตูสุ ภูมีสุ กุสลํ อกุสลํ กามาวจรสฺส วิปากโต อเหตุเก จิตฺตุปฺปาเท ฐเปตฺวา จตูสุ ภูมีสุ วิปาโก กามาวจร\-กิริยโต อเหตุเก จิตฺตุปฺปาเท ฐเปตฺวา ตีสุ ภูมีสุ กิริยาพฺยากตํ เอตฺถุปฺปนฺเน เหตู ฐเปตฺวา. อิเม ธมฺมา เหตุสมฺปยุตฺตา เจว น จ เหตู. เหตุวิปฺปยุตฺตา ธมฺมา น วตฺตพฺพา เหตู เจว เหตุสมฺปยุตฺตา จาติปิ เหตุสมฺปยุตฺตา เจว น จ เหตูติปิ.}

\proseitem{1451}{กตเม ธมฺมา น เหตู สเหตุกา. จตูสุ ภูมีสุ กุสลํ อกุสลํ กามาวจรสฺส วิปากโต อเหตุเก จิตฺตุปฺปาเท ฐเปตฺวา จตูสุ ภูมีสุ วิปาโก กามาวจร\-กิริยโต อเหตุเก จิตฺตุปฺปาเท ฐเปตฺวา ตีสุ ภูมีสุ กิริยาพฺยากตํ เอตฺถุปฺปนฺเน เหตู ฐเปตฺวา. อิเม ธมฺมา น เหตู สเหตุกา.}

\proseitem{1452}{กตเม ธมฺมา น เหตู อเหตุกา. เทฺว\-ปญฺจวิญฺญาณานิ ติสฺโส จ มโนธาตุโย ปญฺจ จ อเหตุก\-มโนวิญฺญาณธาตุโย รูปญฺจ นิพฺพานญฺจ. อิเม ธมฺมา น เหตู อเหตุกา. เหตู ธมฺมา น วตฺตพฺพา น เหตู สเหตุกาติปิ น เหตู อเหตุกาติปิ.}

\csromanmatikaanchor{mtk.121}
\csromantocmarkref{section}{จูฬนฺตรทุก}{mtk.121}
\bookmark[dest={mtk.121},level=1]{จูฬนฺตรทุก}
\csromanheader{1}{4. อฏฺฐกถา\-กณฺฑ จูฬนฺตรทุก}
\proseitem{1453}{\csromanfolio{277}กตเม ธมฺมา สปฺปจฺจยา. จตูสุ ภูมีสุ กุสลํ อกุสลํ จตูสุ ภูมีสุ วิปาโก ตีสุ ภูมีสุ กิริยาพฺยากตํ สพฺพญฺจ รูปํ. อิเม ธมฺมา สปฺปจฺจยา.}

\proseitem{1454}{กตเม ธมฺมา อปฺปจฺจยา. นิพฺพานํ. อิเม ธมฺมา อปฺปจฺจยา.}

\proseitem{1455}{กตเม ธมฺมา สงฺขตา. จตูสุ ภูมีสุ กุสลํ อกุสลํ จตูสุ ภูมีสุ วิปาโก ตีสุ ภูมีสุ กิริยาพฺยากตํ สพฺพญฺจ รูปํ. อิเม ธมฺมา สงฺขตา.}

\proseitem{1456}{กตเม ธมฺมา อสงฺขตา. นิพฺพานํ. อิเม ธมฺมา อสงฺขตา.}

\proseitem{1457}{กตเม ธมฺมา สนิทสฺสนา. รูปายตนํ. อิเม ธมฺมา สนิทสฺสนา.}

\proseitem{1458}{กตเม ธมฺมา อนิทสฺสนา. จกฺขายตนํ ฯเปฯ โผฏฺฐพฺพา\-ยตนํ จตูสุ ภูมีสุ กุสลํ อกุสลํ จตูสุ ภูมีสุ วิปาโก ตีสุ ภูมีสุ กิริยาพฺยากตํ, ยญฺจ รูปํ อนิทสฺสนํ อปฺปฏิฆํ ธมฺมายตน\-ปริยาปนฺนํ นิพฺพานญฺจ. อิเม ธมฺมา อนิทสฺสนา.}

\proseitem{1459}{กตเม ธมฺมา สปฺปฏิฆา. จกฺขายตนํ ฯเปฯ โผฏฺฐพฺพา\-ยตนํ. อิเม ธมฺมา สปฺปฏิฆา.}

\proseitem{1460}{กตเม ธมฺมา อปฺปฏิฆา. จตูสุ ภูมีสุ กุสลํ อกุสลํ จตูสุ ภูมีสุ วิปาโก ตีสุ ภูมีสุ กิริยาพฺยากตํ, ยญฺจ รูปํ อนิทสฺสนํ อปฺปฏิฆํ ธมฺมายตน\-ปริยาปนฺนํ นิพฺพานญฺจ. อิเม ธมฺมา อปฺปฏิฆา.}

\proseitem{1461}{กตเม ธมฺมา รูปิโน. จตฺตาโร จ มหาภูตา จตุนฺนญฺจ มหาภูตานํ อุปาทาย รูปํ. อิเม ธมฺมา รูปิโน.}

\proseitem{1462}{กตเม ธมฺมา อรูปิโน. จตูสุ ภูมีสุ กุสลํ อกุสลํ จตูสุ ภูมีสุ วิปาโก ตีสุ ภูมีสุ กิริยาพฺยากตํ นิพฺพานญฺจ. อิเม ธมฺมา อรูปิโน.}

\proseitem{1463}{\csromanfolio{278}กตเม ธมฺมา โลกิยา. ตีสุ ภูมีสุ กุสลํ อกุสลํ ตีสุ ภูมีสุ วิปาโก ตีสุ ภูมีสุ กิริยาพฺยากตํ สพฺพญฺจ รูปํ. อิเม ธมฺมา โลกิยา.}

\proseitem{1464}{กตเม ธมฺมา โลกุตฺตรา. จตฺตาโร มคฺคา อปริยาปนฺนา จตฺตาริ จ สามญฺญผลานิ นิพฺพานญฺจ. อิเม ธมฺมา โลกุตฺตรา. สพฺเพ\csromansharedfootnote{5d6e586a11729e53}{สพฺเพว (สฺยา)} ธมฺมา เกนจิ วิญฺเญยฺยา เกนจิ น วิญฺเญยฺยา.}

\csromanmatikaanchor{mtk.122}
\csromantocmarkref{section}{อาสวโคจฺฉก}{mtk.122}
\bookmark[dest={mtk.122},level=1]{อาสวโคจฺฉก}
\csromanheader{1}{\textbf{อาสวโคจฺฉก}}
\proseitem{1465}{กตเม ธมฺมา อาสวา. จตฺตาโร อาสวา กามาสโว ภวาสโว ทิฏฺฐาสโว อวิชฺชาสโว. กามาสโว อฏฺฐสุ โลภ\-สหคเตสุ จิตฺตุปฺปาเทสุ อุปฺปชฺชติ, ภวาสโว จตูสุ ทิฏฺฐิคต\-วิปฺปยุตฺต\-โลภ\-สหคเตสุ จิตฺตุปฺปาเทสุ อุปฺปชฺชติ, ทิฏฺฐาสโว จตูสุ ทิฏฺฐิคต\-สมฺปยุตฺเตสุ จิตฺตุปฺปาเทสุ อุปฺปชฺชติ, อวิชฺชาสโว สพฺพากุสเลสุ อุปฺปชฺชติ. อิเม ธมฺมา อาสวา.}

\proseitem{1466}{กตเม ธมฺมา โน อาสวา. ฐเปตฺวา อาสเว อวเสสํ อกุสลํ จตูสุ ภูมีสุ กุสลํ จตูสุ ภูมีสุ วิปาโก ตีสุ ภูมีสุ กิริยาพฺยากตํ รูปญฺจ นิพฺพานญฺจ. อิเม ธมฺมา โน อาสวา.}

\proseitem{1467}{กตเม ธมฺมา สาสวา. ตีสุ ภูมีสุ กุสลํ อกุสลํ ตีสุ ภูมีสุ วิปาโก ตีสุ ภูมีสุ กิริยาพฺยากตํ สพฺพญฺจ รูปํ. อิเม ธมฺมา สาสวา.}

\proseitem{1468}{กตเม ธมฺมา อนาสวา. จตฺตาโร มคฺคา อปริยาปนฺนา จตฺตาริ จ สมญฺญผลานิ นิพฺพานญฺจ. อิเม ธมฺมา อนาสวา.}

\proseitem{1469}{กตเม ธมฺมา อาสว\-สมฺปยุตฺตา. เทฺว โทมนสฺสสหคต\-จิตฺตุปฺปาทา เอตฺถุปฺปนฺนํ โมหํ ฐเปตฺวา, วิจิกิจฺฉา\-สหคตํ อุทฺธจฺจ\-สหคตํ โมหํ ฐเปตฺวา อวเสสํ อกุสลํ. อิเม ธมฺมา อาสว\-สมฺปยุตฺตา.}

\chapterheadpageread{4. อฏฺฐกถา\-กณฺฑ}
\proseitem{1470}{\csromanfolio{279}กตเม ธมฺมา อาสว\-วิปฺปยุตฺตา. ทฺวีสุ โทมนสฺส\-สหคเตสุ จิตฺตุปฺปาเทสุ อุปฺปนฺโน โมโห วิจิกิจฺฉา\-สหคโต โมโห อุทฺธจฺจ\-สหคโต โมโห จตูสุ ภูมีสุ กุสลํ จตูสุ ภูมีสุ วิปาโก ตีสุ ภูมีสุ กิริยาพฺยากตํ รูปญฺจ นิพฺพานญฺจ. อิเม ธมฺมา อาสว\-วิปฺปยุตฺตา.}

\proseitem{1471}{กตเม ธมฺมา อาสโว เจว สาสวา จ. เตว อาสวา อาสวา เจว สาสวา จ.}

\proseitem{1472}{กตเม ธมฺมา สาสวา เจว โน จ อาสวา. ฐเปตฺวา อาสเว อวเสสํ อกุสลํ ตีสุ ภูมีสุ กุสลํ ตีสุ ภูมีสุ วิปาโก ตีสุ ภูมีสุ กิริยาพฺยากตํ สพฺพญฺจ รูปํ. อิเม ธมฺมา สาสวา เจว โน จ อาสวา. อนาสวา ธมฺมา น วตฺตพฺพา อาสวา เจว สาสวา จาติปิ สาสวา เจว โน จ อาสวาติปิ.}

\proseitem{1473}{กตเม ธมฺมา อาสวา เจว อาสว\-สมฺปยุตฺตา จ. ยตฺถ เทฺว ตโย อาสวา เอกโต อุปฺปชฺชนฺติ. อิเม ธมฺมา อาสวา เจว อาสว\-สมฺปยุตฺตา จ.}

\proseitem{1474}{กตเม ธมฺมา อาสว\-สมฺปยุตฺตา เจว โน จ อาสวา. ฐเปตฺวา อาสเว อวเสสํ อกุสลํ. อิเม ธมฺมา อาสว\-สมฺปยุตฺตา เจว โน จ อาสวา. อาสว\-วิปฺปยุตฺตา ธมฺมา น วตฺตพฺพา อาสวา เจว อาสว\-สมฺปยุตฺตา จาติปิ อาสว\-สมฺปยุตฺตา เจว โน จ อาสวาติปิ.}

\proseitem{1475}{กตเม ธมฺมา อาสว\-วิปฺปยุตฺตา สาสวา. ทฺวีสุ โทมนสฺส\-สหคเตสุ จิตฺตุปฺปาเทสุ อุปฺปนฺนา โมโห วิจิกิจฺฉา\-สหคโต โมโห อุทฺธจฺจ\-สหคโต โมโห ตีสุ ภูมีสุ กุสลํ ตีสุ ภูมีสุ วิปาโก ตีสุ ภูมีสุ กิริยาพฺยากตํ สพฺพญฺจ รูปํ. อิเม ธมฺมา อาสว\-วิปฺปยุตฺตา สาสวา.}

\proseitem{1476}{กตเม ธมฺมา อาสว\-วิปฺปยุตฺตา อนาสวา. จตฺตาโร มคฺคา อปริยาปนฺนา จตฺตาริ จ สามญฺญผลานิ นิพฺพานญฺจ. อิเม ธมฺมา อาสว\-วิปฺปยุตฺตา อนาสวา. อาสว\-สมฺปยุตฺตา ธมฺมา น วตฺตพฺพา อาสว\-วิปฺปยุตฺตา สาสวาติปิ อาสว\-วิปฺปยุตฺตา อนาสวาติปิ.}

\csromanmatikaanchor{mtk.123}
\csromantocmarkref{section}{สญฺโญชนโคจฺฉก}{mtk.123}
\bookmark[dest={mtk.123},level=1]{สญฺโญชนโคจฺฉก}
\csromanheader{1}{สญฺโญชน\-โคจฺฉก}
\proseitem{1477}{\csromanfolio{280}กตเม ธมฺมา สญฺโญชนา. ทส สญฺโญชนานิ กามราค\-สญฺโญชนํ ปฏิฆ\-สญฺโญชนํ มานสญฺโญชนํ ทิฏฺฐิ\-สญฺโญชนํ วิจิกิจฺฉา\-สญฺโญชนํ สีลพฺพตปรามาส\-สญฺโญชนํ ภวราค\-สญฺโญชนํ อิสฺสาสญฺโญชนํ มจฺฉริย\-สญฺโญชนํ อวิชฺชา\-สญฺโญชนํ. กามราค\-สญฺโญชนํ อฏฺฐสุ โลภ\-สหคเตสุ จิตฺตุปฺปาเทสุ อุปฺปชฺชติ, ปฏิฆ\-สญฺโญชนํ ทฺวีสุ โทมนสฺส\-สหคเตสุ จิตฺตุปฺปาเทสุ อุปฺปชฺชติ, มานสญฺโญชนํ จตูสุ ทิฏฺฐิคต\-วิปฺปยุตฺต\-โลภ\-สหคเตสุ จิตฺตุปฺปาเทสุ อุปฺปชฺชติ, ทิฏฺฐิ\-สญฺโญชนํ จตูสุ ทิฏฺฐิคต\-สมฺปยุตฺเตสุ จิตฺตุปฺปาเทสุ อุปฺปชฺชติ, วิจิกิจฺฉา\-สญฺโญชนํ วิจิกิจฺฉา\-สหคเตสุ จิตฺตุปฺปาเทสุ อุปฺปชฺชติ, สีลพฺพตปรามาส\-สญฺโญชนํ จตูสุ ทิฏฺฐิคต\-สมฺปยุตฺเตสุ จิตฺตุปฺปาเทสุ อุปฺปชฺชติ, ภวราค\-สญฺโญชนํ จตูสุ ทิฏฺฐิคต\-วิปฺปยุตฺต\-โลภ\-สหคเตสุ จิตฺตุปฺปาเทสุ อุปฺปชฺชติ, อิสฺสา\-สญฺโญชนญฺจ มจฺฉริย\-สญฺโญชนญฺจ ทฺวีสุ โทมนสฺส\-สหคเตสุ จิตฺตุปฺปาเทสุ อุปฺปชฺชนฺติ, อวิชฺชา\-สญฺโญชนํ สพฺพากุสเลสุ อุปฺปชฺชติ. อิเม ธมฺมา สญฺโญชนา.}

\proseitem{1478}{กตเม ธมฺมา โน สญฺโญชนา. ฐเปตฺวา สญฺโญชเน อวเสสํ อกุสลํ จตูสุ ภูมีสุ กุสลํ จตูสุ ภูมีสุ วิปาโก ตีสุ ภูมีสุ กิริยาพฺยากตํ รูปญฺจ นิพฺพานญฺจ. อิเม ธมฺมา โน สญฺโญชนา.}

\proseitem{1479}{กตเม ธมฺมา สญฺโญชนิยา. ตีสุ ภูมีสุ กุสลํ อกุสลํ ตีสุ ภูมีสุ วิปาโก ตีสุ ภูมีสุ กิริยาพฺยากตํ สพฺพญฺจ รูปํ. อิเม ธมฺมา สญฺโญชนิยา.}

\proseitem{1480}{กตเม ธมฺมา อสญฺโญชนิยา. จตฺตาโร มคฺคา อปริยาปนฺนา จตฺตาริ จ สามญฺญผลานิ นิพฺพานญฺจ. อิเม ธมฺมา อสญฺโญชนิยา.}

\proseitem{1481}{กตเม ธมฺมา สญฺโญชน\-สมฺปยุตฺตา. อุทฺธจฺจ\-สหคตํ โมหํ ฐเปตฺวา อวเสสํ อกุสลํ. อิเม ธมฺมา สญฺโญชน\-สมฺปยุตฺตา.}

\proseitem{1482}{กตเม ธมฺมา สญฺโญชน\-วิปฺปยุตฺตา. อุทฺธจฺจ\-สหคโต โมโห จตูสุ ภูมีสุ กุสลํ จตูสุ ภูมีสุ วิปาโก ตีสุ ภูมีสุ กิริยาพฺยากตํ รูปญฺจ นิพฺพานญฺจ. อิเม ธมฺมา สญฺโญชน\-วิปฺปยุตฺตา.}

\proseitem{1483}{กตเม ธมฺมา สญฺโญชนา เจว สญฺโญชนิยา จ. ตาเนว สญฺโญชนานิ สญฺโญชนา เจว สญฺโญชนิยา จ.}

\chapterheadpageread{4. อฏฺฐกถา\-กณฺฑ}
\proseitem{1484}{\csromanfolio{281}กตเม ธมฺมา สญฺโญชนิยา เจว โน จ สญฺโญชนา. ฐเปตฺวา สญฺโญชเน อวเสสํ อกุสลํ ตีสุ ภูมีสุ กุสลํ ตีสุ ภูมีสุ วิปาโก ตีสุ ภูมีสุ กิริยาพฺยากตํ สพฺพญฺจ รูปํ. อิเม ธมฺมา สญฺโญชนิยา เจว โน จ สญฺโญชนา. อสญฺโญชนิยา ธมฺมา น วตฺตพฺพา สญฺโญชนา เจว สญฺโญชนิยา จาติปิ สญฺโญชนิยา เจว โน จ สญฺโญชนาติปิ.}

\proseitem{1485}{กตเม ธมฺมา สญฺโญชนา เจว สญฺโญชน\-สมฺปยุตฺตา จ. ยตฺถ เทฺว ตีณิ สญฺโญชนานิ เอกโต อุปฺปชฺชนฺติ. อิเม ธมฺมา สญฺโญชนา เจว สญฺโญชน\-สมฺปยุตฺตา จ.}

\proseitem{1486}{กตเม ธมฺมา สญฺโญชน\-สมฺปยุตฺตา เจว โน จ สญฺโญชนา. ฐเปตฺวา สญฺโญชเน อวเสสํ อกุสลํ. อิเม ธมฺมา สญฺโญชน\-สมฺปยุตฺตา เจว โน จ สญฺโญชนา. สญฺโญชน\-วิปฺปยุตฺตา ธมฺมา น วตฺตพฺพา สญฺโญชนา เจว สญฺโญชน\-สมฺปยุตฺตา จาติปิ สญฺโญชน\-สมฺปยุตฺตา เจว โน จ สญฺโญชนาติปิ.}

\proseitem{1487}{กตเม ธมฺมา สญฺโญชน\-วิปฺปยุตฺตา สญฺโญชนิยา. อุทฺธจฺจ\-สหคโต โมโห ตีสุ ภูมีสุ กุสลํ ตีสุ ภูมีสุ วิปาโก ตีสุ ภูมีสุ กิริยาพฺยากตํ สพฺพญฺจ รูปํ. อิเม ธมฺมา สญฺโญชน\-วิปฺปยุตฺตา สญฺโญชนิยา.}

\proseitem{1488}{กตเม ธมฺมา สญฺโญชน\-วิปฺปยุตฺตา อสญฺโญชนิยา. จตฺตาโร มคฺคา อปริยาปนฺนา จตฺตาริ จ สามญฺญผลานิ นิพฺพานญฺจ. อิเม ธมฺมา สญฺโญชน\-วิปฺปยุตฺตา อสญฺโญชนิยา. สญฺโญชน\-สมฺปยุตฺตา ธมฺมา น วตฺตพฺพา สญฺโญชน\-วิปฺปยุตฺตา สญฺโญชนิยาติปิ สญฺโญชน\-วิปฺปยุตฺตา อสญฺโญชนิยาติปิ.}

\csromanmatikaanchor{mtk.124}
\csromantocmarkref{section}{คนฺถโคจฺฉก}{mtk.124}
\bookmark[dest={mtk.124},level=1]{คนฺถโคจฺฉก}
\csromanheader{1}{\textbf{คนฺถโคจฺฉก}}
\proseitem{1489}{กตเม ธมฺมา คนฺถา. จตฺตาโร คนฺถา อภิชฺฌา\-กายคนฺโถ พฺยาปาโท กายคนฺโถ สีลพฺพต\-ปรามาโส กายคนฺโถ อิทํสจฺจา\-ภินิเวโส กายคนฺโถ. อภิชฺฌา\-กายคนฺโถ อฏฺฐสุ \csromanfolio{282}โลภ\-สหคเตสุ จิตฺตุปฺปาเทสุ อุปฺปชฺชติ, พฺยาปาโท กายคนฺโถ ทฺวีสุ โทมนสฺส\-สหคเตสุ จิตฺตุปฺปาเทสุ อุปฺปชฺชติ, สีลพฺพต\-ปรามาโส กายคนฺโถ จ อิทํสจฺจา\-ภินิเวโส กายคนฺโถ จ จตูสุ ทิฏฺฐิคต\-สมฺปยุตฺเตสุ จิตฺตุปฺปาเทสุ อุปฺปชฺชนฺติ. อิเม ธมฺมา คนฺถา.}

\proseitem{1490}{กตเม ธมฺมา โน คนฺถา. ฐเปตฺวา คนฺเถ อวเสสํ อกุสลํ จตูสุ ภูมีสุ กุสลํ จตูสุ ภูมีสุ วิปาโก ตีสุ ภูมีสุ กิริยาพฺยากตํ รูปญฺจ นิพฺพานญฺจ. อิเม ธมฺมา โน คนฺถา.}

\proseitem{1491}{กตเม ธมฺมา คนฺถนิยา. ตีสุ ภูมีสุ กุสลํ อกุสลํ ตีสุ ภูมีสุ วิปาโก ตีสุ ภูมีสุ กิริยาพฺยากตํ สพฺพญฺจ รูปํ. อิเม ธมฺมา คนฺถนิยา.}

\proseitem{1492}{กตเม ธมฺมา อคนฺถนิยา. จตฺตาโร มคฺคา อปริยาปนฺนา จตฺตาริ จ สามญฺญผลานิ นิพฺพานญฺจ. อิเม ธมฺมา อคนฺถนิยา.}

\proseitem{1493}{กตเม ธมฺมา คนฺถ\-สมฺปยุตฺตา. จตฺตาโร ทิฏฺฐิคต\-สมฺปยุตฺต\-จิตฺตุปฺปาทา จตฺตาโร ทิฏฺฐิคต\-วิปฺปยุตฺต\-โลภ\-สหคต\-จิตฺตุปฺปาทา เอตฺถุปฺปนฺนํ โลภํ ฐเปตฺวา, เทฺว โทมนสฺสสหคต\-จิตฺตุปฺปาทา เอตฺถุปฺปนฺนํ ปฏิฆํ ฐเปตฺวา. อิเม ธมฺมา คนฺถ\-สมฺปยุตฺตา.}

\proseitem{1494}{กตเม ธมฺมา คนฺถ\-วิปฺปยุตฺตา. จตูสุ ทิฏฺฐิคต\-วิปฺปยุตฺต\-โลภ\-สหคเตสุ จิตฺตุปฺปาเทสุ อุปฺปนฺโน โลโภ ทฺวีสุ โทมนสฺส\-สหคเตสุ จิตฺตุปฺปาเทสุ อุปฺปนฺนํ ปฏิฆํ วิจิกิจฺฉา\-สหคโต จิตฺตุปฺปาโท อุทฺธจฺจ\-สหคโต จิตฺตุปฺปาโท จตูสุ ภูมีสุ กุสลํ จตูสุ ภูมีสุ วิปาโก ตีสุ ภูมีสุ กิริยาพฺยากตํ รูปญฺจ นิพฺพานญฺจ. อิเม ธมฺมา คนฺถ\-วิปฺปยุตฺตา.}

\proseitem{1495}{กตเม ธมฺมา คนฺถา เจว คนฺถนิยา จ. เตว คนฺถา คนฺถา เจว คนฺถนิยา จ.}

\proseitem{1496}{กตเม ธมฺมา คนฺถนิยา เจว โน จ คนฺถา. ฐเปตฺวา คนฺเถ อวเสสํ อกุสลํ ตีสุ ภูมีสุ กุสลํ ตีสุ ภูมีสุ วิปาโก ตีสุ ภูมีสุ กิริยาพฺยากตํ สพฺพญฺจ รูปํ. อิเม ธมฺมา คนฺถนิยา เจว โน จ คนฺถา. อคนฺถนิยา ธมฺมา น วตฺตพฺพา คนฺถา เจว คนฺถนิยา จาติปิ คนฺถนิยา เจว โน จ คนฺถาติปิ.}

\chapterheadpageread{4. อฏฺฐกถา\-กณฺฑ}
\proseitem{1497}{\csromanfolio{283}กตเม ธมฺมา คนฺถา เจว คนฺถ\-สมฺปยุตฺตา จ. ยตฺถ ทิฏฺฐิ จ โลโภ จ เอกโต อุปฺปชฺชนฺติ. อิเม ธมฺมา คนฺถา เจว คนฺถ\-สมฺปยุตฺตา จ.}

\proseitem{1498}{กตเม ธมฺมา คนฺถ\-สมฺปยุตฺตา เจว โน จ คนฺถา. อฏฺฐ โลภ\-สหคต\-จิตฺตุปฺปาทา เทฺว โทมนสฺสสหคต\-จิตฺตุปฺปาทา เอตฺถุปฺปนฺเน คนฺเถ ฐเปตฺวา. อิเม ธมฺมา คนฺถ\-สมฺปยุตฺตา เจว โน จ คนฺถา. คนฺถ\-วิปฺปยุตฺตา ธมฺมา น วตฺตพฺพา คนฺถา เจว คนฺถ\-สมฺปยุตฺตา จาติปิ คนฺถ\-สมฺปยุตฺตา เจว โน จ คนฺถาติปิ.}

\proseitem{1499}{กตเม ธมฺมา คนฺถ\-วิปฺปยุตฺตา คนฺถนิยา. จตูสุ ทิฏฺฐิคต\-วิปฺปยุตฺต\-โลภ\-สหคเตสุ จิตฺตุปฺปาเทสุ อุปฺปนฺโน โลโภ ทฺวีสุ โทมนสฺส\-สหคเตสุ จิตฺตุปฺปาเทสุ อุปฺปนฺนํ ปฏิฆํ วิจิกิจฺฉา\-สหคโต จิตฺตุปฺปาโท อุทฺธจฺจ\-สหคโต จิตฺตุปฺปาโท ตีสุ ภูมีสุ กุสลํ ตีสุ ภูมีสุ วิปาโก ตีสุ ภูมีสุ กิริยาพฺยากตํ สพฺพญฺจ รูปํ. อิเม ธมฺมา คนฺถ\-วิปฺปยุตฺตา คนฺถนิยา.}

\proseitem{1500}{กตเม ธมฺมา คนฺถ\-วิปฺปยุตฺตา อคนฺถนิยา. จตฺตาโร มคฺคา อปริยาปนฺนา จตฺตาริ จ สามญฺญผลานิ นิพฺพานญฺจ. อิเม ธมฺมา คนฺถ\-วิปฺปยุตฺตา อคนฺถนิยา. คนฺถ\-สมฺปยุตฺตา ธมฺมา น วตฺตพฺพา คนฺถ\-วิปฺปยุตฺตา คนฺถนิยาติปิ คนฺถ\-วิปฺปยุตฺตา อคนฺถนิยาติปิ.}

\csromanmatikaanchor{mtk.125}
\csromantocmarkref{section}{โอฆโคจฺฉก}{mtk.125}
\bookmark[dest={mtk.125},level=1]{โอฆโคจฺฉก}
\csromanheader{1}{\textbf{โอฆโคจฺฉก}}
\vspace{\tipitakaparskip}
\csromancenter{กตเม ธมฺมา โอฆา ฯเปฯ.}
\csromanmatikaanchor{mtk.126}
\csromantocmarkref{section}{โยคโคจฺฉก}{mtk.126}
\bookmark[dest={mtk.126},level=1]{โยคโคจฺฉก}
\csromanheader{1}{\textbf{โยคโคจฺฉก}}
\vspace{\tipitakaparskip}
\csromancenter{กตเม ธมฺมา โยคา ฯเปฯ.}
\csromanmatikaanchor{mtk.127}
\csromantocmarkref{section}{นีวรณโคจฺฉก}{mtk.127}
\bookmark[dest={mtk.127},level=1]{นีวรณโคจฺฉก}
\csromanheader{1}{นีวรณ\-โคจฺฉก}
\proseitem{1503}{กตเม ธมฺมา นีวรณา. ฉ นีวรณา กามจฺฉนฺท\-นีวรณํ พฺยาปาท\-นีวรณํ ถินมิทฺธ\-นีวรณํ อุทฺธจฺจ\-กุกฺกุจฺจ\-นีวรณํ วิจิกิจฺฉา\-นีวรณํ \csromanfolio{284}อวิชฺชา\-นีวรณํ. กามจฺฉนฺท\-นีวรณํ อฏฺฐสุ โลภ\-สหคเตสุ จิตฺตุปฺปาเทสุ อุปฺปชฺชติ, พฺยาปาท\-นีวรณํ ทฺวีสุ โทมนสฺส\-สหคเตสุ จิตฺตุปฺปาเทสุ อุปฺปชฺชติ, ถินมิทฺธ\-นีวรณํ สสงฺขาริเกสุ อกุสเลสุ อุปฺปชฺชติ, อุทฺธจฺจ\-นีวรณํ อุทฺธจฺจ\-สหคเตสุ จิตฺตุปฺปาเทสุ อุปฺปชฺชติ, กุกฺกุจฺจ\-นีวรณํ ทฺวีสุ โทมนสฺส\-สหคเตสุ จิตฺตุปฺปาเทสุ อุปฺปชฺชติ, วิจิกิจฺฉา\-นีวรณํ วิจิกิจฺฉา\-สหคเตสุ จิตฺตุปฺปาเทสุ อุปฺปชฺชติ, อวิชฺชา\-นีวรณํ สพฺพากุสเลสุ อุปฺปชฺชติ. อิเม ธมฺมา นีวรณา.}

\proseitem{1504}{กตเม ธมฺมา โน นีวรณา. ฐเปตฺวา นีวรเณ อวเสสํ อกุสลํ จตูสุ ภูมีสุ กุสลํ จตูสุ ภูมีสุ วิปาโก ตีสุ ภูมีสุ กิริยาพฺยากตํ รูปญฺจ นิพฺพานญฺจ. อิเม ธมฺมา โน นีวรณา.}

\proseitem{1505}{กตเม ธมฺมา นีวรณิยา. ตีสุ ภูมีสุ กุสลํ อกุสลํ ตีสุ ภูมีสุ วิปาโก ตีสุ ภูมีสุ กิริยาพฺยากตํ สพฺพญฺจ รูปํ. อิเม ธมฺมา นีวรณิยา.}

\proseitem{1506}{กตเม ธมฺมา อนีวรณิยา. จตฺตาโร มคฺคา อปริยาปนฺนา จตฺตาริ จ สามญฺญผลานิ นิพฺพานญฺจ. อิเม ธมฺมา อนีวรณิยา.}

\proseitem{1507}{กตเม ธมฺมา นีวรณ\-สมฺปยุตฺตา. ทฺวาทส อกุสล\-จิตฺตุปฺปาทา. อิเม ธมฺมา นีวรณ\-สมฺปยุตฺตา.}

\proseitem{1508}{กตเม ธมฺมา นีวรณ\-วิปฺปยุตฺตา. จตูสุ ภูมีสุ กุสลํ จตูสุ ภูมีสุ วิปาโก ตีสุ ภูมีสุ กิริยาพฺยากตํ รูปญฺจ นิพฺพานญฺจ. อิเม ธมฺมา นีวรณ\-วิปฺปยุตฺตา.}

\proseitem{1509}{กตเม ธมฺมา นีวรณา เจว นีวรณิยา จ. ตาเนว นีวรณานิ นีวรณา เจว นีวรณิยา จ.}

\proseitem{1510}{กตเม ธมฺมา นีวรณิยา เจว โน จ นีวรณา. ฐเปตฺวา นีวรเณ อวเสสํ อกุสลํ ตีสุ ภูมีสุ กุสลํ ตีสุ ภูมีสุ วิปาโก ตีสุ ภูมีสุ กิริยาพฺยากตํ สพฺพญฺจ รูปํ. อิเม ธมฺมา นีวรณิยา เจว โน จ นีวรณา. อนีวรณิยา ธมฺมา น วตฺตพฺพา นีวรณา เจว นีวรณิยา จาติปิ นีวรณิยา เจว โน จ นีวรณาติปิ.}

\chapterheadpageread{4. อฏฺฐกถา\-กณฺฑ}
\proseitem{1511}{\csromanfolio{285}กตเม ธมฺมา นีวรณา เจว นีวรณ\-สมฺปยุตฺตา จ. ยตฺถ เทฺว ตีณิ นีวรณานิ เอกโต อุปฺปชฺชนฺติ. อิเม ธมฺมา นีวรณา เจว นีวรณ\-สมฺปยุตฺตา จ.}

\proseitem{1512}{กตเม ธมฺมา นีวรณ\-สมฺปยุตฺตา เจว โน จ นีวรณา. ฐเปตฺวา นีวรเณ อวเสสํ อกุสลํ. อิเม ธมฺมา นีวรณ\-สมฺปยุตฺตา เจว โน จ นีวรณา. นีวรณ\-วิปฺปยุตฺตา ธมฺมา น วตฺตพฺพา นีวรณา เจว นีวรณ\-สมฺปยุตฺตา จาติปิ นีวรณ\-สมฺปยุตฺตา เจว โน จ นีวรณาติปิ.}

\proseitem{1513}{กตเม ธมฺมา นีวรณ\-วิปฺปยุตฺตา นีวรณิยา. ตีสุ ภูมีสุ กุสลํ ตีสุ ภูมีสุ วิปาโก ตีสุ ภูมีสุ กิริยาพฺยากตํ สพฺพญฺจ รูปํ. อิเม ธมฺมา นีวรณ\-วิปฺปยุตฺตา นีวรณิยา.}

\proseitem{1514}{กตเม ธมฺมา นีวรณ\-วิปฺปยุตฺตา อนีวรณิยา. จตฺตาโร มคฺคา อปริยาปนฺนา จตฺตาริ จ สามญฺญผลานิ นิพฺพานญฺจ. อิเม ธมฺมา นีวรณ\-วิปฺปยุตฺตา อนีวรณิยา. นีวรณ\-สมฺปยุตฺตา ธมฺมา น วตฺตพฺพา นีวรณ\-วิปฺปยุตฺตา นีวรณิยาติปิ นีวรณ\-วิปฺปยุตฺตา อนีวรณิยาติปิ.}

\csromanmatikaanchor{mtk.128}
\csromantocmarkref{section}{ปรามาสโคจฺฉก}{mtk.128}
\bookmark[dest={mtk.128},level=1]{ปรามาสโคจฺฉก}
\csromanheader{1}{ปรามาส\-โคจฺฉก}
\proseitem{1515}{กตเม ธมฺมา ปรามาสา. ทิฏฺฐิปรามาโส จตูสุ ทิฏฺฐิคต\-สมฺปยุตฺเตสุ จิตฺตุปฺปาเทสุ อุปฺปชฺชติ. อิเม ธมฺมา ปรามาสา.}

\proseitem{1516}{กตเม ธมฺมา โน ปรามาสา. ฐเปตฺวา ปรามาสํ อวเสสํ อกุสลํ จตูสุ ภูมีสุ กุสลํ จตูสุ ภูมีสุ วิปาโก ตีสุ ภูมีสุ กิริยาพฺยากตํ รูปญฺจ นิพฺพานญฺจ. อิเม ธมฺมา โน ปรามาสา.}

\proseitem{1517}{กตเม ธมฺมา ปรามฏฺฐา. ตีสุ ภูมีสุ กุสลํ อกุสลํ ตีสุ ภูมีสุ วิปาโก ตีสุ ภูมีสุ กิริยาพฺยากตํ สพฺพญฺจ รูปํ. อิเม ธมฺมา ปรามฏฺฐา.}

\proseitem{1518}{\csromanfolio{286}กตเม ธมฺมา อปรามฏฺฐา. จตฺตาโร มคฺคา อปริยาปนฺนา จตฺตาริ จ สามญฺญผลานิ นิพฺพานญฺจ. อิเม ธมฺมา อปรามฏฺฐา.}

\proseitem{1519}{กตเม ธมฺมา ปรามาส\-สมฺปยุตฺตา. จตฺตาโร ทิฏฺฐิคต\-สมฺปยุตฺต\-จิตฺตุปฺปาทา เอตฺถุปฺปนฺนํ ปรามาสํ ฐเปตฺวา. อิเม ธมฺมา ปรามาส\-สมฺปยุตฺตา.}

\proseitem{1520}{กตเม ธมฺมา ปรามาส\-วิปฺปยุตฺตา. จตฺตาโร ทิฏฺฐิคต\-วิปฺปยุตฺต\-โลภ\-สหคต\-จิตฺตุปฺปาทา เทฺว โทมนสฺสสหคต\-จิตฺตุปฺปาทา วิจิกิจฺฉา\-สหคโต จิตฺตุปฺปาโท อุทฺธจฺจ\-สหคโต จิตฺตุปฺปาโท จตูสุ ภูมีสุ กุสลํ จตูสุ ภูมีสุ วิปาโก ตีสุ ภูมีสุ กิริยาพฺยากตํ รูปญฺจ นิพฺพานญฺจ. อิเม ธมฺมา ปรามาส\-วิปฺปยุตฺตา. ปรามาโส น วตฺตพฺโพ ปรามาสสมฺปยุตฺโตติปิ ปรามาสวิปฺปยุตฺโตติปิ.}

\proseitem{1521}{กตเม ธมฺมา ปรามาสา เจว ปรามฏฺฐา จ. โส เอว ปรามาโส ปรามาโส เจว ปรามฏฺโฐ จ.}

\proseitem{1522}{กตเม ธมฺมา ปรามฏฺฐา เจว โน จ ปรามาสา. ฐเปตฺวา ปรามาสํ อวเสสํ อกุสลํ ตีสุ ภูมีสุ กุสลํ ตีสุ ภูมีสุ วิปาโก ตีสุ ภูมีสุ กิริยาพฺยากตํ สพฺพญฺจ รูปํ. อิเม ธมฺมา ปรามฏฺฐา เจว โน จ ปรามาสา. อปรามฏฺฐา ธมฺมา น วตฺตพฺพา ปรามาสา เจว ปรามฏฺฐา จาติปิ ปรามฏฺฐา เจว โน จ ปรามาสาติปิ.}

\proseitem{1523}{กตเม ธมฺมา ปรามสวิปฺปยุตฺตา ปรามฏฺฐา. จตฺตาโร ทิฏฺฐิคต\-วิปฺปยุตฺต\-โลภ\-สหคต\-จิตฺตุปฺปาทา เทฺว โทมนสฺสสหคต\-จิตฺตุปฺปาทา วิจิกิจฺฉา\-สหคโต จิตฺตุปฺปาโท อุทฺธจฺจ\-สหคโต จิตฺตุปฺปาโท ตีสุ ภูมีสุ กุสลํ ตีสุ ภูมีสุ วิปาโก ตีสุ ภูมีสุ กิริยาพฺยากตํ สพฺพญฺจ รูปํ. อิเม ธมฺมา ปรามาส\-วิปฺปยุตฺตา ปรามฏฺฐา.}

\proseitem{1524}{กตเม ธมฺมา ปรามาส\-วิปฺปยุตฺตา อปรามฏฺฐา. จตฺตาโร มคฺคา อปริยาปนฺนา จตฺตาริ จ สามญฺญผลานิ นิพฺพานญฺจ. อิเม ธมฺมา ปรามาส\-วิปฺปยุตฺตา อปรามฏฺฐา. ปรามาสา จ ปรามาส\-สมฺปยุตฺตา จ ธมฺมา น วตฺตพฺพา ปรามาส\-วิปฺปยุตฺตา ปรามฏฺฐาติปิ ปรามาส\-วิปฺปยุตฺตา อปรามฏฺฐาติปิ.}

\csromanmatikaanchor{mtk.129}
\csromantocmarkref{section}{มหนฺตรทุก}{mtk.129}
\bookmark[dest={mtk.129},level=1]{มหนฺตรทุก}
\csromanheader{1}{4. อฏฺฐกถา\-กณฺฑ มหนฺตรทุก}
\proseitem{1525}{\csromanfolio{287}กตเม ธมฺมา สารมฺมณา. จตูสุ ภูมีสุ กุสลํ อกุสลํ จตูสุ ภูมีสุ วิปาโก ตีสุ ภูมีสุ กิริยาพฺยากตํ. อิเม ธมฺมา สารมฺมณา.}

\proseitem{1526}{กตเม ธมฺมา อนารมฺมณา. รูปญฺจ นิพฺพานญฺจ. อิเม ธมฺมา อนารมฺมณา.}

\proseitem{1527}{กตเม ธมฺมา จิตฺตา. จกฺขุวิญฺญาณํ โสตวิญฺญาณํ ฆานวิญฺญาณํ ชิวฺหาวิญฺญาณํ กายวิญฺญาณํ มโนธาตุ มโนวิญฺญาณ\-ธาตุ. อิเม ธมฺมา จิตฺตา.}

\proseitem{1528}{กตเม ธมฺมา โน จิตฺตา. เวทนากฺขนฺโธ สญฺญากฺขนฺโธ สงฺขาร\-กฺขนฺโธ รูปญฺจ นิพฺพานญฺจ. อิเม ธมฺมา โน จิตฺตา.}

\proseitem{1529}{กตเม ธมฺมา เจตสิกา. เวทนากฺขนฺโธ สญฺญากฺขนฺโธ. สงฺขาร\-กฺขนฺโธ. อิเม ธมฺมา เจตสิกา.}

\proseitem{1530}{กตเม ธมฺมา อเจตสิกา. จิตฺตญฺจ รูปญฺจ นิพฺพานญฺจ. อิเม ธมฺมา อเจตสิกา.}

\proseitem{1531}{กตเม ธมฺมา จิตฺต\-สมฺปยุตฺตา. เวทนากฺขนฺโธ สญฺญากฺขนฺโธ สงฺขาร\-กฺขนฺโธ. อิเม ธมฺมา จิตฺต\-สมฺปยุตฺตา.}

\proseitem{1532}{กตเม ธมฺมา จิตฺต\-วิปฺปยุตฺตา. รูปญฺจ นิพฺพานญฺจ. อิเม ธมฺมา จิตฺต\-วิปฺปยุตฺตา. จิตฺตํ น วตฺตพฺพํ จิตฺเตน สมฺปยุตฺตนฺติปิ จิตฺเตน วิปฺปยุตฺตนฺติปิ.}

\proseitem{1533}{กตเม ธมฺมา จิตฺตสํสฏฺฐา. เวทนากฺขนฺโธ สญฺญากฺขนฺโธ สงฺขาร\-กฺขนฺโธ. อิเม ธมฺมา จิตฺตสํสฏฺฐา.}

\proseitem{1534}{กตเม ธมฺมา จิตฺต\-วิสํสฏฺฐา. รูปญฺจ นิพฺพานญฺจ. อิเม ธมฺมา จิตฺต\-วิสํสฏฺฐา. จิตฺตํ น วตฺตพฺพํ จิตฺเตน สํสฏฺฐนฺติปิ จิตฺเตน วิสํสฏฺฐนฺติปิ.}

\proseitem{1535}{กตเม ธมฺมา จิตฺต\-สมุฏฺฐานา. เวทนากฺขนฺโธ สญฺญาขนฺโธ สงฺขาร\-กฺขนฺโธ กายวิญฺญตฺติ วจีวิญฺญตฺติ, ยํ วา ปนญฺญมฺปิ อตฺถิ รูปํ จิตฺตชํ \csromanfolio{288}จิตฺตเหตุกํ จิตฺต\-สมุฏฺฐานํ รูปายตนํ สทฺทายตนํ คนฺธายตนํ รสายตนํ โผฏฺฐพฺพา\-ยตนํ อากาสธาตุ อาโปธาตุ รูปสฺส ลหุตา รูปสฺส มุทุตา รูปสฺส กมฺมญฺญตา รูปสฺส อุปจโย รูปสฺส สนฺตติ กพฬีกาโร อาหาโร. อิเม ธมฺมา จิตฺต\-สมุฏฺฐานา.}

\proseitem{1536}{กตเม ธมฺมา โน จิตฺต\-สมุฏฺฐานา. จิตฺตญฺจ อวเสสญฺจ รูปํ นิพฺพานญฺจ. อิเม ธมฺมา โน จิตฺต\-สมุฏฺฐานา.}

\proseitem{1537}{กตเม ธมฺมา จิตฺตสหภุโน. เวทนากฺขนฺโธ สญฺญากฺขนฺโธ สงฺขาร\-กฺขนฺโธ กายวิญฺญตฺติ วจีวิญฺญตฺติ. อิเม ธมฺมา จิตฺตสหภุโน.}

\proseitem{1538}{กตเม ธมฺมา โน จิตฺตสหภุโน. จิตฺตญฺจ อวเสสญฺจ รูปํ นิพฺพานญฺจ. อิเม ธมฺมา โน จิตฺตสหภุโน.}

\proseitem{1539}{กตเม ธมฺมา จิตฺตา\-นุปริวตฺติโน. เวทนากฺขนฺโธ สญฺญากฺขนฺโธ สงฺขาร\-กฺขนฺโธ กายวิญฺญตฺติ วจีวิญฺญตฺติ. อิเม ธมฺมา จิตฺตา\-นุปริวตฺติโน.}

\proseitem{1540}{กตเม ธมฺมา โน จิตฺตา\-นุปริวตฺติโน. จิตฺตญฺจ อวเสสญฺจ รูปํ นิพฺพานญฺจ. อิเม ธมฺมา โน จิตฺตา\-นุปริวตฺติโน.}

\proseitem{1541}{กตเม ธมฺมา จิตฺต\-สํสฏฺฐ\-สมุฏฺฐานา. เวทนากฺขนฺโธ สญฺญากฺขนฺโธ สงฺขาร\-กฺขนฺโธ. อิเม ธมฺมา จิตฺต\-สํสฏฺฐ\-สมุฏฺฐานา.}

\proseitem{1542}{กตเม ธมฺมา โน จิตฺต\-สํสฏฺฐ\-สมุฏฺฐานา. จิตฺตญฺจ รูปญฺจ นิพฺพานญฺจ. อิเม ธมฺมา โน จิตฺต\-สํสฏฺฐ\-สมุฏฺฐานา.}

\proseitem{1543}{กตเม ธมฺมา จิตฺต\-สํสฏฺฐ\-สมุฏฺฐาน\-สหภุโน. เวทนากฺขนฺโธ สญฺญากฺขนฺโธ สงฺขาร\-กฺขนฺโธ. อิเม ธมฺมา จิตฺต\-สํสฏฺฐ\-สมุฏฺฐาน\-สหภุโน.}

\proseitem{1544}{กตเม ธมฺมา โน จิตฺต\-สํสฏฺฐ\-สมุฏฺฐาน\-สหภุโน. จิตฺตญฺจ รูปญฺจ นิพฺพานญฺจ. อิเม ธมฺมา โน จิตฺต\-สํสฏฺฐ\-สมุฏฺฐาน\-สหภุโน.}

\proseitem{1545}{กตเม ธมฺมา จิตฺต\-สํสฏฺฐ\-สมุฏฺฐานา\-นุปริวตฺติโน. เวทนากฺขนฺโธ สญฺญากฺขนฺโธ สงฺขาร\-กฺขนฺโธ. อิเม ธมฺมา จิตฺต\-สํสฏฺฐ\-สมุฏฺฐานา\-นุปริวตฺติโน.}

\chapterheadpageread{4. อฏฺฐกถา\-กณฺฑ}
\proseitem{1546}{\csromanfolio{289}กตเม ธมฺมา โน จิตฺต\-สํสฏฺฐ\-สมุฏฺฐานา\-นุปริวตฺติโน. จิตฺตญฺจ รูปญฺจ นิพฺพานญฺจ. อิเม ธมฺมา โน จิตฺต\-สํสฏฺฐ\-สมุฏฺฐานา\-นุปริวตฺติโน.}

\proseitem{1547}{กตเม ธมฺมา อชฺฌตฺติกา. จกฺขายตนํ ฯเปฯ มนายตนํ. อิเม ธมฺมา อชฺฌตฺติกา.}

\proseitem{1548}{กตเม ธมฺมา พาหิรา. รูปายตนํ ฯเปฯ ธมฺมายตนํ. อิเม ธมฺมา พาหิรา.}

\proseitem{1549}{กตเม ธมฺมา อุปาทา. จกฺขายตนํ ฯเปฯ กพฬีกาโร อาหาโร. อิเม ธมฺมา อุปาทา.}

\proseitem{1550}{กตเม ธมฺมา โน อุปาทา. จตูสุ ภูมีสุ กุสลํ อกุสลํ จตูสุ ภูมีสุ วิปาโก ตีสุ ภูมีสุ กิริยาพฺยากตํ จตฺตาโร จ มหาภูตา นิพฺพานญฺจ. อิเม ธมฺมา โน อุปาทา.}

\proseitem{1551}{กตเม ธมฺมา อุปาทิณฺณา. ตีสุ ภูมีสุ วิปาโก, ยญฺจ รูปํ กมฺมสฺส กตตฺตา. อิเม ธมฺมา อุปาทิณฺณา.}

\proseitem{1552}{กตเม ธมฺมา อนุปาทิณฺณา. ตีสุ ภูมีสุ กุสลํ อกุสลํ ตีสุ ภูมีสุ กิริยาพฺยากตํ, ยญฺจ รูปํ น กมฺมสฺส กตตฺตา จตฺตาโร มคฺคา อปริยาปนฺนา จตฺตาริ จ สามญฺญผลานิ นิพฺพานญฺจ. อิเม ธมฺมา อนุปาทิณฺณา.}

\csromanmatikaanchor{mtk.130}
\csromantocmarkref{section}{อุปาทานโคจฺฉก}{mtk.130}
\bookmark[dest={mtk.130},level=1]{อุปาทานโคจฺฉก}
\csromanheader{1}{อุปาทาน\-โคจฺฉก}
\proseitem{1553}{กตเม ธมฺมา อุปาทานา. จตฺตาริ อุปาทานานิ กามุปาทานํ ทิฏฺฐุปาทานํ สีลพฺพตุ\-ปาทานํ อตฺตวาทุ\-ปาทานํ. กามุปาทานํ อฏฺฐสุ โลภ\-สหคเตสุ จิตฺตุปฺปาเทสุ อุปฺปชฺชติ, ทิฏฺฐุปาทานญฺจ สีลพฺพตุปาทานญฺจ อตฺตวาทุปาทานญฺจ จตูสุ ทิฏฺฐิคต\-สมฺปยุตฺเตสุ จิตฺตุปฺปาเทสุ อุปฺปชฺชนฺติ. อิเม ธมฺมา อุปาทานา.}

\proseitem{1554}{กตเม ธมฺมา โน อุปาทานา. ฐเปตฺวา อุปาทาเน อวเสสํ อกุสลํ จตูสุ ภูมีสุ กุสลํ จตูสุ ภูมีสุ วิปาโก ตีสุ ภูมีสุ กิริยาพฺยากตํ รูปญฺจ นิพฺพานญฺจ. อิเม ธมฺมา โน อุปาทานา.}

\proseitem{1555}{\csromanfolio{290}กตเม ธมฺมา อุปาทานิยา. ตีสุ ภูมีสุ กุสลํ อกุสลํ ตีสุ ภูมีสุ วิปาโก ตีสุ ภูมีสุ กิริยาพฺยากตํ สพฺพญฺจ รูปํ. อิเม ธมฺมา อุปาทานิยา.}

\proseitem{1556}{กตเม ธมฺมา อนุปาทานิยา. จตฺตาโร มคฺคา อปริยาปนฺนา จตฺตาริ จ สามญฺญผลานิ นิพฺพานญฺจ. อิเม ธมฺมา อนุปาทานิยา.}

\proseitem{1557}{กตเม ธมฺมา อุปาทาน\-สมฺปยุตฺตา. จตฺตาโร ทิฏฺฐิคต\-สมฺปยุตฺต\-โลภ\-สหคต\-จิตฺตุปฺปาทา จตฺตาโร ทิฏฺฐิคต\-วิปฺปยุตฺต\-โลภ\-สหคต\-จิตฺตุปฺปาทา เอตฺถุปฺปนฺนํ โลภํ ฐเปตฺวา. อิเม ธมฺมา อุปาทาน\-สมฺปยุตฺตา.}

\proseitem{1558}{กตเม ธมฺมา อุปาทาน\-วิปฺปยุตฺตา. จตูสุ ทิฏฺฐิคต\-วิปฺปยุตฺต\-โลภ\-สหคเตสุ จิตฺตุปฺปาเทสุ อุปฺปนฺโน โลโภ เทฺว โทมนสฺสสหคต\-จิตฺตุปฺปาทา วิจิกิจฺฉา\-สหคโต จิตฺตุปฺปาโท อุทฺธจฺจ\-สหคโต จิตฺตุปฺปาโท จตูสุ ภูมีสุ กุสลํ จตูสุ ภูมีสุ วิปาโก ตีสุ ภูมีสุ กิริยาพฺยากตํ รูปญฺจ นิพฺพานญฺจ. อิเม ธมฺมา อุปาทาน\-วิปฺปยุตฺตา.}

\proseitem{1559}{กตเม ธมฺมา อุปาทานา เจว อุปาทานิยา จ. ตาเนว อุปาทานานิ อุปาทานา เจว อุปาทานิยา จ.}

\proseitem{1560}{กตเม ธมฺมา อุปาทานิยา เจว โน จ อุปาทานา. ฐเปตฺวา อุปาทาเน อวเสสํ อกุสลํ ตีสุ ภูมีสุ กุสลํ ตีสุ ภูมีสุ วิปาโก ตีสุ ภูมีสุ กิริยาพฺยากตํ สพฺพญฺจ รูปํ. อิเม ธมฺมา อุปาทานิยา เจว โน จ อุปาทานา. อนุปาทานิยา ธมฺมา น วตฺตพฺพา อุปาทานา เจว อุปาทานิยา จาติปิ อุปาทานิยา เจว โน จ อุปาทานาติปิ.}

\proseitem{1561}{กตเม ธมฺมา อุปาทานา เจว อุปาทาน\-สมฺปยุตฺตา จ. ยตฺถ ทิฏฺฐิ จ โลโภ จ เอกโต อุปฺปชฺชนฺติ. อิเม ธมฺมา อุปาทานา เจว อุปาทาน\-สมฺปยุตฺตา จ.}

\proseitem{1562}{กตเม ธมฺมา อุปาทาน\-สมฺปยุตฺตา เจว โน จ อุปาทานา. อฏฺฐ\-โลภ\-สหคต\-จิตฺตุปฺปาทา เอตฺถุปฺปนฺเน อุปาทาเน ฐเปตฺวา. อิเม ธมฺมา อุปาทาน\-สมฺปยุตฺตา เจว โน จ อุปาทานา. อุปาทาน\-วิปฺปยุตฺตา ธมฺมา น วตฺตพฺพา อุปาทานา เจว อุปาทาน\-สมฺปยุตฺตา จาติปิ อุปาทาน\-สมฺปยุตฺตา เจว โน จ อุปาทานาติปิ.}

\chapterheadpageread{4. อฏฺฐกถา\-กณฺฑ}
\proseitem{1563}{\csromanfolio{291}กตเม ธมฺมา อุปาทาน\-วิปฺปยุตฺตา อุปาทานิยา. จตูสุ ทิฏฺฐิคต\-วิปฺปยุตฺต\-โลภ\-สหคเตสุ จิตฺตุปฺปาเทสุ อุปฺปนฺโน โลโภ เทฺว โทมนสฺสสหคต\-จิตฺตุปฺปาทา วิจิกิจฺฉา\-สหคโต จิตฺตุปฺปาโท อุทฺธจฺจ\-สหคโต จิตฺตุปฺปาโท ตีสุ ภูมีสุ กุสลํ ตีสุ ภูมีสุ วิปาโก ตีสุ ภูมีสุ กิริยาพฺยากตํ สพฺพญฺจ รูปํ. อิเม ธมฺมา อุปาทาน\-วิปฺปยุตฺตา อุปาทานิยา.}

\proseitem{1564}{กตเม ธมฺมา อุปาทาน\-วิปฺปยุตฺตา อนุปาทานิยา. จตฺตาโร มคฺคา อปริยาปนฺนา จตฺตาริ จ สามญฺญผลานิ นิพฺพานญฺจ. อิเม ธมฺมา อุปาทาน\-วิปฺปยุตฺตา อนุปาทานิยา. อุปาทาน\-สมฺปยุตฺตา ธมฺมา น วตฺตพฺพา อุปาทาน\-วิปฺปยุตฺตา อุปาทานิยาติปิ อุปาทาน\-วิปฺปยุตฺตา อนุปาทานิยาติปิ.}

\csromanmatikaanchor{mtk.131}
\csromantocmarkref{section}{กิเลสโคจฺฉก}{mtk.131}
\bookmark[dest={mtk.131},level=1]{กิเลสโคจฺฉก}
\csromanheader{1}{\textbf{กิเลสโคจฺฉก}}
\proseitem{1565}{กตเม ธมฺมา กิเลสา. ทส กิเลสวตฺถูนิ โลโภ โทโส โมโห มาโน ทิฏฺฐิ วิจิกิจฺฉา ถินํ อุทฺธจฺจํ อหิริกํ อโนตฺตปฺปํ. โลโภ อฏฺฐสุ โลภ\-สหคเตสุ จิตฺตุปฺปาเทสุ อุปฺปชฺชติ, โทโส ทฺวีสุ โทมนสฺส\-สหคเตสุ จิตฺตุปฺปาเทสุ อุปฺปชฺชติ, โมโห สพฺพากุสเลสุ อุปฺปชฺชติ, มาโน จตูสุ ทิฏฺฐิคต\-วิปฺปยุตฺต\-โลภ\-สหคเตสุ จิตฺตุปฺปาเทสุ อุปฺปชฺชติ, ทิฏฺฐิ จตูสุ ทิฏฺฐิคต\-สมฺปยุตฺเตสุ จิตฺตุปฺปาเทสุ อุปฺปชฺชติ, วิจิกิจฺฉา วิจิกิจฺฉา\-สหคเตสุ จิตฺตุปฺปาเทสุ อุปฺปชฺชติ, ถินํ สสงฺขาริเกสุ อกุสเลสุ อุปฺปชฺชติ, อุทฺธจฺจญฺจ อหิริกญฺจ อโนตฺตปฺปญฺจ สพฺพากุสเลสุ อุปฺปชฺชนฺติ. อิเม ธมฺมา กิเลสา.}

\proseitem{1566}{กตเม ธมฺมา โน กิเลสา. ฐเปตฺวา กิเลเส อวเสสํ อกุสลํ จตูสุ ภูมีสุ กุสลํ จตูสุ ภูมีสุ วิปาโก ตีสุ ภูมีสุ กิริยาพฺยากตํ รูปญฺจ นิพฺพานญฺจ. อิเม ธมฺมา โน กิเลสา.}

\proseitem{1567}{กตเม ธมฺมา สํกิเลสิกา. ตีสุ ภูมีสุ กุสลํ อกุสลํ ตีสุ ภูมีสุ วิปาโก ตีสุ ภูมีสุ กิริยาพฺยากตํ สพฺพญฺจ รูปํ. อิเม ธมฺมา สํกิเลสิกา.}

\proseitem{1568}{กตเม ธมฺมา อสํกิเลสิกา. จตฺตาโร มคฺคา อปริยาปนฺนา จตฺตาริ จ สามญฺญผลานิ นิพฺพานญฺจ. อิเม ธมฺมา อสํกิเลสิกา.}

\proseitem{1569}{\csromanfolio{292}กตเม ธมฺมา สํกิลิฏฺฐา. ทฺวาทส อกุสล\-จิตฺตุปฺปาทา. อิเม ธมฺมา สํกิลิฏฺฐา.}

\proseitem{1570}{กตเม ธมฺมา อสํกิลิฏฺฐา. จตูสุ ภูมีสุ กุสลํ จตูสุ ภูมีสุ วิปาโก ตีสุ ภูมีสุ กิริยาพฺยากตํ รูปญฺจ นิพฺพานญฺจ. อิเม ธมฺมา อสํกิลิฏฺฐา.}

\proseitem{1571}{กตเม ธมฺมา กิเลส\-สมฺปยุตฺตา. ทฺวาทส อกุสล\-จิตฺตุปฺปาทา. อิเม ธมฺมา กิเลส\-สมฺปยุตฺตา.}

\proseitem{1572}{กตเม ธมฺมา กิเลส\-วิปฺปยุตฺตา. จตูสุ ภูมีสุ กุสลํ จตูสุ ภูมีสุ วิปาโก ตีสุ ภูมีสุ กิริยาพฺยากตํ รูปญฺจ นิพฺพานญฺจ. อิเม ธมฺมา กิเลส\-วิปฺปยุตฺตา.}

\proseitem{1573}{กตเม ธมฺมา กิเลสา เจว สํกิเลสิกา จ. เตว กิเลสา กิเลสา เจว สํกิเลสิกา จ.}

\proseitem{1574}{กตเม ธมฺมา สํกิเลสิกา เจว โน จ กิเลสา. ฐเปตฺวา กิเลเส อวเสสํ อกุสลํ ตีสุ ภูมีสุ กุสลํ ตีสุ ภูมีสุ วิปาโก ตีสุ ภูมีสุ กิริยาพฺยากตํ สพฺพญฺจ รูปํ. อิเม ธมฺมา สํกิเลสิกา เจว โน จ กิเลสา. อสํกิเลสิกา ธมฺมา น วตฺตพฺพา กิเลสา เจว สํกิเลสิกา จาติปิ สํกิเลสิกา เจว โน จ กิเลสาติปิ.}

\proseitem{1575}{กตเม ธมฺมา กิเลสา เจว สํกิลิฏฺฐา จ. เตว กิเลสา กิเลสา เจว สํกิลิฏฺฐา จ.}

\proseitem{1576}{กตเม ธมฺมา สํกิลิฏฺฐา เจว โน จ กิเลสา. ฐเปตฺวา กิเลเส อวเสสํ อกุสลํ. อิเม ธมฺมา สํกิลิฏฺฐา เจว โน จ กิเลสา. อสํกิลิฏฺฐา ธมฺมา น วตฺตพฺพา กิเลสา เจว สํกิลิฏฺฐา จาติปิ สํกิลิฏฺฐา เจว โน จ กิเลสาติปิ.}

\proseitem{1577}{กตเม ธมฺมา กิเลสา เจว กิเลส\-สมฺปยุตฺตา จ. ยตฺถ เทฺว ตโย กิเลสา เอกโต อุปฺปชฺชนฺติ. อิเม ธมฺมา กิเลสา เจว กิเลส\-สมฺปยุตฺตา จ.}

\chapterheadpageread{4. อฏฺฐกถา\-กณฺฑ}
\proseitem{1578}{\csromanfolio{293}กตเม ธมฺมา กิเลส\-สมฺปยุตฺตา เจว โน จ กิเลสา. ฐเปตฺวา กิเลเส อวเสสํ อกุสลํ. อิเม ธมฺมา กิเลส\-สมฺปยุตฺตา เจว โน จ กิเลสา. กิเลส\-วิปฺปยุตฺตา ธมฺมา น วตฺตพฺพา กิเลสา เจว กิเลส\-สมฺปยุตฺตา จาติปิ กิเลส\-สมฺปยุตฺตา เจว โน จ กิเลสาติปิ.}

\proseitem{1579}{กตเม ธมฺมา กิเลส\-วิปฺปยุตฺตา สํกิเลสิกา. ตีสุ ภูมีสุ กุสลํ ตีสุ ภูมีสุ วิปาโก ตีสุ ภูมีสุ กิริยาพฺยากตํ สพฺพญฺจ รูปํ. อิเม ธมฺมา กิเลส\-วิปฺปยุตฺตา สํกิเลสิกา.}

\proseitem{1580}{กตเม ธมฺมา กิเลส\-วิปฺปยุตฺตา อสํกิเลสิกา. จตฺตาโร มคฺคา อปริยาปนฺนา จตฺตาริ จ สามญฺญผลานิ นิพฺพานญฺจ. อิเม ธมฺมา กิเลส\-วิปฺปยุตฺตา อสํกิเลสิกา. กิเลส\-สมฺปยุตฺตา ธมฺมา น วตฺตพฺพา กิเลส\-วิปฺปยุตฺตา สํกิเลสิกาติปิ กิเลส\-วิปฺปยุตฺตา อสํกิเลสิกาติปิ.}

\csromanmatikaanchor{mtk.132}
\csromantocmarkref{section}{ปิฏฺฐิทุก}{mtk.132}
\bookmark[dest={mtk.132},level=1]{ปิฏฺฐิทุก}
\csromanheader{1}{\textbf{ปิฏฺฐิทุก}}
\proseitem{1581}{กตเม ธมฺมา ทสฺสเนน ปหาตพฺพา. จตฺตาโร ทิฏฺฐิคต\-สมฺปยุตฺต\-จิตฺตุปฺปาทา วิจิกิจฺฉา\-สหคโต จิตฺตุปฺปาโท. อิเม ธมฺมา ทสฺสเนน ปหาตพฺพา. จตฺตาโร ทิฏฺฐิคต\-วิปฺปยุตฺต\-โลภ\-สหคต\-จิตฺตุปฺปาทา เทฺว โทมนสฺสสหคต\-จิตฺตุปฺปาทา. อิเม ธมฺมา สิยา ทสฺสเนน ปหาตพฺพา, สิยา น ทสฺสเนน ปหาตพฺพา.}

\proseitem{1582}{กตเม ธมฺมา น ทสฺสเนน ปหาตพฺพา. อุทฺธจฺจ\-สหคโต จิตฺตุปฺปาโท จตูสุ ภูมีสุ กุสลํ จตูสุ ภูมีสุ วิปาโก ตีสุ ภูมีสุ กิริยาพฺยากตํ รูปญฺจ นิพฺพานญฺจ. อิเม ธมฺมา น ทสฺสเนน ปหาตพฺพา.}

\proseitem{1583}{กตเม ธมฺมา ภาวนาย ปหาตพฺพา. อุทฺธจฺจ\-สหคโต จิตฺตุปฺปาโท. อิเม ธมฺมา ภาวนาย ปหาตพฺพา. จตฺตาโร ทิฏฺฐิคต\-วิปฺปยุตฺต\-โลภ\-สหคต\-จิตฺตุปฺปาทา เทฺว โทมนสฺสสหคต\-จิตฺตุปฺปาทา. อิเม ธมฺมา สิยา ภาวนาย ปหาตพฺพา, สิยา น ภาวนาย ปหาตพฺพา.}

\proseitem{1584}{\csromanfolio{294}กตเม ธมฺมา น ภาวนาย ปหาตพฺพา. จตฺตาโร ทิฏฺฐิคต\-สมฺปยุตฺต\-จิตฺตุปฺปาทา วิจิกิจฺฉา\-สหคโต จิตฺตุปาโท จตูสุ ภูมีสุ กุสลํ จตูสุ ภูมีสุ วิปาโก ตีสุ ภูมีสุ กิริยาพฺยากตํ รูปญฺจ นิพฺพานญฺจ. อิเม ธมฺมา น ภาวนาย ปหาตพฺพา.}

\proseitem{1585}{กตเม ธมฺมา ทสฺสเนน ปหาตพฺพ\-เหตุกา. จตฺตาโร ทิฏฺฐิคต\-สมฺปยุตฺต\-จิตฺตุปฺปาทา วิจิกิจฺฉา\-สหคโต จิตฺตุปฺปาโท เอตฺถุปฺปนฺนํ โมหํ ฐเปตฺวา. อิเม ธมฺมา ทสฺสเนน ปหาตพฺพ\-เหตุกา. จตฺตาโร ทิฏฺฐิคต\-วิปฺปยุตฺต\-โลภ\-สหคต\-จิตฺตุปฺปาทา เทฺว โทมนสฺสสหคต\-จิตฺตุปฺปาทา. อิเม ธมฺมา สิยา ทสฺสเนน ปหาตพฺพ\-เหตุกา, สิยา น ทสฺสเนน ปหาตพฺพ\-เหตุกา.}

\proseitem{1586}{กตเม ธมฺมา น ทสฺสเนน ปหาตพฺพ\-เหตุกา. วิจิกิจฺฉา\-สหคโต โมโห อุทฺธจฺจ\-สหคโต จิตฺตุปฺปาโท จตูสุ ภูมีสุ กุสลํ จตูสุ ภูมีสุ วิปาโก ตีสุ ภูมีสุ กิริยาพฺยากตํ รูปญฺจ นิพฺพานญฺจ. อิเม ธมฺมา น ทสฺสเนน ปหาตพฺพ\-เหตุกา.}

\proseitem{1587}{กตเม ธมฺมา ภาวนาย ปหาตพฺพ\-เหตุกา. อุทฺธจฺจ\-สหคโต จิตฺตุปฺปาโท เอตฺถุปฺปนฺนํ โมหํ ฐเปตฺวา. อิเม ธมฺมา ภาวนาย ปหาตพฺพ\-เหตุกา. จตฺตาโร ทิฏฺฐิคต\-วิปฺปยุตฺต\-โลภ\-สหคต\-จิตฺตุปฺปาทา เทฺว โทมนสฺสสหคต\-จิตฺตุปฺปาทา. อิเม ธมฺมา สิยา ภาวนาย ปหาตพฺพ\-เหตุกา, สิยา น ภาวนาย ปหาตพฺพ\-เหตุกา.}

\proseitem{1588}{กตเม ธมฺมา น ภาวนาย ปหาตพฺพ\-เหตุกา. จตฺตาโร ทิฏฺฐิคต\-สมฺปยุตฺต\-จิตฺตุปฺปาทา วิจิกิจฺฉา\-สหคโต จิตฺตุปฺปาโท อุทฺธจฺจ\-สหคโต โมโห จตูสุ ภูมีสุ กุสลํ จตูสุ ภูมีสุ วิปาโก ตีสุ ภูมีสุ กิริยาพฺยากตํ รูปญฺจ นิพฺพานญฺจ. อิเม ธมฺมา น ภาวนาย ปหาตพฺพ\-เหตุกา.}

\proseitem{1589}{กตเม ธมฺมา สวิตกฺกา. กามาวจร\-กุสลํ อกุสลํ กามาวจร\-กุสลสฺส วิปากโต เอกาทส จิตฺตุปฺปาทา อกุสลสฺส วิปากโต เทฺว กิริยโต เอกาทส รูปาวจรํ ปฐมํ ฌานํ กุสลโต จ วิปากโต จ กิริยโต จ โลกุตฺตรํ ปฐมํ ฌานํ กุสลโต จ วิปากโต จ เอตฺถุปฺปนฺนํ วิตกฺกํ ฐเปตฺวา. อิเม ธมฺมา สวิตกฺกา.}

\chapterheadpageread{4. อฏฺฐกถา\-กณฺฑ}
\proseitem{1590}{\csromanfolio{295}กตเม ธมฺมา อวิตกฺกา. เทฺว\-ปญฺจวิญฺญาณานิ รูปาวจร\-ติกจตุกฺกชฺฌานา กุสลโต จ วิปากโต จ กิริยโต จ จตฺตาโร อรูปาวจรา กุสลโต จ วิปากโต จ กิริยโต จ โลกุตฺตร\-ติกจตุกฺกชฺฌานา กุสลโต จ วิปากโต จ วิตกฺโก จ รูปญฺจ นิพฺพานญฺจ. อิเม ธมฺมา อวิตกฺกา.}

\proseitem{1591}{กตเม ธมฺมา สวิจารา. กามาวจร\-กุสลํ อกุสลํ กามาวจร\-กุสลสฺส วิปากโต เอกาทส จิตฺตุปฺปาทา อกุสลสฺส วิปากโต เทฺว กิริยโต เอกาทส รูปาวจรเอกกทุกชฺฌานา กุสลโต จ วิปากโต จ กิริยโต จ โลกุตฺตรเอกกทุกชฺฌานา กุสลโต จ วิปากโต จ เอตฺถุปฺปนฺนํ วิจารํ ฐเปตฺวา. อิเม ธมฺมา สวิจารา.}

\proseitem{1592}{กตเม ธมฺมา อวิจารา. เทฺว\-ปญฺจวิญฺญาณานิ รูปาวจร\-ติก\-ติก\-ชฺฌานา กุสลโต จ วิปากโต จ กิริยโต จ จตฺตาโร อารุปฺปา กุสลโต จ วิปากโต จ กิริยโต จ โลกุตฺตร\-ติก\-ติก\-ชฺฌานา กุสลโต จ วิปากโต จ วิจาโร จ รูปญฺจ นิพฺพานญฺจ. อิเม ธมฺมา อวิจารา.}

\proseitem{1593}{กตเม ธมฺมา สปฺปีติกา. กามาวจร\-กุสลโต จตฺตาโร โสมนสฺสสหคต\-จิตฺตุปฺปาทา อกุสลโต จตฺตาโร กามาวจร\-กุสลสฺส วิปากโต ปญฺจ กิริยโต ปญฺจ รูปาวจร\-ทุ\-กติก\-ชฺฌานา กุสลโต จ วิปากโต จ กิริยโต จ โลกุตฺตร\-ทุ\-กติก\-ชฺฌานา กุสลโต จ วิปากโต จ เอตฺถุปฺปนฺนํ ปีติํ ฐเปตฺวา. อิเม ธมฺมา สปฺปีติกา.}

\proseitem{1594}{กตเม ธมฺมา อปฺปีติกา. กามาวจร\-กุสลโต จตฺตาโร อุเปกฺขาสหคต\-จิตฺตุปฺปาทา อกุสลโต อฏฺฐ กามาวจร\-กุสลสฺส วิปากโต เอกาทส อกุสลสฺส วิปากโต สตฺต กิริยโต ฉ รูปาวจร\-ทุก\-ทุก\-ชฺฌานา กุสลโต จ วิปากโต จ กิริยโต จ จตฺตาโร อารุปฺปา กุสลโต จ วิปากโต จ กิริยโต จ โลกุตฺตร\-ทุก\-ทุก\-ชฺฌานา กุสลโต จ วิปากโต จ ปีติ จ รูปญฺจ นิพฺพานญฺจ. อิเม ธมฺมา อปฺปีติกา.}

\proseitem{1595}{\csromanfolio{296}กตเม ธมฺมา ปีติสหคตา. กามาวจร\-กุสลโต จตฺตาโร โสมนสฺสสหคต\-จิตฺตุปฺปาทา อกุสลโต จตฺตาโร กามาวจร\-กุสลสฺส วิปากโต ปญฺจ กิริยโต ปญฺจ รูปาวจร\-ทุ\-กติก\-ชฺฌานา กุสลโต จ วิปากโต จ กิริยโต จ โลกุตฺตร\-ทุ\-กติก\-ชฺฌานา กุสลโต จ วิปากโต จ เอตฺถุปฺปนฺนํ ปีติํ ฐเปตฺวา. อิเม ธมฺมา ปีติสหคตา.}

\proseitem{1596}{กตเม ธมฺมา น ปีติสหคตา. กามาวจร\-กุสลโต จตฺตาโร อุเปกฺขาสหคต\-จิตฺตุปฺปาทา อกุสลโต อฏฺฐ กามาวจร\-กุสลสฺส วิปากโต เอกาทส อกุสลสฺส วิปากโต สตฺต กิริยโต ฉ รูปาวจร\-ทุก\-ทุก\-ชฺฌานา กุสลโต จ วิปากโต จ กิริยโต จ จตฺตาโร อารุปฺปา กุสลโต จ วิปากโต จ กิริยโต จ โลกุตฺตร\-ทุก\-ทุก\-ชฺฌานา กุสลโต จ วิปากโต จ ปีติ จ รูปญฺจ นิพฺพานญฺจ. อิเม ธมฺมา น ปีติสหคตา.}

\proseitem{1597}{กตเม ธมฺมา สุขสหคตา. กามาวจร\-กุสลโต จตฺตาโร โสมนสฺสสหคต\-จิตฺตุปฺปาทา อกุสลโต จตฺตาโร กามาวจร\-กุสลสฺส วิปากโต ฉ กิริยโต ปญฺจ รูปาวจร\-ติกจตุกฺกชฺฌานา กุสลโต จ วิปากโต จ กิริยโต จ โลกุตฺตร\-ติกจตุกฺกชฺฌานา กุสลโต จ วิปากโต จ เอตฺถุปฺปนฺนํ สุขํ ฐเปตฺวา. อิเม ธมฺมา สุขสหคตา.}

\proseitem{1598}{กตเม ธมฺมา น สุขสหคตา. กามาวจร\-กุสลโต จตฺตาโร อุเปกฺขาสหคต\-จิตฺตุปฺปาทา อกุสลโต อฏฺฐ กามาวจร\-กุสลสฺส วิปากโต ทส อกุสลสฺส วิปากโต สตฺต กิริยโต ฉ รูปาวจรํ จตุตฺถํ ฌานํ กุสลโต จ วิปากโต จ กิริยโต จ จตฺตาโร อารุปฺปา กุสลโต จ วิปากโต จ กิริยโต จ โลกุตฺตรํ จตุตฺถํ ฌานํ กุสลโต จ วิปากโต จ สุขญฺจ รูปญฺจ นิพฺพานญฺจ. อิเม ธมฺมา น สุขสหคตา.}

\proseitem{1599}{กตเม ธมฺมา อุเปกฺขา\-สหคตา. กามาวจร\-กุสลโต จตฺตาโร อุเปกฺขาสหคต\-จิตฺตุปฺปาทา อกุสลโต ฉ กามาวจร\-กุสลสฺส วิปากโต ทส อกุสลสฺส วิปากโต ฉ กิริยโต \csromanfolio{297}ฉ รูปาวจรํ จตุตฺถํ ฌานํ กุสลโต จ วิปากโต จ กิริยโต จ จตฺตาโร อารุปฺปา กุสลโต จ วิปากโต จ กิริยโต จ โลกุตฺตรํ จตุตฺถํ ฌานํ กุสลโต จ วิปากโต จ เอตฺถุปฺปนฺนํ อุเปกฺขํ ฐเปตฺวา. อิเม ธมฺมา อุเปกฺขา\-สหคตา.}

\proseitem{1600}{กตเม ธมฺมา น อุเปกฺขา\-สหคตา. กามาวจร\-กุสลโต จตฺตาโร โสมนสฺสสหคต\-จิตฺตุปฺปาทา อกุสลโต ฉ กามาวจร\-กุสลสฺส วิปากโต ฉ อกุสลสฺส วิปากโต เอโก กิริยโต ปญฺจ รูปาวจร\-ติกจตุกฺกชฺฌานา กุสลโต จ วิปากโต จ กิริยโต จ โลกุตฺตร\-ติกจตุกฺกชฺฌานา กุสลโต จ วิปากโต จ อุเปกฺขา จ รูปญฺจ นิพฺพานญฺจ. อิเม ธมฺมา น อุเปกฺขา\-สหคตา.}

\proseitem{1601}{กตเม ธมฺมา กามาวจรา. กามาวจร\-กุสลํ อกุสลํ สพฺโพ กามาวจรสฺส วิปาโก กามาวจร\-กิริยา\-พฺยากตํ สพฺพญฺจ รูปํ อิเม ธมฺมา กามาวจรา.}

\proseitem{1602}{กตเม ธมฺมา น กามาวจรา. รูปาวจรา อรูปาวจรา อปริยาปนฺนา. อิเม ธมฺมา น กามาวจรา.}

\proseitem{1603}{กตเม ธมฺมา รูปาวจรา. รูปาวจร\-จตุกฺกปญฺจกชฺฌานา กุสลโต จ วิปากโต จ กิริยโต จ. อิเม ธมฺมา รูปาวจรา.}

\proseitem{1604}{กตเม ธมฺมา น รูปาวจรา. กามาวจรา อรูปาวจรา อปริยาปนฺนา. อิเม ธมฺมา น รูปาวจรา.}

\proseitem{1605}{กตเม ธมฺมา น อรูปาวจรา. จตฺตาโร อารุปฺปา กุสลโต จ วิปากโต จ กิริยโต จ. อิเม ธมฺมา อรูปาวจรา.}

\proseitem{1606}{กตเม ธมฺมา น อรูปาวจรา. กามาวจรา รูปาวจรา อปริยาปนฺนา. อิเม ธมฺมา น อรูปาวจรา.}

\proseitem{1607}{กตเม ธมฺมา ปริยาปนฺนา. ตีสุ ภูมีสุ กุสลํ อกุสลํ ตีสุ ภูมีสุ วิปาโก ตีสุ ภูมีสุ กิริยาพฺยากตํ สพฺพญฺจ รูปํ. อิเม ธมฺมา ปริยาปนฺนา.}

\proseitem{1608}{กตเม ธมฺมา อปริยาปนฺนา. จตฺตาโร มคฺคา อปริยาปนฺนา จตฺตาริ จ สามญฺญผลานิ นิพฺพานญฺจ. อิเม ธมฺมา อปริยาปนฺนา.}

\proseitem{1609}{\csromanfolio{298}กตเม ธมฺมา นิยฺยานิกา. จตฺตาโร มคฺคา อปริยาปนฺนา. อิเม ธมฺมา นิยฺยานิกา.}

\proseitem{1610}{กตเม ธมฺมา อนิยฺยานิกา. ตีสุ ภูมีสุ กุสลํ อกุสลํ จตูสุ ภูมีสุ วิปาโก ตีสุ ภูมีสุ กิริยาพฺยากตํ รูปญฺจ นิพฺพานญฺจ. อิเม ธมฺมา อนิยฺยานิกา.}

\proseitem{1611}{กตเม ธมฺมา นิยตา. จตฺตาโร ทิฏฺฐิคต\-สมฺปยุตฺต\-จิตฺตุปฺปาทา เทฺว โทมนสฺสสหคต\-จิตฺตุปฺปาทา, อิเม ธมฺมา สิยา นิยตา สิยา อนิยตา. จตฺตาโร มคฺคา อปริยาปนฺนา. อิเม ธมฺมา นิยตา.}

\proseitem{1612}{กตเม ธมฺมา อนิยตา. จตฺตาโร ทิฏฺฐิคต\-วิปฺปยุตฺต\-โลภ\-สหคต\-จิตฺตุปฺปาทา วิจิกิจฺฉา\-สหคโต จิตฺตุปฺปาโท อุทฺธจฺจ\-สหคโต จิตฺตุปฺปาโท ตีสุ ภูมีสุ กุสลํ จตูสุ ภูมีสุ วิปาโก ตีสุ ภูมีสุ กิริยาพฺยากตํ รูปญฺจ นิพฺพานญฺจ. อิเม ธมฺมา อนิยตา.}

\proseitem{1613}{กตเม ธมฺมา สอุตฺตรา. ตีสุ ภูมีสุ กุสลํ อกุสลํ ตีสุ ภูมีสุ วิปาโก ตีสุ ภูมีสุ กิริยาพฺยากตํ สพฺพญฺจ รูปํ. อิเม ธมฺมา สอุตฺตรา.}

\proseitem{1614}{กตเม ธมฺมา อนุตฺตรา. จตฺตาโร มคฺคา อปริยาปนฺนา จตฺตาริ จ สามญฺญผลานิ นิพฺพานญฺจ. อิเม ธมฺมา อนุตฺตรา.}

\proseitem{1615}{กตเม ธมฺมา สรณา. ทฺวาทส อกุสล\-จิตฺตุปฺปาทา. อิเม ธมฺมา สรณา.}

\proseitemwithcloser{1616}{กตเม ธมฺมา อรณา. จตูสุ ภูมีสุ กุสลํ จตูสุ ภูมีสุ วิปาโก ตีสุ ภูมีสุ กิริยาพฺยากตํ รูปญฺจ นิพฺพานญฺจ. อิเม ธมฺมา อรณา.}{อตฺถุทฺธาโร นิฏฺฐิโต.}
\nitthitam{ธมฺมสงฺคณี\-ปกรณํ นิฏฺฐิตํ.}
